% Options for packages loaded elsewhere
\PassOptionsToPackage{unicode}{hyperref}
\PassOptionsToPackage{hyphens}{url}
\PassOptionsToPackage{dvipsnames,svgnames,x11names}{xcolor}
\documentclass[
  portuguese,
  11pt,
  openany]{book}
\usepackage{xcolor}
\usepackage[margin=2.8cm]{geometry}
\usepackage{amsmath,amssymb}
\setcounter{secnumdepth}{5}
\usepackage{iftex}
\ifPDFTeX
  \usepackage[T1]{fontenc}
  \usepackage[utf8]{inputenc}
  \usepackage{textcomp} % provide euro and other symbols
\else % if luatex or xetex
  \usepackage{unicode-math} % this also loads fontspec
  \defaultfontfeatures{Scale=MatchLowercase}
  \defaultfontfeatures[\rmfamily]{Ligatures=TeX,Scale=1}
\fi
\usepackage{lmodern}
\ifPDFTeX\else
  % xetex/luatex font selection
\fi
% Use upquote if available, for straight quotes in verbatim environments
\IfFileExists{upquote.sty}{\usepackage{upquote}}{}
\IfFileExists{microtype.sty}{% use microtype if available
  \usepackage[]{microtype}
  \UseMicrotypeSet[protrusion]{basicmath} % disable protrusion for tt fonts
}{}
\usepackage{setspace}
\makeatletter
\@ifundefined{KOMAClassName}{% if non-KOMA class
  \IfFileExists{parskip.sty}{%
    \usepackage{parskip}
  }{% else
    \setlength{\parindent}{0pt}
    \setlength{\parskip}{6pt plus 2pt minus 1pt}}
}{% if KOMA class
  \KOMAoptions{parskip=half}}
\makeatother
\usepackage{color}
\usepackage{fancyvrb}
\newcommand{\VerbBar}{|}
\newcommand{\VERB}{\Verb[commandchars=\\\{\}]}
\DefineVerbatimEnvironment{Highlighting}{Verbatim}{commandchars=\\\{\}}
% Add ',fontsize=\small' for more characters per line
\usepackage{framed}
\definecolor{shadecolor}{RGB}{248,248,248}
\newenvironment{Shaded}{\begin{snugshade}}{\end{snugshade}}
\newcommand{\AlertTok}[1]{\textcolor[rgb]{0.94,0.16,0.16}{#1}}
\newcommand{\AnnotationTok}[1]{\textcolor[rgb]{0.56,0.35,0.01}{\textbf{\textit{#1}}}}
\newcommand{\AttributeTok}[1]{\textcolor[rgb]{0.13,0.29,0.53}{#1}}
\newcommand{\BaseNTok}[1]{\textcolor[rgb]{0.00,0.00,0.81}{#1}}
\newcommand{\BuiltInTok}[1]{#1}
\newcommand{\CharTok}[1]{\textcolor[rgb]{0.31,0.60,0.02}{#1}}
\newcommand{\CommentTok}[1]{\textcolor[rgb]{0.56,0.35,0.01}{\textit{#1}}}
\newcommand{\CommentVarTok}[1]{\textcolor[rgb]{0.56,0.35,0.01}{\textbf{\textit{#1}}}}
\newcommand{\ConstantTok}[1]{\textcolor[rgb]{0.56,0.35,0.01}{#1}}
\newcommand{\ControlFlowTok}[1]{\textcolor[rgb]{0.13,0.29,0.53}{\textbf{#1}}}
\newcommand{\DataTypeTok}[1]{\textcolor[rgb]{0.13,0.29,0.53}{#1}}
\newcommand{\DecValTok}[1]{\textcolor[rgb]{0.00,0.00,0.81}{#1}}
\newcommand{\DocumentationTok}[1]{\textcolor[rgb]{0.56,0.35,0.01}{\textbf{\textit{#1}}}}
\newcommand{\ErrorTok}[1]{\textcolor[rgb]{0.64,0.00,0.00}{\textbf{#1}}}
\newcommand{\ExtensionTok}[1]{#1}
\newcommand{\FloatTok}[1]{\textcolor[rgb]{0.00,0.00,0.81}{#1}}
\newcommand{\FunctionTok}[1]{\textcolor[rgb]{0.13,0.29,0.53}{\textbf{#1}}}
\newcommand{\ImportTok}[1]{#1}
\newcommand{\InformationTok}[1]{\textcolor[rgb]{0.56,0.35,0.01}{\textbf{\textit{#1}}}}
\newcommand{\KeywordTok}[1]{\textcolor[rgb]{0.13,0.29,0.53}{\textbf{#1}}}
\newcommand{\NormalTok}[1]{#1}
\newcommand{\OperatorTok}[1]{\textcolor[rgb]{0.81,0.36,0.00}{\textbf{#1}}}
\newcommand{\OtherTok}[1]{\textcolor[rgb]{0.56,0.35,0.01}{#1}}
\newcommand{\PreprocessorTok}[1]{\textcolor[rgb]{0.56,0.35,0.01}{\textit{#1}}}
\newcommand{\RegionMarkerTok}[1]{#1}
\newcommand{\SpecialCharTok}[1]{\textcolor[rgb]{0.81,0.36,0.00}{\textbf{#1}}}
\newcommand{\SpecialStringTok}[1]{\textcolor[rgb]{0.31,0.60,0.02}{#1}}
\newcommand{\StringTok}[1]{\textcolor[rgb]{0.31,0.60,0.02}{#1}}
\newcommand{\VariableTok}[1]{\textcolor[rgb]{0.00,0.00,0.00}{#1}}
\newcommand{\VerbatimStringTok}[1]{\textcolor[rgb]{0.31,0.60,0.02}{#1}}
\newcommand{\WarningTok}[1]{\textcolor[rgb]{0.56,0.35,0.01}{\textbf{\textit{#1}}}}
\usepackage{longtable,booktabs,array}
\newcounter{none} % for unnumbered tables
\usepackage{calc} % for calculating minipage widths
% Correct order of tables after \paragraph or \subparagraph
\usepackage{etoolbox}
\makeatletter
\patchcmd\longtable{\par}{\if@noskipsec\mbox{}\fi\par}{}{}
\makeatother
% Allow footnotes in longtable head/foot
\IfFileExists{footnotehyper.sty}{\usepackage{footnotehyper}}{\usepackage{footnote}}
\makesavenoteenv{longtable}
\usepackage{graphicx}
\makeatletter
\newsavebox\pandoc@box
\newcommand*\pandocbounded[1]{% scales image to fit in text height/width
  \sbox\pandoc@box{#1}%
  \Gscale@div\@tempa{\textheight}{\dimexpr\ht\pandoc@box+\dp\pandoc@box\relax}%
  \Gscale@div\@tempb{\linewidth}{\wd\pandoc@box}%
  \ifdim\@tempb\p@<\@tempa\p@\let\@tempa\@tempb\fi% select the smaller of both
  \ifdim\@tempa\p@<\p@\scalebox{\@tempa}{\usebox\pandoc@box}%
  \else\usebox{\pandoc@box}%
  \fi%
}
% Set default figure placement to htbp
\def\fps@figure{htbp}
\makeatother
\ifLuaTeX
\usepackage[bidi=basic,shorthands=off]{babel}
\else
\usepackage[bidi=default,shorthands=off]{babel}
\fi
\ifLuaTeX
  \usepackage{selnolig} % disable illegal ligatures
\fi
\setlength{\emergencystretch}{3em} % prevent overfull lines
\providecommand{\tightlist}{%
  \setlength{\itemsep}{0pt}\setlength{\parskip}{0pt}}
\usepackage{bookmark}
\IfFileExists{xurl.sty}{\usepackage{xurl}}{} % add URL line breaks if available
\urlstyle{same}
\hypersetup{
  pdftitle={Introdução ao R e à Simulação Estatística},
  pdfauthor={Renato de Paula},
  pdflang={pt-PT},
  colorlinks=true,
  linkcolor={Maroon},
  filecolor={Maroon},
  citecolor={Blue},
  urlcolor={Blue},
  pdfcreator={LaTeX via pandoc}}

\title{Introdução ao R e à Simulação Estatística}
\usepackage{etoolbox}
\makeatletter
\providecommand{\subtitle}[1]{% add subtitle to \maketitle
  \apptocmd{\@title}{\par {\large #1 \par}}{}{}
}
\makeatother
\subtitle{Um Laboratório de Estatística Computacional}
\author{Renato de Paula}
\date{06 de July de 2026}

\begin{document}
\maketitle

{
\setcounter{tocdepth}{1}
\tableofcontents
}
\setstretch{1.15}
\chapter*{Prefácio}\label{prefuxe1cio}


Este livro nasceu de uma convicção simples: aprende-se estatística a fazer
estatística. Durante muito tempo, o ensino da disciplina viveu dividido entre,
por um lado, a teoria --- distribuições, teoremas, demonstrações --- e, por outro,
a prática, muitas vezes reduzida a exercícios de papel e lápis com números
pequenos e cuidadosamente escolhidos para ``dar certo''. O computador estava
ausente, ou aparecia apenas no fim, como uma calculadora sofisticada.

Hoje sabemos que essa separação é artificial. A estatística moderna é, em larga
medida, uma ciência computacional. Quando queremos perceber como se comporta a
média de uma amostra, não precisamos de recitar o Teorema do Limite Central de
memória: podemos gerar milhares de amostras no computador e \emph{ver} a distribuição
a formar-se diante dos nossos olhos. Quando um integral não tem solução
fechada, podemos estimá-lo por simulação. Quando um modelo é demasiado complexo
para a análise tradicional, simulamos. A capacidade de programar deixou de ser
um extra para quem estuda estatística --- passou a ser uma competência central.

\textbf{O que este livro se propõe fazer.} Levar o leitor, sem pressupor qualquer
experiência prévia de programação, desde os primeiros comandos em R até à
utilização de métodos de simulação para resolver problemas de probabilidade e
estatística. Fazemo-lo mantendo sempre juntas as duas metades da disciplina: a
teoria que dá sentido aos métodos e a aplicação que os torna úteis.

\section{Porquê o R?}\label{porquuxea-o-r}

Existem muitas linguagens de programação. Porquê estudar precisamente o R?

O R foi concebido, desde a raiz, \emph{por} estatísticos e \emph{para} fazer estatística.
Não é uma linguagem de propósito geral à qual se acrescentaram, mais tarde,
ferramentas estatísticas; é o contrário. Trabalhar com vetores, matrizes,
tabelas de dados, distribuições de probabilidade e gráficos é, em R, natural e
imediato --- porque foi para isso que a linguagem foi desenhada. Um cálculo que
noutras linguagens exigiria vários passos, em R resume-se muitas vezes a uma
única linha legível.

A isto junta-se uma comunidade científica enorme e ativa. O R é gratuito, de
código aberto, e conta com milhares de pacotes que estendem as suas capacidades
a praticamente todas as áreas da ciência de dados. É a lingua franca da
estatística académica e é largamente usado na indústria, na investigação
biomédica, nas ciências sociais, na economia e na finança. Aprender R não é
aprender uma ferramenta de nicho: é adquirir uma competência transferível e
valorizada.

\section{Porquê a simulação?}\label{porquuxea-a-simulauxe7uxe3o}

A simulação é o fio condutor deste livro. A ideia é poderosa e, ao mesmo tempo,
desarmante de tão simples: em vez de calcular uma probabilidade ou um valor
esperado por métodos analíticos, deixamos o computador \emph{repetir a experiência}
muitas vezes e observamos o que acontece.

Suponha que quer saber a probabilidade de, ao lançar dez moedas, obter exatamente
cinco caras. Pode deduzir a fórmula binomial --- e deve aprender a fazê-lo. Mas
pode também pedir ao R que ``lance'' dez moedas cem mil vezes e conte a proporção
de vezes em que saíram cinco caras. Os dois valores hão de coincidir, e essa
coincidência não é um acaso: é a \textbf{Lei dos Grandes Números} em ação. A simulação
transforma teoremas abstratos em fenómenos que podemos ver, medir e experimentar.

Ao longo do livro, a simulação servir-nos-á para três propósitos: \textbf{compreender}
(tornando visíveis conceitos como a variabilidade amostral ou a convergência),
\textbf{estimar} (aproximando probabilidades, integrais e valores esperados que não
sabemos calcular de outra forma) e \textbf{validar} (confrontando resultados teóricos
com evidência gerada por computador).

\section{Para quem é este livro}\label{para-quem-uxe9-este-livro}

Escrevemos este livro a pensar em estudantes do primeiro ano de licenciaturas em
que a estatística tem um papel central --- matemática aplicada, estatística,
ciência de dados, economia, engenharia, biologia --- e, mais geralmente, em
qualquer pessoa que queira aprender a programar em R com o objetivo de fazer
estatística.

\textbf{Não é preciso saber programar.} Começamos do princípio, com a instalação do
software e os primeiros cálculos, e avançamos gradualmente. Pressupomos apenas os
conhecimentos de matemática do ensino secundário e uma curiosidade genuína.

Do lado da estatística e da probabilidade, o livro é razoavelmente
autossuficiente: os conceitos necessários são introduzidos à medida que fazem
falta. Ainda assim, os leitores que estejam a estudar simultaneamente uma
disciplina de probabilidades encontrarão aqui o complemento computacional
ideal --- o lugar onde a teoria ganha corpo.

\section{Como está organizado}\label{como-estuxe1-organizado}

O livro segue uma progressão deliberada, em que cada capítulo prepara o seguinte.

Na \textbf{primeira parte} aprendemos a linguagem: instalamos o R e o RStudio, usamos
o R como calculadora, conhecemos as estruturas de dados fundamentais (vetores,
matrizes, tabelas e listas) e as ferramentas que dão vida à programação ---
estruturas de decisão, funções, ciclos e leitura de dados.

Na \textbf{segunda parte} viramo-nos para a visualização e a manipulação de dados,
aprendendo a produzir gráficos informativos e a transformar conjuntos de dados
reais.

Na \textbf{terceira parte} chegamos ao coração do livro: probabilidade, geração de
números aleatórios e métodos de simulação --- o método da transformada inversa, o
método da aceitação-rejeição, a ilustração dos teoremas limite e o cálculo de
integrais por simulação (Monte Carlo).

Cada capítulo alterna explicação, código comentado e exemplos aplicados, e
termina com uma coleção de exercícios. As soluções encontram-se no final do
livro.

\section{Como tirar o máximo partido}\label{como-tirar-o-muxe1ximo-partido}

\textbf{Programe enquanto lê.} Não há forma de aprender a programar apenas lendo
sobre programação, tal como não se aprende a nadar a ver vídeos de natação.
Tenha o RStudio aberto ao lado do livro, escreva o código à medida que avança,
experimente variações e não tenha medo dos erros --- em programação, um erro é
apenas informação sobre o que ajustar.

As caixas coloridas ao longo do texto assinalam informação de tipos diferentes:

\begin{itemize}
\tightlist
\item
  As caixas de \textbf{nota} contêm observações e esclarecimentos úteis.
\item
  As caixas de \textbf{dica} oferecem sugestões práticas e boas práticas.
\item
  As caixas de \textbf{atenção} avisam para erros comuns e armadilhas frequentes.
\end{itemize}

O código aparece em caixas próprias. Quando uma linha começa por \texttt{\#\#}, trata-se
do \textbf{resultado} que o R devolve --- foi assim que combinámos representar aquilo
que apareceria no ecrã depois de executar o comando anterior.

\section{Uma nota final}\label{uma-nota-final}

Um livro assim não substitui a prática paciente nem a discussão com colegas e
docentes; pretende ser um companheiro fiável nessa aprendizagem. Se, ao chegar
ao fim, o leitor olhar para um problema estatístico e a sua primeira reação for
``posso simular isto para ver o que acontece'', então este livro terá cumprido o
seu propósito.

Bom trabalho --- e bem-vindo ao laboratório.

\chapter{Estatística, computação e simulação}\label{intro}

Antes de escrevermos a primeira linha de código, vale a pena parar um momento
para responder a uma pergunta legítima: \emph{porque é que um estudante de
estatística precisa de aprender a programar?} A resposta a esta pergunta é, em
boa medida, a justificação para todo este livro --- e perceber bem a resposta
tornará tudo o que se segue mais motivado e mais fácil.

\section{A estatística tornou-se computacional}\label{a-estatuxedstica-tornou-se-computacional}

Durante grande parte do século XX, a estatística viveu com um travão
permanente: a dificuldade de fazer contas. Os métodos que se ensinavam e usavam
eram, muitas vezes, aqueles que se conseguiam calcular à mão ou com uma
calculadora --- tabelas de distribuições impressas no fim dos manuais,
aproximações engenhosas, fórmulas fechadas cuidadosamente deduzidas para evitar
cálculos impraticáveis.

O computador mudou tudo. Aquilo que antes era impossível de calcular passou a
ser uma questão de segundos. E, com isso, mudou também aquilo que faz sentido
ensinar e aprender. Já não precisamos de decorar valores de tabelas nem de nos
limitarmos aos poucos modelos que admitem solução analítica simples. Podemos
explorar dados reais, ajustar modelos complexos e, sobretudo, \textbf{simular}.

Não estamos a dizer que a teoria deixou de importar --- pelo contrário. Continuamos
a precisar de compreender \emph{porque é que} os métodos funcionam. O que mudou é que
a teoria e o computador passaram a trabalhar em conjunto: a teoria dá-nos a
compreensão, o computador dá-nos o alcance.

\section{O que é o R}\label{o-que-uxe9-o-r}

O \textbf{R} é uma linguagem de programação e um ambiente de trabalho pensados
especificamente para a computação estatística e para a produção de gráficos.
Voltaremos a esta descrição, com mais detalhe, no próximo capítulo; para já,
basta reter três ideias.

Primeiro, o R é \textbf{gratuito e de código aberto}. Qualquer pessoa o pode
descarregar, usar, inspecionar e modificar --- o que faz dele uma ferramenta
democrática, sem barreiras de custo entre quem quer aprender e o conhecimento.

Segundo, o R foi feito \textbf{por estatísticos, para estatísticos}. Manipular
conjuntos de dados, trabalhar com distribuições de probabilidade, ajustar
modelos e desenhar gráficos são tarefas para as quais a linguagem está
naturalmente vocacionada.

Terceiro, o R tem por trás uma \textbf{comunidade científica vastíssima}. Milhares de
pacotes, escritos e partilhados livremente por investigadores de todo o mundo,
estendem as suas capacidades a praticamente qualquer domínio.

\section{A ideia central: simular}\label{a-ideia-central-simular}

A palavra ``simulação'' percorre este livro de uma ponta à outra. Vale a pena
introduzi-la desde já, com um exemplo minúsculo, para que o leitor perceba de
imediato o poder da ideia.

Imagine que lançamos um dado equilibrado e queremos saber qual é o valor médio
que sai a longo prazo. A teoria responde sem hesitar: a média é
\((1+2+3+4+5+6)/6 = 3{,}5\). Mas há uma segunda maneira de chegar à mesma resposta,
sem fazer a conta: podemos \emph{lançar o dado muitas vezes} e calcular a média dos
resultados. Em R, isso escreve-se assim:

\begin{Shaded}
\begin{Highlighting}[]
\CommentTok{\# Lançar um dado equilibrado 1 000 000 de vezes}
\NormalTok{lancamentos }\OtherTok{\textless{}{-}} \FunctionTok{sample}\NormalTok{(}\DecValTok{1}\SpecialCharTok{:}\DecValTok{6}\NormalTok{, }\AttributeTok{size =} \DecValTok{1000000}\NormalTok{, }\AttributeTok{replace =} \ConstantTok{TRUE}\NormalTok{)}

\CommentTok{\# Média dos resultados}
\FunctionTok{mean}\NormalTok{(lancamentos)}
\DocumentationTok{\#\# [1] 3.499}
\end{Highlighting}
\end{Shaded}

O valor obtido --- próximo de \(3.5\) --- não é imposto por nenhuma fórmula: \emph{emerge}
da repetição da experiência. Esta é a essência da simulação. Em vez de deduzir o
resultado, construímos uma versão computacional do fenómeno aleatório e
observamos o seu comportamento.

Ainda não é preciso perceber cada detalhe do código acima --- a função \texttt{sample()},
o argumento \texttt{replace}, e assim por diante. Tudo isto será explicado nos próximos
capítulos. O importante, por agora, é ficar com a intuição: \textbf{o computador pode
repetir uma experiência aleatória tantas vezes quantas quisermos, e a partir
dessas repetições estimar aquilo que nos interessa.}

Este pequeno exemplo já contém, em embrião, tudo o que faremos mais tarde de
forma mais ambiciosa. Quando não soubermos calcular uma probabilidade,
estimá-la-emos por simulação. Quando um integral não tiver solução fechada,
aproximá-lo-emos por simulação. Quando quisermos \emph{ver} um teorema a funcionar,
ilustrá-lo-emos por simulação. Mas nada disto será possível sem primeiro
dominarmos a linguagem --- e é para isso que servem os capítulos que se seguem.

\section{O caminho que vamos percorrer}\label{o-caminho-que-vamos-percorrer}

O livro está organizado como uma subida gradual, em que cada degrau assenta no
anterior.

Começamos por \textbf{instalar e conhecer as ferramentas} (Capítulo~\ref{rrstudio})
e por usar o R como uma \textbf{calculadora} que também guarda resultados e toma
decisões (Capítulo~\ref{calculadora}). Aprendemos depois a organizar
informação em \textbf{estruturas de dados} --- vetores, matrizes, tabelas e listas
(Capítulo~\ref{estruturas}) --- e a controlar o fluxo dos programas com
\textbf{estruturas de decisão} (Capítulo~\ref{selecao}) e \textbf{funções}
(Capítulo~\ref{funcoes}).

A partir daí, o livro ganha ritmo: automatizamos tarefas repetitivas com ciclos,
lemos dados a partir de ficheiros, visualizamos e manipulamos conjuntos de dados
reais e, finalmente, entramos no território da \textbf{probabilidade e da simulação},
onde tudo o que aprendemos converge num objetivo comum.

Se está a ler este livro no âmbito de uma disciplina, use os capítulos como um
guia de laboratório: leia com o RStudio aberto, execute cada exemplo e resolva
os exercícios no fim de cada capítulo antes de avançar. A programação aprende-se
com as mãos, não apenas com os olhos.

Comecemos, então, pelo princípio: pôr as ferramentas a funcionar.

\chapter{O R e o RStudio}\label{rrstudio}

Toda a aprendizagem começa por preparar as ferramentas. Neste capítulo damos os
primeiros passos práticos: percebemos o que é o R, instalamo-lo no computador,
instalamos também o RStudio --- o ambiente que torna o trabalho com R muito mais
confortável --- e fazemos uma primeira visita guiada ao seu ecrã. No fim, o leitor
terá um laboratório de estatística a funcionar e estará pronto para escrever os
primeiros comandos.

\section{O que é o R}\label{o-que-uxe9-o-r-1}

O R é um software de código aberto, desenvolvido como uma implementação gratuita
da linguagem \textbf{S}, concebida especificamente para computação estatística,
programação e produção de gráficos. A intenção original era dar a quem trabalha
com dados a capacidade de os explorar de forma intuitiva e interativa,
recorrendo a representações gráficas que facilitassem a sua compreensão. O R foi
criado por Ross Ihaka e Robert Gentleman, da Universidade de Auckland, na Nova
Zelândia --- e o nome, com o seu toque de humor, brinca com as iniciais dos dois
autores e com a linguagem S que lhe deu origem.

Mais importante do que a história é perceber o que o R oferece a quem faz
estatística. Em termos práticos, o R reúne num só ambiente:

\begin{itemize}
\tightlist
\item
  manipulação eficiente de dados e formas flexíveis de os armazenar;
\item
  operadores poderosos para cálculos com vetores e matrizes;
\item
  uma coleção abrangente e coerente de ferramentas para análise de dados;
\item
  capacidades gráficas avançadas, tanto para visualização no ecrã como para
  produção de figuras de qualidade de publicação;
\item
  uma linguagem de programação completa, com estruturas de decisão, ciclos,
  funções definidas pelo utilizador e ferramentas de entrada e saída de dados.
\end{itemize}

Repare que a última destas capacidades --- o R ser uma \textbf{linguagem de programação}
completa --- é o que distingue o R de uma folha de cálculo. Numa folha de cálculo,
descrevemos \emph{onde} estão os dados; em R, descrevemos \emph{o que queremos fazer} com
eles. É essa mudança de perspetiva que dá ao R a sua flexibilidade, e é o que
vamos aprender ao longo do livro.

\section{Instalação}\label{instalauxe7uxe3o}

Trabalhar com R implica, na prática, instalar \textbf{dois} programas
complementares. O primeiro é o R propriamente dito --- o motor que executa os
cálculos. O segundo é o RStudio, uma interface que torna a escrita e a
organização do código muito mais agradáveis.

\subsection{Instalar o R}\label{instalar-o-r}

A versão base do R, que inclui o conjunto fundamental de funções, descarrega-se
do sítio oficial do projeto:

\begin{quote}
\url{https://www.r-project.org/}
\end{quote}

Basta seguir a ligação para o CRAN (a rede de repositórios do R), escolher o
sistema operativo (Windows, macOS ou Linux) e instalar como qualquer outro
programa.

\subsection{Instalar o RStudio}\label{instalar-o-rstudio}

Poderíamos trabalhar apenas com o R base, mas escrever código dessa forma é
pouco confortável. Por isso, instala-se quase sempre um \textbf{ambiente de
desenvolvimento integrado} (IDE) --- um programa que ajuda a escrever, guardar e
organizar os \emph{scripts}, a executar comandos e a gerir os resultados. A escolha
habitual é o \textbf{RStudio}, gratuito e disponível em:

\begin{quote}
\url{https://posit.co/download/rstudio-desktop/}
\end{quote}

A ordem importa: instale \textbf{primeiro} o R e \textbf{só depois} o RStudio. O RStudio
não substitui o R --- precisa de o encontrar já instalado no computador para
funcionar. Se instalar o RStudio sem ter o R, ele não terá motor para executar
os comandos.

\subsection{Pacotes: estender o R}\label{pacotes-estender-o-r}

Além do conjunto base de funções, o R oferece milhares de \textbf{pacotes} adicionais,
desenvolvidos e partilhados pela comunidade. Um pacote é uma coleção de funções e
dados que acrescenta novas capacidades ao R --- desde gráficos elegantes até
modelos estatísticos sofisticados.

Instalar um pacote faz-se uma única vez, com a função \texttt{install.packages()}
(é preciso estar ligado à Internet):

\begin{Shaded}
\begin{Highlighting}[]
\FunctionTok{install.packages}\NormalTok{(}\StringTok{"ggplot2"}\NormalTok{)}
\end{Highlighting}
\end{Shaded}

Para ver todos os pacotes já instalados no seu sistema:

\begin{Shaded}
\begin{Highlighting}[]
\FunctionTok{installed.packages}\NormalTok{()}
\end{Highlighting}
\end{Shaded}

Distinga bem dois passos que são frequentemente confundidos. \textbf{Instalar} um
pacote (com \texttt{install.packages()}) faz-se uma vez só --- descarrega o pacote para o
computador. \textbf{Carregar} um pacote (com \texttt{library()}) faz-se em cada sessão em que
o queremos usar --- torna as suas funções disponíveis. É como comprar um livro
(uma vez) e depois tirá-lo da estante sempre que precisamos dele.

\begin{Shaded}
\begin{Highlighting}[]
\FunctionTok{library}\NormalTok{(ggplot2)   }\CommentTok{\# torna as funções do pacote disponíveis nesta sessão}
\end{Highlighting}
\end{Shaded}

\section{Uma visita guiada ao RStudio}\label{uma-visita-guiada-ao-rstudio}

Ao abrir o RStudio pela primeira vez, o ecrã pode parecer intimidante, com
vários painéis ao mesmo tempo. Na verdade, são quatro áreas, cada uma com uma
função clara. Vale a pena conhecê-las, porque vamos passar muitas horas neste
ambiente.

\textbf{O Editor de \emph{scripts}} (por omissão, no canto superior esquerdo) é onde
escrevemos e guardamos o nosso código. Um \emph{script} é um ficheiro de texto com uma
sequência de comandos --- o ``guião'' da nossa análise. O editor oferece realce de
sintaxe (que colore o código para o tornar legível), completamento automático e a
possibilidade de executar linhas ou blocos de código diretamente. É aqui que
deve escrever o código que quer conservar.

\textbf{A Consola} (canto inferior esquerdo) é onde os comandos são efetivamente
executados e onde aparecem os resultados. Podemos escrever comandos diretamente
na consola para testes rápidos, mas, ao contrário do editor, o que escrevemos na
consola não fica guardado. Pense na consola como uma folha de rascunho e no
editor como o caderno definitivo.

\textbf{O painel \emph{Environment}/\emph{History}} (canto superior direito) mostra, no
separador \emph{Environment}, todos os objetos atualmente na memória da sessão --- os
dados e variáveis que fomos criando. O separador \emph{History} guarda o histórico de
comandos executados, permitindo recuperá-los facilmente.

\textbf{O painel inferior direito} é multifuncional e reúne vários separadores úteis:
\emph{Files} (os ficheiros do diretório de trabalho), \emph{Plots} (os gráficos que
produzimos), \emph{Packages} (os pacotes instalados e carregados), \emph{Help} (o sistema
de ajuda) e \emph{Viewer} (para visualizar documentos HTML e outros conteúdos).

O separador \textbf{Help} será um dos seus melhores amigos. Todo o R está documentado,
e aprender a ler essa documentação é uma competência que compensa cultivar desde
cedo. Veremos, já no próximo capítulo, como pedir ajuda sobre qualquer função.

\section{Atalhos que poupam tempo}\label{atalhos-que-poupam-tempo}

Programar é uma atividade repetitiva, e os bons atalhos de teclado fazem uma
diferença real ao fim de algumas horas. Estes são os mais úteis para começar:

{\def\LTcaptype{none} % do not increment counter
\begin{longtable}[]{@{}
  >{\raggedright\arraybackslash}p{(\linewidth - 2\tabcolsep) * \real{0.5000}}
  >{\raggedright\arraybackslash}p{(\linewidth - 2\tabcolsep) * \real{0.5000}}@{}}
\toprule\noalign{}
\begin{minipage}[b]{\linewidth}\raggedright
Atalho
\end{minipage} & \begin{minipage}[b]{\linewidth}\raggedright
O que faz
\end{minipage} \\
\midrule\noalign{}
\endhead
\bottomrule\noalign{}
\endlastfoot
\texttt{CTRL} + \texttt{ENTER} & Executa a(s) linha(s) selecionada(s) no \emph{script} \\
\texttt{ALT} + \texttt{-} & Insere o operador de atribuição \texttt{\textless{}-} \\
\texttt{CTRL} + \texttt{SHIFT} + \texttt{M} & Insere o operador \emph{pipe} (\texttt{\%\textgreater{}\%} ou \texttt{\textbar{}\textgreater{}}) \\
\texttt{CTRL} + \texttt{1} & Move o cursor para o painel do \emph{script} \\
\texttt{CTRL} + \texttt{2} & Move o cursor para a consola \\
\texttt{CTRL} + \texttt{ALT} + \texttt{I} & Insere um novo bloco de código num documento R Markdown \\
\texttt{CTRL} + \texttt{SHIFT} + \texttt{K} & Compila um documento R Markdown \\
\texttt{ALT} + \texttt{SHIFT} + \texttt{K} & Abre a lista completa de atalhos \\
\end{longtable}
}

Num Mac, os atalhos são geralmente equivalentes aos usados em Windows ou
Linux, mas a tecla \textbf{CTRL} é normalmente substituída por \textbf{Command} e a
tecla \textbf{ALT} por \textbf{Option}.

Com as ferramentas instaladas e o ecrã já familiar, estamos prontos para o que
realmente interessa: começar a dar ordens ao R. É o que faremos no próximo
capítulo, onde o usaremos primeiro como uma calculadora e, depois, como algo bem
mais poderoso.

\chapter{O R como calculadora}\label{calculadora}

A melhor forma de começar a programar em R é usá-lo para aquilo que ele faz de
forma mais imediata: contas. Neste capítulo tratamos o R como uma calculadora
muito poderosa --- mas, à medida que avançamos, vamos descobrir que é bem mais do
que isso. Uma calculadora dá-nos um resultado e esquece-o; o R permite-nos
\emph{guardar} valores, \emph{dar-lhes nomes} e \emph{reutilizá-los}, e é esse pequeno salto ---
de calcular para nomear e reutilizar --- que abre as portas à programação.

Há ainda uma razão de fundo para começar por aqui. A estatística está
intimamente ligada à álgebra: representamos conjuntos de dados e variáveis
através de vetores e matrizes, e manipulamo-los com operações aritméticas. Antes
de avançarmos para comandos estatísticos mais específicos, convém perceber bem
como o R lida com números, valores e operações elementares.

\section{A consola e o primeiro comando}\label{a-consola-e-o-primeiro-comando}

O R aceita comandos numa interface de linha de comando, assinalada pelo símbolo
\texttt{\textgreater{}}, a que chamamos \textbf{prompt}. Escrevemos um comando a seguir ao \emph{prompt},
pressionamos \texttt{Enter}, e o R interpreta-o, executa-o e mostra o resultado.

\begin{Shaded}
\begin{Highlighting}[]
\FunctionTok{print}\NormalTok{(}\StringTok{"O meu primeiro comando em R!"}\NormalTok{)}
\DocumentationTok{\#\# [1] "O meu primeiro comando em R!"}
\end{Highlighting}
\end{Shaded}

Ao longo do livro, o código que escrevemos aparece em caixas próprias. Quando uma
linha começa por \texttt{\#\#}, trata-se do \textbf{resultado} devolvido pelo R --- aquilo que
apareceria no ecrã depois de executar o comando anterior. Assim, pode ler cada
bloco como um diálogo: primeiro o que pedimos, depois o que o R respondeu.

O \texttt{{[}1{]}} que precede o resultado indica a posição do primeiro elemento mostrado
nessa linha. Para um único valor parece supérfluo, mas quando o R devolve muitos
valores, esta numeração ajuda-nos a localizá-los.

\section{Comentários: escrever para humanos}\label{comentuxe1rios-escrever-para-humanos}

O carácter \texttt{\#} inicia um \textbf{comentário}: tudo o que vem depois dele, na mesma
linha, é ignorado pelo R. Os comentários não servem para o computador --- servem
para nós e para quem, mais tarde, ler o nosso código (muitas vezes, esse alguém
somos nós próprios, uns meses depois).

\begin{Shaded}
\begin{Highlighting}[]
\CommentTok{\# Este é um comentário: o R ignora{-}o}
\DecValTok{2} \SpecialCharTok{+} \DecValTok{2}   \CommentTok{\# o R calcula 2 mais 2}
\DocumentationTok{\#\# [1] 4}
\end{Highlighting}
\end{Shaded}

Comente o \emph{porquê}, não o \emph{quê}. Uma linha como \texttt{x\ \textless{}-\ x\ +\ 1\ \ \#\ incrementa\ x} é
inútil, porque o comentário apenas repete o código. Um bom comentário explica a
intenção: \texttt{x\ \textless{}-\ x\ +\ 1\ \ \#\ passamos\ ao\ dia\ seguinte}.

\section{Pedir ajuda}\label{pedir-ajuda}

Ninguém decora todas as funções do R --- e não precisa. O que é essencial é saber
consultar a documentação. Se conhece o nome de uma função, precede-o de um ponto
de interrogação para abrir a respetiva página de ajuda:

\begin{Shaded}
\begin{Highlighting}[]
\NormalTok{?sum}
\end{Highlighting}
\end{Shaded}

A ajuda descreve o que a função faz, quais os seus argumentos e o que devolve.
Para além disso, muitas funções trazem exemplos prontos a correr, que se veem
com \texttt{example()}:

\begin{Shaded}
\begin{Highlighting}[]
\FunctionTok{example}\NormalTok{(sum)}
\end{Highlighting}
\end{Shaded}

Aprender a ler estas páginas é uma das competências mais rentáveis para quem
programa em R.

\section{Objetos e variáveis}\label{objetos-e-variuxe1veis}

Chegamos ao salto de que falámos na abertura. Em R, tudo o que criamos ou
carregamos numa sessão --- um número, um texto, um vetor, uma tabela, até uma
função --- é um \textbf{objeto}, uma unidade de armazenamento na memória.

Uma \textbf{variável} é simplesmente um \emph{nome} que associamos a um objeto. Ao criar
uma variável, criamos um ``rótulo'' que aponta para o objeto guardado na memória e
que nos permite voltar a ele sempre que quisermos, sem ter de o recalcular ou
reescrever.

\subsection{Atribuição}\label{atribuiuxe7uxe3o}

Atribuir um valor a uma variável faz-se com o operador \texttt{\textless{}-} (uma seta que se lê
``recebe''). A expressão seguinte cria uma variável chamada \texttt{x} e guarda-lhe o
valor 10:

\begin{Shaded}
\begin{Highlighting}[]
\NormalTok{x }\OtherTok{\textless{}{-}} \DecValTok{10}
\end{Highlighting}
\end{Shaded}

O operador \texttt{\textless{}-} guarda o valor do lado direito na variável do lado esquerdo. Do
lado esquerdo deve estar apenas um nome de variável. E, uma vez criada, a
variável pode ser reutilizada:

\begin{Shaded}
\begin{Highlighting}[]
\NormalTok{a }\OtherTok{\textless{}{-}} \DecValTok{10}
\NormalTok{b }\OtherTok{\textless{}{-}}\NormalTok{ a }\SpecialCharTok{+} \DecValTok{1}   \CommentTok{\# b recebe o valor de a mais 1}
\NormalTok{b}
\DocumentationTok{\#\# [1] 11}
\end{Highlighting}
\end{Shaded}

Uma atribuição \textbf{não é} uma equação matemática. Escrever \texttt{a\ +\ 2\ \textless{}-\ 10} não faz
sentido: do lado esquerdo tem de estar um nome, não uma expressão a calcular.
Embora o R também aceite \texttt{=} e \texttt{-\textgreater{}} para atribuir, recomendamos o uso de \texttt{\textless{}-},
para não confundir a atribuição (\texttt{\textless{}-}) com o teste de igualdade (\texttt{==}), que
veremos mais adiante.

Podemos guardar não só um número, mas uma sequência de valores. A função \texttt{c()}
(de \emph{combine}, combinar) junta vários valores num \textbf{vetor}:

\begin{Shaded}
\begin{Highlighting}[]
\NormalTok{X }\OtherTok{\textless{}{-}} \FunctionTok{c}\NormalTok{(}\DecValTok{2}\NormalTok{, }\DecValTok{12}\NormalTok{, }\DecValTok{22}\NormalTok{, }\DecValTok{32}\NormalTok{)}
\NormalTok{X}
\DocumentationTok{\#\# [1]  2 12 22 32}
\end{Highlighting}
\end{Shaded}

O R \textbf{distingue maiúsculas de minúsculas}. \texttt{X} e \texttt{x} são dois objetos
diferentes. Este é um dos erros mais comuns de quem começa, por isso convém
gravá-lo desde já.

Para ver todos os objetos que existem atualmente na memória da sessão, usamos
\texttt{ls()} (ou consultamos o separador \emph{Environment} no RStudio):

\begin{Shaded}
\begin{Highlighting}[]
\FunctionTok{ls}\NormalTok{()}
\end{Highlighting}
\end{Shaded}

\subsection{Regras para nomear variáveis}\label{regras-para-nomear-variuxe1veis}

A escolha dos nomes das variáveis não é só uma questão de gosto: bons nomes
tornam o código legível, e o R impõe algumas regras.

Um nome de variável deve começar por uma letra (ou por um ponto seguido de
letra) e pode conter letras, números, pontos e sublinhados. Não pode conter
espaços nem carateres especiais como \texttt{@}, \texttt{\#}, \texttt{\$} ou \texttt{\%}. Assim,
\texttt{nome\_cliente2} é válido; \texttt{2nome} ou \texttt{nome\ cliente} não são.

Também não podemos usar \textbf{palavras reservadas} --- termos com significado especial
no R, como \texttt{if}, \texttt{else}, \texttt{for}, \texttt{while}, \texttt{TRUE}, \texttt{FALSE}, \texttt{function}.

E, como já vimos, o R distingue maiúsculas de minúsculas:

\begin{Shaded}
\begin{Highlighting}[]
\NormalTok{idade }\OtherTok{\textless{}{-}} \DecValTok{20}
\NormalTok{Idade }\OtherTok{\textless{}{-}} \DecValTok{30}
\end{Highlighting}
\end{Shaded}

Aqui, \texttt{idade} e \texttt{Idade} são duas variáveis distintas.

Prefira nomes que descrevam o conteúdo. \texttt{nota\_final} diz-nos muito mais do que
\texttt{x}. Uma convenção comum e legível é usar minúsculas e separar palavras com
sublinhados: \texttt{faculdade\_de\_ciencias}. Código escrito para ser lido é código que
se consegue corrigir e reutilizar.

\section{Tipos de dados}\label{tipos-de-dados}

Nem todos os valores são da mesma natureza. Um número comporta-se de forma
diferente de um texto, e ambos diferem de um valor lógico. Em R, os principais
tipos de dados que vamos encontrar são:

\begin{itemize}
\tightlist
\item
  \textbf{Numeric} --- números, inteiros ou decimais. Exemplos: \texttt{42}, \texttt{3.14}.
\item
  \textbf{Character} --- texto (cadeias de carateres). Exemplo: \texttt{"Olá"}.
\item
  \textbf{Logical} --- valores lógicos, \texttt{TRUE} ou \texttt{FALSE}.
\item
  \textbf{Vectors} --- coleções de elementos do mesmo tipo, como \texttt{c(1,\ 2,\ 3)}.
\item
  \textbf{Factors} --- variáveis categóricas (por exemplo, ``baixo'', ``médio'', ``alto'').
\item
  \textbf{Data frames} --- tabelas com linhas e colunas, ao estilo de uma folha de
  cálculo.
\item
  \textbf{Lists} --- coleções que podem misturar elementos de tipos diferentes.
\end{itemize}

Os quatro primeiros aparecem já neste capítulo; os restantes serão o tema do
Capítulo~\ref{estruturas}. Um primeiro contacto:

\begin{Shaded}
\begin{Highlighting}[]
\NormalTok{a }\OtherTok{\textless{}{-}} \FloatTok{3.14}                \CommentTok{\# numeric}
\NormalTok{b }\OtherTok{\textless{}{-}} \StringTok{"Programação R"}     \CommentTok{\# character}
\NormalTok{c }\OtherTok{\textless{}{-}} \DecValTok{3} \SpecialCharTok{\textless{}} \DecValTok{2}               \CommentTok{\# logical (resultado de uma comparação)}
\NormalTok{d }\OtherTok{\textless{}{-}} \FunctionTok{c}\NormalTok{(}\DecValTok{1}\NormalTok{, }\DecValTok{2}\NormalTok{, }\DecValTok{3}\NormalTok{)          }\CommentTok{\# vector}
\end{Highlighting}
\end{Shaded}

\subsection{Tipagem dinâmica}\label{tipagem-dinuxe2mica}

Uma característica cómoda do R é a \textbf{tipagem dinâmica}: não precisamos de
declarar o tipo de uma variável --- o R deduz-o a partir do valor que lhe
atribuímos. A função \texttt{class()} diz-nos o tipo (a \emph{classe}) de um objeto:

\begin{Shaded}
\begin{Highlighting}[]
\NormalTok{x }\OtherTok{\textless{}{-}} \DecValTok{5}
\FunctionTok{class}\NormalTok{(x)}
\DocumentationTok{\#\# [1] "numeric"}

\NormalTok{y }\OtherTok{\textless{}{-}} \StringTok{"Cinco"}
\FunctionTok{class}\NormalTok{(y)}
\DocumentationTok{\#\# [1] "character"}

\NormalTok{z }\OtherTok{\textless{}{-}} \ConstantTok{TRUE}
\FunctionTok{class}\NormalTok{(z)}
\DocumentationTok{\#\# [1] "logical"}
\end{Highlighting}
\end{Shaded}

Existe ainda a função \texttt{typeof()}, que revela o tipo de armazenamento interno ---
o modo como o R guarda o objeto na memória. Nem sempre coincide com \texttt{class()}, e
essa distinção é por vezes reveladora:

\begin{Shaded}
\begin{Highlighting}[]
\NormalTok{x }\OtherTok{\textless{}{-}} \DecValTok{1}\SpecialCharTok{:}\DecValTok{10}
\FunctionTok{class}\NormalTok{(x)}
\DocumentationTok{\#\# [1] "integer"}
\FunctionTok{typeof}\NormalTok{(x)}
\DocumentationTok{\#\# [1] "integer"}

\NormalTok{y }\OtherTok{\textless{}{-}} \FunctionTok{c}\NormalTok{(}\FloatTok{1.1}\NormalTok{, }\FloatTok{2.2}\NormalTok{, }\FloatTok{3.3}\NormalTok{)}
\FunctionTok{class}\NormalTok{(y)}
\DocumentationTok{\#\# [1] "numeric"}
\FunctionTok{typeof}\NormalTok{(y)}
\DocumentationTok{\#\# [1] "double"}

\NormalTok{z }\OtherTok{\textless{}{-}} \FunctionTok{data.frame}\NormalTok{(}\AttributeTok{a =} \DecValTok{1}\SpecialCharTok{:}\DecValTok{3}\NormalTok{, }\AttributeTok{b =} \FunctionTok{c}\NormalTok{(}\StringTok{"A"}\NormalTok{, }\StringTok{"B"}\NormalTok{, }\StringTok{"C"}\NormalTok{))}
\FunctionTok{class}\NormalTok{(z)}
\DocumentationTok{\#\# [1] "data.frame"}
\FunctionTok{typeof}\NormalTok{(z)}
\DocumentationTok{\#\# [1] "list"}
\end{Highlighting}
\end{Shaded}

Repare no último caso: uma tabela (\texttt{data.frame}) tem classe \texttt{"data.frame"} mas,
internamente, é armazenada como uma \textbf{lista}. Voltaremos a este facto quando
estudarmos as estruturas de dados.

\subsection{Converter entre tipos}\label{converter-entre-tipos}

Muitas vezes, sobretudo quando importamos dados de fontes externas, um valor
chega-nos com o tipo ``errado'' --- um número guardado como texto, por exemplo. As
funções da família \texttt{as.*()} convertem entre tipos:

\begin{Shaded}
\begin{Highlighting}[]
\NormalTok{b }\OtherTok{\textless{}{-}} \FunctionTok{as.character}\NormalTok{(}\DecValTok{15}\NormalTok{)   }\CommentTok{\# de número para texto}
\NormalTok{b}
\DocumentationTok{\#\# [1] "15"}

\NormalTok{y }\OtherTok{\textless{}{-}} \FunctionTok{as.integer}\NormalTok{(}\FloatTok{1.5}\NormalTok{)    }\CommentTok{\# de decimal para inteiro (trunca)}
\NormalTok{y}
\DocumentationTok{\#\# [1] 1}

\NormalTok{w }\OtherTok{\textless{}{-}} \FunctionTok{as.numeric}\NormalTok{(}\StringTok{"10"}\NormalTok{)   }\CommentTok{\# de texto para número}
\NormalTok{w}
\DocumentationTok{\#\# [1] 10}
\end{Highlighting}
\end{Shaded}

\section{Operadores aritméticos}\label{operadores-aritmuxe9ticos}

Com valores para manipular, vejamos as operações que os combinam. Os operadores
aritméticos do R são os que seria de esperar, com dois acréscimos úteis: a
divisão inteira e o resto.

{\def\LTcaptype{none} % do not increment counter
\begin{longtable}[]{@{}lll@{}}
\toprule\noalign{}
Operador & Descrição & Exemplo \\
\midrule\noalign{}
\endhead
\bottomrule\noalign{}
\endlastfoot
\texttt{+} & Adição & \texttt{5\ +\ 2} dá \texttt{7} \\
\texttt{-} & Subtração & \texttt{5\ -\ 2} dá \texttt{3} \\
\texttt{*} & Multiplicação & \texttt{5\ *\ 2} dá \texttt{10} \\
\texttt{/} & Divisão & \texttt{5\ /\ 2} dá \texttt{2.5} \\
\texttt{\%/\%} & Divisão inteira & \texttt{5\ \%/\%\ 2} dá \texttt{2} \\
\texttt{\%\%} & Resto da divisão & \texttt{5\ \%\%\ 2} dá \texttt{1} \\
\texttt{\^{}} & Potência & \texttt{5\ \^{}\ 2} dá \texttt{25} \\
\end{longtable}
}

Alguns exemplos, incluindo funções matemáticas comuns:

\begin{Shaded}
\begin{Highlighting}[]
\DecValTok{5} \SpecialCharTok{*} \DecValTok{21}
\DocumentationTok{\#\# [1] 105}

\FunctionTok{sqrt}\NormalTok{(}\DecValTok{9}\NormalTok{)      }\CommentTok{\# raiz quadrada}
\DocumentationTok{\#\# [1] 3}

\DecValTok{3}\SpecialCharTok{\^{}}\DecValTok{3}
\DocumentationTok{\#\# [1] 27}

\FunctionTok{log}\NormalTok{(}\DecValTok{9}\NormalTok{)       }\CommentTok{\# logaritmo natural}
\DocumentationTok{\#\# [1] 2.197}

\FunctionTok{log10}\NormalTok{(}\DecValTok{9}\NormalTok{)     }\CommentTok{\# logaritmo de base 10}
\DocumentationTok{\#\# [1] 0.9542}

\FunctionTok{exp}\NormalTok{(}\DecValTok{1}\NormalTok{)       }\CommentTok{\# número de Euler, e}
\DocumentationTok{\#\# [1] 2.718}
\end{Highlighting}
\end{Shaded}

Os operadores \texttt{\%\%} e \texttt{\%/\%} são mais úteis do que parecem. O resto \texttt{\%\%} responde a
perguntas como ``este número é par?'' (\texttt{n\ \%\%\ 2\ ==\ 0}) ou ``esta hora, somadas
algumas horas, dá que horas no relógio de 24?'' --- questões que aparecerão nos
exercícios. Vale a pena habituar-se a eles.

\subsection{Prioridade das operações}\label{prioridade-das-operauxe7uxf5es}

Como na matemática, as operações têm uma ordem de prioridade: primeiro potências,
depois multiplicações e divisões, e só então adições e subtrações. Os parênteses
alteram essa ordem e, na dúvida, tornam o código mais claro:

\begin{Shaded}
\begin{Highlighting}[]
\DecValTok{19} \SpecialCharTok{+} \DecValTok{26}\SpecialCharTok{/}\DecValTok{4} \SpecialCharTok{{-}} \DecValTok{2}\SpecialCharTok{*}\DecValTok{10}
\DocumentationTok{\#\# [1] 5.5}

\NormalTok{((}\DecValTok{19} \SpecialCharTok{+} \DecValTok{26}\NormalTok{) }\SpecialCharTok{/}\NormalTok{ (}\DecValTok{4} \SpecialCharTok{{-}} \DecValTok{2}\NormalTok{)) }\SpecialCharTok{*} \DecValTok{10}
\DocumentationTok{\#\# [1] 225}
\end{Highlighting}
\end{Shaded}

\subsection{Controlar os dígitos apresentados}\label{controlar-os-duxedgitos-apresentados}

Por vezes queremos ver mais (ou menos) casas decimais. A opção global \texttt{digits}
controla quantos algarismos significativos o R apresenta:

\begin{Shaded}
\begin{Highlighting}[]
\FunctionTok{exp}\NormalTok{(}\DecValTok{1}\NormalTok{)}
\DocumentationTok{\#\# [1] 2.718282}

\FunctionTok{options}\NormalTok{(}\AttributeTok{digits =} \DecValTok{20}\NormalTok{)}
\FunctionTok{exp}\NormalTok{(}\DecValTok{1}\NormalTok{)}
\DocumentationTok{\#\# [1] 2.718281828459045091}

\FunctionTok{options}\NormalTok{(}\AttributeTok{digits =} \DecValTok{7}\NormalTok{)   }\CommentTok{\# o valor por omissão}
\FunctionTok{exp}\NormalTok{(}\DecValTok{1}\NormalTok{)}
\DocumentationTok{\#\# [1] 2.718282}
\end{Highlighting}
\end{Shaded}

Isto altera apenas a \emph{apresentação}: internamente, o R continua a guardar o valor
com toda a precisão disponível. Mudar \texttt{digits} nunca arredonda os cálculos ---
apenas o que vemos no ecrã.

\section{Objetos predefinidos e valores especiais}\label{objetos-predefinidos-e-valores-especiais}

O R traz consigo constantes matemáticas úteis, como \texttt{pi} (o número \(\pi\)), e
representações para situações que a aritmética habitual não cobre.

\begin{Shaded}
\begin{Highlighting}[]
\NormalTok{pi}
\DocumentationTok{\#\# [1] 3.141593}

\DecValTok{1}\SpecialCharTok{/}\DecValTok{0}
\DocumentationTok{\#\# [1] Inf}

\SpecialCharTok{{-}}\DecValTok{1}\SpecialCharTok{/}\DecValTok{0}
\DocumentationTok{\#\# [1] {-}Inf}

\DecValTok{0}\SpecialCharTok{/}\DecValTok{0}
\DocumentationTok{\#\# [1] NaN}

\FunctionTok{sqrt}\NormalTok{(}\SpecialCharTok{{-}}\DecValTok{1}\NormalTok{)}
\DocumentationTok{\#\# Warning: NaNs produced}
\DocumentationTok{\#\# [1] NaN}
\end{Highlighting}
\end{Shaded}

Estes resultados especiais têm nomes e significados precisos, que importa
distinguir:

\begin{itemize}
\tightlist
\item
  \textbf{\texttt{Inf}} (e \texttt{-Inf}) representam o infinito, resultado de operações como
  dividir por zero.
\item
  \textbf{\texttt{NaN}} (\emph{Not a Number}) representa um resultado numérico \textbf{indefinido},
  como \texttt{0/0} ou a raiz quadrada de um número negativo. É um número que não é
  número.
\item
  \textbf{\texttt{NA}} (\emph{Not Available}) representa um valor \textbf{ausente} --- um dado que não foi
  medido ou que está em falta. Ao contrário de \texttt{NaN}, o \texttt{NA} não é
  necessariamente numérico: pode representar qualquer tipo de dado em falta.
\end{itemize}

A distinção entre \texttt{NaN} e \texttt{NA} é subtil mas importante: o primeiro é fruto de uma
conta impossível; o segundo, de informação que simplesmente não temos.

\subsection{Lidar com valores ausentes}\label{lidar-com-valores-ausentes}

Os valores ausentes são a regra, não a exceção, em dados reais --- questionários
com perguntas em branco, sensores que falharam, registos incompletos. Por isso,
o R trata-os com cuidado: qualquer cálculo que envolva um \texttt{NA} devolve, por
omissão, \texttt{NA}. Faz sentido --- se não sabemos um dos valores, não podemos saber a
média.

\begin{Shaded}
\begin{Highlighting}[]
\FunctionTok{mean}\NormalTok{(}\FunctionTok{c}\NormalTok{(}\DecValTok{1}\NormalTok{, }\DecValTok{2}\NormalTok{, }\DecValTok{3}\NormalTok{, }\ConstantTok{NA}\NormalTok{, }\DecValTok{5}\NormalTok{))}
\DocumentationTok{\#\# [1] NA}
\end{Highlighting}
\end{Shaded}

Quando queremos, ainda assim, calcular ignorando os valores em falta, usamos o
argumento \texttt{na.rm\ =\ TRUE} (\emph{remove NA}), disponível na maioria das funções de
resumo:

\begin{Shaded}
\begin{Highlighting}[]
\FunctionTok{mean}\NormalTok{(}\FunctionTok{c}\NormalTok{(}\DecValTok{1}\NormalTok{, }\DecValTok{2}\NormalTok{, }\DecValTok{3}\NormalTok{, }\ConstantTok{NA}\NormalTok{, }\DecValTok{5}\NormalTok{), }\AttributeTok{na.rm =} \ConstantTok{TRUE}\NormalTok{)}
\DocumentationTok{\#\# [1] 2.75}

\FunctionTok{sum}\NormalTok{(}\FunctionTok{c}\NormalTok{(}\DecValTok{1}\NormalTok{, }\DecValTok{2}\NormalTok{, }\DecValTok{3}\NormalTok{, }\ConstantTok{NA}\NormalTok{, }\DecValTok{5}\NormalTok{), }\AttributeTok{na.rm =} \ConstantTok{TRUE}\NormalTok{)}
\DocumentationTok{\#\# [1] 11}
\end{Highlighting}
\end{Shaded}

Não se pode testar a ausência com o operador de igualdade. \texttt{x\ ==\ NA} \textbf{não}
devolve \texttt{TRUE}/\texttt{FALSE} --- devolve \texttt{NA}, porque comparar algo com um valor
desconhecido é, ele próprio, desconhecido. Para detetar valores em falta, usamos
a função \texttt{is.na()}:

\begin{Shaded}
\begin{Highlighting}[]
\NormalTok{y }\OtherTok{\textless{}{-}} \FunctionTok{c}\NormalTok{(}\DecValTok{1}\NormalTok{, }\DecValTok{2}\NormalTok{, }\ConstantTok{NA}\NormalTok{, }\DecValTok{4}\NormalTok{, }\DecValTok{5}\NormalTok{)}

\NormalTok{y }\SpecialCharTok{==} \ConstantTok{NA}        \CommentTok{\# não funciona como se poderia esperar}
\DocumentationTok{\#\# [1] NA NA NA NA NA}

\FunctionTok{is.na}\NormalTok{(y)       }\CommentTok{\# a forma correta}
\DocumentationTok{\#\# [1] FALSE FALSE  TRUE FALSE FALSE}
\end{Highlighting}
\end{Shaded}

Existe ainda \texttt{is.nan()}, específica para detetar \texttt{NaN}. Note a relação entre as
duas: \texttt{is.na()} deteta tanto \texttt{NA} como \texttt{NaN} (afinal, um \texttt{NaN} também não é um
valor disponível), ao passo que \texttt{is.nan()} só deteta \texttt{NaN}.

\section{Exibir e combinar texto}\label{exibir-e-combinar-texto}

Para além de calcular, precisamos frequentemente de \textbf{comunicar} resultados ---
mostrá-los ao utilizador, combiná-los com texto explicativo, ou pedir dados a
quem executa o programa. Algumas funções essenciais para isso:

\begin{itemize}
\tightlist
\item
  \texttt{print()} --- exibe valores e resultados na consola.
\item
  \texttt{readline()} --- recebe uma entrada escrita pelo utilizador via teclado.
\item
  \texttt{paste()} --- junta (concatena) textos, com um separador (por omissão, um
  espaço).
\item
  \texttt{paste0()} --- como \texttt{paste()}, mas sem separador.
\item
  \texttt{cat()} --- junta e exibe textos e valores de forma direta, sem as aspas de
  \texttt{print()}.
\end{itemize}

Um primeiro exemplo, combinando duas palavras:

\begin{Shaded}
\begin{Highlighting}[]
\NormalTok{nome1 }\OtherTok{\textless{}{-}} \StringTok{"faculdade"}
\NormalTok{nome2 }\OtherTok{\textless{}{-}} \StringTok{"ciências"}
\FunctionTok{print}\NormalTok{(}\FunctionTok{paste}\NormalTok{(nome1, nome2))}
\DocumentationTok{\#\# [1] "faculdade ciências"}
\end{Highlighting}
\end{Shaded}

A função \texttt{readline()} permite escrever programas interativos, que pedem
informação a quem os executa. O que o utilizador escreve chega sempre como
\textbf{texto}, pelo que muitas vezes é preciso convertê-lo com \texttt{as.numeric()} ou
\texttt{as.integer()}:

\begin{Shaded}
\begin{Highlighting}[]
\NormalTok{n }\OtherTok{\textless{}{-}} \FunctionTok{readline}\NormalTok{(}\AttributeTok{prompt =} \StringTok{"Digite um número: "}\NormalTok{)}
\NormalTok{n }\OtherTok{\textless{}{-}} \FunctionTok{as.integer}\NormalTok{(n)   }\CommentTok{\# converte o texto num inteiro}
\FunctionTok{print}\NormalTok{(n }\SpecialCharTok{+} \DecValTok{1}\NormalTok{)}
\end{Highlighting}
\end{Shaded}

Um exemplo mais completo, que usa \texttt{cat()} para uma saudação e trata uma entrada
inválida:

\begin{Shaded}
\begin{Highlighting}[]
\NormalTok{idade }\OtherTok{\textless{}{-}} \FunctionTok{readline}\NormalTok{(}\AttributeTok{prompt =} \StringTok{"Digite a sua idade: "}\NormalTok{)}
\NormalTok{idade }\OtherTok{\textless{}{-}} \FunctionTok{as.numeric}\NormalTok{(idade)}

\ControlFlowTok{if}\NormalTok{ (}\FunctionTok{is.na}\NormalTok{(idade)) \{}
  \FunctionTok{cat}\NormalTok{(}\StringTok{"Entrada inválida. Insira um valor numérico.}\SpecialCharTok{\textbackslash{}n}\StringTok{"}\NormalTok{)}
\NormalTok{\} }\ControlFlowTok{else}\NormalTok{ \{}
  \FunctionTok{cat}\NormalTok{(}\StringTok{"Tem"}\NormalTok{, idade, }\StringTok{"anos.}\SpecialCharTok{\textbackslash{}n}\StringTok{"}\NormalTok{)}
\NormalTok{\}}
\end{Highlighting}
\end{Shaded}

\texttt{readline()} só funciona de forma interativa (na consola). Num documento
compilado ou num \emph{script} executado de uma vez, não há ninguém para responder ao
pedido, pelo que estes exemplos se destinam a ser experimentados diretamente na
consola.

\subsection{\texorpdfstring{A função \texttt{paste()} com mais detalhe}{A função paste() com mais detalhe}}\label{a-funuxe7uxe3o-paste-com-mais-detalhe}

A função \texttt{paste()} é surpreendentemente versátil. Aceita vários argumentos,
combina vetores elemento a elemento e deixa-nos escolher o separador:

\begin{Shaded}
\begin{Highlighting}[]
\CommentTok{\# Vários argumentos}
\FunctionTok{paste}\NormalTok{(}\StringTok{"Data"}\NormalTok{, }\StringTok{"Science"}\NormalTok{, }\StringTok{"com"}\NormalTok{, }\StringTok{"R"}\NormalTok{)}
\DocumentationTok{\#\# [1] "Data Science com R"}

\CommentTok{\# Separador específico}
\FunctionTok{paste}\NormalTok{(}\StringTok{"2024"}\NormalTok{, }\StringTok{"04"}\NormalTok{, }\StringTok{"28"}\NormalTok{, }\AttributeTok{sep =} \StringTok{"{-}"}\NormalTok{)}
\DocumentationTok{\#\# [1] "2024{-}04{-}28"}

\CommentTok{\# Combinação elemento a elemento de dois vetores}
\NormalTok{nomes  }\OtherTok{\textless{}{-}} \FunctionTok{c}\NormalTok{(}\StringTok{"Ana"}\NormalTok{, }\StringTok{"Bruno"}\NormalTok{, }\StringTok{"Carlos"}\NormalTok{)}
\NormalTok{apelidos }\OtherTok{\textless{}{-}} \FunctionTok{c}\NormalTok{(}\StringTok{"Silva"}\NormalTok{, }\StringTok{"Oliveira"}\NormalTok{, }\StringTok{"Santos"}\NormalTok{)}
\FunctionTok{paste}\NormalTok{(nomes, apelidos)}
\DocumentationTok{\#\# [1] "Ana Silva"      "Bruno Oliveira" "Carlos Santos"}

\CommentTok{\# Colar texto a cada elemento de um vetor}
\FunctionTok{paste}\NormalTok{(}\StringTok{"Número"}\NormalTok{, }\DecValTok{1}\SpecialCharTok{:}\DecValTok{3}\NormalTok{)}
\DocumentationTok{\#\# [1] "Número 1" "Número 2" "Número 3"}

\CommentTok{\# Sem separador, com paste0()}
\FunctionTok{paste0}\NormalTok{(}\StringTok{"Olá"}\NormalTok{, }\StringTok{"Mundo"}\NormalTok{)}
\DocumentationTok{\#\# [1] "OláMundo"}
\end{Highlighting}
\end{Shaded}

\section{Operadores lógicos e relacionais}\label{operadores-luxf3gicos-e-relacionais}

Terminamos com os operadores que permitem ao R \textbf{tomar decisões} --- comparar
valores e combinar condições. São a base das estruturas de decisão do
Capítulo~\ref{selecao}, por isso este é o primeiro passo para escrever
programas que reagem aos dados.

\subsection{Operadores relacionais}\label{operadores-relacionais}

Os operadores relacionais comparam valores e devolvem um valor lógico
(\texttt{TRUE} ou \texttt{FALSE}):

{\def\LTcaptype{none} % do not increment counter
\begin{longtable}[]{@{}ll@{}}
\toprule\noalign{}
Operador & Testa se\ldots{} \\
\midrule\noalign{}
\endhead
\bottomrule\noalign{}
\endlastfoot
\texttt{a\ ==\ b} & \texttt{a} é igual a \texttt{b} \\
\texttt{a\ !=\ b} & \texttt{a} é diferente de \texttt{b} \\
\texttt{a\ \textgreater{}\ b} & \texttt{a} é maior que \texttt{b} \\
\texttt{a\ \textless{}\ b} & \texttt{a} é menor que \texttt{b} \\
\texttt{a\ \textgreater{}=\ b} & \texttt{a} é maior ou igual a \texttt{b} \\
\texttt{a\ \textless{}=\ b} & \texttt{a} é menor ou igual a \texttt{b} \\
\end{longtable}
}

\begin{Shaded}
\begin{Highlighting}[]
\DecValTok{2} \SpecialCharTok{\textgreater{}} \DecValTok{1}
\DocumentationTok{\#\# [1] TRUE}

\DecValTok{1} \SpecialCharTok{==} \DecValTok{2}\SpecialCharTok{/}\DecValTok{2}
\DocumentationTok{\#\# [1] TRUE}

\DecValTok{1} \SpecialCharTok{!=} \DecValTok{2}
\DocumentationTok{\#\# [1] TRUE}
\end{Highlighting}
\end{Shaded}

Não confunda \texttt{=} com \texttt{==}. Um só sinal de igual (\texttt{=}) é uma \textbf{atribuição}; dois
sinais (\texttt{==}) são um \textbf{teste de igualdade}. Trocar um pelo outro é uma fonte
inesgotável de erros --- e, felizmente, dos mais fáceis de corrigir quando os
reconhecemos.

\subsection{Operadores lógicos}\label{operadores-luxf3gicos}

Os operadores lógicos combinam ou invertem condições:

\begin{itemize}
\tightlist
\item
  \texttt{\&} (\textbf{e} lógico): devolve \texttt{TRUE} apenas se \textbf{todas} as condições forem
  verdadeiras.
\item
  \texttt{\textbar{}} (\textbf{ou} lógico): devolve \texttt{TRUE} se \textbf{pelo menos uma} condição for
  verdadeira.
\item
  \texttt{!} (\textbf{não} lógico): inverte o valor lógico.
\end{itemize}

\begin{Shaded}
\begin{Highlighting}[]
\NormalTok{(}\DecValTok{5} \SpecialCharTok{\textgreater{}} \DecValTok{3}\NormalTok{) }\SpecialCharTok{\&}\NormalTok{ (}\DecValTok{4} \SpecialCharTok{\textgreater{}} \DecValTok{2}\NormalTok{)   }\CommentTok{\# ambas verdadeiras}
\DocumentationTok{\#\# [1] TRUE}

\NormalTok{(}\DecValTok{5} \SpecialCharTok{\textless{}} \DecValTok{3}\NormalTok{) }\SpecialCharTok{|}\NormalTok{ (}\DecValTok{4} \SpecialCharTok{\textgreater{}} \DecValTok{2}\NormalTok{)   }\CommentTok{\# basta uma ser verdadeira}
\DocumentationTok{\#\# [1] TRUE}

\SpecialCharTok{!}\NormalTok{(}\DecValTok{5} \SpecialCharTok{\textgreater{}} \DecValTok{3}\NormalTok{)            }\CommentTok{\# inverte TRUE em FALSE}
\DocumentationTok{\#\# [1] FALSE}
\end{Highlighting}
\end{Shaded}

Com estas peças --- valores, nomes, operações e comparações --- temos já o vocabulário
elementar de um programa. Nos exercícios que se seguem, vamos combiná-las para
resolver pequenos problemas. E, no próximo capítulo, aprenderemos a organizar os
valores em estruturas mais ricas, que nos permitirão trabalhar com conjuntos de
dados a sério.

\section{Exercícios}\label{exercuxedcios}

\textbf{1.} Escreva um programa em R que leia dois inteiros inseridos pelo utilizador
e imprima: a soma dos dois números; o produto; a diferença entre o primeiro e o
segundo; a divisão do primeiro pelo segundo; o resto dessa divisão; e o primeiro
número elevado à potência do segundo.
\emph{Sugestão:} use \texttt{readline()} para a entrada e \texttt{as.integer()} para a conversão.

\textbf{2.} Escreva um programa em R que leia dois números decimais e imprima a soma,
a diferença, o produto e o primeiro número elevado à potência do segundo.
\emph{Sugestão:} use \texttt{as.numeric()} para a conversão.

\textbf{3.} Escreva um programa em R que leia uma distância em milhas e a converta
para quilómetros, usando a fórmula \(K = M \times 1.609344\).

\textbf{4.} Escreva um programa em R que leia três inteiros correspondentes ao
comprimento, à largura e à altura de um paralelepípedo, e imprima o seu volume.

\textbf{5.} Escreva um programa em R que leia três inteiros e imprima a sua média.

\textbf{6.} Escreva um programa em R que leia uma temperatura em graus Fahrenheit e a
converta para graus Celsius, usando a fórmula \(C = \dfrac{F - 32}{1.8}\).

\textbf{7.} Escreva um programa em R que leia uma hora no formato de 24 horas e
imprima a hora correspondente no formato de 12 horas.

\textbf{8.} São exatamente 14h e definiu um alarme para tocar dentro de 51 horas. A
que horas tocará o alarme?
\emph{Sugestão:} o operador \texttt{\%\%} é seu amigo.

\textbf{9.} Generalize o problema anterior: peça ao utilizador a hora atual e o número
de horas de espera, e imprima a hora a que o alarme tocará.

\textbf{10.} Escreva um programa em R que leia um número inteiro e imprima \texttt{TRUE} se
for maior que 10, ou \texttt{FALSE} caso contrário.

\textbf{11.} Escreva um programa em R que leia dois números e imprima \texttt{TRUE} se forem
iguais, ou \texttt{FALSE} caso contrário.

\textbf{12.} Escreva um programa em R que leia dois números e imprima \texttt{TRUE} se o
primeiro for maior ou igual ao segundo, ou \texttt{FALSE} caso contrário.

\textbf{13.} Escreva um programa em R que leia um número e imprima \texttt{TRUE} se ele
estiver entre 0 e 100 (inclusive), ou \texttt{FALSE} caso contrário.

\textbf{14.} Escreva um programa em R que leia três números e imprima uma mensagem
indicando se o primeiro é menor que o segundo \emph{e} o segundo é menor que o
terceiro.

\chapter{Estruturas de dados}\label{estruturas}

Até aqui, trabalhámos sobretudo com valores isolados --- um número, um texto, um
resultado lógico. Mas a estatística raramente lida com valores isolados: lida com
\emph{conjuntos} de dados. As alturas de uma turma, as vendas de um ano, as respostas
de um inquérito --- tudo isto são coleções de valores que queremos guardar,
organizar e analisar em conjunto. É aqui que entram as \textbf{estruturas de dados}.

Uma estrutura de dados é uma forma de organizar vários valores num único objeto,
para que possamos manipulá-los como um todo. Escolher a estrutura certa para cada
problema é meio caminho andado: cada uma tem a sua vocação, e conhecê-las é
essencial para programar bem em R.

Antes de as estudarmos uma a uma, convém reter que, em R, existem dois grandes
tipos de objetos:

\begin{itemize}
\tightlist
\item
  \textbf{Funções} --- objetos que \emph{executam} operações, como \texttt{cos()}, \texttt{print()},
  \texttt{plot()} ou \texttt{integrate()}.
\item
  \textbf{Dados} --- objetos que \emph{contêm} informação, como \texttt{23}, \texttt{"Olá"}, \texttt{TRUE}, ou as
  estruturas que vamos conhecer agora.
\end{itemize}

Neste capítulo percorremos as quatro estruturas de dados fundamentais, das mais
simples às mais flexíveis: os \textbf{vetores} (coleções homogéneas), os \textbf{fatores}
(vetores para dados categóricos), as \textbf{matrizes e arrays} (coleções
bidimensionais e multidimensionais), os \textbf{data frames} (tabelas) e as \textbf{listas}
(coleções heterogéneas). Cada uma responde a uma necessidade diferente.

\section{Vetor}\label{vetor}

O \textbf{vetor} é a estrutura mais elementar e, ao mesmo tempo, a mais importante do
R --- a tal ponto que, mesmo um único número, é internamente um vetor de comprimento
um. Um vetor é uma sequência ordenada de elementos \textbf{do mesmo tipo}: todos
números, ou todos texto, ou todos valores lógicos.

Três ideias a reter desde já:

\begin{itemize}
\tightlist
\item
  todos os elementos de um vetor têm de ser do mesmo tipo;
\item
  os elementos são indexados a partir de \textbf{1} (o primeiro elemento está na
  posição 1, não 0, como em muitas outras linguagens);
\item
  criam-se com a função \texttt{c()} e manipulam-se com uma vasta coleção de funções.
\end{itemize}

\subsection{Tipos de vetores}\label{tipos-de-vetores}

\begin{Shaded}
\begin{Highlighting}[]
\CommentTok{\# Vetor numérico}
\FunctionTok{c}\NormalTok{(}\FloatTok{1.1}\NormalTok{, }\FloatTok{2.2}\NormalTok{, }\FloatTok{3.3}\NormalTok{)}
\DocumentationTok{\#\# [1] 1.1 2.2 3.3}

\CommentTok{\# Vetor de carateres}
\FunctionTok{c}\NormalTok{(}\StringTok{"a"}\NormalTok{, }\StringTok{"b"}\NormalTok{, }\StringTok{"c"}\NormalTok{)}
\DocumentationTok{\#\# [1] "a" "b" "c"}

\CommentTok{\# Vetor lógico}
\FunctionTok{c}\NormalTok{(}\ConstantTok{TRUE}\NormalTok{, }\DecValTok{1} \SpecialCharTok{==} \DecValTok{2}\NormalTok{)}
\DocumentationTok{\#\# [1]  TRUE FALSE}

\CommentTok{\# O que acontece se misturarmos tipos?}
\FunctionTok{c}\NormalTok{(}\DecValTok{3}\NormalTok{, }\DecValTok{1} \SpecialCharTok{==} \DecValTok{2}\NormalTok{, }\StringTok{"a"}\NormalTok{)}
\DocumentationTok{\#\# [1] "3"     "FALSE" "a"}
\end{Highlighting}
\end{Shaded}

Repare no último exemplo. Ao tentarmos misturar um número, um valor lógico e um
texto, o R \textbf{não} dá erro: converte tudo para o tipo mais ``geral'' --- neste caso,
texto. A este fenómeno chamamos \textbf{coerção}. É cómodo, mas traiçoeiro: se
esperava um vetor de números e o R o converteu silenciosamente em texto, os
cálculos seguintes falharão. Convém, por isso, garantir que cada vetor contém um
só tipo de dados.

\subsection{Construir vetores}\label{construir-vetores}

Escrever todos os elementos à mão seria impraticável para vetores longos. O R
oferece atalhos poderosos para os construir:

\begin{Shaded}
\begin{Highlighting}[]
\CommentTok{\# Inteiros de 1 a 10}
\NormalTok{x }\OtherTok{\textless{}{-}} \DecValTok{1}\SpecialCharTok{:}\DecValTok{10}
\NormalTok{x}
\DocumentationTok{\#\#  [1]  1  2  3  4  5  6  7  8  9 10}

\CommentTok{\# Em ordem decrescente}
\DecValTok{10}\SpecialCharTok{:}\DecValTok{1}
\DocumentationTok{\#\#  [1] 10  9  8  7  6  5  4  3  2  1}

\CommentTok{\# Sequência de 0 a 50, de 10 em 10}
\FunctionTok{seq}\NormalTok{(}\AttributeTok{from =} \DecValTok{0}\NormalTok{, }\AttributeTok{to =} \DecValTok{50}\NormalTok{, }\AttributeTok{by =} \DecValTok{10}\NormalTok{)}
\DocumentationTok{\#\# [1]  0 10 20 30 40 50}

\CommentTok{\# 15 números igualmente espaçados entre 0 e 1}
\FunctionTok{seq}\NormalTok{(}\DecValTok{0}\NormalTok{, }\DecValTok{1}\NormalTok{, }\AttributeTok{length =} \DecValTok{15}\NormalTok{)}
\DocumentationTok{\#\#  [1] 0.00000 0.07143 0.14286 0.21429 0.28571 0.35714 0.42857 0.50000}
\DocumentationTok{\#\#  [9] 0.57143 0.64286 0.71429 0.78571 0.85714 0.92857 1.00000}

\CommentTok{\# Repetir um vetor inteiro várias vezes}
\FunctionTok{rep}\NormalTok{(}\DecValTok{1}\SpecialCharTok{:}\DecValTok{3}\NormalTok{, }\AttributeTok{times =} \DecValTok{4}\NormalTok{)}
\DocumentationTok{\#\#  [1] 1 2 3 1 2 3 1 2 3 1 2 3}

\CommentTok{\# Repetir cada elemento várias vezes}
\FunctionTok{rep}\NormalTok{(}\DecValTok{1}\SpecialCharTok{:}\DecValTok{3}\NormalTok{, }\AttributeTok{each =} \DecValTok{4}\NormalTok{)}
\DocumentationTok{\#\#  [1] 1 1 1 1 2 2 2 2 3 3 3 3}

\CommentTok{\# Combinar vetores e números num novo vetor}
\FunctionTok{c}\NormalTok{(}\DecValTok{1}\SpecialCharTok{:}\DecValTok{3}\NormalTok{, }\DecValTok{100}\NormalTok{, }\DecValTok{1}\SpecialCharTok{:}\DecValTok{3}\NormalTok{)}
\DocumentationTok{\#\# [1]   1   2   3 100   1   2   3}
\end{Highlighting}
\end{Shaded}

O operador \texttt{:} e as funções \texttt{seq()} e \texttt{rep()} vão aparecer constantemente. Note a
diferença subtil entre \texttt{times} e \texttt{each} em \texttt{rep()}: \texttt{times} repete a sequência
inteira; \texttt{each} repete cada elemento antes de passar ao seguinte. Confundir os
dois é comum --- experimente ambos até a distinção ficar natural.

\subsection{Aceder a elementos}\label{aceder-a-elementos}

Guardar um vetor só é útil se depois conseguirmos extrair partes dele. Fazemo-lo
com \textbf{colchetes} \texttt{{[}\ {]}} e índices. Esta é uma das operações que mais vamos usar,
por isso vale a pena dominá-la bem. Há três formas principais de indexar.

\textbf{Por posição.} Indicamos os índices dos elementos que queremos:

\begin{Shaded}
\begin{Highlighting}[]
\NormalTok{x }\OtherTok{\textless{}{-}}\NormalTok{ (}\SpecialCharTok{{-}}\DecValTok{5}\NormalTok{)}\SpecialCharTok{:}\DecValTok{5}
\NormalTok{x}
\DocumentationTok{\#\#  [1] {-}5 {-}4 {-}3 {-}2 {-}1  0  1  2  3  4  5}

\NormalTok{x[}\DecValTok{4}\SpecialCharTok{:}\DecValTok{8}\NormalTok{]           }\CommentTok{\# do quarto ao oitavo elemento}
\DocumentationTok{\#\# [1] {-}2 {-}1  0  1  2}

\NormalTok{x[}\SpecialCharTok{{-}}\FunctionTok{length}\NormalTok{(x)]    }\CommentTok{\# todos menos o último}
\DocumentationTok{\#\# [1] {-}5 {-}4 {-}3 {-}2 {-}1  0  1  2  3  4}
\end{Highlighting}
\end{Shaded}

\textbf{Por exclusão.} Um índice negativo \emph{remove} elementos:

\begin{Shaded}
\begin{Highlighting}[]
\NormalTok{x[}\SpecialCharTok{{-}}\FunctionTok{c}\NormalTok{(}\DecValTok{1}\SpecialCharTok{:}\DecValTok{3}\NormalTok{, }\DecValTok{9}\SpecialCharTok{:}\DecValTok{11}\NormalTok{)]   }\CommentTok{\# remove os três primeiros e os três últimos}
\DocumentationTok{\#\# [1] {-}2 {-}1  0  1  2}
\end{Highlighting}
\end{Shaded}

\textbf{Por condição lógica.} Esta é a forma mais poderosa: fornecemos um vetor de
\texttt{TRUE}/\texttt{FALSE}, e o R devolve os elementos nas posições onde há \texttt{TRUE}. Como as
comparações produzem exatamente esses vetores lógicos, podemos filtrar dados de
forma muito expressiva:

\begin{Shaded}
\begin{Highlighting}[]
\CommentTok{\# Selecionar os elementos com valor absoluto menor que 3}
\NormalTok{index }\OtherTok{\textless{}{-}} \FunctionTok{abs}\NormalTok{(x) }\SpecialCharTok{\textless{}} \DecValTok{3}
\NormalTok{index}
\DocumentationTok{\#\#  [1] FALSE FALSE FALSE  TRUE  TRUE  TRUE  TRUE  TRUE FALSE FALSE FALSE}

\NormalTok{x[index]}
\DocumentationTok{\#\# [1] {-}2 {-}1  0  1  2}

\CommentTok{\# ... ou tudo de uma vez, sem variável auxiliar}
\NormalTok{x[}\FunctionTok{abs}\NormalTok{(x) }\SpecialCharTok{\textless{}} \DecValTok{3}\NormalTok{]}
\DocumentationTok{\#\# [1] {-}2 {-}1  0  1  2}
\end{Highlighting}
\end{Shaded}

Mais exemplos desta indexação condicional, que convém estudar com calma:

\begin{Shaded}
\begin{Highlighting}[]
\NormalTok{y }\OtherTok{\textless{}{-}} \DecValTok{1}\SpecialCharTok{:}\DecValTok{10}
\NormalTok{y[y }\SpecialCharTok{\textgreater{}} \DecValTok{5}\NormalTok{]           }\CommentTok{\# os maiores que 5}
\DocumentationTok{\#\# [1]  6  7  8  9 10}

\NormalTok{y[y }\SpecialCharTok{\%\%} \DecValTok{2} \SpecialCharTok{==} \DecValTok{0}\NormalTok{]     }\CommentTok{\# os pares}
\DocumentationTok{\#\# [1]  2  4  6  8 10}

\NormalTok{y[}\DecValTok{5}\NormalTok{] }\OtherTok{\textless{}{-}} \ConstantTok{NA}
\NormalTok{y[}\SpecialCharTok{!}\FunctionTok{is.na}\NormalTok{(y)]       }\CommentTok{\# todos os que não são NA}
\DocumentationTok{\#\# [1]  1  2  3  4  6  7  8  9 10}
\end{Highlighting}
\end{Shaded}

O R traz ainda alguns vetores predefinidos convenientes, como \texttt{letters} e
\texttt{LETTERS} (o alfabeto em minúsculas e maiúsculas):

\begin{Shaded}
\begin{Highlighting}[]
\NormalTok{letters[}\DecValTok{1}\SpecialCharTok{:}\DecValTok{3}\NormalTok{]}
\DocumentationTok{\#\# [1] "a" "b" "c"}

\NormalTok{LETTERS[}\FunctionTok{c}\NormalTok{(}\DecValTok{2}\NormalTok{, }\DecValTok{4}\NormalTok{, }\DecValTok{6}\NormalTok{)]}
\DocumentationTok{\#\# [1] "B" "D" "F"}
\end{Highlighting}
\end{Shaded}

\subsection{Funções úteis para vetores}\label{funuxe7uxf5es-uxfateis-para-vetores}

Grande parte do valor dos vetores está nas dezenas de funções que operam sobre
eles de uma só vez. Eis um catálogo das mais frequentes, aplicadas a um vetor
numérico:

\begin{Shaded}
\begin{Highlighting}[]
\NormalTok{v }\OtherTok{\textless{}{-}} \FunctionTok{c}\NormalTok{(}\FloatTok{2.2}\NormalTok{, }\FloatTok{1.1}\NormalTok{, }\FloatTok{3.3}\NormalTok{)}

\FunctionTok{length}\NormalTok{(v)     }\CommentTok{\# número de elementos}
\DocumentationTok{\#\# [1] 3}
\FunctionTok{max}\NormalTok{(v)        }\CommentTok{\# máximo}
\DocumentationTok{\#\# [1] 3.3}
\FunctionTok{min}\NormalTok{(v)        }\CommentTok{\# mínimo}
\DocumentationTok{\#\# [1] 1.1}
\FunctionTok{sum}\NormalTok{(v)        }\CommentTok{\# soma}
\DocumentationTok{\#\# [1] 6.6}
\FunctionTok{mean}\NormalTok{(v)       }\CommentTok{\# média}
\DocumentationTok{\#\# [1] 2.2}
\FunctionTok{median}\NormalTok{(v)     }\CommentTok{\# mediana}
\DocumentationTok{\#\# [1] 2.2}
\FunctionTok{range}\NormalTok{(v)      }\CommentTok{\# mínimo e máximo}
\DocumentationTok{\#\# [1] 1.1 3.3}
\FunctionTok{var}\NormalTok{(v)        }\CommentTok{\# variância amostral}
\DocumentationTok{\#\# [1] 1.21}
\FunctionTok{sort}\NormalTok{(v)       }\CommentTok{\# ordenar (crescente)}
\DocumentationTok{\#\# [1] 1.1 2.2 3.3}
\FunctionTok{sort}\NormalTok{(v, }\AttributeTok{decreasing =} \ConstantTok{TRUE}\NormalTok{)   }\CommentTok{\# ordenar (decrescente)}
\DocumentationTok{\#\# [1] 3.3 2.2 1.1}
\FunctionTok{cumsum}\NormalTok{(v)     }\CommentTok{\# somas acumuladas}
\DocumentationTok{\#\# [1] 2.2 3.3 6.6}
\end{Highlighting}
\end{Shaded}

Uma função particularmente útil é \texttt{which()}, que devolve as \textbf{posições} dos
elementos que satisfazem uma condição --- e não os valores em si:

\begin{Shaded}
\begin{Highlighting}[]
\NormalTok{y }\OtherTok{\textless{}{-}} \FunctionTok{c}\NormalTok{(}\DecValTok{8}\NormalTok{, }\DecValTok{3}\NormalTok{, }\DecValTok{5}\NormalTok{, }\DecValTok{7}\NormalTok{, }\DecValTok{6}\NormalTok{, }\DecValTok{6}\NormalTok{, }\DecValTok{8}\NormalTok{, }\DecValTok{9}\NormalTok{, }\DecValTok{2}\NormalTok{, }\DecValTok{3}\NormalTok{, }\DecValTok{9}\NormalTok{, }\DecValTok{4}\NormalTok{, }\DecValTok{10}\NormalTok{, }\DecValTok{4}\NormalTok{, }\DecValTok{11}\NormalTok{)}

\FunctionTok{which}\NormalTok{(y }\SpecialCharTok{\textgreater{}} \DecValTok{5}\NormalTok{)     }\CommentTok{\# as posições dos elementos maiores que 5}
\DocumentationTok{\#\# [1]  1  4  5  6  7  8 11 13 15}

\NormalTok{y[y }\SpecialCharTok{\textgreater{}} \DecValTok{5}\NormalTok{]         }\CommentTok{\# os valores em si}
\DocumentationTok{\#\# [1]  8  7  6  6  8  9  9 10 11}
\end{Highlighting}
\end{Shaded}

A distinção entre \emph{posições} (\texttt{which(y\ \textgreater{}\ 5)}) e \emph{valores} (\texttt{y{[}y\ \textgreater{}\ 5{]}}) é uma
fonte comum de confusão. Guarde-a: uma pergunta como ``em que dias o consumo
excedeu a média?'' pede posições; ``quais foram os consumos acima da média?'' pede
valores.

\subsection{Operações e reciclagem}\label{operauxe7uxf5es-e-reciclagem}

As operações aritméticas em R são \textbf{vetorizadas}: aplicam-se a cada elemento do
vetor, sem necessidade de escrever ciclos. Isto torna o código curto e legível:

\begin{Shaded}
\begin{Highlighting}[]
\NormalTok{v }\OtherTok{\textless{}{-}} \FunctionTok{c}\NormalTok{(}\FloatTok{2.2}\NormalTok{, }\FloatTok{1.1}\NormalTok{, }\FloatTok{3.3}\NormalTok{)}

\NormalTok{v }\SpecialCharTok{+} \DecValTok{1}        \CommentTok{\# soma 1 a cada elemento}
\DocumentationTok{\#\# [1] 3.2 2.1 4.3}

\NormalTok{v }\SpecialCharTok{*} \DecValTok{2}        \CommentTok{\# multiplica cada elemento por 2}
\DocumentationTok{\#\# [1] 4.4 2.2 6.6}

\NormalTok{v }\SpecialCharTok{\textgreater{}} \DecValTok{2}        \CommentTok{\# compara cada elemento com 2}
\DocumentationTok{\#\# [1]  TRUE FALSE  TRUE}

\FunctionTok{c}\NormalTok{(}\DecValTok{2}\NormalTok{, }\DecValTok{3}\NormalTok{, }\DecValTok{5}\NormalTok{, }\DecValTok{7}\NormalTok{)}\SpecialCharTok{\^{}}\DecValTok{2}   \CommentTok{\# eleva cada elemento ao quadrado}
\DocumentationTok{\#\# [1]  4  9 25 49}
\end{Highlighting}
\end{Shaded}

Mas o que acontece quando operamos com dois vetores de \textbf{comprimentos
diferentes}? O R aplica a chamada \textbf{regra da reciclagem}: repete o vetor mais
curto até igualar o comprimento do mais longo.

\begin{Shaded}
\begin{Highlighting}[]
\FunctionTok{c}\NormalTok{(}\DecValTok{2}\NormalTok{, }\DecValTok{3}\NormalTok{, }\DecValTok{5}\NormalTok{, }\DecValTok{7}\NormalTok{)}\SpecialCharTok{\^{}}\FunctionTok{c}\NormalTok{(}\DecValTok{2}\NormalTok{, }\DecValTok{3}\NormalTok{)}
\DocumentationTok{\#\# [1]   4  27  25 343}
\end{Highlighting}
\end{Shaded}

Internamente, o vetor \texttt{c(2,\ 3)} foi reciclado para \texttt{c(2,\ 3,\ 2,\ 3)}, de modo que a
operação passou a ser \texttt{c(2,\ 3,\ 5,\ 7)\^{}c(2,\ 3,\ 2,\ 3)}.

A reciclagem é cómoda, mas pode esconder erros. Se o comprimento do vetor maior
não for múltiplo do menor, a reciclagem fica ``incompleta'' e o R emite um aviso:

\begin{Shaded}
\begin{Highlighting}[]
\FunctionTok{c}\NormalTok{(}\DecValTok{2}\NormalTok{, }\DecValTok{3}\NormalTok{, }\DecValTok{5}\NormalTok{, }\DecValTok{7}\NormalTok{)}\SpecialCharTok{\^{}}\FunctionTok{c}\NormalTok{(}\DecValTok{2}\NormalTok{, }\DecValTok{3}\NormalTok{, }\DecValTok{4}\NormalTok{)}
\DocumentationTok{\#\# Warning: longer object length is not a multiple of shorter object length}
\DocumentationTok{\#\# [1]   4  27 625  49}
\end{Highlighting}
\end{Shaded}

Aqui, \texttt{c(2,\ 3,\ 4)} foi reciclado para \texttt{c(2,\ 3,\ 4,\ 2)} --- o último elemento repetiu
apenas parcialmente. \textbf{Leia sempre os avisos do R:} um aviso de reciclagem é,
muitas vezes, o sintoma de um vetor com o comprimento errado.

\subsection{Exercícios}\label{exercuxedcios-1}

\textbf{1.} Crie os vetores:

\begin{enumerate}
\def\labelenumi{(\alph{enumi})}
\tightlist
\item
  \((1, 2, 3, \ldots, 19, 20)\)
\item
  \((20, 19, \ldots, 2, 1)\)
\item
  \((1, 2, \ldots, 19, 20, 19, 18, \ldots, 2, 1)\)
\item
  \((10, 20, 30, \ldots, 90, 10)\)
\item
  \((1, 1, \ldots, 1, 2, 2, \ldots, 2, 3, 3, \ldots, 3)\), com 10 ocorrências do
  1, 20 do 2 e 30 do 3.
\end{enumerate}

\textbf{2.} Use a função \texttt{paste()} para criar o vetor de carateres de tamanho 20:
\texttt{"nome\ 1"}, \texttt{"nome\ 2"}, \ldots, \texttt{"nome\ 20"}.

\textbf{3.} Crie um vetor \(x_1\) igual a ``A'' ``A'' ``B'' ``B'' ``C'' ``C'' ``D'' ``D'' ``E'' ``E''.

\textbf{4.} Crie um vetor \(x_2\) igual a ``a'' ``b'' ``c'' ``d'' ``e'' ``a'' ``b'' ``c'' ``d'' ``e''.

\textbf{5.} Crie um vetor \(x_3\) com a palavra ``uva'' 10 vezes, ``maçã'' 9 vezes,
``laranja'' 6 vezes e ``banana'' 1 vez.

\textbf{6.} Crie o vetor \texttt{notas\ \textless{}-\ c(7,\ 8,\ 9)} e atribua-lhe os nomes dos alunos
``Ana'', ``João'' e ``Pedro''. Aceda à nota do ``João'' usando o nome.
\emph{Sugestão:} use a função \texttt{names()}.

\textbf{7.} Crie um vetor de 15 números aleatórios entre 1 e 100 (use \texttt{sample()}).
Ordene-o por ordem crescente e depois decrescente. Encontre o menor e o maior
valor.

\textbf{8.} Crie um vetor de 20 números aleatórios entre 1 e 50 (use \texttt{sample()}).
Calcule a soma, a média, o desvio padrão (\texttt{sd()}) e o produto (\texttt{prod()}) dos seus
elementos.

\textbf{9.} Crie um vetor de 10 números aleatórios entre 1 e 100 (use \texttt{sample()}).
Extraia os elementos maiores que 50. De seguida, substitua por 0 os valores
menores que 30.

\textbf{10.} Crie um vetor de 10 números. Verifique quais elementos são maiores que 5
e quais são pares. Crie um novo vetor apenas com os que satisfazem \textbf{ambas} as
condições.

\textbf{11.} Crie um vetor com 10 números inteiros. Multiplique por 2 os elementos nas
posições 2, 4 e 6. Substitua o último elemento por 100.

\textbf{12.} Calcule a média dos vetores:

\begin{enumerate}
\def\labelenumi{(\alph{enumi})}
\tightlist
\item
  \(x = (1, 0, NA, 5, 7)\)
\item
  \(y = (-\infty, 0, 1, NA, 2, 7, \infty)\)
\end{enumerate}

\textbf{13.} Crie:

\begin{enumerate}
\def\labelenumi{(\alph{enumi})}
\tightlist
\item
  um vetor com os valores \(e^{x}\sin(x)\) nos pontos \(x = 2, 2{,}1, 2{,}2,
  \ldots, 6\);
\item
  o vetor \(\left(3, \dfrac{3^2}{2}, \dfrac{3^3}{3}, \ldots,
  \dfrac{3^{30}}{30}\right)\).
\end{enumerate}

\textbf{14.} Calcule:

\begin{enumerate}
\def\labelenumi{(\alph{enumi})}
\tightlist
\item
  \(\displaystyle\sum_{i=1}^{1000} \frac{9}{10^{i}}\)
\item
  \(\displaystyle\sum_{i=10}^{100} i^{3} + 4i^{2}\)
\item
  \(\displaystyle\sum_{i=1}^{25} \frac{2^{i}}{i} + \frac{3^{i}}{i^2}\)
\end{enumerate}

\textbf{15.} Uma empresa monitorizou o consumo diário de energia (em kWh) de uma casa
durante 30 dias, guardando-os num vetor \texttt{consumo}.

\begin{enumerate}
\def\labelenumi{(\alph{enumi})}
\tightlist
\item
  Defina \texttt{consumo} com 30 valores gerados aleatoriamente entre 10 e 40 kWh,
  usando \texttt{runif()} com semente 2025 (\texttt{set.seed(2025)}). Consulte \texttt{?runif} para
  os parâmetros.
\item
  Calcule o consumo médio e guarde-o em \texttt{media}.
\item
  Crie um vetor lógico \texttt{acima\_da\_media} indicando os dias com consumo superior à
  média.
\item
  Em quantos dias o consumo superou a média?
\item
  Quais foram os dias com consumo de pico (superior a 35 kWh)? Devolva os
  valores e as posições.
\item
  Substitua por \texttt{NA} todos os valores abaixo de 15 kWh (simulando falhas de
  medição).
\item
  Recalcule a média ignorando os \texttt{NA}.
\item
  Crie um vetor com os consumos dos dias pares do mês.
\end{enumerate}

\textbf{16.} Uma cooperativa agrícola registou a produção anual de soja (em toneladas)
de 12 propriedades.

\begin{enumerate}
\def\labelenumi{(\alph{enumi})}
\tightlist
\item
  Crie o vetor \texttt{producao}:
\end{enumerate}

\begin{Shaded}
\begin{Highlighting}[]
\NormalTok{producao }\OtherTok{\textless{}{-}} \FunctionTok{c}\NormalTok{(}\DecValTok{80}\NormalTok{, }\DecValTok{95}\NormalTok{, }\DecValTok{70}\NormalTok{, }\DecValTok{68}\NormalTok{, }\DecValTok{105}\NormalTok{, }\DecValTok{120}\NormalTok{, }\DecValTok{73}\NormalTok{, }\DecValTok{89}\NormalTok{, }\DecValTok{110}\NormalTok{, }\DecValTok{66}\NormalTok{, }\DecValTok{92}\NormalTok{, }\DecValTok{78}\NormalTok{)}
\end{Highlighting}
\end{Shaded}

\begin{enumerate}
\def\labelenumi{(\alph{enumi})}
\setcounter{enumi}{1}
\tightlist
\item
  Calcule a produção média por propriedade e guarde-a em \texttt{media\_total}.
\item
  Identifique as propriedades com produção abaixo da média e crie um vetor com
  os seus índices.
\item
  Uma seca afetou as propriedades 4, 10 e 12, cuja produção foi subestimada em
  10\%. Corrija esses valores no vetor \texttt{producao}.
\item
  Recalcule a média com os dados corrigidos e compare com a anterior.
\item
  Crie um vetor lógico \texttt{alta\_produtividade} para as propriedades com produção
  superior a 100 toneladas.
\item
  Quantas tiveram alta produtividade? Que percentagem do total?
\item
  Ordene a produção do menor para o maior, guardando o resultado em \texttt{ordenado}.
\item
  Quais são as 3 propriedades com maior produção?
\end{enumerate}

\section{Fatores}\label{fatores}

Nem todos os dados são números ou textos livres. Muitas variáveis são
\textbf{categóricas}: assumem um número limitado de valores possíveis --- o sexo, o
estado civil, o grau de satisfação, o nível de escolaridade. Poderíamos
guardá-las como vetores de texto, mas o R oferece uma estrutura feita à medida
para elas: o \textbf{fator}.

Um fator é um vetor especial, concebido para representar dados categóricos. À
primeira vista parece um vetor de texto, mas por baixo esconde uma organização
mais inteligente, que torna as análises estatísticas mais eficientes e mais
corretas. Três conceitos definem um fator:

\begin{itemize}
\tightlist
\item
  \textbf{Níveis} (\emph{levels}): os valores distintos que a categoria pode assumir. Para
  uma variável ``tamanho'', os níveis poderiam ser ``pequeno'', ``médio'' e ``grande''.
\item
  \textbf{Armazenamento interno}: internamente, o R guarda cada valor como um inteiro
  que aponta para um nível. É mais eficiente e evita erros de escrita --- se ``médio''
  é um nível, não há forma de acidentalmente escrever ``medio'' numa observação.
\item
  \textbf{Ordem}: os fatores podem ser \textbf{ordenados} (quando há uma hierarquia natural,
  como ``baixo'' \textless{} ``médio'' \textless{} ``alto'') ou \textbf{não ordenados} (quando os níveis não têm
  ordem intrínseca, como cores).
\end{itemize}

\subsection{Criar fatores}\label{criar-fatores}

Criamos um fator com a função \texttt{factor()}:

\begin{Shaded}
\begin{Highlighting}[]
\NormalTok{dados }\OtherTok{\textless{}{-}} \FunctionTok{c}\NormalTok{(}\StringTok{"baixo"}\NormalTok{, }\StringTok{"medio"}\NormalTok{, }\StringTok{"alto"}\NormalTok{, }\StringTok{"medio"}\NormalTok{, }\StringTok{"baixo"}\NormalTok{, }\StringTok{"alto"}\NormalTok{)}
\NormalTok{factor\_dados }\OtherTok{\textless{}{-}} \FunctionTok{factor}\NormalTok{(dados)}
\NormalTok{factor\_dados}
\DocumentationTok{\#\# [1] baixo medio alto  medio baixo alto }
\DocumentationTok{\#\# Levels: alto baixo medio}
\end{Highlighting}
\end{Shaded}

Por omissão, o R ordena os níveis \textbf{alfabeticamente} --- repare que ficaram ``alto'',
``baixo'', ``medio'', o que quase nunca é a ordem que queremos. Para variáveis com uma
hierarquia natural, isto é um problema: um gráfico ou uma tabela apareceriam por
ordem alfabética em vez da ordem lógica. A solução é indicar explicitamente a
ordem dos níveis.

\begin{Shaded}
\begin{Highlighting}[]
\CommentTok{\# Especificar a ordem dos níveis}
\NormalTok{factor\_dados }\OtherTok{\textless{}{-}} \FunctionTok{factor}\NormalTok{(dados, }\AttributeTok{levels =} \FunctionTok{c}\NormalTok{(}\StringTok{"baixo"}\NormalTok{, }\StringTok{"medio"}\NormalTok{, }\StringTok{"alto"}\NormalTok{))}
\NormalTok{factor\_dados}
\DocumentationTok{\#\# [1] baixo medio alto  medio baixo alto }
\DocumentationTok{\#\# Levels: baixo medio alto}
\end{Highlighting}
\end{Shaded}

Se, além da ordem de apresentação, quisermos que o R reconheça uma verdadeira
\textbf{hierarquia} entre os níveis (para poder responder a perguntas do tipo ``é este
nível maior do que aquele?''), usamos \texttt{ordered\ =\ TRUE}:

\begin{Shaded}
\begin{Highlighting}[]
\NormalTok{fator\_ordenado }\OtherTok{\textless{}{-}} \FunctionTok{factor}\NormalTok{(dados,}
                         \AttributeTok{levels =} \FunctionTok{c}\NormalTok{(}\StringTok{"baixo"}\NormalTok{, }\StringTok{"medio"}\NormalTok{, }\StringTok{"alto"}\NormalTok{),}
                         \AttributeTok{ordered =} \ConstantTok{TRUE}\NormalTok{)}
\NormalTok{fator\_ordenado[}\DecValTok{1}\NormalTok{] }\SpecialCharTok{\textless{}}\NormalTok{ fator\_ordenado[}\DecValTok{3}\NormalTok{]   }\CommentTok{\# "baixo" \textless{} "alto"?}
\DocumentationTok{\#\# [1] TRUE}
\end{Highlighting}
\end{Shaded}

\subsection{Manipular fatores}\label{manipular-fatores}

Podemos consultar e modificar os níveis de um fator:

\begin{Shaded}
\begin{Highlighting}[]
\CommentTok{\# Consultar os níveis}
\FunctionTok{levels}\NormalTok{(factor\_dados)}
\DocumentationTok{\#\# [1] "baixo" "medio" "alto"}

\CommentTok{\# Renomear os níveis}
\FunctionTok{levels}\NormalTok{(factor\_dados) }\OtherTok{\textless{}{-}} \FunctionTok{c}\NormalTok{(}\StringTok{"Baixo"}\NormalTok{, }\StringTok{"Medio"}\NormalTok{, }\StringTok{"Alto"}\NormalTok{)}
\NormalTok{factor\_dados}
\DocumentationTok{\#\# [1] Baixo Medio Alto  Medio Baixo Alto }
\DocumentationTok{\#\# Levels: Baixo Medio Alto}
\end{Highlighting}
\end{Shaded}

\subsection{\texorpdfstring{Gerar fatores com \texttt{gl()}}{Gerar fatores com gl()}}\label{gerar-fatores-com-gl}

Quando precisamos de fatores com padrões regulares --- típico em desenhos
experimentais, onde há grupos de igual dimensão --- a função \texttt{gl()} (\emph{generate
levels}) poupa trabalho:

\begin{Shaded}
\begin{Highlighting}[]
\CommentTok{\# gl(n, k): n níveis, cada um repetido k vezes}
\FunctionTok{gl}\NormalTok{(}\AttributeTok{n =} \DecValTok{4}\NormalTok{, }\AttributeTok{k =} \DecValTok{3}\NormalTok{)}
\DocumentationTok{\#\#  [1] 1 1 1 2 2 2 3 3 3 4 4 4}
\DocumentationTok{\#\# Levels: 1 2 3 4}

\CommentTok{\# Com rótulos personalizados}
\FunctionTok{gl}\NormalTok{(}\AttributeTok{n =} \DecValTok{2}\NormalTok{, }\AttributeTok{k =} \DecValTok{10}\NormalTok{, }\AttributeTok{labels =} \FunctionTok{c}\NormalTok{(}\StringTok{"mulher"}\NormalTok{, }\StringTok{"homem"}\NormalTok{))}
\DocumentationTok{\#\#  [1] mulher mulher mulher mulher mulher mulher mulher mulher mulher mulher}
\DocumentationTok{\#\# [11] homem  homem  homem  homem  homem  homem  homem  homem  homem  homem }
\DocumentationTok{\#\# Levels: mulher homem}
\end{Highlighting}
\end{Shaded}

\subsection{Exercícios}\label{exercuxedcios-2}

\textbf{1.} Crie o vetor \texttt{c("pequeno",\ "médio",\ "grande",\ "pequeno",\ "grande")} e
transforme-o num fator com níveis ordenados ``pequeno'', ``médio'', ``grande''. Exiba os
níveis.

\textbf{2.} Crie um fator com as cores de carros \texttt{c("vermelho",\ "azul",\ "preto",\ "vermelho",\ "branco",\ "azul")}. Use \texttt{table()} para contar quantos carros há de
cada cor.

\textbf{3.} Crie um fator com avaliações ``bom'', ``médio'', ``mau'' e modifique os níveis
para ``excelente'', ``razoável'' e ``insatisfatório''. Exiba o fator modificado.

\textbf{4.} Crie o vetor \texttt{c("solteiro",\ "casado",\ "divorciado")}. Verifique se é um
fator com \texttt{is.factor()} e depois converta-o em fator.

\textbf{5.} Crie um fator de escolaridade com ``ensino superior'', ``ensino secundário'',
``doutoramento'', ``mestrado'', ``ensino secundário'', ``ensino superior''. Defina uma
ordem lógica para os níveis e ordene o fator.

\textbf{6.} Crie o fator \texttt{c("baixo",\ "alto",\ "médio",\ "baixo",\ "alto")} com níveis
``baixo'', ``médio'', ``alto''. Converta-o num vetor numérico em que ``baixo'' seja 1,
``médio'' 2 e ``alto'' 3.

\textbf{7.} Crie um fator com as categorias ``verde'', ``amarelo'', ``vermelho''. Substitua o
nível ``amarelo'' por ``laranja'' e exiba o fator.

\textbf{8.} Uma empresa de Internet inquiriu 20 clientes sobre o atendimento,
classificando as respostas em ``Muito insatisfeito'', ``Insatisfeito'', ``Neutro'',
``Satisfeito'' e ``Muito satisfeito''.

\begin{enumerate}
\def\labelenumi{(\alph{enumi})}
\tightlist
\item
  Crie o vetor \texttt{respostas}:
\end{enumerate}

\begin{Shaded}
\begin{Highlighting}[]
\NormalTok{respostas }\OtherTok{\textless{}{-}} \FunctionTok{c}\NormalTok{(}\StringTok{"Satisfeito"}\NormalTok{, }\StringTok{"Muito satisfeito"}\NormalTok{, }\StringTok{"Neutro"}\NormalTok{, }\StringTok{"Insatisfeito"}\NormalTok{, }\StringTok{"Satisfeito"}\NormalTok{,}
  \StringTok{"Muito insatisfeito"}\NormalTok{, }\StringTok{"Satisfeito"}\NormalTok{, }\StringTok{"Neutro"}\NormalTok{, }\StringTok{"Insatisfeito"}\NormalTok{, }\StringTok{"Satisfeito"}\NormalTok{,}
  \StringTok{"Muito satisfeito"}\NormalTok{, }\StringTok{"Neutro"}\NormalTok{, }\StringTok{"Insatisfeito"}\NormalTok{, }\StringTok{"Satisfeito"}\NormalTok{, }\StringTok{"Neutro"}\NormalTok{,}
  \StringTok{"Satisfeito"}\NormalTok{, }\StringTok{"Muito satisfeito"}\NormalTok{, }\StringTok{"Muito insatisfeito"}\NormalTok{, }\StringTok{"Neutro"}\NormalTok{, }\StringTok{"Satisfeito"}\NormalTok{)}
\end{Highlighting}
\end{Shaded}

\begin{enumerate}
\def\labelenumi{(\alph{enumi})}
\setcounter{enumi}{1}
\tightlist
\item
  Converta \texttt{respostas} num fator \texttt{respostas\_fator}, com níveis ordenados do mais
  negativo para o mais positivo.
\item
  Verifique a estrutura com a função \texttt{str()}. O R reconhece a ordem?
\item
  Quantos clientes ficaram insatisfeitos ou muito insatisfeitos?
\item
  Qual a resposta mais frequente?
  \emph{Sugestão:} Use a função \texttt{table()} ou \texttt{which.max(table(...))}.
\item
  Calcule a frequência relativa (\%) de cada nível.
\item
  Substitua as respostas ``Neutro'' por ``Indiferente'', mantendo a ordem dos
  níveis.
\end{enumerate}

\textbf{9.} A Perturbação de Hiperactividade e Défice de Atenção (PHDA) é uma condição do neurodesenvolvimento que influencia a forma como o cérebro regula a atenção, a impulsividade e o nível de actividade. Estima-se que cerca de 5\% dos adultos vivam com PHDA, muitos deles sem diagnóstico durante a infância. No ensino superior, estudantes com PHDA referem frequentemente dificuldades de organização, concentração e gestão do tempo, embora também possam apresentar características como hiperfoco em áreas de interesse, criatividade e pensamento divergente.

\begin{enumerate}
\def\labelenumi{(\alph{enumi})}
\tightlist
\item
  Crie o vetor \texttt{phda\_desafios}:
\end{enumerate}

\begin{Shaded}
\begin{Highlighting}[]
\NormalTok{phda\_desafios }\OtherTok{\textless{}{-}} \FunctionTok{c}\NormalTok{(}\StringTok{"Concentração"}\NormalTok{, }\StringTok{"Procrastinação"}\NormalTok{, }\StringTok{"Gestão do tempo"}\NormalTok{, }\StringTok{"Organização"}\NormalTok{,}
  \StringTok{"Nenhum"}\NormalTok{, }\StringTok{"Concentração"}\NormalTok{, }\StringTok{"Concentração"}\NormalTok{, }\StringTok{"Organização"}\NormalTok{, }\StringTok{"Organização"}\NormalTok{, }\StringTok{"Organização"}\NormalTok{,}
  \StringTok{"Gestão do tempo"}\NormalTok{, }\StringTok{"Gestão do tempo"}\NormalTok{, }\StringTok{"Procrastinação"}\NormalTok{, }\StringTok{"Procrastinação"}\NormalTok{,}
  \StringTok{"Procrastinação"}\NormalTok{, }\StringTok{"Procrastinação"}\NormalTok{, }\StringTok{"Concentração"}\NormalTok{, }\StringTok{"Concentração"}\NormalTok{, }\StringTok{"Concentração"}\NormalTok{)}
\end{Highlighting}
\end{Shaded}

Converta-o num fator não ordenado \texttt{phda\_fator}.

\begin{enumerate}
\def\labelenumi{(\alph{enumi})}
\setcounter{enumi}{1}
\tightlist
\item
  Obtenha as frequências absolutas e relativas de cada desafio.
\item
  Quantos estudantes indicaram ``Concentração'' ou ``Procrastinação'' como principal
  dificuldade?
\item
  Recodifique ``Nenhum'' para ``Sem dificuldade declarada''.
\item
  Crie um gráfico de barras dos desafios, com rótulos de frequência.
  \emph{Sugestão:} Use a função \texttt{barplot()}.
\end{enumerate}

\textbf{10.} A dislexia é uma condição neurológica que afeta o processamento da
linguagem escrita, com desafios na leitura fluente, escrita e ortografia que não
se relacionam com o nível de inteligência. Estima-se que entre 5\% e 10\% da
população seja disléxica. Reconhecer a dislexia como parte da neurodiversidade é
essencial para ambientes de ensino inclusivos.

\begin{enumerate}
\def\labelenumi{(\alph{enumi})}
\tightlist
\item
  Crie o vetor \texttt{dificuldades} com as respostas de 10 crianças:
\end{enumerate}

\begin{Shaded}
\begin{Highlighting}[]
\NormalTok{dificuldades }\OtherTok{\textless{}{-}} \FunctionTok{c}\NormalTok{(}\StringTok{"Leitura"}\NormalTok{, }\StringTok{"Escrita"}\NormalTok{, }\StringTok{"Ortografia"}\NormalTok{, }\StringTok{"Leitura"}\NormalTok{, }\StringTok{"Compreensão"}\NormalTok{,}
  \StringTok{"Ortografia"}\NormalTok{, }\StringTok{"Memória auditiva"}\NormalTok{, }\StringTok{"Leitura"}\NormalTok{, }\StringTok{"Escrita"}\NormalTok{, }\StringTok{"Ortografia"}\NormalTok{)}
\end{Highlighting}
\end{Shaded}

Converta-o num fator não ordenado \texttt{dislexia\_fator}.

\begin{enumerate}
\def\labelenumi{(\alph{enumi})}
\setcounter{enumi}{1}
\tightlist
\item
  Crie uma tabela de frequências e identifique a dificuldade mais prevalente.
\item
  Recodifique ``Ortografia'' para ``Erros ortográficos'', mantendo o fator.
\item
  Visualize as frequências com \texttt{barplot()}.
\end{enumerate}

\section{Matriz e array}\label{matriz-e-array}

Os vetores organizam dados em \textbf{uma} dimensão --- uma linha de valores. Mas muitos
dados são naturalmente \textbf{bidimensionais}: uma tabela de números, os pixéis de
uma imagem, os coeficientes de um sistema de equações. Para estes, temos a
\textbf{matriz}.

Uma matriz é uma estrutura bidimensional cujos elementos são \textbf{todos do mesmo
tipo}, organizados em linhas e colunas. É, no fundo, um vetor a que se deu forma
retangular. O \textbf{array} generaliza a ideia a três ou mais dimensões.

\subsection{Criar matrizes}\label{criar-matrizes}

Criamos uma matriz com \texttt{matrix()}, indicando o número de linhas (\texttt{nrow}) e de
colunas (\texttt{ncol}):

\begin{Shaded}
\begin{Highlighting}[]
\FunctionTok{matrix}\NormalTok{(}\FunctionTok{c}\NormalTok{(}\DecValTok{1}\NormalTok{, }\DecValTok{2}\NormalTok{, }\DecValTok{3}\NormalTok{, }\DecValTok{4}\NormalTok{, }\DecValTok{5}\NormalTok{, }\DecValTok{6}\NormalTok{), }\AttributeTok{nrow =} \DecValTok{2}\NormalTok{, }\AttributeTok{ncol =} \DecValTok{3}\NormalTok{, }\AttributeTok{byrow =} \ConstantTok{FALSE}\NormalTok{)}
\DocumentationTok{\#\#      [,1] [,2] [,3]}
\DocumentationTok{\#\# [1,]    1    3    5}
\DocumentationTok{\#\# [2,]    2    4    6}
\end{Highlighting}
\end{Shaded}

Por omissão (\texttt{byrow\ =\ FALSE}), o R preenche a matriz \textbf{por colunas}: primeiro a
coluna 1 de cima a baixo, depois a coluna 2, e assim por diante. Se quiser
preencher por linhas, use \texttt{byrow\ =\ TRUE}. Esta é outra fonte frequente de
surpresa para quem começa --- se a matriz ``sai trocada'', verifique este argumento.

As operações vetorizadas também se aplicam a matrizes:

\begin{Shaded}
\begin{Highlighting}[]
\FunctionTok{matrix}\NormalTok{(}\FunctionTok{c}\NormalTok{(}\DecValTok{1}\NormalTok{, }\DecValTok{2}\NormalTok{, }\DecValTok{3}\NormalTok{, }\DecValTok{4}\NormalTok{, }\DecValTok{5}\NormalTok{, }\DecValTok{6}\NormalTok{), }\AttributeTok{nrow =} \DecValTok{2}\NormalTok{, }\AttributeTok{ncol =} \DecValTok{3}\NormalTok{) }\SpecialCharTok{\textgreater{}} \DecValTok{5}
\DocumentationTok{\#\#       [,1]  [,2]  [,3]}
\DocumentationTok{\#\# [1,] FALSE FALSE FALSE}
\DocumentationTok{\#\# [2,] FALSE FALSE  TRUE}
\end{Highlighting}
\end{Shaded}

\subsection{Construir matrizes a partir de vetores}\label{construir-matrizes-a-partir-de-vetores}

Frequentemente, temos vetores que queremos juntar como linhas ou colunas de uma
matriz. As funções \texttt{rbind()} (\emph{row bind}, juntar por linhas) e \texttt{cbind()} (\emph{column
bind}, juntar por colunas) fazem exatamente isso:

\begin{Shaded}
\begin{Highlighting}[]
\NormalTok{v1 }\OtherTok{\textless{}{-}} \FunctionTok{c}\NormalTok{(}\DecValTok{1}\NormalTok{, }\DecValTok{2}\NormalTok{, }\DecValTok{3}\NormalTok{)}
\NormalTok{v2 }\OtherTok{\textless{}{-}} \FunctionTok{c}\NormalTok{(}\DecValTok{4}\NormalTok{, }\DecValTok{5}\NormalTok{, }\DecValTok{6}\NormalTok{)}

\FunctionTok{rbind}\NormalTok{(v1, v2)     }\CommentTok{\# como linhas}
\DocumentationTok{\#\#    [,1] [,2] [,3]}
\DocumentationTok{\#\# v1    1    2    3}
\DocumentationTok{\#\# v2    4    5    6}

\FunctionTok{cbind}\NormalTok{(v1, v2)     }\CommentTok{\# como colunas}
\DocumentationTok{\#\#      v1 v2}
\DocumentationTok{\#\# [1,]  1  4}
\DocumentationTok{\#\# [2,]  2  5}
\DocumentationTok{\#\# [3,]  3  6}
\end{Highlighting}
\end{Shaded}

\subsection{Aceder a elementos de uma matriz}\label{aceder-a-elementos-de-uma-matriz}

A indexação estende naturalmente a dos vetores, agora com \textbf{dois} índices,
\texttt{{[}linha,\ coluna{]}}. Deixar um deles em branco significa ``todas'':

\begin{Shaded}
\begin{Highlighting}[]
\NormalTok{A }\OtherTok{\textless{}{-}} \FunctionTok{matrix}\NormalTok{((}\SpecialCharTok{{-}}\DecValTok{4}\NormalTok{)}\SpecialCharTok{:}\DecValTok{5}\NormalTok{, }\AttributeTok{nrow =} \DecValTok{2}\NormalTok{, }\AttributeTok{ncol =} \DecValTok{5}\NormalTok{)}
\NormalTok{A}
\DocumentationTok{\#\#      [,1] [,2] [,3] [,4] [,5]}
\DocumentationTok{\#\# [1,]   {-}4   {-}2    0    2    4}
\DocumentationTok{\#\# [2,]   {-}3   {-}1    1    3    5}

\NormalTok{A[}\DecValTok{1}\NormalTok{, }\DecValTok{2}\NormalTok{]      }\CommentTok{\# elemento na linha 1, coluna 2}
\DocumentationTok{\#\# [1] {-}2}

\NormalTok{A[A }\SpecialCharTok{\textless{}} \DecValTok{0}\NormalTok{]     }\CommentTok{\# todos os elementos negativos (indexação lógica)}
\DocumentationTok{\#\# [1] {-}4 {-}3 {-}2 {-}1}

\NormalTok{A[A }\SpecialCharTok{\textless{}} \DecValTok{0}\NormalTok{] }\OtherTok{\textless{}{-}} \DecValTok{0}    \CommentTok{\# substituir os negativos por 0}
\NormalTok{A}
\DocumentationTok{\#\#      [,1] [,2] [,3] [,4] [,5]}
\DocumentationTok{\#\# [1,]    0    0    0    2    4}
\DocumentationTok{\#\# [2,]    0    0    1    3    5}

\NormalTok{A[}\DecValTok{2}\NormalTok{, ]       }\CommentTok{\# toda a linha 2}
\DocumentationTok{\#\# [1] 0 0 1 3 5}

\NormalTok{A[, }\FunctionTok{c}\NormalTok{(}\DecValTok{2}\NormalTok{, }\DecValTok{4}\NormalTok{)] }\CommentTok{\# as colunas 2 e 4}
\DocumentationTok{\#\#      [,1] [,2]}
\DocumentationTok{\#\# [1,]    0    2}
\DocumentationTok{\#\# [2,]    0    3}
\end{Highlighting}
\end{Shaded}

\subsection{Nomear linhas e colunas}\label{nomear-linhas-e-colunas}

Tal como os vetores, as matrizes podem ter nomes nas linhas e colunas, o que
torna os resultados muito mais legíveis:

\begin{Shaded}
\begin{Highlighting}[]
\NormalTok{x }\OtherTok{\textless{}{-}} \FunctionTok{matrix}\NormalTok{(}\FunctionTok{rnorm}\NormalTok{(}\DecValTok{12}\NormalTok{), }\AttributeTok{nrow =} \DecValTok{4}\NormalTok{)}
\FunctionTok{colnames}\NormalTok{(x) }\OtherTok{\textless{}{-}} \FunctionTok{paste0}\NormalTok{(}\StringTok{"dados"}\NormalTok{, }\DecValTok{1}\SpecialCharTok{:}\DecValTok{3}\NormalTok{)}
\FunctionTok{rownames}\NormalTok{(x) }\OtherTok{\textless{}{-}}\NormalTok{ letters[}\DecValTok{1}\SpecialCharTok{:}\DecValTok{4}\NormalTok{]}
\end{Highlighting}
\end{Shaded}

\subsection{Álgebra de matrizes}\label{uxe1lgebra-de-matrizes}

Aqui reside uma das grandes vantagens do R para a estatística: a álgebra linear é
imediata. A \textbf{multiplicação de matrizes} faz-se com o operador \texttt{\%*\%} (repare que
\texttt{*} faria a multiplicação elemento a elemento, que é outra coisa):

\begin{Shaded}
\begin{Highlighting}[]
\NormalTok{M }\OtherTok{\textless{}{-}} \FunctionTok{matrix}\NormalTok{(}\FunctionTok{rnorm}\NormalTok{(}\DecValTok{20}\NormalTok{), }\AttributeTok{nrow =} \DecValTok{4}\NormalTok{, }\AttributeTok{ncol =} \DecValTok{5}\NormalTok{)}
\NormalTok{N }\OtherTok{\textless{}{-}} \FunctionTok{matrix}\NormalTok{(}\FunctionTok{rnorm}\NormalTok{(}\DecValTok{15}\NormalTok{), }\AttributeTok{nrow =} \DecValTok{5}\NormalTok{, }\AttributeTok{ncol =} \DecValTok{3}\NormalTok{)}

\NormalTok{M }\SpecialCharTok{\%*\%}\NormalTok{ N        }\CommentTok{\# produto matricial (4x5 por 5x3 = 4x3)}
\end{Highlighting}
\end{Shaded}

E as operações clássicas da álgebra linear têm, cada uma, a sua função. Sendo \texttt{M}
uma matriz quadrada:

{\def\LTcaptype{none} % do not increment counter
\begin{longtable}[]{@{}ll@{}}
\toprule\noalign{}
Operação & Função \\
\midrule\noalign{}
\endhead
\bottomrule\noalign{}
\endlastfoot
Dimensão & \texttt{dim(M)} \\
Transposta & \texttt{t(M)} \\
Diagonal principal & \texttt{diag(M)} \\
Determinante & \texttt{det(M)} \\
Inversa & \texttt{solve(M)} \\
Valores e vetores próprios & \texttt{eigen(M)} \\
Soma de todos os elementos & \texttt{sum(M)} \\
Média de todos os elementos & \texttt{mean(M)} \\
\end{longtable}
}

Há ainda a poderosa função \texttt{apply()}, que aplica uma função a cada linha ou a cada
coluna. O segundo argumento indica a margem: \texttt{1} para linhas, \texttt{2} para colunas:

\begin{Shaded}
\begin{Highlighting}[]
\FunctionTok{apply}\NormalTok{(M, }\DecValTok{1}\NormalTok{, sum)    }\CommentTok{\# soma de cada linha}
\FunctionTok{apply}\NormalTok{(M, }\DecValTok{2}\NormalTok{, mean)   }\CommentTok{\# média de cada coluna}
\end{Highlighting}
\end{Shaded}

\subsection{Adicionar linhas e colunas}\label{adicionar-linhas-e-colunas}

\texttt{rbind()} e \texttt{cbind()} também servem para acrescentar linhas e colunas a uma matriz
já existente:

\begin{Shaded}
\begin{Highlighting}[]
\NormalTok{X }\OtherTok{\textless{}{-}} \FunctionTok{matrix}\NormalTok{(}\FunctionTok{c}\NormalTok{(}\DecValTok{1}\NormalTok{, }\DecValTok{2}\NormalTok{, }\DecValTok{3}\NormalTok{, }\DecValTok{4}\NormalTok{, }\DecValTok{5}\NormalTok{, }\DecValTok{6}\NormalTok{), }\AttributeTok{nrow =} \DecValTok{2}\NormalTok{)}

\FunctionTok{cbind}\NormalTok{(X, }\FunctionTok{c}\NormalTok{(}\DecValTok{7}\NormalTok{, }\DecValTok{8}\NormalTok{))          }\CommentTok{\# nova coluna à direita}
\FunctionTok{rbind}\NormalTok{(X, }\FunctionTok{c}\NormalTok{(}\DecValTok{7}\NormalTok{, }\DecValTok{8}\NormalTok{, }\DecValTok{9}\NormalTok{))       }\CommentTok{\# nova linha em baixo}
\end{Highlighting}
\end{Shaded}

\subsection{Arrays}\label{arrays}

Quando duas dimensões não bastam, generalizamos para o \textbf{array}. Um array
tridimensional pode imaginar-se como um ``empilhamento'' de matrizes (camadas).
Criamo-lo com \texttt{array()}, indicando as dimensões em \texttt{dim}:

\begin{Shaded}
\begin{Highlighting}[]
\CommentTok{\# 4 linhas, 3 colunas, 2 camadas}
\NormalTok{A }\OtherTok{\textless{}{-}} \FunctionTok{array}\NormalTok{(}\DecValTok{1}\SpecialCharTok{:}\DecValTok{24}\NormalTok{, }\AttributeTok{dim =} \FunctionTok{c}\NormalTok{(}\DecValTok{4}\NormalTok{, }\DecValTok{3}\NormalTok{, }\DecValTok{2}\NormalTok{))}
\end{Highlighting}
\end{Shaded}

A indexação usa agora três índices, \texttt{{[}linha,\ coluna,\ camada{]}}:

\begin{Shaded}
\begin{Highlighting}[]
\NormalTok{A[}\DecValTok{1}\NormalTok{, }\DecValTok{2}\NormalTok{, }\DecValTok{1}\NormalTok{]     }\CommentTok{\# elemento na linha 1, coluna 2, camada 1}
\DocumentationTok{\#\# [1] 5}

\NormalTok{A[}\DecValTok{2}\NormalTok{, , }\DecValTok{1}\NormalTok{]      }\CommentTok{\# toda a linha 2 da camada 1}
\DocumentationTok{\#\# [1]  2  6 10}

\NormalTok{A[, , }\DecValTok{1}\NormalTok{]       }\CommentTok{\# toda a primeira camada (uma matriz)}
\end{Highlighting}
\end{Shaded}

\subsection{Exercícios}\label{exercuxedcios-3}

\textbf{1.} Crie uma matriz \(3 \times 4\) com os números de 1 a 12, preenchida por
colunas. Exiba-a e determine a soma dos elementos da segunda coluna.

\textbf{2.} Crie duas matrizes \(2 \times 3\), \(A\) e \(B\), com inteiros aleatórios de 1 a
10. Some-as e multiplique-as elemento a elemento.
\emph{Sugestão:} Use a função \texttt{sample()}.

\textbf{3.} Crie uma matriz \(3 \times 3\), \(M\), com números aleatórios entre 1 e 9, e um
vetor \(v\) de comprimento 3. Calcule \(M \times v\) e depois \(v \times M\).

\textbf{4.} Crie uma matriz \(4 \times 2\), \(N\), com valores de 1 a 8. Transponha-a e
calcule a soma dos elementos de cada linha da transposta. Some depois a primeira e
a segunda linha da transposta.

\textbf{5.} Crie uma matriz \(3 \times 3\), \(P\), à sua escolha. Calcule o seu
determinante e, se for não nulo, a sua inversa.

\textbf{6.} Crie uma matriz \(5 \times 5\), \(Q\), com números aleatórios de 1 a 25.
Extraia a submatriz das 2.ª, 3.ª e 4.ª linhas e colunas.

\textbf{7.} Crie uma matriz \(3 \times 3\) com valores de 1 a 9. Calcule a soma da
primeira linha e a soma da terceira coluna.

\textbf{8.} Considere a matriz
\[A = \begin{bmatrix} 1 & 1 & 3 \\ 5 & 2 & 6 \\ -2 & -1 & -3 \end{bmatrix}.\]

\begin{enumerate}
\def\labelenumi{(\alph{enumi})}
\tightlist
\item
  Verifique que \(A^3 = 0\).
\item
  Substitua a terceira coluna pela soma das colunas 1 e 3.
\item
  Na matriz
  \[M = \begin{bmatrix} 20 & 22 & 23 \\ 34 & 55 & 57 \\ 99 & 97 & 71 \\
  12 & 16 & 19 \\ 10 & 53 & 24 \\ 14 & 21 & 28 \end{bmatrix},\]
  substitua por 0 todos os números pares.
\end{enumerate}

\textbf{9.} Crie uma matriz \(3 \times 3\) e verifique se é simétrica (igual à
transposta). Teste com
\[A = \begin{bmatrix} 1 & 2 & 3 \\ 2 & 4 & 5 \\ 3 & 5 & 6 \end{bmatrix}.\]
\emph{Sugestão:} Use a função \texttt{all()}.

\textbf{10.} Crie uma matriz \(4 \times 4\) e calcule o seu traço (a soma da diagonal
principal).

\textbf{11.} Crie uma matriz \(5 \times 5\) com valores de 1 a 25 e extraia as colunas de
índice par (2 e 4).

\textbf{12.} Uma turma usou o ChatGPT como apoio ao estudo. Para cada estudante
registou-se o número de perguntas feitas numa semana, o tempo médio de resposta
(segundos) e o nível de satisfação (1 a 5). Os dados de 3 alunos:
\[M = \begin{bmatrix} 12 & 8 & 5 \\ 20 & 10 & 4 \\ 15 & 9 & 5 \end{bmatrix},\]
com as linhas a corresponder aos alunos A, B e C, e as colunas a (Perguntas,
Tempo, Satisfação).

\begin{enumerate}
\def\labelenumi{(\alph{enumi})}
\tightlist
\item
  Crie a matriz \texttt{M} e nomeie linhas e colunas.
\item
  Use \texttt{rbind()} para acrescentar o aluno D: 18 perguntas, 11 segundos,
  satisfação 3.
\item
  Use \texttt{cbind()} para acrescentar a coluna ``Horas de uso'' com os valores
  \texttt{c(5,\ 7,\ 6,\ 8)}.
\item
  Calcule a média de cada métrica com \texttt{colMeans()}.
\item
  Que aluno fez mais perguntas?
  \emph{Sugestão:} Use a função \texttt{which.max()}.
\end{enumerate}

\textbf{13.} Um ecrã em tons de cinzento pode representar-se por uma matriz, em que
cada entrada é a intensidade de um pixel (0 = preto, 9 = branco). Considere a
imagem \(10 \times 10\):

\begin{enumerate}
\def\labelenumi{(\alph{enumi})}
\tightlist
\item
  Crie a matriz \texttt{I}:
\end{enumerate}

\begin{Shaded}
\begin{Highlighting}[]
\NormalTok{I }\OtherTok{\textless{}{-}} \FunctionTok{matrix}\NormalTok{(}\FunctionTok{c}\NormalTok{(}
  \DecValTok{0}\NormalTok{,}\DecValTok{0}\NormalTok{,}\DecValTok{0}\NormalTok{,}\DecValTok{0}\NormalTok{,}\DecValTok{0}\NormalTok{,}\DecValTok{0}\NormalTok{,}\DecValTok{0}\NormalTok{,}\DecValTok{0}\NormalTok{,}\DecValTok{0}\NormalTok{,}\DecValTok{0}\NormalTok{,}
  \DecValTok{0}\NormalTok{,}\DecValTok{9}\NormalTok{,}\DecValTok{9}\NormalTok{,}\DecValTok{9}\NormalTok{,}\DecValTok{0}\NormalTok{,}\DecValTok{0}\NormalTok{,}\DecValTok{9}\NormalTok{,}\DecValTok{9}\NormalTok{,}\DecValTok{9}\NormalTok{,}\DecValTok{0}\NormalTok{,}
  \DecValTok{0}\NormalTok{,}\DecValTok{9}\NormalTok{,}\DecValTok{0}\NormalTok{,}\DecValTok{9}\NormalTok{,}\DecValTok{0}\NormalTok{,}\DecValTok{0}\NormalTok{,}\DecValTok{9}\NormalTok{,}\DecValTok{0}\NormalTok{,}\DecValTok{9}\NormalTok{,}\DecValTok{0}\NormalTok{,}
  \DecValTok{0}\NormalTok{,}\DecValTok{9}\NormalTok{,}\DecValTok{9}\NormalTok{,}\DecValTok{9}\NormalTok{,}\DecValTok{0}\NormalTok{,}\DecValTok{0}\NormalTok{,}\DecValTok{9}\NormalTok{,}\DecValTok{9}\NormalTok{,}\DecValTok{9}\NormalTok{,}\DecValTok{0}\NormalTok{,}
  \DecValTok{0}\NormalTok{,}\DecValTok{9}\NormalTok{,}\DecValTok{0}\NormalTok{,}\DecValTok{0}\NormalTok{,}\DecValTok{0}\NormalTok{,}\DecValTok{0}\NormalTok{,}\DecValTok{0}\NormalTok{,}\DecValTok{0}\NormalTok{,}\DecValTok{9}\NormalTok{,}\DecValTok{0}\NormalTok{,}
  \DecValTok{0}\NormalTok{,}\DecValTok{9}\NormalTok{,}\DecValTok{0}\NormalTok{,}\DecValTok{0}\NormalTok{,}\DecValTok{0}\NormalTok{,}\DecValTok{0}\NormalTok{,}\DecValTok{0}\NormalTok{,}\DecValTok{0}\NormalTok{,}\DecValTok{9}\NormalTok{,}\DecValTok{0}\NormalTok{,}
  \DecValTok{0}\NormalTok{,}\DecValTok{9}\NormalTok{,}\DecValTok{0}\NormalTok{,}\DecValTok{0}\NormalTok{,}\DecValTok{0}\NormalTok{,}\DecValTok{0}\NormalTok{,}\DecValTok{0}\NormalTok{,}\DecValTok{0}\NormalTok{,}\DecValTok{9}\NormalTok{,}\DecValTok{0}\NormalTok{,}
  \DecValTok{0}\NormalTok{,}\DecValTok{9}\NormalTok{,}\DecValTok{0}\NormalTok{,}\DecValTok{0}\NormalTok{,}\DecValTok{0}\NormalTok{,}\DecValTok{0}\NormalTok{,}\DecValTok{0}\NormalTok{,}\DecValTok{0}\NormalTok{,}\DecValTok{9}\NormalTok{,}\DecValTok{0}\NormalTok{,}
  \DecValTok{0}\NormalTok{,}\DecValTok{9}\NormalTok{,}\DecValTok{9}\NormalTok{,}\DecValTok{9}\NormalTok{,}\DecValTok{9}\NormalTok{,}\DecValTok{9}\NormalTok{,}\DecValTok{9}\NormalTok{,}\DecValTok{9}\NormalTok{,}\DecValTok{9}\NormalTok{,}\DecValTok{0}\NormalTok{,}
  \DecValTok{0}\NormalTok{,}\DecValTok{0}\NormalTok{,}\DecValTok{0}\NormalTok{,}\DecValTok{0}\NormalTok{,}\DecValTok{0}\NormalTok{,}\DecValTok{0}\NormalTok{,}\DecValTok{0}\NormalTok{,}\DecValTok{0}\NormalTok{,}\DecValTok{0}\NormalTok{,}\DecValTok{0}
\NormalTok{), }\AttributeTok{nrow =} \DecValTok{10}\NormalTok{, }\AttributeTok{byrow =} \ConstantTok{TRUE}\NormalTok{)}
\end{Highlighting}
\end{Shaded}

\begin{enumerate}
\def\labelenumi{(\alph{enumi})}
\setcounter{enumi}{1}
\tightlist
\item
  Visualize-a como imagem com \texttt{image()}.
  \emph{Sugestão:} Use \texttt{col\ =\ gray.colors(10)}.
\item
  Inverta as cores (substitua cada valor \(x\) por \(9 - x\)).
\item
  Calcule a intensidade média de todos os pixéis (o brilho médio da imagem).
\item
  Acrescente uma linha de zeros no fim da matriz.
\end{enumerate}

\textbf{14.} Numa região há Humanos (H) e Zumbis (Z). A cada hora, 20\% dos humanos
tornam-se zumbis e 10\% dos zumbis voltam a humanos. Organizamos estas
probabilidades numa \textbf{matriz de transição} \(P\), e se \(\mathbf{x}_t\) é o vetor com
o número de humanos e zumbis no instante \(t\), então
\(\mathbf{x}_{t+1} = P \cdot \mathbf{x}_t\). Inicialmente há 150 humanos e 150
zumbis.

\begin{enumerate}
\def\labelenumi{(\alph{enumi})}
\tightlist
\item
  Crie a matriz \(P\) e o vetor inicial \(\mathbf{x}_0 = (150, 150)^\top\).
\item
  Calcule \(\mathbf{x}_1\), \(\mathbf{x}_2\) e \(\mathbf{x}_{10}\). Quantos zumbis há
  à 10.ª hora?
  \emph{Sugestão:} use \texttt{\%\^{}\%} (do pacote \texttt{expm}) para potências de
  matrizes.
\end{enumerate}

Este é, na verdade, um exemplo simples de \textbf{cadeia de Markov} --- um sistema que
evolui em passos sucessivos, em que o estado seguinte depende apenas do estado
atual e de uma matriz de transição. Reencontraremos esta ideia mais adiante, no
estudo da simulação.

\textbf{15.} Uma forma de cifrar mensagens é através da multiplicação por matrizes.
Associando as letras a números (sem a letra K, com espaço = 0), a mensagem
converte-se numa matriz \(M\), cifra-se multiplicando por uma \textbf{matriz-chave}
invertível \(C\) (\(M_{\text{cod}} = M \cdot C\)), e decifra-se multiplicando pela
inversa. Com a chave
\[C = \begin{bmatrix} 1 & 0 & 1 \\ -1 & 3 & 1 \\ 0 & 1 & 1 \end{bmatrix}:\]

\begin{enumerate}
\def\labelenumi{(\alph{enumi})}
\tightlist
\item
  Recebeu-se (em blocos \(3 \times 3\), lidos por linhas): \(-12, 48, 23, -2, 42,
  26, 1, 42, 29\). Decifre a mensagem.
\item
  Escolha uma matriz-chave para palavras até 16 letras e cifre e decifre à
  vontade.
\end{enumerate}

\textbf{16.} Crie um array \(3 \times 3 \times 3\) com valores de 1 a 27. Aceda ao
elemento \texttt{{[}2,\ 3,\ 1{]}}, à segunda camada e à terceira linha da segunda camada.

\textbf{17.} Crie dois arrays \(2 \times 2 \times 2\) com valores aleatórios de 1 a 10.
Some-os, multiplique-os elemento a elemento, calcule a média da soma e multiplique
a primeira matriz do primeiro array pela segunda matriz do segundo.

\section{Data frame}\label{data-frame}

Chegamos, talvez, à estrutura mais importante de todas para o trabalho
estatístico do dia a dia. As matrizes obrigam a que todos os elementos sejam do
mesmo tipo --- mas os dados reais quase nunca são assim. Uma tabela de estudantes
tem nomes (texto), idades (números inteiros), notas (números decimais) e, talvez,
uma indicação de aprovação (lógica). Precisamos de uma estrutura que combine
colunas de tipos diferentes. Essa estrutura é o \textbf{data frame}.

Um data frame é uma tabela bidimensional em que \textbf{cada coluna pode ter o seu
próprio tipo}, embora dentro de cada coluna todos os elementos sejam do mesmo
tipo. É a estrutura mais próxima de uma folha de cálculo ou de uma tabela de base
de dados --- e é o formato em que a esmagadora maioria dos conjuntos de dados chega
até nós.

\subsection{Criar um data frame}\label{criar-um-data-frame}

Criamo-lo com \texttt{data.frame()}, indicando cada coluna pelo seu nome:

\begin{Shaded}
\begin{Highlighting}[]
\NormalTok{df }\OtherTok{\textless{}{-}} \FunctionTok{data.frame}\NormalTok{(}
  \AttributeTok{id =} \DecValTok{1}\SpecialCharTok{:}\DecValTok{4}\NormalTok{,}
  \AttributeTok{nome =} \FunctionTok{c}\NormalTok{(}\StringTok{"Ana"}\NormalTok{, }\StringTok{"Bruno"}\NormalTok{, }\StringTok{"Carlos"}\NormalTok{, }\StringTok{"Diana"}\NormalTok{),}
  \AttributeTok{idade =} \FunctionTok{c}\NormalTok{(}\DecValTok{23}\NormalTok{, }\DecValTok{35}\NormalTok{, }\DecValTok{31}\NormalTok{, }\DecValTok{28}\NormalTok{),}
  \AttributeTok{salario =} \FunctionTok{c}\NormalTok{(}\DecValTok{5000}\NormalTok{, }\DecValTok{6000}\NormalTok{, }\DecValTok{7000}\NormalTok{, }\DecValTok{8000}\NormalTok{)}
\NormalTok{)}
\NormalTok{df}
\DocumentationTok{\#\#   id   nome idade salario}
\DocumentationTok{\#\# 1  1    Ana    23    5000}
\DocumentationTok{\#\# 2  2  Bruno    35    6000}
\DocumentationTok{\#\# 3  3 Carlos    31    7000}
\DocumentationTok{\#\# 4  4  Diana    28    8000}
\end{Highlighting}
\end{Shaded}

Para perceber porque não basta uma matriz, tentemos criar a ``mesma'' tabela com
\texttt{cbind()}:

\begin{Shaded}
\begin{Highlighting}[]
\FunctionTok{cbind}\NormalTok{(}\AttributeTok{id =} \DecValTok{1}\SpecialCharTok{:}\DecValTok{4}\NormalTok{, }\AttributeTok{nome =} \FunctionTok{c}\NormalTok{(}\StringTok{"Ana"}\NormalTok{, }\StringTok{"Bruno"}\NormalTok{, }\StringTok{"Carlos"}\NormalTok{, }\StringTok{"Diana"}\NormalTok{),}
      \AttributeTok{idade =} \FunctionTok{c}\NormalTok{(}\DecValTok{23}\NormalTok{, }\DecValTok{35}\NormalTok{, }\DecValTok{31}\NormalTok{, }\DecValTok{28}\NormalTok{))}
\DocumentationTok{\#\#      id  nome     idade }
\DocumentationTok{\#\# [1,] "1" "Ana"    "23"  }
\DocumentationTok{\#\# [2,] "2" "Bruno"  "35"  }
\DocumentationTok{\#\# ...}
\end{Highlighting}
\end{Shaded}

Repare que a matriz converteu \textbf{tudo} para texto --- as idades ficaram \texttt{"23"},
\texttt{"35"}, entre aspas, e deixaram de ser números com os quais podemos calcular. É a
coerção em ação, e é exatamente o problema que o data frame resolve: nele, cada
coluna preserva o seu tipo.

\subsection{Aceder a colunas}\label{aceder-a-colunas}

Um data frame oferece várias formas de aceder às suas colunas, e vale a pena
conhecê-las porque devolvem coisas subtilmente diferentes. A forma mais legível é
o operador \texttt{\$}:

\begin{Shaded}
\begin{Highlighting}[]
\NormalTok{df}\SpecialCharTok{$}\NormalTok{id           }\CommentTok{\# a coluna \textquotesingle{}id\textquotesingle{} como vetor}
\DocumentationTok{\#\# [1] 1 2 3 4}

\NormalTok{df[[}\StringTok{"id"}\NormalTok{]]      }\CommentTok{\# equivalente, útil quando o nome está numa variável}
\DocumentationTok{\#\# [1] 1 2 3 4}

\NormalTok{df[, }\DecValTok{1}\NormalTok{]         }\CommentTok{\# por posição, também devolve um vetor}
\DocumentationTok{\#\# [1] 1 2 3 4}

\NormalTok{df[}\StringTok{"id"}\NormalTok{]        }\CommentTok{\# atenção: devolve um data frame de uma coluna, não um vetor!}
\DocumentationTok{\#\#   id}
\DocumentationTok{\#\# 1  1}
\DocumentationTok{\#\# 2  2}
\DocumentationTok{\#\# 3  3}
\DocumentationTok{\#\# 4  4}
\end{Highlighting}
\end{Shaded}

A distinção entre ``devolver um vetor'' e ``devolver um data frame de uma coluna'' é
uma subtileza que confunde muitos iniciantes, mas que tem uma regra simples: os
colchetes duplos \texttt{{[}{[}\ {]}{]}} e o \texttt{\$} extraem o \textbf{conteúdo} (um vetor); os colchetes
simples \texttt{{[}\ {]}} preservam a \textbf{estrutura} (um data frame).

\subsection{Aceder a linhas}\label{aceder-a-linhas}

As linhas acedem-se com a notação \texttt{{[}linha,\ coluna{]}}, tal como nas matrizes,
deixando em branco a dimensão que queremos por inteiro:

\begin{Shaded}
\begin{Highlighting}[]
\NormalTok{df[}\DecValTok{1}\NormalTok{, ]         }\CommentTok{\# a primeira linha (todas as colunas)}
\NormalTok{df[}\DecValTok{1}\SpecialCharTok{:}\DecValTok{2}\NormalTok{, ]       }\CommentTok{\# as duas primeiras linhas}
\NormalTok{df[}\FunctionTok{c}\NormalTok{(}\DecValTok{1}\NormalTok{, }\DecValTok{4}\NormalTok{), ]   }\CommentTok{\# as linhas 1 e 4}
\NormalTok{df[}\SpecialCharTok{{-}}\DecValTok{1}\NormalTok{, ]        }\CommentTok{\# todas menos a primeira}
\NormalTok{df[}\DecValTok{1}\NormalTok{, }\StringTok{"nome"}\NormalTok{]   }\CommentTok{\# o valor na linha 1, coluna \textquotesingle{}nome\textquotesingle{}}
\DocumentationTok{\#\# [1] "Ana"}
\end{Highlighting}
\end{Shaded}

\subsection{Adicionar e remover colunas}\label{adicionar-e-remover-colunas}

Acrescentar uma coluna é tão simples como atribuir-lhe um valor com \texttt{\$}; muitas
vezes, a nova coluna é calculada a partir das existentes:

\begin{Shaded}
\begin{Highlighting}[]
\CommentTok{\# Nova coluna calculada (aumento de 10\% no salário)}
\NormalTok{df}\SpecialCharTok{$}\NormalTok{novo\_salario }\OtherTok{\textless{}{-}}\NormalTok{ df}\SpecialCharTok{$}\NormalTok{salario }\SpecialCharTok{*} \FloatTok{1.1}

\CommentTok{\# Remover uma coluna: atribuir{-}lhe NULL}
\NormalTok{df}\SpecialCharTok{$}\NormalTok{id }\OtherTok{\textless{}{-}} \ConstantTok{NULL}
\end{Highlighting}
\end{Shaded}

\subsection{Fundir data frames}\label{fundir-data-frames}

Trabalhar com dados reais implica muitas vezes juntar tabelas. Há três situações
típicas, cada uma com a sua ferramenta.

\textbf{Acrescentar observações (linhas)} com \texttt{rbind()}, quando as tabelas têm as
mesmas colunas:

\begin{Shaded}
\begin{Highlighting}[]
\NormalTok{df1 }\OtherTok{\textless{}{-}} \FunctionTok{data.frame}\NormalTok{(}\AttributeTok{id =} \DecValTok{1}\SpecialCharTok{:}\DecValTok{2}\NormalTok{, }\AttributeTok{nome =} \FunctionTok{c}\NormalTok{(}\StringTok{"Ana"}\NormalTok{, }\StringTok{"Bruno"}\NormalTok{))}
\NormalTok{df2 }\OtherTok{\textless{}{-}} \FunctionTok{data.frame}\NormalTok{(}\AttributeTok{id =} \DecValTok{3}\SpecialCharTok{:}\DecValTok{4}\NormalTok{, }\AttributeTok{nome =} \FunctionTok{c}\NormalTok{(}\StringTok{"Carlos"}\NormalTok{, }\StringTok{"Diana"}\NormalTok{))}
\FunctionTok{rbind}\NormalTok{(df1, df2)}
\end{Highlighting}
\end{Shaded}

\textbf{Acrescentar variáveis (colunas)} com \texttt{cbind()}, quando as tabelas têm o mesmo
número de linhas, alinhadas pela ordem.

\textbf{Cruzar tabelas por uma chave comum} com \texttt{merge()} --- a operação mais poderosa e
a mais próxima do que se faz numa base de dados. Cruzamos duas tabelas por uma
variável partilhada (como \texttt{id} ou \texttt{curso}):

\begin{Shaded}
\begin{Highlighting}[]
\NormalTok{df1 }\OtherTok{\textless{}{-}} \FunctionTok{data.frame}\NormalTok{(}\AttributeTok{id =} \DecValTok{1}\SpecialCharTok{:}\DecValTok{3}\NormalTok{, }\AttributeTok{nome =} \FunctionTok{c}\NormalTok{(}\StringTok{"Ana"}\NormalTok{, }\StringTok{"Bruno"}\NormalTok{, }\StringTok{"Carlos"}\NormalTok{))}
\NormalTok{df2 }\OtherTok{\textless{}{-}} \FunctionTok{data.frame}\NormalTok{(}\AttributeTok{id =} \DecValTok{2}\SpecialCharTok{:}\DecValTok{4}\NormalTok{, }\AttributeTok{idade =} \FunctionTok{c}\NormalTok{(}\DecValTok{35}\NormalTok{, }\DecValTok{40}\NormalTok{, }\DecValTok{28}\NormalTok{))}

\FunctionTok{merge}\NormalTok{(df1, df2, }\AttributeTok{by =} \StringTok{"id"}\NormalTok{)                 }\CommentTok{\# só as chaves em comum (inner join)}
\DocumentationTok{\#\#   id   nome idade}
\DocumentationTok{\#\# 1  2  Bruno    35}
\DocumentationTok{\#\# 2  3 Carlos    40}

\FunctionTok{merge}\NormalTok{(df1, df2, }\AttributeTok{by =} \StringTok{"id"}\NormalTok{, }\AttributeTok{all.x =} \ConstantTok{TRUE}\NormalTok{)   }\CommentTok{\# todas as linhas de df1 (left join)}
\FunctionTok{merge}\NormalTok{(df1, df2, }\AttributeTok{by =} \StringTok{"id"}\NormalTok{, }\AttributeTok{all =} \ConstantTok{TRUE}\NormalTok{)     }\CommentTok{\# todas as linhas de ambos (full join)}
\end{Highlighting}
\end{Shaded}

Estes cruzamentos --- \emph{inner}, \emph{left}, \emph{right} e \emph{full join} --- são conceitos
centrais em bases de dados. O pacote \texttt{dplyr} oferece versões mais legíveis com os
mesmos nomes (\texttt{inner\_join()}, \texttt{left\_join()}, etc.), que reencontraremos no
capítulo de manipulação de dados.

\subsection{Explorar um data frame}\label{explorar-um-data-frame}

Perante um conjunto de dados novo, o primeiro passo é sempre \emph{conhecê-lo}. Um
punhado de funções permite fazê-lo rapidamente. Usemos como exemplo o \texttt{iris}, um
famoso conjunto de dados incluído no R:

\begin{Shaded}
\begin{Highlighting}[]
\NormalTok{df }\OtherTok{\textless{}{-}}\NormalTok{ iris}

\FunctionTok{names}\NormalTok{(df)     }\CommentTok{\# nomes das colunas}
\FunctionTok{dim}\NormalTok{(df)       }\CommentTok{\# dimensões (linhas e colunas)}
\DocumentationTok{\#\# [1] 150   5}
\FunctionTok{nrow}\NormalTok{(df); }\FunctionTok{ncol}\NormalTok{(df)}
\FunctionTok{str}\NormalTok{(df)       }\CommentTok{\# a estrutura: tipos e primeiros valores de cada coluna}
\FunctionTok{head}\NormalTok{(df, }\DecValTok{3}\NormalTok{)   }\CommentTok{\# as primeiras linhas}
\FunctionTok{tail}\NormalTok{(df, }\DecValTok{3}\NormalTok{)   }\CommentTok{\# as últimas linhas}
\end{Highlighting}
\end{Shaded}

Faça de \texttt{str()} e \texttt{head()} um reflexo. Sempre que carrega um conjunto de dados,
comece por aí: \texttt{str()} diz-lhe quantas observações e variáveis tem e de que tipo
é cada uma; \texttt{head()} mostra-lhe como se parecem os dados. Estes dois olhares
poupam horas de confusão mais à frente.

\subsection{Filtrar dados}\label{filtrar-dados}

Uma das tarefas mais comuns é selecionar as linhas que satisfazem uma condição.
A função \texttt{subset()} fá-lo de forma clara, permitindo filtrar linhas e escolher
colunas ao mesmo tempo:

\begin{Shaded}
\begin{Highlighting}[]
\NormalTok{df }\OtherTok{\textless{}{-}} \FunctionTok{data.frame}\NormalTok{(}
  \AttributeTok{id =} \DecValTok{1}\SpecialCharTok{:}\DecValTok{4}\NormalTok{,}
  \AttributeTok{nome =} \FunctionTok{c}\NormalTok{(}\StringTok{"Ana"}\NormalTok{, }\StringTok{"Bruno"}\NormalTok{, }\StringTok{"Carlos"}\NormalTok{, }\StringTok{"Diana"}\NormalTok{),}
  \AttributeTok{idade =} \FunctionTok{c}\NormalTok{(}\DecValTok{23}\NormalTok{, }\DecValTok{35}\NormalTok{, }\DecValTok{31}\NormalTok{, }\DecValTok{28}\NormalTok{),}
  \AttributeTok{salario =} \FunctionTok{c}\NormalTok{(}\DecValTok{5000}\NormalTok{, }\DecValTok{6000}\NormalTok{, }\DecValTok{7000}\NormalTok{, }\DecValTok{8000}\NormalTok{)}
\NormalTok{)}

\FunctionTok{subset}\NormalTok{(df, idade }\SpecialCharTok{\textgreater{}} \DecValTok{30}\NormalTok{)                             }\CommentTok{\# linhas com idade \textgreater{} 30}
\FunctionTok{subset}\NormalTok{(df, idade }\SpecialCharTok{\textgreater{}} \DecValTok{30}\NormalTok{, }\AttributeTok{select =} \FunctionTok{c}\NormalTok{(nome, salario))  }\CommentTok{\# essas linhas, só duas colunas}
\FunctionTok{subset}\NormalTok{(df, idade }\SpecialCharTok{\textgreater{}} \DecValTok{25}\NormalTok{, }\AttributeTok{select =} \SpecialCharTok{{-}}\NormalTok{id)               }\CommentTok{\# todas as colunas menos \textquotesingle{}id\textquotesingle{}}
\end{Highlighting}
\end{Shaded}

\subsection{\texorpdfstring{A função \texttt{summary()}}{A função summary()}}\label{a-funuxe7uxe3o-summary}

Para um resumo estatístico rápido, a função \texttt{summary()} é insuperável.
Adapta-se ao objeto: para uma variável numérica, dá os quartis, a média e os
extremos; para um fator, conta as frequências; para um data frame inteiro,
resume cada coluna à sua maneira:

\begin{Shaded}
\begin{Highlighting}[]
\FunctionTok{summary}\NormalTok{(iris)}
\DocumentationTok{\#\#   Sepal.Length   Sepal.Width    Petal.Length   Petal.Width         Species  }
\DocumentationTok{\#\#  Min.   :4.30   Min.   :2.00   Min.   :1.00   Min.   :0.1   setosa    :50  }
\DocumentationTok{\#\#  1st Qu.:5.10   1st Qu.:2.80   1st Qu.:1.60   1st Qu.:0.3   versicolor:50  }
\DocumentationTok{\#\#  Median :5.80   Median :3.00   Median :4.35   Median :1.3   virginica :50  }
\DocumentationTok{\#\#  Mean   :5.84   Mean   :3.06   Mean   :3.76   Mean   :1.2                  }
\DocumentationTok{\#\#  3rd Qu.:6.40   3rd Qu.:3.30   3rd Qu.:5.10   3rd Qu.:1.8                  }
\DocumentationTok{\#\#  Max.   :7.90   Max.   :4.40   Max.   :6.90   Max.   :2.5}
\end{Highlighting}
\end{Shaded}

\subsection{Valores ausentes em data frames}\label{valores-ausentes-em-data-frames}

Como já vimos com os vetores, os valores em falta exigem cuidado. O conjunto de
dados \texttt{airquality}, por exemplo, tem \texttt{NA} em algumas colunas, o que faz com que a
sua média direta seja \texttt{NA}:

\begin{Shaded}
\begin{Highlighting}[]
\FunctionTok{colMeans}\NormalTok{(airquality)}
\DocumentationTok{\#\#   Ozone Solar.R    Wind    Temp   Month     Day }
\DocumentationTok{\#\#      NA      NA   9.958  77.882   6.993  15.804}
\end{Highlighting}
\end{Shaded}

A solução é a mesma: o argumento \texttt{na.rm\ =\ TRUE} ignora os valores em falta:

\begin{Shaded}
\begin{Highlighting}[]
\FunctionTok{mean}\NormalTok{(airquality}\SpecialCharTok{$}\NormalTok{Ozone, }\AttributeTok{na.rm =} \ConstantTok{TRUE}\NormalTok{)}
\DocumentationTok{\#\# [1] 42.13}

\FunctionTok{colMeans}\NormalTok{(airquality, }\AttributeTok{na.rm =} \ConstantTok{TRUE}\NormalTok{)}
\end{Highlighting}
\end{Shaded}

\subsection{Exercícios}\label{exercuxedcios-4}

\textbf{1.} Crie um data frame \texttt{estudantes} com as colunas \texttt{Nome}, \texttt{Idade}, \texttt{Curso} e
\texttt{Nota}, com dados de cinco estudantes fictícios. Acrescente a estudante Filipa,
de 20 anos, de Matemática, com nota 16. Use \texttt{order()} para ordenar por \texttt{Nota}.

\textbf{2.} No data frame \texttt{estudantes}, aumente todas as notas em 0,5. Depois, mude o
\texttt{Curso} para ``Estatística'' em todos os estudantes com nota superior a 8,0.

\textbf{3.} No data frame \texttt{mtcars}, filtre as linhas com \texttt{cyl\ ==\ 6} e \texttt{hp\ \textgreater{}\ 100},
criando um novo data frame com o resultado.

\textbf{4.} Acrescente ao \texttt{mtcars} uma coluna \texttt{Peso\_KG} que converta o peso (\texttt{wt}, em
milhares de libras) para quilogramas (multiplicando por 453,592).

\textbf{5.} No \texttt{mtcars}, calcule a média, o mínimo e o máximo da potência (\texttt{hp}) para
cada número de cilindros (\texttt{cyl}). Use a função \texttt{aggregate()}.

\textbf{6.} No \texttt{mtcars}, use \texttt{sapply()} para calcular o desvio padrão de todas as
colunas numéricas.

\textbf{7.} Transforme a coluna \texttt{am} do \texttt{mtcars} num fator com os rótulos ``Automático''
(0) e ``Manual'' (1). Conte quantos carros há de cada tipo.

\textbf{8.} Ordene o \texttt{mtcars} por \texttt{mpg} decrescente e mostre os 5 carros mais e os 5
menos económicos.
\emph{Sugestão:} Use a função \texttt{order()}.

\textbf{9.} A partir do \texttt{mtcars}, acrescente uma coluna \texttt{car} com os nomes das linhas.
Crie \texttt{carros1} com as colunas \texttt{car}, \texttt{mpg}, \texttt{cyl} e \texttt{carros2} com \texttt{car}, \texttt{hp},
\texttt{wt}. Junte-os com \texttt{merge()} pela coluna \texttt{car}.

\textbf{10.} Carregue o \texttt{airquality} e mostre as primeiras e as últimas 6 linhas.

\textbf{11.} Conte quantos \texttt{NA} existem em cada coluna do \texttt{airquality}.
\emph{Sugestão:} Use a função \texttt{colSums()}.

\textbf{12.} Calcule a média de \texttt{Solar.R} no \texttt{airquality}, ignorando os \texttt{NA}.

\textbf{13.} Crie um subconjunto do \texttt{airquality} com \texttt{Wind\ \textgreater{}\ 10} e \texttt{Temp\ \textgreater{}\ 80}. Mostre
as primeiras 6 linhas.

\textbf{14.} Acrescente ao \texttt{airquality} uma coluna \texttt{Temp\_Wind\_Sum} com a soma de \texttt{Temp}
e \texttt{Wind}. Mostre as primeiras 6 linhas.

\textbf{15.} No \texttt{iris}, filtre a espécie ``setosa''. Crie a coluna \texttt{Petal.Area}
(comprimento × largura da pétala), filtre \texttt{Petal.Area\ \textgreater{}\ 0.2} e ordene por
\texttt{Petal.Length} decrescente.

\textbf{16.} No \texttt{airquality}, remova as linhas com \texttt{NA} (use \texttt{na.omit()}). Acrescente
\texttt{Temp.Celsius} = \(\frac{(Temp - 32)}{1{,}8}\). Agrupe por \texttt{Month} e calcule a média
de \texttt{Ozone} com \texttt{aggregate()}. Faça um gráfico de dispersão de \texttt{Ozone} contra
\texttt{Temp.Celsius}, colorindo por \texttt{Month}.

\textbf{17.} O naufrágio do RMS Titanic, em 1912, deu origem a um dos conjuntos de
dados mais famosos da ciência de dados. Trabalhe com uma versão reduzida e
fictícia:

\begin{enumerate}
\def\labelenumi{(\alph{enumi})}
\tightlist
\item
  Crie o data frame:
\end{enumerate}

\begin{Shaded}
\begin{Highlighting}[]
\NormalTok{titanic }\OtherTok{\textless{}{-}} \FunctionTok{data.frame}\NormalTok{(}
  \AttributeTok{nome =} \FunctionTok{c}\NormalTok{(}\StringTok{"Ana"}\NormalTok{, }\StringTok{"Bruno"}\NormalTok{, }\StringTok{"Clara"}\NormalTok{, }\StringTok{"Daniel"}\NormalTok{, }\StringTok{"Eva"}\NormalTok{, }\StringTok{"Filipe"}\NormalTok{),}
  \AttributeTok{sexo =} \FunctionTok{c}\NormalTok{(}\StringTok{"F"}\NormalTok{, }\StringTok{"M"}\NormalTok{, }\StringTok{"F"}\NormalTok{, }\StringTok{"M"}\NormalTok{, }\StringTok{"F"}\NormalTok{, }\StringTok{"M"}\NormalTok{),}
  \AttributeTok{idade =} \FunctionTok{c}\NormalTok{(}\DecValTok{22}\NormalTok{, }\DecValTok{35}\NormalTok{, }\DecValTok{18}\NormalTok{, }\DecValTok{50}\NormalTok{, }\DecValTok{28}\NormalTok{, }\DecValTok{40}\NormalTok{),}
  \AttributeTok{classe =} \FunctionTok{c}\NormalTok{(}\StringTok{"1ª"}\NormalTok{, }\StringTok{"3ª"}\NormalTok{, }\StringTok{"3ª"}\NormalTok{, }\StringTok{"1ª"}\NormalTok{, }\StringTok{"2ª"}\NormalTok{, }\StringTok{"2ª"}\NormalTok{),}
  \AttributeTok{sobreviveu =} \FunctionTok{c}\NormalTok{(}\ConstantTok{TRUE}\NormalTok{, }\ConstantTok{FALSE}\NormalTok{, }\ConstantTok{TRUE}\NormalTok{, }\ConstantTok{TRUE}\NormalTok{, }\ConstantTok{FALSE}\NormalTok{, }\ConstantTok{FALSE}\NormalTok{)}
\NormalTok{)}
\end{Highlighting}
\end{Shaded}

\begin{enumerate}
\def\labelenumi{(\alph{enumi})}
\setcounter{enumi}{1}
\tightlist
\item
  Quantas passageiras (sexo feminino) sobreviveram?
\item
  Qual a idade média das pessoas que sobreviveram?
\item
  Quantas pessoas viajavam na 3.ª classe? Qual a taxa de sobrevivência nesse
  grupo?
\end{enumerate}

\section{Listas}\label{listas}

Vimos que os vetores e as matrizes são homogéneos (um só tipo) e que o data frame
é uma tabela com colunas de tipos diferentes, mas ainda com uma forma retangular
rígida. Falta-nos uma estrutura verdadeiramente livre --- capaz de guardar, num só
objeto, coisas tão diferentes como um número, um vetor, uma matriz e até uma
função. Essa estrutura é a \textbf{lista}.

A lista é a estrutura mais flexível do R. Ao contrário do vetor, os seus
elementos podem ser de tipos e tamanhos completamente diferentes. É a ferramenta
ideal para agrupar informação heterogénea --- e, não por acaso, é a estrutura sobre
a qual o próprio data frame é construído (recorde que \texttt{typeof()} de um data frame
devolvia \texttt{"list"}).

\subsection{Criar e manipular listas}\label{criar-e-manipular-listas}

Criamos uma lista com \texttt{list()}, dando (opcionalmente) um nome a cada elemento:

\begin{Shaded}
\begin{Highlighting}[]
\NormalTok{minha\_lista }\OtherTok{\textless{}{-}} \FunctionTok{list}\NormalTok{(}
  \AttributeTok{nome =} \StringTok{"Estudante"}\NormalTok{,}
  \AttributeTok{idade =} \DecValTok{21}\NormalTok{,}
  \AttributeTok{notas =} \FunctionTok{c}\NormalTok{(}\DecValTok{85}\NormalTok{, }\DecValTok{90}\NormalTok{, }\DecValTok{92}\NormalTok{),}
  \AttributeTok{disciplinas =} \FunctionTok{c}\NormalTok{(}\StringTok{"Matemática"}\NormalTok{, }\StringTok{"Estatística"}\NormalTok{, }\StringTok{"Computação"}\NormalTok{),}
  \AttributeTok{matriz\_exemplo =} \FunctionTok{matrix}\NormalTok{(}\DecValTok{1}\SpecialCharTok{:}\DecValTok{9}\NormalTok{, }\AttributeTok{nrow =} \DecValTok{3}\NormalTok{, }\AttributeTok{byrow =} \ConstantTok{TRUE}\NormalTok{),}
  \AttributeTok{media =} \ControlFlowTok{function}\NormalTok{(x) }\FunctionTok{mean}\NormalTok{(x)}
\NormalTok{)}
\end{Highlighting}
\end{Shaded}

Repare na variedade: um texto, um número, dois vetores, uma matriz e uma função ---
tudo dentro de um único objeto. Aceder aos elementos faz-se de três formas, cuja
distinção é importante:

\begin{Shaded}
\begin{Highlighting}[]
\CommentTok{\# [[ ]] ou $ extraem o CONTEÚDO do elemento (o vetor, o número, etc.)}
\NormalTok{minha\_lista[[}\DecValTok{1}\NormalTok{]]}
\DocumentationTok{\#\# [1] "Estudante"}
\NormalTok{minha\_lista}\SpecialCharTok{$}\NormalTok{nome}
\DocumentationTok{\#\# [1] "Estudante"}

\CommentTok{\# [ ] devolve uma SUB{-}LISTA (mantém a estrutura de lista)}
\NormalTok{minha\_lista[}\DecValTok{1}\NormalTok{]}
\DocumentationTok{\#\# $nome}
\DocumentationTok{\#\# [1] "Estudante"}

\CommentTok{\# Podemos combinar: aceder a um elemento e depois indexá{-}lo}
\NormalTok{minha\_lista}\SpecialCharTok{$}\NormalTok{notas[}\DecValTok{2}\NormalTok{]}
\DocumentationTok{\#\# [1] 90}
\end{Highlighting}
\end{Shaded}

A diferença entre \texttt{minha\_lista{[}{[}1{]}{]}} e \texttt{minha\_lista{[}1{]}} é a mesma que já
encontrámos nos data frames, e é uma das subtilezas mais características do R.
\texttt{{[}{[}\ {]}{]}} ``abre a caixa'' e devolve o que está lá dentro; \texttt{{[}\ {]}} devolve uma caixa
mais pequena com o elemento ainda embrulhado. Uma imagem útil: se a lista é uma
mala com compartimentos, \texttt{{[}{[}\ {]}{]}} tira o objeto do compartimento, ao passo que
\texttt{{[}\ {]}} dá-nos uma mala mais pequena só com esse compartimento.

Acrescentar e remover elementos segue a lógica habitual:

\begin{Shaded}
\begin{Highlighting}[]
\CommentTok{\# Acrescentar um elemento}
\NormalTok{minha\_lista}\SpecialCharTok{$}\NormalTok{nota\_final }\OtherTok{\textless{}{-}} \FunctionTok{mean}\NormalTok{(minha\_lista}\SpecialCharTok{$}\NormalTok{notas)}

\CommentTok{\# Remover um elemento}
\NormalTok{minha\_lista}\SpecialCharTok{$}\NormalTok{idade }\OtherTok{\textless{}{-}} \ConstantTok{NULL}
\end{Highlighting}
\end{Shaded}

\subsection{Exercícios}\label{exercuxedcios-5}

\textbf{1.} Crie uma lista \texttt{dados\_estudante} com: \texttt{nome} = ``João'', \texttt{idade} = 21 e
\texttt{notas} = um vetor com as notas de Estatística (85), Matemática (90) e Computação
(95). Aceda e imprima o nome, a idade e a nota de Computação.

\textbf{2.} Na lista \texttt{dados\_estudante}, acrescente um elemento \texttt{status} = ``Aprovado''.
Substitua a nota de Estatística por 88. Imprima a lista completa.

\textbf{3.} Crie duas listas, \texttt{estudante1} e \texttt{estudante2}, com a mesma estrutura:

\begin{itemize}
\tightlist
\item
  \texttt{estudante1}: ``Maria'', 22, notas 78, 85, 90, status ``Aprovado''.
\item
  \texttt{estudante2}: ``Carlos'', 23, notas 70, 75, 80, status ``Aprovado''.
\end{itemize}

Combine-as numa lista \texttt{turma} e imprima-a.

\textbf{4.} Usando a lista \texttt{turma}:

\begin{enumerate}
\def\labelenumi{(\alph{enumi})}
\tightlist
\item
  Extraia e imprima o nome do segundo estudante.
\item
  Calcule a média das notas do primeiro estudante.
\item
  Mude o status do segundo estudante para ``Reprovado'' e imprima a lista.
\end{enumerate}

\textbf{5.} Crie uma lista \texttt{estatistica\_aplicada} com duas listas internas, \texttt{turma1} e
\texttt{turma2}, cada uma com dois estudantes (mesma estrutura de antes). \texttt{turma1} deve
conter \texttt{estudante1} e \texttt{estudante2}; \texttt{turma2}, dois novos estudantes. Aceda e
imprima o nome do primeiro estudante da \texttt{turma2}, a média das notas do segundo
estudante da \texttt{turma1} e a lista completa.

\textbf{6.} Uma fábrica produz 4 produtos e regista, para cada um, o nome, a quantidade
produzida mensalmente, o custo de produção unitário (€) e o preço de venda
unitário (€). Crie a lista \texttt{produtos\_fabrica} e:

\begin{enumerate}
\def\labelenumi{(\alph{enumi})}
\tightlist
\item
  Calcule e acrescente à lista a margem de lucro mensal de cada produto (preço
  menos custo, vezes a quantidade). Imprima-as.
\item
  Identifique o produto com maior margem de lucro mensal.
\item
  Aumente em 5\% o preço de venda do segundo produto e imprima a sua nova margem.
\item
  Crie uma lista com os produtos cuja margem mensal supera 10\% do custo de
  produção mensal.
\end{enumerate}

\textbf{7.} Investigadores registaram, para três estudantes, o número de mensagens
trocadas com o ChatGPT e o tempo total de uso (minutos) numa semana:

\begin{Shaded}
\begin{Highlighting}[]
\NormalTok{chatgpt\_logs }\OtherTok{\textless{}{-}} \FunctionTok{list}\NormalTok{(}
  \AttributeTok{estudante1 =} \FunctionTok{list}\NormalTok{(}\AttributeTok{nome =} \StringTok{"Ana"}\NormalTok{, }\AttributeTok{mensagens =} \FunctionTok{c}\NormalTok{(}\DecValTok{12}\NormalTok{, }\DecValTok{15}\NormalTok{, }\DecValTok{9}\NormalTok{), }\AttributeTok{tempo\_total =} \DecValTok{35}\NormalTok{),}
  \AttributeTok{estudante2 =} \FunctionTok{list}\NormalTok{(}\AttributeTok{nome =} \StringTok{"Bruno"}\NormalTok{, }\AttributeTok{mensagens =} \FunctionTok{c}\NormalTok{(}\DecValTok{7}\NormalTok{, }\DecValTok{6}\NormalTok{, }\DecValTok{5}\NormalTok{, }\DecValTok{12}\NormalTok{), }\AttributeTok{tempo\_total =} \DecValTok{42}\NormalTok{),}
  \AttributeTok{estudante3 =} \FunctionTok{list}\NormalTok{(}\AttributeTok{nome =} \StringTok{"Carla"}\NormalTok{, }\AttributeTok{mensagens =} \FunctionTok{c}\NormalTok{(}\DecValTok{10}\NormalTok{, }\DecValTok{20}\NormalTok{), }\AttributeTok{tempo\_total =} \DecValTok{28}\NormalTok{)}
\NormalTok{)}
\end{Highlighting}
\end{Shaded}

\begin{enumerate}
\def\labelenumi{(\alph{enumi})}
\tightlist
\item
  Quantas mensagens enviou o Bruno na terceira sessão?
\item
  Calcule o total de mensagens de cada estudante.
\item
  Acrescente a cada estudante um elemento \texttt{media\_tempo\_sessao} com o tempo médio
  por sessão.
\item
  Calcule a média de mensagens por sessão de cada um. Quem usou o ChatGPT de
  forma mais intensiva?
\end{enumerate}

\textbf{8.} Três perfis de Instagram sobre estatística publicaram vários conteúdos numa
semana:

\begin{Shaded}
\begin{Highlighting}[]
\NormalTok{instagram }\OtherTok{\textless{}{-}} \FunctionTok{list}\NormalTok{(}
  \AttributeTok{biostats\_girl =} \FunctionTok{list}\NormalTok{(}\AttributeTok{temas =} \FunctionTok{c}\NormalTok{(}\StringTok{"café"}\NormalTok{, }\StringTok{"dados"}\NormalTok{, }\StringTok{"meme"}\NormalTok{),}
                       \AttributeTok{visualizacoes =} \FunctionTok{c}\NormalTok{(}\DecValTok{1200}\NormalTok{, }\DecValTok{1800}\NormalTok{, }\DecValTok{2400}\NormalTok{)),}
  \AttributeTok{r\_nerd =} \FunctionTok{list}\NormalTok{(}\AttributeTok{temas =} \FunctionTok{c}\NormalTok{(}\StringTok{"dica\_R"}\NormalTok{, }\StringTok{"ggplot2"}\NormalTok{, }\StringTok{"shiny"}\NormalTok{),}
                \AttributeTok{visualizacoes =} \FunctionTok{c}\NormalTok{(}\DecValTok{800}\NormalTok{, }\DecValTok{1900}\NormalTok{, }\DecValTok{2100}\NormalTok{)),}
  \AttributeTok{estatistica\_hoje =} \FunctionTok{list}\NormalTok{(}\AttributeTok{temas =} \FunctionTok{c}\NormalTok{(}\StringTok{"história"}\NormalTok{, }\StringTok{"carreira"}\NormalTok{, }\StringTok{"código"}\NormalTok{),}
                          \AttributeTok{visualizacoes =} \FunctionTok{c}\NormalTok{(}\DecValTok{1500}\NormalTok{, }\DecValTok{1600}\NormalTok{, }\DecValTok{2000}\NormalTok{))}
\NormalTok{)}
\end{Highlighting}
\end{Shaded}

\begin{enumerate}
\def\labelenumi{(\alph{enumi})}
\tightlist
\item
  Qual foi o tema mais visto do perfil \texttt{r\_nerd}?
\item
  Calcule o total de visualizações de cada perfil e guarde-o num novo elemento
  \texttt{total\_views}.
\item
  Ordene os perfis do mais para o menos visto.
\item
  Calcule a média de visualizações por publicação de cada perfil. Qual teve
  melhor desempenho médio?
\end{enumerate}

\textbf{9.} Os resultados de três turmas na disciplina de Probabilidades incluem, por
turma, as notas dos alunos em dois testes (T1 e T2) e se entregaram o projeto:

\begin{Shaded}
\begin{Highlighting}[]
\NormalTok{avaliacoes }\OtherTok{\textless{}{-}} \FunctionTok{list}\NormalTok{(}
  \AttributeTok{turmaA =} \FunctionTok{list}\NormalTok{(}\AttributeTok{notas\_t1 =} \FunctionTok{c}\NormalTok{(}\DecValTok{14}\NormalTok{, }\DecValTok{12}\NormalTok{, }\DecValTok{16}\NormalTok{), }\AttributeTok{notas\_t2 =} \FunctionTok{c}\NormalTok{(}\DecValTok{15}\NormalTok{, }\DecValTok{13}\NormalTok{, }\DecValTok{14}\NormalTok{),}
                \AttributeTok{projeto =} \FunctionTok{c}\NormalTok{(}\ConstantTok{TRUE}\NormalTok{, }\ConstantTok{TRUE}\NormalTok{, }\ConstantTok{FALSE}\NormalTok{)),}
  \AttributeTok{turmaB =} \FunctionTok{list}\NormalTok{(}\AttributeTok{notas\_t1 =} \FunctionTok{c}\NormalTok{(}\DecValTok{10}\NormalTok{, }\DecValTok{11}\NormalTok{, }\DecValTok{9}\NormalTok{), }\AttributeTok{notas\_t2 =} \FunctionTok{c}\NormalTok{(}\DecValTok{13}\NormalTok{, }\DecValTok{12}\NormalTok{, }\DecValTok{10}\NormalTok{),}
                \AttributeTok{projeto =} \FunctionTok{c}\NormalTok{(}\ConstantTok{TRUE}\NormalTok{, }\ConstantTok{FALSE}\NormalTok{, }\ConstantTok{FALSE}\NormalTok{)),}
  \AttributeTok{turmaC =} \FunctionTok{list}\NormalTok{(}\AttributeTok{notas\_t1 =} \FunctionTok{c}\NormalTok{(}\DecValTok{17}\NormalTok{, }\DecValTok{18}\NormalTok{, }\DecValTok{19}\NormalTok{), }\AttributeTok{notas\_t2 =} \FunctionTok{c}\NormalTok{(}\DecValTok{18}\NormalTok{, }\DecValTok{17}\NormalTok{, }\DecValTok{20}\NormalTok{),}
                \AttributeTok{projeto =} \FunctionTok{c}\NormalTok{(}\ConstantTok{TRUE}\NormalTok{, }\ConstantTok{TRUE}\NormalTok{, }\ConstantTok{TRUE}\NormalTok{))}
\NormalTok{)}
\end{Highlighting}
\end{Shaded}

\begin{enumerate}
\def\labelenumi{(\alph{enumi})}
\tightlist
\item
  Calcule a média final de cada aluno em cada turma.
\item
  Quantos alunos, por turma, entregaram o projeto e tiveram média final \(\geq
  13\)?
\item
  Para cada turma, calcule a média geral das notas finais e a percentagem de
  aprovados (quem entregou o projeto e teve média \(\geq 10\)).
\item
  Qual foi a turma com melhor desempenho global?
\end{enumerate}

\chapter{Estruturas de decisão}\label{selecao}

Até agora, os nossos programas executavam sempre a mesma sequência de comandos, do
princípio ao fim, sem alternativas. Mas um programa útil precisa, muitas vezes, de
\emph{decidir}: se a nota é suficiente, aprova; caso contrário, reprova. Se o número é
par, faz uma coisa; se é ímpar, faz outra. É esta capacidade de escolher caminhos
diferentes consoante as circunstâncias que separa uma simples lista de contas de
um verdadeiro programa.

As \textbf{estruturas de decisão} (ou de seleção) controlam o fluxo de execução com
base em condições. E, felizmente, já temos as peças de que precisamos: as
condições são exatamente as expressões lógicas que aprendemos no
Capítulo~\ref{calculadora} --- aquelas que devolvem \texttt{TRUE} ou \texttt{FALSE}. Vamos
agora usá-las para dirigir o programa. As estruturas mais comuns são o \texttt{if}, o
\texttt{else} e o \texttt{else\ if}.

\section{\texorpdfstring{A condicional \texttt{if}}{A condicional if}}\label{a-condicional-if}

O \texttt{if} é a estrutura mais elementar: executa um bloco de código \textbf{apenas se} uma
condição for verdadeira. Se a condição for falsa, o bloco é simplesmente ignorado
e o programa continua.

\begin{Shaded}
\begin{Highlighting}[]
\ControlFlowTok{if}\NormalTok{ (condição) \{}
  \CommentTok{\# código executado apenas se a condição for TRUE}
\NormalTok{\}}
\end{Highlighting}
\end{Shaded}

Um exemplo concreto: verificar se alguém é maior de idade.

\begin{Shaded}
\begin{Highlighting}[]
\NormalTok{idade }\OtherTok{\textless{}{-}} \DecValTok{18}

\ControlFlowTok{if}\NormalTok{ (idade }\SpecialCharTok{\textgreater{}=} \DecValTok{18}\NormalTok{) \{}
  \FunctionTok{print}\NormalTok{(}\StringTok{"Maior de idade"}\NormalTok{)}
\NormalTok{\}}
\DocumentationTok{\#\# [1] "Maior de idade"}
\end{Highlighting}
\end{Shaded}

Como \texttt{idade} vale 18, a condição \texttt{idade\ \textgreater{}=\ 18} é verdadeira, e o R executa o
código dentro das chavetas.

Repare na anatomia da instrução: a condição vai entre \textbf{parênteses} \texttt{(\ )} e o
bloco de código a executar vai entre \textbf{chavetas} \texttt{\{\ \}}. Habitue-se a esta
pontuação desde já --- esquecer uma chaveta é um dos erros mais frequentes.

\section{\texorpdfstring{A estrutura \texttt{if...else}}{A estrutura if...else}}\label{a-estrutura-if...else}

O \texttt{if} sozinho responde a ``faz isto \emph{se}\ldots{}''. Mas muitas vezes queremos também
dizer ``\ldots{}\emph{caso contrário}, faz aquilo''. Para isso serve o \texttt{else}, que
acrescenta um caminho alternativo, executado quando a condição é falsa.

\begin{Shaded}
\begin{Highlighting}[]
\ControlFlowTok{if}\NormalTok{ (condição) \{}
  \CommentTok{\# código executado se a condição for TRUE}
\NormalTok{\} }\ControlFlowTok{else}\NormalTok{ \{}
  \CommentTok{\# código executado se a condição for FALSE}
\NormalTok{\}}
\end{Highlighting}
\end{Shaded}

Reformulando o exemplo anterior para cobrir os dois casos:

\begin{Shaded}
\begin{Highlighting}[]
\NormalTok{idade }\OtherTok{\textless{}{-}} \DecValTok{16}

\ControlFlowTok{if}\NormalTok{ (idade }\SpecialCharTok{\textgreater{}=} \DecValTok{18}\NormalTok{) \{}
  \FunctionTok{print}\NormalTok{(}\StringTok{"Maior de idade"}\NormalTok{)}
\NormalTok{\} }\ControlFlowTok{else}\NormalTok{ \{}
  \FunctionTok{print}\NormalTok{(}\StringTok{"Menor de idade"}\NormalTok{)}
\NormalTok{\}}
\DocumentationTok{\#\# [1] "Menor de idade"}
\end{Highlighting}
\end{Shaded}

Agora, como \texttt{idade} vale 16, a condição é falsa e o R executa o bloco do \texttt{else}.
Note que exatamente um dos dois blocos é executado --- nunca ambos, nunca nenhum.

\section{\texorpdfstring{A condicional \texttt{else\ if}}{A condicional else if}}\label{a-condicional-else-if}

E quando há mais de dois caminhos? Uma nota pode ser ``Excelente'', ``Bom'' ou
``Insuficiente'' --- três hipóteses, não duas. O \texttt{else\ if} permite encadear várias
condições. O R avalia-as por ordem e, assim que uma é verdadeira, executa o
respetivo bloco e ignora todos os restantes.

\begin{Shaded}
\begin{Highlighting}[]
\ControlFlowTok{if}\NormalTok{ (condição1) \{}
  \CommentTok{\# se condição1 for TRUE}
\NormalTok{\} }\ControlFlowTok{else} \ControlFlowTok{if}\NormalTok{ (condição2) \{}
  \CommentTok{\# se condição2 for TRUE (e a primeira for FALSE)}
\NormalTok{\} }\ControlFlowTok{else}\NormalTok{ \{}
  \CommentTok{\# se nenhuma das anteriores for TRUE}
\NormalTok{\}}
\end{Highlighting}
\end{Shaded}

Classifiquemos uma nota:

\begin{Shaded}
\begin{Highlighting}[]
\NormalTok{nota }\OtherTok{\textless{}{-}} \DecValTok{85}

\ControlFlowTok{if}\NormalTok{ (nota }\SpecialCharTok{\textgreater{}=} \DecValTok{90}\NormalTok{) \{}
  \FunctionTok{print}\NormalTok{(}\StringTok{"Excelente"}\NormalTok{)}
\NormalTok{\} }\ControlFlowTok{else} \ControlFlowTok{if}\NormalTok{ (nota }\SpecialCharTok{\textgreater{}=} \DecValTok{70}\NormalTok{) \{}
  \FunctionTok{print}\NormalTok{(}\StringTok{"Bom"}\NormalTok{)}
\NormalTok{\} }\ControlFlowTok{else}\NormalTok{ \{}
  \FunctionTok{print}\NormalTok{(}\StringTok{"Insuficiente"}\NormalTok{)}
\NormalTok{\}}
\DocumentationTok{\#\# [1] "Bom"}
\end{Highlighting}
\end{Shaded}

A \textbf{ordem} das condições importa. O R testa-as de cima para baixo e pára na
primeira verdadeira. Por isso, escrevemos primeiro \texttt{nota\ \textgreater{}=\ 90} e só depois
\texttt{nota\ \textgreater{}=\ 70}: uma nota de 95 satisfaz ambas, mas queremos que seja ``Excelente''.
Se invertêssemos a ordem, uma nota de 95 satisfaria logo \texttt{nota\ \textgreater{}=\ 70} e seria
erradamente rotulada de ``Bom''.

\section{\texorpdfstring{A função \texttt{ifelse()}}{A função ifelse()}}\label{a-funuxe7uxe3o-ifelse}

Todas as estruturas anteriores decidem sobre \textbf{um único valor} de cada vez. Mas
lembremo-nos de que, em R, trabalhamos constantemente com vetores. E se
quiséssemos classificar \emph{cada elemento} de um vetor? Escrever um \texttt{if} para cada um
seria absurdo.

É aqui que entra a função \texttt{ifelse()}, a versão \textbf{vetorizada} da decisão. Aplica
uma condição a cada elemento de um vetor de uma só vez, devolvendo um valor
consoante a condição seja verdadeira ou falsa:

\begin{Shaded}
\begin{Highlighting}[]
\NormalTok{resultado }\OtherTok{\textless{}{-}} \FunctionTok{ifelse}\NormalTok{(condição, valor\_se\_true, valor\_se\_false)}
\end{Highlighting}
\end{Shaded}

Vejamos:

\begin{Shaded}
\begin{Highlighting}[]
\NormalTok{valores }\OtherTok{\textless{}{-}} \FunctionTok{c}\NormalTok{(}\DecValTok{4}\NormalTok{, }\DecValTok{6}\NormalTok{, }\DecValTok{9}\NormalTok{, }\DecValTok{3}\NormalTok{)}
\FunctionTok{ifelse}\NormalTok{(valores }\SpecialCharTok{\textgreater{}} \DecValTok{5}\NormalTok{, }\StringTok{"maior que 5"}\NormalTok{, }\StringTok{"não é maior que 5"}\NormalTok{)}
\DocumentationTok{\#\# [1] "não é maior que 5" "maior que 5"       "maior que 5"      }
\DocumentationTok{\#\# [4] "não é maior que 5"}
\end{Highlighting}
\end{Shaded}

A função avaliou \texttt{valores\ \textgreater{}\ 5} para cada elemento e devolveu, elemento a elemento,
a etiqueta correspondente. É uma ferramenta indispensável para criar novas
variáveis a partir de condições --- algo que faremos muitas vezes ao trabalhar com
dados.

Não confunda \texttt{if} com \texttt{ifelse()}. O \texttt{if} decide sobre \textbf{uma} condição e executa
um \textbf{bloco} de código; espera uma condição de comprimento 1. O \texttt{ifelse()} é uma
\textbf{função} que percorre um vetor e devolve um vetor. Regra prática: se está a
decidir o rumo do programa, use \texttt{if}; se está a transformar um vetor elemento a
elemento, use \texttt{ifelse()}.

\section{Erros comuns e exemplos de análise}\label{erros-comuns-e-exemplos-de-anuxe1lise}

Ler e depurar código de decisão é tão importante como escrevê-lo. Os exemplos
seguintes treinam esse olhar crítico --- começando por identificar erros.

\textbf{Exemplo 1.} Que erro há no código abaixo?

\begin{Shaded}
\begin{Highlighting}[]
\ControlFlowTok{if}\NormalTok{ (x }\SpecialCharTok{\%\%} \DecValTok{2} \OtherTok{=} \DecValTok{0}\NormalTok{) \{}
  \FunctionTok{print}\NormalTok{(}\StringTok{"Par"}\NormalTok{)}
\NormalTok{\} }\ControlFlowTok{else}\NormalTok{ \{}
  \FunctionTok{print}\NormalTok{(}\StringTok{"Ímpar"}\NormalTok{)}
\NormalTok{\}}
\end{Highlighting}
\end{Shaded}

O erro está no \texttt{=}: dentro da condição queremos \textbf{testar} a igualdade, não
atribuir. A forma correta usa \texttt{==}:

\begin{Shaded}
\begin{Highlighting}[]
\ControlFlowTok{if}\NormalTok{ (x }\SpecialCharTok{\%\%} \DecValTok{2} \SpecialCharTok{==} \DecValTok{0}\NormalTok{) \{}
  \FunctionTok{print}\NormalTok{(}\StringTok{"Par"}\NormalTok{)}
\NormalTok{\} }\ControlFlowTok{else}\NormalTok{ \{}
  \FunctionTok{print}\NormalTok{(}\StringTok{"Ímpar"}\NormalTok{)}
\NormalTok{\}}
\end{Highlighting}
\end{Shaded}

\textbf{Exemplo 2.} Que erros há neste código?

\begin{Shaded}
\begin{Highlighting}[]
\ControlFlowTok{if}\NormalTok{ (a }\SpecialCharTok{\textgreater{}} \DecValTok{0}\NormalTok{) \{}
  \FunctionTok{print}\NormalTok{(}\StringTok{"Positivo"}\NormalTok{)}
  \ControlFlowTok{if}\NormalTok{ (a }\SpecialCharTok{\%\%} \DecValTok{5} \OtherTok{=} \DecValTok{0}\NormalTok{)}
    \FunctionTok{print}\NormalTok{(}\StringTok{"Divisível por 5"}\NormalTok{)}
\NormalTok{\} }\ControlFlowTok{else} \ControlFlowTok{if}\NormalTok{ (a }\SpecialCharTok{==} \DecValTok{0}\NormalTok{)}
    \FunctionTok{print}\NormalTok{(}\StringTok{"Zero"}\NormalTok{)}
  \ControlFlowTok{else} \ControlFlowTok{if}\NormalTok{ \{}
    \FunctionTok{print}\NormalTok{(}\StringTok{"Negativo"}\NormalTok{)}
\NormalTok{\}}
\end{Highlighting}
\end{Shaded}

Há vários problemas: o \texttt{=} em vez de \texttt{==}, chavetas em falta e um \texttt{else\ if} sem
condição (que devia ser apenas \texttt{else}). A versão corrigida:

\begin{Shaded}
\begin{Highlighting}[]
\ControlFlowTok{if}\NormalTok{ (a }\SpecialCharTok{\textgreater{}} \DecValTok{0}\NormalTok{) \{}
  \FunctionTok{print}\NormalTok{(}\StringTok{"Positivo"}\NormalTok{)}
  \ControlFlowTok{if}\NormalTok{ (a }\SpecialCharTok{\%\%} \DecValTok{5} \SpecialCharTok{==} \DecValTok{0}\NormalTok{) \{}
    \FunctionTok{print}\NormalTok{(}\StringTok{"Divisível por 5"}\NormalTok{)}
\NormalTok{  \}}
\NormalTok{\} }\ControlFlowTok{else} \ControlFlowTok{if}\NormalTok{ (a }\SpecialCharTok{==} \DecValTok{0}\NormalTok{) \{}
  \FunctionTok{print}\NormalTok{(}\StringTok{"Zero"}\NormalTok{)}
\NormalTok{\} }\ControlFlowTok{else}\NormalTok{ \{}
  \FunctionTok{print}\NormalTok{(}\StringTok{"Negativo"}\NormalTok{)}
\NormalTok{\}}
\end{Highlighting}
\end{Shaded}

Use sempre chavetas \texttt{\{\ \}}, mesmo quando o bloco tem uma só linha. É verdade que o
R as dispensa nesse caso, mas incluí-las torna o código mais claro, evita
surpresas quando mais tarde acrescentamos linhas ao bloco, e previne muitos dos
erros que vimos acima.

\textbf{Exemplo 3.} Quais são os valores de \texttt{x} e \texttt{y} no fim da execução?

\begin{Shaded}
\begin{Highlighting}[]
\NormalTok{x }\OtherTok{\textless{}{-}} \DecValTok{1}
\NormalTok{y }\OtherTok{\textless{}{-}} \DecValTok{0}
\ControlFlowTok{if}\NormalTok{ (x }\SpecialCharTok{==} \DecValTok{1}\NormalTok{) \{}
\NormalTok{  y }\OtherTok{\textless{}{-}}\NormalTok{ y }\SpecialCharTok{{-}} \DecValTok{1}
\NormalTok{\}}
\ControlFlowTok{if}\NormalTok{ (y }\SpecialCharTok{==} \DecValTok{1}\NormalTok{) \{}
\NormalTok{  x }\OtherTok{\textless{}{-}}\NormalTok{ x }\SpecialCharTok{+} \DecValTok{1}
\NormalTok{\}}
\end{Highlighting}
\end{Shaded}

Vale a pena seguir o raciocínio passo a passo, como faria o R. A primeira condição
(\texttt{x\ ==\ 1}) é verdadeira, logo \texttt{y} passa a \(-1\). A segunda (\texttt{y\ ==\ 1}) é falsa
(porque \texttt{y} vale agora \(-1\)), logo \texttt{x} mantém-se em 1. No fim: \texttt{x\ =\ 1} e \texttt{y\ =\ -1}.

\textbf{Exemplo 4.} Considere o programa seguinte.

\begin{Shaded}
\begin{Highlighting}[]
\ControlFlowTok{if}\NormalTok{ (x }\SpecialCharTok{==} \DecValTok{1}\NormalTok{) \{}
\NormalTok{  x }\OtherTok{\textless{}{-}}\NormalTok{ x }\SpecialCharTok{+} \DecValTok{1}
  \ControlFlowTok{if}\NormalTok{ (x }\SpecialCharTok{==} \DecValTok{1}\NormalTok{) \{}
\NormalTok{    x }\OtherTok{\textless{}{-}}\NormalTok{ x }\SpecialCharTok{+} \DecValTok{1}
\NormalTok{  \} }\ControlFlowTok{else}\NormalTok{ \{}
\NormalTok{    x }\OtherTok{\textless{}{-}}\NormalTok{ x }\SpecialCharTok{{-}} \DecValTok{1}
\NormalTok{  \}}
\NormalTok{\} }\ControlFlowTok{else}\NormalTok{ \{}
\NormalTok{  x }\OtherTok{\textless{}{-}}\NormalTok{ x }\SpecialCharTok{{-}} \DecValTok{1}
\NormalTok{\}}
\end{Highlighting}
\end{Shaded}

\begin{itemize}
\tightlist
\item
  Se \texttt{x\ =\ 1}, qual será o valor final de \texttt{x}?
\item
  Que valor teria de ter \texttt{x} para que, no fim, valesse -1?
\item
  Há uma parte do programa que \textbf{nunca} é executada. Qual, e porquê?
\end{itemize}

Sugestão para o terceiro ponto: dentro do primeiro \texttt{if}, começamos por fazer
\texttt{x\ \textless{}-\ x\ +\ 1}. Que valor tem \texttt{x} imediatamente a seguir? Pode a condição
\texttt{if\ (x\ ==\ 1)} interior alguma vez ser verdadeira nesse ponto? Percorrer o código
``à mão'', anotando o valor das variáveis a cada passo, é a técnica mais fiável para
responder a estas perguntas --- e um hábito que vale a pena cultivar.

\section{Exercícios}\label{exercuxedcios-6}

\textbf{1.} Escreva um programa que leia um número e exiba ``O número é positivo'' se for
maior que zero.

\textbf{2.} Escreva um programa que leia um inteiro e imprima ``Par'' se for par e
``Ímpar'' caso contrário.

\textbf{3.} Escreva um programa que leia um número não negativo e indique se está no
intervalo \([0, 10]\), \([11, 20]\), ou se é maior que 20 (intervalos inclusivos).
Por exemplo, para 15, deve exibir ``O número está no intervalo {[}11, 20{]}''.

\textbf{4.} Escreva um programa que leia uma idade e a classifique como ``Criança''
(0--12), ``Adolescente'' (13--17), ``Adulto'' (18--64) ou ``Idoso'' (\(\geq 65\)). Se a idade for
negativa, exiba ``Idade inválida''.

\textbf{5.} Escreva um programa que leia o preço de um produto e aplique um desconto:
5\% se o preço for inferior a 100 €; 10\% se estiver entre 100 € e 500 €
(inclusive); 15\% se for superior a 500 €. Se o preço for negativo, exiba ``Preço
inválido''. Por exemplo, para 350, deve exibir ``O preço com desconto é 315 €''.

\textbf{6.} Escreva um programa que leia as coordenadas \((x, y)\) de um ponto e
determine o quadrante em que se encontra. Se estiver sobre um eixo, indique ``Eixo
X'' ou ``Eixo Y''. Por exemplo, para \((-3, 2)\), deve exibir ``Segundo quadrante''.

\textbf{7.} Escreva um programa que leia um ano e determine se é bissexto (divisível
por 4, mas não por 100, exceto se for divisível por 400). Se o ano for negativo,
exiba ``Ano inválido''. Por exemplo, para 2000, deve exibir ``Ano bissexto''.

\textbf{8.} Escreva um programa que leia um inteiro entre 0 e 999 e o descreva em
centenas, dezenas e unidades. Por exemplo, para 304, deve exibir ``3 centenas 0
dezenas e 4 unidades''. Fora de \([0, 999]\), exiba ``Número fora do intervalo''.

\textbf{9.} Escreva um programa que leia um inteiro positivo (\(\leq 5\)) e escreva a sua
representação em numeração romana. Fora de \([1, 5]\), exiba ``Número fora do
intervalo''. Por exemplo, para 3, deve exibir ``III''.

\textbf{10.} Escreva um programa que leia quatro inteiros (um de cada vez) e exiba,
após cada leitura, o menor número lido até ao momento.

\textbf{11.} Crie um vetor com os números de 1 a 10 e use \texttt{ifelse()} para rotular cada
número como ``par'' ou ``ímpar''.

\textbf{12.} Crie o vetor de notas \texttt{c(50,\ 85,\ 70,\ 40,\ 60,\ 90)}. Use \texttt{ifelse()} para
marcar como ``Aprovado'' quem tem nota \(\geq 50\) e ``Reprovado'' os restantes. Depois,
modifique o programa para que as notas entre 50 e 59 sejam ``Aprovado com
Restrição''.

\textbf{13.} Dado um vetor de idades, use \texttt{ifelse()} (encadeado) para as categorizar em
``Criança'' (\(\leq 12\)), ``Adolescente'' (13-18) e ``Adulto'' (\textgreater{} 18).

\chapter{Funções}\label{funcoes}

Ao longo dos capítulos anteriores usámos dezenas de funções --- \texttt{mean()}, \texttt{sum()},
\texttt{c()}, \texttt{factor()}, \texttt{subset()}. Chegou o momento de virar a mesa e aprender a
escrever as \emph{nossas próprias} funções. Esta é, talvez, a competência mais
transformadora de toda a programação: a partir do momento em que sabemos criar
funções, deixamos de ser apenas utilizadores da linguagem e passamos a
estendê-la.

A motivação é simples e profunda. Imagine que precisa de calcular o índice de
massa corporal de dez pessoas. Poderia repetir a fórmula dez vezes --- mas, se se
enganasse na fórmula, teria de corrigir dez lugares, e se quisesse mudar algo,
outros dez. A alternativa é escrever a fórmula \textbf{uma vez}, dentro de uma função,
e chamá-la dez vezes. A este princípio dá-se por vezes o nome de ``não te
repitas'': tudo o que se repete deve ser encapsulado num só lugar.

Uma função é, então, um bloco de código com nome, que realiza uma tarefa
específica e só corre quando é chamado. As suas vantagens resumem-se a três
palavras:

\begin{itemize}
\tightlist
\item
  \textbf{Reutilização} --- uma vez definida, chama-se as vezes que forem precisas, em
  qualquer parte do programa.
\item
  \textbf{Modularização} --- divide um problema grande em peças pequenas, cada uma com
  uma responsabilidade clara, tornando o código mais organizado e legível.
\item
  \textbf{Abstração} --- depois de escrita e testada, podemos usar a função sem voltar a
  pensar em como funciona por dentro. \texttt{mean()} calcula uma média sem nos obrigar a
  saber como.
\end{itemize}

\section{Anatomia de uma função}\label{anatomia-de-uma-funuxe7uxe3o}

A estrutura de uma função em R é a seguinte:

\begin{Shaded}
\begin{Highlighting}[]
\NormalTok{nome\_da\_funcao }\OtherTok{\textless{}{-}} \ControlFlowTok{function}\NormalTok{(argumentos) \{}
  \CommentTok{\# corpo: o código que realiza a tarefa}
\NormalTok{  resultado }\OtherTok{\textless{}{-}}\NormalTok{ ...}
  \FunctionTok{return}\NormalTok{(resultado)   }\CommentTok{\# o valor devolvido}
\NormalTok{\}}
\end{Highlighting}
\end{Shaded}

Quatro elementos a identificar:

\begin{itemize}
\tightlist
\item
  \textbf{Nome} --- o identificador pelo qual chamaremos a função (segue as regras dos
  nomes de variáveis do Capítulo~\ref{calculadora}).
\item
  \textbf{Argumentos} --- os valores de entrada que a função recebe. Podem ser
  obrigatórios ou opcionais.
\item
  \textbf{Corpo} --- o bloco de código, entre chavetas, que faz o trabalho.
\item
  \textbf{Valor devolvido} --- o resultado que a função entrega, indicado por
  \texttt{return()}.
\end{itemize}

Vejamos um primeiro exemplo completo: uma função que calcula a área de um
retângulo.

\begin{Shaded}
\begin{Highlighting}[]
\NormalTok{calcula\_area }\OtherTok{\textless{}{-}} \ControlFlowTok{function}\NormalTok{(largura, altura) \{}
\NormalTok{  area }\OtherTok{\textless{}{-}}\NormalTok{ largura }\SpecialCharTok{*}\NormalTok{ altura}
  \FunctionTok{return}\NormalTok{(area)}
\NormalTok{\}}

\CommentTok{\# Chamar a função}
\FunctionTok{calcula\_area}\NormalTok{(}\DecValTok{4}\NormalTok{, }\DecValTok{3}\NormalTok{)}
\DocumentationTok{\#\# [1] 12}
\end{Highlighting}
\end{Shaded}

Definimos a função uma vez e, a partir daí, calcular a área de qualquer retângulo
é uma questão de a chamar com novos valores. É este o ganho.

\subsection{Variáveis locais e globais}\label{variuxe1veis-locais-e-globais}

Há aqui uma subtileza importante, que evita muita confusão mais tarde. Dentro da
função \texttt{calcula\_area}, as variáveis \texttt{largura}, \texttt{altura} e \texttt{area} são \textbf{locais}:
existem apenas enquanto a função corre e desaparecem quando ela termina. Não são
visíveis fora da função, nem interferem com variáveis de igual nome que existam no
resto do programa.

\begin{Shaded}
\begin{Highlighting}[]
\NormalTok{x }\OtherTok{\textless{}{-}} \DecValTok{10}

\NormalTok{minha\_funcao }\OtherTok{\textless{}{-}} \ControlFlowTok{function}\NormalTok{() \{}
\NormalTok{  x }\OtherTok{\textless{}{-}} \DecValTok{5}        \CommentTok{\# este \textquotesingle{}x\textquotesingle{} é local: só existe dentro da função}
  \FunctionTok{return}\NormalTok{(x)}
\NormalTok{\}}

\FunctionTok{minha\_funcao}\NormalTok{()}
\DocumentationTok{\#\# [1] 5}
\NormalTok{x               }\CommentTok{\# o \textquotesingle{}x\textquotesingle{} global não foi afetado}
\DocumentationTok{\#\# [1] 10}
\end{Highlighting}
\end{Shaded}

Este isolamento é uma vantagem, não um obstáculo. Significa que podemos escrever
uma função sem receio de que os seus nomes internos colidam com o resto do
programa --- cada função tem o seu ``mundo'' próprio. É uma das razões pelas quais as
funções tornam o código mais seguro e mais fácil de manter.

\section{Argumentos: mais flexibilidade}\label{argumentos-mais-flexibilidade}

As funções tornam-se muito mais úteis quando aprendemos a controlar como os seus
argumentos são passados.

\subsection{Valores por omissão}\label{valores-por-omissuxe3o}

Podemos dar a um argumento um \textbf{valor por omissão}, usado sempre que quem chama a
função não o especifica. Isto torna certos argumentos opcionais:

\begin{Shaded}
\begin{Highlighting}[]
\NormalTok{calcula\_area }\OtherTok{\textless{}{-}} \ControlFlowTok{function}\NormalTok{(largura, }\AttributeTok{altura =} \DecValTok{2}\NormalTok{) \{}
  \FunctionTok{return}\NormalTok{(largura }\SpecialCharTok{*}\NormalTok{ altura)}
\NormalTok{\}}

\FunctionTok{calcula\_area}\NormalTok{(}\DecValTok{4}\NormalTok{)     }\CommentTok{\# \textquotesingle{}altura\textquotesingle{} não foi dado, usa{-}se o valor por omissão (2)}
\DocumentationTok{\#\# [1] 8}
\end{Highlighting}
\end{Shaded}

\subsection{Argumentos por nome}\label{argumentos-por-nome}

Ao chamar uma função, podemos passar os argumentos \textbf{pelo nome}, e não apenas
pela posição. Isto liberta-nos da ordem e torna as chamadas mais legíveis:

\begin{Shaded}
\begin{Highlighting}[]
\NormalTok{calcula\_area }\OtherTok{\textless{}{-}} \ControlFlowTok{function}\NormalTok{(largura, }\AttributeTok{altura =} \DecValTok{2}\NormalTok{) \{}
  \FunctionTok{return}\NormalTok{(largura }\SpecialCharTok{*}\NormalTok{ altura)}
\NormalTok{\}}

\FunctionTok{calcula\_area}\NormalTok{(}\AttributeTok{altura =} \DecValTok{4}\NormalTok{, }\AttributeTok{largura =} \DecValTok{3}\NormalTok{)   }\CommentTok{\# a ordem já não importa}
\DocumentationTok{\#\# [1] 12}
\end{Highlighting}
\end{Shaded}

Passar argumentos pelo nome é uma excelente prática quando uma função tem muitos
argumentos, ou quando queremos especificar apenas alguns dos opcionais. Compare
\texttt{rnorm(100,\ 0,\ 1)} com \texttt{rnorm(n\ =\ 100,\ mean\ =\ 0,\ sd\ =\ 1)}: a segunda forma
dispensa quem lê o código de decorar a ordem dos argumentos.

\section{\texorpdfstring{Validar entradas com \texttt{stop()}}{Validar entradas com stop()}}\label{validar-entradas-com-stop}

Uma boa função protege-se contra usos indevidos. A função \texttt{stop()} interrompe a
execução e emite uma mensagem de erro personalizada, útil quando uma condição
necessária não é satisfeita:

\begin{Shaded}
\begin{Highlighting}[]
\NormalTok{f }\OtherTok{\textless{}{-}} \ControlFlowTok{function}\NormalTok{(x) \{}
  \ControlFlowTok{if}\NormalTok{ (x }\SpecialCharTok{\textless{}} \DecValTok{0}\NormalTok{) \{}
    \FunctionTok{stop}\NormalTok{(}\StringTok{"Erro: x não pode ser negativo"}\NormalTok{)}
\NormalTok{  \}}
  \FunctionTok{return}\NormalTok{(}\FunctionTok{sqrt}\NormalTok{(x))}
\NormalTok{\}}

\FunctionTok{f}\NormalTok{(}\SpecialCharTok{{-}}\DecValTok{2}\NormalTok{)}
\DocumentationTok{\#\# Error in f({-}2): Erro: x não pode ser negativo}
\end{Highlighting}
\end{Shaded}

Ao verificar \texttt{x\ \textless{}\ 0} logo no início, a função recusa liminarmente uma entrada
inválida (a raiz quadrada de um negativo), em vez de devolver um \texttt{NaN} silencioso
que poderia contaminar cálculos posteriores. Antecipar e sinalizar erros é a marca
de uma função bem escrita.

\section{Exercícios de análise}\label{exercuxedcios-de-anuxe1lise}

Antes dos exercícios de construção, alguns exercícios de leitura --- para
consolidar a compreensão do que as funções fazem.

\textbf{Exemplo 1.} Qual é o resultado da chamada?

\begin{Shaded}
\begin{Highlighting}[]
\NormalTok{dados\_estudante }\OtherTok{\textless{}{-}} \ControlFlowTok{function}\NormalTok{(nome, }\AttributeTok{altura =} \DecValTok{167}\NormalTok{) \{}
  \FunctionTok{print}\NormalTok{(}\FunctionTok{paste}\NormalTok{(}\StringTok{"O(A) estudante"}\NormalTok{, nome, }\StringTok{"tem"}\NormalTok{, altura, }\StringTok{"centímetros de altura."}\NormalTok{))}
\NormalTok{\}}

\FunctionTok{dados\_estudante}\NormalTok{(}\StringTok{"Joana"}\NormalTok{, }\DecValTok{160}\NormalTok{)}
\end{Highlighting}
\end{Shaded}

\textbf{Exemplo 2.} Qual é a sintaxe correta para definir uma função que soma dois
números?

\begin{enumerate}
\def\labelenumi{(\alph{enumi})}
\tightlist
\item
  \texttt{soma\ \textless{}-\ function(x,\ y)\ \{\ return(x\ +\ y)\ \}}
\item
  \texttt{function\ soma(x,\ y)\ \{\ return(x\ +\ y)\ \}}
\item
  \texttt{def\ soma(x,\ y)\ \{\ return(x\ +\ y)\ \}}
\item
  \texttt{soma(x,\ y)\ =\ function\ \{\ return(x\ +\ y)\ \}}
\end{enumerate}

\textbf{Exemplo 3.} Considerando a função \texttt{dados\_estudante} do Exemplo 1, quais das
chamadas seguintes estão corretas?

\begin{Shaded}
\begin{Highlighting}[]
\FunctionTok{dados\_estudante}\NormalTok{(}\StringTok{"Joana"}\NormalTok{, }\DecValTok{160}\NormalTok{)}
\FunctionTok{dados\_estudante}\NormalTok{(}\AttributeTok{altura =} \DecValTok{160}\NormalTok{, }\AttributeTok{nome =} \StringTok{"Joana"}\NormalTok{)}
\FunctionTok{dados\_estudante}\NormalTok{(}\AttributeTok{nome =} \StringTok{"Joana"}\NormalTok{, }\DecValTok{160}\NormalTok{)}
\FunctionTok{dados\_estudante}\NormalTok{(}\AttributeTok{altura =} \DecValTok{160}\NormalTok{, }\StringTok{"Joana"}\NormalTok{)}
\FunctionTok{dados\_estudante}\NormalTok{(}\DecValTok{160}\NormalTok{)}
\end{Highlighting}
\end{Shaded}

A última chamada, \texttt{dados\_estudante(160)}, é sintaticamente válida mas
logicamente absurda: 160 será interpretado como o \texttt{nome} (por ser o primeiro
argumento) e a \texttt{altura} usará o valor por omissão, 167 --- resultando em ``O(A)
estudante 160 tem 167 centímetros de altura.'' Um bom lembrete de que o R faz o que
lhe pedimos, não o que queremos dizer.

\textbf{Exemplo 4.} Qual é o valor de \texttt{resultado}?

\begin{Shaded}
\begin{Highlighting}[]
\NormalTok{adi }\OtherTok{\textless{}{-}} \ControlFlowTok{function}\NormalTok{(a, b) \{}
  \FunctionTok{return}\NormalTok{(}\FunctionTok{c}\NormalTok{(a }\SpecialCharTok{+} \DecValTok{5}\NormalTok{, b }\SpecialCharTok{+} \DecValTok{5}\NormalTok{))}
\NormalTok{\}}

\NormalTok{resultado }\OtherTok{\textless{}{-}} \FunctionTok{adi}\NormalTok{(}\DecValTok{3}\NormalTok{, }\DecValTok{2}\NormalTok{)}
\end{Highlighting}
\end{Shaded}

\section{Exercícios}\label{exercuxedcios-7}

\textbf{1.} Crie uma função \texttt{quadrado()} que receba um número \texttt{x} e devolva o seu
quadrado. Teste-a com 3, 5 e 7.

\textbf{2.} Crie uma função \texttt{hipotenusa(a,\ b)} que calcule a hipotenusa de um triângulo
retângulo a partir dos dois catetos, usando o Teorema de Pitágoras.

\textbf{3.} Crie uma função \texttt{divisao\_segura(numerador,\ denominador)} que devolva a
divisão, ou uma mensagem de erro apropriada se o denominador for zero.

\textbf{4.} Crie uma função que receba dois inteiros e devolva o maior. Teste-a lendo
dois números do utilizador. Depois, crie uma segunda função que devolva o menor,
aproveitando a primeira.

\textbf{5.} Crie uma função que devolva o maior de três inteiros, reutilizando a função
\texttt{maior()} do exercício anterior.

\textbf{6.} Crie uma função \texttt{analise\_vetor()} que receba um vetor numérico e devolva um
vetor com a soma, a média, o desvio padrão e o máximo. Aplique-a a
\texttt{c(4,\ 8,\ 15,\ 16,\ 23,\ 42)}. \emph{Funções úteis:} \texttt{sum()}, \texttt{mean()}, \texttt{sd()}, \texttt{max()}.

\textbf{7.} Crie uma função \texttt{classificar\_numero()} que classifique um inteiro como
``positivo'', ``negativo'' ou ``zero''.

\textbf{8.} Crie uma função que calcule o Índice de Massa Corporal a partir do peso
(kg) e da altura (m) e o classifique segundo a tabela:

{\def\LTcaptype{none} % do not increment counter
\begin{longtable}[]{@{}ll@{}}
\toprule\noalign{}
IMC (kg/m²) & Classificação \\
\midrule\noalign{}
\endhead
\bottomrule\noalign{}
\endlastfoot
\textless{} 18,5 & Baixo peso \\
18,5 -- 24,9 & Peso normal \\
25 -- 29,9 & Excesso de peso \\
30 -- 34,9 & Obesidade grau I \\
35 -- 39,9 & Obesidade grau II \\
\(\geq\) 40 & Obesidade grau III \\
\end{longtable}
}

Recorde que \(\text{IMC} = \dfrac{\text{Peso}}{\text{Altura}^2}\).

\textbf{9.} Crie funções para converter euros noutras moedas, com base em taxas de
câmbio. Escreva depois um programa que peça um valor em euros e imprima a
conversão para dólar, libra, iene e real.

\textbf{10.} Crie uma função que calcule o intervalo de confiança a
\((1 - \alpha) \times 100\%\) para o valor médio (para amostras grandes, de dimensão
\(n \geq 30\)), dado por
\[\left(\bar{x} - z_{1-\frac{\alpha}{2}}\frac{s}{\sqrt{n}},\;
\bar{x} + z_{1-\frac{\alpha}{2}}\frac{s}{\sqrt{n}}\right),\]
onde \(s\) é o desvio padrão amostral e \(z_{1-\frac{\alpha}{2}}\) é o quantil de
probabilidade \(1 - \alpha/2\) (use \texttt{qnorm()}).

\begin{Shaded}
\begin{Highlighting}[]
\NormalTok{IC }\OtherTok{\textless{}{-}} \ControlFlowTok{function}\NormalTok{(amostra, alpha) \{}
  \CommentTok{\# o seu código aqui}
\NormalTok{\}}
\end{Highlighting}
\end{Shaded}

\textbf{11.} Generalize a função anterior para calcular o intervalo de confiança para
vários valores de \(\alpha\) de uma só vez (por exemplo, 0,001, 0,01, 0,05 e 0,10).

\begin{Shaded}
\begin{Highlighting}[]
\NormalTok{ICs }\OtherTok{\textless{}{-}} \ControlFlowTok{function}\NormalTok{(amostra, alphas) \{}
  \CommentTok{\# o seu código aqui}
\NormalTok{\}}
\end{Highlighting}
\end{Shaded}

\textbf{12.} Crie uma função \texttt{calcula\_moda()} que receba um vetor (numérico ou de
carateres) e devolva a moda --- o valor mais frequente. Se houver mais do que uma
moda, deve devolvê-las todas.
\emph{Sugestão:} Use a função \texttt{table()} para contar ocorrências.

Teste com \texttt{c(1,\ 2,\ 2,\ 3,\ 3,\ 4,\ 4)} e com \texttt{c("A",\ "B",\ "A",\ "C",\ "C",\ "C",\ "B")}.

\textbf{13.} A distância euclidiana entre dois pontos \(A(x_1, y_1)\) e \(B(x_2, y_2)\) é
\[d(A, B) = \sqrt{(x_2 - x_1)^2 + (y_2 - y_1)^2}.\]
Crie uma função \texttt{distancia\_euclidiana()} que a calcule, e teste-a com \(A(2, 3)\) e
\(B(6, 9)\). Generalize-a depois para pontos em qualquer dimensão e teste-a com
\(A(1, 2, 3)\) e \(B(4, 5, 6)\).

\textbf{14.} Crie uma função \texttt{cria\_matriz\_aleatoria(linhas,\ colunas)} que devolva uma
matriz com números aleatórios uniformes entre 0 e 1. Crie depois uma função
\texttt{analisa\_matriz()} que devolva a soma, a média e o traço (se a matriz for
quadrada) de uma matriz dada. Teste-as com uma matriz \(5 \times 5\). Por fim,
generalize \texttt{analisa\_matriz()} para devolver também o determinante, apenas quando a
matriz for quadrada; caso contrário, deve informar que não é possível calculá-lo.

Com as funções, fechamos a caixa de ferramentas fundamental da programação em R.
Temos agora tudo o que é preciso --- valores, estruturas de dados, decisões e
funções --- para, nos próximos capítulos, automatizar tarefas com ciclos, ler dados
reais e, finalmente, mergulhar na probabilidade e na simulação.

\chapter{Scripts em R}\label{scripts}

Nos capítulos anteriores escrevemos código diretamente na consola e definimos as
nossas próprias funções (Capítulo~\ref{funcoes}). Enquanto experimentamos, a
consola é cómoda; mas o trabalho que fica apenas na consola perde-se quando
fechamos a sessão e é difícil de repetir com exatidão. Chegou o momento de guardar
o código num lugar permanente: um \emph{script}.

Um \textbf{script} é simplesmente um ficheiro de texto, com a extensão \texttt{.R}, que reúne
uma sequência de instruções --- variáveis, funções, blocos de código --- pensada para
ser guardada, lida e executada de novo sempre que for preciso. Os scripts são a
base de duas ideias centrais em ciência de dados: a \textbf{automatização} de tarefas
repetitivas e a \textbf{reprodutibilidade} da análise, isto é, a garantia de que
qualquer pessoa --- incluindo o nosso próprio ``eu'' de daqui a seis meses --- consegue
obter exatamente os mesmos resultados a partir dos mesmos dados.

Neste capítulo aprendemos a guardar código num ficheiro \texttt{.R}, a executá-lo com
\texttt{source()}, a compreender o papel do diretório de trabalho e a adotar hábitos que
tornam a nossa análise reprodutível.

\section{Criar um script}\label{criar-um-script}

No RStudio, cria-se um novo script pelo menu \emph{File \textgreater{} New File \textgreater{} R Script} (ou com
o atalho \texttt{Ctrl+Shift+N} / \texttt{Cmd+Shift+N}). Escreve-se o código no editor e guarda-se
o ficheiro com a extensão \texttt{.R}.

Guarde, por exemplo, o código seguinte num ficheiro chamado \texttt{meu\_script.R}:

\begin{Shaded}
\begin{Highlighting}[]
\NormalTok{produto }\OtherTok{\textless{}{-}} \ControlFlowTok{function}\NormalTok{(x, y) \{}
  \FunctionTok{return}\NormalTok{(x }\SpecialCharTok{*}\NormalTok{ y)}
\NormalTok{\}}
\end{Highlighting}
\end{Shaded}

Repare que guardar o script \textbf{não executa} o código --- apenas o preserva. Definir a
função e poder usá-la são dois passos distintos.

\section{Executar um script}\label{executar-um-script}

Para correr um script guardado, usamos a função \texttt{source()}, na consola ou dentro de
outro script. O R lê o ficheiro e executa, uma a uma, todas as suas linhas.

Se o ficheiro \texttt{meu\_script.R} estiver no \textbf{diretório de trabalho atual}, basta:

\begin{Shaded}
\begin{Highlighting}[]
\FunctionTok{source}\NormalTok{(}\StringTok{"meu\_script.R"}\NormalTok{)}
\end{Highlighting}
\end{Shaded}

Se estiver noutra pasta, indicamos o caminho completo:

\begin{Shaded}
\begin{Highlighting}[]
\FunctionTok{source}\NormalTok{(}\StringTok{"/caminho/para/o/seu/script/meu\_script.R"}\NormalTok{)}
\end{Highlighting}
\end{Shaded}

\subsection{O que acontece quando executamos um script?}\label{o-que-acontece-quando-executamos-um-script}

Ao chamar \texttt{source()}, todas as funções, variáveis e resultados definidos no script
passam a estar disponíveis no ambiente de trabalho. Depois de executar o
\texttt{meu\_script.R} acima, a função \texttt{produto()} fica pronta a usar:

\begin{Shaded}
\begin{Highlighting}[]
\FunctionTok{produto}\NormalTok{(}\DecValTok{2}\NormalTok{, }\DecValTok{3}\NormalTok{)}
\DocumentationTok{\#\# [1] 6}
\end{Highlighting}
\end{Shaded}

É esta a lógica que torna os scripts tão úteis: escrevemos e testamos um conjunto de
funções uma vez, guardamo-las num ficheiro e, em qualquer análise futura, basta uma
linha --- \texttt{source("as\_minhas\_funcoes.R")} --- para as ter todas à disposição.

No RStudio, o botão \textbf{Source} (canto superior direito do editor) executa o script
inteiro, e \texttt{Ctrl+Enter} / \texttt{Cmd+Enter} corre apenas a linha ou o bloco selecionado.
Correr linha a linha é a melhor forma de perceber e depurar um script enquanto o
escrevemos.

\section{O diretório de trabalho}\label{o-diretuxf3rio-de-trabalho}

O \textbf{diretório de trabalho} (\emph{working directory}) é a pasta que o R toma como ponto
de referência: é aí que procura os ficheiros a ler e onde guarda os resultados, a
menos que indiquemos um caminho completo. Duas funções bastam para o gerir:

\begin{itemize}
\tightlist
\item
  \texttt{getwd()} --- devolve o caminho do diretório de trabalho atual.
\item
  \texttt{setwd("caminho/da/pasta")} --- muda o diretório de trabalho para o caminho
  indicado.
\end{itemize}

\begin{Shaded}
\begin{Highlighting}[]
\FunctionTok{getwd}\NormalTok{()}
\DocumentationTok{\#\# [1] "C:/Users/Utilizador/Documents"}
\end{Highlighting}
\end{Shaded}

Com o diretório de trabalho bem definido, referimo-nos aos ficheiros pelo nome, sem
escrever o caminho todo. Se \texttt{dados.csv} estiver no diretório de trabalho:

\begin{Shaded}
\begin{Highlighting}[]
\NormalTok{dados }\OtherTok{\textless{}{-}} \FunctionTok{read.csv}\NormalTok{(}\StringTok{"dados.csv"}\NormalTok{)}
\end{Highlighting}
\end{Shaded}

Evite espalhar chamadas a \texttt{setwd("C:/Users/...")} pelos seus scripts. Um caminho
absoluto só existe no seu computador: quando o script chega às mãos de outra pessoa
(ou muda de máquina), deixa de funcionar --- e a reprodutibilidade, que era o objetivo,
perde-se. A boa prática é trabalhar dentro de um \textbf{Projeto do RStudio} (\emph{File \textgreater{} New
Project}), que fixa automaticamente o diretório de trabalho na pasta do projeto e
permite usar apenas caminhos \textbf{relativos}, como \texttt{"dados/dados.csv"}.

\section{Exercícios}\label{exercuxedcios-8}

\textbf{1.} Crie os seguintes scripts passo a passo:

\begin{enumerate}
\def\labelenumi{(\alph{enumi})}
\item
  Crie um script \texttt{quadrado.R} com funções para calcular o perímetro e a área de
  um quadrado, dado o comprimento do lado. Noutro script, use \texttt{quadrado.R} para
  calcular o perímetro e a área de um quadrado de lado 5.
\item
  Crie um script \texttt{estatistica.R} com funções para ordenar uma amostra e calcular a
  média, a variância e o desvio padrão de um vetor numérico.
\item
  Crie um programa que use o script \texttt{estatistica.R} para calcular a mediana, a
  variância e o desvio padrão da amostra \texttt{amostra\ \textless{}-\ c(1,\ 2,\ 3,\ 4,\ 5,\ 6,\ 7)}.
\end{enumerate}

\textbf{2.} Crie um script \texttt{circulo.R} com duas funções:

\begin{itemize}
\tightlist
\item
  \texttt{area\_circulo(r)} --- calcula a área de um círculo de raio \(r\);
\item
  \texttt{circunferencia(r)} --- calcula o perímetro (circunferência) de um círculo de raio
  \(r\).
\end{itemize}

Depois, noutro script, use as funções de \texttt{circulo.R} para calcular a área e o
perímetro de um círculo de raio 10.

\textbf{3.} Crie um script \texttt{analise.R} com funções para calcular:

\begin{itemize}
\tightlist
\item
  o valor mínimo (\texttt{minimo}) e máximo (\texttt{maximo}) de um vetor numérico;
\item
  a amplitude (\texttt{amplitude}), definida como a diferença entre o máximo e o mínimo.
\end{itemize}

Depois, num script separado, use \texttt{analise.R} para calcular o mínimo, o máximo e a
amplitude do vetor \texttt{dados\ \textless{}-\ c(8,\ 15,\ 22,\ 3,\ 18,\ 7)}.

\textbf{4.} Crie um script \texttt{categorias.R} com uma função \texttt{contar\_categorias()} que
receba um vetor categórico e devolva uma tabela de frequências de cada categoria.
Use-o depois para analisar o vetor
\texttt{cidades\ \textless{}-\ c("Lisboa",\ "Porto",\ "Lisboa",\ "Coimbra",\ "Porto",\ "Porto",\ "Lisboa")}
e apresente a frequência de cada cidade.

\textbf{5.} Crie um script \texttt{escore\_z.R} com uma função
\texttt{calcular\_escore\_z(valor,\ media,\ desvio\_padrao)} que calcule o valor padronizado
(\emph{escore z}) de um valor, dado por \(z = \dfrac{x - \mu}{\sigma}\). Num script
separado, use-o para calcular o \emph{escore z} dos valores
\texttt{valores\ \textless{}-\ c(5,\ 6,\ 7,\ 8,\ 9)}, assumindo média 7 e desvio padrão 1,5.

\textbf{6.} Crie um script \texttt{binomial.R} com uma função \texttt{probabilidade\_binomial(n,\ p,\ x)}
que calcule a probabilidade de ocorrerem exatamente \(x\) sucessos em \(n\) provas
independentes de um ensaio de Bernoulli:
\[P(X = x) = \binom{n}{x}\, p^{x}\, (1 - p)^{\,n - x},\]
onde \(\binom{n}{x}\) é o coeficiente binomial (em R, \texttt{choose(n,\ x)}) e \(p\) é a
probabilidade de sucesso em cada prova. Num script separado, use-o para calcular a
probabilidade de obter exatamente 3 sucessos em 5 provas, com \(p = 0{,}6\).

Com os scripts guardámos o nosso código de forma permanente e reutilizável. No
próximo capítulo damos o passo seguinte e igualmente decisivo: trazer para dentro do
R \textbf{dados reais}, guardados em ficheiros de texto, CSV ou Excel.

\chapter{Leitura de dados}\label{leitura}

Até aqui os dados com que trabalhámos foram sempre escritos à mão, dentro do próprio
código. Mas a estatística vive de dados \textbf{reais} --- recolhidos em inquéritos,
sensores, registos administrativos ou plataformas online --- e esses dados chegam-nos
guardados em ficheiros. Saber importá-los para o R é, por isso, a ponte entre a
linguagem que aprendemos e a análise que queremos fazer. É neste capítulo que o
livro se torna verdadeiramente aplicado.

Seja qual for o formato, o processo de leitura segue quase sempre os mesmos três
passos:

\begin{enumerate}
\def\labelenumi{\arabic{enumi}.}
\tightlist
\item
  \textbf{Indicar onde está o ficheiro} --- o seu nome ou caminho.
\item
  \textbf{Descrever as suas características} --- separador de campos, separador decimal,
  presença de cabeçalho, codificação, valores em falta.
\item
  \textbf{Ler e guardar os dados num objeto} --- quase sempre uma \emph{data frame}
  (Capítulo~\ref{estruturas}), a estrutura ideal para manipular e analisar
  dados tabulares.
\end{enumerate}

Neste capítulo aprendemos a ler dados a partir do teclado e de ficheiros de texto,
CSV e Excel, com atenção às particularidades dos ficheiros portugueses (vírgula
decimal, acentos), e a guardar resultados de volta em disco.

\section{Ler dados do teclado}\label{ler-dados-do-teclado}

A função \texttt{scan()} lê valores da entrada padrão (o teclado) ou de um ficheiro e
devolve-os num vetor. É útil para introduzir manualmente pequenas quantidades de
dados:

\begin{Shaded}
\begin{Highlighting}[]
\CommentTok{\# Pede ao utilizador uma série de números}
\NormalTok{numeros }\OtherTok{\textless{}{-}} \FunctionTok{scan}\NormalTok{()}
\end{Highlighting}
\end{Shaded}

Depois de escrever \texttt{scan()} e premir \emph{Enter}, introduzem-se os valores separados por
espaços; uma linha vazia (novo \emph{Enter}) termina a entrada.

\section{O diretório de trabalho}\label{o-diretuxf3rio-de-trabalho-1}

Como vimos no Capítulo~\ref{scripts}, o \textbf{diretório de trabalho} é a pasta que o
R usa como referência para ler e guardar ficheiros. Vale a pena recordar as duas
funções que o gerem:

\begin{itemize}
\tightlist
\item
  \texttt{getwd()} --- devolve o caminho do diretório de trabalho atual;
\item
  \texttt{setwd("caminho/da/pasta")} --- define um novo diretório de trabalho.
\end{itemize}

\begin{Shaded}
\begin{Highlighting}[]
\FunctionTok{getwd}\NormalTok{()}
\DocumentationTok{\#\# [1] "C:/Users/Utilizador/Documents"}
\end{Highlighting}
\end{Shaded}

Com o diretório bem definido, se o ficheiro \texttt{dados.csv} estiver nessa pasta, lê-se
simplesmente pelo nome, sem o caminho completo:

\begin{Shaded}
\begin{Highlighting}[]
\NormalTok{dados }\OtherTok{\textless{}{-}} \FunctionTok{read.csv}\NormalTok{(}\StringTok{"dados.csv"}\NormalTok{)}
\FunctionTok{head}\NormalTok{(dados)   }\CommentTok{\# mostra as primeiras linhas}
\end{Highlighting}
\end{Shaded}

No RStudio não é sequer necessário escrever \texttt{setwd()}: o menu \emph{Session \textgreater{} Set Working
Directory \textgreater{} To Source File Location} fixa o diretório na pasta do script. Ainda
melhor, trabalhar dentro de um \textbf{Projeto do RStudio} dispensa por completo caminhos
absolutos.

\section{\texorpdfstring{A função \texttt{read.table()}}{A função read.table()}}\label{a-funuxe7uxe3o-read.table}

A \texttt{read.table()} é a função mais versátil para ler ficheiros de texto em formato
tabular. Importa os dados e permite configurá-los ao pormenor. (Pode descarregar o
ficheiro \texttt{dados.txt} a partir do
\href{https://renatorpaula.wixsite.com/renato/laboratorio-estatistica}{site da disciplina}
para praticar.)

\begin{Shaded}
\begin{Highlighting}[]
\NormalTok{dados }\OtherTok{\textless{}{-}} \FunctionTok{read.table}\NormalTok{(}\AttributeTok{file =} \StringTok{"dados.txt"}\NormalTok{, }\AttributeTok{header =} \ConstantTok{TRUE}\NormalTok{, }\AttributeTok{sep =} \StringTok{""}\NormalTok{)}
\FunctionTok{print}\NormalTok{(dados)}
\end{Highlighting}
\end{Shaded}

Os seus argumentos mais usados são:

\begin{itemize}
\tightlist
\item
  \texttt{file} --- nome ou caminho do ficheiro a ler.
\item
  \texttt{header} --- indica se a primeira linha contém os nomes das colunas
  (\texttt{header\ =\ TRUE}). \textbf{Por omissão é \texttt{FALSE}}.
\item
  \texttt{sep} --- o símbolo que separa os campos. \textbf{Por omissão é o espaço em branco}
  (\texttt{sep\ =\ ""}, que trata qualquer sequência de espaços ou tabulações como
  separador); pode ser \texttt{";"}, \texttt{","}, \texttt{"\textbackslash{}t"} (tabulação), etc.
\item
  \texttt{dec} --- o símbolo que separa a parte inteira da decimal (\texttt{"."} ou \texttt{","}).
\item
  \texttt{nrows} --- número máximo de linhas a ler.
\item
  \texttt{na.strings} --- que valores no ficheiro devem ser interpretados como \texttt{NA} (por
  exemplo, \texttt{"NA"}, \texttt{"N/A"}).
\item
  \texttt{skip} --- número de linhas a ignorar no início do ficheiro.
\item
  \texttt{comment.char} --- símbolo que marca comentários; linhas iniciadas por ele são
  ignoradas.
\end{itemize}

É frequente confundir os separadores por omissão. Ao contrário do que se poderia
pensar, o separador por omissão de \texttt{read.table()} \textbf{não é a vírgula, mas sim o
espaço em branco}. A vírgula é o separador por omissão da \texttt{read.csv()}, que veremos
já a seguir. Confundir os dois é uma das causas mais comuns de leituras que ``correm
bem'' mas produzem uma única coluna com tudo lá dentro.

Para consultar todos os argumentos e os seus valores por omissão:

\begin{Shaded}
\begin{Highlighting}[]
\FunctionTok{args}\NormalTok{(read.table)}
\DocumentationTok{\#\# function (file, header = FALSE, sep = "", quote = "\textbackslash{}"\textquotesingle{}", dec = ".",}
\DocumentationTok{\#\#     numerals = c("allow.loss", "warn.loss", "no.loss"), row.names,}
\DocumentationTok{\#\#     col.names, as.is = !stringsAsFactors, na.strings = "NA", ...)}
\DocumentationTok{\#\# NULL}
\end{Highlighting}
\end{Shaded}

Se o ficheiro for um CSV, basta indicar o separador apropriado:

\begin{Shaded}
\begin{Highlighting}[]
\NormalTok{dados\_csv }\OtherTok{\textless{}{-}} \FunctionTok{read.table}\NormalTok{(}\StringTok{"dados.csv"}\NormalTok{, }\AttributeTok{header =} \ConstantTok{TRUE}\NormalTok{, }\AttributeTok{sep =} \StringTok{","}\NormalTok{)}
\end{Highlighting}
\end{Shaded}

\section{\texorpdfstring{A função \texttt{read.csv()}}{A função read.csv()}}\label{a-funuxe7uxe3o-read.csv}

A \texttt{read.csv()} está otimizada para ficheiros CSV (\emph{Comma-Separated Values}),
ficheiros de texto em que cada linha é um registo e os valores são separados por
vírgulas. Na prática, é a \texttt{read.table()} com outros valores por omissão: separador
\texttt{","} e \texttt{header\ =\ TRUE}.

\begin{Shaded}
\begin{Highlighting}[]
\NormalTok{dados\_csv }\OtherTok{\textless{}{-}} \FunctionTok{read.csv}\NormalTok{(}\StringTok{"dados.csv"}\NormalTok{, }\AttributeTok{header =} \ConstantTok{TRUE}\NormalTok{)}
\end{Highlighting}
\end{Shaded}

\section{\texorpdfstring{A função \texttt{read.csv2()}}{A função read.csv2()}}\label{a-funuxe7uxe3o-read.csv2}

A \texttt{read.csv2()} destina-se a ficheiros CSV ``à europeia'': separador de campos \textbf{ponto
e vírgula} (\texttt{;}) e separador decimal \textbf{vírgula} (\texttt{,}). É o formato que a maioria
dos programas exporta em Portugal, precisamente porque cá a vírgula é o separador
decimal.

\begin{Shaded}
\begin{Highlighting}[]
\NormalTok{dados\_csv2 }\OtherTok{\textless{}{-}} \FunctionTok{read.csv2}\NormalTok{(}\StringTok{"dados.csv"}\NormalTok{, }\AttributeTok{header =} \ConstantTok{TRUE}\NormalTok{)}
\end{Highlighting}
\end{Shaded}

Esta distinção é essencial em Portugal. Um valor como \texttt{1.234,56} num ficheiro
português significa ``mil duzentos e trinta e quatro vírgula cinquenta e seis''. Se o
lermos com \texttt{read.csv()} (que espera a vírgula como separador de campos e o ponto como
decimal), o R interpreta tudo mal. A regra prática: ficheiros gerados por sistemas
portugueses leem-se, quase sempre, com \texttt{read.csv2()}.

\section{\texorpdfstring{A função \texttt{read\_excel()}}{A função read\_excel()}}\label{a-funuxe7uxe3o-read_excel}

Para ler ficheiros Excel (\texttt{.xls} e \texttt{.xlsx}) usamos a função \texttt{read\_excel()} do pacote
\texttt{readxl}, que precisa de ser instalado e carregado antes da primeira utilização:

\begin{Shaded}
\begin{Highlighting}[]
\FunctionTok{install.packages}\NormalTok{(}\StringTok{"readxl"}\NormalTok{)   }\CommentTok{\# apenas uma vez}
\FunctionTok{library}\NormalTok{(readxl)}

\CommentTok{\# Sintaxe básica}
\FunctionTok{read\_excel}\NormalTok{(path, }\AttributeTok{sheet =} \ConstantTok{NULL}\NormalTok{, }\AttributeTok{range =} \ConstantTok{NULL}\NormalTok{, }\AttributeTok{col\_names =} \ConstantTok{TRUE}\NormalTok{,}
           \AttributeTok{col\_types =} \ConstantTok{NULL}\NormalTok{, }\AttributeTok{na =} \StringTok{""}\NormalTok{, }\AttributeTok{trim\_ws =} \ConstantTok{TRUE}\NormalTok{, }\AttributeTok{skip =} \DecValTok{0}\NormalTok{, }\AttributeTok{n\_max =} \ConstantTok{Inf}\NormalTok{)}
\end{Highlighting}
\end{Shaded}

Os seus parâmetros principais são:

\begin{itemize}
\tightlist
\item
  \texttt{path} --- caminho do ficheiro Excel (obrigatório); pode ser local ou um URL.
\item
  \texttt{sheet} --- nome ou número da folha a ler. Por omissão, lê a primeira.
\item
  \texttt{range} --- intervalo específico a ler, como \texttt{"A1:D10"}. Por omissão (\texttt{NULL}), lê a
  folha inteira.
\item
  \texttt{col\_names} --- se a primeira linha deve ser usada como nomes das colunas
  (\texttt{TRUE}, por omissão).
\item
  \texttt{col\_types} --- tipos das colunas (\texttt{"numeric"}, \texttt{"date"}, \texttt{"text"}, \texttt{"guess"}, \ldots).
  Por omissão (\texttt{NULL}), a função adivinha.
\item
  \texttt{na} --- cadeias a interpretar como valores em falta (\texttt{NA}).
\item
  \texttt{trim\_ws} --- se remove espaços no início e no fim dos valores de texto (\texttt{TRUE}, por
  omissão).
\item
  \texttt{skip}, \texttt{n\_max} --- linhas a ignorar no início e número máximo de linhas a ler.
\end{itemize}

Suponha um ficheiro \texttt{instagram.xlsx}. Para ler a folha \texttt{Sheet1} e guardá-la numa
\emph{data frame} chamada \texttt{dados\_instagram}:

\begin{Shaded}
\begin{Highlighting}[]
\NormalTok{dados\_instagram }\OtherTok{\textless{}{-}} \FunctionTok{read\_excel}\NormalTok{(}\AttributeTok{path =} \StringTok{"instagram.xlsx"}\NormalTok{,}
                              \AttributeTok{sheet =} \StringTok{"Sheet1"}\NormalTok{,}
                              \AttributeTok{col\_names =} \ConstantTok{TRUE}\NormalTok{)}
\end{Highlighting}
\end{Shaded}

No RStudio, o botão \textbf{Import Dataset} (separador \emph{Environment}) permite importar
ficheiros de texto e Excel através de uma janela, mostrando uma pré-visualização e
gerando automaticamente o código de leitura. É uma excelente forma de acertar os
argumentos (separador, decimal, codificação) sem tentativa e erro.

\section{Ler dados a partir de um URL}\label{ler-dados-a-partir-de-um-url}

O R também lê dados diretamente da internet. Basta passar o endereço do ficheiro à
função de leitura apropriada ao formato:

\begin{Shaded}
\begin{Highlighting}[]
\NormalTok{url }\OtherTok{\textless{}{-}} \StringTok{"https://exemplo.com/dados.csv"}
\NormalTok{dados\_online }\OtherTok{\textless{}{-}} \FunctionTok{read.table}\NormalTok{(url, }\AttributeTok{header =} \ConstantTok{TRUE}\NormalTok{, }\AttributeTok{sep =} \StringTok{","}\NormalTok{)}
\end{Highlighting}
\end{Shaded}

\section{Guardar resultados em disco}\label{guardar-resultados-em-disco}

Ler dados é só metade da história: depois de os analisar, queremos muitas vezes
\textbf{guardar} o resultado. As funções \texttt{write.table()} e \texttt{write.csv()} fazem o percurso
inverso das que vimos, escrevendo uma \emph{data frame} num ficheiro de texto:

\begin{Shaded}
\begin{Highlighting}[]
\FunctionTok{write.csv}\NormalTok{(dados, }\AttributeTok{file =} \StringTok{"resultado.csv"}\NormalTok{, }\AttributeTok{row.names =} \ConstantTok{FALSE}\NormalTok{)}
\end{Highlighting}
\end{Shaded}

O argumento \texttt{row.names\ =\ FALSE} evita que o R acrescente uma coluna extra com os
números das linhas --- um pormenor pequeno, mas que poupa confusão a quem voltar a ler
o ficheiro.

Os acentos do português dependem da \textbf{codificação} do ficheiro. Se, ao ler, os
caracteres acentuados aparecerem trocados (por exemplo, \texttt{ç} em vez de \texttt{ç}),
especifique a codificação com o argumento \texttt{fileEncoding}, tipicamente
\texttt{fileEncoding\ =\ "UTF-8"} ou \texttt{fileEncoding\ =\ "latin1"}, conforme a origem do ficheiro.

\section{Exercícios}\label{exercuxedcios-9}

\textbf{1.} Um instituto de climatologia registou as temperaturas mínima e máxima durante
3 dias. Os dados estão no ficheiro
\href{docs/temperaturas.txt}{temperaturas.txt}, com separador de tabulação e cabeçalho.

\begin{enumerate}
\def\labelenumi{(\alph{enumi})}
\item
  Leia \texttt{temperaturas.txt} com \texttt{read.table()}.
\item
  Veja as 6 primeiras linhas com \texttt{head()}.
\item
  Qual foi a temperatura máxima no segundo dia?
\item
  Crie uma coluna \texttt{diferenca\_temp} com a diferença entre a temperatura máxima e a
  mínima e guarde o resultado num ficheiro \texttt{temperaturas\_nova.txt}. \emph{Sugestão:}
  \texttt{write.table(x\ =\ ,\ file\ =\ ,\ row.names\ =\ )}.
\end{enumerate}

\textbf{2.} Um ficheiro \texttt{.csv} exportado de um sistema português contém despesas e
receitas mensais: \href{docs/financas.csv}{financas.csv}. Os valores usam vírgula como
separador decimal e ponto e vírgula como separador de campos.

\begin{enumerate}
\def\labelenumi{(\alph{enumi})}
\item
  Importe os dados com \texttt{read.csv2()}.
\item
  Calcule a soma total das receitas.
\item
  Crie uma coluna \texttt{saldo\ =\ receitas\ -\ despesas}.
\item
  Guarde o resultado num ficheiro \texttt{financas\_nova.csv}. \emph{Sugestão:}
  \texttt{write.csv2(x\ =\ ,\ file\ =\ ,\ row.names\ =\ )}.
\end{enumerate}

\textbf{3.} Um laboratório gerou um ficheiro sem cabeçalho com 4 colunas: experimento,
tempo, concentração e erro, separadas por espaços:
\href{docs/resultados_lab.txt}{resultados\_lab.txt}.

\begin{enumerate}
\def\labelenumi{(\alph{enumi})}
\item
  Leia os dados com \texttt{read.table(...,\ header\ =\ FALSE)}.
\item
  Atribua nomes às colunas: \texttt{"experimento"}, \texttt{"tempo"}, \texttt{"concentracao"},
  \texttt{"erro"}.
\item
  Qual a concentração do 5.º experimento?
\end{enumerate}

\textbf{4.} Um ficheiro \texttt{.csv} contém uma entrada incorreta na coluna \texttt{salario}
(``quatro mil''), o que impede que a coluna seja lida como numérica:
\href{docs/dados_com_erros.csv}{dados\_com\_erros.csv}.

\begin{enumerate}
\def\labelenumi{(\alph{enumi})}
\item
  Importe com \texttt{read.csv()} e verifique a classe da coluna \texttt{salario}.
\item
  Corrija os dados e converta \texttt{salario} em \texttt{numeric}.
\end{enumerate}

\textbf{5.} Leia o ficheiro \texttt{instagram.xlsx}
(\href{https://renatorpaula.wixsite.com/renato/laboratorio-estatistica}{aqui}; lembre-se
de instalar o pacote \texttt{readxl}) e responda:

\begin{enumerate}
\def\labelenumi{(\alph{enumi})}
\item
  Mostre as 6 primeiras linhas do conjunto de dados.
\item
  \textbf{Tabela de frequências.} Construa uma tabela de frequências para cada variável.
  Para as variáveis numéricas, agrupe em classes. O código seguinte fá-lo para o
  \texttt{indice\_influencia}, usando a Regra de Sturges para o número de classes:
\end{enumerate}

\begin{Shaded}
\begin{Highlighting}[]
\CommentTok{\# {-}{-}{-}{-}{-}{-}{-}{-}{-}{-}{-}{-}{-}{-}{-}{-}{-}{-}{-}{-}{-}{-}{-}{-}{-}{-}{-}{-}{-}{-}{-}{-}{-}}
\CommentTok{\# Índice de influência}
\CommentTok{\# {-}{-}{-}{-}{-}{-}{-}{-}{-}{-}{-}{-}{-}{-}{-}{-}{-}{-}{-}{-}{-}{-}{-}{-}{-}{-}{-}{-}{-}{-}{-}{-}{-}}

\CommentTok{\# Número de classes pela Regra de Sturges}
\NormalTok{n }\OtherTok{\textless{}{-}} \FunctionTok{length}\NormalTok{(dados\_instagram}\SpecialCharTok{$}\NormalTok{indice\_influencia)     }\CommentTok{\# número de observações}
\NormalTok{k }\OtherTok{\textless{}{-}} \FunctionTok{ceiling}\NormalTok{(}\DecValTok{1} \SpecialCharTok{+} \FloatTok{3.322} \SpecialCharTok{*} \FunctionTok{log10}\NormalTok{(n))                 }\CommentTok{\# número de classes (arredondado)}

\CommentTok{\# Limites das classes}
\NormalTok{min\_val }\OtherTok{\textless{}{-}} \FunctionTok{min}\NormalTok{(dados\_instagram}\SpecialCharTok{$}\NormalTok{indice\_influencia)}
\NormalTok{max\_val }\OtherTok{\textless{}{-}} \FunctionTok{max}\NormalTok{(dados\_instagram}\SpecialCharTok{$}\NormalTok{indice\_influencia)}
\NormalTok{intervalos }\OtherTok{\textless{}{-}} \FunctionTok{seq}\NormalTok{(min\_val, max\_val, }\AttributeTok{length.out =}\NormalTok{ k }\SpecialCharTok{+} \DecValTok{1}\NormalTok{)}

\CommentTok{\# Tabela de frequências absolutas}
\NormalTok{tabela\_if }\OtherTok{\textless{}{-}} \FunctionTok{cut}\NormalTok{(dados\_instagram}\SpecialCharTok{$}\NormalTok{indice\_influencia,}
                 \AttributeTok{breaks =}\NormalTok{ intervalos, }\AttributeTok{include.lowest =} \ConstantTok{TRUE}\NormalTok{)}
\NormalTok{tabela\_if }\OtherTok{\textless{}{-}} \FunctionTok{table}\NormalTok{(tabela\_if)}

\CommentTok{\# Converter em data frame para uma leitura mais organizada}
\NormalTok{tabela\_if\_df }\OtherTok{\textless{}{-}} \FunctionTok{as.data.frame}\NormalTok{(tabela\_if)}
\FunctionTok{names}\NormalTok{(tabela\_if\_df) }\OtherTok{\textless{}{-}} \FunctionTok{c}\NormalTok{(}\StringTok{"Intervalo"}\NormalTok{, }\StringTok{"Frequencia"}\NormalTok{)}

\CommentTok{\# Frequência relativa}
\NormalTok{tabela\_if\_df}\SpecialCharTok{$}\NormalTok{Freq\_Relativa }\OtherTok{\textless{}{-}} \FunctionTok{round}\NormalTok{(tabela\_if\_df}\SpecialCharTok{$}\NormalTok{Frequencia }\SpecialCharTok{/}\NormalTok{ n, }\AttributeTok{digits =} \DecValTok{3}\NormalTok{)}

\CommentTok{\# Frequências acumuladas (absoluta e relativa)}
\NormalTok{tabela\_if\_df}\SpecialCharTok{$}\NormalTok{Freq\_Abs\_Acum  }\OtherTok{\textless{}{-}} \FunctionTok{cumsum}\NormalTok{(tabela\_if\_df}\SpecialCharTok{$}\NormalTok{Frequencia)}
\NormalTok{tabela\_if\_df}\SpecialCharTok{$}\NormalTok{Freq\_Rela\_Acum }\OtherTok{\textless{}{-}} \FunctionTok{cumsum}\NormalTok{(tabela\_if\_df}\SpecialCharTok{$}\NormalTok{Freq\_Relativa)}

\NormalTok{tabela\_if\_df}
\end{Highlighting}
\end{Shaded}

\begin{enumerate}
\def\labelenumi{(\alph{enumi})}
\setcounter{enumi}{2}
\item
  \textbf{Representação gráfica.} Crie o gráfico adequado a cada variável: histogramas
  (\texttt{hist()}) para variáveis contínuas, gráficos de barras (\texttt{barplot()}) para
  variáveis categóricas, e caixas de bigodes (\texttt{boxplot()}) para observar a
  dispersão.
\item
  \textbf{Medidas descritivas.} Determine, para cada variável, a média, a mediana, a
  variância, o desvio padrão, o coeficiente de variação, o primeiro e o terceiro
  quartis (\(Q_1\) e \(Q_3\)), a amplitude e a amplitude interquartil. Indique quando
  alguma destas medidas não se aplica.
\item
  Calcule a média e a mediana do índice de influência \textbf{por género}. Comente as
  diferenças e discuta se sugerem tendências distintas de influência.
\item
  Calcule a média e a mediana da média de gostos por género. Avalie se há
  diferenças relevantes e discuta possíveis implicações.
\item
  \textbf{Comparação de dispersão.} Compare a dispersão do índice de influência e do
  número de publicações usando a variância e o desvio padrão. Justifique qual
  apresenta maior dispersão e interprete o resultado.
\item
  \textbf{Percentis.} Identifique os percentis 10\%, 25\%, 75\% e 90\% do número de
  seguidores. Comente como podem servir para categorizar influenciadores por nível
  de popularidade.
\item
  \textbf{Outliers.} Investigue a presença de valores atípicos nas variáveis número de
  publicações, número de seguidores e média de gostos.
\item
  \textbf{Correlação.} Calcule a correlação entre o índice de influência e a média de
  gostos. Comente se um índice de influência elevado tende a associar-se a mais
  gostos. \emph{Sugestão:} Use a função \texttt{cor()}.
\end{enumerate}

\textbf{6.} Importe o ficheiro \texttt{dataset.csv}
(\href{https://renatorpaula.wixsite.com/renato/laboratorio-estatistica}{aqui}) e
guarde-o numa \emph{data frame}. Em seguida:

\begin{enumerate}
\def\labelenumi{(\alph{enumi})}
\item
  Mostre as 6 primeiras linhas.
\item
  Escolha uma variável categórica, construa a sua tabela de frequências e
  represente a distribuição num gráfico de barras.
\item
  Para cada variável numérica, crie um histograma (distribuição) e uma caixa de
  bigodes (dispersão e outliers).
\item
  Calcule as principais medidas descritivas de uma variável numérica à sua
  escolha: média, mediana, variância, desvio padrão, coeficiente de variação,
  quartis (\(Q_1\) e \(Q_3\)), amplitude e amplitude interquartil.
\item
  Escolha uma variável numérica e uma categórica e compare a numérica entre os
  grupos da categórica:
\end{enumerate}

\begin{itemize}
\tightlist
\item
  calcule a média e a mediana da variável numérica em cada grupo;
\item
  use caixas de bigodes para comparar visualmente as distribuições.
\end{itemize}

Com os dados dentro do R, o passo seguinte é transformá-los e resumi-los. No próximo
capítulo conhecemos o operador \emph{pipe}, que torna essas sequências de operações
muito mais legíveis.

\chapter{\texorpdfstring{O operador \emph{pipe}}{O operador pipe}}\label{pipe}

Depois de trazer os dados para o R (Capítulo~\ref{leitura}), raramente lhes
aplicamos uma única operação. O mais comum é encadear várias --- filtrar, transformar,
resumir --- umas a seguir às outras. Escrever esse encadeamento de forma legível é o
problema que o operador \emph{pipe} resolve com elegância.

Considere uma tarefa simples: calcular a raiz quadrada da soma de um vetor. Sem
\emph{pipe}, escrevemos as funções aninhadas umas dentro das outras:

\begin{Shaded}
\begin{Highlighting}[]
\NormalTok{x }\OtherTok{\textless{}{-}} \FunctionTok{c}\NormalTok{(}\DecValTok{1}\NormalTok{, }\DecValTok{2}\NormalTok{, }\DecValTok{3}\NormalTok{, }\DecValTok{4}\NormalTok{)}
\FunctionTok{sqrt}\NormalTok{(}\FunctionTok{sum}\NormalTok{(x))}
\DocumentationTok{\#\# [1] 3.16}
\end{Highlighting}
\end{Shaded}

O problema é que lemos este código \textbf{de dentro para fora}: primeiro \texttt{sum}, depois
\texttt{sqrt}, embora a escrita seja a ordem inversa. Com duas funções ainda se gere; com
cinco, torna-se um quebra-cabeças. O \emph{pipe} permite escrever e ler as operações pela
ordem em que acontecem, \textbf{da esquerda para a direita}.

\section{\texorpdfstring{O \emph{pipe} do \texttt{magrittr} (\texttt{\%\textgreater{}\%})}{O pipe do magrittr (\%\textgreater\%)}}\label{o-pipe-do-magrittr}

O \emph{pipe} foi popularizado pelo pacote \texttt{magrittr}, de Stefan Milton Bache (2014). O
nome é uma homenagem ao pintor surrealista René Magritte e à sua obra \emph{Ceci n'est pas
une pipe} (``Isto não é um cachimbo'') --- um trocadilho com a ideia de fazer os dados
``fluírem'' por uma sequência de operações.

\begin{Shaded}
\begin{Highlighting}[]
\FunctionTok{install.packages}\NormalTok{(}\StringTok{"magrittr"}\NormalTok{)}
\FunctionTok{library}\NormalTok{(magrittr)}
\end{Highlighting}
\end{Shaded}

O operador \texttt{\%\textgreater{}\%} pega no \textbf{valor à sua esquerda} e passa-o como \textbf{primeiro
argumento} da função à sua direita. O exemplo anterior reescreve-se assim:

\begin{Shaded}
\begin{Highlighting}[]
\FunctionTok{library}\NormalTok{(magrittr)}
\NormalTok{x }\SpecialCharTok{\%\textgreater{}\%} \FunctionTok{sum}\NormalTok{() }\SpecialCharTok{\%\textgreater{}\%} \FunctionTok{sqrt}\NormalTok{()}
\DocumentationTok{\#\# [1] 3.16}
\end{Highlighting}
\end{Shaded}

Lê-se de forma quase literal: ``pega em \texttt{x}, soma, e depois tira a raiz quadrada''. A
vantagem cresce com o número de passos. Para calcular a raiz quadrada da soma dos
logaritmos dos valores absolutos de um vetor:

\begin{Shaded}
\begin{Highlighting}[]
\CommentTok{\# Sem pipe: lê{-}se de dentro para fora}
\NormalTok{resultado }\OtherTok{\textless{}{-}} \FunctionTok{sqrt}\NormalTok{(}\FunctionTok{sum}\NormalTok{(}\FunctionTok{log}\NormalTok{(}\FunctionTok{abs}\NormalTok{(x))))}

\CommentTok{\# Com pipe: lê{-}se na ordem das operações}
\NormalTok{resultado }\OtherTok{\textless{}{-}}\NormalTok{ x }\SpecialCharTok{\%\textgreater{}\%}
  \FunctionTok{abs}\NormalTok{() }\SpecialCharTok{\%\textgreater{}\%}
  \FunctionTok{log}\NormalTok{() }\SpecialCharTok{\%\textgreater{}\%}
  \FunctionTok{sum}\NormalTok{() }\SpecialCharTok{\%\textgreater{}\%}
  \FunctionTok{sqrt}\NormalTok{()}
\end{Highlighting}
\end{Shaded}

\subsection{\texorpdfstring{O ponto (\texttt{.}) como marcador de posição}{O ponto (.) como marcador de posição}}\label{o-ponto-.-como-marcador-de-posiuxe7uxe3o}

Por vezes queremos passar o valor da esquerda para um argumento que \textbf{não é o
primeiro}. Nesses casos, usamos um ponto (\texttt{.}) para marcar a posição exata onde o
valor deve entrar.

Por exemplo, na função \texttt{lm()}, que ajusta um modelo linear, os dados vão no argumento
\texttt{data\ =} (o segundo). Usamos o \texttt{.} para lá os colocar:

\begin{Shaded}
\begin{Highlighting}[]
\NormalTok{airquality }\SpecialCharTok{\%\textgreater{}\%}
  \FunctionTok{na.omit}\NormalTok{() }\SpecialCharTok{\%\textgreater{}\%}                                     \CommentTok{\# remove valores em falta}
  \FunctionTok{lm}\NormalTok{(Ozone }\SpecialCharTok{\textasciitilde{}}\NormalTok{ Wind }\SpecialCharTok{+}\NormalTok{ Temp }\SpecialCharTok{+}\NormalTok{ Solar.R, }\AttributeTok{data =}\NormalTok{ .) }\SpecialCharTok{\%\textgreater{}\%}   \CommentTok{\# ajusta o modelo linear}
  \FunctionTok{summary}\NormalTok{()                                         }\CommentTok{\# resume o modelo}
\DocumentationTok{\#\# }
\DocumentationTok{\#\# Coefficients:}
\DocumentationTok{\#\#             Estimate Std. Error t value Pr(\textgreater{}|t|)    }
\DocumentationTok{\#\# (Intercept) {-}64.3421    23.0547   {-}2.79   0.0062 ** }
\DocumentationTok{\#\# Wind         {-}3.3336     0.6544   {-}5.09  1.5e{-}06 ***}
\DocumentationTok{\#\# Temp          1.6521     0.2535    6.52  2.4e{-}09 ***}
\DocumentationTok{\#\# Solar.R       0.0598     0.0232    2.58   0.0112 *  }
\DocumentationTok{\#\# {-}{-}{-}}
\DocumentationTok{\#\# Residual standard error: 21.2 on 107 degrees of freedom}
\DocumentationTok{\#\# Multiple R{-}squared:  0.606,  Adjusted R{-}squared:  0.595}
\end{Highlighting}
\end{Shaded}

O conjunto \texttt{airquality} é limpo de valores em falta com \texttt{na.omit()}, o resultado
entra como \texttt{data} em \texttt{lm()}, e o modelo é finalmente resumido --- toda a cadeia se lê
de uma só vez, de cima para baixo.

\section{\texorpdfstring{O \emph{pipe} nativo do R (\texttt{\textbar{}\textgreater{}})}{O pipe nativo do R (\textbar\textgreater)}}\label{o-pipe-nativo-do-r}

Desde a versão 4.1 (2021), o próprio R passou a incluir um \emph{pipe} nativo, escrito
\texttt{\textbar{}\textgreater{}}, que \textbf{não exige nenhum pacote}. O seu comportamento essencial é o mesmo:
passa o valor da esquerda como primeiro argumento da função à direita.

\begin{Shaded}
\begin{Highlighting}[]
\NormalTok{x }\SpecialCharTok{|\textgreater{}} \FunctionTok{sum}\NormalTok{() }\SpecialCharTok{|\textgreater{}} \FunctionTok{sqrt}\NormalTok{()}
\DocumentationTok{\#\# [1] 3.16}
\end{Highlighting}
\end{Shaded}

Como faz parte da base do R e é hoje a forma recomendada em código novo, vale a pena
conhecer as duas diferenças principais em relação ao \texttt{\%\textgreater{}\%}:

\begin{itemize}
\tightlist
\item
  O \texttt{\textbar{}\textgreater{}} \textbf{exige sempre os parênteses} da função: escreve-se \texttt{x\ \textbar{}\textgreater{}\ sqrt()}, e não
  \texttt{x\ \textbar{}\textgreater{}\ sqrt}.
\item
  Para colocar o valor num argumento que não o primeiro, o \texttt{\textbar{}\textgreater{}} usa o marcador \texttt{\_}
  (sublinhado), em vez do \texttt{.}, e esse argumento tem de ser \textbf{nomeado} (a partir da
  versão 4.2). O exemplo do modelo linear fica:
\end{itemize}

\begin{Shaded}
\begin{Highlighting}[]
\NormalTok{airquality }\SpecialCharTok{|\textgreater{}}
  \FunctionTok{na.omit}\NormalTok{() }\SpecialCharTok{|\textgreater{}}
  \FunctionTok{lm}\NormalTok{(Ozone }\SpecialCharTok{\textasciitilde{}}\NormalTok{ Wind }\SpecialCharTok{+}\NormalTok{ Temp }\SpecialCharTok{+}\NormalTok{ Solar.R, }\AttributeTok{data =}\NormalTok{ \_) }\SpecialCharTok{|\textgreater{}}
  \FunctionTok{summary}\NormalTok{()}
\end{Highlighting}
\end{Shaded}

Qual usar? Para código novo e simples, o \texttt{\textbar{}\textgreater{}} é a escolha natural: é nativo, não
depende de pacotes e será compreendido por qualquer instalação recente do R. O \texttt{\%\textgreater{}\%}
do \texttt{magrittr} continua indispensável quando se usam os verbos do \texttt{dplyr} (como
veremos) ou os seus atalhos (\texttt{add()}, \texttt{extract()}, \texttt{divide\_by()}, \ldots), e oferece mais
flexibilidade com o marcador \texttt{.}. Neste capítulo, os exercícios que recorrem a essas
funções auxiliares usam, por isso, o \texttt{\%\textgreater{}\%}.

No RStudio, o atalho \texttt{Ctrl+Shift+M} / \texttt{Cmd+Shift+M} insere o operador \emph{pipe}. Em
\emph{Tools \textgreater{} Global Options \textgreater{} Code} pode escolher se esse atalho insere \texttt{\%\textgreater{}\%} ou \texttt{\textbar{}\textgreater{}}.

\section{Exercícios}\label{exercuxedcios-10}

\textbf{1.} Reescreva a expressão seguinte usando o \texttt{\%\textgreater{}\%}.

\begin{Shaded}
\begin{Highlighting}[]
\FunctionTok{round}\NormalTok{(}\FunctionTok{mean}\NormalTok{(}\FunctionTok{sum}\NormalTok{(}\DecValTok{1}\SpecialCharTok{:}\DecValTok{10}\NormalTok{) }\SpecialCharTok{/} \DecValTok{3}\NormalTok{), }\AttributeTok{digits =} \DecValTok{1}\NormalTok{)}
\end{Highlighting}
\end{Shaded}

\emph{Sugestão:} use a função \texttt{magrittr::divide\_by()}. Consulte a ajuda da função.

\textbf{2.} Reescreva o código seguinte usando o \texttt{\%\textgreater{}\%}.

\begin{Shaded}
\begin{Highlighting}[]
\NormalTok{x }\OtherTok{\textless{}{-}} \FunctionTok{rnorm}\NormalTok{(}\DecValTok{100}\NormalTok{)}
\NormalTok{x.pos }\OtherTok{\textless{}{-}}\NormalTok{ x[x }\SpecialCharTok{\textgreater{}} \DecValTok{0}\NormalTok{]}
\NormalTok{media }\OtherTok{\textless{}{-}} \FunctionTok{mean}\NormalTok{(x.pos)}
\NormalTok{saida }\OtherTok{\textless{}{-}} \FunctionTok{round}\NormalTok{(media, }\DecValTok{2}\NormalTok{)}
\end{Highlighting}
\end{Shaded}

\emph{Sugestão:} use a função \texttt{magrittr::extract()}. Consulte a ajuda da função.

\textbf{3.} Sem executar, indique qual é o resultado do código seguinte. Consulte a ajuda
das funções, se precisar.

\begin{Shaded}
\begin{Highlighting}[]
\DecValTok{2} \SpecialCharTok{\%\textgreater{}\%}
  \FunctionTok{add}\NormalTok{(}\DecValTok{2}\NormalTok{) }\SpecialCharTok{\%\textgreater{}\%}
  \FunctionTok{c}\NormalTok{(}\DecValTok{6}\NormalTok{, }\ConstantTok{NA}\NormalTok{) }\SpecialCharTok{\%\textgreater{}\%}
  \FunctionTok{mean}\NormalTok{(}\AttributeTok{na.rm =} \ConstantTok{TRUE}\NormalTok{) }\SpecialCharTok{\%\textgreater{}\%}
  \FunctionTok{equals}\NormalTok{(}\DecValTok{5}\NormalTok{)}
\end{Highlighting}
\end{Shaded}

\textbf{4.} No conjunto \texttt{mtcars}, selecione as colunas \texttt{mpg}, \texttt{hp} e \texttt{wt}, filtre os
carros com mais de 100 cavalos (\texttt{hp\ \textgreater{}\ 100}) e ordene-os por \texttt{wt} (peso) de forma
crescente. Reescreva o código seguinte usando o \emph{pipe}:

\begin{Shaded}
\begin{Highlighting}[]
\FunctionTok{install.packages}\NormalTok{(}\StringTok{"dplyr"}\NormalTok{)}
\FunctionTok{library}\NormalTok{(dplyr)}
\NormalTok{result }\OtherTok{\textless{}{-}} \FunctionTok{arrange}\NormalTok{(}\FunctionTok{filter}\NormalTok{(}\FunctionTok{select}\NormalTok{(mtcars, mpg, hp, wt), hp }\SpecialCharTok{\textgreater{}} \DecValTok{100}\NormalTok{), wt)}
\FunctionTok{print}\NormalTok{(result)}
\end{Highlighting}
\end{Shaded}

\textbf{5.} No conjunto \texttt{iris}, filtre as linhas com \texttt{Sepal.Length\ \textgreater{}\ 5}, selecione as
colunas \texttt{Sepal.Length} e \texttt{Species} e calcule a média de \texttt{Sepal.Length} por espécie.
Reescreva o código seguinte usando o \emph{pipe}:

\begin{Shaded}
\begin{Highlighting}[]
\NormalTok{result }\OtherTok{\textless{}{-}} \FunctionTok{aggregate}\NormalTok{(Sepal.Length }\SpecialCharTok{\textasciitilde{}}\NormalTok{ Species,}
                    \AttributeTok{data =} \FunctionTok{select}\NormalTok{(}\FunctionTok{filter}\NormalTok{(iris, Sepal.Length }\SpecialCharTok{\textgreater{}} \DecValTok{5}\NormalTok{),}
\NormalTok{                                  Sepal.Length, Species),}
\NormalTok{                    mean)}
\FunctionTok{print}\NormalTok{(result)}
\end{Highlighting}
\end{Shaded}

\textbf{6.} No conjunto \texttt{airquality}, remova as linhas com valores em falta, selecione as
colunas \texttt{Ozone} e \texttt{Wind} e calcule a média de \texttt{Ozone} nas observações em que
\texttt{Wind\ \textgreater{}\ 10}. Reescreva o código seguinte usando o \emph{pipe}.
\emph{Sugestão:} Use a função \texttt{summarise()} para calcular a média.

\begin{Shaded}
\begin{Highlighting}[]
\NormalTok{cleaned\_data }\OtherTok{\textless{}{-}} \FunctionTok{na.omit}\NormalTok{(airquality)}
\NormalTok{filtered\_data }\OtherTok{\textless{}{-}} \FunctionTok{filter}\NormalTok{(cleaned\_data, Wind }\SpecialCharTok{\textgreater{}} \DecValTok{10}\NormalTok{)}
\NormalTok{result }\OtherTok{\textless{}{-}} \FunctionTok{mean}\NormalTok{(filtered\_data}\SpecialCharTok{$}\NormalTok{Ozone)}
\FunctionTok{print}\NormalTok{(result)}
\end{Highlighting}
\end{Shaded}

\textbf{7.} No conjunto \texttt{mtcars}, crie uma coluna \texttt{hp\_per\_wt} igual à razão entre \texttt{hp} e
\texttt{wt}, filtre os carros com essa razão maior que 30 e selecione as colunas \texttt{mpg},
\texttt{hp\_per\_wt} e \texttt{gear}. Reescreva o código seguinte usando o \emph{pipe}.
\emph{Sugestão:} Use a função \texttt{mutate()} para criar a nova coluna.

\begin{Shaded}
\begin{Highlighting}[]
\NormalTok{mtcars}\SpecialCharTok{$}\NormalTok{hp\_per\_wt }\OtherTok{\textless{}{-}}\NormalTok{ mtcars}\SpecialCharTok{$}\NormalTok{hp }\SpecialCharTok{/}\NormalTok{ mtcars}\SpecialCharTok{$}\NormalTok{wt}
\NormalTok{filtered\_data }\OtherTok{\textless{}{-}} \FunctionTok{filter}\NormalTok{(mtcars, hp\_per\_wt }\SpecialCharTok{\textgreater{}} \DecValTok{30}\NormalTok{)}
\NormalTok{result }\OtherTok{\textless{}{-}} \FunctionTok{select}\NormalTok{(filtered\_data, mpg, hp\_per\_wt, gear)}
\FunctionTok{print}\NormalTok{(result)}
\end{Highlighting}
\end{Shaded}

O \emph{pipe} encadeia operações que corremos \textbf{uma vez}. Mas muitas tarefas exigem
\textbf{repetir} o mesmo cálculo várias vezes --- e é aí que entram os ciclos, tema dos
dois capítulos seguintes.

\chapter{O ciclo while}\label{while}

Até aqui, cada instrução do nosso código corria uma só vez. Mas há tarefas que
exigem \textbf{repetição}: somar os números até um certo limite, pedir dados ao
utilizador até ele decidir parar, refinar um cálculo até atingir a precisão desejada.
Para estas situações, o R oferece os \textbf{ciclos} (ou \emph{loops}), e o primeiro que
estudamos é o \texttt{while}.

O \texttt{while} repete um bloco de código \textbf{enquanto} uma condição lógica se mantiver
verdadeira. É a estrutura ideal quando não sabemos de antemão quantas repetições
serão precisas --- sabemos apenas a \emph{condição} que faz o ciclo parar.

\begin{Shaded}
\begin{Highlighting}[]
\ControlFlowTok{while}\NormalTok{ (condição) \{}
  \CommentTok{\# bloco de código a repetir}
\NormalTok{\}}
\end{Highlighting}
\end{Shaded}

\begin{itemize}
\tightlist
\item
  \textbf{condição} --- uma expressão lógica, avaliada antes de cada repetição. Enquanto for
  \texttt{TRUE}, o bloco corre; assim que passar a \texttt{FALSE}, o ciclo termina.
\item
  \textbf{bloco de código} --- as instruções que se repetem.
\end{itemize}

O funcionamento é um ciclo de três passos: avalia a condição; se for verdadeira,
executa o bloco; volta a avaliar a condição. Repete-se assim até a condição se tornar
falsa, momento em que o programa segue para a instrução seguinte ao \texttt{while}.

Neste capítulo aprendemos a repetir código com \texttt{while}, a garantir que o ciclo
termina, a interrompê-lo com \texttt{break} e a reconhecer e evitar os ciclos infinitos.

\section{Exemplos}\label{exemplos}

\textbf{Somar números até um limite.} O ciclo seguinte soma todos os inteiros de 1 a 10:

\begin{Shaded}
\begin{Highlighting}[]
\NormalTok{limite }\OtherTok{\textless{}{-}} \DecValTok{10}
\NormalTok{soma }\OtherTok{\textless{}{-}} \DecValTok{0}
\NormalTok{contador }\OtherTok{\textless{}{-}} \DecValTok{1}

\ControlFlowTok{while}\NormalTok{ (contador }\SpecialCharTok{\textless{}=}\NormalTok{ limite) \{}
\NormalTok{  soma }\OtherTok{\textless{}{-}}\NormalTok{ soma }\SpecialCharTok{+}\NormalTok{ contador}
\NormalTok{  contador }\OtherTok{\textless{}{-}}\NormalTok{ contador }\SpecialCharTok{+} \DecValTok{1}
\NormalTok{\}}

\FunctionTok{print}\NormalTok{(}\FunctionTok{paste}\NormalTok{(}\StringTok{"A soma dos números de 1 a"}\NormalTok{, limite, }\StringTok{"é:"}\NormalTok{, soma))}
\DocumentationTok{\#\# [1] "A soma dos números de 1 a 10 é: 55"}
\end{Highlighting}
\end{Shaded}

Em cada repetição, o valor de \texttt{contador} é somado a \texttt{soma} e depois incrementado.
Quando \texttt{contador} ultrapassa \texttt{limite}, a condição torna-se falsa e o ciclo termina. A
linha \texttt{contador\ \textless{}-\ contador\ +\ 1} é o que garante o progresso rumo à paragem --- e é
precisamente aí que reside o erro mais comum, como o próximo exemplo mostra.

\textbf{Escrever a tabuada de um número.} O programa seguinte deveria imprimir a tabuada
de um número dado pelo utilizador --- mas \textbf{nunca termina}. Consegue ver porquê?

\begin{Shaded}
\begin{Highlighting}[]
\NormalTok{n }\OtherTok{\textless{}{-}} \FunctionTok{as.numeric}\NormalTok{(}\FunctionTok{readline}\NormalTok{(}\StringTok{"Introduza um número inteiro: "}\NormalTok{))}

\FunctionTok{print}\NormalTok{(}\FunctionTok{paste}\NormalTok{(}\StringTok{"Tabuada do"}\NormalTok{, n, }\StringTok{":"}\NormalTok{))}
\NormalTok{i }\OtherTok{\textless{}{-}} \DecValTok{1}
\ControlFlowTok{while}\NormalTok{ (i }\SpecialCharTok{\textless{}=} \DecValTok{10}\NormalTok{) \{}
  \FunctionTok{print}\NormalTok{(}\FunctionTok{paste}\NormalTok{(n, }\StringTok{"x"}\NormalTok{, i, }\StringTok{"="}\NormalTok{, n }\SpecialCharTok{*}\NormalTok{ i))}
\NormalTok{  i }\SpecialCharTok{+} \DecValTok{1}          \CommentTok{\# \textless{}{-} o erro está aqui}
\NormalTok{\}}
\end{Highlighting}
\end{Shaded}

A expressão \texttt{i\ +\ 1} \textbf{calcula} o novo valor, mas não o \textbf{guarda} em \texttt{i}: o
resultado é simplesmente descartado. Como \texttt{i} nunca muda, a condição \texttt{i\ \textless{}=\ 10}
mantém-se eternamente verdadeira e o ciclo corre para sempre. A correção é atribuir o
resultado de volta à variável:

\begin{Shaded}
\begin{Highlighting}[]
\NormalTok{i }\OtherTok{\textless{}{-}}\NormalTok{ i }\SpecialCharTok{+} \DecValTok{1}
\end{Highlighting}
\end{Shaded}

Sempre que um \texttt{while} parece não terminar, o primeiro suspeito é uma variável de
controlo que não está a ser atualizada.

\textbf{Interromper com \texttt{break}.} A instrução \texttt{break} sai imediatamente do ciclo, mesmo
que a condição ainda seja verdadeira:

\begin{Shaded}
\begin{Highlighting}[]
\NormalTok{limite }\OtherTok{\textless{}{-}} \DecValTok{10}
\NormalTok{soma }\OtherTok{\textless{}{-}} \DecValTok{0}
\NormalTok{contador }\OtherTok{\textless{}{-}} \DecValTok{1}

\ControlFlowTok{while}\NormalTok{ (contador }\SpecialCharTok{\textless{}=}\NormalTok{ limite) \{}
\NormalTok{  soma }\OtherTok{\textless{}{-}}\NormalTok{ soma }\SpecialCharTok{+}\NormalTok{ contador}
  \FunctionTok{print}\NormalTok{(contador)}
  \ControlFlowTok{if}\NormalTok{ (contador }\SpecialCharTok{==} \DecValTok{3}\NormalTok{) \{}
    \ControlFlowTok{break}            \CommentTok{\# sai do ciclo quando contador chega a 3}
\NormalTok{  \}}
\NormalTok{  contador }\OtherTok{\textless{}{-}}\NormalTok{ contador }\SpecialCharTok{+} \DecValTok{1}
\NormalTok{\}}
\DocumentationTok{\#\# [1] 1}
\DocumentationTok{\#\# [1] 2}
\DocumentationTok{\#\# [1] 3}
\end{Highlighting}
\end{Shaded}

Há duas instruções irmãs do \texttt{break}. A \texttt{next} não termina o ciclo --- apenas \textbf{salta}
para a repetição seguinte, ignorando o resto do bloco na iteração atual. E, além do
\texttt{while}, o R tem o \texttt{repeat}, que repete indefinidamente e só para quando encontra um
\texttt{break}; é útil quando queremos executar o bloco \textbf{pelo menos uma vez} antes de
testar a condição de paragem.

\section{Considerações importantes}\label{considerauxe7uxf5es-importantes}

\begin{itemize}
\tightlist
\item
  \textbf{Condição de paragem.} Garanta que a condição acaba por se tornar falsa. Um
  \texttt{while} cuja condição nunca falha é um \textbf{ciclo infinito}, que bloqueia o
  programa (no RStudio, interrompe-se com \emph{Esc} ou o botão de \emph{stop}).
\item
  \textbf{Atualizar a variável de controlo.} A variável que a condição testa tem de ser
  atualizada dentro do ciclo, no sentido correto --- foi o que faltou no exemplo da
  tabuada.
\item
  \textbf{Desempenho.} Em R, os ciclos são muitas vezes mais lentos do que as operações
  \textbf{vetorizadas}. Para grandes volumes de dados, prefira funções como \texttt{sum()},
  \texttt{mean()} ou a família \texttt{apply} a percorrer os elementos um a um.
\end{itemize}

\section{Exercícios}\label{exercuxedcios-11}

Nos exercícios seguintes, use um ciclo \texttt{while}.

\textbf{1.} Peça ao utilizador um inteiro \(n\) e imprima todos os inteiros de 0 até \(n\), um
por linha.

\textbf{2.} Peça um inteiro, calcule e mostre o seu quadrado; depois pergunte ao
utilizador se quer continuar ou terminar. O ciclo deve prosseguir até o utilizador
escolher terminar.

\textbf{3.} Peça um inteiro \(n\) e calcule a soma dos primeiros \(n\) inteiros (de 1 a \(n\)).

\textbf{4.} Peça um inteiro \(n\) e conte quantos números pares existem no intervalo de 0 a
\(n\).

\textbf{5.} Some inteiros positivos introduzidos pelo utilizador até que ele insira um
número negativo, momento em que o programa para e apresenta a soma dos positivos.

\textbf{6.} Peça um inteiro positivo e faça uma contagem decrescente até zero, imprimindo
cada número.

\textbf{7.} Peça um inteiro negativo e conte de forma crescente até zero, imprimindo cada
número.

\textbf{8.} Peça um número ao utilizador: se for positivo, conte de forma decrescente até
zero; se for negativo, conte de forma crescente até zero. Imprima cada número.

\textbf{9.} Peça um inteiro não negativo e calcule o seu fatorial. Recorde que \(0! = 1\) e
que, por exemplo, \(4! = 4 \times 3 \times 2 \times 1 = 24\).

\textbf{10.} Leia um inteiro positivo e calcule a soma dos seus algarismos. Por exemplo,
para 23, o resultado é \(2 + 3 = 5\).
\emph{Sugestão:} Os operadores \texttt{\%/\%} (divisão inteira) e \texttt{\%\%} (resto) permitem extrair os algarismos um a um.

\textbf{11.} Gere e imprima os termos da sucessão de Fibonacci até atingir um valor
máximo indicado pelo utilizador. A sucessão começa em \(1, 1\) e cada termo é a soma
dos dois anteriores: \(1, 1, 2, 3, 5, 8, 13, \ldots\)

O \texttt{while} repete enquanto uma condição se verifica. Quando, pelo contrário, sabemos
\textbf{exatamente} sobre que valores queremos iterar --- os elementos de um vetor, as
linhas de uma matriz ---, há uma estrutura mais natural: o ciclo \texttt{for}, tema do próximo
capítulo.

\chapter{O ciclo for}\label{for}

O ciclo \texttt{while} do capítulo anterior repete-se enquanto uma condição se mantém --- ideal
quando não sabemos quantas repetições serão precisas. Muitas vezes, porém, sabemos
exatamente sobre que valores queremos trabalhar: os elementos de um vetor, as linhas
de uma matriz, as colunas de uma \emph{data frame}. Para esses casos existe o ciclo \texttt{for},
que percorre, um a um, os elementos de uma sequência.

\begin{Shaded}
\begin{Highlighting}[]
\ControlFlowTok{for}\NormalTok{ (variável }\ControlFlowTok{in}\NormalTok{ sequência) \{}
  \CommentTok{\# bloco de código a executar em cada iteração}
\NormalTok{\}}
\end{Highlighting}
\end{Shaded}

\begin{itemize}
\tightlist
\item
  \textbf{variável} --- a variável de controlo, que em cada iteração assume o valor do
  elemento seguinte da sequência;
\item
  \textbf{sequência} --- a estrutura a percorrer (um vetor, uma lista, um intervalo como
  \texttt{1:10}, \ldots);
\item
  \textbf{bloco de código} --- as instruções executadas em cada passagem.
\end{itemize}

O funcionamento é direto: a variável de controlo toma sucessivamente cada valor da
sequência, o bloco corre uma vez para cada um, e o ciclo termina quando todos os
elementos tiverem sido processados. Ao contrário do \texttt{while}, o número de iterações
fica determinado à partida.

Neste capítulo aprendemos a percorrer sequências com \texttt{for}, a iterar por índices, a
trabalhar sobre matrizes e \emph{data frames}, e a fazê-lo de forma correta e eficiente
(percorrendo com \texttt{seq\_along()} e pré-alocando os resultados).

\section{Exemplos}\label{exemplos-1}

\textbf{Imprimir uma sequência.} O caso mais simples percorre um intervalo:

\begin{Shaded}
\begin{Highlighting}[]
\ControlFlowTok{for}\NormalTok{ (i }\ControlFlowTok{in} \DecValTok{0}\SpecialCharTok{:}\DecValTok{10}\NormalTok{) \{}
  \FunctionTok{print}\NormalTok{(i)}
\NormalTok{\}}
\DocumentationTok{\#\# [1] 0}
\DocumentationTok{\#\# [1] 1}
\DocumentationTok{\#\# [1] 2}
\DocumentationTok{\#\# ... (até 10)}
\end{Highlighting}
\end{Shaded}

\textbf{Somar os elementos de um vetor.} A variável de controlo assume cada valor do
vetor:

\begin{Shaded}
\begin{Highlighting}[]
\NormalTok{numeros }\OtherTok{\textless{}{-}} \FunctionTok{c}\NormalTok{(}\DecValTok{1}\NormalTok{, }\DecValTok{2}\NormalTok{, }\DecValTok{3}\NormalTok{, }\DecValTok{4}\NormalTok{, }\DecValTok{5}\NormalTok{)}
\NormalTok{soma }\OtherTok{\textless{}{-}} \DecValTok{0}

\ControlFlowTok{for}\NormalTok{ (num }\ControlFlowTok{in}\NormalTok{ numeros) \{}
\NormalTok{  soma }\OtherTok{\textless{}{-}}\NormalTok{ soma }\SpecialCharTok{+}\NormalTok{ num}
\NormalTok{\}}

\FunctionTok{print}\NormalTok{(}\FunctionTok{paste}\NormalTok{(}\StringTok{"A soma dos números é:"}\NormalTok{, soma))}
\DocumentationTok{\#\# [1] "A soma dos números é: 15"}
\end{Highlighting}
\end{Shaded}

\textbf{Iterar por índices.} Quando precisamos de \textbf{modificar} os elementos nas suas
posições, iteramos sobre os índices em vez dos valores. Aqui, \texttt{seq\_along(numeros)}
gera as posições \texttt{1,\ 2,\ …,\ length(numeros)}:

\begin{Shaded}
\begin{Highlighting}[]
\NormalTok{numeros }\OtherTok{\textless{}{-}} \FunctionTok{c}\NormalTok{(}\DecValTok{10}\NormalTok{, }\DecValTok{20}\NormalTok{, }\DecValTok{30}\NormalTok{, }\DecValTok{40}\NormalTok{, }\DecValTok{50}\NormalTok{)}

\ControlFlowTok{for}\NormalTok{ (i }\ControlFlowTok{in} \FunctionTok{seq\_along}\NormalTok{(numeros)) \{}
\NormalTok{  numeros[i] }\OtherTok{\textless{}{-}}\NormalTok{ numeros[i] }\SpecialCharTok{*} \DecValTok{2}
\NormalTok{\}}

\FunctionTok{print}\NormalTok{(numeros)}
\DocumentationTok{\#\# [1]  20  40  60  80 100}
\end{Highlighting}
\end{Shaded}

É tentador escrever \texttt{for\ (i\ in\ 1:length(numeros))}, e é assim que aparece em muito
código por aí --- mas esconde uma armadilha. Se o vetor estiver \textbf{vazio},
\texttt{length(numeros)} é \texttt{0} e \texttt{1:0} não produz uma sequência vazia, como seria de
esperar, mas sim o vetor \texttt{c(1,\ 0)} --- pelo que o ciclo corre duas vezes, com índices
inválidos, gerando erros difíceis de diagnosticar. A função \texttt{seq\_along(numeros)} (ou
\texttt{seq\_len(n)}) evita o problema, devolvendo corretamente uma sequência vazia quando não
há elementos. Prefira sempre \texttt{seq\_along()} a \texttt{1:length()}.

\textbf{Iterar sobre uma matriz.} O \texttt{for} percorre matrizes por linha ou por coluna. Aqui
calculamos a soma de cada linha, guardando o resultado num vetor \textbf{pré-alocado}:

\begin{Shaded}
\begin{Highlighting}[]
\NormalTok{matriz }\OtherTok{\textless{}{-}} \FunctionTok{matrix}\NormalTok{(}\DecValTok{1}\SpecialCharTok{:}\DecValTok{9}\NormalTok{, }\AttributeTok{nrow =} \DecValTok{3}\NormalTok{, }\AttributeTok{ncol =} \DecValTok{3}\NormalTok{)}
\NormalTok{soma\_linhas }\OtherTok{\textless{}{-}} \FunctionTok{numeric}\NormalTok{(}\FunctionTok{nrow}\NormalTok{(matriz))     }\CommentTok{\# vetor de zeros, do tamanho certo}

\ControlFlowTok{for}\NormalTok{ (i }\ControlFlowTok{in} \DecValTok{1}\SpecialCharTok{:}\FunctionTok{nrow}\NormalTok{(matriz)) \{}
\NormalTok{  soma\_linhas[i] }\OtherTok{\textless{}{-}} \FunctionTok{sum}\NormalTok{(matriz[i, ])}
\NormalTok{\}}

\FunctionTok{print}\NormalTok{(soma\_linhas)}
\DocumentationTok{\#\# [1] 12 15 18}
\end{Highlighting}
\end{Shaded}

\textbf{Médias das colunas de uma \emph{data frame}.}

\begin{Shaded}
\begin{Highlighting}[]
\NormalTok{dados }\OtherTok{\textless{}{-}} \FunctionTok{data.frame}\NormalTok{(}
  \AttributeTok{A =} \FunctionTok{c}\NormalTok{(}\DecValTok{1}\NormalTok{, }\DecValTok{2}\NormalTok{, }\DecValTok{3}\NormalTok{),}
  \AttributeTok{B =} \FunctionTok{c}\NormalTok{(}\DecValTok{4}\NormalTok{, }\DecValTok{5}\NormalTok{, }\DecValTok{6}\NormalTok{),}
  \AttributeTok{C =} \FunctionTok{c}\NormalTok{(}\DecValTok{7}\NormalTok{, }\DecValTok{8}\NormalTok{, }\DecValTok{9}\NormalTok{)}
\NormalTok{)}

\NormalTok{medias }\OtherTok{\textless{}{-}} \FunctionTok{numeric}\NormalTok{(}\FunctionTok{ncol}\NormalTok{(dados))}

\ControlFlowTok{for}\NormalTok{ (col }\ControlFlowTok{in} \DecValTok{1}\SpecialCharTok{:}\FunctionTok{ncol}\NormalTok{(dados)) \{}
\NormalTok{  medias[col] }\OtherTok{\textless{}{-}} \FunctionTok{mean}\NormalTok{(dados[, col])}
\NormalTok{\}}

\FunctionTok{print}\NormalTok{(medias)}
\DocumentationTok{\#\# [1] 2 5 8}
\end{Highlighting}
\end{Shaded}

Repare que, nos dois últimos exemplos, criámos o vetor de resultados \textbf{antes} do
ciclo, já com o tamanho final (\texttt{numeric(nrow(matriz))}). Este hábito --- a
\textbf{pré-alocação} --- é importante: fazer o vetor ``crescer'' dentro do ciclo (por
exemplo, com \texttt{c()} a cada iteração) obriga o R a copiá-lo repetidamente e torna o
código muito mais lento em problemas grandes.

\section{Exercícios}\label{exercuxedcios-12}

Nos exercícios seguintes, use um ciclo \texttt{for}.

\textbf{1.} Imprima todos os números de 1 a 10.

\textbf{2.} Imprima o quadrado de cada número de 1 a 5.

\textbf{3.} Some todos os elementos do vetor \texttt{numeros\ \textless{}-\ c(2,\ 4,\ 6,\ 8,\ 10)}.

\textbf{4.} Calcule a média das notas do vetor \texttt{notas\ \textless{}-\ c(85,\ 90,\ 78,\ 92,\ 88)}.

\textbf{5.} Imprima todos os números ímpares do vetor
\texttt{valores\ \textless{}-\ c(1,\ 3,\ 4,\ 6,\ 9,\ 11,\ 14)}.

\textbf{6.} Leia um inteiro e imprima a sua tabuada de 1 a 10.

\textbf{7.} Leia um inteiro \(n\) e construa o vetor \texttt{c(1,\ 2,\ ...,\ n-1,\ n)}.

\textbf{8.} Leia um inteiro e determine se é um \textbf{número perfeito} --- aquele cuja soma dos
divisores próprios é igual ao próprio número (por exemplo, 6 tem divisores próprios
1, 2 e 3, e \(1 + 2 + 3 = 6\)).

\textbf{9.} Encontre e imprima todos os números perfeitos entre 1 e 1000.

\textbf{10.} Leia dois inteiros \(n\) e \(m\) e imprima os números entre 0 e \(n\)
(inclusive), com um intervalo de \(m\) entre eles.

\textbf{11.} Calcule a média de um conjunto de notas dadas pelo utilizador. O número de
notas é indicado antes da sua introdução. Eis um exemplo da interação:

\begin{verbatim}
Introduza o número de notas: 4
Introduza a nota: 12
Introduza a nota: 13
Introduza a nota: 15.5
Introduza a nota: 16
A média das notas é 14.125
\end{verbatim}

\textbf{12.} Crie uma matriz \(4 \times 4\) com os números de 1 a 16. Usando o ciclo \texttt{for},
calcule:

\begin{enumerate}
\def\labelenumi{(\alph{enumi})}
\tightlist
\item
  a média da terceira coluna;
\item
  a média da terceira linha;
\item
  a média de cada coluna;
\item
  a média de cada linha;
\item
  a média dos números pares de cada coluna;
\item
  a média dos números ímpares de cada linha.
\end{enumerate}

\emph{Sugestão:} Use estruturas de decisão (\texttt{if}\ldots{}\texttt{else}\ldots, Capítulo~\ref{selecao}) sempre
que precise de condicionar a lógica dentro do ciclo.

Os ciclos \texttt{for} e \texttt{while} dão-nos controlo total sobre a repetição. No entanto, em R,
muitas dessas repetições podem ser escritas de forma mais curta e rápida com a família
de funções \texttt{apply} --- o tema que exploramos a seguir.

\chapter{\texorpdfstring{A família \texttt{apply}}{A família apply}}\label{apply}

Os ciclos \texttt{for} e \texttt{while} (Capítulos~\ref{while} e~\ref{for}) dão-nos controlo
total sobre a repetição, mas em R há muitas vezes um caminho mais curto --- e mais
rápido. A família de funções \texttt{apply} foi desenhada precisamente para \textbf{aplicar uma
função a cada elemento de um objeto} sem escrever um ciclo explícito. O resultado é
código mais conciso, mais legível e, tipicamente, mais eficiente.

A ideia é sempre a mesma: em vez de descrevermos \emph{como} percorrer os elementos
(inicializar um contador, incrementá-lo, indexar), dizemos apenas \emph{o que} queremos
fazer a cada um. Esta forma de pensar --- separar a operação da mecânica da iteração ---
chama-se \textbf{programação funcional}, e é uma das marcas do estilo idiomático em R.

As quatro funções mais usadas resumem-se assim:

{\def\LTcaptype{none} % do not increment counter
\begin{longtable}[]{@{}
  >{\raggedright\arraybackslash}p{(\linewidth - 8\tabcolsep) * \real{0.2000}}
  >{\raggedright\arraybackslash}p{(\linewidth - 8\tabcolsep) * \real{0.2000}}
  >{\raggedright\arraybackslash}p{(\linewidth - 8\tabcolsep) * \real{0.2000}}
  >{\raggedright\arraybackslash}p{(\linewidth - 8\tabcolsep) * \real{0.2000}}
  >{\raggedright\arraybackslash}p{(\linewidth - 8\tabcolsep) * \real{0.2000}}@{}}
\toprule\noalign{}
\begin{minipage}[b]{\linewidth}\raggedright
Função
\end{minipage} & \begin{minipage}[b]{\linewidth}\raggedright
Sintaxe
\end{minipage} & \begin{minipage}[b]{\linewidth}\raggedright
O que faz
\end{minipage} & \begin{minipage}[b]{\linewidth}\raggedright
Entrada
\end{minipage} & \begin{minipage}[b]{\linewidth}\raggedright
Saída
\end{minipage} \\
\midrule\noalign{}
\endhead
\bottomrule\noalign{}
\endlastfoot
\texttt{apply} & \texttt{apply(X,\ MARGIN,\ FUN)} & Aplica \texttt{FUN} às linhas ou colunas & Matriz ou \emph{data frame} & vetor, lista ou \emph{array} \\
\texttt{lapply} & \texttt{lapply(X,\ FUN)} & Aplica \texttt{FUN} a cada elemento & Lista, vetor ou \emph{data frame} & lista \\
\texttt{sapply} & \texttt{sapply(X,\ FUN)} & Como \texttt{lapply}, mas simplifica & Lista, vetor ou \emph{data frame} & vetor ou matriz \\
\texttt{tapply} & \texttt{tapply(X,\ INDEX,\ FUN)} & Aplica \texttt{FUN} a cada grupo & Vetor & \emph{array} \\
\end{longtable}
}

Neste capítulo aprendemos a substituir ciclos por \texttt{apply} (linhas/colunas de uma
matriz), \texttt{lapply} e \texttt{sapply} (elementos de uma lista ou vetor) e \texttt{tapply} (grupos
definidos por um fator), escrevendo código mais curto e mais expressivo.

\section{\texorpdfstring{A função \texttt{apply()}}{A função apply()}}\label{a-funuxe7uxe3o-apply}

A \texttt{apply()} aplica uma função às \textbf{margens} de uma matriz ou \emph{data frame}: às linhas,
às colunas, ou a ambas.

\begin{Shaded}
\begin{Highlighting}[]
\FunctionTok{apply}\NormalTok{(X, MARGIN, FUN)}
\end{Highlighting}
\end{Shaded}

\begin{itemize}
\tightlist
\item
  \texttt{X} --- a matriz ou \emph{data frame};
\item
  \texttt{MARGIN} --- \texttt{1} para as linhas, \texttt{2} para as colunas, \texttt{c(1,\ 2)} para cada célula;
\item
  \texttt{FUN} --- a função a aplicar.
\end{itemize}

Para calcular a soma, a média e a raiz quadrada de cada coluna de uma matriz:

\begin{Shaded}
\begin{Highlighting}[]
\NormalTok{matriz }\OtherTok{\textless{}{-}} \FunctionTok{matrix}\NormalTok{(}\DecValTok{1}\SpecialCharTok{:}\DecValTok{9}\NormalTok{, }\AttributeTok{nrow =} \DecValTok{3}\NormalTok{)}

\FunctionTok{apply}\NormalTok{(matriz, }\DecValTok{2}\NormalTok{, sum)}
\DocumentationTok{\#\# [1]  6 15 24}

\FunctionTok{apply}\NormalTok{(matriz, }\DecValTok{2}\NormalTok{, mean)}
\DocumentationTok{\#\# [1] 2 5 8}

\FunctionTok{apply}\NormalTok{(matriz, }\DecValTok{2}\NormalTok{, sqrt)}
\DocumentationTok{\#\#      [,1] [,2] [,3]}
\DocumentationTok{\#\# [1,] 1.00 2.00 2.65}
\DocumentationTok{\#\# [2,] 1.41 2.24 2.83}
\DocumentationTok{\#\# [3,] 1.73 2.45 3.00}
\end{Highlighting}
\end{Shaded}

O segundo argumento é a chave: \texttt{apply(matriz,\ 1,\ sum)} somaria cada \textbf{linha};
\texttt{apply(matriz,\ 2,\ sum)}, cada \textbf{coluna}.

\section{\texorpdfstring{A função \texttt{lapply()}}{A função lapply()}}\label{a-funuxe7uxe3o-lapply}

A \texttt{lapply()} aplica uma função a cada elemento de uma lista ou vetor e devolve
\textbf{sempre uma lista} --- daí o \texttt{l} do nome.

\begin{Shaded}
\begin{Highlighting}[]
\FunctionTok{lapply}\NormalTok{(X, FUN, ...)}
\end{Highlighting}
\end{Shaded}

Por exemplo, para pôr uma série de nomes em minúsculas:

\begin{Shaded}
\begin{Highlighting}[]
\NormalTok{nomes }\OtherTok{\textless{}{-}} \FunctionTok{c}\NormalTok{(}\StringTok{"ANA"}\NormalTok{, }\StringTok{"JOAO"}\NormalTok{, }\StringTok{"PAULO"}\NormalTok{, }\StringTok{"FILIPA"}\NormalTok{)}
\NormalTok{nomes\_minusc }\OtherTok{\textless{}{-}} \FunctionTok{lapply}\NormalTok{(nomes, tolower)}
\FunctionTok{str}\NormalTok{(nomes\_minusc)}
\DocumentationTok{\#\# List of 4}
\DocumentationTok{\#\#  $ : chr "ana"}
\DocumentationTok{\#\#  $ : chr "joao"}
\DocumentationTok{\#\#  $ : chr "paulo"}
\DocumentationTok{\#\#  $ : chr "filipa"}
\end{Highlighting}
\end{Shaded}

A função aplica-se também a cada elemento de uma lista --- aqui, a raiz quadrada de dois
vetores:

\begin{Shaded}
\begin{Highlighting}[]
\NormalTok{dados }\OtherTok{\textless{}{-}} \FunctionTok{list}\NormalTok{(}\AttributeTok{a =} \DecValTok{1}\SpecialCharTok{:}\DecValTok{4}\NormalTok{, }\AttributeTok{b =} \DecValTok{5}\SpecialCharTok{:}\DecValTok{8}\NormalTok{)}
\FunctionTok{lapply}\NormalTok{(dados, sqrt)}
\DocumentationTok{\#\# $a}
\DocumentationTok{\#\# [1] 1.00 1.41 1.73 2.00}
\DocumentationTok{\#\#}
\DocumentationTok{\#\# $b}
\DocumentationTok{\#\# [1] 2.24 2.45 2.65 2.83}
\end{Highlighting}
\end{Shaded}

\section{\texorpdfstring{A função \texttt{sapply()}}{A função sapply()}}\label{a-funuxe7uxe3o-sapply}

A \texttt{sapply()} faz o mesmo que a \texttt{lapply()}, mas \textbf{simplifica} o resultado sempre que
possível --- para um vetor ou uma matriz --- em vez de devolver uma lista. O \texttt{s} vem de
\emph{simplify}.

\begin{Shaded}
\begin{Highlighting}[]
\NormalTok{dados }\OtherTok{\textless{}{-}} \DecValTok{1}\SpecialCharTok{:}\DecValTok{5}
\NormalTok{f }\OtherTok{\textless{}{-}} \ControlFlowTok{function}\NormalTok{(x) x}\SpecialCharTok{\^{}}\DecValTok{2}

\FunctionTok{lapply}\NormalTok{(dados, f)     }\CommentTok{\# devolve uma lista com cinco elementos}
\DocumentationTok{\#\# [[1]]}
\DocumentationTok{\#\# [1] 1}
\DocumentationTok{\#\# ... (até [[5]] com 25)}

\FunctionTok{sapply}\NormalTok{(dados, f)     }\CommentTok{\# devolve um vetor}
\DocumentationTok{\#\# [1]  1  4  9 16 25}
\end{Highlighting}
\end{Shaded}

A simplificação da \texttt{sapply()} é cómoda, mas tem um senão: quando os resultados não
têm todos o mesmo comprimento, a função devolve silenciosamente uma lista, e o nosso
código a jusante pode partir-se sem aviso claro. Em programas onde a robustez importa,
prefira a \texttt{vapply()}, que obriga a declarar de antemão o tipo e o tamanho do
resultado de cada elemento e falha imediatamente se algo não corresponder --- tornando
o erro visível e fácil de localizar.

\section{\texorpdfstring{A função \texttt{tapply()}}{A função tapply()}}\label{a-funuxe7uxe3o-tapply}

A \texttt{tapply()} aplica uma função a \textbf{subconjuntos} de um vetor, definidos por um fator.
É a ferramenta natural para calcular estatísticas por grupo.

\begin{Shaded}
\begin{Highlighting}[]
\FunctionTok{tapply}\NormalTok{(X, INDEX, FUN, ...)}
\end{Highlighting}
\end{Shaded}

\begin{itemize}
\tightlist
\item
  \texttt{X} --- o vetor de dados;
\item
  \texttt{INDEX} --- o fator (ou lista de fatores) que define os grupos;
\item
  \texttt{FUN} --- a função a aplicar a cada grupo.
\end{itemize}

Usemos o conjunto \texttt{iris}, um clássico introduzido por Ronald A. Fisher em 1936. Contém
medições morfológicas (comprimento e largura de sépalas e pétalas) de três espécies de
íris (\emph{setosa}, \emph{versicolor} e \emph{virginica}). Vejamos as primeiras linhas:

\begin{Shaded}
\begin{Highlighting}[]
\FunctionTok{head}\NormalTok{(iris)}
\DocumentationTok{\#\#   Sepal.Length Sepal.Width Petal.Length Petal.Width Species}
\DocumentationTok{\#\# 1          5.1         3.5          1.4         0.2  setosa}
\DocumentationTok{\#\# 2          4.9         3.0          1.4         0.2  setosa}
\DocumentationTok{\#\# 3          4.7         3.2          1.3         0.2  setosa}
\DocumentationTok{\#\# 4          4.6         3.1          1.5         0.2  setosa}
\DocumentationTok{\#\# 5          5.0         3.6          1.4         0.2  setosa}
\DocumentationTok{\#\# 6          5.4         3.9          1.7         0.4  setosa}
\end{Highlighting}
\end{Shaded}

Para obter o comprimento médio da pétala \textbf{em cada espécie}, uma única linha basta:

\begin{Shaded}
\begin{Highlighting}[]
\FunctionTok{tapply}\NormalTok{(iris}\SpecialCharTok{$}\NormalTok{Petal.Length, iris}\SpecialCharTok{$}\NormalTok{Species, mean)}
\DocumentationTok{\#\#     setosa versicolor  virginica}
\DocumentationTok{\#\#       1.46       4.26       5.55}
\end{Highlighting}
\end{Shaded}

O resultado revela de imediato o que distingue as espécies: as pétalas da \emph{virginica}
são, em média, quase quatro vezes mais compridas do que as da \emph{setosa}.

A família tem outros membros úteis. A \texttt{vapply()} é uma \texttt{sapply()} mais segura (com tipo
de saída declarado); a \texttt{mapply()} aplica uma função a \textbf{vários} vetores em paralelo,
elemento a elemento; e a \texttt{Map()} faz o mesmo devolvendo uma lista. Dominadas as
quatro funções deste capítulo, estas variantes seguem-se com naturalidade.

\section{Exercícios}\label{exercuxedcios-13}

\textbf{1.} Crie um vetor \texttt{nomes} com ``Alice'', ``Bruno'', ``Carla'' e ``Daniel''. Use \texttt{lapply()}
para os converter em maiúsculas e verifique o resultado.

\begin{Shaded}
\begin{Highlighting}[]
\NormalTok{nomes }\OtherTok{\textless{}{-}} \FunctionTok{c}\NormalTok{(}\StringTok{"Alice"}\NormalTok{, }\StringTok{"Bruno"}\NormalTok{, }\StringTok{"Carla"}\NormalTok{, }\StringTok{"Daniel"}\NormalTok{)}
\CommentTok{\# A sua solução aqui}
\end{Highlighting}
\end{Shaded}

\textbf{2.} Crie um vetor \texttt{palavras} com ``cão'', ``gato'', ``elefante'' e ``leão''. Use \texttt{sapply()}
para calcular o comprimento (número de carateres) de cada palavra. \emph{Sugestão:} veja
\texttt{?nchar}.

\begin{Shaded}
\begin{Highlighting}[]
\NormalTok{palavras }\OtherTok{\textless{}{-}} \FunctionTok{c}\NormalTok{(}\StringTok{"cão"}\NormalTok{, }\StringTok{"gato"}\NormalTok{, }\StringTok{"elefante"}\NormalTok{, }\StringTok{"leão"}\NormalTok{)}
\CommentTok{\# A sua solução aqui}
\end{Highlighting}
\end{Shaded}

\textbf{3.} Crie uma matriz \(3 \times 3\) chamada \texttt{matriz\_exemplo} com os números de 1 a 9.
Use \texttt{apply()} para calcular a soma de cada linha e, depois, a média de cada coluna.

\begin{Shaded}
\begin{Highlighting}[]
\NormalTok{matriz\_exemplo }\OtherTok{\textless{}{-}} \FunctionTok{matrix}\NormalTok{(}\DecValTok{1}\SpecialCharTok{:}\DecValTok{9}\NormalTok{, }\AttributeTok{nrow =} \DecValTok{3}\NormalTok{)}
\CommentTok{\# A sua solução aqui}
\end{Highlighting}
\end{Shaded}

\textbf{4.} No conjunto \texttt{mtcars}, use \texttt{tapply()} para calcular o consumo médio (\texttt{mpg}) para
cada número de cilindros (\texttt{cyl}).

\begin{Shaded}
\begin{Highlighting}[]
\FunctionTok{data}\NormalTok{(mtcars)}
\CommentTok{\# A sua solução aqui}
\end{Highlighting}
\end{Shaded}

\textbf{5.} Crie uma lista \texttt{minha\_lista} com três elementos: um vetor numérico de 1 a 5, um
vetor de carateres \texttt{c("A",\ "B",\ "C")} e uma matriz \(2 \times 2\) com os valores de 1 a
4. Use \texttt{lapply()} para calcular o comprimento de cada elemento e, depois, \texttt{sapply()}
para simplificar o resultado num vetor.

\begin{Shaded}
\begin{Highlighting}[]
\NormalTok{minha\_lista }\OtherTok{\textless{}{-}} \FunctionTok{list}\NormalTok{(}\DecValTok{1}\SpecialCharTok{:}\DecValTok{5}\NormalTok{, }\FunctionTok{c}\NormalTok{(}\StringTok{"A"}\NormalTok{, }\StringTok{"B"}\NormalTok{, }\StringTok{"C"}\NormalTok{), }\FunctionTok{matrix}\NormalTok{(}\DecValTok{1}\SpecialCharTok{:}\DecValTok{4}\NormalTok{, }\AttributeTok{nrow =} \DecValTok{2}\NormalTok{))}
\CommentTok{\# A sua solução aqui}
\end{Highlighting}
\end{Shaded}

\textbf{6.} Crie um vetor \texttt{numeros} de 1 a 10. Use \texttt{sapply()} com uma função personalizada
que devolva \texttt{"par"} se o número for par e \texttt{"ímpar"} caso contrário.

\begin{Shaded}
\begin{Highlighting}[]
\NormalTok{numeros }\OtherTok{\textless{}{-}} \DecValTok{1}\SpecialCharTok{:}\DecValTok{10}
\CommentTok{\# A sua solução aqui}
\end{Highlighting}
\end{Shaded}

\textbf{7.} No conjunto \texttt{iris}, use \texttt{tapply()} para calcular a média do comprimento das
pétalas (\texttt{Petal.Length}) por espécie e a variância da largura das sépalas
(\texttt{Sepal.Width}) por espécie.

\begin{Shaded}
\begin{Highlighting}[]
\FunctionTok{data}\NormalTok{(iris)}
\CommentTok{\# A sua solução aqui}
\end{Highlighting}
\end{Shaded}

\textbf{8.} No conjunto \texttt{mtcars}, use \texttt{apply()} para normalizar (escala de 0 a 1) todas as
colunas numéricas e converta o resultado numa \emph{data frame}. Depois, use \texttt{lapply()}
para aplicar \texttt{summary()} a cada coluna normalizada. O valor normalizado é
\(\dfrac{x - \min(x)}{\max(x) - \min(x)}\).

\begin{Shaded}
\begin{Highlighting}[]
\FunctionTok{data}\NormalTok{(mtcars)}
\CommentTok{\# A sua solução aqui}
\end{Highlighting}
\end{Shaded}

\textbf{9.} Um professor quer analisar o desempenho de 30 alunos em quatro exames.

\begin{itemize}
\tightlist
\item
  Crie uma \emph{data frame} \texttt{notas\_alunos} com 30 linhas (alunos) e 4 colunas (exames),
  preenchida com valores aleatórios entre 0 e 100 (use \texttt{runif()}).
\item
  Use \texttt{apply()} para calcular a média de cada aluno nos quatro exames.
\item
  Use \texttt{sapply()} para calcular a nota mínima e máxima de cada exame.
\end{itemize}

\textbf{10.} Um gestor de \emph{marketing} analisa o número de clientes de quatro lojas ao longo
de um mês.

\begin{itemize}
\tightlist
\item
  Crie uma matriz \texttt{frequencia\_lojas} com 4 linhas (lojas) e 30 colunas (dias),
  preenchida com valores aleatórios entre 0 e 200 (use \texttt{sample()}).
\item
  Use \texttt{apply()} para calcular a média diária de clientes em cada loja.
\item
  Use \texttt{apply()} para calcular o total de clientes de cada loja no mês.
\item
  Converta a matriz numa \emph{data frame} e use \texttt{sapply()} para calcular o desvio padrão
  do número de clientes em cada dia (cada coluna).
\end{itemize}

Com a família \texttt{apply} fechamos as ferramentas essenciais de programação. O próximo
capítulo é uma pausa para consolidar tudo o que vimos até aqui, através de exercícios
que combinam estruturas de dados, decisões, ciclos e funções.

\chapter{Exercícios de revisão}\label{revisao}

Chegámos ao fim da primeira parte do livro. Até aqui construímos, peça a peça, as
ferramentas fundamentais da programação em R: os tipos e as estruturas de dados
(Capítulo~\ref{estruturas}), as estruturas de decisão (Capítulo~\ref{selecao}),
as funções (Capítulo~\ref{funcoes}), os ciclos (Capítulos~\ref{while} e~\ref{for}) e a família \texttt{apply} (Capítulo~\ref{apply}). Isoladamente, cada uma
destas ferramentas é simples; a sua força vem de as combinarmos para resolver
problemas reais.

Este capítulo é uma pausa para consolidar. Os problemas que se seguem são deliberadamente
integradores --- cada um exige juntar várias ferramentas --- e inspiram-se em contextos
concretos: satisfação de clientes, produtividade de equipas, simulação de amostras,
temperaturas urbanas e até o efeito da radiação na memória de um computador. Resolvê-los
sem consultar as soluções é a melhor forma de medir o que ficou realmente aprendido.

Consolidar, através de problemas aplicados, a combinação de estruturas de dados,
funções, estruturas de decisão, ciclos e a família \texttt{apply}.

\textbf{1.} Um investigador analisa uma sondagem que avalia a satisfação de clientes com
três serviços (A, B e C). Cada linha representa um cliente e cada coluna a pontuação
(de 1 a 10) atribuída a cada serviço:

\begin{Shaded}
\begin{Highlighting}[]
\NormalTok{satisfacao }\OtherTok{\textless{}{-}} \FunctionTok{data.frame}\NormalTok{(}
  \AttributeTok{servico\_A =} \FunctionTok{c}\NormalTok{(}\DecValTok{8}\NormalTok{, }\DecValTok{5}\NormalTok{, }\DecValTok{9}\NormalTok{, }\DecValTok{6}\NormalTok{, }\DecValTok{7}\NormalTok{, }\DecValTok{4}\NormalTok{, }\DecValTok{10}\NormalTok{, }\DecValTok{6}\NormalTok{, }\DecValTok{3}\NormalTok{, }\DecValTok{7}\NormalTok{),}
  \AttributeTok{servico\_B =} \FunctionTok{c}\NormalTok{(}\DecValTok{7}\NormalTok{, }\DecValTok{6}\NormalTok{, }\DecValTok{8}\NormalTok{, }\DecValTok{5}\NormalTok{, }\DecValTok{8}\NormalTok{, }\DecValTok{5}\NormalTok{, }\DecValTok{9}\NormalTok{, }\DecValTok{7}\NormalTok{, }\DecValTok{4}\NormalTok{, }\DecValTok{6}\NormalTok{),}
  \AttributeTok{servico\_C =} \FunctionTok{c}\NormalTok{(}\DecValTok{9}\NormalTok{, }\DecValTok{4}\NormalTok{, }\DecValTok{10}\NormalTok{, }\DecValTok{6}\NormalTok{, }\DecValTok{6}\NormalTok{, }\DecValTok{3}\NormalTok{, }\DecValTok{10}\NormalTok{, }\DecValTok{5}\NormalTok{, }\DecValTok{5}\NormalTok{, }\DecValTok{8}\NormalTok{)}
\NormalTok{)}
\end{Highlighting}
\end{Shaded}

\begin{enumerate}
\def\labelenumi{(\alph{enumi})}
\item
  Escreva uma função \texttt{calcula\_media()} que receba a \emph{data frame} \texttt{satisfacao} e
  devolva um vetor com a média de satisfação de cada cliente (média das 3 colunas).
\item
  Acrescente à \emph{data frame} uma coluna \texttt{media} com os valores calculados.
\item
  Com base na coluna \texttt{media}, use um ciclo \texttt{for} com \texttt{if}\ldots{}\texttt{else} para classificar
  cada cliente como \texttt{"Satisfeito"} (média \(\geq 7\)) ou \texttt{"Insatisfeito"} (caso
  contrário). Acrescente essa informação numa nova coluna \texttt{classificacao}.
\end{enumerate}

\textbf{2.} Uma empresa registou o número de tarefas concluídas por 15 colaboradores ao
longo de 5 dias úteis. A meta é realizar pelo menos 20 tarefas por dia; pretende-se
saber quem atingiu a meta em pelo menos 3 dos 5 dias. Cada linha é um colaborador,
cada coluna um dia:

\begin{Shaded}
\begin{Highlighting}[]
\NormalTok{tarefa\_funcionarios }\OtherTok{\textless{}{-}} \FunctionTok{data.frame}\NormalTok{(}
  \AttributeTok{dia1 =} \FunctionTok{c}\NormalTok{(}\DecValTok{22}\NormalTok{, }\DecValTok{18}\NormalTok{, }\DecValTok{19}\NormalTok{, }\DecValTok{25}\NormalTok{, }\DecValTok{30}\NormalTok{, }\DecValTok{15}\NormalTok{, }\DecValTok{21}\NormalTok{, }\DecValTok{24}\NormalTok{, }\DecValTok{20}\NormalTok{, }\DecValTok{19}\NormalTok{, }\DecValTok{28}\NormalTok{, }\DecValTok{26}\NormalTok{, }\DecValTok{12}\NormalTok{, }\DecValTok{20}\NormalTok{, }\DecValTok{18}\NormalTok{),}
  \AttributeTok{dia2 =} \FunctionTok{c}\NormalTok{(}\DecValTok{18}\NormalTok{, }\DecValTok{21}\NormalTok{, }\DecValTok{23}\NormalTok{, }\DecValTok{20}\NormalTok{, }\DecValTok{27}\NormalTok{, }\DecValTok{18}\NormalTok{, }\DecValTok{20}\NormalTok{, }\DecValTok{25}\NormalTok{, }\DecValTok{16}\NormalTok{, }\DecValTok{21}\NormalTok{, }\DecValTok{22}\NormalTok{, }\DecValTok{29}\NormalTok{, }\DecValTok{17}\NormalTok{, }\DecValTok{19}\NormalTok{, }\DecValTok{20}\NormalTok{),}
  \AttributeTok{dia3 =} \FunctionTok{c}\NormalTok{(}\DecValTok{25}\NormalTok{, }\DecValTok{17}\NormalTok{, }\DecValTok{15}\NormalTok{, }\DecValTok{18}\NormalTok{, }\DecValTok{26}\NormalTok{, }\DecValTok{16}\NormalTok{, }\DecValTok{19}\NormalTok{, }\DecValTok{30}\NormalTok{, }\DecValTok{21}\NormalTok{, }\DecValTok{18}\NormalTok{, }\DecValTok{24}\NormalTok{, }\DecValTok{20}\NormalTok{, }\DecValTok{14}\NormalTok{, }\DecValTok{23}\NormalTok{, }\DecValTok{22}\NormalTok{),}
  \AttributeTok{dia4 =} \FunctionTok{c}\NormalTok{(}\DecValTok{20}\NormalTok{, }\DecValTok{19}\NormalTok{, }\DecValTok{22}\NormalTok{, }\DecValTok{26}\NormalTok{, }\DecValTok{28}\NormalTok{, }\DecValTok{20}\NormalTok{, }\DecValTok{23}\NormalTok{, }\DecValTok{22}\NormalTok{, }\DecValTok{17}\NormalTok{, }\DecValTok{25}\NormalTok{, }\DecValTok{19}\NormalTok{, }\DecValTok{18}\NormalTok{, }\DecValTok{19}\NormalTok{, }\DecValTok{21}\NormalTok{, }\DecValTok{24}\NormalTok{),}
  \AttributeTok{dia5 =} \FunctionTok{c}\NormalTok{(}\DecValTok{19}\NormalTok{, }\DecValTok{22}\NormalTok{, }\DecValTok{18}\NormalTok{, }\DecValTok{27}\NormalTok{, }\DecValTok{29}\NormalTok{, }\DecValTok{14}\NormalTok{, }\DecValTok{24}\NormalTok{, }\DecValTok{28}\NormalTok{, }\DecValTok{20}\NormalTok{, }\DecValTok{23}\NormalTok{, }\DecValTok{25}\NormalTok{, }\DecValTok{27}\NormalTok{, }\DecValTok{16}\NormalTok{, }\DecValTok{17}\NormalTok{, }\DecValTok{18}\NormalTok{)}
\NormalTok{)}
\end{Highlighting}
\end{Shaded}

\begin{enumerate}
\def\labelenumi{(\alph{enumi})}
\item
  Escreva uma função \texttt{verifica\_meta()} que receba a \emph{data frame} e devolva um vetor
  lógico indicando, para cada colaborador, se atingiu a meta em pelo menos 3 dias.
\item
  Com base nesse vetor lógico, use um \texttt{for} com \texttt{if}\ldots{}\texttt{else} para criar um vetor
  \texttt{classificacao} com os valores \texttt{"Atingiu\ Meta"} (3 ou mais dias) ou
  \texttt{"Não\ Atingiu\ Meta"} (caso contrário).
\item
  Acrescente o vetor \texttt{classificacao} como nova coluna da \emph{data frame}.
\end{enumerate}

\textbf{3.} Um estatístico estuda o comportamento da média amostral. Simulou 6 amostras de
tamanho 15 (valores entre 0 e 100) e quer calcular a média de cada uma, interrompendo
o processo assim que encontrar uma média superior a 75. Os dados estão na matriz:

\begin{Shaded}
\begin{Highlighting}[]
\NormalTok{amostras }\OtherTok{\textless{}{-}} \FunctionTok{matrix}\NormalTok{(}
  \FunctionTok{c}\NormalTok{(}
    \DecValTok{52}\NormalTok{, }\DecValTok{68}\NormalTok{, }\DecValTok{75}\NormalTok{, }\DecValTok{48}\NormalTok{, }\DecValTok{60}\NormalTok{, }\DecValTok{55}\NormalTok{, }\DecValTok{72}\NormalTok{, }\DecValTok{81}\NormalTok{, }\DecValTok{64}\NormalTok{, }\DecValTok{70}\NormalTok{, }\DecValTok{45}\NormalTok{, }\DecValTok{69}\NormalTok{, }\DecValTok{74}\NormalTok{, }\DecValTok{60}\NormalTok{, }\DecValTok{71}\NormalTok{,}
    \DecValTok{49}\NormalTok{, }\DecValTok{50}\NormalTok{, }\DecValTok{63}\NormalTok{, }\DecValTok{41}\NormalTok{, }\DecValTok{59}\NormalTok{, }\DecValTok{55}\NormalTok{, }\DecValTok{67}\NormalTok{, }\DecValTok{62}\NormalTok{, }\DecValTok{58}\NormalTok{, }\DecValTok{66}\NormalTok{, }\DecValTok{61}\NormalTok{, }\DecValTok{54}\NormalTok{, }\DecValTok{65}\NormalTok{, }\DecValTok{60}\NormalTok{, }\DecValTok{57}\NormalTok{,}
    \DecValTok{80}\NormalTok{, }\DecValTok{82}\NormalTok{, }\DecValTok{85}\NormalTok{, }\DecValTok{79}\NormalTok{, }\DecValTok{90}\NormalTok{, }\DecValTok{88}\NormalTok{, }\DecValTok{92}\NormalTok{, }\DecValTok{85}\NormalTok{, }\DecValTok{87}\NormalTok{, }\DecValTok{84}\NormalTok{, }\DecValTok{86}\NormalTok{, }\DecValTok{81}\NormalTok{, }\DecValTok{89}\NormalTok{, }\DecValTok{83}\NormalTok{, }\DecValTok{91}\NormalTok{,}
    \DecValTok{30}\NormalTok{, }\DecValTok{42}\NormalTok{, }\DecValTok{51}\NormalTok{, }\DecValTok{39}\NormalTok{, }\DecValTok{35}\NormalTok{, }\DecValTok{48}\NormalTok{, }\DecValTok{46}\NormalTok{, }\DecValTok{52}\NormalTok{, }\DecValTok{41}\NormalTok{, }\DecValTok{37}\NormalTok{, }\DecValTok{45}\NormalTok{, }\DecValTok{43}\NormalTok{, }\DecValTok{38}\NormalTok{, }\DecValTok{44}\NormalTok{, }\DecValTok{36}\NormalTok{,}
    \DecValTok{69}\NormalTok{, }\DecValTok{74}\NormalTok{, }\DecValTok{70}\NormalTok{, }\DecValTok{66}\NormalTok{, }\DecValTok{68}\NormalTok{, }\DecValTok{72}\NormalTok{, }\DecValTok{71}\NormalTok{, }\DecValTok{73}\NormalTok{, }\DecValTok{65}\NormalTok{, }\DecValTok{67}\NormalTok{, }\DecValTok{75}\NormalTok{, }\DecValTok{64}\NormalTok{, }\DecValTok{63}\NormalTok{, }\DecValTok{62}\NormalTok{, }\DecValTok{76}\NormalTok{,}
    \DecValTok{55}\NormalTok{, }\DecValTok{58}\NormalTok{, }\DecValTok{61}\NormalTok{, }\DecValTok{56}\NormalTok{, }\DecValTok{60}\NormalTok{, }\DecValTok{63}\NormalTok{, }\DecValTok{57}\NormalTok{, }\DecValTok{62}\NormalTok{, }\DecValTok{64}\NormalTok{, }\DecValTok{59}\NormalTok{, }\DecValTok{65}\NormalTok{, }\DecValTok{66}\NormalTok{, }\DecValTok{68}\NormalTok{, }\DecValTok{67}\NormalTok{, }\DecValTok{69}
\NormalTok{  ),}
  \AttributeTok{nrow =} \DecValTok{6}\NormalTok{,}
  \AttributeTok{byrow =} \ConstantTok{TRUE}
\NormalTok{)}
\end{Highlighting}
\end{Shaded}

\begin{enumerate}
\def\labelenumi{(\alph{enumi})}
\item
  Crie uma função \texttt{calcula\_media\_amostra()} que receba a matriz \texttt{amostras} e um
  índice, e devolva a média dessa linha.
\item
  Escreva um ciclo \texttt{while} que percorra as linhas, calculando a média de cada
  amostra. Se alguma média exceder 75, o programa imprime \texttt{"Média\ Excedeu\ Limite"} e
  para; caso contrário, continua até processar todas as amostras.
\end{enumerate}

\textbf{4.} Crie uma lista \texttt{listas\_transformadas} de comprimento 7, em que o elemento de
índice \(i\) contém a sequência
\[x_j = j^2 + i, \qquad j = 1, 2, \ldots, i.\]
Por exemplo, para \(i = 3\) o vetor é \([\,1^2 + 3,\; 2^2 + 3,\; 3^2 + 3\,] = [4, 7, 12]\).

\begin{enumerate}
\def\labelenumi{(\alph{enumi})}
\item
  Construa a lista \texttt{listas\_transformadas} com um ciclo \texttt{for}.
\item
  Verifique o comprimento do último vetor da lista.
\item
  Escreva uma função \texttt{soma\_elementos()} que receba uma lista e devolva um vetor com
  a soma de cada vetor nela contido.
\item
  Use \texttt{soma\_elementos()} para somar os elementos de cada vetor de
  \texttt{listas\_transformadas}.
\end{enumerate}

\textbf{5.} Crie dois vetores \(x\) e \(y\), ambos de comprimento 60, com inteiros aleatórios
escolhidos com reposição entre 0 e 999.

\begin{enumerate}
\def\labelenumi{(\alph{enumi})}
\item
  Gere \(x\) e \(y\) com \texttt{sample()} e \texttt{replace\ =\ TRUE}.
\item
  Selecione os valores de \(y\) menores que 500 cujos correspondentes em \(x\) (mesma
  posição) sejam múltiplos de 3.
\item
  Crie o vetor \texttt{dif\_absoluta} com o valor absoluto da diferença \(|x_i - y_i|\) em
  cada posição.
\item
  Calcule a média e o desvio padrão de \texttt{dif\_absoluta}.
\end{enumerate}

\textbf{6.} Seja \(x\) um vetor de 80 observações de uma distribuição normal de média 5 e
desvio padrão 2.

\begin{enumerate}
\def\labelenumi{(\alph{enumi})}
\item
  Gere \(x\) com \texttt{rnorm()}.
\item
  Calcule a média amostral \(\bar{x}\).
\item
  Crie o vetor \texttt{resultado} cujo \(i\)-ésimo elemento é
  \(z_i = \sqrt{\left| x_i^{3} - \bar{x}^{2} \right|}\).
\item
  Calcule o máximo, o mínimo e a média de \texttt{resultado}.
\end{enumerate}

\textbf{7.} Um climatologista estuda as temperaturas médias diárias em 12 cidades ao longo
de uma semana. Quer a temperatura média semanal de cada cidade e verificar se está
acima ou abaixo de 25~ºC. Cada linha é uma cidade, cada coluna um dia:

\begin{Shaded}
\begin{Highlighting}[]
\NormalTok{temperaturas\_cidades }\OtherTok{\textless{}{-}} \FunctionTok{data.frame}\NormalTok{(}
  \AttributeTok{segunda  =} \FunctionTok{c}\NormalTok{(}\DecValTok{28}\NormalTok{, }\DecValTok{22}\NormalTok{, }\DecValTok{25}\NormalTok{, }\DecValTok{19}\NormalTok{, }\DecValTok{30}\NormalTok{, }\DecValTok{24}\NormalTok{, }\DecValTok{21}\NormalTok{, }\DecValTok{26}\NormalTok{, }\DecValTok{27}\NormalTok{, }\DecValTok{23}\NormalTok{, }\DecValTok{31}\NormalTok{, }\DecValTok{20}\NormalTok{),}
  \AttributeTok{terca    =} \FunctionTok{c}\NormalTok{(}\DecValTok{29}\NormalTok{, }\DecValTok{24}\NormalTok{, }\DecValTok{23}\NormalTok{, }\DecValTok{20}\NormalTok{, }\DecValTok{32}\NormalTok{, }\DecValTok{25}\NormalTok{, }\DecValTok{22}\NormalTok{, }\DecValTok{27}\NormalTok{, }\DecValTok{26}\NormalTok{, }\DecValTok{24}\NormalTok{, }\DecValTok{33}\NormalTok{, }\DecValTok{21}\NormalTok{),}
  \AttributeTok{quarta   =} \FunctionTok{c}\NormalTok{(}\DecValTok{30}\NormalTok{, }\DecValTok{25}\NormalTok{, }\DecValTok{22}\NormalTok{, }\DecValTok{18}\NormalTok{, }\DecValTok{33}\NormalTok{, }\DecValTok{26}\NormalTok{, }\DecValTok{23}\NormalTok{, }\DecValTok{28}\NormalTok{, }\DecValTok{25}\NormalTok{, }\DecValTok{22}\NormalTok{, }\DecValTok{34}\NormalTok{, }\DecValTok{22}\NormalTok{),}
  \AttributeTok{quinta   =} \FunctionTok{c}\NormalTok{(}\DecValTok{31}\NormalTok{, }\DecValTok{23}\NormalTok{, }\DecValTok{24}\NormalTok{, }\DecValTok{17}\NormalTok{, }\DecValTok{31}\NormalTok{, }\DecValTok{27}\NormalTok{, }\DecValTok{20}\NormalTok{, }\DecValTok{29}\NormalTok{, }\DecValTok{28}\NormalTok{, }\DecValTok{21}\NormalTok{, }\DecValTok{32}\NormalTok{, }\DecValTok{23}\NormalTok{),}
  \AttributeTok{sexta    =} \FunctionTok{c}\NormalTok{(}\DecValTok{27}\NormalTok{, }\DecValTok{21}\NormalTok{, }\DecValTok{26}\NormalTok{, }\DecValTok{16}\NormalTok{, }\DecValTok{29}\NormalTok{, }\DecValTok{28}\NormalTok{, }\DecValTok{19}\NormalTok{, }\DecValTok{30}\NormalTok{, }\DecValTok{29}\NormalTok{, }\DecValTok{25}\NormalTok{, }\DecValTok{30}\NormalTok{, }\DecValTok{24}\NormalTok{),}
  \AttributeTok{sabado   =} \FunctionTok{c}\NormalTok{(}\DecValTok{26}\NormalTok{, }\DecValTok{22}\NormalTok{, }\DecValTok{27}\NormalTok{, }\DecValTok{15}\NormalTok{, }\DecValTok{28}\NormalTok{, }\DecValTok{26}\NormalTok{, }\DecValTok{18}\NormalTok{, }\DecValTok{28}\NormalTok{, }\DecValTok{27}\NormalTok{, }\DecValTok{26}\NormalTok{, }\DecValTok{29}\NormalTok{, }\DecValTok{25}\NormalTok{),}
  \AttributeTok{domingo  =} \FunctionTok{c}\NormalTok{(}\DecValTok{25}\NormalTok{, }\DecValTok{23}\NormalTok{, }\DecValTok{28}\NormalTok{, }\DecValTok{14}\NormalTok{, }\DecValTok{27}\NormalTok{, }\DecValTok{25}\NormalTok{, }\DecValTok{17}\NormalTok{, }\DecValTok{27}\NormalTok{, }\DecValTok{26}\NormalTok{, }\DecValTok{27}\NormalTok{, }\DecValTok{28}\NormalTok{, }\DecValTok{26}\NormalTok{)}
\NormalTok{)}
\end{Highlighting}
\end{Shaded}

\begin{enumerate}
\def\labelenumi{(\alph{enumi})}
\item
  Escreva uma função \texttt{calcula\_media\_temperaturas()} que receba a \emph{data frame} e
  devolva um vetor \texttt{media\_semanal} com a temperatura média semanal de cada cidade.
\item
  Use um ciclo \texttt{while} para percorrer \texttt{media\_semanal} e criar um vetor
  \texttt{classificacao} que classifique cada cidade como \texttt{"Acima\ de\ 25ºC"} (média \(> 25\))
  ou \texttt{"Abaixo\ de\ 25ºC"} (caso contrário).
\item
  Acrescente \texttt{media\_semanal} e \texttt{classificacao} como novas colunas da \emph{data frame}.
\end{enumerate}

\textbf{8.} Os computadores podem sofrer alterações aleatórias em \emph{bits} da memória
(\emph{bit flips}) devido à ação de partículas de alta energia, como os raios gama --- um
fenómeno relevante em contextos críticos como a aviação, a medicina ou a criptografia.
Um engenheiro simula o impacto da radiação em blocos de memória RAM, procurando os
blocos mais afetados:

\begin{Shaded}
\begin{Highlighting}[]
\FunctionTok{set.seed}\NormalTok{(}\DecValTok{2025}\NormalTok{)}

\CommentTok{\# Estado original de 1000 bits (0 ou 1)}
\NormalTok{original }\OtherTok{\textless{}{-}} \FunctionTok{sample}\NormalTok{(}\FunctionTok{c}\NormalTok{(}\DecValTok{0}\NormalTok{, }\DecValTok{1}\NormalTok{), }\AttributeTok{size =} \DecValTok{1000}\NormalTok{, }\AttributeTok{replace =} \ConstantTok{TRUE}\NormalTok{)}

\CommentTok{\# Estado final após exposição à radiação:}
\CommentTok{\# cada bit tem 3\% de probabilidade de ser alterado (invertido)}
\NormalTok{flipado }\OtherTok{\textless{}{-}} \FunctionTok{ifelse}\NormalTok{(}\FunctionTok{runif}\NormalTok{(}\DecValTok{1000}\NormalTok{) }\SpecialCharTok{\textless{}} \FloatTok{0.03}\NormalTok{, }\DecValTok{1} \SpecialCharTok{{-}}\NormalTok{ original, original)}

\CommentTok{\# Divisão dos bits em blocos de memória (100 bits por bloco)}
\NormalTok{blocos }\OtherTok{\textless{}{-}} \FunctionTok{rep}\NormalTok{(}\DecValTok{1}\SpecialCharTok{:}\DecValTok{10}\NormalTok{, }\AttributeTok{each =} \DecValTok{100}\NormalTok{)}
\end{Highlighting}
\end{Shaded}

\begin{enumerate}
\def\labelenumi{(\alph{enumi})}
\item
  Crie um vetor \texttt{erros} com o número de \emph{bits} alterados, comparando \texttt{original} com
  \texttt{flipado}.
\item
  Crie um fator \texttt{bloco\_memoria} com os valores de 1 a 10, indicando o bloco de cada
  \emph{bit} (como em \texttt{blocos}).
\item
  Calcule o número de \emph{bit flips} em cada bloco de memória.
  \emph{Sugestão:} Use a função \texttt{tapply()}.
\end{enumerate}

Consolidadas as bases, a segunda parte do livro vira-se para a \textbf{visualização}: no
próximo capítulo aprendemos a representar dados graficamente com o R base.

\chapter{Gráficos com o R base}\label{graficos}

Começa aqui a segunda parte do livro, dedicada à \textbf{visualização}. Um gráfico bem
feito faz num instante o que uma tabela de números raramente consegue: revela um
padrão, uma tendência, uma assimetria, um valor atípico. A visualização não é um
enfeite que se acrescenta no fim de uma análise --- é uma ferramenta de raciocínio, e
muitas vezes o primeiro passo para compreender os dados.

O R inclui, de origem, um sistema de gráficos poderoso e flexível, a que se costuma
chamar \emph{R base}. Neste capítulo percorremos os tipos de gráfico mais comuns ---
barras, circular, histograma, caixa de bigodes, dispersão e linhas --- associando cada
um ao tipo de variável e à pergunta que responde. Mais adiante (Capítulo~\ref{ggplot}) conheceremos o \texttt{ggplot2}, uma alternativa moderna; mas dominar o R
base é indispensável, pela sua ubiquidade e pela rapidez com que produz um gráfico
exploratório.

Neste capítulo aprendemos a construir e personalizar os gráficos essenciais do R base
e, sobretudo, a escolher o gráfico certo para cada tipo de dados: barras e gráficos
circulares para variáveis categóricas; histogramas e caixas de bigodes para variáveis
quantitativas; dispersão e linhas para relações entre variáveis.

\section{Gráfico de barras}\label{gruxe1fico-de-barras}

Um \textbf{gráfico de barras} representa uma variável categórica: cada barra corresponde a
uma categoria e a sua altura indica a frequência (absoluta ou relativa) dessa
categoria. A largura das barras é constante e não tem significado estatístico.

\begin{Shaded}
\begin{Highlighting}[]
\FunctionTok{barplot}\NormalTok{(altura,}
        \AttributeTok{names.arg =} \ConstantTok{NULL}\NormalTok{,   }\CommentTok{\# rótulos do eixo x}
        \AttributeTok{main =} \ConstantTok{NULL}\NormalTok{,        }\CommentTok{\# título}
        \AttributeTok{xlab =} \ConstantTok{NULL}\NormalTok{, }\AttributeTok{ylab =} \ConstantTok{NULL}\NormalTok{,}
        \AttributeTok{col =} \ConstantTok{NULL}\NormalTok{,         }\CommentTok{\# cor(es) das barras}
        \AttributeTok{beside =} \ConstantTok{FALSE}\NormalTok{)     }\CommentTok{\# barras lado a lado (TRUE) ou empilhadas (FALSE)}
\end{Highlighting}
\end{Shaded}

O argumento \texttt{altura} pode ser um vetor de frequências (gráfico simples) ou uma matriz
(barras agrupadas ou empilhadas). O caminho habitual é obter as frequências com
\texttt{table()} e passá-las a \texttt{barplot()}:

\begin{Shaded}
\begin{Highlighting}[]
\CommentTok{\# Cores preferidas de nove pessoas}
\NormalTok{cores }\OtherTok{\textless{}{-}} \FunctionTok{c}\NormalTok{(}\StringTok{"Azul"}\NormalTok{, }\StringTok{"Vermelho"}\NormalTok{, }\StringTok{"Verde"}\NormalTok{, }\StringTok{"Azul"}\NormalTok{, }\StringTok{"Verde"}\NormalTok{,}
           \StringTok{"Vermelho"}\NormalTok{, }\StringTok{"Azul"}\NormalTok{, }\StringTok{"Verde"}\NormalTok{, }\StringTok{"Azul"}\NormalTok{)}

\NormalTok{frequencia\_absoluta }\OtherTok{\textless{}{-}} \FunctionTok{table}\NormalTok{(cores)}

\FunctionTok{barplot}\NormalTok{(frequencia\_absoluta,}
        \AttributeTok{main =} \StringTok{"Cores preferidas"}\NormalTok{,}
        \AttributeTok{xlab =} \StringTok{"Cor"}\NormalTok{,}
        \AttributeTok{ylab =} \StringTok{"Frequência absoluta"}\NormalTok{,}
        \AttributeTok{col =} \FunctionTok{c}\NormalTok{(}\StringTok{"steelblue"}\NormalTok{, }\StringTok{"seagreen"}\NormalTok{, }\StringTok{"firebrick"}\NormalTok{),}
        \AttributeTok{border =} \StringTok{"black"}\NormalTok{)}
\end{Highlighting}
\end{Shaded}

\begin{figure}

{\centering \includegraphics[width=0.8\linewidth]{introR_files/figure-latex/graf-barras-abs-1} 

}

\caption{Gráfico de barras com frequências absolutas.}\label{fig:graf-barras-abs}
\end{figure}

Para representar as \textbf{frequências relativas}, basta dividir pela dimensão da amostra:

\begin{Shaded}
\begin{Highlighting}[]
\NormalTok{frequencia\_relativa }\OtherTok{\textless{}{-}}\NormalTok{ frequencia\_absoluta }\SpecialCharTok{/} \FunctionTok{length}\NormalTok{(cores)}

\FunctionTok{barplot}\NormalTok{(frequencia\_relativa,}
        \AttributeTok{main =} \StringTok{"Cores preferidas"}\NormalTok{,}
        \AttributeTok{xlab =} \StringTok{"Cor"}\NormalTok{,}
        \AttributeTok{ylab =} \StringTok{"Frequência relativa"}\NormalTok{,}
        \AttributeTok{col =} \FunctionTok{c}\NormalTok{(}\StringTok{"steelblue"}\NormalTok{, }\StringTok{"seagreen"}\NormalTok{, }\StringTok{"firebrick"}\NormalTok{),}
        \AttributeTok{border =} \StringTok{"black"}\NormalTok{)}
\end{Highlighting}
\end{Shaded}

\begin{figure}

{\centering \includegraphics[width=0.8\linewidth]{introR_files/figure-latex/graf-barras-rel-1} 

}

\caption{O mesmo gráfico com frequências relativas.}\label{fig:graf-barras-rel}
\end{figure}

Repare que os dois gráficos têm exatamente a mesma forma --- muda apenas a escala do
eixo vertical. As frequências relativas são úteis quando se querem comparar amostras
de dimensões diferentes, para as quais as frequências absolutas seriam enganadoras.

O argumento \texttt{horiz\ =\ TRUE} desenha as barras na horizontal, o que ajuda quando os
rótulos das categorias são longos:

\begin{Shaded}
\begin{Highlighting}[]
\FunctionTok{barplot}\NormalTok{(frequencia\_absoluta,}
        \AttributeTok{main =} \StringTok{"Cores preferidas"}\NormalTok{,}
        \AttributeTok{xlab =} \StringTok{"Frequência absoluta"}\NormalTok{,}
        \AttributeTok{col =} \FunctionTok{c}\NormalTok{(}\StringTok{"steelblue"}\NormalTok{, }\StringTok{"seagreen"}\NormalTok{, }\StringTok{"firebrick"}\NormalTok{),}
        \AttributeTok{horiz =} \ConstantTok{TRUE}\NormalTok{)}
\end{Highlighting}
\end{Shaded}

\begin{figure}

{\centering \includegraphics[width=0.8\linewidth]{introR_files/figure-latex/graf-barras-horiz-1} 

}

\caption{Gráfico de barras horizontal.}\label{fig:graf-barras-horiz}
\end{figure}

\subsection{Barras agrupadas e empilhadas}\label{barras-agrupadas-e-empilhadas}

Quando queremos comparar \textbf{duas variáveis} ao mesmo tempo, organizamos os dados numa
matriz. Comparemos, por exemplo, a quantidade de amostras de dois tipos de rocha em
três regiões:

\begin{Shaded}
\begin{Highlighting}[]
\NormalTok{regioes }\OtherTok{\textless{}{-}} \FunctionTok{c}\NormalTok{(}\StringTok{"Região 1"}\NormalTok{, }\StringTok{"Região 2"}\NormalTok{, }\StringTok{"Região 3"}\NormalTok{)}
\NormalTok{granito }\OtherTok{\textless{}{-}} \FunctionTok{c}\NormalTok{(}\DecValTok{15}\NormalTok{, }\DecValTok{10}\NormalTok{, }\DecValTok{25}\NormalTok{)}
\NormalTok{basalto }\OtherTok{\textless{}{-}} \FunctionTok{c}\NormalTok{(}\DecValTok{20}\NormalTok{, }\DecValTok{15}\NormalTok{, }\DecValTok{10}\NormalTok{)}

\NormalTok{dados\_rochas }\OtherTok{\textless{}{-}} \FunctionTok{rbind}\NormalTok{(granito, basalto)}

\FunctionTok{barplot}\NormalTok{(dados\_rochas,}
        \AttributeTok{beside =} \ConstantTok{TRUE}\NormalTok{,}
        \AttributeTok{names.arg =}\NormalTok{ regioes,}
        \AttributeTok{col =} \FunctionTok{c}\NormalTok{(}\StringTok{"lightblue"}\NormalTok{, }\StringTok{"lightcoral"}\NormalTok{),}
        \AttributeTok{legend.text =} \FunctionTok{c}\NormalTok{(}\StringTok{"Granito"}\NormalTok{, }\StringTok{"Basalto"}\NormalTok{),}
        \AttributeTok{main =} \StringTok{"Amostras de rocha por região"}\NormalTok{,}
        \AttributeTok{xlab =} \StringTok{"Região"}\NormalTok{,}
        \AttributeTok{ylab =} \StringTok{"Número de amostras"}\NormalTok{)}
\end{Highlighting}
\end{Shaded}

\begin{figure}

{\centering \includegraphics[width=0.8\linewidth]{introR_files/figure-latex/graf-barras-agrup-1} 

}

\caption{Barras agrupadas: cada região mostra os dois tipos de rocha lado a lado.}\label{fig:graf-barras-agrup}
\end{figure}

Usámos \texttt{rbind()} para juntar os dois vetores numa matriz e \texttt{beside\ =\ TRUE} para pôr as
barras lado a lado; \texttt{legend.text} gera a legenda. Removendo \texttt{beside\ =\ TRUE}, as barras
passam a estar \textbf{empilhadas} --- cada barra mostra o total da região, repartido pelos
dois tipos de rocha:

\begin{Shaded}
\begin{Highlighting}[]
\FunctionTok{barplot}\NormalTok{(dados\_rochas,}
        \AttributeTok{names.arg =}\NormalTok{ regioes,}
        \AttributeTok{col =} \FunctionTok{c}\NormalTok{(}\StringTok{"lightblue"}\NormalTok{, }\StringTok{"lightcoral"}\NormalTok{),}
        \AttributeTok{legend.text =} \FunctionTok{c}\NormalTok{(}\StringTok{"Granito"}\NormalTok{, }\StringTok{"Basalto"}\NormalTok{),}
        \AttributeTok{main =} \StringTok{"Amostras de rocha por região"}\NormalTok{,}
        \AttributeTok{xlab =} \StringTok{"Região"}\NormalTok{,}
        \AttributeTok{ylab =} \StringTok{"Número de amostras"}\NormalTok{)}
\end{Highlighting}
\end{Shaded}

\begin{figure}

{\centering \includegraphics[width=0.8\linewidth]{introR_files/figure-latex/graf-barras-empilhado-1} 

}

\caption{As mesmas barras, agora empilhadas.}\label{fig:graf-barras-empilhado}
\end{figure}

\section{Gráfico circular}\label{gruxe1fico-circular}

O \textbf{gráfico circular} (ou de setores) mostra a proporção de cada categoria face ao
total: cada fatia é proporcional à percentagem que representa.

\begin{Shaded}
\begin{Highlighting}[]
\FunctionTok{pie}\NormalTok{(x, }\AttributeTok{labels =} \ConstantTok{NULL}\NormalTok{, }\AttributeTok{main =} \ConstantTok{NULL}\NormalTok{, }\AttributeTok{col =} \ConstantTok{NULL}\NormalTok{, }\AttributeTok{radius =} \DecValTok{1}\NormalTok{)}
\end{Highlighting}
\end{Shaded}

\begin{Shaded}
\begin{Highlighting}[]
\FunctionTok{pie}\NormalTok{(frequencia\_relativa,}
    \AttributeTok{main =} \StringTok{"Cores preferidas"}\NormalTok{,}
    \AttributeTok{col =} \FunctionTok{c}\NormalTok{(}\StringTok{"steelblue"}\NormalTok{, }\StringTok{"seagreen"}\NormalTok{, }\StringTok{"firebrick"}\NormalTok{),}
    \AttributeTok{labels =} \FunctionTok{paste}\NormalTok{(}\FunctionTok{names}\NormalTok{(frequencia\_relativa),}
                   \FunctionTok{round}\NormalTok{(frequencia\_relativa }\SpecialCharTok{*} \DecValTok{100}\NormalTok{, }\DecValTok{1}\NormalTok{), }\StringTok{"\%"}\NormalTok{))}
\end{Highlighting}
\end{Shaded}

\begin{figure}

{\centering \includegraphics[width=0.8\linewidth]{introR_files/figure-latex/graf-circular-1} 

}

\caption{Gráfico circular das cores preferidas, com percentagens.}\label{fig:graf-circular}
\end{figure}

O gráfico circular é popular, mas deve usar-se com moderação. O olho humano compara
comprimentos (barras) com muito mais precisão do que ângulos ou áreas (fatias), pelo
que diferenças pequenas entre categorias são difíceis de ler num gráfico circular, e
tornam-se ilegíveis quando há muitas fatias. Como regra prática: para comparar
frequências, um gráfico de barras é quase sempre a escolha mais clara.

\section{Histograma}\label{histograma}

O \textbf{histograma} é o gráfico por excelência para uma variável quantitativa contínua.
Divide-se o eixo horizontal em classes (intervalos) e desenha-se, sobre cada classe,
um retângulo cuja altura traduz a frequência das observações nela contidas. É a forma
mais direta de ver a \textbf{distribuição} dos dados: onde se concentram, se são simétricos,
se têm valores atípicos.

\begin{Shaded}
\begin{Highlighting}[]
\FunctionTok{hist}\NormalTok{(x, }\AttributeTok{breaks =} \StringTok{"Sturges"}\NormalTok{,}
     \AttributeTok{main =} \ConstantTok{NULL}\NormalTok{, }\AttributeTok{xlab =} \ConstantTok{NULL}\NormalTok{, }\AttributeTok{ylab =} \ConstantTok{NULL}\NormalTok{,}
     \AttributeTok{col =} \ConstantTok{NULL}\NormalTok{, }\AttributeTok{border =} \ConstantTok{NULL}\NormalTok{)}
\end{Highlighting}
\end{Shaded}

O argumento \texttt{breaks} controla as classes: pode ser um número, ou um método automático
(\texttt{"Sturges"}, \texttt{"Scott"}, \texttt{"FD"}). Por omissão, o R usa a regra de Sturges e classes do
tipo \((a, b]\) (abertas à esquerda, fechadas à direita).

\begin{Shaded}
\begin{Highlighting}[]
\CommentTok{\# Massa (kg) de 40 bicicletas}
\NormalTok{bicicletas }\OtherTok{\textless{}{-}} \FunctionTok{c}\NormalTok{(}\FloatTok{4.3}\NormalTok{, }\FloatTok{6.8}\NormalTok{, }\FloatTok{9.2}\NormalTok{, }\FloatTok{7.2}\NormalTok{, }\FloatTok{8.7}\NormalTok{, }\FloatTok{8.6}\NormalTok{, }\FloatTok{6.6}\NormalTok{, }\FloatTok{5.2}\NormalTok{, }\FloatTok{8.1}\NormalTok{, }\FloatTok{10.9}\NormalTok{,}
                \FloatTok{7.4}\NormalTok{, }\FloatTok{4.5}\NormalTok{, }\FloatTok{3.8}\NormalTok{, }\FloatTok{7.6}\NormalTok{, }\FloatTok{6.8}\NormalTok{, }\FloatTok{7.8}\NormalTok{, }\FloatTok{8.4}\NormalTok{, }\FloatTok{7.5}\NormalTok{, }\FloatTok{10.5}\NormalTok{, }\FloatTok{6.0}\NormalTok{,}
                \FloatTok{7.7}\NormalTok{, }\FloatTok{8.1}\NormalTok{, }\FloatTok{7.0}\NormalTok{, }\FloatTok{8.2}\NormalTok{, }\FloatTok{8.4}\NormalTok{, }\FloatTok{8.8}\NormalTok{, }\FloatTok{6.7}\NormalTok{, }\FloatTok{8.2}\NormalTok{, }\FloatTok{9.4}\NormalTok{, }\FloatTok{7.7}\NormalTok{,}
                \FloatTok{6.3}\NormalTok{, }\FloatTok{7.7}\NormalTok{, }\FloatTok{9.1}\NormalTok{, }\FloatTok{7.9}\NormalTok{, }\FloatTok{7.9}\NormalTok{, }\FloatTok{9.4}\NormalTok{, }\FloatTok{8.2}\NormalTok{, }\FloatTok{6.7}\NormalTok{, }\FloatTok{8.2}\NormalTok{, }\FloatTok{6.5}\NormalTok{)}

\NormalTok{h }\OtherTok{\textless{}{-}} \FunctionTok{hist}\NormalTok{(bicicletas,}
          \AttributeTok{breaks =} \StringTok{"Sturges"}\NormalTok{,}
          \AttributeTok{main =} \StringTok{"Massa das bicicletas"}\NormalTok{,}
          \AttributeTok{xlab =} \StringTok{"Massa (kg)"}\NormalTok{,}
          \AttributeTok{ylab =} \StringTok{"Frequência absoluta"}\NormalTok{,}
          \AttributeTok{col =} \StringTok{"lightblue"}\NormalTok{,}
          \AttributeTok{border =} \StringTok{"black"}\NormalTok{,}
          \AttributeTok{ylim =} \FunctionTok{c}\NormalTok{(}\DecValTok{0}\NormalTok{, }\DecValTok{12}\NormalTok{),}
          \AttributeTok{labels =} \ConstantTok{TRUE}\NormalTok{)}
\end{Highlighting}
\end{Shaded}

\begin{figure}

{\centering \includegraphics[width=0.8\linewidth]{introR_files/figure-latex/graf-hist-abs-1} 

}

\caption{Histograma da massa de 40 bicicletas (frequências absolutas).}\label{fig:graf-hist-abs}
\end{figure}

Guardar o resultado de \texttt{hist()} num objeto (aqui, \texttt{h}) dá acesso aos valores usados na
sua construção --- os limites das classes e as contagens:

\begin{Shaded}
\begin{Highlighting}[]
\NormalTok{h}\SpecialCharTok{$}\NormalTok{breaks    }\CommentTok{\# limites das classes}
\DocumentationTok{\#\# [1]  3  4  5  6  7  8  9 10 11}

\NormalTok{h}\SpecialCharTok{$}\NormalTok{counts    }\CommentTok{\# frequência absoluta em cada classe}
\DocumentationTok{\#\# [1]  1  2  2  8 10 11  4  2}
\end{Highlighting}
\end{Shaded}

\subsection{Densidade e curvas suaves}\label{densidade-e-curvas-suaves}

Com o argumento \texttt{freq\ =\ FALSE}, o histograma passa a representar a \textbf{densidade} em
vez das contagens:

\begin{Shaded}
\begin{Highlighting}[]
\FunctionTok{hist}\NormalTok{(bicicletas,}
     \AttributeTok{main =} \StringTok{"Massa das bicicletas"}\NormalTok{,}
     \AttributeTok{xlab =} \StringTok{"Massa (kg)"}\NormalTok{,}
     \AttributeTok{ylab =} \StringTok{"Densidade"}\NormalTok{,}
     \AttributeTok{freq =} \ConstantTok{FALSE}\NormalTok{,}
     \AttributeTok{col =} \StringTok{"lightblue"}\NormalTok{,}
     \AttributeTok{border =} \StringTok{"black"}\NormalTok{)}

\CommentTok{\# Sobrepor uma estimativa suave da densidade}
\FunctionTok{lines}\NormalTok{(}\FunctionTok{density}\NormalTok{(bicicletas), }\AttributeTok{col =} \StringTok{"red"}\NormalTok{, }\AttributeTok{lwd =} \DecValTok{2}\NormalTok{)}
\end{Highlighting}
\end{Shaded}

\begin{figure}

{\centering \includegraphics[width=0.8\linewidth]{introR_files/figure-latex/graf-hist-dens-1} 

}

\caption{Histograma em escala de densidade, com a curva de densidade estimada.}\label{fig:graf-hist-dens}
\end{figure}

É comum ler-se que \texttt{freq\ =\ FALSE} produz ``frequências relativas'', mas isso não é
rigorosamente exato. O que a opção desenha é a \textbf{densidade}: as alturas são escaladas
de modo a que a \emph{área total} do histograma seja igual a 1. Só quando todas as classes
têm a mesma amplitude é que a altura fica proporcional à frequência relativa. A
distinção importa porque é a escala de densidade --- e não a de frequências --- que
permite sobrepor ao histograma uma curva de densidade, como fizemos com
\texttt{density()}, ou a curva de um modelo teórico.

\section{\texorpdfstring{Caixa de bigodes (\emph{boxplot})}{Caixa de bigodes (boxplot)}}\label{caixa-de-bigodes-boxplot}

A \textbf{caixa de bigodes} resume a distribuição de uma variável a partir dos seus
quartis. A caixa vai do primeiro quartil (\(Q_1\)) ao terceiro (\(Q_3\)), com um traço na
mediana (\(Q_2\)); os ``bigodes'' estendem-se a partir da caixa e os pontos para lá deles
assinalam possíveis valores atípicos.

\begin{Shaded}
\begin{Highlighting}[]
\FunctionTok{boxplot}\NormalTok{(x,}
        \AttributeTok{main =} \ConstantTok{NULL}\NormalTok{, }\AttributeTok{xlab =} \ConstantTok{NULL}\NormalTok{, }\AttributeTok{ylab =} \ConstantTok{NULL}\NormalTok{,}
        \AttributeTok{col =} \ConstantTok{NULL}\NormalTok{, }\AttributeTok{border =} \ConstantTok{NULL}\NormalTok{,}
        \AttributeTok{horizontal =} \ConstantTok{FALSE}\NormalTok{)}
\end{Highlighting}
\end{Shaded}

Ao contrário do que por vezes se diz, os bigodes \textbf{não} marcam necessariamente o
mínimo e o máximo. Por convenção, cada bigode estende-se até à observação mais extrema
que ainda esteja dentro de \(1.5 \times \text{AIQ}\) da caixa, onde
\(\text{AIQ} = Q_3 - Q_1\) é a amplitude interquartil. As observações que caem para além
dos bigodes são desenhadas individualmente e interpretadas como candidatas a valores
atípicos. É esta regra que torna a caixa de bigodes tão útil para detetar \emph{outliers}.

\begin{Shaded}
\begin{Highlighting}[]
\FunctionTok{boxplot}\NormalTok{(bicicletas,}
        \AttributeTok{main =} \StringTok{"Massa das bicicletas"}\NormalTok{,}
        \AttributeTok{ylab =} \StringTok{"Massa (kg)"}\NormalTok{,}
        \AttributeTok{col =} \StringTok{"lightblue"}\NormalTok{,}
        \AttributeTok{border =} \StringTok{"darkblue"}\NormalTok{)}
\end{Highlighting}
\end{Shaded}

\begin{figure}

{\centering \includegraphics[width=0.8\linewidth]{introR_files/figure-latex/graf-box-vert-1} 

}

\caption{Caixa de bigodes da massa das bicicletas.}\label{fig:graf-box-vert}
\end{figure}

O argumento \texttt{horizontal\ =\ TRUE} roda o gráfico, e a função \texttt{par(mfrow\ =\ ...)} permite
dispor vários gráficos no mesmo painel:

\begin{Shaded}
\begin{Highlighting}[]
\FunctionTok{par}\NormalTok{(}\AttributeTok{mfrow =} \FunctionTok{c}\NormalTok{(}\DecValTok{1}\NormalTok{, }\DecValTok{2}\NormalTok{))                                    }\CommentTok{\# painel de 1 linha, 2 colunas}

\FunctionTok{boxplot}\NormalTok{(bicicletas, }\AttributeTok{main =} \StringTok{"Vertical"}\NormalTok{, }\AttributeTok{col =} \StringTok{"lightblue"}\NormalTok{)}
\FunctionTok{boxplot}\NormalTok{(bicicletas, }\AttributeTok{main =} \StringTok{"Horizontal"}\NormalTok{, }\AttributeTok{col =} \StringTok{"lightgreen"}\NormalTok{, }\AttributeTok{horizontal =} \ConstantTok{TRUE}\NormalTok{)}
\end{Highlighting}
\end{Shaded}

\begin{figure}

{\centering \includegraphics[width=0.8\linewidth]{introR_files/figure-latex/graf-box-par-1} 

}

\caption{Dois gráficos no mesmo painel, com par(mfrow).}\label{fig:graf-box-par}
\end{figure}

\begin{Shaded}
\begin{Highlighting}[]

\FunctionTok{par}\NormalTok{(}\AttributeTok{mfrow =} \FunctionTok{c}\NormalTok{(}\DecValTok{1}\NormalTok{, }\DecValTok{1}\NormalTok{))                                    }\CommentTok{\# repor o painel original}
\end{Highlighting}
\end{Shaded}

O verdadeiro valor da caixa de bigodes revela-se ao \textbf{comparar grupos}. Usando a
notação de fórmula \texttt{variável\ \textasciitilde{}\ grupo}, comparamos o comprimento da sépala entre as
três espécies de íris:

\begin{Shaded}
\begin{Highlighting}[]
\FunctionTok{boxplot}\NormalTok{(Sepal.Length }\SpecialCharTok{\textasciitilde{}}\NormalTok{ Species,}
        \AttributeTok{data =}\NormalTok{ iris,}
        \AttributeTok{main =} \StringTok{"Comprimento da sépala por espécie"}\NormalTok{,}
        \AttributeTok{xlab =} \StringTok{"Espécie"}\NormalTok{,}
        \AttributeTok{ylab =} \StringTok{"Comprimento da sépala (cm)"}\NormalTok{,}
        \AttributeTok{col =} \FunctionTok{c}\NormalTok{(}\StringTok{"lightblue"}\NormalTok{, }\StringTok{"lightgreen"}\NormalTok{, }\StringTok{"lightpink"}\NormalTok{))}
\end{Highlighting}
\end{Shaded}

\begin{figure}

{\centering \includegraphics[width=0.8\linewidth]{introR_files/figure-latex/graf-box-iris-1} 

}

\caption{Comprimento da sépala por espécie no conjunto iris.}\label{fig:graf-box-iris}
\end{figure}

A leitura é imediata: a \emph{virginica} tem sépalas tipicamente mais compridas, a \emph{setosa}
as mais curtas, e a dispersão dentro de cada espécie é relativamente semelhante.

\section{Gráfico de dispersão}\label{gruxe1fico-de-dispersuxe3o}

O \textbf{gráfico de dispersão} representa duas variáveis quantitativas, uma em cada eixo, e
é a ferramenta natural para investigar a \textbf{relação} entre elas: se crescem juntas, se
uma decresce quando a outra cresce, se há agrupamentos ou pontos afastados.

\begin{Shaded}
\begin{Highlighting}[]
\FunctionTok{plot}\NormalTok{(x, y,}
     \AttributeTok{main =} \ConstantTok{NULL}\NormalTok{, }\AttributeTok{xlab =} \ConstantTok{NULL}\NormalTok{, }\AttributeTok{ylab =} \ConstantTok{NULL}\NormalTok{,}
     \AttributeTok{col =} \ConstantTok{NULL}\NormalTok{,   }\CommentTok{\# cor dos pontos}
     \AttributeTok{pch =} \ConstantTok{NULL}\NormalTok{,   }\CommentTok{\# símbolo (ex.: 19 = ponto cheio)}
     \AttributeTok{cex =} \DecValTok{1}\NormalTok{)      }\CommentTok{\# tamanho dos pontos}
\end{Highlighting}
\end{Shaded}

\begin{Shaded}
\begin{Highlighting}[]
\FunctionTok{set.seed}\NormalTok{(}\DecValTok{1}\NormalTok{)}
\NormalTok{x }\OtherTok{\textless{}{-}} \FunctionTok{rnorm}\NormalTok{(}\DecValTok{100}\NormalTok{)}
\NormalTok{y }\OtherTok{\textless{}{-}}\NormalTok{ x }\SpecialCharTok{+} \FunctionTok{rnorm}\NormalTok{(}\DecValTok{100}\NormalTok{, }\AttributeTok{sd =} \FloatTok{0.5}\NormalTok{)}

\FunctionTok{plot}\NormalTok{(x, y,}
     \AttributeTok{main =} \StringTok{"Dispersão de X e Y"}\NormalTok{,}
     \AttributeTok{xlab =} \StringTok{"Variável X"}\NormalTok{,}
     \AttributeTok{ylab =} \StringTok{"Variável Y"}\NormalTok{,}
     \AttributeTok{col =} \StringTok{"steelblue"}\NormalTok{,}
     \AttributeTok{pch =} \DecValTok{19}\NormalTok{)}
\end{Highlighting}
\end{Shaded}

\begin{figure}

{\centering \includegraphics[width=0.8\linewidth]{introR_files/figure-latex/graf-disp-1} 

}

\caption{Gráfico de dispersão de duas variáveis relacionadas.}\label{fig:graf-disp}
\end{figure}

Podemos acrescentar a \textbf{reta de regressão} dos mínimos quadrados com \texttt{abline()}
aplicada a um modelo linear \texttt{lm()}, o que ajuda a ver a tendência global:

\begin{Shaded}
\begin{Highlighting}[]
\FunctionTok{plot}\NormalTok{(x, y,}
     \AttributeTok{main =} \StringTok{"Dispersão com reta de regressão"}\NormalTok{,}
     \AttributeTok{xlab =} \StringTok{"Variável X"}\NormalTok{,}
     \AttributeTok{ylab =} \StringTok{"Variável Y"}\NormalTok{,}
     \AttributeTok{col =} \StringTok{"steelblue"}\NormalTok{,}
     \AttributeTok{pch =} \DecValTok{19}\NormalTok{)}

\FunctionTok{abline}\NormalTok{(}\FunctionTok{lm}\NormalTok{(y }\SpecialCharTok{\textasciitilde{}}\NormalTok{ x), }\AttributeTok{col =} \StringTok{"seagreen"}\NormalTok{, }\AttributeTok{lwd =} \DecValTok{2}\NormalTok{)}
\end{Highlighting}
\end{Shaded}

\begin{figure}

{\centering \includegraphics[width=0.8\linewidth]{introR_files/figure-latex/graf-disp-reg-1} 

}

\caption{O mesmo gráfico com a reta de regressão dos mínimos quadrados.}\label{fig:graf-disp-reg}
\end{figure}

Quando há \textbf{grupos} nos dados, podemos distingui-los por cor, acrescentando pontos com
\texttt{points()} e identificando-os com uma legenda:

\begin{Shaded}
\begin{Highlighting}[]
\FunctionTok{set.seed}\NormalTok{(}\DecValTok{123}\NormalTok{)}
\NormalTok{x1 }\OtherTok{\textless{}{-}} \FunctionTok{rnorm}\NormalTok{(}\DecValTok{50}\NormalTok{, }\AttributeTok{mean =} \DecValTok{5}\NormalTok{); y1 }\OtherTok{\textless{}{-}} \DecValTok{3} \SpecialCharTok{*}\NormalTok{ x1 }\SpecialCharTok{+} \FunctionTok{rnorm}\NormalTok{(}\DecValTok{50}\NormalTok{)}
\NormalTok{x2 }\OtherTok{\textless{}{-}} \FunctionTok{rnorm}\NormalTok{(}\DecValTok{50}\NormalTok{, }\AttributeTok{mean =} \DecValTok{7}\NormalTok{); y2 }\OtherTok{\textless{}{-}} \DecValTok{3} \SpecialCharTok{*}\NormalTok{ x2 }\SpecialCharTok{+} \FunctionTok{rnorm}\NormalTok{(}\DecValTok{50}\NormalTok{)}

\FunctionTok{plot}\NormalTok{(x1, y1,}
     \AttributeTok{main =} \StringTok{"Dispersão com dois grupos"}\NormalTok{,}
     \AttributeTok{xlab =} \StringTok{"Eixo X"}\NormalTok{, }\AttributeTok{ylab =} \StringTok{"Eixo Y"}\NormalTok{,}
     \AttributeTok{col =} \StringTok{"firebrick"}\NormalTok{, }\AttributeTok{pch =} \DecValTok{16}\NormalTok{,}
     \AttributeTok{xlim =} \FunctionTok{range}\NormalTok{(}\FunctionTok{c}\NormalTok{(x1, x2)), }\AttributeTok{ylim =} \FunctionTok{range}\NormalTok{(}\FunctionTok{c}\NormalTok{(y1, y2)))}
\FunctionTok{points}\NormalTok{(x2, y2, }\AttributeTok{col =} \StringTok{"steelblue"}\NormalTok{, }\AttributeTok{pch =} \DecValTok{16}\NormalTok{)}
\FunctionTok{legend}\NormalTok{(}\StringTok{"topleft"}\NormalTok{,}
       \AttributeTok{legend =} \FunctionTok{c}\NormalTok{(}\StringTok{"Grupo 1"}\NormalTok{, }\StringTok{"Grupo 2"}\NormalTok{),}
       \AttributeTok{col =} \FunctionTok{c}\NormalTok{(}\StringTok{"firebrick"}\NormalTok{, }\StringTok{"steelblue"}\NormalTok{), }\AttributeTok{pch =} \DecValTok{16}\NormalTok{)}
\end{Highlighting}
\end{Shaded}

\begin{figure}

{\centering \includegraphics[width=0.8\linewidth]{introR_files/figure-latex/graf-disp-grupos-1} 

}

\caption{Dois grupos distinguidos por cor.}\label{fig:graf-disp-grupos}
\end{figure}

Ao sobrepor um segundo conjunto de pontos com \texttt{points()}, os limites dos eixos ficam
fixados pelo \textbf{primeiro} \texttt{plot()}. Se os pontos do segundo grupo caírem fora desses
limites, ficam invisíveis. Por isso definimos \texttt{xlim} e \texttt{ylim} à partida, com
\texttt{range()}, de modo a abranger todos os dados.

\section{Gráfico de linhas}\label{gruxe1fico-de-linhas}

O \textbf{gráfico de linhas} liga os pontos por segmentos e é indicado para mostrar a
evolução de uma variável ao longo de uma sequência ordenada --- em especial o tempo
(séries temporais). Obtém-se com \texttt{plot()} e o argumento \texttt{type}:

\begin{Shaded}
\begin{Highlighting}[]
\FunctionTok{plot}\NormalTok{(x, y,}
     \AttributeTok{type =} \StringTok{"l"}\NormalTok{,   }\CommentTok{\# "l" = linha, "p" = pontos, "b"/"o" = ambos}
     \AttributeTok{lty =} \DecValTok{1}\NormalTok{,      }\CommentTok{\# tipo de traço (1 = contínuo, 2 = tracejado)}
     \AttributeTok{lwd =} \DecValTok{1}\NormalTok{)      }\CommentTok{\# espessura da linha}
\end{Highlighting}
\end{Shaded}

\begin{Shaded}
\begin{Highlighting}[]
\NormalTok{tempo }\OtherTok{\textless{}{-}} \DecValTok{1}\SpecialCharTok{:}\DecValTok{10}
\NormalTok{valores }\OtherTok{\textless{}{-}} \FunctionTok{c}\NormalTok{(}\DecValTok{3}\NormalTok{, }\DecValTok{5}\NormalTok{, }\DecValTok{7}\NormalTok{, }\DecValTok{9}\NormalTok{, }\DecValTok{11}\NormalTok{, }\DecValTok{10}\NormalTok{, }\DecValTok{8}\NormalTok{, }\DecValTok{6}\NormalTok{, }\DecValTok{4}\NormalTok{, }\DecValTok{2}\NormalTok{)}

\FunctionTok{plot}\NormalTok{(tempo, valores,}
     \AttributeTok{type =} \StringTok{"o"}\NormalTok{,}
     \AttributeTok{main =} \StringTok{"Evolução ao longo do tempo"}\NormalTok{,}
     \AttributeTok{xlab =} \StringTok{"Tempo"}\NormalTok{,}
     \AttributeTok{ylab =} \StringTok{"Valor"}\NormalTok{,}
     \AttributeTok{col =} \StringTok{"steelblue"}\NormalTok{,}
     \AttributeTok{lwd =} \DecValTok{2}\NormalTok{,}
     \AttributeTok{pch =} \DecValTok{19}\NormalTok{)}
\end{Highlighting}
\end{Shaded}

\begin{figure}

{\centering \includegraphics[width=0.8\linewidth]{introR_files/figure-latex/graf-linhas-1} 

}

\caption{Gráfico de linhas com pontos e linha.}\label{fig:graf-linhas}
\end{figure}

\section{Exercícios}\label{exercuxedcios-14}

\textbf{1.} No conjunto \texttt{iris}, crie um gráfico de barras com a contagem de cada espécie
(título ``Contagem de espécies'', eixo x ``Espécie'', eixo y ``Frequência absoluta'', cores
azul, verde e amarelo). Depois, refaça-o com frequências relativas.

\textbf{2.} No conjunto \texttt{airquality}, crie um gráfico de barras com a temperatura média
(\texttt{Temp}) por mês.
\emph{Sugestão:} Use a função \texttt{tapply()} para obter a média mensal. (Título ``Média
de temperatura por mês'', eixo x ``Mês'', eixo y ``Temperatura média (°F)'', cor azul.)

\textbf{3.} Crie um gráfico circular das espécies de flores do conjunto \texttt{iris},
identificando cada fatia com o nome da espécie e a respetiva percentagem.

\textbf{4.} No conjunto \texttt{Titanic}, crie um gráfico circular com a proporção de sobreviventes
e não sobreviventes.
\emph{Sugestão:} transforme \texttt{Titanic} numa \emph{data frame} com \texttt{as.data.frame()}.

\textbf{5.} Crie um histograma da variável \texttt{Sepal.Length} (conjunto \texttt{iris}). Personalize-o
com cor, ajuste o número de classes (\texttt{breaks}) e acrescente título e rótulos aos eixos.

\textbf{6.} Um investigador estuda a distribuição de alturas de plantas de duas espécies (A
e B), com 100 amostras de cada. Gere os dados com \texttt{rnorm()} (médias e desvios padrão
diferentes), desenhe um histograma \textbf{sobreposto} das duas distribuições (com
\texttt{add\ =\ TRUE}) e acrescente uma legenda. Use cores translúcidas para distinguir as
espécies:

\begin{Shaded}
\begin{Highlighting}[]
\NormalTok{altura\_A }\OtherTok{\textless{}{-}} \FunctionTok{rnorm}\NormalTok{(}\DecValTok{100}\NormalTok{, }\AttributeTok{mean =} \DecValTok{150}\NormalTok{, }\AttributeTok{sd =} \DecValTok{10}\NormalTok{)}
\NormalTok{altura\_B }\OtherTok{\textless{}{-}} \FunctionTok{rnorm}\NormalTok{(}\DecValTok{100}\NormalTok{, }\AttributeTok{mean =} \DecValTok{160}\NormalTok{, }\AttributeTok{sd =} \DecValTok{15}\NormalTok{)}

\CommentTok{\# Cores translúcidas: rgb(vermelho, verde, azul, transparência)}
\CommentTok{\# col = rgb(1, 0, 0, 0.5)  e  col = rgb(0, 0, 1, 0.5)}
\FunctionTok{legend}\NormalTok{(}\StringTok{"topright"}\NormalTok{,}
       \AttributeTok{legend =} \FunctionTok{c}\NormalTok{(}\StringTok{"Espécie A"}\NormalTok{, }\StringTok{"Espécie B"}\NormalTok{),}
       \AttributeTok{fill =} \FunctionTok{c}\NormalTok{(}\FunctionTok{rgb}\NormalTok{(}\DecValTok{1}\NormalTok{, }\DecValTok{0}\NormalTok{, }\DecValTok{0}\NormalTok{, }\FloatTok{0.5}\NormalTok{), }\FunctionTok{rgb}\NormalTok{(}\DecValTok{0}\NormalTok{, }\DecValTok{0}\NormalTok{, }\DecValTok{1}\NormalTok{, }\FloatTok{0.5}\NormalTok{)))}
\end{Highlighting}
\end{Shaded}

\textbf{7.} Um professor aplicou um exame a 200 alunos (notas de 0 a 20) e quer analisar a
distribuição das notas. Gere as notas com \texttt{rnorm(200,\ mean\ =\ 14,\ sd\ =\ 4)} (use
\texttt{set.seed(1234)}), desenhe o histograma, calcule a média, a mediana e o desvio padrão,
e marque a média e a mediana no gráfico com \texttt{abline()}. Junte uma legenda:

\begin{Shaded}
\begin{Highlighting}[]
\FunctionTok{legend}\NormalTok{(}\StringTok{"topright"}\NormalTok{,}
       \AttributeTok{legend =} \FunctionTok{c}\NormalTok{(}\FunctionTok{paste}\NormalTok{(}\StringTok{"Média:"}\NormalTok{, }\FunctionTok{round}\NormalTok{(media\_nota, }\DecValTok{2}\NormalTok{)),}
                  \FunctionTok{paste}\NormalTok{(}\StringTok{"Mediana:"}\NormalTok{, }\FunctionTok{round}\NormalTok{(mediana\_nota, }\DecValTok{2}\NormalTok{)),}
                  \FunctionTok{paste}\NormalTok{(}\StringTok{"Desvio padrão:"}\NormalTok{, }\FunctionTok{round}\NormalTok{(desvio\_nota, }\DecValTok{2}\NormalTok{))),}
       \AttributeTok{bty =} \StringTok{"n"}\NormalTok{)}
\end{Highlighting}
\end{Shaded}

\textbf{8.} No conjunto \texttt{airquality}, crie histogramas da variável \texttt{Ozono} (\texttt{Ozone})
separadamente para cada mês (de maio a setembro). Use \texttt{par(mfrow\ =\ c(3,\ 2))} para os
dispor no mesmo painel e compare a distribuição dos níveis de ozono ao longo dos meses.

\textbf{9.} No conjunto \texttt{iris}, crie uma caixa de bigodes que compare o comprimento da
pétala (\texttt{Petal.Length}) entre as três espécies. Personalize com cores diferentes por
grupo e acrescente título e rótulos.

\textbf{10.} No conjunto \texttt{airquality}, crie uma caixa de bigodes dos níveis de ozono
(\texttt{Ozone}) por mês (\texttt{Month}) e acrescente uma linha horizontal na mediana global com
\texttt{abline()}. Comente como o ozono varia ao longo do tempo.

\textbf{11.} Estude a relação entre a altura e o peso de um grupo de pessoas.

\begin{itemize}
\tightlist
\item
  Crie um vetor \texttt{altura} com 30 valores aleatórios entre 150 e 200 cm (use \texttt{runif()}).
\item
  Crie um vetor \texttt{peso} aproximadamente proporcional à altura:
  \texttt{peso\ \textless{}-\ altura\ *\ 0.4\ +\ runif(30,\ -5,\ 5)}.
\item
  Desenhe um gráfico de dispersão (altura no eixo x, peso no eixo y), com título e
  rótulos.
\item
  Que tendência observa? Comente a correlação entre as duas variáveis.
\end{itemize}

\textbf{12.} Analise o crescimento de uma população ao longo de 20 anos.

\begin{itemize}
\tightlist
\item
  Crie um vetor \texttt{anos} de 2000 a 2019.
\item
  Crie um vetor \texttt{populacao} a partir de 100 000 pessoas em 2000, com crescimento de 2\%
  ao ano e uma pequena flutuação aleatória:
  \texttt{populacao\ \textless{}-\ 100000\ *\ cumprod(1.02\ +\ runif(20,\ -0.01,\ 0.01))}.
\item
  Desenhe um gráfico de linhas com a evolução da população ao longo dos anos.
\end{itemize}

Os gráficos do R base são rápidos e diretos, ideais para a exploração. No próximo
capítulo aprofundamos a \textbf{manipulação de dados} --- o passo que quase sempre precede um
bom gráfico --- e, mais adiante, conheceremos o \texttt{ggplot2}, que leva a visualização a
outro nível de expressividade.

\chapter{Manipulação de dados}\label{manipulacao}

Ler dados para dentro do R (Capítulo~\ref{leitura}) é apenas o começo. Antes de os
analisar ou representar graficamente, quase sempre é preciso \emph{prepará-los}: escolher
colunas, filtrar linhas, criar variáveis novas, resumir por grupos, juntar tabelas de
origens diferentes. A este trabalho --- que costuma ocupar a maior parte do tempo de uma
análise --- chama-se \textbf{manipulação de dados}.

O R base já permite fazer tudo isto (vimo-lo no Capítulo~\ref{estruturas}), mas de
forma por vezes verbosa. Neste capítulo conhecemos o \texttt{dplyr}, um pacote do
\textbf{\texttt{tidyverse}} --- a coleção de pacotes, criada em torno de Hadley Wickham, que
uniformizou a ciência de dados em R. O \texttt{dplyr} oferece um pequeno conjunto de ``verbos'',
cada um com um nome que descreve o que faz, que se combinam com o operador \emph{pipe}
(Capítulo~\ref{pipe}) para exprimir manipulações complexas de forma legível.

Este capítulo usa vários pacotes do \texttt{tidyverse}. Instale-os uma vez com
\texttt{install.packages("tidyverse")} (ou, individualmente,
\texttt{install.packages(c("tibble",\ "dplyr",\ "stringr"))}) e carregue-os no início da sessão.

\section{Tibbles}\label{tibbles}

O \texttt{tidyverse} trabalha com uma versão modernizada do \emph{data frame} chamada \textbf{tibble}.
Um \emph{tibble} guarda os mesmos dados de um \emph{data frame} --- é, na verdade, um \emph{data frame} ---
mas com alguns comportamentos mais previsíveis e cómodos. As diferenças principais são:

{\def\LTcaptype{none} % do not increment counter
\begin{longtable}[]{@{}
  >{\raggedright\arraybackslash}p{(\linewidth - 4\tabcolsep) * \real{0.3333}}
  >{\raggedright\arraybackslash}p{(\linewidth - 4\tabcolsep) * \real{0.3333}}
  >{\raggedright\arraybackslash}p{(\linewidth - 4\tabcolsep) * \real{0.3333}}@{}}
\toprule\noalign{}
\begin{minipage}[b]{\linewidth}\raggedright
Característica
\end{minipage} & \begin{minipage}[b]{\linewidth}\raggedright
\emph{Data frame}
\end{minipage} & \begin{minipage}[b]{\linewidth}\raggedright
\emph{Tibble}
\end{minipage} \\
\midrule\noalign{}
\endhead
\bottomrule\noalign{}
\endlastfoot
Impressão na consola & Mostra todas as linhas & Mostra um resumo (10 linhas) e o tipo de cada coluna \\
Texto → fatores & Mantém texto como texto (desde o R 4.0) & Mantém sempre texto como texto \\
Nomes das colunas & Torna-os únicos e sem espaços & Aceita nomes não convencionais (com crases) \\
Indexação com \texttt{{[}} & Pode devolver um vetor ou um \emph{data frame} & Devolve sempre um \emph{tibble} \\
Integração & Base do R & Otimizado para o \texttt{tidyverse} \\
\end{longtable}
}

Lê-se com frequência que ``os \emph{data frames} convertem texto em fatores por omissão''.
Isto foi verdade durante muitos anos, mas \textbf{mudou no R 4.0} (2020): desde então,
\texttt{data.frame()} e \texttt{read.table()} usam \texttt{stringsAsFactors\ =\ FALSE}, ou seja, mantêm o
texto como texto, tal como os \emph{tibbles}. É um bom exemplo de como convém confirmar
afirmações sobre a linguagem na documentação atual, e não confiar em material antigo.

Criar um \emph{tibble} é como criar um \emph{data frame}, mas a função permite referir colunas
acabadas de definir:

\begin{Shaded}
\begin{Highlighting}[]
\FunctionTok{library}\NormalTok{(tibble)}

\NormalTok{tb }\OtherTok{\textless{}{-}} \FunctionTok{tibble}\NormalTok{(}
  \AttributeTok{x =} \DecValTok{1}\SpecialCharTok{:}\DecValTok{5}\NormalTok{,}
  \AttributeTok{y =} \FunctionTok{c}\NormalTok{(}\StringTok{"A"}\NormalTok{, }\StringTok{"B"}\NormalTok{, }\StringTok{"C"}\NormalTok{, }\StringTok{"D"}\NormalTok{, }\StringTok{"E"}\NormalTok{),}
  \AttributeTok{z =}\NormalTok{ x }\SpecialCharTok{*} \DecValTok{2}                       \CommentTok{\# pode usar{-}se a coluna \textquotesingle{}x\textquotesingle{} recém{-}criada}
\NormalTok{)}

\NormalTok{tb}
\DocumentationTok{\#\# \# A tibble: 5 × 3}
\DocumentationTok{\#\#       x y         z}
\DocumentationTok{\#\#   \textless{}int\textgreater{} \textless{}chr\textgreater{} \textless{}dbl\textgreater{}}
\DocumentationTok{\#\# 1     1 A         2}
\DocumentationTok{\#\# 2     2 B         4}
\DocumentationTok{\#\# 3     3 C         6}
\DocumentationTok{\#\# 4     4 D         8}
\DocumentationTok{\#\# 5     5 E        10}
\end{Highlighting}
\end{Shaded}

Para converter um \emph{data frame} já existente, usa-se \texttt{as\_tibble()}.

\section{\texorpdfstring{Os verbos do \texttt{dplyr}}{Os verbos do dplyr}}\label{os-verbos-do-dplyr}

O \texttt{dplyr} organiza a manipulação em torno de um punhado de verbos, cada um com uma
função clara:

\begin{itemize}
\tightlist
\item
  \texttt{select()} --- escolhe colunas;
\item
  \texttt{arrange()} --- ordena linhas;
\item
  \texttt{filter()} --- filtra linhas por condições;
\item
  \texttt{mutate()} --- cria ou modifica colunas;
\item
  \texttt{group\_by()} --- agrupa por categorias;
\item
  \texttt{summarize()} --- resume os dados (em geral, após \texttt{group\_by()}).
\end{itemize}

Neste capítulo aprendemos a selecionar, ordenar, filtrar, transformar, agrupar e
resumir dados com o \texttt{dplyr}, e a combinar tabelas com as funções de junção --- sempre com
o operador \emph{pipe} a tornar as sequências de operações legíveis.

Vamos explorá-los com o conjunto \texttt{starwars}, que acompanha o \texttt{dplyr} e descreve as
personagens da saga. Guardamo-lo em \texttt{sw}:

\begin{Shaded}
\begin{Highlighting}[]
\FunctionTok{library}\NormalTok{(dplyr)}

\NormalTok{sw }\OtherTok{\textless{}{-}}\NormalTok{ starwars}
\FunctionTok{glimpse}\NormalTok{(sw)   }\CommentTok{\# visão geral das colunas e dos seus tipos}
\DocumentationTok{\#\# Rows: 87}
\DocumentationTok{\#\# Columns: 14}
\DocumentationTok{\#\# $ name       \textless{}chr\textgreater{} "Luke Skywalker", "C{-}3PO", "R2{-}D2", "Darth Vader", "Leia Or\textasciitilde{}}
\DocumentationTok{\#\# $ height     \textless{}int\textgreater{} 172, 167, 96, 202, 150, 178, 165, 97, 183, 182, 188, 180, 2\textasciitilde{}}
\DocumentationTok{\#\# $ mass       \textless{}dbl\textgreater{} 77.0, 75.0, 32.0, 136.0, 49.0, 120.0, 75.0, 32.0, 84.0, 77.\textasciitilde{}}
\DocumentationTok{\#\# $ hair\_color \textless{}chr\textgreater{} "blond", NA, NA, "none", "brown", "brown, grey", "brown", N\textasciitilde{}}
\DocumentationTok{\#\# $ skin\_color \textless{}chr\textgreater{} "fair", "gold", "white, blue", "white", "light", "light", "\textasciitilde{}}
\DocumentationTok{\#\# $ eye\_color  \textless{}chr\textgreater{} "blue", "yellow", "red", "yellow", "brown", "blue", "blue",\textasciitilde{}}
\DocumentationTok{\#\# $ birth\_year \textless{}dbl\textgreater{} 19.0, 112.0, 33.0, 41.9, 19.0, 52.0, 47.0, NA, 24.0, 57.0, \textasciitilde{}}
\DocumentationTok{\#\# $ sex        \textless{}chr\textgreater{} "male", "none", "none", "male", "female", "male", "female",\textasciitilde{}}
\DocumentationTok{\#\# $ gender     \textless{}chr\textgreater{} "masculine", "masculine", "masculine", "masculine", "femini\textasciitilde{}}
\DocumentationTok{\#\# $ homeworld  \textless{}chr\textgreater{} "Tatooine", "Tatooine", "Naboo", "Tatooine", "Alderaan", "T\textasciitilde{}}
\DocumentationTok{\#\# $ species    \textless{}chr\textgreater{} "Human", "Droid", "Droid", "Human", "Human", "Human", "Huma\textasciitilde{}}
\DocumentationTok{\#\# $ films      \textless{}list\textgreater{} \textless{}"A New Hope", "The Empire Strikes Back", "Return of the J\textasciitilde{}}
\DocumentationTok{\#\# $ vehicles   \textless{}list\textgreater{} \textless{}"Snowspeeder", "Imperial Speeder Bike"\textgreater{}, \textless{}\textgreater{}, \textless{}\textgreater{}, \textless{}\textgreater{}, "Imp\textasciitilde{}}
\DocumentationTok{\#\# $ starships  \textless{}list\textgreater{} \textless{}"X{-}wing", "Imperial shuttle"\textgreater{}, \textless{}\textgreater{}, \textless{}\textgreater{}, "TIE Advanced x1",\textasciitilde{}}
\end{Highlighting}
\end{Shaded}

\subsection{\texorpdfstring{Selecionar colunas com \texttt{select()}}{Selecionar colunas com select()}}\label{selecionar-colunas-com-select}

A \texttt{select()} escolhe colunas pelo nome, sem aspas:

\begin{Shaded}
\begin{Highlighting}[]
\FunctionTok{select}\NormalTok{(sw, name)                       }\CommentTok{\# uma coluna}
\DocumentationTok{\#\# \# A tibble: 87 x 1}
\DocumentationTok{\#\#    name              }
\DocumentationTok{\#\#    \textless{}chr\textgreater{}             }
\DocumentationTok{\#\#  1 Luke Skywalker    }
\DocumentationTok{\#\#  2 C{-}3PO             }
\DocumentationTok{\#\#  3 R2{-}D2             }
\DocumentationTok{\#\#  4 Darth Vader       }
\DocumentationTok{\#\#  5 Leia Organa       }
\DocumentationTok{\#\#  6 Owen Lars         }
\DocumentationTok{\#\#  7 Beru Whitesun Lars}
\DocumentationTok{\#\#  8 R5{-}D4             }
\DocumentationTok{\#\#  9 Biggs Darklighter }
\DocumentationTok{\#\# 10 Obi{-}Wan Kenobi    }
\DocumentationTok{\#\# \# i 77 more rows}
\FunctionTok{select}\NormalTok{(sw, name, mass, hair\_color)     }\CommentTok{\# várias colunas}
\DocumentationTok{\#\# \# A tibble: 87 x 3}
\DocumentationTok{\#\#    name                mass hair\_color   }
\DocumentationTok{\#\#    \textless{}chr\textgreater{}              \textless{}dbl\textgreater{} \textless{}chr\textgreater{}        }
\DocumentationTok{\#\#  1 Luke Skywalker        77 blond        }
\DocumentationTok{\#\#  2 C{-}3PO                 75 \textless{}NA\textgreater{}         }
\DocumentationTok{\#\#  3 R2{-}D2                 32 \textless{}NA\textgreater{}         }
\DocumentationTok{\#\#  4 Darth Vader          136 none         }
\DocumentationTok{\#\#  5 Leia Organa           49 brown        }
\DocumentationTok{\#\#  6 Owen Lars            120 brown, grey  }
\DocumentationTok{\#\#  7 Beru Whitesun Lars    75 brown        }
\DocumentationTok{\#\#  8 R5{-}D4                 32 \textless{}NA\textgreater{}         }
\DocumentationTok{\#\#  9 Biggs Darklighter     84 black        }
\DocumentationTok{\#\# 10 Obi{-}Wan Kenobi        77 auburn, white}
\DocumentationTok{\#\# \# i 77 more rows}
\FunctionTok{select}\NormalTok{(sw, name}\SpecialCharTok{:}\NormalTok{hair\_color)            }\CommentTok{\# um intervalo de colunas}
\DocumentationTok{\#\# \# A tibble: 87 x 4}
\DocumentationTok{\#\#    name               height  mass hair\_color   }
\DocumentationTok{\#\#    \textless{}chr\textgreater{}               \textless{}int\textgreater{} \textless{}dbl\textgreater{} \textless{}chr\textgreater{}        }
\DocumentationTok{\#\#  1 Luke Skywalker        172    77 blond        }
\DocumentationTok{\#\#  2 C{-}3PO                 167    75 \textless{}NA\textgreater{}         }
\DocumentationTok{\#\#  3 R2{-}D2                  96    32 \textless{}NA\textgreater{}         }
\DocumentationTok{\#\#  4 Darth Vader           202   136 none         }
\DocumentationTok{\#\#  5 Leia Organa           150    49 brown        }
\DocumentationTok{\#\#  6 Owen Lars             178   120 brown, grey  }
\DocumentationTok{\#\#  7 Beru Whitesun Lars    165    75 brown        }
\DocumentationTok{\#\#  8 R5{-}D4                  97    32 \textless{}NA\textgreater{}         }
\DocumentationTok{\#\#  9 Biggs Darklighter     183    84 black        }
\DocumentationTok{\#\# 10 Obi{-}Wan Kenobi        182    77 auburn, white}
\DocumentationTok{\#\# \# i 77 more rows}
\end{Highlighting}
\end{Shaded}

Há funções auxiliares muito úteis: \texttt{starts\_with()}, \texttt{ends\_with()} e \texttt{contains()}
selecionam colunas pelo nome. Um sinal \texttt{-} remove colunas:

\begin{Shaded}
\begin{Highlighting}[]
\FunctionTok{select}\NormalTok{(sw, }\FunctionTok{ends\_with}\NormalTok{(}\StringTok{"color"}\NormalTok{))         }\CommentTok{\# colunas terminadas em "color"}
\DocumentationTok{\#\# \# A tibble: 87 x 3}
\DocumentationTok{\#\#    hair\_color    skin\_color  eye\_color}
\DocumentationTok{\#\#    \textless{}chr\textgreater{}         \textless{}chr\textgreater{}       \textless{}chr\textgreater{}    }
\DocumentationTok{\#\#  1 blond         fair        blue     }
\DocumentationTok{\#\#  2 \textless{}NA\textgreater{}          gold        yellow   }
\DocumentationTok{\#\#  3 \textless{}NA\textgreater{}          white, blue red      }
\DocumentationTok{\#\#  4 none          white       yellow   }
\DocumentationTok{\#\#  5 brown         light       brown    }
\DocumentationTok{\#\#  6 brown, grey   light       blue     }
\DocumentationTok{\#\#  7 brown         light       blue     }
\DocumentationTok{\#\#  8 \textless{}NA\textgreater{}          white, red  red      }
\DocumentationTok{\#\#  9 black         light       brown    }
\DocumentationTok{\#\# 10 auburn, white fair        blue{-}gray}
\DocumentationTok{\#\# \# i 77 more rows}
\FunctionTok{select}\NormalTok{(sw, }\SpecialCharTok{{-}}\FunctionTok{ends\_with}\NormalTok{(}\StringTok{"color"}\NormalTok{))        }\CommentTok{\# todas menos essas}
\DocumentationTok{\#\# \# A tibble: 87 x 11}
\DocumentationTok{\#\#    name    height  mass birth\_year sex   gender homeworld species films vehicles}
\DocumentationTok{\#\#    \textless{}chr\textgreater{}    \textless{}int\textgreater{} \textless{}dbl\textgreater{}      \textless{}dbl\textgreater{} \textless{}chr\textgreater{} \textless{}chr\textgreater{}  \textless{}chr\textgreater{}     \textless{}chr\textgreater{}   \textless{}lis\textgreater{} \textless{}list\textgreater{}  }
\DocumentationTok{\#\#  1 Luke S\textasciitilde{}    172    77       19   male  mascu\textasciitilde{} Tatooine  Human   \textless{}chr\textgreater{} \textless{}chr\textgreater{}   }
\DocumentationTok{\#\#  2 C{-}3PO      167    75      112   none  mascu\textasciitilde{} Tatooine  Droid   \textless{}chr\textgreater{} \textless{}chr\textgreater{}   }
\DocumentationTok{\#\#  3 R2{-}D2       96    32       33   none  mascu\textasciitilde{} Naboo     Droid   \textless{}chr\textgreater{} \textless{}chr\textgreater{}   }
\DocumentationTok{\#\#  4 Darth \textasciitilde{}    202   136       41.9 male  mascu\textasciitilde{} Tatooine  Human   \textless{}chr\textgreater{} \textless{}chr\textgreater{}   }
\DocumentationTok{\#\#  5 Leia O\textasciitilde{}    150    49       19   fema\textasciitilde{} femin\textasciitilde{} Alderaan  Human   \textless{}chr\textgreater{} \textless{}chr\textgreater{}   }
\DocumentationTok{\#\#  6 Owen L\textasciitilde{}    178   120       52   male  mascu\textasciitilde{} Tatooine  Human   \textless{}chr\textgreater{} \textless{}chr\textgreater{}   }
\DocumentationTok{\#\#  7 Beru W\textasciitilde{}    165    75       47   fema\textasciitilde{} femin\textasciitilde{} Tatooine  Human   \textless{}chr\textgreater{} \textless{}chr\textgreater{}   }
\DocumentationTok{\#\#  8 R5{-}D4       97    32       NA   none  mascu\textasciitilde{} Tatooine  Droid   \textless{}chr\textgreater{} \textless{}chr\textgreater{}   }
\DocumentationTok{\#\#  9 Biggs \textasciitilde{}    183    84       24   male  mascu\textasciitilde{} Tatooine  Human   \textless{}chr\textgreater{} \textless{}chr\textgreater{}   }
\DocumentationTok{\#\# 10 Obi{-}Wa\textasciitilde{}    182    77       57   male  mascu\textasciitilde{} Stewjon   Human   \textless{}chr\textgreater{} \textless{}chr\textgreater{}   }
\DocumentationTok{\#\# \# i 77 more rows}
\DocumentationTok{\#\# \# i 1 more variable: starships \textless{}list\textgreater{}}
\end{Highlighting}
\end{Shaded}

\subsubsection{Exercícios}\label{exercuxedcios-15}

Use a base \texttt{sw} nos exercícios seguintes.

\textbf{1.} Aplique \texttt{glimpse()} à base \texttt{sw}. O que faz esta função?

\textbf{2.} Crie uma tabela \texttt{sw\_simples} apenas com as colunas \texttt{name}, \texttt{gender} e \texttt{films}.

\textbf{3.} Selecione as colunas \texttt{hair\_color}, \texttt{skin\_color} e \texttt{eye\_color} usando a função
auxiliar \texttt{contains()}.

\textbf{4.} Escreva todas as formas diferentes que conseguir de devolver a base \texttt{sw} \textbf{sem}
as colunas \texttt{hair\_color}, \texttt{skin\_color} e \texttt{eye\_color}.

\subsection{\texorpdfstring{Ordenar com \texttt{arrange()}}{Ordenar com arrange()}}\label{ordenar-com-arrange}

A \texttt{arrange()} ordena as linhas por uma ou mais colunas; \texttt{desc()} inverte a ordem:

\begin{Shaded}
\begin{Highlighting}[]
\FunctionTok{arrange}\NormalTok{(sw, mass)                        }\CommentTok{\# massa crescente}
\DocumentationTok{\#\# \# A tibble: 87 x 14}
\DocumentationTok{\#\#    name     height  mass hair\_color skin\_color eye\_color birth\_year sex   gender}
\DocumentationTok{\#\#    \textless{}chr\textgreater{}     \textless{}int\textgreater{} \textless{}dbl\textgreater{} \textless{}chr\textgreater{}      \textless{}chr\textgreater{}      \textless{}chr\textgreater{}          \textless{}dbl\textgreater{} \textless{}chr\textgreater{} \textless{}chr\textgreater{} }
\DocumentationTok{\#\#  1 Ratts T\textasciitilde{}     79    15 none       grey, blue unknown           NA male  mascu\textasciitilde{}}
\DocumentationTok{\#\#  2 Yoda         66    17 white      green      brown            896 male  mascu\textasciitilde{}}
\DocumentationTok{\#\#  3 Wicket \textasciitilde{}     88    20 brown      brown      brown              8 male  mascu\textasciitilde{}}
\DocumentationTok{\#\#  4 R2{-}D2        96    32 \textless{}NA\textgreater{}       white, bl\textasciitilde{} red               33 none  mascu\textasciitilde{}}
\DocumentationTok{\#\#  5 R5{-}D4        97    32 \textless{}NA\textgreater{}       white, red red               NA none  mascu\textasciitilde{}}
\DocumentationTok{\#\#  6 Sebulba     112    40 none       grey, red  orange            NA male  mascu\textasciitilde{}}
\DocumentationTok{\#\#  7 Padmé A\textasciitilde{}    185    45 brown      light      brown             46 fema\textasciitilde{} femin\textasciitilde{}}
\DocumentationTok{\#\#  8 Dud Bolt     94    45 none       blue, grey yellow            NA male  mascu\textasciitilde{}}
\DocumentationTok{\#\#  9 Wat Tam\textasciitilde{}    193    48 none       green, gr\textasciitilde{} unknown           NA male  mascu\textasciitilde{}}
\DocumentationTok{\#\# 10 Sly Moo\textasciitilde{}    178    48 none       pale       white             NA \textless{}NA\textgreater{}  \textless{}NA\textgreater{}  }
\DocumentationTok{\#\# \# i 77 more rows}
\DocumentationTok{\#\# \# i 5 more variables: homeworld \textless{}chr\textgreater{}, species \textless{}chr\textgreater{}, films \textless{}list\textgreater{},}
\DocumentationTok{\#\# \#   vehicles \textless{}list\textgreater{}, starships \textless{}list\textgreater{}}
\FunctionTok{arrange}\NormalTok{(sw, }\FunctionTok{desc}\NormalTok{(mass))                  }\CommentTok{\# massa decrescente}
\DocumentationTok{\#\# \# A tibble: 87 x 14}
\DocumentationTok{\#\#    name     height  mass hair\_color skin\_color eye\_color birth\_year sex   gender}
\DocumentationTok{\#\#    \textless{}chr\textgreater{}     \textless{}int\textgreater{} \textless{}dbl\textgreater{} \textless{}chr\textgreater{}      \textless{}chr\textgreater{}      \textless{}chr\textgreater{}          \textless{}dbl\textgreater{} \textless{}chr\textgreater{} \textless{}chr\textgreater{} }
\DocumentationTok{\#\#  1 Jabba D\textasciitilde{}    175  1358 \textless{}NA\textgreater{}       green{-}tan\textasciitilde{} orange         600   herm\textasciitilde{} mascu\textasciitilde{}}
\DocumentationTok{\#\#  2 Grievous    216   159 none       brown, wh\textasciitilde{} green, y\textasciitilde{}       NA   male  mascu\textasciitilde{}}
\DocumentationTok{\#\#  3 IG{-}88       200   140 none       metal      red             15   none  mascu\textasciitilde{}}
\DocumentationTok{\#\#  4 Darth V\textasciitilde{}    202   136 none       white      yellow          41.9 male  mascu\textasciitilde{}}
\DocumentationTok{\#\#  5 Tarfful     234   136 brown      brown      blue            NA   male  mascu\textasciitilde{}}
\DocumentationTok{\#\#  6 Owen La\textasciitilde{}    178   120 brown, gr\textasciitilde{} light      blue            52   male  mascu\textasciitilde{}}
\DocumentationTok{\#\#  7 Bossk       190   113 none       green      red             53   male  mascu\textasciitilde{}}
\DocumentationTok{\#\#  8 Chewbac\textasciitilde{}    228   112 brown      unknown    blue           200   male  mascu\textasciitilde{}}
\DocumentationTok{\#\#  9 Jek Ton\textasciitilde{}    180   110 brown      fair       blue            NA   \textless{}NA\textgreater{}  \textless{}NA\textgreater{}  }
\DocumentationTok{\#\# 10 Dexter \textasciitilde{}    198   102 none       brown      yellow          NA   male  mascu\textasciitilde{}}
\DocumentationTok{\#\# \# i 77 more rows}
\DocumentationTok{\#\# \# i 5 more variables: homeworld \textless{}chr\textgreater{}, species \textless{}chr\textgreater{}, films \textless{}list\textgreater{},}
\DocumentationTok{\#\# \#   vehicles \textless{}list\textgreater{}, starships \textless{}list\textgreater{}}
\FunctionTok{arrange}\NormalTok{(sw, }\FunctionTok{desc}\NormalTok{(height), }\FunctionTok{desc}\NormalTok{(mass))    }\CommentTok{\# por altura e, em empate, por massa}
\DocumentationTok{\#\# \# A tibble: 87 x 14}
\DocumentationTok{\#\#    name     height  mass hair\_color skin\_color eye\_color birth\_year sex   gender}
\DocumentationTok{\#\#    \textless{}chr\textgreater{}     \textless{}int\textgreater{} \textless{}dbl\textgreater{} \textless{}chr\textgreater{}      \textless{}chr\textgreater{}      \textless{}chr\textgreater{}          \textless{}dbl\textgreater{} \textless{}chr\textgreater{} \textless{}chr\textgreater{} }
\DocumentationTok{\#\#  1 Yarael \textasciitilde{}    264    NA none       white      yellow          NA   male  mascu\textasciitilde{}}
\DocumentationTok{\#\#  2 Tarfful     234   136 brown      brown      blue            NA   male  mascu\textasciitilde{}}
\DocumentationTok{\#\#  3 Lama Su     229    88 none       grey       black           NA   male  mascu\textasciitilde{}}
\DocumentationTok{\#\#  4 Chewbac\textasciitilde{}    228   112 brown      unknown    blue           200   male  mascu\textasciitilde{}}
\DocumentationTok{\#\#  5 Roos Ta\textasciitilde{}    224    82 none       grey       orange          NA   male  mascu\textasciitilde{}}
\DocumentationTok{\#\#  6 Grievous    216   159 none       brown, wh\textasciitilde{} green, y\textasciitilde{}       NA   male  mascu\textasciitilde{}}
\DocumentationTok{\#\#  7 Taun We     213    NA none       grey       black           NA   fema\textasciitilde{} femin\textasciitilde{}}
\DocumentationTok{\#\#  8 Tion Me\textasciitilde{}    206    80 none       grey       black           NA   male  mascu\textasciitilde{}}
\DocumentationTok{\#\#  9 Rugor N\textasciitilde{}    206    NA none       green      orange          NA   male  mascu\textasciitilde{}}
\DocumentationTok{\#\# 10 Darth V\textasciitilde{}    202   136 none       white      yellow          41.9 male  mascu\textasciitilde{}}
\DocumentationTok{\#\# \# i 77 more rows}
\DocumentationTok{\#\# \# i 5 more variables: homeworld \textless{}chr\textgreater{}, species \textless{}chr\textgreater{}, films \textless{}list\textgreater{},}
\DocumentationTok{\#\# \#   vehicles \textless{}list\textgreater{}, starships \textless{}list\textgreater{}}
\end{Highlighting}
\end{Shaded}

\subsubsection{Exercícios}\label{exercuxedcios-16}

\textbf{1.} Ordene \texttt{sw} por \texttt{mass} (crescente) e, em empate, por \texttt{birth\_year}
(decrescente). Guarde em \texttt{sw\_ordenados}.

\textbf{2.} Selecione as colunas \texttt{name} e \texttt{birth\_year} e ordene por \texttt{birth\_year}
decrescente. Faça-o de três formas: com funções aninhadas, com um objeto intermédio e
com o \emph{pipe}.

\subsection{\texorpdfstring{Filtrar linhas com \texttt{filter()}}{Filtrar linhas com filter()}}\label{filtrar-linhas-com-filter}

A \texttt{filter()} seleciona linhas que satisfazem uma ou mais condições lógicas:

\begin{Shaded}
\begin{Highlighting}[]
\NormalTok{sw }\SpecialCharTok{\%\textgreater{}\%} \FunctionTok{filter}\NormalTok{(height }\SpecialCharTok{\textgreater{}} \DecValTok{170}\NormalTok{)                       }\CommentTok{\# uma condição}
\DocumentationTok{\#\# \# A tibble: 55 x 14}
\DocumentationTok{\#\#    name     height  mass hair\_color skin\_color eye\_color birth\_year sex   gender}
\DocumentationTok{\#\#    \textless{}chr\textgreater{}     \textless{}int\textgreater{} \textless{}dbl\textgreater{} \textless{}chr\textgreater{}      \textless{}chr\textgreater{}      \textless{}chr\textgreater{}          \textless{}dbl\textgreater{} \textless{}chr\textgreater{} \textless{}chr\textgreater{} }
\DocumentationTok{\#\#  1 Luke Sk\textasciitilde{}    172    77 blond      fair       blue            19   male  mascu\textasciitilde{}}
\DocumentationTok{\#\#  2 Darth V\textasciitilde{}    202   136 none       white      yellow          41.9 male  mascu\textasciitilde{}}
\DocumentationTok{\#\#  3 Owen La\textasciitilde{}    178   120 brown, gr\textasciitilde{} light      blue            52   male  mascu\textasciitilde{}}
\DocumentationTok{\#\#  4 Biggs D\textasciitilde{}    183    84 black      light      brown           24   male  mascu\textasciitilde{}}
\DocumentationTok{\#\#  5 Obi{-}Wan\textasciitilde{}    182    77 auburn, w\textasciitilde{} fair       blue{-}gray       57   male  mascu\textasciitilde{}}
\DocumentationTok{\#\#  6 Anakin \textasciitilde{}    188    84 blond      fair       blue            41.9 male  mascu\textasciitilde{}}
\DocumentationTok{\#\#  7 Wilhuff\textasciitilde{}    180    NA auburn, g\textasciitilde{} fair       blue            64   male  mascu\textasciitilde{}}
\DocumentationTok{\#\#  8 Chewbac\textasciitilde{}    228   112 brown      unknown    blue           200   male  mascu\textasciitilde{}}
\DocumentationTok{\#\#  9 Han Solo    180    80 brown      fair       brown           29   male  mascu\textasciitilde{}}
\DocumentationTok{\#\# 10 Greedo      173    74 \textless{}NA\textgreater{}       green      black           44   male  mascu\textasciitilde{}}
\DocumentationTok{\#\# \# i 45 more rows}
\DocumentationTok{\#\# \# i 5 more variables: homeworld \textless{}chr\textgreater{}, species \textless{}chr\textgreater{}, films \textless{}list\textgreater{},}
\DocumentationTok{\#\# \#   vehicles \textless{}list\textgreater{}, starships \textless{}list\textgreater{}}
\NormalTok{sw }\SpecialCharTok{\%\textgreater{}\%} \FunctionTok{filter}\NormalTok{(height }\SpecialCharTok{\textgreater{}} \DecValTok{170}\NormalTok{) }\SpecialCharTok{\%\textgreater{}\%} \FunctionTok{select}\NormalTok{(name, height)}
\DocumentationTok{\#\# \# A tibble: 55 x 2}
\DocumentationTok{\#\#    name              height}
\DocumentationTok{\#\#    \textless{}chr\textgreater{}              \textless{}int\textgreater{}}
\DocumentationTok{\#\#  1 Luke Skywalker       172}
\DocumentationTok{\#\#  2 Darth Vader          202}
\DocumentationTok{\#\#  3 Owen Lars            178}
\DocumentationTok{\#\#  4 Biggs Darklighter    183}
\DocumentationTok{\#\#  5 Obi{-}Wan Kenobi       182}
\DocumentationTok{\#\#  6 Anakin Skywalker     188}
\DocumentationTok{\#\#  7 Wilhuff Tarkin       180}
\DocumentationTok{\#\#  8 Chewbacca            228}
\DocumentationTok{\#\#  9 Han Solo             180}
\DocumentationTok{\#\# 10 Greedo               173}
\DocumentationTok{\#\# \# i 45 more rows}
\NormalTok{sw }\SpecialCharTok{\%\textgreater{}\%} \FunctionTok{filter}\NormalTok{(height }\SpecialCharTok{\textgreater{}} \DecValTok{170}\NormalTok{, mass }\SpecialCharTok{\textgreater{}} \DecValTok{80}\NormalTok{)            }\CommentTok{\# duas condições (E lógico)}
\DocumentationTok{\#\# \# A tibble: 21 x 14}
\DocumentationTok{\#\#    name     height  mass hair\_color skin\_color eye\_color birth\_year sex   gender}
\DocumentationTok{\#\#    \textless{}chr\textgreater{}     \textless{}int\textgreater{} \textless{}dbl\textgreater{} \textless{}chr\textgreater{}      \textless{}chr\textgreater{}      \textless{}chr\textgreater{}          \textless{}dbl\textgreater{} \textless{}chr\textgreater{} \textless{}chr\textgreater{} }
\DocumentationTok{\#\#  1 Darth V\textasciitilde{}    202   136 none       white      yellow          41.9 male  mascu\textasciitilde{}}
\DocumentationTok{\#\#  2 Owen La\textasciitilde{}    178   120 brown, gr\textasciitilde{} light      blue            52   male  mascu\textasciitilde{}}
\DocumentationTok{\#\#  3 Biggs D\textasciitilde{}    183    84 black      light      brown           24   male  mascu\textasciitilde{}}
\DocumentationTok{\#\#  4 Anakin \textasciitilde{}    188    84 blond      fair       blue            41.9 male  mascu\textasciitilde{}}
\DocumentationTok{\#\#  5 Chewbac\textasciitilde{}    228   112 brown      unknown    blue           200   male  mascu\textasciitilde{}}
\DocumentationTok{\#\#  6 Jabba D\textasciitilde{}    175  1358 \textless{}NA\textgreater{}       green{-}tan\textasciitilde{} orange         600   herm\textasciitilde{} mascu\textasciitilde{}}
\DocumentationTok{\#\#  7 Jek Ton\textasciitilde{}    180   110 brown      fair       blue            NA   \textless{}NA\textgreater{}  \textless{}NA\textgreater{}  }
\DocumentationTok{\#\#  8 IG{-}88       200   140 none       metal      red             15   none  mascu\textasciitilde{}}
\DocumentationTok{\#\#  9 Bossk       190   113 none       green      red             53   male  mascu\textasciitilde{}}
\DocumentationTok{\#\# 10 Ackbar      180    83 none       brown mot\textasciitilde{} orange          41   male  mascu\textasciitilde{}}
\DocumentationTok{\#\# \# i 11 more rows}
\DocumentationTok{\#\# \# i 5 more variables: homeworld \textless{}chr\textgreater{}, species \textless{}chr\textgreater{}, films \textless{}list\textgreater{},}
\DocumentationTok{\#\# \#   vehicles \textless{}list\textgreater{}, starships \textless{}list\textgreater{}}
\end{Highlighting}
\end{Shaded}

Para variáveis categóricas, comparamos com \texttt{==} ou testamos a pertença a um conjunto
com \texttt{\%in\%}:

\begin{Shaded}
\begin{Highlighting}[]
\NormalTok{sw }\SpecialCharTok{\%\textgreater{}\%} \FunctionTok{filter}\NormalTok{(hair\_color }\SpecialCharTok{==} \StringTok{"black"} \SpecialCharTok{|}\NormalTok{ hair\_color }\SpecialCharTok{==} \StringTok{"brown"}\NormalTok{)}
\DocumentationTok{\#\# \# A tibble: 31 x 14}
\DocumentationTok{\#\#    name     height  mass hair\_color skin\_color eye\_color birth\_year sex   gender}
\DocumentationTok{\#\#    \textless{}chr\textgreater{}     \textless{}int\textgreater{} \textless{}dbl\textgreater{} \textless{}chr\textgreater{}      \textless{}chr\textgreater{}      \textless{}chr\textgreater{}          \textless{}dbl\textgreater{} \textless{}chr\textgreater{} \textless{}chr\textgreater{} }
\DocumentationTok{\#\#  1 Leia Or\textasciitilde{}    150  49   brown      light      brown           19   fema\textasciitilde{} femin\textasciitilde{}}
\DocumentationTok{\#\#  2 Beru Wh\textasciitilde{}    165  75   brown      light      blue            47   fema\textasciitilde{} femin\textasciitilde{}}
\DocumentationTok{\#\#  3 Biggs D\textasciitilde{}    183  84   black      light      brown           24   male  mascu\textasciitilde{}}
\DocumentationTok{\#\#  4 Chewbac\textasciitilde{}    228 112   brown      unknown    blue           200   male  mascu\textasciitilde{}}
\DocumentationTok{\#\#  5 Han Solo    180  80   brown      fair       brown           29   male  mascu\textasciitilde{}}
\DocumentationTok{\#\#  6 Wedge A\textasciitilde{}    170  77   brown      fair       hazel           21   male  mascu\textasciitilde{}}
\DocumentationTok{\#\#  7 Jek Ton\textasciitilde{}    180 110   brown      fair       blue            NA   \textless{}NA\textgreater{}  \textless{}NA\textgreater{}  }
\DocumentationTok{\#\#  8 Boba Fe\textasciitilde{}    183  78.2 black      fair       brown           31.5 male  mascu\textasciitilde{}}
\DocumentationTok{\#\#  9 Lando C\textasciitilde{}    177  79   black      dark       brown           31   male  mascu\textasciitilde{}}
\DocumentationTok{\#\# 10 Arvel C\textasciitilde{}     NA  NA   brown      fair       brown           NA   male  mascu\textasciitilde{}}
\DocumentationTok{\#\# \# i 21 more rows}
\DocumentationTok{\#\# \# i 5 more variables: homeworld \textless{}chr\textgreater{}, species \textless{}chr\textgreater{}, films \textless{}list\textgreater{},}
\DocumentationTok{\#\# \#   vehicles \textless{}list\textgreater{}, starships \textless{}list\textgreater{}}
\NormalTok{sw }\SpecialCharTok{\%\textgreater{}\%} \FunctionTok{filter}\NormalTok{(hair\_color }\SpecialCharTok{\%in\%} \FunctionTok{c}\NormalTok{(}\StringTok{"black"}\NormalTok{, }\StringTok{"brown"}\NormalTok{))   }\CommentTok{\# forma equivalente, mais curta}
\DocumentationTok{\#\# \# A tibble: 31 x 14}
\DocumentationTok{\#\#    name     height  mass hair\_color skin\_color eye\_color birth\_year sex   gender}
\DocumentationTok{\#\#    \textless{}chr\textgreater{}     \textless{}int\textgreater{} \textless{}dbl\textgreater{} \textless{}chr\textgreater{}      \textless{}chr\textgreater{}      \textless{}chr\textgreater{}          \textless{}dbl\textgreater{} \textless{}chr\textgreater{} \textless{}chr\textgreater{} }
\DocumentationTok{\#\#  1 Leia Or\textasciitilde{}    150  49   brown      light      brown           19   fema\textasciitilde{} femin\textasciitilde{}}
\DocumentationTok{\#\#  2 Beru Wh\textasciitilde{}    165  75   brown      light      blue            47   fema\textasciitilde{} femin\textasciitilde{}}
\DocumentationTok{\#\#  3 Biggs D\textasciitilde{}    183  84   black      light      brown           24   male  mascu\textasciitilde{}}
\DocumentationTok{\#\#  4 Chewbac\textasciitilde{}    228 112   brown      unknown    blue           200   male  mascu\textasciitilde{}}
\DocumentationTok{\#\#  5 Han Solo    180  80   brown      fair       brown           29   male  mascu\textasciitilde{}}
\DocumentationTok{\#\#  6 Wedge A\textasciitilde{}    170  77   brown      fair       hazel           21   male  mascu\textasciitilde{}}
\DocumentationTok{\#\#  7 Jek Ton\textasciitilde{}    180 110   brown      fair       blue            NA   \textless{}NA\textgreater{}  \textless{}NA\textgreater{}  }
\DocumentationTok{\#\#  8 Boba Fe\textasciitilde{}    183  78.2 black      fair       brown           31.5 male  mascu\textasciitilde{}}
\DocumentationTok{\#\#  9 Lando C\textasciitilde{}    177  79   black      dark       brown           31   male  mascu\textasciitilde{}}
\DocumentationTok{\#\# 10 Arvel C\textasciitilde{}     NA  NA   brown      fair       brown           NA   male  mascu\textasciitilde{}}
\DocumentationTok{\#\# \# i 21 more rows}
\DocumentationTok{\#\# \# i 5 more variables: homeworld \textless{}chr\textgreater{}, species \textless{}chr\textgreater{}, films \textless{}list\textgreater{},}
\DocumentationTok{\#\# \#   vehicles \textless{}list\textgreater{}, starships \textless{}list\textgreater{}}
\end{Highlighting}
\end{Shaded}

Quando queremos procurar um padrão \emph{dentro} do texto --- e não uma correspondência
exata --- recorremos à \texttt{str\_detect()}, do pacote \texttt{stringr}, que indica, para cada
elemento, se contém o padrão:

\begin{Shaded}
\begin{Highlighting}[]
\FunctionTok{library}\NormalTok{(stringr)}

\NormalTok{sw }\SpecialCharTok{\%\textgreater{}\%} \FunctionTok{filter}\NormalTok{(}\FunctionTok{str\_detect}\NormalTok{(hair\_color, }\StringTok{"grey"}\NormalTok{))   }\CommentTok{\# cabelo que contenha "grey"}
\DocumentationTok{\#\# \# A tibble: 3 x 14}
\DocumentationTok{\#\#   name      height  mass hair\_color skin\_color eye\_color birth\_year sex   gender}
\DocumentationTok{\#\#   \textless{}chr\textgreater{}      \textless{}int\textgreater{} \textless{}dbl\textgreater{} \textless{}chr\textgreater{}      \textless{}chr\textgreater{}      \textless{}chr\textgreater{}          \textless{}dbl\textgreater{} \textless{}chr\textgreater{} \textless{}chr\textgreater{} }
\DocumentationTok{\#\# 1 Owen Lars    178   120 brown, gr\textasciitilde{} light      blue              52 male  mascu\textasciitilde{}}
\DocumentationTok{\#\# 2 Wilhuff \textasciitilde{}    180    NA auburn, g\textasciitilde{} fair       blue              64 male  mascu\textasciitilde{}}
\DocumentationTok{\#\# 3 Palpatine    170    75 grey       pale       yellow            82 male  mascu\textasciitilde{}}
\DocumentationTok{\#\# \# i 5 more variables: homeworld \textless{}chr\textgreater{}, species \textless{}chr\textgreater{}, films \textless{}list\textgreater{},}
\DocumentationTok{\#\# \#   vehicles \textless{}list\textgreater{}, starships \textless{}list\textgreater{}}
\end{Highlighting}
\end{Shaded}

\subsubsection{Exercícios}\label{exercuxedcios-17}

Use a base \texttt{sw}.

\textbf{1.} Crie um objeto \texttt{humanos} apenas com personagens humanas.

\textbf{2.} Crie um objeto \texttt{altos\_fortes} com personagens de mais de 200 cm e mais de
100 kg.

\textbf{3.} Devolva tabelas (\emph{tibbles}) apenas com:

\begin{enumerate}
\def\labelenumi{(\alph{enumi})}
\tightlist
\item
  personagens humanas nascidas mais de 100 anos antes da batalha de Yavin
  (\texttt{birth\_year\ \textgreater{}\ 100});
\item
  personagens cuja cor (de pele ou olhos) contenha \texttt{"light"} ou \texttt{"red"};
\item
  personagens com mais de 100 kg, ordenadas por altura decrescente, mostrando só
  \texttt{name}, \texttt{mass} e \texttt{height};
\item
  personagens da espécie \texttt{"Human"} ou \texttt{"Droid"} com mais de 170 cm;
\item
  personagens sem informação de altura \textbf{e} de massa (\texttt{NA} em ambas).
\end{enumerate}

\subsection{\texorpdfstring{Criar e modificar colunas com \texttt{mutate()}}{Criar e modificar colunas com mutate()}}\label{criar-e-modificar-colunas-com-mutate}

A \texttt{mutate()} cria colunas novas ou altera colunas existentes. Se o nome já existir, a
coluna é substituída; se for novo, a coluna é acrescentada no fim. Por exemplo, para
converter a altura de centímetros para metros:

\begin{Shaded}
\begin{Highlighting}[]
\NormalTok{sw }\SpecialCharTok{\%\textgreater{}\%} \FunctionTok{mutate}\NormalTok{(}\AttributeTok{height =}\NormalTok{ height }\SpecialCharTok{/} \DecValTok{100}\NormalTok{)          }\CommentTok{\# substitui a coluna existente}
\DocumentationTok{\#\# \# A tibble: 87 x 14}
\DocumentationTok{\#\#    name     height  mass hair\_color skin\_color eye\_color birth\_year sex   gender}
\DocumentationTok{\#\#    \textless{}chr\textgreater{}     \textless{}dbl\textgreater{} \textless{}dbl\textgreater{} \textless{}chr\textgreater{}      \textless{}chr\textgreater{}      \textless{}chr\textgreater{}          \textless{}dbl\textgreater{} \textless{}chr\textgreater{} \textless{}chr\textgreater{} }
\DocumentationTok{\#\#  1 Luke Sk\textasciitilde{}   1.72    77 blond      fair       blue            19   male  mascu\textasciitilde{}}
\DocumentationTok{\#\#  2 C{-}3PO      1.67    75 \textless{}NA\textgreater{}       gold       yellow         112   none  mascu\textasciitilde{}}
\DocumentationTok{\#\#  3 R2{-}D2      0.96    32 \textless{}NA\textgreater{}       white, bl\textasciitilde{} red             33   none  mascu\textasciitilde{}}
\DocumentationTok{\#\#  4 Darth V\textasciitilde{}   2.02   136 none       white      yellow          41.9 male  mascu\textasciitilde{}}
\DocumentationTok{\#\#  5 Leia Or\textasciitilde{}   1.5     49 brown      light      brown           19   fema\textasciitilde{} femin\textasciitilde{}}
\DocumentationTok{\#\#  6 Owen La\textasciitilde{}   1.78   120 brown, gr\textasciitilde{} light      blue            52   male  mascu\textasciitilde{}}
\DocumentationTok{\#\#  7 Beru Wh\textasciitilde{}   1.65    75 brown      light      blue            47   fema\textasciitilde{} femin\textasciitilde{}}
\DocumentationTok{\#\#  8 R5{-}D4      0.97    32 \textless{}NA\textgreater{}       white, red red             NA   none  mascu\textasciitilde{}}
\DocumentationTok{\#\#  9 Biggs D\textasciitilde{}   1.83    84 black      light      brown           24   male  mascu\textasciitilde{}}
\DocumentationTok{\#\# 10 Obi{-}Wan\textasciitilde{}   1.82    77 auburn, w\textasciitilde{} fair       blue{-}gray       57   male  mascu\textasciitilde{}}
\DocumentationTok{\#\# \# i 77 more rows}
\DocumentationTok{\#\# \# i 5 more variables: homeworld \textless{}chr\textgreater{}, species \textless{}chr\textgreater{}, films \textless{}list\textgreater{},}
\DocumentationTok{\#\# \#   vehicles \textless{}list\textgreater{}, starships \textless{}list\textgreater{}}
\NormalTok{sw }\SpecialCharTok{\%\textgreater{}\%} \FunctionTok{mutate}\NormalTok{(}\AttributeTok{height\_meters =}\NormalTok{ height }\SpecialCharTok{/} \DecValTok{100}\NormalTok{)   }\CommentTok{\# cria uma nova coluna}
\DocumentationTok{\#\# \# A tibble: 87 x 15}
\DocumentationTok{\#\#    name     height  mass hair\_color skin\_color eye\_color birth\_year sex   gender}
\DocumentationTok{\#\#    \textless{}chr\textgreater{}     \textless{}int\textgreater{} \textless{}dbl\textgreater{} \textless{}chr\textgreater{}      \textless{}chr\textgreater{}      \textless{}chr\textgreater{}          \textless{}dbl\textgreater{} \textless{}chr\textgreater{} \textless{}chr\textgreater{} }
\DocumentationTok{\#\#  1 Luke Sk\textasciitilde{}    172    77 blond      fair       blue            19   male  mascu\textasciitilde{}}
\DocumentationTok{\#\#  2 C{-}3PO       167    75 \textless{}NA\textgreater{}       gold       yellow         112   none  mascu\textasciitilde{}}
\DocumentationTok{\#\#  3 R2{-}D2        96    32 \textless{}NA\textgreater{}       white, bl\textasciitilde{} red             33   none  mascu\textasciitilde{}}
\DocumentationTok{\#\#  4 Darth V\textasciitilde{}    202   136 none       white      yellow          41.9 male  mascu\textasciitilde{}}
\DocumentationTok{\#\#  5 Leia Or\textasciitilde{}    150    49 brown      light      brown           19   fema\textasciitilde{} femin\textasciitilde{}}
\DocumentationTok{\#\#  6 Owen La\textasciitilde{}    178   120 brown, gr\textasciitilde{} light      blue            52   male  mascu\textasciitilde{}}
\DocumentationTok{\#\#  7 Beru Wh\textasciitilde{}    165    75 brown      light      blue            47   fema\textasciitilde{} femin\textasciitilde{}}
\DocumentationTok{\#\#  8 R5{-}D4        97    32 \textless{}NA\textgreater{}       white, red red             NA   none  mascu\textasciitilde{}}
\DocumentationTok{\#\#  9 Biggs D\textasciitilde{}    183    84 black      light      brown           24   male  mascu\textasciitilde{}}
\DocumentationTok{\#\# 10 Obi{-}Wan\textasciitilde{}    182    77 auburn, w\textasciitilde{} fair       blue{-}gray       57   male  mascu\textasciitilde{}}
\DocumentationTok{\#\# \# i 77 more rows}
\DocumentationTok{\#\# \# i 6 more variables: homeworld \textless{}chr\textgreater{}, species \textless{}chr\textgreater{}, films \textless{}list\textgreater{},}
\DocumentationTok{\#\# \#   vehicles \textless{}list\textgreater{}, starships \textless{}list\textgreater{}, height\_meters \textless{}dbl\textgreater{}}
\end{Highlighting}
\end{Shaded}

Podemos combinar várias colunas numa expressão. Calculemos o índice de massa corporal
(IMC) de cada personagem, \(\text{IMC} = \text{massa} / \text{altura}^2\) (com a altura
em metros), devolvendo \texttt{NA} sempre que falte a altura ou a massa:

\begin{Shaded}
\begin{Highlighting}[]
\NormalTok{sw }\SpecialCharTok{\%\textgreater{}\%}
  \FunctionTok{mutate}\NormalTok{(}\AttributeTok{IMC =}\NormalTok{ mass }\SpecialCharTok{/}\NormalTok{ (height }\SpecialCharTok{/} \DecValTok{100}\NormalTok{)}\SpecialCharTok{\^{}}\DecValTok{2}\NormalTok{) }\SpecialCharTok{\%\textgreater{}\%}
  \FunctionTok{select}\NormalTok{(name, height, mass, IMC)}
\DocumentationTok{\#\# \# A tibble: 87 x 4}
\DocumentationTok{\#\#    name               height  mass   IMC}
\DocumentationTok{\#\#    \textless{}chr\textgreater{}               \textless{}int\textgreater{} \textless{}dbl\textgreater{} \textless{}dbl\textgreater{}}
\DocumentationTok{\#\#  1 Luke Skywalker        172    77  26.0}
\DocumentationTok{\#\#  2 C{-}3PO                 167    75  26.9}
\DocumentationTok{\#\#  3 R2{-}D2                  96    32  34.7}
\DocumentationTok{\#\#  4 Darth Vader           202   136  33.3}
\DocumentationTok{\#\#  5 Leia Organa           150    49  21.8}
\DocumentationTok{\#\#  6 Owen Lars             178   120  37.9}
\DocumentationTok{\#\#  7 Beru Whitesun Lars    165    75  27.5}
\DocumentationTok{\#\#  8 R5{-}D4                  97    32  34.0}
\DocumentationTok{\#\#  9 Biggs Darklighter     183    84  25.1}
\DocumentationTok{\#\# 10 Obi{-}Wan Kenobi        182    77  23.2}
\DocumentationTok{\#\# \# i 77 more rows}
\end{Highlighting}
\end{Shaded}

Não foi preciso tratar os \texttt{NA} explicitamente: qualquer operação aritmética que
envolva um \texttt{NA} devolve \texttt{NA}, pelo que as personagens sem altura ou sem massa ficam
automaticamente com \texttt{IMC} igual a \texttt{NA}. Esta propagação silenciosa dos valores em falta
é uma característica do R que convém ter sempre presente.

\subsubsection{Exercícios}\label{exercuxedcios-18}

\textbf{1.} Crie a coluna \texttt{dif\_altura\_peso} (diferença entre altura e massa) e guarde a
tabela em \texttt{starwars\_dif}. Depois, filtre as personagens cuja altura excede a massa e
ordene por \texttt{dif\_altura\_peso} crescente.

\textbf{2.} Crie as colunas: (a) \texttt{indice\_massa\_altura\ =\ mass\ /\ height};
(b) \texttt{indice\_massa\_medio\ =\ mean(mass,\ na.rm\ =\ TRUE)};
(c) \texttt{indice\_relativo\ =\ (indice\_massa\_altura\ -\ indice\_massa\_medio)\ /\ indice\_massa\_medio};
(d) \texttt{acima\_media\ =\ ifelse(indice\_massa\_altura\ \textgreater{}\ indice\_massa\_medio,\ "sim",\ "não")}.

\textbf{3.} Crie uma coluna que classifique cada personagem como \texttt{"recente"} (nascida menos
de 100 anos antes da batalha de Yavin) ou \texttt{"antiga"} (100 anos ou mais).

\subsection{\texorpdfstring{Resumir com \texttt{summarize()} e \texttt{group\_by()}}{Resumir com summarize() e group\_by()}}\label{resumir-com-summarize-e-group_by}

A \texttt{summarize()} reduz uma ou mais colunas a estatísticas --- uma média, uma mediana, uma
variância. Combinada com \texttt{group\_by()}, calcula essas estatísticas \textbf{para cada grupo},
o que é a essência da análise por categorias:

\begin{Shaded}
\begin{Highlighting}[]
\CommentTok{\# Uma estatística sobre toda a tabela}
\NormalTok{sw }\SpecialCharTok{\%\textgreater{}\%} \FunctionTok{summarize}\NormalTok{(}\AttributeTok{media\_massa =} \FunctionTok{mean}\NormalTok{(mass, }\AttributeTok{na.rm =} \ConstantTok{TRUE}\NormalTok{))}
\DocumentationTok{\#\# \# A tibble: 1 x 1}
\DocumentationTok{\#\#   media\_massa}
\DocumentationTok{\#\#         \textless{}dbl\textgreater{}}
\DocumentationTok{\#\# 1        97.3}

\CommentTok{\# A mesma estatística, agora por espécie}
\NormalTok{sw }\SpecialCharTok{\%\textgreater{}\%}
  \FunctionTok{group\_by}\NormalTok{(species) }\SpecialCharTok{\%\textgreater{}\%}
  \FunctionTok{summarize}\NormalTok{(}\AttributeTok{media\_massa =} \FunctionTok{mean}\NormalTok{(mass, }\AttributeTok{na.rm =} \ConstantTok{TRUE}\NormalTok{))}
\DocumentationTok{\#\# \# A tibble: 38 x 2}
\DocumentationTok{\#\#    species   media\_massa}
\DocumentationTok{\#\#    \textless{}chr\textgreater{}           \textless{}dbl\textgreater{}}
\DocumentationTok{\#\#  1 Aleena           15  }
\DocumentationTok{\#\#  2 Besalisk        102  }
\DocumentationTok{\#\#  3 Cerean           82  }
\DocumentationTok{\#\#  4 Chagrian        NaN  }
\DocumentationTok{\#\#  5 Clawdite         55  }
\DocumentationTok{\#\#  6 Droid            69.8}
\DocumentationTok{\#\#  7 Dug              40  }
\DocumentationTok{\#\#  8 Ewok             20  }
\DocumentationTok{\#\#  9 Geonosian        80  }
\DocumentationTok{\#\# 10 Gungan           74  }
\DocumentationTok{\#\# \# i 28 more rows}
\end{Highlighting}
\end{Shaded}

Podemos calcular várias estatísticas ao mesmo tempo:

\begin{Shaded}
\begin{Highlighting}[]
\NormalTok{sw }\SpecialCharTok{\%\textgreater{}\%} \FunctionTok{summarize}\NormalTok{(}
  \AttributeTok{media\_massa     =} \FunctionTok{mean}\NormalTok{(mass, }\AttributeTok{na.rm =} \ConstantTok{TRUE}\NormalTok{),}
  \AttributeTok{mediana\_massa   =} \FunctionTok{median}\NormalTok{(mass, }\AttributeTok{na.rm =} \ConstantTok{TRUE}\NormalTok{),}
  \AttributeTok{variancia\_massa =} \FunctionTok{var}\NormalTok{(mass, }\AttributeTok{na.rm =} \ConstantTok{TRUE}\NormalTok{)}
\NormalTok{)}
\DocumentationTok{\#\# \# A tibble: 1 x 3}
\DocumentationTok{\#\#   media\_massa mediana\_massa variancia\_massa}
\DocumentationTok{\#\#         \textless{}dbl\textgreater{}         \textless{}dbl\textgreater{}           \textless{}dbl\textgreater{}}
\DocumentationTok{\#\# 1        97.3            79          28716.}
\end{Highlighting}
\end{Shaded}

E resumir uma variável dentro de cada categoria de outra --- aqui, a altura média por cor
de cabelo:

\begin{Shaded}
\begin{Highlighting}[]
\NormalTok{sw }\SpecialCharTok{\%\textgreater{}\%}
  \FunctionTok{filter}\NormalTok{(}\SpecialCharTok{!}\FunctionTok{is.na}\NormalTok{(hair\_color), }\SpecialCharTok{!}\FunctionTok{is.na}\NormalTok{(height)) }\SpecialCharTok{\%\textgreater{}\%}
  \FunctionTok{group\_by}\NormalTok{(hair\_color) }\SpecialCharTok{\%\textgreater{}\%}
  \FunctionTok{summarize}\NormalTok{(}\AttributeTok{media\_altura =} \FunctionTok{mean}\NormalTok{(height))}
\DocumentationTok{\#\# \# A tibble: 11 x 2}
\DocumentationTok{\#\#    hair\_color    media\_altura}
\DocumentationTok{\#\#    \textless{}chr\textgreater{}                \textless{}dbl\textgreater{}}
\DocumentationTok{\#\#  1 auburn                150 }
\DocumentationTok{\#\#  2 auburn, grey          180 }
\DocumentationTok{\#\#  3 auburn, white         182 }
\DocumentationTok{\#\#  4 black                 174.}
\DocumentationTok{\#\#  5 blond                 177.}
\DocumentationTok{\#\#  6 blonde                168 }
\DocumentationTok{\#\#  7 brown                 177.}
\DocumentationTok{\#\#  8 brown, grey           178 }
\DocumentationTok{\#\#  9 grey                  170 }
\DocumentationTok{\#\# 10 none                  181.}
\DocumentationTok{\#\# 11 white                 156}
\end{Highlighting}
\end{Shaded}

Repare no papel de \texttt{na.rm\ =\ TRUE}. Funções como \texttt{mean()} ou \texttt{var()} devolvem \texttt{NA} se
houver um único valor em falta nos dados; \texttt{na.rm\ =\ TRUE} manda ignorá-los no cálculo.
Ignorar valores em falta nem sempre é inócuo --- pode enviesar resultados ---, mas é
frequentemente o que se pretende ao resumir dados reais, que quase nunca estão
completos.

\subsubsection{Exercícios}\label{exercuxedcios-19}

Use a base \texttt{sw}.

\textbf{1.} Calcule a altura média e mediana das personagens.

\textbf{2.} Calcule a massa média das personagens com mais de 175 cm.

\textbf{3.} Na mesma tabela, apresente a massa média das personagens com menos de 175 cm e
a das personagens com 175 cm ou mais.

\textbf{4.} Devolva tabelas (\emph{tibbles}) com: (a) a altura média por espécie; (b) a massa
média e mediana por espécie; (c) apenas o nome das personagens que participaram em mais
de 2 filmes.

\section{Juntar tabelas}\label{juntar-tabelas}

Em ciência de dados, a informação está muitas vezes repartida por várias tabelas que
partilham uma \textbf{coluna-chave} --- um identificador comum. As funções de junção do
\texttt{dplyr} combinam-nas, e distinguem-se pela forma como tratam as linhas sem
correspondência:

\begin{itemize}
\tightlist
\item
  \texttt{left\_join()} --- mantém todas as linhas da tabela da esquerda;
\item
  \texttt{right\_join()} --- mantém todas as linhas da tabela da direita;
\item
  \texttt{full\_join()} --- mantém todas as linhas de ambas.
\end{itemize}

Onde não há correspondência, as colunas em falta são preenchidas com \texttt{NA}.

Para ilustrar, criamos dois subconjuntos de \texttt{sw} --- as personagens altas e as
personagens humanas --- ligados pela chave \texttt{name}:

\begin{Shaded}
\begin{Highlighting}[]
\NormalTok{personagens\_altos }\OtherTok{\textless{}{-}}\NormalTok{ sw }\SpecialCharTok{\%\textgreater{}\%}
  \FunctionTok{filter}\NormalTok{(height }\SpecialCharTok{\textgreater{}} \DecValTok{180}\NormalTok{) }\SpecialCharTok{\%\textgreater{}\%}
  \FunctionTok{select}\NormalTok{(name, height)}

\NormalTok{personagens\_humanos }\OtherTok{\textless{}{-}}\NormalTok{ sw }\SpecialCharTok{\%\textgreater{}\%}
  \FunctionTok{filter}\NormalTok{(species }\SpecialCharTok{==} \StringTok{"Human"}\NormalTok{) }\SpecialCharTok{\%\textgreater{}\%}
  \FunctionTok{select}\NormalTok{(name, species)}
\end{Highlighting}
\end{Shaded}

A \texttt{left\_join()} mantém \textbf{todas} as personagens altas e acrescenta a espécie apenas às
que também são humanas (as restantes ficam com \texttt{NA} em \texttt{species}):

\begin{Shaded}
\begin{Highlighting}[]
\FunctionTok{left\_join}\NormalTok{(personagens\_altos, personagens\_humanos, }\AttributeTok{by =} \StringTok{"name"}\NormalTok{)}
\DocumentationTok{\#\# \# A tibble: 39 x 3}
\DocumentationTok{\#\#    name              height species}
\DocumentationTok{\#\#    \textless{}chr\textgreater{}              \textless{}int\textgreater{} \textless{}chr\textgreater{}  }
\DocumentationTok{\#\#  1 Darth Vader          202 Human  }
\DocumentationTok{\#\#  2 Biggs Darklighter    183 Human  }
\DocumentationTok{\#\#  3 Obi{-}Wan Kenobi       182 Human  }
\DocumentationTok{\#\#  4 Anakin Skywalker     188 Human  }
\DocumentationTok{\#\#  5 Chewbacca            228 \textless{}NA\textgreater{}   }
\DocumentationTok{\#\#  6 Boba Fett            183 Human  }
\DocumentationTok{\#\#  7 IG{-}88                200 \textless{}NA\textgreater{}   }
\DocumentationTok{\#\#  8 Bossk                190 \textless{}NA\textgreater{}   }
\DocumentationTok{\#\#  9 Qui{-}Gon Jinn         193 Human  }
\DocumentationTok{\#\# 10 Nute Gunray          191 \textless{}NA\textgreater{}   }
\DocumentationTok{\#\# \# i 29 more rows}
\end{Highlighting}
\end{Shaded}

A \texttt{right\_join()} faz o inverso: mantém todas as personagens humanas e acrescenta a
altura às que também são altas:

\begin{Shaded}
\begin{Highlighting}[]
\FunctionTok{right\_join}\NormalTok{(personagens\_altos, personagens\_humanos, }\AttributeTok{by =} \StringTok{"name"}\NormalTok{)}
\DocumentationTok{\#\# \# A tibble: 35 x 3}
\DocumentationTok{\#\#    name              height species}
\DocumentationTok{\#\#    \textless{}chr\textgreater{}              \textless{}int\textgreater{} \textless{}chr\textgreater{}  }
\DocumentationTok{\#\#  1 Darth Vader          202 Human  }
\DocumentationTok{\#\#  2 Biggs Darklighter    183 Human  }
\DocumentationTok{\#\#  3 Obi{-}Wan Kenobi       182 Human  }
\DocumentationTok{\#\#  4 Anakin Skywalker     188 Human  }
\DocumentationTok{\#\#  5 Boba Fett            183 Human  }
\DocumentationTok{\#\#  6 Qui{-}Gon Jinn         193 Human  }
\DocumentationTok{\#\#  7 Padmé Amidala        185 Human  }
\DocumentationTok{\#\#  8 Ric Olié             183 Human  }
\DocumentationTok{\#\#  9 Quarsh Panaka        183 Human  }
\DocumentationTok{\#\# 10 Mace Windu           188 Human  }
\DocumentationTok{\#\# \# i 25 more rows}
\end{Highlighting}
\end{Shaded}

A \texttt{full\_join()} mantém todas as personagens de ambas as tabelas:

\begin{Shaded}
\begin{Highlighting}[]
\FunctionTok{full\_join}\NormalTok{(personagens\_altos, personagens\_humanos, }\AttributeTok{by =} \StringTok{"name"}\NormalTok{)}
\DocumentationTok{\#\# \# A tibble: 59 x 3}
\DocumentationTok{\#\#    name              height species}
\DocumentationTok{\#\#    \textless{}chr\textgreater{}              \textless{}int\textgreater{} \textless{}chr\textgreater{}  }
\DocumentationTok{\#\#  1 Darth Vader          202 Human  }
\DocumentationTok{\#\#  2 Biggs Darklighter    183 Human  }
\DocumentationTok{\#\#  3 Obi{-}Wan Kenobi       182 Human  }
\DocumentationTok{\#\#  4 Anakin Skywalker     188 Human  }
\DocumentationTok{\#\#  5 Chewbacca            228 \textless{}NA\textgreater{}   }
\DocumentationTok{\#\#  6 Boba Fett            183 Human  }
\DocumentationTok{\#\#  7 IG{-}88                200 \textless{}NA\textgreater{}   }
\DocumentationTok{\#\#  8 Bossk                190 \textless{}NA\textgreater{}   }
\DocumentationTok{\#\#  9 Qui{-}Gon Jinn         193 Human  }
\DocumentationTok{\#\# 10 Nute Gunray          191 \textless{}NA\textgreater{}   }
\DocumentationTok{\#\# \# i 49 more rows}
\end{Highlighting}
\end{Shaded}

As três junções são a versão em \texttt{dplyr} das combinações de tabelas que já tínhamos
visto com a função \texttt{merge()} do R base (Capítulo~\ref{estruturas}). A vantagem aqui é
a clareza: o nome de cada função diz exatamente que linhas serão preservadas.

\subsubsection{Exercícios}\label{exercuxedcios-20}

\textbf{1.} Crie \texttt{personagens\_altos} (altura \textgreater{} 180 cm) e \texttt{personagens\_leves} (massa \textless{} 75 kg)
e combine-as com \texttt{left\_join()} pela coluna \texttt{name}. Observe que as personagens altas sem
correspondência ficam com \texttt{NA} na massa.

\textbf{2.} Crie \texttt{personagens\_humanos} (espécie \texttt{"Human"}) e \texttt{cor\_cabelo} (personagens com
\texttt{hair\_color} não \texttt{NA}) e combine-as com \texttt{right\_join()} por \texttt{name}.

\textbf{3.} Crie \texttt{especies\_personagens} (espécie não \texttt{NA}) e \texttt{cor\_olhos} (\texttt{eye\_color} não
\texttt{NA}) e combine-as com \texttt{full\_join()} por \texttt{name}.

Com os dados lidos, limpos e organizados, estamos prontos para o passo mais
recompensador: representá-los graficamente com toda a expressividade. É o que faremos no
próximo capítulo, com o \texttt{ggplot2}.

\chapter{\texorpdfstring{O pacote \texttt{ggplot2}}{O pacote ggplot2}}\label{ggplot}

No Capítulo~\ref{graficos} aprendemos a desenhar gráficos com o R base, acrescentando
elementos ao gráfico um comando de cada vez. Neste capítulo conhecemos uma abordagem
diferente e muito influente: a do \texttt{ggplot2}, que não nos pede que desenhemos o gráfico,
mas que o \textbf{descrevamos}.

\section{Um pouco de história}\label{um-pouco-de-histuxf3ria}

A ideia que está por trás do \texttt{ggplot2} nasceu num livro: \emph{The Grammar of Graphics}, do
estatístico e informático Leland Wilkinson, publicado em 1999 (com segunda edição em
2005). Wilkinson propôs que qualquer gráfico estatístico, por mais complexo, pode ser
descrito como a combinação de um pequeno número de componentes independentes --- os
dados, os mapeamentos para propriedades visuais, as formas geométricas, as escalas, o
sistema de coordenadas. Tal como a gramática de uma língua permite formar infinitas
frases a partir de poucas regras, esta ``gramática dos gráficos'' permite construir uma
enorme variedade de visualizações a partir de poucos elementos.

Em meados da década de 2000, Hadley Wickham implementou e adaptou estas ideias num
pacote de R, a que chamou \texttt{ggplot2} --- o ``gg'' vem precisamente de \emph{grammar of graphics}.
A sua adaptação, publicada em 2010 no artigo \emph{A Layered Grammar of Graphics}, organiza a
construção de um gráfico em \textbf{camadas} que se somam umas às outras. O \texttt{ggplot2}
tornou-se, desde então, um dos pacotes mais usados do R e uma das pedras angulares do
\texttt{tidyverse}.

\section{Porque é importante}\label{porque-uxe9-importante}

O que distingue o \texttt{ggplot2} da abordagem do R base é a mudança de perspetiva. Com o R
base, pensamos em \emph{instruções}: ``desenha estes pontos, agora acrescenta uma legenda,
agora uma linha''. Com o \texttt{ggplot2}, pensamos em \emph{relações}: ``a variável \texttt{wt} mapeia-se no
eixo horizontal, a \texttt{mpg} no eixo vertical, e cada observação é um ponto''. Descrevemos a
correspondência entre os dados e as propriedades visuais, e o pacote encarrega-se do
resto --- escalas, eixos, legendas.

Desta filosofia decorrem as suas grandes vantagens: a \textbf{consistência} (a mesma
gramática serve para todos os tipos de gráfico), a facilidade de construir gráficos
sofisticados por acumulação de camadas, resultados com \textbf{qualidade de publicação} logo
por omissão, e um vasto ecossistema de extensões. É, por tudo isto, a ferramenta de
eleição para a visualização de dados em R.

Este capítulo usa o \texttt{ggplot2}, incluído no \texttt{tidyverse}. Instale-o com
\texttt{install.packages("ggplot2")} (ou \texttt{install.packages("tidyverse")}) e carregue-o com
\texttt{library(ggplot2)}.

\section{A gramática, camada a camada}\label{a-gramuxe1tica-camada-a-camada}

Um gráfico em \texttt{ggplot2} compõe-se dos seguintes elementos:

\begin{enumerate}
\def\labelenumi{\arabic{enumi}.}
\tightlist
\item
  \textbf{Dados} (\emph{data}) --- o conjunto de dados a visualizar (um \emph{data frame} ou \emph{tibble}).
\item
  \textbf{Estética} (\emph{aes}) --- os mapeamentos das variáveis para propriedades visuais: a
  posição (\texttt{x}, \texttt{y}), a cor (\texttt{colour}), o tamanho (\texttt{size}), a forma (\texttt{shape}), o
  preenchimento (\texttt{fill}).
\item
  \textbf{Geometrias} (\emph{geoms}) --- o tipo de objeto que representa os dados: pontos
  (\texttt{geom\_point()}), barras (\texttt{geom\_bar()}), linhas (\texttt{geom\_line()}), caixas
  (\texttt{geom\_boxplot()}), etc.
\item
  \textbf{Escalas} (\emph{scales}) --- como os valores dos dados se traduzem em propriedades
  visuais (que cores, que intervalos nos eixos).
\item
  \textbf{Facetas} (\emph{facets}) --- a divisão do gráfico em subgráficos, um por grupo.
\item
  \textbf{Temas} (\emph{themes}) --- a aparência geral (tipos de letra, fundos, grelhas).
\end{enumerate}

As camadas somam-se literalmente com o operador \texttt{+}. A estrutura mínima é sempre a
mesma: \texttt{ggplot()} recebe os dados e a estética, e uma geometria acrescenta os objetos
visuais.

\begin{Shaded}
\begin{Highlighting}[]
\FunctionTok{library}\NormalTok{(ggplot2)}

\FunctionTok{ggplot}\NormalTok{(}\AttributeTok{data =}\NormalTok{ mtcars, }\FunctionTok{aes}\NormalTok{(}\AttributeTok{x =}\NormalTok{ wt, }\AttributeTok{y =}\NormalTok{ mpg)) }\SpecialCharTok{+}
  \FunctionTok{geom\_point}\NormalTok{()}
\end{Highlighting}
\end{Shaded}

\begin{figure}

{\centering \includegraphics[width=0.8\linewidth]{introR_files/figure-latex/gg-base-1} 

}

\caption{A estrutura mínima: dados, estética e uma geometria.}\label{fig:gg-base}
\end{figure}

Aqui, \texttt{mtcars} são os dados, \texttt{aes(x\ =\ wt,\ y\ =\ mpg)} diz que o peso vai no eixo
horizontal e o consumo no vertical, e \texttt{geom\_point()} representa cada carro por um ponto.

\section{Exemplos}\label{exemplos-2}

\textbf{Dispersão com uma terceira variável.} A força da estética revela-se quando mapeamos
\emph{mais} variáveis. Ao colorir os pontos pelo número de cilindros, acrescentamos uma
terceira dimensão de informação sem mudar de gráfico. Como \texttt{cyl} é, na verdade, uma
variável categórica, convertemo-la num fator:

\begin{Shaded}
\begin{Highlighting}[]
\FunctionTok{ggplot}\NormalTok{(mtcars, }\FunctionTok{aes}\NormalTok{(}\AttributeTok{x =}\NormalTok{ wt, }\AttributeTok{y =}\NormalTok{ mpg, }\AttributeTok{colour =} \FunctionTok{factor}\NormalTok{(cyl))) }\SpecialCharTok{+}
  \FunctionTok{geom\_point}\NormalTok{(}\AttributeTok{size =} \DecValTok{2}\NormalTok{) }\SpecialCharTok{+}
  \FunctionTok{labs}\NormalTok{(}\AttributeTok{title =} \StringTok{"Consumo em função do peso"}\NormalTok{,}
       \AttributeTok{x =} \StringTok{"Peso (1000 lb)"}\NormalTok{, }\AttributeTok{y =} \StringTok{"Consumo (milhas/galão)"}\NormalTok{,}
       \AttributeTok{colour =} \StringTok{"Cilindros"}\NormalTok{)}
\end{Highlighting}
\end{Shaded}

\begin{figure}

{\centering \includegraphics[width=0.8\linewidth]{introR_files/figure-latex/gg-disp-cor-1} 

}

\caption{Consumo vs peso, com a cor a codificar o número de cilindros.}\label{fig:gg-disp-cor}
\end{figure}

\textbf{Gráfico de barras.} Com uma variável categórica, \texttt{geom\_bar()} conta automaticamente
as observações de cada categoria:

\begin{Shaded}
\begin{Highlighting}[]
\FunctionTok{ggplot}\NormalTok{(mpg, }\FunctionTok{aes}\NormalTok{(}\AttributeTok{x =}\NormalTok{ class)) }\SpecialCharTok{+}
  \FunctionTok{geom\_bar}\NormalTok{(}\AttributeTok{fill =} \StringTok{"steelblue"}\NormalTok{) }\SpecialCharTok{+}
  \FunctionTok{labs}\NormalTok{(}\AttributeTok{title =} \StringTok{"Modelos por classe"}\NormalTok{, }\AttributeTok{x =} \StringTok{"Classe"}\NormalTok{, }\AttributeTok{y =} \StringTok{"Contagem"}\NormalTok{)}
\end{Highlighting}
\end{Shaded}

\begin{figure}

{\centering \includegraphics[width=0.8\linewidth]{introR_files/figure-latex/gg-barras-1} 

}

\caption{Número de modelos por classe de automóvel (conjunto mpg).}\label{fig:gg-barras}
\end{figure}

\textbf{Barras proporcionais.} Mapeando \texttt{fill} para uma segunda variável categórica e usando
\texttt{position\ =\ "fill"}, cada barra passa a mostrar \textbf{proporções} --- ideal para comparar a
composição entre grupos. Usamos o conjunto \texttt{diamonds}, que acompanha o \texttt{ggplot2}:

\begin{Shaded}
\begin{Highlighting}[]
\FunctionTok{ggplot}\NormalTok{(diamonds, }\FunctionTok{aes}\NormalTok{(}\AttributeTok{x =}\NormalTok{ cut, }\AttributeTok{fill =}\NormalTok{ clarity)) }\SpecialCharTok{+}
  \FunctionTok{geom\_bar}\NormalTok{(}\AttributeTok{position =} \StringTok{"fill"}\NormalTok{) }\SpecialCharTok{+}
  \FunctionTok{labs}\NormalTok{(}\AttributeTok{title =} \StringTok{"Pureza por qualidade de corte"}\NormalTok{,}
       \AttributeTok{x =} \StringTok{"Corte"}\NormalTok{, }\AttributeTok{y =} \StringTok{"Proporção"}\NormalTok{, }\AttributeTok{fill =} \StringTok{"Pureza"}\NormalTok{)}
\end{Highlighting}
\end{Shaded}

\begin{figure}

{\centering \includegraphics[width=0.8\linewidth]{introR_files/figure-latex/gg-barras-prop-1} 

}

\caption{Composição da pureza (clarity) em cada qualidade de corte (cut).}\label{fig:gg-barras-prop}
\end{figure}

\textbf{Caixas de bigodes por grupo.} A \texttt{geom\_boxplot()} compara a distribuição de uma
variável quantitativa entre categorias:

\begin{Shaded}
\begin{Highlighting}[]
\FunctionTok{ggplot}\NormalTok{(mpg, }\FunctionTok{aes}\NormalTok{(}\AttributeTok{x =}\NormalTok{ class, }\AttributeTok{y =}\NormalTok{ hwy)) }\SpecialCharTok{+}
  \FunctionTok{geom\_boxplot}\NormalTok{(}\AttributeTok{fill =} \StringTok{"lightblue"}\NormalTok{) }\SpecialCharTok{+}
  \FunctionTok{labs}\NormalTok{(}\AttributeTok{title =} \StringTok{"Consumo em autoestrada por classe"}\NormalTok{,}
       \AttributeTok{x =} \StringTok{"Classe"}\NormalTok{, }\AttributeTok{y =} \StringTok{"Milhas por galão (autoestrada)"}\NormalTok{)}
\end{Highlighting}
\end{Shaded}

\begin{figure}

{\centering \includegraphics[width=0.8\linewidth]{introR_files/figure-latex/gg-box-1} 

}

\caption{Consumo em autoestrada por classe de automóvel.}\label{fig:gg-box}
\end{figure}

\textbf{Tendência com regressão.} Somando uma camada \texttt{geom\_smooth()}, o \texttt{ggplot2} ajusta e
desenha uma tendência --- com \texttt{method\ =\ "lm"}, a reta de regressão dos mínimos quadrados,
rodeada da sua banda de confiança:

\begin{Shaded}
\begin{Highlighting}[]
\FunctionTok{ggplot}\NormalTok{(mtcars, }\FunctionTok{aes}\NormalTok{(}\AttributeTok{x =}\NormalTok{ wt, }\AttributeTok{y =}\NormalTok{ mpg)) }\SpecialCharTok{+}
  \FunctionTok{geom\_point}\NormalTok{() }\SpecialCharTok{+}
  \FunctionTok{geom\_smooth}\NormalTok{(}\AttributeTok{method =} \StringTok{"lm"}\NormalTok{) }\SpecialCharTok{+}
  \FunctionTok{labs}\NormalTok{(}\AttributeTok{title =} \StringTok{"Consumo vs peso, com reta de regressão"}\NormalTok{,}
       \AttributeTok{x =} \StringTok{"Peso (1000 lb)"}\NormalTok{, }\AttributeTok{y =} \StringTok{"Consumo (milhas/galão)"}\NormalTok{)}
\end{Highlighting}
\end{Shaded}

\begin{figure}

{\centering \includegraphics[width=0.8\linewidth]{introR_files/figure-latex/gg-smooth-1} 

}

\caption{Dispersão com reta de regressão e banda de confiança.}\label{fig:gg-smooth}
\end{figure}

\textbf{Facetas.} As facetas dividem o gráfico numa grelha de subgráficos, um por grupo,
partilhando os eixos --- o que torna as comparações imediatas. Aqui, a relação entre
cilindrada e consumo, separada por classe:

\begin{Shaded}
\begin{Highlighting}[]
\FunctionTok{ggplot}\NormalTok{(mpg, }\FunctionTok{aes}\NormalTok{(}\AttributeTok{x =}\NormalTok{ displ, }\AttributeTok{y =}\NormalTok{ hwy)) }\SpecialCharTok{+}
  \FunctionTok{geom\_point}\NormalTok{() }\SpecialCharTok{+}
  \FunctionTok{facet\_wrap}\NormalTok{(}\SpecialCharTok{\textasciitilde{}}\NormalTok{ class) }\SpecialCharTok{+}
  \FunctionTok{labs}\NormalTok{(}\AttributeTok{x =} \StringTok{"Cilindrada (l)"}\NormalTok{, }\AttributeTok{y =} \StringTok{"Consumo (autoestrada)"}\NormalTok{)}
\end{Highlighting}
\end{Shaded}

\begin{figure}

{\centering \includegraphics[width=0.8\linewidth]{introR_files/figure-latex/gg-facet-1} 

}

\caption{Subgráfico por classe, usando a função facet\_wrap().}\label{fig:gg-facet}
\end{figure}

\textbf{Temas.} Um tema controla toda a aparência do gráfico. Trocar \texttt{theme\_minimal()} por
outro tema muda o aspeto sem tocar nos dados nem na estética:

\begin{Shaded}
\begin{Highlighting}[]
\FunctionTok{ggplot}\NormalTok{(mpg, }\FunctionTok{aes}\NormalTok{(}\AttributeTok{x =}\NormalTok{ displ, }\AttributeTok{y =}\NormalTok{ hwy)) }\SpecialCharTok{+}
  \FunctionTok{geom\_point}\NormalTok{(}\AttributeTok{colour =} \StringTok{"steelblue"}\NormalTok{) }\SpecialCharTok{+}
  \FunctionTok{theme\_minimal}\NormalTok{() }\SpecialCharTok{+}
  \FunctionTok{labs}\NormalTok{(}\AttributeTok{x =} \StringTok{"Cilindrada (l)"}\NormalTok{, }\AttributeTok{y =} \StringTok{"Consumo (autoestrada)"}\NormalTok{)}
\end{Highlighting}
\end{Shaded}

\begin{figure}

{\centering \includegraphics[width=0.8\linewidth]{introR_files/figure-latex/gg-tema-1} 

}

\caption{O mesmo gráfico com o tema theme\_minimal().}\label{fig:gg-tema}
\end{figure}

Alguns hábitos tornam qualquer gráfico melhor: dar títulos e rótulos claros com
\texttt{labs(title\ =\ ...,\ x\ =\ ...,\ y\ =\ ...)}; evitar sobrecarregar o gráfico com informação a
mais; escolher cores com bom contraste (e legíveis por quem tem daltonismo, por exemplo
com \texttt{scale\_colour\_viridis\_d()}); e usar facetas em vez de amontoar tudo num só painel
quando há muitos grupos.

\section{Exercícios}\label{exercuxedcios-21}

Nos exercícios seguintes, use o \texttt{ggplot2} e os conjuntos de dados indicados.

\textbf{1.} Com o conjunto \texttt{mtcars}, faça um gráfico de dispersão entre \texttt{wt} e \texttt{mpg},
colorindo os pontos pela variável \texttt{cyl}.

\textbf{2.} Com o conjunto \texttt{mtcars}, faça um gráfico de dispersão entre \texttt{mpg} e \texttt{hp},
colorindo os pontos por \texttt{cyl} e variando a forma (\texttt{shape}) dos pontos por \texttt{am}.

\textbf{3.} Com o conjunto \texttt{mpg}, faça um gráfico de barras da variável \texttt{drv} e altere a cor
das barras com \texttt{fill}.

\textbf{4.} Com o conjunto \texttt{diamonds}, faça um gráfico de barras da variável \texttt{cut},
preenchendo (\texttt{fill}) pelas categorias de \texttt{clarity}. Use \texttt{position\ =\ "fill"} para
comparar as proporções relativas.

\textbf{5.} Com o conjunto \texttt{mpg}, desenhe um histograma da variável \texttt{hwy} e sobreponha-lhe
uma curva de densidade com \texttt{geom\_density()}. Explore o argumento \texttt{alpha} para controlar
a transparência.

\textbf{6.} Faça uma caixa de bigodes de \texttt{hwy} por \texttt{manufacturer} (conjunto \texttt{mpg}),
ordenando os fabricantes pela média de \texttt{hwy}.
\emph{Sugestão:} \texttt{forcats::fct\_reorder()}.

\textbf{7.} Com o conjunto \texttt{iris}, faça caixas de bigodes de \texttt{Sepal.Length} por espécie e
sobreponha um ponto com a média de cada grupo, usando \texttt{stat\_summary()}.

\textbf{8.} Com o conjunto \texttt{mpg}, ajuste uma curva \texttt{LOESS} a \texttt{displ} vs \texttt{hwy} com
\texttt{geom\_smooth()}, separando por \texttt{drv} com \texttt{facet\_grid()}.

\textbf{9.} Exporte um dos seus gráficos para um ficheiro \texttt{.png} com
\texttt{ggsave("grafico.png")}.

\textbf{10.} Escolha um conjunto de dados à sua escolha (ou um \emph{tibble} pequeno) e crie uma
visualização que combine pelo menos \textbf{dois} tipos de geometria (\texttt{geom}) diferentes
--- por exemplo, pontos e uma linha de tendência.

Dominadas a leitura, a manipulação e a visualização de dados, encerra-se a parte
aplicada de tratamento de dados. A partir daqui, o livro entra no seu tema central: a
\textbf{aleatoriedade} e a \textbf{simulação}. Começamos, no próximo capítulo, pela ferramenta
mais elementar para gerar acaso no R --- a função \texttt{sample()}.

\chapter{\texorpdfstring{A função \texttt{sample()}}{A função sample()}}\label{sample}

Começa aqui a parte central do livro: a \textbf{aleatoriedade} e a \textbf{simulação}. Simular é
imitar, no computador, um fenómeno sujeito ao acaso --- lançar um dado, extrair uma
amostra, esperar numa fila --- e observar o que acontece quando o repetimos muitas vezes.
É uma forma poderosa de estudar a probabilidade: em vez de resolver um problema apenas
com fórmulas, podemos \emph{fazê-lo acontecer} milhares de vezes e medir o resultado.

A ferramenta mais elementar para gerar acaso no R é a função \texttt{sample()}, que extrai
elementos ao acaso de um conjunto. Com ela lançam-se dados e moedas, baralham-se
cartas, sorteiam-se amostras --- os alicerces de quase todas as simulações que se seguem.

\begin{Shaded}
\begin{Highlighting}[]
\FunctionTok{sample}\NormalTok{(x, size, }\AttributeTok{replace =} \ConstantTok{FALSE}\NormalTok{, }\AttributeTok{prob =} \ConstantTok{NULL}\NormalTok{)}
\end{Highlighting}
\end{Shaded}

\begin{itemize}
\tightlist
\item
  \texttt{x} --- o vetor de onde se extraem os elementos;
\item
  \texttt{size} --- a dimensão da amostra a extrair;
\item
  \texttt{replace} --- se a extração é \textbf{com} reposição (\texttt{TRUE}) ou \textbf{sem} reposição
  (\texttt{FALSE}, por omissão);
\item
  \texttt{prob} --- um vetor de probabilidades associadas a cada elemento de \texttt{x}; se omitido,
  todos são igualmente prováveis.
\end{itemize}

Neste capítulo aprendemos a extrair amostras aleatórias com \texttt{sample()} --- com e sem
reposição, com probabilidades iguais ou diferentes --- e a tornar as simulações
reprodutíveis com \texttt{set.seed()}.

\section{Extrair amostras}\label{extrair-amostras}

\textbf{Sem reposição.} Cada elemento sai no máximo uma vez, como ao retirar bolas de um saco
sem as repor. Extraiamos 5 números de uma população de 1 a 10:

\begin{Shaded}
\begin{Highlighting}[]
\NormalTok{pop }\OtherTok{\textless{}{-}} \DecValTok{1}\SpecialCharTok{:}\DecValTok{10}
\FunctionTok{sample}\NormalTok{(pop, }\AttributeTok{size =} \DecValTok{5}\NormalTok{, }\AttributeTok{replace =} \ConstantTok{FALSE}\NormalTok{)}
\DocumentationTok{\#\# [1]  5  2  4 10  8      \# um resultado possível}
\end{Highlighting}
\end{Shaded}

\textbf{Com reposição.} Cada elemento pode sair mais do que uma vez, como ao lançar
repetidamente o mesmo dado. É \texttt{replace\ =\ TRUE} que permite repetições:

\begin{Shaded}
\begin{Highlighting}[]
\FunctionTok{sample}\NormalTok{(pop, }\AttributeTok{size =} \DecValTok{5}\NormalTok{, }\AttributeTok{replace =} \ConstantTok{TRUE}\NormalTok{)}
\DocumentationTok{\#\# [1] 6 8 9 9 6           \# o 6 e o 9 repetem{-}se}
\end{Highlighting}
\end{Shaded}

Sem reposição, a dimensão da amostra não pode exceder o número de elementos disponíveis:
\texttt{sample(1:10,\ size\ =\ 15)} gera um erro, porque não há 15 valores distintos para
extrair. Com \texttt{replace\ =\ TRUE}, pelo contrário, podemos extrair amostras tão grandes
quanto quisermos.

\textbf{Com probabilidades diferentes.} O argumento \texttt{prob} permite que uns elementos sejam
mais prováveis do que outros --- indispensável para simular dados ou moedas
enviesados. Aqui, os dois últimos valores têm metade da probabilidade dos restantes:

\begin{Shaded}
\begin{Highlighting}[]
\NormalTok{prob }\OtherTok{\textless{}{-}} \FunctionTok{c}\NormalTok{(}\FloatTok{0.1}\NormalTok{, }\FloatTok{0.1}\NormalTok{, }\FloatTok{0.1}\NormalTok{, }\FloatTok{0.1}\NormalTok{, }\FloatTok{0.1}\NormalTok{, }\FloatTok{0.1}\NormalTok{, }\FloatTok{0.1}\NormalTok{, }\FloatTok{0.1}\NormalTok{, }\FloatTok{0.05}\NormalTok{, }\FloatTok{0.05}\NormalTok{)}
\FunctionTok{sample}\NormalTok{(pop, }\AttributeTok{size =} \DecValTok{5}\NormalTok{, }\AttributeTok{prob =}\NormalTok{ prob)}
\DocumentationTok{\#\# [1] 2 4 6 5 3           \# um resultado possível}
\end{Highlighting}
\end{Shaded}

As probabilidades em \texttt{prob} não têm de somar 1: o R normaliza-as automaticamente,
dividindo cada uma pela soma total. Assim, \texttt{prob\ =\ c(2,\ 1,\ 1)} é equivalente a
\texttt{prob\ =\ c(0.5,\ 0.25,\ 0.25)} --- o que conta são as proporções relativas entre os
elementos.

\section{\texorpdfstring{Reprodutibilidade: \texttt{set.seed()}}{Reprodutibilidade: set.seed()}}\label{reprodutibilidade-set.seed}

Há uma subtileza fundamental em tudo o que envolve o acaso no computador. Os números
que o R diz serem ``aleatórios'' são, na verdade, \textbf{pseudoaleatórios}: resultam de um
algoritmo determinístico que, partindo de um valor inicial (a \emph{semente}), produz uma
sequência que \emph{parece} aleatória. Isto tem uma consequência muito útil: se fixarmos a
semente, obtemos sempre a mesma sequência.

A função \texttt{set.seed()} fixa essa semente. Executar \texttt{set.seed()} com o mesmo número antes
de uma simulação garante que ela produz exatamente o mesmo resultado de cada vez --- em
qualquer computador:

\begin{Shaded}
\begin{Highlighting}[]
\FunctionTok{set.seed}\NormalTok{(}\DecValTok{42}\NormalTok{)}
\FunctionTok{sample}\NormalTok{(}\DecValTok{1}\SpecialCharTok{:}\DecValTok{10}\NormalTok{, }\AttributeTok{size =} \DecValTok{5}\NormalTok{)}
\DocumentationTok{\#\# [1] 1 5 10 8 2}

\FunctionTok{set.seed}\NormalTok{(}\DecValTok{42}\NormalTok{)     }\CommentTok{\# a mesma semente...}
\FunctionTok{sample}\NormalTok{(}\DecValTok{1}\SpecialCharTok{:}\DecValTok{10}\NormalTok{, }\AttributeTok{size =} \DecValTok{5}\NormalTok{)}
\DocumentationTok{\#\# [1] 1 5 10 8 2   \# ...produz exatamente o mesmo resultado}
\end{Highlighting}
\end{Shaded}

Adquira o hábito de começar qualquer simulação com \texttt{set.seed()}. Não torna os resultados
menos aleatórios --- a sequência continua a comportar-se como se fosse acaso ---, mas torna
o seu trabalho \textbf{reprodutível}: quem executar o seu código obterá os mesmos números,
poderá confirmar as suas conclusões e depurar erros. A reprodutibilidade é uma exigência
central da ciência, e no computador consegue-se com uma linha.

Cuidado com uma armadilha de \texttt{sample()}: quando o primeiro argumento é um \textbf{único
número} positivo \texttt{n}, o R interpreta-o como \texttt{1:n}. Assim, \texttt{sample(10)} devolve uma
permutação de 1 a 10, e não o número 10. Isto raramente é o que se quer quando \texttt{x} é
calculado dentro de um programa e pode, por acaso, ter comprimento 1. Para amostrar com
segurança dos \emph{elementos} de um vetor \texttt{x}, use \texttt{x{[}sample(length(x),\ size){]}}, que indexa
sempre pelas posições.

\section{Exercícios}\label{exercuxedcios-22}

\textbf{1.} Crie um vetor com os números de 1 a 20 e extraia uma amostra aleatória de 5
elementos, sem reposição.

\textbf{2.} Considere uma população representada pelos números de 1 a 10. Extraia uma
amostra de 10 elementos com reposição.

\textbf{3.} Crie um vetor com as letras A, B, C, D, E e extraia uma amostra de 3 letras, em
que a probabilidade de cada uma é dada por \texttt{c(0.1,\ 0.2,\ 0.3,\ 0.25,\ 0.15)}.

\textbf{4.} Crie um vetor com os números de 1 a 10 e use \texttt{sample()} para o reordenar
aleatoriamente (uma permutação).

\textbf{5.} Crie um vetor com cinco frutas (``Maçã'', ``Banana'', ``Laranja'', ``Uva'', ``Pera'').
Selecione aleatoriamente uma fruta e, depois, uma amostra de 3 frutas.

\textbf{6.} É responsável pelo controlo de qualidade numa fábrica com 1000 produtos,
numerados de 1 a 1000. Extraia uma amostra aleatória de 50 produtos para inspeção, sem
reposição.

\textbf{7.} Simule o lançamento de dois dados equilibrados 10 000 vezes e registe a soma das
faces. Use \texttt{sample()} e, no fim, faça um histograma das somas obtidas.

\textbf{8.} Tem 200 estudantes distribuídos por três turmas --- A, B e C --- com 50, 100 e 50
alunos, respetivamente. Usando \texttt{sample()}, extraia uma amostra de 20 alunos que mantenha
a proporção das turmas.

\textbf{9.} Um cartão de Bingo contém 24 números aleatórios entre 1 e 75 (excluindo a casa
central ``livre''). Crie 5 cartões de Bingo distintos com \texttt{sample()}.

\textbf{10.} Num ensaio clínico, 30 pacientes têm de ser aleatorizados em dois grupos:
tratamento (20 pacientes) e controlo (10). Faça a aleatorização com \texttt{sample()}.
\emph{Sugestão:} selecione os 20 do tratamento e obtenha os restantes com \texttt{setdiff()}.

Com a \texttt{sample()} sabemos gerar acaso. Falta-nos a linguagem para o descrever
rigorosamente: as \textbf{variáveis aleatórias} e a \textbf{probabilidade}, tema do próximo
capítulo.

\chapter{Probabilidade e variáveis aleatórias}\label{probabilidade}

No capítulo anterior aprendemos a \emph{gerar} acaso. Agora precisamos da linguagem
matemática para o \emph{descrever} com rigor --- a linguagem das \textbf{variáveis aleatórias}, que
está no coração de toda a estatística e de toda a simulação.

Uma \textbf{variável aleatória} é uma função que atribui um valor numérico a cada resultado
possível de uma experiência aleatória. Formalmente, é uma função que faz corresponder a
cada elemento do espaço de resultados \(\Omega\) um número real. Ao lançar um dado, por
exemplo, podemos definir a variável aleatória \(X\) como o número da face voltada para
cima, que assume valores de 1 a 6.

As variáveis aleatórias são fundamentais porque permitem quantificar e analisar
matematicamente fenómenos incertos. Classificam-se em dois grandes tipos:

\begin{itemize}
\tightlist
\item
  \textbf{discretas} --- assumem valores num conjunto finito ou numerável, como o número de
  caras em três lançamentos de uma moeda;
\item
  \textbf{contínuas} --- assumem qualquer valor num intervalo, como a altura de uma pessoa.
\end{itemize}

Neste capítulo revemos as ferramentas que descrevem uma variável aleatória --- função de
probabilidade, função densidade, função de distribuição, valor esperado e variância --- e
aprendemos a calculá-las no R, com particular destaque para a integração numérica com
\texttt{integrate()}.

\section{Como se descreve uma variável aleatória}\label{como-se-descreve-uma-variuxe1vel-aleatuxf3ria}

A distribuição de uma variável aleatória descreve-se de formas diferentes consoante ela
seja discreta ou contínua.

Uma variável \textbf{discreta} descreve-se pela sua \textbf{função massa de probabilidade}
(f.m.p.), \(f_X(x) = P(X = x)\), que dá a probabilidade de cada valor. Para ser uma
f.m.p. válida, tem de satisfazer duas condições:
\[f_X(x) \geq 0 \quad \forall x, \qquad \sum_{x} f_X(x) = 1.\]

Uma variável \textbf{contínua} descreve-se pela sua \textbf{função densidade de probabilidade}
(f.d.p.), \(f_X(x)\). Aqui, a probabilidade não está em cada ponto (que é nula), mas na
\textbf{área sob a curva}: \(P(a < X < b) = \int_a^b f_X(x)\,dx\). As condições análogas são:
\[f_X(x) \geq 0 \quad \forall x, \qquad \int_{-\infty}^{+\infty} f_X(x)\,dx = 1.\]

Em qualquer dos casos, a \textbf{função de distribuição} (acumulada), \(F_X(x) = P(X \leq x)\),
acumula a probabilidade até \(x\) --- somando, no caso discreto, ou integrando, no contínuo:
\[F_X(x) = \sum_{t \leq x} f_X(t) \qquad \text{ou} \qquad F_X(x) = \int_{-\infty}^{x} f_X(t)\,dt.\]

Finalmente, dois números resumem a distribuição: o \textbf{valor esperado} \(E(X)\) (o centro,
a ``média teórica'') e a \textbf{variância} \(V(X)\) (a dispersão em torno desse centro). As suas
definições espelham a natureza discreta ou contínua da variável:

\[E(X) = \sum_x x\, f_X(x) \quad \text{ou} \quad E(X) = \int_{-\infty}^{+\infty} x\, f_X(x)\,dx,\]
\[V(X) = E\!\left[(X - E(X))^2\right] = \sum_x (x - E(X))^2 f_X(x) \quad \text{ou} \quad \int_{-\infty}^{+\infty} (x - E(X))^2 f_X(x)\,dx.\]

Repare no paralelo perfeito entre os dois mundos: onde a variável discreta \textbf{soma}
(\(\sum\)), a contínua \textbf{integra} (\(\int\)). É a mesma ideia --- acumular probabilidade ---
expressa nas duas linguagens. Guardar este paralelo em mente torna toda a matéria mais
simples.

\section{Exemplo: uma variável discreta}\label{exemplo-uma-variuxe1vel-discreta}

Numa empresa, o número de reclamações recebidas por dia é uma variável aleatória
discreta \(X\) com a seguinte f.m.p.:

\[f_X(x) = P(X = x) =
\begin{cases}
0{,}10, & x = 0 \\
0{,}25, & x = 1 \\
0{,}30, & x = 2 \\
0{,}20, & x = 3 \\
0{,}15, & x = 4
\end{cases}\]

Queremos (a) confirmar que \(f_X\) é uma f.m.p., (b) calcular o número esperado e a
variância das reclamações, e (c) a probabilidade de haver mais de 2 reclamações num dia.

Para (a), verificamos que as probabilidades são não negativas e somam 1:

\begin{Shaded}
\begin{Highlighting}[]
\NormalTok{X\_valores }\OtherTok{\textless{}{-}} \FunctionTok{c}\NormalTok{(}\DecValTok{0}\NormalTok{, }\DecValTok{1}\NormalTok{, }\DecValTok{2}\NormalTok{, }\DecValTok{3}\NormalTok{, }\DecValTok{4}\NormalTok{)}
\NormalTok{P\_X }\OtherTok{\textless{}{-}} \FunctionTok{c}\NormalTok{(}\FloatTok{0.10}\NormalTok{, }\FloatTok{0.25}\NormalTok{, }\FloatTok{0.30}\NormalTok{, }\FloatTok{0.20}\NormalTok{, }\FloatTok{0.15}\NormalTok{)}

\FunctionTok{sum}\NormalTok{(P\_X)          }\CommentTok{\# deve dar 1}
\DocumentationTok{\#\# [1] 1}
\end{Highlighting}
\end{Shaded}

Para (b), aplicamos diretamente as definições de \(E(X)\) e \(V(X)\):

\begin{Shaded}
\begin{Highlighting}[]
\NormalTok{E\_X   }\OtherTok{\textless{}{-}} \FunctionTok{sum}\NormalTok{(X\_valores }\SpecialCharTok{*}\NormalTok{ P\_X)                 }\CommentTok{\# valor esperado}
\NormalTok{Var\_X }\OtherTok{\textless{}{-}} \FunctionTok{sum}\NormalTok{((X\_valores }\SpecialCharTok{{-}}\NormalTok{ E\_X)}\SpecialCharTok{\^{}}\DecValTok{2} \SpecialCharTok{*}\NormalTok{ P\_X)       }\CommentTok{\# variância}
\NormalTok{E\_X}
\DocumentationTok{\#\# [1] 2.05}
\NormalTok{Var\_X}
\DocumentationTok{\#\# [1] 1.448}
\end{Highlighting}
\end{Shaded}

E para (c), somamos as probabilidades dos valores superiores a 2:

\begin{Shaded}
\begin{Highlighting}[]
\NormalTok{P\_X\_maior\_2 }\OtherTok{\textless{}{-}} \FunctionTok{sum}\NormalTok{(P\_X[X\_valores }\SpecialCharTok{\textgreater{}} \DecValTok{2}\NormalTok{])}
\NormalTok{P\_X\_maior\_2}
\DocumentationTok{\#\# [1] 0.35}
\end{Highlighting}
\end{Shaded}

Em média, a empresa recebe cerca de 2 reclamações por dia, e a probabilidade de receber
mais de 2 é de 0,35.

\section{Exemplo: uma variável contínua}\label{exemplo-uma-variuxe1vel-contuxednua}

Seja \(X\) uma variável contínua com f.d.p.

\[f(x) = \begin{cases} 2e^{-2x}, & x \geq 0 \\ 0, & x < 0 \end{cases}\]

(trata-se de uma distribuição exponencial). Queremos (a) confirmar que é uma f.d.p.,
(b) calcular \(P(X > 2)\) e (c) \(P(0{,}3 < X < 0{,}7)\).

Começamos por definir a função em R:

\begin{Shaded}
\begin{Highlighting}[]
\NormalTok{f1 }\OtherTok{\textless{}{-}} \ControlFlowTok{function}\NormalTok{(x) \{}
  \FunctionTok{ifelse}\NormalTok{(x }\SpecialCharTok{\textless{}} \DecValTok{0}\NormalTok{, }\DecValTok{0}\NormalTok{, }\DecValTok{2} \SpecialCharTok{*} \FunctionTok{exp}\NormalTok{(}\SpecialCharTok{{-}}\DecValTok{2} \SpecialCharTok{*}\NormalTok{ x))}
\NormalTok{\}}
\end{Highlighting}
\end{Shaded}

O seu gráfico mostra a forma típica de decaimento da exponencial:

\begin{Shaded}
\begin{Highlighting}[]
\FunctionTok{plot}\NormalTok{(f1, }\AttributeTok{from =} \DecValTok{0}\NormalTok{, }\AttributeTok{to =} \DecValTok{5}\NormalTok{,}
     \AttributeTok{main =} \StringTok{"Densidade exponencial"}\NormalTok{,}
     \AttributeTok{xlab =} \StringTok{"x"}\NormalTok{, }\AttributeTok{ylab =} \StringTok{"f(x)"}\NormalTok{, }\AttributeTok{col =} \StringTok{"steelblue"}\NormalTok{, }\AttributeTok{lwd =} \DecValTok{2}\NormalTok{)}
\end{Highlighting}
\end{Shaded}

\begin{figure}

{\centering \includegraphics[width=0.8\linewidth]{introR_files/figure-latex/prob-exp-1} 

}

\caption{Função densidade da distribuição exponencial de parâmetro 2.}\label{fig:prob-exp}
\end{figure}

Para (a), a condição \(\int_{-\infty}^{+\infty} f(x)\,dx = 1\) verifica-se por integração
numérica com \texttt{integrate()}:

\begin{Shaded}
\begin{Highlighting}[]
\FunctionTok{integrate}\NormalTok{(f1, }\AttributeTok{lower =} \DecValTok{0}\NormalTok{, }\AttributeTok{upper =} \ConstantTok{Inf}\NormalTok{)}
\DocumentationTok{\#\# 1 with absolute error \textless{} 5e{-}07}
\end{Highlighting}
\end{Shaded}

As alíneas (b) e (c) correspondem a áreas sob a curva, ou seja, a integrais nos
intervalos pedidos:
\[P(X > 2) = \int_2^{+\infty} 2e^{-2x}\,dx, \qquad P(0{,}3 < X < 0{,}7) = \int_{0{,}3}^{0{,}7} 2e^{-2x}\,dx.\]

\begin{Shaded}
\begin{Highlighting}[]
\FunctionTok{integrate}\NormalTok{(f1, }\AttributeTok{lower =} \DecValTok{2}\NormalTok{, }\AttributeTok{upper =} \ConstantTok{Inf}\NormalTok{)      }\CommentTok{\# P(X \textgreater{} 2)}
\DocumentationTok{\#\# 0.01832 with absolute error \textless{} 2.8e{-}06}
\FunctionTok{integrate}\NormalTok{(f1, }\AttributeTok{lower =} \FloatTok{0.3}\NormalTok{, }\AttributeTok{upper =} \FloatTok{0.7}\NormalTok{)    }\CommentTok{\# P(0.3 \textless{} X \textless{} 0.7)}
\DocumentationTok{\#\# 0.3022 with absolute error \textless{} 3.4e{-}15}
\end{Highlighting}
\end{Shaded}

A \texttt{integrate()} devolve um objeto com vários elementos; o valor do integral está em
\texttt{\$value}. Assim, para usar o resultado num cálculo posterior, escrevemos, por exemplo,
\texttt{integrate(f1,\ 0,\ Inf)\$value}. Guardar este pormenor evita muita frustração quando se
encadeiam contas.

Quando a forma da distribuição é conhecida, a sua função de distribuição pode ser
escrita de forma fechada, sem integrar. Para esta exponencial,
\(F_X(x) = 1 - e^{-2x}\) para \(x \geq 0\), e \(P(X > 2) = 1 - F_X(2) = e^{-4} \approx 0{,}018\)
--- exatamente o valor que a integração numérica devolveu. Comparar o cálculo analítico
com o numérico é uma excelente forma de confirmar que ambos estão certos.

\section{Exemplo: uma densidade definida por troços}\label{exemplo-uma-densidade-definida-por-trouxe7os}

A procura diária de arroz num supermercado, em centenas de quilogramas, é uma variável
aleatória \(X\) com f.d.p.

\[f(x) = \begin{cases}
\frac{2}{3}x, & 0 \leq x < 1 \\[4pt]
-\frac{x}{3} + 1, & 1 \leq x < 3 \\[4pt]
0, & x < 0 \ \text{ou} \ x \geq 3
\end{cases}\]

Queremos (a) a probabilidade de se venderem mais de 150 kg num dia e (b) a venda
esperada em 30 dias. Definimos a função por troços com \texttt{ifelse()} encadeados:

\begin{Shaded}
\begin{Highlighting}[]
\NormalTok{f2 }\OtherTok{\textless{}{-}} \ControlFlowTok{function}\NormalTok{(x) \{}
  \FunctionTok{ifelse}\NormalTok{(x }\SpecialCharTok{\textgreater{}=} \DecValTok{0} \SpecialCharTok{\&}\NormalTok{ x }\SpecialCharTok{\textless{}} \DecValTok{1}\NormalTok{, (}\DecValTok{2}\SpecialCharTok{/}\DecValTok{3}\NormalTok{) }\SpecialCharTok{*}\NormalTok{ x,}
         \FunctionTok{ifelse}\NormalTok{(x }\SpecialCharTok{\textgreater{}=} \DecValTok{1} \SpecialCharTok{\&}\NormalTok{ x }\SpecialCharTok{\textless{}} \DecValTok{3}\NormalTok{, }\SpecialCharTok{{-}}\NormalTok{x}\SpecialCharTok{/}\DecValTok{3} \SpecialCharTok{+} \DecValTok{1}\NormalTok{, }\DecValTok{0}\NormalTok{))}
\NormalTok{\}}
\end{Highlighting}
\end{Shaded}

Confirmamos que é uma f.d.p. (o integral vale 1) e vemos o seu gráfico --- uma densidade
triangular:

\begin{Shaded}
\begin{Highlighting}[]
\FunctionTok{integrate}\NormalTok{(f2, }\DecValTok{0}\NormalTok{, }\DecValTok{3}\NormalTok{)}
\DocumentationTok{\#\# 1 with absolute error \textless{} 1.1e{-}15}

\FunctionTok{plot}\NormalTok{(f2, }\AttributeTok{from =} \SpecialCharTok{{-}}\DecValTok{1}\NormalTok{, }\AttributeTok{to =} \DecValTok{4}\NormalTok{,}
     \AttributeTok{main =} \StringTok{"Densidade da procura de arroz"}\NormalTok{,}
     \AttributeTok{xlab =} \StringTok{"x (centenas de kg)"}\NormalTok{, }\AttributeTok{ylab =} \StringTok{"f(x)"}\NormalTok{, }\AttributeTok{col =} \StringTok{"steelblue"}\NormalTok{, }\AttributeTok{lwd =} \DecValTok{2}\NormalTok{)}
\end{Highlighting}
\end{Shaded}

\begin{figure}

{\centering \includegraphics[width=0.8\linewidth]{introR_files/figure-latex/prob-troços-1} 

}

\caption{Densidade da procura de arroz, definida por dois troços.}\label{fig:prob-troços}
\end{figure}

Para (a), 150 kg correspondem a \(x = 1{,}5\) centenas de quilogramas, pelo que queremos
\(P(X > 1{,}5) = \int_{1{,}5}^{3} f(x)\,dx\). Podemos assinalar essa área no gráfico:

\begin{Shaded}
\begin{Highlighting}[]
\FunctionTok{plot}\NormalTok{(f2, }\AttributeTok{from =} \SpecialCharTok{{-}}\DecValTok{1}\NormalTok{, }\AttributeTok{to =} \DecValTok{4}\NormalTok{,}
     \AttributeTok{main =} \StringTok{"P(X \textgreater{} 1,5)"}\NormalTok{, }\AttributeTok{xlab =} \StringTok{"x"}\NormalTok{, }\AttributeTok{ylab =} \StringTok{"f(x)"}\NormalTok{, }\AttributeTok{col =} \StringTok{"steelblue"}\NormalTok{, }\AttributeTok{lwd =} \DecValTok{2}\NormalTok{)}
\FunctionTok{polygon}\NormalTok{(}\AttributeTok{x =} \FunctionTok{c}\NormalTok{(}\FloatTok{1.5}\NormalTok{, }\FloatTok{1.5}\NormalTok{, }\DecValTok{3}\NormalTok{), }\AttributeTok{y =} \FunctionTok{c}\NormalTok{(}\DecValTok{0}\NormalTok{, }\FunctionTok{f2}\NormalTok{(}\FloatTok{1.5}\NormalTok{), }\DecValTok{0}\NormalTok{), }\AttributeTok{density =} \DecValTok{15}\NormalTok{, }\AttributeTok{col =} \StringTok{"firebrick"}\NormalTok{)}
\end{Highlighting}
\end{Shaded}

\begin{figure}

{\centering \includegraphics[width=0.8\linewidth]{introR_files/figure-latex/prob-area-1} 

}

\caption{A área sombreada representa P(X > 1,5).}\label{fig:prob-area}
\end{figure}

\begin{Shaded}
\begin{Highlighting}[]
\FunctionTok{integrate}\NormalTok{(f2, }\FloatTok{1.5}\NormalTok{, }\DecValTok{3}\NormalTok{)      }\CommentTok{\# P(X \textgreater{} 1,5)}
\DocumentationTok{\#\# 0.375 with absolute error \textless{} 4.2e{-}15}
\end{Highlighting}
\end{Shaded}

Para (b), a venda esperada num dia é \(E(X) = \int_{-\infty}^{+\infty} x\, f(x)\,dx\);
definimos a função integranda \(x\,f(x)\) e integramos. A venda esperada em 30 dias é 30
vezes esse valor:

\begin{Shaded}
\begin{Highlighting}[]
\NormalTok{ef2 }\OtherTok{\textless{}{-}} \ControlFlowTok{function}\NormalTok{(x) x }\SpecialCharTok{*} \FunctionTok{f2}\NormalTok{(x)}

\NormalTok{E\_X }\OtherTok{\textless{}{-}} \FunctionTok{integrate}\NormalTok{(ef2, }\DecValTok{0}\NormalTok{, }\DecValTok{3}\NormalTok{)}\SpecialCharTok{$}\NormalTok{value}
\NormalTok{E\_X            }\CommentTok{\# venda esperada num dia (centenas de kg)}
\DocumentationTok{\#\# [1] 1.333}
\DecValTok{30} \SpecialCharTok{*}\NormalTok{ E\_X       }\CommentTok{\# venda esperada em 30 dias}
\DocumentationTok{\#\# [1] 40}
\end{Highlighting}
\end{Shaded}

Note-se que os exemplos deste capítulo são simples e poderiam resolver-se
analiticamente, sem métodos numéricos --- foram escolhidos apenas para ilustrar as
funções do R. O verdadeiro valor da integração numérica revela-se nos problemas
complexos, em que a solução analítica é difícil ou impossível de obter. É aí que
ferramentas como a \texttt{integrate()} se tornam indispensáveis.

\section{Exercícios}\label{exercuxedcios-23}

\textbf{1.} O número de produtos defeituosos numa linha de produção, ao fim de cada hora, é
uma variável aleatória discreta \(X\) com f.m.p.:

\[P(X = x) = \begin{cases}
1/10, & x = 0 \\ 2/10, & x = 1 \\ 3/10, & x = 2 \\ 2/10, & x = 3 \\ 2/10, & x = 4
\end{cases}\]

\begin{enumerate}
\def\labelenumi{(\alph{enumi})}
\tightlist
\item
  Verifique se é uma f.m.p.
\item
  Calcule o valor esperado e a variância do número de produtos defeituosos por hora.
\item
  Calcule a probabilidade de haver, no máximo, 2 produtos defeituosos numa hora.
\end{enumerate}

\textbf{2.} Um inquérito numa universidade registou o número de cadeiras em que cada aluno
está inscrito. A variável aleatória \(Z\) tem f.m.p.:

\[P(Z = z) = \begin{cases}
0{,}05, & z = 1 \\ 0{,}15, & z = 2 \\ 0{,}35, & z = 3 \\ 0{,}25, & z = 4 \\ 0{,}20, & z = 5
\end{cases}\]

\begin{enumerate}
\def\labelenumi{(\alph{enumi})}
\tightlist
\item
  Verifique se é uma f.m.p.
\item
  Calcule o valor esperado e a variância do número de cadeiras.
\item
  Calcule a probabilidade de um aluno estar inscrito em mais de 3 cadeiras.
\end{enumerate}

\textbf{3.} Numa dada localidade, a distribuição do rendimento, em u.m. (unidades
monetárias), é uma variável aleatória \(X\) com \textbf{função densidade de probabilidade}:

\[f_X(x) = \begin{cases}
\frac{1}{10}x + \frac{1}{10}, & 0 \leq x \leq 2 \\[4pt]
-\frac{3}{40}x + \frac{9}{20}, & 2 < x \leq 6 \\[4pt]
0, & x < 0 \ \text{ou} \ x > 6
\end{cases}\]

\begin{enumerate}
\def\labelenumi{(\alph{enumi})}
\tightlist
\item
  Mostre que \(f_X\) é uma f.d.p.
\item
  Qual o rendimento médio nesta localidade?
\item
  Calcule a probabilidade de encontrar uma pessoa com rendimento acima de 4,5 u.m. e
  assinale o resultado no gráfico da densidade.
\end{enumerate}

\textbf{4.} Uma variável aleatória contínua \(X\) tem distribuição uniforme no intervalo
\([10, 20]\).

\begin{enumerate}
\def\labelenumi{(\alph{enumi})}
\tightlist
\item
  Apresente o gráfico da f.d.p.
\item
  Calcule \(P(X < 15)\).
\item
  Calcule \(P(12 \leq X \leq 18)\).
\item
  Calcule \(E(X)\) e \(V(X)\).
\end{enumerate}

\textbf{5.} O tempo de espera (em minutos) num balcão de atendimento é uma variável aleatória
contínua \(Y\) com f.d.p.

\[f_Y(y) = \begin{cases} 0{,}125\,y, & 0 \leq y \leq 4 \\ 0, & \text{caso contrário} \end{cases}\]

\begin{enumerate}
\def\labelenumi{(\alph{enumi})}
\tightlist
\item
  Crie uma função para a função de distribuição acumulada de \(Y\).
\item
  Calcule \(P(Y > 3)\).
\item
  Calcule \(P(1 < Y < 3)\).
\item
  Calcule \(P(Y < 2)\).
\item
  Calcule \(P(Y < 3 \mid Y > 1)\) (probabilidade condicional).
\end{enumerate}

\textbf{6.} A duração, em anos, de uma lâmpada especial é uma variável aleatória contínua \(X\)
com densidade

\[f_X(x) = \begin{cases} 3e^{-3x}, & x \geq 0 \\ 0, & x < 0 \end{cases}\]

\begin{enumerate}
\def\labelenumi{(\alph{enumi})}
\tightlist
\item
  Crie uma função para a função de distribuição acumulada de \(X\).
\item
  Calcule \(P(X > 2)\).
\item
  Calcule \(P(0{,}5 < X < 1{,}2)\).
\item
  Calcule \(P(X > 3)\).
\item
  Calcule \(P(X < 3 \mid X > 1)\).
\end{enumerate}

Descrita a probabilidade de uma variável aleatória, o passo natural é conhecer as
\textbf{distribuições} mais importantes --- a binomial, a normal, a exponencial --- e as funções
que o R lhes dedica. É esse o tema do próximo capítulo.

\chapter{Simulação}\label{simulacao}

Chegámos ao coração do livro. Simular é usar o computador para \emph{imitar} um fenómeno
aleatório --- repeti-lo milhares ou milhões de vezes --- e, a partir dos resultados
observados, responder a perguntas que seriam difíceis, ou mesmo impossíveis, de
resolver só com papel e caneta.

O fundamento teórico da simulação é a \textbf{Lei dos Grandes Números}: se observarmos uma
grande amostra de variáveis aleatórias independentes e identicamente distribuídas
(i.i.d.) com média finita, a média dessas observações converge para a média verdadeira
da distribuição à medida que a dimensão da amostra aumenta. Por outras palavras, em vez
de calcular analiticamente uma média (ou uma probabilidade, que é a média de uma
variável indicadora), podemos gerar uma amostra grande e simplesmente calcular a média
observada --- que será uma boa estimativa do valor verdadeiro.

A eficácia deste método depende de três fatores:

\begin{enumerate}
\def\labelenumi{\arabic{enumi}.}
\tightlist
\item
  \textbf{Identificar as variáveis aleatórias certas} --- que distribuições descrevem
  corretamente o fenómeno que estamos a modelar.
\item
  \textbf{Gerá-las com fidelidade} --- os algoritmos de geração de números pseudoaleatórios
  têm de produzir amostras que representem bem as distribuições pretendidas.
\item
  \textbf{Escolher uma dimensão de amostra adequada} --- grande o suficiente para que a
  estimativa seja fiável, sabendo que amostras maiores reduzem a variabilidade.
\end{enumerate}

Neste capítulo compreendemos o princípio da simulação através da Lei dos Grandes
Números, vemos como o R gera números pseudoaleatórios (e o que isso significa) e
resolvemos, por simulação, um problema cuja solução analítica seria muito trabalhosa.

\section{A simulação reproduz a teoria}\label{a-simulauxe7uxe3o-reproduz-a-teoria}

\textbf{Exemplo (estimar a média de uma uniforme).} A média da distribuição uniforme em
\([0, 1]\) é, sabidamente, \(\frac{1}{2}\). Pela Lei dos Grandes Números, a média de uma
amostra \(X_1, \ldots, X_n\) dessa distribuição, \(\overline{X} = \frac{1}{n}\sum_{i=1}^n X_i\),
deve aproximar-se de 0,5 à medida que \(n\) cresce.

A tabela seguinte mostra as médias de várias amostras simuladas, para diferentes
dimensões \(n\). Note-se como as médias estão sempre próximas de 0,5, mas com uma variação
que \textbf{diminui} claramente à medida que \(n\) aumenta:

{\def\LTcaptype{none} % do not increment counter
\begin{longtable}[]{@{}llllll@{}}
\toprule\noalign{}
\(n\) & Rép. 1 & Rép. 2 & Rép. 3 & Rép. 4 & Rép. 5 \\
\midrule\noalign{}
\endhead
\bottomrule\noalign{}
\endlastfoot
100 & 0,485 & 0,481 & 0,484 & 0,569 & 0,441 \\
1 000 & 0,497 & 0,506 & 0,480 & 0,498 & 0,499 \\
10 000 & 0,502 & 0,501 & 0,499 & 0,498 & 0,498 \\
100 000 & 0,502 & 0,499 & 0,500 & 0,498 & 0,499 \\
\end{longtable}
}

Para gerar uma amostra uniforme em \([0, 1]\) usa-se a função \texttt{runif()}:

\begin{Shaded}
\begin{Highlighting}[]
\FunctionTok{runif}\NormalTok{(}\DecValTok{100}\NormalTok{, }\AttributeTok{min =} \DecValTok{0}\NormalTok{, }\AttributeTok{max =} \DecValTok{1}\NormalTok{)}
\end{Highlighting}
\end{Shaded}

Aqui não havia necessidade real de simular, pois conhecíamos a média analiticamente ---
o exemplo serve apenas para mostrar que a simulação reproduz a teoria. A lição a reter
é esta: por maior que seja a amostra, a média amostral \textbf{não} será exatamente igual à
média verdadeira; a simulação introduz uma variabilidade que temos de considerar, e que
só se reduz aumentando \(n\).

\section{Quando a simulação vale mesmo a pena}\label{quando-a-simulauxe7uxe3o-vale-mesmo-a-pena}

Vejamos agora um problema fácil de descrever, mas cuja solução analítica seria penosa.

\textbf{O problema.} Dois empregados de balcão, A e B, começam a atender clientes ao mesmo
tempo e combinam fazer uma pausa assim que ambos tiverem atendido 10 clientes.
Provavelmente um terminará antes do outro e terá de esperar. Quanto tempo, em média,
terá um deles de esperar pelo outro?

Modelamos os tempos de atendimento de cada cliente como variáveis i.i.d. com
distribuição exponencial de taxa \(\lambda = 0{,}3\) clientes por minuto. Então, o tempo
total que cada empregado leva a atender 10 clientes é a soma de 10 exponenciais i.i.d.,
que segue uma distribuição \textbf{gama} de forma \(k = 10\) e taxa \(\lambda = 0{,}3\). Sejam
\(X\) e \(Y\) esses tempos para A e B; queremos \(E(|X - Y|)\).

A solução analítica exigiria um integral duplo sobre uma região não retangular. A
simulação resolve o problema em poucas linhas: geramos muitos pares \((X, Y)\), calculamos
\(Z = |X - Y|\) em cada um, e estimamos \(E(|X - Y|)\) pela média dos valores de \(Z\).

\begin{Shaded}
\begin{Highlighting}[]
\FunctionTok{set.seed}\NormalTok{(}\DecValTok{123}\NormalTok{)}

\CommentTok{\# Tempos para atender 10 clientes: soma de 10 exponenciais = Gama(forma=10, taxa=0.3)}
\NormalTok{x }\OtherTok{\textless{}{-}} \FunctionTok{rgamma}\NormalTok{(}\DecValTok{10000}\NormalTok{, }\AttributeTok{shape =} \DecValTok{10}\NormalTok{, }\AttributeTok{rate =} \FloatTok{0.3}\NormalTok{)}
\NormalTok{y }\OtherTok{\textless{}{-}} \FunctionTok{rgamma}\NormalTok{(}\DecValTok{10000}\NormalTok{, }\AttributeTok{shape =} \DecValTok{10}\NormalTok{, }\AttributeTok{rate =} \FloatTok{0.3}\NormalTok{)}

\NormalTok{z }\OtherTok{\textless{}{-}} \FunctionTok{abs}\NormalTok{(x }\SpecialCharTok{{-}}\NormalTok{ y)      }\CommentTok{\# tempo de espera em cada réplica}
\FunctionTok{mean}\NormalTok{(z)              }\CommentTok{\# estimativa de E(|X {-} Y|)}
\DocumentationTok{\#\# [1] 11.76}
\end{Highlighting}
\end{Shaded}

A simulação estima um tempo médio de espera de cerca de 11,7 minutos. O histograma
mostra a distribuição completa desses tempos:

\begin{Shaded}
\begin{Highlighting}[]
\FunctionTok{hist}\NormalTok{(z,}
     \AttributeTok{main =} \StringTok{"Tempo de espera"}\NormalTok{,}
     \AttributeTok{xlab =} \StringTok{"Tempo de espera (minutos)"}\NormalTok{,}
     \AttributeTok{ylab =} \StringTok{"Frequência absoluta"}\NormalTok{,}
     \AttributeTok{col =} \StringTok{"lightblue"}\NormalTok{, }\AttributeTok{border =} \StringTok{"black"}\NormalTok{)}
\end{Highlighting}
\end{Shaded}

\begin{figure}

{\centering \includegraphics[width=0.8\linewidth]{introR_files/figure-latex/sim-espera-1} 

}

\caption{Distribuição do tempo de espera de um empregado pelo outro.}\label{fig:sim-espera}
\end{figure}

\section{Números pseudoaleatórios}\label{nuxfameros-pseudoaleatuxf3rios}

Vale a pena perceber o que o computador faz quando lhe pedimos ``acaso''.

Números \textbf{verdadeiramente aleatórios} derivam de processos físicos imprevisíveis --- o
ruído térmico de um circuito, o decaimento radioativo, a radiação cósmica. Obtê-los é
difícil e, para a maioria das aplicações, desnecessário.

Números \textbf{pseudoaleatórios}, pelo contrário, são produzidos por um algoritmo
determinístico que, a partir de um valor inicial --- a \textbf{semente} (\emph{seed}) ---, gera uma
sequência que \emph{parece} aleatória. Como o algoritmo é determinístico, a mesma semente
produz sempre a mesma sequência. Não sendo verdadeiramente aleatórios, são rápidos de
gerar, suficientemente imprevisíveis para quase tudo e, sobretudo, \textbf{reprodutíveis} ---
uma vantagem decisiva para a investigação e para a depuração de código.

\subsection{O método congruencial multiplicativo}\label{o-muxe9todo-congruencial-multiplicativo}

Um dos geradores clássicos é o \textbf{método congruencial multiplicativo} (ou de Lehmer). A
partir de uma semente \(x_0\), os valores seguintes obtêm-se pela recorrência
\[x_n = (a \cdot x_{n-1}) \bmod m,\]
onde \(a\) e \(m\) são inteiros positivos dados e \(x_n\) é o resto da divisão inteira de
\(a \cdot x_{n-1}\) por \(m\). Cada quociente \(x_n / m\) é um número pseudoaleatório,
aproximando uma variável uniforme em \((0, 1)\).

Para servir numa simulação, as constantes \(a\) e \(m\) devem garantir que a sequência
\emph{pareça} aleatória, tenha um \textbf{período longo} antes de se repetir e seja eficiente de
calcular. Com \(m = 17\), \(a = 5\) e \(x_0 = 7\):

\[x_1 = (5 \times 7) \bmod 17 = 1, \quad x_2 = (5 \times 1) \bmod 17 = 5, \quad x_3 = (5 \times 5) \bmod 17 = 8, \ \ldots\]

Continuando, obtém-se a sequência
\[7,\, 1,\, 5,\, 8,\, 6,\, 13,\, 14,\, 2,\, 10,\, 16,\, 12,\, 9,\, 11,\, 4,\, 3,\, 15,\, \underbrace{7}_{\text{repete}},\, 1,\, 5,\, 8, \ldots\]
com período 16.

Não é por acaso que o período é exatamente 16. Como \(m = 17\) é primo e \(a = 5\) é uma
\emph{raiz primitiva} módulo 17, a sequência percorre \textbf{todos} os valores de 1 a \(m - 1\)
antes de repetir --- o chamado período completo, \(m - 1\). Note-se ainda que, sendo o
gerador multiplicativo, o valor 0 nunca aparece: os \(x_n\) estão sempre entre 1 e
\(m - 1\), e os \(x_n/m\) em \((0, 1)\). Escolher bem \(a\) e \(m\) é, por isto, uma questão
delicada --- os geradores usados na prática têm períodos astronomicamente longos.

\section{\texorpdfstring{A função \texttt{set.seed()}}{A função set.seed()}}\label{a-funuxe7uxe3o-set.seed}

No R, a semente do gerador fixa-se com \texttt{set.seed()}. Executar \texttt{set.seed()} com o mesmo
número antes de uma operação aleatória garante que ela produz sempre os mesmos
resultados:

\begin{Shaded}
\begin{Highlighting}[]
\FunctionTok{set.seed}\NormalTok{(}\DecValTok{123}\NormalTok{)}
\FunctionTok{sample}\NormalTok{(}\DecValTok{1}\SpecialCharTok{:}\DecValTok{100}\NormalTok{, }\DecValTok{5}\NormalTok{)}
\DocumentationTok{\#\# [1] 31 79 51 14 67}
\end{Highlighting}
\end{Shaded}

Reexecutando o mesmo código com a mesma semente, obtêm-se exatamente os mesmos cinco
números. O gerador que o R usa por omissão é o \emph{Mersenne Twister}, com um período
verdadeiramente enorme; quando não fixamos a semente, o R constrói uma a partir da hora
e da sessão, pelo que cada execução parte de um ponto diferente.

Fixar a semente é um hábito, não uma limitação: os resultados continuam a comportar-se
como acaso, mas tornam-se reprodutíveis. Todas as simulações deste livro que devam dar
sempre o mesmo resultado começam por \texttt{set.seed()}.

\section{Exercícios}\label{exercuxedcios-24}

\textbf{1.} Fixe uma semente à sua escolha com \texttt{set.seed()}.

\begin{itemize}
\tightlist
\item
  Gere duas amostras de dimensão 10 dos números de 1 a 50, sem reposição.
\item
  Mude a semente e gere novamente as amostras.
\item
  Compare os resultados antes e depois de mudar a semente.
\end{itemize}

\textbf{2.} Fixe a semente com \texttt{set.seed(123)}.

\begin{itemize}
\tightlist
\item
  Simule 100 lançamentos de um dado (1 a 6) e conte a frequência de cada face.
\item
  Repita o processo sem redefinir a semente.
\item
  Volte a fixar \texttt{set.seed(123)} e repita. Compare com a primeira simulação.
\end{itemize}

\textbf{3.} Use \texttt{set.seed()} para gerar 50 valores de uma normal \(N(\mu = 10, \sigma = 2)\)
com \texttt{rnorm(50,\ mean\ =\ 10,\ sd\ =\ 2)}.

\begin{itemize}
\tightlist
\item
  Com a mesma semente, gere 50 valores de uma uniforme \(U(0, 1)\) com \texttt{runif(50)}.
\item
  Mude a semente e repita. Compare os conjuntos de dados gerados.
\end{itemize}

\textbf{4.} Fixe a semente com \texttt{set.seed(42)}.

\begin{itemize}
\tightlist
\item
  Divida os números de 1 a 100 em dois grupos aleatórios de 50 (use \texttt{sample()} e
  \texttt{setdiff()}).
\item
  Repita a divisão sem fixar a semente.
\item
  Volte a fixar \texttt{set.seed(42)} e repita. Verifique se os grupos são os mesmos.
\end{itemize}

\textbf{5.} Use \texttt{set.seed()} para simular as respostas de 50 participantes a uma pergunta de
escolha múltipla com 4 opções (A, B, C, D).

\begin{itemize}
\tightlist
\item
  Mude a semente e repita a simulação.
\item
  Verifique se a distribuição das respostas muda ao alterar, ou ao manter, a semente.
\end{itemize}

Sabemos agora simular e reproduzir o acaso. No próximo capítulo passamos em revista as
\textbf{distribuições de probabilidade} mais importantes e as funções que o R lhes dedica ---
o vocabulário indispensável de qualquer simulação.

\chapter{Distribuições de probabilidade}\label{distribuicoes}

No capítulo anterior vimos que simular consiste, no essencial, em gerar variáveis
aleatórias com a distribuição certa. Chega, por isso, o momento de conhecer as
\textbf{distribuições de probabilidade} mais importantes e as funções que o R lhes dedica.
Cada uma modela um tipo de fenómeno --- contagens, tempos de espera, medições --- e para
cada uma o R oferece, de forma sistemática, quatro funções.

O R inclui as distribuições univariadas mais comuns. Para cada distribuição existem
quatro funções, distinguidas por uma letra inicial:

{\def\LTcaptype{none} % do not increment counter
\begin{longtable}[]{@{}
  >{\raggedright\arraybackslash}p{(\linewidth - 2\tabcolsep) * \real{0.5000}}
  >{\raggedright\arraybackslash}p{(\linewidth - 2\tabcolsep) * \real{0.5000}}@{}}
\toprule\noalign{}
\begin{minipage}[b]{\linewidth}\raggedright
Função
\end{minipage} & \begin{minipage}[b]{\linewidth}\raggedright
O que devolve
\end{minipage} \\
\midrule\noalign{}
\endhead
\bottomrule\noalign{}
\endlastfoot
\texttt{d}\emph{nome}(\ldots) & a função de probabilidade (discreta) ou densidade (contínua) \\
\texttt{p}\emph{nome}(\ldots) & a função de distribuição, \(F(x) = P(X \leq x)\) \\
\texttt{q}\emph{nome}(\ldots) & o quantil correspondente a uma dada probabilidade (o inverso de \texttt{p}) \\
\texttt{r}\emph{nome}(\ldots) & uma amostra aleatória da distribuição \\
\end{longtable}
}

O \emph{nome} é a abreviatura da distribuição: \texttt{binom}, \texttt{geom}, \texttt{pois}, \texttt{unif}, \texttt{exp},
\texttt{norm}, e assim por diante.

Vale a pena fixar esta lógica de quatro letras, porque se aplica a \emph{todas} as
distribuições do R. Se sabe usar \texttt{dnorm}, \texttt{pnorm}, \texttt{qnorm} e \texttt{rnorm}, sabe usar as
funções de qualquer outra distribuição --- muda apenas a raiz do nome e, eventualmente,
os parâmetros. É uma das grandes comodidades da linguagem.

Neste capítulo percorremos as principais distribuições --- uniforme, Bernoulli, binomial,
geométrica, Poisson (discretas), uniforme, exponencial e normal (contínuas) ---, vendo
para cada uma a definição, o valor esperado e a variância, o cálculo de probabilidades
com \texttt{d}/\texttt{p}/\texttt{q} e a geração de amostras com \texttt{r}, comparando sempre a teoria com a
simulação.

\section{Primeiros exemplos}\label{primeiros-exemplos}

\textbf{Exemplo 1 (lançar três moedas).} Simulemos o lançamento de três moedas equilibradas
e a contagem do número de caras, \(X\). Queremos estimar \(P(X = 1)\) e \(E(X)\); e depois
repetir para uma moeda enviesada com \(P(\text{cara}) = 3/4\).

\begin{Shaded}
\begin{Highlighting}[]
\FunctionTok{set.seed}\NormalTok{(}\DecValTok{123}\NormalTok{)}
\NormalTok{n }\OtherTok{\textless{}{-}} \DecValTok{10000}
\NormalTok{sim1 }\OtherTok{\textless{}{-}} \FunctionTok{numeric}\NormalTok{(n)   }\CommentTok{\# indicador de "saiu exatamente 1 cara"}
\NormalTok{sim2 }\OtherTok{\textless{}{-}} \FunctionTok{numeric}\NormalTok{(n)   }\CommentTok{\# número de caras}
\ControlFlowTok{for}\NormalTok{ (i }\ControlFlowTok{in} \DecValTok{1}\SpecialCharTok{:}\NormalTok{n) \{}
\NormalTok{  moedas }\OtherTok{\textless{}{-}} \FunctionTok{sample}\NormalTok{(}\DecValTok{0}\SpecialCharTok{:}\DecValTok{1}\NormalTok{, }\DecValTok{3}\NormalTok{, }\AttributeTok{replace =} \ConstantTok{TRUE}\NormalTok{)}
\NormalTok{  sim1[i] }\OtherTok{\textless{}{-}} \ControlFlowTok{if}\NormalTok{ (}\FunctionTok{sum}\NormalTok{(moedas) }\SpecialCharTok{==} \DecValTok{1}\NormalTok{) }\DecValTok{1} \ControlFlowTok{else} \DecValTok{0}
\NormalTok{  sim2[i] }\OtherTok{\textless{}{-}} \FunctionTok{sum}\NormalTok{(moedas)}
\NormalTok{\}}
\FunctionTok{mean}\NormalTok{(sim1)   }\CommentTok{\# estimativa de P(X = 1)}
\DocumentationTok{\#\# [1] 0.3821}
\FunctionTok{mean}\NormalTok{(sim2)   }\CommentTok{\# estimativa de E(X)}
\DocumentationTok{\#\# [1] 1.4928}
\end{Highlighting}
\end{Shaded}

Para a moeda enviesada, basta indicar as probabilidades em \texttt{sample()}:

\begin{Shaded}
\begin{Highlighting}[]
\FunctionTok{set.seed}\NormalTok{(}\DecValTok{123}\NormalTok{)}
\NormalTok{n }\OtherTok{\textless{}{-}} \DecValTok{10000}
\NormalTok{sim1 }\OtherTok{\textless{}{-}} \FunctionTok{numeric}\NormalTok{(n)}
\NormalTok{sim2 }\OtherTok{\textless{}{-}} \FunctionTok{numeric}\NormalTok{(n)}
\ControlFlowTok{for}\NormalTok{ (i }\ControlFlowTok{in} \DecValTok{1}\SpecialCharTok{:}\NormalTok{n) \{}
\NormalTok{  moedas }\OtherTok{\textless{}{-}} \FunctionTok{sample}\NormalTok{(}\FunctionTok{c}\NormalTok{(}\DecValTok{0}\NormalTok{, }\DecValTok{1}\NormalTok{), }\DecValTok{3}\NormalTok{, }\AttributeTok{prob =} \FunctionTok{c}\NormalTok{(}\DecValTok{1}\SpecialCharTok{/}\DecValTok{4}\NormalTok{, }\DecValTok{3}\SpecialCharTok{/}\DecValTok{4}\NormalTok{), }\AttributeTok{replace =} \ConstantTok{TRUE}\NormalTok{)}
\NormalTok{  sim1[i] }\OtherTok{\textless{}{-}} \ControlFlowTok{if}\NormalTok{ (}\FunctionTok{sum}\NormalTok{(moedas) }\SpecialCharTok{==} \DecValTok{1}\NormalTok{) }\DecValTok{1} \ControlFlowTok{else} \DecValTok{0}
\NormalTok{  sim2[i] }\OtherTok{\textless{}{-}} \FunctionTok{sum}\NormalTok{(moedas)}
\NormalTok{\}}
\FunctionTok{mean}\NormalTok{(sim1)}
\DocumentationTok{\#\# [1] 0.1384}
\FunctionTok{mean}\NormalTok{(sim2)}
\DocumentationTok{\#\# [1] 2.2503}
\end{Highlighting}
\end{Shaded}

Mas o número de caras em três lançamentos é, por definição, uma
\(\text{Binomial}(n = 3, p)\). Podemos, por isso, gerar as observações diretamente com
\texttt{rbinom()}, sem o ciclo:

\begin{Shaded}
\begin{Highlighting}[]
\FunctionTok{set.seed}\NormalTok{(}\DecValTok{123}\NormalTok{)}
\NormalTok{valores }\OtherTok{\textless{}{-}} \FunctionTok{rbinom}\NormalTok{(}\DecValTok{10000}\NormalTok{, }\AttributeTok{size =} \DecValTok{3}\NormalTok{, }\AttributeTok{prob =} \FloatTok{0.5}\NormalTok{)}
\FunctionTok{sum}\NormalTok{(valores }\SpecialCharTok{==} \DecValTok{1}\NormalTok{) }\SpecialCharTok{/} \FunctionTok{length}\NormalTok{(valores)   }\CommentTok{\# P(X = 1)}
\DocumentationTok{\#\# [1] 0.383}
\FunctionTok{mean}\NormalTok{(valores)                         }\CommentTok{\# E(X)}
\DocumentationTok{\#\# [1] 1.4897}
\end{Highlighting}
\end{Shaded}

E, para a moeda enviesada, com \(X \sim \text{Binomial}(3, 3/4)\):

\begin{Shaded}
\begin{Highlighting}[]
\FunctionTok{set.seed}\NormalTok{(}\DecValTok{123}\NormalTok{)}
\NormalTok{valores }\OtherTok{\textless{}{-}} \FunctionTok{rbinom}\NormalTok{(}\DecValTok{10000}\NormalTok{, }\AttributeTok{size =} \DecValTok{3}\NormalTok{, }\AttributeTok{prob =} \DecValTok{3}\SpecialCharTok{/}\DecValTok{4}\NormalTok{)}
\FunctionTok{sum}\NormalTok{(valores }\SpecialCharTok{==} \DecValTok{1}\NormalTok{) }\SpecialCharTok{/} \FunctionTok{length}\NormalTok{(valores)}
\DocumentationTok{\#\# [1] 0.1365}
\FunctionTok{mean}\NormalTok{(valores)}
\DocumentationTok{\#\# [1] 2.2558}
\end{Highlighting}
\end{Shaded}

\textbf{Exemplo 2 (tempo até ao autocarro).} O tempo até à chegada de um autocarro segue uma
distribuição exponencial de média 30 minutos. Estimemos por simulação a probabilidade
de o autocarro chegar nos primeiros 20 minutos e comparemos com o valor exato dado por
\texttt{pexp()}:

\begin{Shaded}
\begin{Highlighting}[]
\FunctionTok{set.seed}\NormalTok{(}\DecValTok{123}\NormalTok{)}
\NormalTok{valores }\OtherTok{\textless{}{-}} \FunctionTok{rexp}\NormalTok{(}\DecValTok{10000}\NormalTok{, }\AttributeTok{rate =} \DecValTok{1}\SpecialCharTok{/}\DecValTok{30}\NormalTok{)}
\FunctionTok{sum}\NormalTok{(valores }\SpecialCharTok{\textless{}} \DecValTok{20}\NormalTok{) }\SpecialCharTok{/} \FunctionTok{length}\NormalTok{(valores)   }\CommentTok{\# estimativa de P(X \textless{} 20)}
\DocumentationTok{\#\# [1] 0.4832}
\FunctionTok{pexp}\NormalTok{(}\DecValTok{20}\NormalTok{, }\AttributeTok{rate =} \DecValTok{1}\SpecialCharTok{/}\DecValTok{30}\NormalTok{)                  }\CommentTok{\# valor exato}
\DocumentationTok{\#\# [1] 0.4865829}
\end{Highlighting}
\end{Shaded}

\textbf{Exemplo 3 (cartas até ao primeiro ás).} Retiram-se cartas de um baralho, com
reposição, até sair um ás. Simulemos a média e a variância do número de cartas
necessárias:

\begin{Shaded}
\begin{Highlighting}[]
\FunctionTok{set.seed}\NormalTok{(}\DecValTok{123}\NormalTok{)}
\NormalTok{n }\OtherTok{\textless{}{-}} \DecValTok{10000}
\NormalTok{simlist }\OtherTok{\textless{}{-}} \FunctionTok{numeric}\NormalTok{(n)}
\ControlFlowTok{for}\NormalTok{ (i }\ControlFlowTok{in} \DecValTok{1}\SpecialCharTok{:}\NormalTok{n) \{}
\NormalTok{  ct }\OtherTok{\textless{}{-}} \DecValTok{0}
\NormalTok{  as }\OtherTok{\textless{}{-}} \DecValTok{0}
  \ControlFlowTok{while}\NormalTok{ (as }\SpecialCharTok{==} \DecValTok{0}\NormalTok{) \{}
\NormalTok{    carta }\OtherTok{\textless{}{-}} \FunctionTok{sample}\NormalTok{(}\DecValTok{1}\SpecialCharTok{:}\DecValTok{52}\NormalTok{, }\DecValTok{1}\NormalTok{, }\AttributeTok{replace =} \ConstantTok{TRUE}\NormalTok{)   }\CommentTok{\# ases: 1, 2, 3, 4}
\NormalTok{    ct }\OtherTok{\textless{}{-}}\NormalTok{ ct }\SpecialCharTok{+} \DecValTok{1}
    \ControlFlowTok{if}\NormalTok{ (carta }\SpecialCharTok{\textless{}=} \DecValTok{4}\NormalTok{) as }\OtherTok{\textless{}{-}} \DecValTok{1}
\NormalTok{  \}}
\NormalTok{  simlist[i] }\OtherTok{\textless{}{-}}\NormalTok{ ct}
\NormalTok{\}}
\FunctionTok{mean}\NormalTok{(simlist)}
\DocumentationTok{\#\# [1] 12.8081}
\FunctionTok{var}\NormalTok{(simlist)}
\DocumentationTok{\#\# [1] 147.5318}
\end{Highlighting}
\end{Shaded}

O número de provas de Bernoulli até ao primeiro sucesso (aparecer um ás, com
\(p = 4/52\)) segue uma distribuição \textbf{geométrica}. Como veremos, o R conta os
\emph{insucessos} até ao primeiro sucesso, pelo que somamos 1 para obter o número de cartas:

\begin{Shaded}
\begin{Highlighting}[]
\FunctionTok{set.seed}\NormalTok{(}\DecValTok{123}\NormalTok{)}
\NormalTok{valores }\OtherTok{\textless{}{-}} \FunctionTok{rgeom}\NormalTok{(}\DecValTok{10000}\NormalTok{, }\AttributeTok{prob =} \DecValTok{4}\SpecialCharTok{/}\DecValTok{52}\NormalTok{) }\SpecialCharTok{+} \DecValTok{1}
\FunctionTok{mean}\NormalTok{(valores)}
\DocumentationTok{\#\# [1] 13.0108}
\FunctionTok{var}\NormalTok{(valores)}
\DocumentationTok{\#\# [1] 152.0335}
\end{Highlighting}
\end{Shaded}

\section{Função de distribuição empírica}\label{funuxe7uxe3o-de-distribuiuxe7uxe3o-empuxedrica}

A \textbf{função de distribuição empírica} é a função de distribuição construída a partir
dos \emph{dados observados}: para cada valor \(x\), dá a proporção de observações da amostra
que são menores ou iguais a \(x\). É uma ferramenta natural para visualizar uma
distribuição amostral e compará-la com um modelo teórico.

Formalmente, para uma amostra de dimensão \(n\),
\[F_n(x) = \frac{1}{n} \sum_{i=1}^n \mathbb{I}(X_i \leq x),\]
onde \(\mathbb{I}(X_i \leq x)\) é a função indicadora, que vale 1 se \(X_i \leq x\) e 0 caso
contrário. Ordenando os dados por ordem crescente, \(x_{(1)} \leq \cdots \leq x_{(n)}\), a
função é uma escada que vale 0 antes de \(x_{(1)}\), sobe \(\frac{1}{n}\) em cada
observação (ou \(\frac{k}{n}\) se um valor se repetir \(k\) vezes) e chega a 1 em \(x_{(n)}\).

\subsection{\texorpdfstring{A função \texttt{ecdf()}}{A função ecdf()}}\label{a-funuxe7uxe3o-ecdf}

No R, a função de distribuição empírica calcula-se com \texttt{ecdf()}:

\begin{Shaded}
\begin{Highlighting}[]
\NormalTok{dados }\OtherTok{\textless{}{-}} \FunctionTok{c}\NormalTok{(}\DecValTok{3}\NormalTok{, }\DecValTok{1}\NormalTok{, }\DecValTok{4}\NormalTok{, }\DecValTok{1}\NormalTok{, }\DecValTok{5}\NormalTok{, }\DecValTok{9}\NormalTok{, }\DecValTok{2}\NormalTok{, }\DecValTok{6}\NormalTok{, }\DecValTok{5}\NormalTok{, }\DecValTok{3}\NormalTok{, }\DecValTok{5}\NormalTok{)}

\NormalTok{Fn }\OtherTok{\textless{}{-}} \FunctionTok{ecdf}\NormalTok{(dados)                 }\CommentTok{\# cria a função de distribuição empírica}
\FunctionTok{plot}\NormalTok{(Fn, }\AttributeTok{main =} \StringTok{"Função de distribuição empírica"}\NormalTok{,}
     \AttributeTok{xlab =} \StringTok{"x"}\NormalTok{, }\AttributeTok{ylab =} \StringTok{"Fn(x)"}\NormalTok{, }\AttributeTok{col =} \StringTok{"blue"}\NormalTok{, }\AttributeTok{lwd =} \DecValTok{2}\NormalTok{)}
\end{Highlighting}
\end{Shaded}

\begin{figure}

{\centering \includegraphics[width=0.8\linewidth]{introR_files/figure-latex/ecdf-exemplo-1} 

}

\caption{Função de distribuição empírica de um pequeno conjunto de dados.}\label{fig:ecdf-exemplo}
\end{figure}

O objeto devolvido por \texttt{ecdf()} é, ele próprio, uma função: \texttt{Fn(20)} devolve a proporção
de observações \(\leq 20\). Retomando o Exemplo 2, podemos estimar \(P(X \leq 20)\) e
compará-la com o valor exato:

\begin{Shaded}
\begin{Highlighting}[]
\FunctionTok{set.seed}\NormalTok{(}\DecValTok{123}\NormalTok{)}
\NormalTok{valores }\OtherTok{\textless{}{-}} \FunctionTok{rexp}\NormalTok{(}\DecValTok{10000}\NormalTok{, }\AttributeTok{rate =} \DecValTok{1}\SpecialCharTok{/}\DecValTok{30}\NormalTok{)}
\NormalTok{Fn }\OtherTok{\textless{}{-}} \FunctionTok{ecdf}\NormalTok{(valores)}
\FunctionTok{Fn}\NormalTok{(}\DecValTok{20}\NormalTok{)                }\CommentTok{\# estimativa de P(X \textless{}= 20)}
\DocumentationTok{\#\# [1] 0.4832}
\FunctionTok{pexp}\NormalTok{(}\DecValTok{20}\NormalTok{, }\AttributeTok{rate =} \DecValTok{1}\SpecialCharTok{/}\DecValTok{30}\NormalTok{) }\CommentTok{\# valor exato}
\DocumentationTok{\#\# [1] 0.4865829}
\end{Highlighting}
\end{Shaded}

\section{Modelos discretos}\label{modelos-discretos}

Um modelo probabilístico discreto descreve fenómenos em que a variável aleatória assume
apenas valores isolados --- num conjunto finito ou numerável --- e especifica a
probabilidade de cada um. Vejamos os principais.

\section{Distribuição uniforme discreta}\label{distribuiuxe7uxe3o-uniforme-discreta}

\textbf{Definição.} A variável aleatória \(X\) tem \textbf{distribuição uniforme discreta} no
conjunto \(\{x_1, x_2, \ldots, x_n\}\) se todos os valores forem igualmente prováveis:
\[P(X = x) = \begin{cases} \frac{1}{n}, & x \in \{x_1, \ldots, x_n\} \\ 0, & \text{caso contrário} \end{cases}\]

\textbf{Notação e momentos.} \(X \sim \text{Uniforme Discreta}(\{x_1, \ldots, x_n\})\), com
\[E(X) = \frac{1}{n}\sum_{i=1}^n x_i, \qquad V(X) = \frac{1}{n}\sum_{i=1}^n x_i^2 - \left(\frac{1}{n}\sum_{i=1}^n x_i\right)^2.\]

No caso particular, muito comum, em que os valores são os inteiros consecutivos
\(\{1, 2, \ldots, n\}\), estas fórmulas simplificam-se para
\[E(X) = \frac{n+1}{2}, \qquad V(X) = \frac{n^2 - 1}{12}.\]
Atenção, porém: estas versões simplificadas só valem para \(\{1, \ldots, n\}\). Para um
conjunto qualquer de valores, há que usar as fórmulas gerais acima.

O R não tem uma função dedicada a esta distribuição --- dada a sua simplicidade ---, mas a
\texttt{sample()} gera-a diretamente. Por exemplo, 15 valores de uma uniforme discreta em
\(\{1, \ldots, 10\}\):

\begin{Shaded}
\begin{Highlighting}[]
\FunctionTok{sample}\NormalTok{(}\DecValTok{1}\SpecialCharTok{:}\DecValTok{10}\NormalTok{, }\AttributeTok{size =} \DecValTok{15}\NormalTok{, }\AttributeTok{replace =} \ConstantTok{TRUE}\NormalTok{)}
\DocumentationTok{\#\# [1] 2 3 4 4 4 5 7 9 4 2 6 7 1 6 9}
\end{Highlighting}
\end{Shaded}

\subsection{Exercícios}\label{exercuxedcios-25}

\textbf{1.} Crie uma variável aleatória uniforme discreta \(X\) com valores de 1 a 10.

\begin{itemize}
\tightlist
\item
  Simule 1000 realizações de \(X\).
\item
  Verifique se a frequência relativa de cada valor se aproxima da probabilidade teórica
  \(P(X = x) = \frac{1}{10}\).
\item
  Represente os resultados num gráfico de barras.
\end{itemize}

\textbf{2.} Defina \(X\) uniforme discreta com valores de 5 a 15.

\begin{itemize}
\tightlist
\item
  Calcule \(E(X)\) e \(V(X)\).
\item
  Simule 10 000 realizações e compare a média e a variância empíricas com as teóricas.
\end{itemize}

\textbf{3.} Considere duas variáveis uniformes discretas independentes \(X\) e \(Y\), com valores
de 1 a 6 (um par de dados).

\begin{itemize}
\tightlist
\item
  Gere 5000 pares \((X, Y)\).
\item
  Calcule a média e a variância da soma \(Z = X + Y\).
\item
  Estime \(P(Z > 8)\).
\end{itemize}

\textbf{4.} Defina \(X\) uniforme discreta com valores de \(-3\) a 3.

\begin{itemize}
\tightlist
\item
  Crie a variável \(Y = X^2 + 2X\).
\item
  Simule 5000 valores de \(X\) e calcule \(E(Y)\) e \(V(Y)\).
\end{itemize}

\textbf{5.} Defina \(X\) uniforme discreta com valores de \(-5\) a 5.

\begin{itemize}
\tightlist
\item
  Crie \(Y = 3X^3 - 2X^2 + X\).
\item
  Simule 5000 valores de \(X\) e calcule a média empírica de \(Y\) e a proporção de valores
  de \(Y\) positivos.
\end{itemize}

\section{Distribuição de Bernoulli}\label{distribuiuxe7uxe3o-de-bernoulli}

\textbf{Prova de Bernoulli.} Uma experiência diz-se uma \textbf{prova de Bernoulli} se tiver
apenas dois resultados: um sucesso, com probabilidade \(p\) (\(0 \leq p \leq 1\)), e um
insucesso, com probabilidade \(1 - p\). Lançar uma moeda, verificar se uma peça é
defeituosa ou se uma perfuração encontra petróleo são exemplos típicos.

\textbf{Definição.} A variável \(X\) que conta o número de sucessos numa prova de Bernoulli
tem \textbf{distribuição de Bernoulli} de parâmetro \(p\), com f.m.p.
\[P(X = x) = p^x (1-p)^{1-x}, \quad x \in \{0, 1\}.\]

\textbf{Notação e momentos.} \(X \sim \text{Bernoulli}(p)\), com \(E(X) = p\) e \(V(X) = p(1-p)\).

\subsection{Cálculo de probabilidades}\label{cuxe1lculo-de-probabilidades}

A Bernoulli é o caso \(n = 1\) da binomial, pelo que se usam as funções \texttt{*binom} com
\texttt{size\ =\ 1}. Seja \(X \sim \text{Bernoulli}(p = 0{,}5)\):

\begin{itemize}
\tightlist
\item
  \(P(X = 0) \to\) \texttt{dbinom(0,\ size\ =\ 1,\ prob\ =\ 0.5)} \(= 0{,}5\)
\item
  \(P(X = 1) \to\) \texttt{dbinom(1,\ size\ =\ 1,\ prob\ =\ 0.5)} \(= 0{,}5\)
\item
  \(P(X \leq 1) \to\) \texttt{pbinom(1,\ size\ =\ 1,\ prob\ =\ 0.5)} \(= 1\)
\item
  \(P(X > 0) \to\) \texttt{pbinom(0,\ size\ =\ 1,\ prob\ =\ 0.5,\ lower.tail\ =\ FALSE)} \(= 0{,}5\)
\item
  amostra de dimensão 5 \(\to\) \texttt{rbinom(5,\ size\ =\ 1,\ prob\ =\ 0.5)}
\end{itemize}

\subsection{Exercícios}\label{exercuxedcios-26}

\textbf{1.} Considere o lançamento de uma moeda equilibrada e a variável \$X = \$ ``número de
caras em 1 lançamento''.

\begin{enumerate}
\def\labelenumi{(\alph{enumi})}
\tightlist
\item
  Simule a experiência, determinando a percentagem de caras para \(n_1 = 5\),
  \(n_2 = 10\), \(n_3 = 100\) e \(n_4 = 1000\) lançamentos.
\item
  Para cada amostra, calcule a média e a variância e compare com \(E(X)\) e \(V(X)\).
\end{enumerate}

\textbf{2.} Seja \(X \sim \text{Bernoulli}\) com \(P(X = 1) = 0{,}7\).

\begin{itemize}
\tightlist
\item
  Simule 1000 valores de \(X\).
\item
  Calcule a frequência relativa de \(X = 1\) e \(X = 0\) e compare com as probabilidades
  teóricas.
\end{itemize}

\textbf{3.} Defina \(X \sim \text{Bernoulli}\) com \(P(X = 1) = 0{,}4\).

\begin{itemize}
\tightlist
\item
  Simule 10 000 valores.
\item
  Compare a média empírica com \(E(X) = p\) e a variância empírica com \(V(X) = p(1-p)\).
\end{itemize}

\textbf{4.} Considere 5 variáveis \(X_1, \ldots, X_5\), cada uma Bernoulli com \(p = 0{,}5\).

\begin{itemize}
\tightlist
\item
  Simule 5000 realizações de cada.
\item
  Calcule a soma \(S = X_1 + \cdots + X_5\).
\item
  Compare a frequência relativa de cada valor de \(S\) (de 0 a 5) com a binomial teórica,
  \texttt{dbinom(0:5,\ size\ =\ 5,\ prob\ =\ 0.5)}.
\end{itemize}

\section{Distribuição binomial}\label{distribuiuxe7uxe3o-binomial}

\textbf{Definição.} A variável \(X\) que conta o número de sucessos em \(n\) provas de Bernoulli
independentes, todas com probabilidade de sucesso \(p\), tem \textbf{distribuição binomial} de
parâmetros \((n, p)\), com f.m.p.
\[P(X = x) = \binom{n}{x} p^x (1-p)^{n-x}, \quad x = 0, 1, \ldots, n,\]
onde \(\binom{n}{x} = \dfrac{n!}{x!\,(n-x)!}\) é o número de combinações (em R, \texttt{choose(n,\ x)}).

\textbf{Notação e momentos.} \(X \sim \text{Binomial}(n, p)\), com \(E(X) = np\) e
\(V(X) = np(1-p)\).

\subsection{Cálculo de probabilidades}\label{cuxe1lculo-de-probabilidades-1}

Seja \(X \sim \text{Binomial}(n = 20, p = 0{,}1)\):

\begin{itemize}
\tightlist
\item
  \(P(X = 4) \to\) \texttt{dbinom(4,\ size\ =\ 20,\ prob\ =\ 0.1)} \(= 0{,}0898\)
\item
  \(P(X \leq 4) \to\) \texttt{pbinom(4,\ size\ =\ 20,\ prob\ =\ 0.1)} \(= 0{,}9568\)
\item
  \(P(X > 4) \to\) \texttt{pbinom(4,\ size\ =\ 20,\ prob\ =\ 0.1,\ lower.tail\ =\ FALSE)} \(= 0{,}0432\)
\item
  amostra de dimensão 5 \(\to\) \texttt{rbinom(5,\ size\ =\ 20,\ prob\ =\ 0.1)}
\end{itemize}

\subsection{Função massa de probabilidade (teórica)}\label{funuxe7uxe3o-massa-de-probabilidade-teuxf3rica}

\begin{Shaded}
\begin{Highlighting}[]
\NormalTok{n }\OtherTok{\textless{}{-}} \DecValTok{20}\NormalTok{; p }\OtherTok{\textless{}{-}} \FloatTok{0.1}
\NormalTok{x }\OtherTok{\textless{}{-}} \DecValTok{0}\SpecialCharTok{:}\DecValTok{20}
\NormalTok{teorico }\OtherTok{\textless{}{-}} \FunctionTok{data.frame}\NormalTok{(}\AttributeTok{x =}\NormalTok{ x, }\AttributeTok{y =} \FunctionTok{dbinom}\NormalTok{(x, }\AttributeTok{size =}\NormalTok{ n, }\AttributeTok{prob =}\NormalTok{ p))}

\FunctionTok{plot}\NormalTok{(teorico}\SpecialCharTok{$}\NormalTok{x, teorico}\SpecialCharTok{$}\NormalTok{y,}
     \AttributeTok{main =} \StringTok{"Binomial(n = 20, p = 0.1)"}\NormalTok{,}
     \AttributeTok{xlab =} \StringTok{"Número de sucessos"}\NormalTok{, }\AttributeTok{ylab =} \StringTok{"Probabilidade"}\NormalTok{,}
     \AttributeTok{pch =} \DecValTok{19}\NormalTok{, }\AttributeTok{col =} \StringTok{"blue"}\NormalTok{)}
\FunctionTok{grid}\NormalTok{(}\AttributeTok{nx =} \DecValTok{21}\NormalTok{, }\AttributeTok{ny =} \ConstantTok{NULL}\NormalTok{)}
\end{Highlighting}
\end{Shaded}

\begin{figure}

{\centering \includegraphics[width=0.8\linewidth]{introR_files/figure-latex/bin-teorica-1} 

}

\caption{Função massa de probabilidade da Binomial(20, 0.1).}\label{fig:bin-teorica}
\end{figure}

\subsection{Função massa de probabilidade (simulação)}\label{funuxe7uxe3o-massa-de-probabilidade-simulauxe7uxe3o}

\begin{Shaded}
\begin{Highlighting}[]
\FunctionTok{set.seed}\NormalTok{(}\DecValTok{1234}\NormalTok{)}
\NormalTok{n }\OtherTok{\textless{}{-}} \DecValTok{20}\NormalTok{; p }\OtherTok{\textless{}{-}} \FloatTok{0.1}\NormalTok{; k }\OtherTok{\textless{}{-}} \DecValTok{1000}            \CommentTok{\# k = número de simulações}

\NormalTok{dados }\OtherTok{\textless{}{-}} \FunctionTok{rbinom}\NormalTok{(k, }\AttributeTok{size =}\NormalTok{ n, }\AttributeTok{prob =}\NormalTok{ p)}
\NormalTok{frequencia\_relativa }\OtherTok{\textless{}{-}} \FunctionTok{table}\NormalTok{(dados) }\SpecialCharTok{/} \FunctionTok{length}\NormalTok{(dados)}

\FunctionTok{barplot}\NormalTok{(frequencia\_relativa,}
        \AttributeTok{main =} \StringTok{"Amostra da Binomial(20, 0.1)"}\NormalTok{,}
        \AttributeTok{col =} \StringTok{"lightblue"}\NormalTok{,}
        \AttributeTok{xlab =} \StringTok{"Número de sucessos"}\NormalTok{, }\AttributeTok{ylab =} \StringTok{"Frequência relativa"}\NormalTok{,}
        \AttributeTok{ylim =} \FunctionTok{c}\NormalTok{(}\DecValTok{0}\NormalTok{, }\FloatTok{0.3}\NormalTok{))}
\FunctionTok{grid}\NormalTok{()}
\end{Highlighting}
\end{Shaded}

\begin{figure}

{\centering \includegraphics[width=0.8\linewidth]{introR_files/figure-latex/bin-sim-1} 

}

\caption{Frequências relativas de uma amostra de 1000 valores da Binomial(20, 0.1).}\label{fig:bin-sim}
\end{figure}

\subsection{Comparação entre teoria e simulação}\label{comparauxe7uxe3o-entre-teoria-e-simulauxe7uxe3o}

\begin{Shaded}
\begin{Highlighting}[]
\FunctionTok{set.seed}\NormalTok{(}\DecValTok{1234}\NormalTok{)}
\NormalTok{n }\OtherTok{\textless{}{-}} \DecValTok{20}\NormalTok{; p }\OtherTok{\textless{}{-}} \FloatTok{0.1}\NormalTok{; k }\OtherTok{\textless{}{-}} \DecValTok{1000}

\NormalTok{dados }\OtherTok{\textless{}{-}} \FunctionTok{rbinom}\NormalTok{(k, }\AttributeTok{size =}\NormalTok{ n, }\AttributeTok{prob =}\NormalTok{ p)}
\NormalTok{frequencia\_relativa }\OtherTok{\textless{}{-}} \FunctionTok{table}\NormalTok{(}\FunctionTok{factor}\NormalTok{(dados, }\AttributeTok{levels =} \DecValTok{0}\SpecialCharTok{:}\NormalTok{n)) }\SpecialCharTok{/} \FunctionTok{length}\NormalTok{(dados)}
\NormalTok{teorico }\OtherTok{\textless{}{-}} \FunctionTok{data.frame}\NormalTok{(}\AttributeTok{x =} \DecValTok{0}\SpecialCharTok{:}\NormalTok{n, }\AttributeTok{y =} \FunctionTok{dbinom}\NormalTok{(}\DecValTok{0}\SpecialCharTok{:}\NormalTok{n, }\AttributeTok{size =}\NormalTok{ n, }\AttributeTok{prob =}\NormalTok{ p))}

\NormalTok{posicoes\_barra }\OtherTok{\textless{}{-}} \FunctionTok{barplot}\NormalTok{(frequencia\_relativa,}
        \AttributeTok{main =} \StringTok{"Amostra da Binomial(20, 0.1)"}\NormalTok{,}
        \AttributeTok{col =} \StringTok{"lightblue"}\NormalTok{,}
        \AttributeTok{xlab =} \StringTok{"Número de sucessos"}\NormalTok{, }\AttributeTok{ylab =} \StringTok{"Frequência relativa"}\NormalTok{,}
        \AttributeTok{ylim =} \FunctionTok{c}\NormalTok{(}\SpecialCharTok{{-}}\FloatTok{0.01}\NormalTok{, }\FloatTok{0.3}\NormalTok{))}
\FunctionTok{points}\NormalTok{(posicoes\_barra, teorico}\SpecialCharTok{$}\NormalTok{y, }\AttributeTok{col =} \StringTok{"magenta"}\NormalTok{, }\AttributeTok{pch =} \DecValTok{19}\NormalTok{)}
\FunctionTok{grid}\NormalTok{()}
\end{Highlighting}
\end{Shaded}

\begin{figure}

{\centering \includegraphics[width=0.8\linewidth]{introR_files/figure-latex/bin-comp-1} 

}

\caption{Frequências relativas simuladas (barras) e probabilidades teóricas (pontos).}\label{fig:bin-comp}
\end{figure}

\subsection{Função de distribuição}\label{funuxe7uxe3o-de-distribuiuxe7uxe3o}

\begin{Shaded}
\begin{Highlighting}[]
\NormalTok{n }\OtherTok{\textless{}{-}} \DecValTok{20}\NormalTok{; p }\OtherTok{\textless{}{-}} \FloatTok{0.1}
\NormalTok{x }\OtherTok{\textless{}{-}} \DecValTok{0}\SpecialCharTok{:}\NormalTok{n}
\NormalTok{cdf\_values }\OtherTok{\textless{}{-}} \FunctionTok{pbinom}\NormalTok{(x, }\AttributeTok{size =}\NormalTok{ n, }\AttributeTok{prob =}\NormalTok{ p)}

\FunctionTok{plot}\NormalTok{(x, cdf\_values, }\AttributeTok{type =} \StringTok{"s"}\NormalTok{, }\AttributeTok{lwd =} \DecValTok{2}\NormalTok{, }\AttributeTok{col =} \StringTok{"blue"}\NormalTok{,}
     \AttributeTok{xlab =} \StringTok{"Número de sucessos"}\NormalTok{, }\AttributeTok{ylab =} \StringTok{"F(x)"}\NormalTok{,}
     \AttributeTok{main =} \StringTok{"Função de distribuição da Binomial(20, 0.1)"}\NormalTok{)}
\end{Highlighting}
\end{Shaded}

\begin{figure}

{\centering \includegraphics[width=0.8\linewidth]{introR_files/figure-latex/bin-fd-1} 

}

\caption{Função de distribuição da Binomial(20, 0.1).}\label{fig:bin-fd}
\end{figure}

\subsection{Função de distribuição empírica}\label{funuxe7uxe3o-de-distribuiuxe7uxe3o-empuxedrica-1}

\begin{Shaded}
\begin{Highlighting}[]
\NormalTok{n }\OtherTok{\textless{}{-}} \DecValTok{20}\NormalTok{; p }\OtherTok{\textless{}{-}} \FloatTok{0.1}
\FunctionTok{set.seed}\NormalTok{(}\DecValTok{1234}\NormalTok{)}
\NormalTok{amostra }\OtherTok{\textless{}{-}} \FunctionTok{rbinom}\NormalTok{(}\DecValTok{1000}\NormalTok{, }\AttributeTok{size =}\NormalTok{ n, }\AttributeTok{prob =}\NormalTok{ p)}

\NormalTok{Fn }\OtherTok{\textless{}{-}} \FunctionTok{ecdf}\NormalTok{(amostra)}
\FunctionTok{plot}\NormalTok{(Fn, }\AttributeTok{main =} \StringTok{"Função de distribuição empírica"}\NormalTok{,}
     \AttributeTok{xlab =} \StringTok{"x"}\NormalTok{, }\AttributeTok{ylab =} \StringTok{"Fn(x)"}\NormalTok{, }\AttributeTok{col =} \StringTok{"blue"}\NormalTok{)}
\end{Highlighting}
\end{Shaded}

\begin{figure}

{\centering \includegraphics[width=0.8\linewidth]{introR_files/figure-latex/bin-ecdf-1} 

}

\caption{Função de distribuição empírica de uma amostra da Binomial(20, 0.1).}\label{fig:bin-ecdf}
\end{figure}

Assim, para \(X \sim \text{Binomial}(20, 0{,}1)\), o valor teórico
\(P(X \leq 4) =\) \texttt{pbinom(4,\ 20,\ 0.1)} \(= 0{,}9568\) é bem aproximado pela estimativa
empírica \texttt{Fn(4)} \(\approx 0{,}956\).

\subsection{Exercícios}\label{exercuxedcios-27}

\textbf{1.} Considere 10 lançamentos de uma moeda equilibrada (\(p = 0{,}5\)).

\begin{enumerate}
\def\labelenumi{(\alph{enumi})}
\tightlist
\item
  Calcule a probabilidade de obter exatamente 6 caras.
\item
  Calcule a probabilidade de obter no máximo 4 caras.
\item
  Calcule a probabilidade de obter mais de 7 caras.
\end{enumerate}

\textbf{2.} Um dado equilibrado é lançado 12 vezes; o sucesso é ``sair 6'' (\(p = 1/6\)).

\begin{enumerate}
\def\labelenumi{(\alph{enumi})}
\tightlist
\item
  Calcule a probabilidade de, no máximo, cinco 6.
\item
  Gere uma amostra de 1000 experiências e registe o número de sucessos em cada uma.
\item
  Faça um gráfico de barras das frequências relativas e sobreponha a distribuição
  teórica de \(X\).
\item
  Use a função de distribuição empírica para estimar a probabilidade da alínea (a) e
  compare com o valor teórico.
\item
  Calcule a média e a variância da amostra.
\item
  Compare com os valores teóricos \(E(X) = np\) e \(V(X) = np(1-p)\).
\end{enumerate}

\textbf{3.} Considere \$X = \$ ``número de caras em 7 lançamentos'' de uma moeda equilibrada.

\begin{enumerate}
\def\labelenumi{(\alph{enumi})}
\tightlist
\item
  Calcule a probabilidade de sair 2 vezes cara.
\item
  Simule a experiência para \(n_1 = 5\), \(n_2 = 10\), \(n_3 = 100\) e \(n_4 = 1000\)
  repetições; para cada caso, determine a percentagem de vezes em que saíram 3 caras.
\item
  Para cada amostra, calcule a média e a variância e compare com \(E(X)\) e \(V(X)\).
\end{enumerate}

\textbf{4.} Um teste de escolha múltipla tem 10 questões, cada uma com 4 alternativas e
apenas uma correta (\(p = 0{,}25\)).

\begin{enumerate}
\def\labelenumi{(\alph{enumi})}
\tightlist
\item
  Calcule a probabilidade de acertar exatamente \(k\) questões, para \(k = 0, \ldots, 10\).
\item
  Represente graficamente a distribuição de probabilidades.
\item
  Identifique o valor de \(k\) mais provável.
\end{enumerate}

\textbf{5.} Uma equipa de basquetebol acerta cada lance livre com probabilidade 0,6 e tenta
15 lances num jogo.

\begin{enumerate}
\def\labelenumi{(\alph{enumi})}
\tightlist
\item
  Calcule a probabilidade de acertar exatamente 9.
\item
  Calcule a probabilidade de acertar entre 8 e 12 (inclusive).
\item
  Gere 500 jogos e estime a proporção em que a equipa acerta entre 8 e 12.
\end{enumerate}

\textbf{6.} Num processo de fabrico, \(X\) é o número de peças defeituosas num lote de 40, com
\(p = 0{,}05\) de uma peça ser defeituosa.

\begin{enumerate}
\def\labelenumi{(\alph{enumi})}
\tightlist
\item
  Com \texttt{set.seed(123)}, gere 10 000 observações de \(X\).
\item
  Conte os lotes com exatamente 2 peças defeituosas.
\item
  Compare a proporção correspondente com \(P(X = 2)\) para \(X \sim \text{Binomial}(40, 0{,}05)\).
\end{enumerate}

\textbf{7.} Numa loja, cada cliente compra com probabilidade \(p = 0{,}3\); entram 25 clientes.
Seja \(X\) o número de compradores.

\begin{enumerate}
\def\labelenumi{(\alph{enumi})}
\tightlist
\item
  Com \texttt{set.seed(456)}, gere 5000 observações de \(X\).
\item
  Conte os períodos com pelo menos 10 compradores.
\item
  Compare a proporção com \(P(X \geq 10)\) para \(X \sim \text{Binomial}(25, 0{,}3)\).
\item
  Use a função de distribuição empírica para estimar essa probabilidade e compare.
\item
  Calcule o número médio de compradores na amostra.
\item
  Compare a média amostral com \(E(X)\).
\end{enumerate}

\textbf{8.} O número de acertos num alvo em 30 tentativas, com probabilidade de acerto 0,4, é
\(X \sim \text{Binomial}(30, 0{,}4)\). Com \texttt{set.seed(123)}, gere uma amostra de dimensão
700.

\begin{enumerate}
\def\labelenumi{(\alph{enumi})}
\tightlist
\item
  Faça um gráfico de barras das frequências relativas e sobreponha a distribuição
  teórica de \(X\).
\item
  Com a função de distribuição empírica, estime \(P(X > 15)\) e calcule também o valor
  teórico.
\end{enumerate}

\section{Distribuição geométrica}\label{distribuiuxe7uxe3o-geomuxe9trica}

\textbf{Definição.} A variável \$X = \$ ``número de provas de Bernoulli independentes (com
probabilidade de sucesso \(p\)) realizadas até ao primeiro sucesso'' tem \textbf{distribuição
geométrica} de parâmetro \(p\), com f.m.p.
\[P(X = x) = p(1-p)^{x-1}, \quad x = 1, 2, 3, \ldots\]

\textbf{Notação e momentos.} \(X \sim \text{Geométrica}(p)\), com \(E(X) = \dfrac{1}{p}\) e
\(V(X) = \dfrac{1-p}{p^2}\).

Há uma segunda forma, equivalente, de definir a distribuição --- e é a que o R adota.

\textbf{Definição (convenção do R).} A variável \$Y = X - 1 = \$ ``número de \emph{insucessos} antes
do primeiro sucesso'' tem f.m.p.
\[P(Y = y) = p(1-p)^{y}, \quad y = 0, 1, 2, \ldots,\]
com \(E(Y) = \dfrac{1-p}{p}\) e \(V(Y) = \dfrac{1-p}{p^2}\).

Esta é uma fonte frequente de erros. As funções \texttt{dgeom()}, \texttt{pgeom()} e \texttt{rgeom()} do R
usam a convenção de \(Y\) --- o \textbf{número de insucessos} antes do primeiro sucesso, com
suporte \(0, 1, 2, \ldots\) --- e não o número de provas \(X\), com suporte \(1, 2, \ldots\). A
relação é simples, \(X = Y + 1\): para obter o número de provas, some 1 ao resultado de
\texttt{rgeom()}; e para calcular \(P(X = k)\), use \texttt{dgeom(k\ -\ 1,\ p)}.

\subsection{Cálculo de probabilidades}\label{cuxe1lculo-de-probabilidades-2}

Nas funções do R (convenção \(Y\), número de insucessos), seja \(Y \sim \text{Geométrica}(p = 0{,}5)\):

\begin{itemize}
\tightlist
\item
  \(P(Y = 0) \to\) \texttt{dgeom(0,\ prob\ =\ 0.5)} \(= 0{,}5\)
\item
  \(P(Y = 1) \to\) \texttt{dgeom(1,\ prob\ =\ 0.5)} \(= 0{,}25\)
\item
  \(P(Y \leq 1) \to\) \texttt{pgeom(1,\ prob\ =\ 0.5)} \(= 0{,}75\)
\item
  \(P(Y > 1) \to\) \texttt{pgeom(1,\ prob\ =\ 0.5,\ lower.tail\ =\ FALSE)} \(= 0{,}25\)
\end{itemize}

\textbf{Exemplo.} Seja \(X\) o número de lançamentos de um dado equilibrado até sair a primeira
face 2. Qual a probabilidade de a face 2 sair no terceiro lançamento? Como no R a
geométrica conta insucessos, \(P(X = 3) = P(Y = 2)\):

\begin{Shaded}
\begin{Highlighting}[]
\FunctionTok{dgeom}\NormalTok{(}\DecValTok{2}\NormalTok{, }\AttributeTok{prob =} \DecValTok{1}\SpecialCharTok{/}\DecValTok{6}\NormalTok{)      }\CommentTok{\# P(X = 3) = P(Y = 2)}
\DocumentationTok{\#\# [1] 0.09645062}
\end{Highlighting}
\end{Shaded}

\subsection{Exercícios}\label{exercuxedcios-28}

\textbf{1.} Um dado equilibrado é lançado até sair um 6 (\(p = 1/6\)).

\begin{itemize}
\tightlist
\item
  Simule 1000 experiências e registe o número de tentativas em cada uma.
\item
  Compare as frequências relativas com as probabilidades teóricas (\texttt{dgeom()}, atendendo
  à convenção do R).
\end{itemize}

\textbf{2.} Com \texttt{set.seed(123)}, simule 5000 valores de uma geométrica com \(p = 0{,}2\).

\begin{itemize}
\tightlist
\item
  Estime \(P(\text{número de tentativas} \leq 5)\).
\item
  Compare com o valor teórico (\texttt{pgeom()}).
\end{itemize}

\textbf{3.} Gere a distribuição geométrica para \(p = 0{,}1\), \(p = 0{,}5\) e \(p = 0{,}9\).

\begin{itemize}
\tightlist
\item
  Faça gráficos de barras para comparar.
\item
  Descreva como \(p\) afeta a forma e o decaimento da distribuição.
\end{itemize}

\textbf{4.} Defina uma geométrica com \(p = 0{,}3\).

\begin{itemize}
\tightlist
\item
  Com \texttt{set.seed(123)}, simule 10 000 valores e calcule a média e a variância empíricas.
\item
  Compare com \(E(X) = \frac{1-p}{p}\) e \(V(X) = \frac{1-p}{p^2}\) (convenção do R).
\end{itemize}

\textbf{5.} Numa sondagem, 80\% das pessoas concordam com uma afirmação (\(p = 0{,}8\)).

\begin{itemize}
\tightlist
\item
  Simule o número de entrevistas até encontrar uma pessoa que discorde (\(1 - p = 0{,}2\)).
\item
  Em 5000 simulações, calcule a média e a variância do número de entrevistas.
\item
  Visualize a distribuição empírica.
\end{itemize}

\textbf{6.} Numa central de apoio técnico, cada tentativa resolve o problema com \(p = 0{,}3\).

\begin{enumerate}
\def\labelenumi{(\alph{enumi})}
\tightlist
\item
  Qual a probabilidade de o resolver em no máximo 5 tentativas?
\item
  Qual a probabilidade de precisar de mais de 8 tentativas?
\item
  Simule 500 casos com \texttt{rgeom()} e estime a frequência relativa de casos com \(\leq 5\)
  tentativas; compare com (a).
\item
  Calcule a média e o desvio padrão da amostra e compare com \(\mu = \frac{1-p}{p}\) e
  \(\sigma = \sqrt{\frac{1-p}{p^2}}\).
\item
  Compare graficamente a distribuição teórica com a frequência relativa observada.
\end{enumerate}

\textbf{7.} Num sistema de autenticação biométrica, cada tentativa tem \(p = 0{,}2\) de
sucesso. O tempo de processamento associado a \(X\) tentativas é \(T = 2X^2 + 3X + 5\) (ms).

\begin{enumerate}
\def\labelenumi{(\alph{enumi})}
\tightlist
\item
  Simule 1000 valores de \(X \sim \text{Geom}(0{,}2)\).
\item
  Calcule o vetor de tempos \(T\).
\item
  Faça um histograma dos tempos.
\item
  Calcule a média e o desvio padrão de \(T\).
\item
  Estime a probabilidade de um acesso demorar mais de 150 ms.
\end{enumerate}

\textbf{8.} Um medicamento elimina uma célula-alvo com \(p = 0{,}1\) por tentativa. Seja \(X\) o
número de falhas antes do sucesso; o custo de um ensaio é \(C = 0{,}5X^2 + 20\) (euros).

\begin{enumerate}
\def\labelenumi{(\alph{enumi})}
\tightlist
\item
  Simule 5000 ensaios, com \(X \sim \text{Geom}(0{,}1)\).
\item
  Calcule o vetor de custos \(C\).
\item
  Faça uma caixa de bigodes dos custos.
\item
  Calcule a média e o desvio padrão dos custos, e a percentagem de ensaios com custo
  superior a 200 € e inferior a 50 €.
\end{enumerate}

\textbf{9.} Num sistema de alta disponibilidade, cada servidor está operacional com
probabilidade 0,1; procura-se sequencialmente até encontrar um funcional. O tempo de
resposta é \(T = \sqrt{X^2 + 3X + 1} + \log(X + 1)\) (ms).

\begin{enumerate}
\def\labelenumi{(\alph{enumi})}
\tightlist
\item
  Simule 3000 valores de \(X \sim \text{Geom}(0{,}1)\).
\item
  Calcule \(T\) para cada valor.
\item
  Faça um gráfico de dispersão entre \(X\) e \(T\).
\item
  Calcule a média, a mediana e o desvio padrão de \(T\).
\item
  Estime a probabilidade de \(T > 10\) ms.
\item
  Calcule os quantis 25\%, 50\% e 75\% de \(T\).
\end{enumerate}

\section{Distribuição de Poisson}\label{distribuiuxe7uxe3o-de-poisson}

Considere-se a contagem do número de ocorrências de um acontecimento num intervalo (de
tempo, comprimento, área, volume) que verifica três propriedades: as ocorrências em
intervalos disjuntos são independentes (ausência de memória); a probabilidade de
ocorrência é a mesma para intervalos de igual amplitude; e a probabilidade de mais de
uma ocorrência num intervalo muito pequeno é desprezável. Uma tal experiência é um
\textbf{processo de Poisson}.

\textbf{Definição.} A variável \$X = \$ ``número de ocorrências por unidade de tempo ou espaço''
tem \textbf{distribuição de Poisson} de parâmetro \(\lambda > 0\), com f.m.p.
\[P(X = x) = \frac{e^{-\lambda} \lambda^x}{x!}, \quad x = 0, 1, 2, \ldots\]

\textbf{Notação e momentos.} \(X \sim \text{Poisson}(\lambda)\), com \(E(X) = V(X) = \lambda\)
(o parâmetro é o número médio de ocorrências por unidade).

\textbf{Aditividade.} Se \(X_1, \ldots, X_n\) são independentes, com
\(X_i \sim \text{Poisson}(\lambda_i)\), então
\(\sum_i X_i \sim \text{Poisson}\!\left(\sum_i \lambda_i\right)\). Esta propriedade é o que
permite, por exemplo, mudar a unidade de tempo: se ocorrem em média 3 acontecimentos por
mês, ocorrem \(3/2\) por quinzena --- a mesma família, com o parâmetro reescalado.

\subsection{Cálculo de probabilidades}\label{cuxe1lculo-de-probabilidades-3}

Seja \(X \sim \text{Poisson}(\lambda = 5)\):

\begin{itemize}
\tightlist
\item
  \(P(X = 4) \to\) \texttt{dpois(4,\ 5)} \(= 0{,}1755\)
\item
  \(P(X \leq 4) \to\) \texttt{ppois(4,\ 5)} \(= 0{,}4405\)
\item
  \(P(X > 4) \to\) \texttt{ppois(4,\ 5,\ lower.tail\ =\ FALSE)} \(= 0{,}5595\)
\end{itemize}

\subsection{Função massa de probabilidade (teórica)}\label{funuxe7uxe3o-massa-de-probabilidade-teuxf3rica-1}

O painel seguinte mostra como a forma da Poisson muda com \(\lambda\) --- assimétrica para
valores pequenos, cada vez mais simétrica à medida que \(\lambda\) cresce. (Usa os pacotes
\texttt{ggplot2}, \texttt{latex2exp} e \texttt{gridExtra}.)

\begin{Shaded}
\begin{Highlighting}[]
\FunctionTok{library}\NormalTok{(ggplot2); }\FunctionTok{library}\NormalTok{(latex2exp); }\FunctionTok{library}\NormalTok{(gridExtra)}

\NormalTok{lambdas }\OtherTok{\textless{}{-}} \FunctionTok{c}\NormalTok{(}\FloatTok{0.1}\NormalTok{, }\DecValTok{1}\NormalTok{, }\FloatTok{2.5}\NormalTok{, }\DecValTok{5}\NormalTok{, }\DecValTok{15}\NormalTok{, }\DecValTok{30}\NormalTok{)}
\NormalTok{x }\OtherTok{\textless{}{-}} \DecValTok{0}\SpecialCharTok{:}\DecValTok{50}
\NormalTok{plots }\OtherTok{\textless{}{-}} \FunctionTok{list}\NormalTok{()}
\ControlFlowTok{for}\NormalTok{ (i }\ControlFlowTok{in} \FunctionTok{seq\_along}\NormalTok{(lambdas)) \{}
\NormalTok{  teorico }\OtherTok{\textless{}{-}} \FunctionTok{data.frame}\NormalTok{(}\AttributeTok{x =}\NormalTok{ x, }\AttributeTok{y =} \FunctionTok{dpois}\NormalTok{(x, }\AttributeTok{lambda =}\NormalTok{ lambdas[i]))}
\NormalTok{  plots[[i]] }\OtherTok{\textless{}{-}} \FunctionTok{ggplot}\NormalTok{(teorico) }\SpecialCharTok{+}
    \FunctionTok{geom\_point}\NormalTok{(}\FunctionTok{aes}\NormalTok{(}\AttributeTok{x =}\NormalTok{ x, }\AttributeTok{y =}\NormalTok{ y), }\AttributeTok{color =} \StringTok{"blue"}\NormalTok{) }\SpecialCharTok{+}
    \FunctionTok{scale\_x\_continuous}\NormalTok{(}\AttributeTok{breaks =} \FunctionTok{seq}\NormalTok{(}\DecValTok{0}\NormalTok{, }\DecValTok{50}\NormalTok{, }\AttributeTok{by =} \DecValTok{10}\NormalTok{)) }\SpecialCharTok{+}
    \FunctionTok{labs}\NormalTok{(}\AttributeTok{title =} \FunctionTok{TeX}\NormalTok{(}\FunctionTok{paste0}\NormalTok{(}\StringTok{"$Poisson(}\SpecialCharTok{\textbackslash{}\textbackslash{}}\StringTok{lambda="}\NormalTok{, lambdas[i], }\StringTok{")$"}\NormalTok{)),}
         \AttributeTok{x =} \StringTok{"x"}\NormalTok{, }\AttributeTok{y =} \StringTok{"Probabilidade"}\NormalTok{) }\SpecialCharTok{+}
    \FunctionTok{theme\_light}\NormalTok{()}
\NormalTok{\}}
\FunctionTok{grid.arrange}\NormalTok{(}\AttributeTok{grobs =}\NormalTok{ plots, }\AttributeTok{nrow =} \DecValTok{2}\NormalTok{, }\AttributeTok{ncol =} \DecValTok{3}\NormalTok{)}
\end{Highlighting}
\end{Shaded}

\begin{figure}

{\centering \includegraphics[width=0.8\linewidth]{introR_files/figure-latex/pois-teorica-1} 

}

\caption{Função massa de probabilidade da Poisson para vários valores de lambda.}\label{fig:pois-teorica}
\end{figure}

\subsection{Função massa de probabilidade (simulação)}\label{funuxe7uxe3o-massa-de-probabilidade-simulauxe7uxe3o-1}

\begin{Shaded}
\begin{Highlighting}[]
\FunctionTok{set.seed}\NormalTok{(}\DecValTok{1234}\NormalTok{)}
\NormalTok{lambdas }\OtherTok{\textless{}{-}} \FunctionTok{c}\NormalTok{(}\FloatTok{0.1}\NormalTok{, }\DecValTok{1}\NormalTok{, }\FloatTok{2.5}\NormalTok{, }\DecValTok{5}\NormalTok{, }\DecValTok{15}\NormalTok{, }\DecValTok{30}\NormalTok{)}
\NormalTok{n }\OtherTok{\textless{}{-}} \DecValTok{1000}
\NormalTok{plots }\OtherTok{\textless{}{-}} \FunctionTok{list}\NormalTok{()}
\ControlFlowTok{for}\NormalTok{ (i }\ControlFlowTok{in} \FunctionTok{seq\_along}\NormalTok{(lambdas)) \{}
\NormalTok{  dados }\OtherTok{\textless{}{-}} \FunctionTok{data.frame}\NormalTok{(}\AttributeTok{X =} \FunctionTok{rpois}\NormalTok{(n, }\AttributeTok{lambda =}\NormalTok{ lambdas[i]))}
\NormalTok{  plots[[i]] }\OtherTok{\textless{}{-}} \FunctionTok{ggplot}\NormalTok{(dados) }\SpecialCharTok{+}
    \FunctionTok{geom\_bar}\NormalTok{(}\FunctionTok{aes}\NormalTok{(}\AttributeTok{x =}\NormalTok{ X, }\AttributeTok{y =} \FunctionTok{after\_stat}\NormalTok{(prop)), }\AttributeTok{fill =} \StringTok{"lightblue"}\NormalTok{) }\SpecialCharTok{+}
    \FunctionTok{labs}\NormalTok{(}\AttributeTok{title =} \FunctionTok{TeX}\NormalTok{(}\FunctionTok{paste0}\NormalTok{(}\StringTok{"$Poisson(}\SpecialCharTok{\textbackslash{}\textbackslash{}}\StringTok{lambda="}\NormalTok{, lambdas[i], }\StringTok{")$"}\NormalTok{)),}
         \AttributeTok{x =} \StringTok{"x"}\NormalTok{, }\AttributeTok{y =} \StringTok{"Frequência relativa"}\NormalTok{) }\SpecialCharTok{+}
    \FunctionTok{theme\_light}\NormalTok{()}
\NormalTok{\}}
\FunctionTok{grid.arrange}\NormalTok{(}\AttributeTok{grobs =}\NormalTok{ plots, }\AttributeTok{nrow =} \DecValTok{2}\NormalTok{, }\AttributeTok{ncol =} \DecValTok{3}\NormalTok{)}
\end{Highlighting}
\end{Shaded}

\begin{figure}

{\centering \includegraphics[width=0.8\linewidth]{introR_files/figure-latex/pois-sim-1} 

}

\caption{Amostras simuladas da Poisson para vários lambda.}\label{fig:pois-sim}
\end{figure}

\subsection{Função de distribuição}\label{funuxe7uxe3o-de-distribuiuxe7uxe3o-1}

\begin{Shaded}
\begin{Highlighting}[]
\NormalTok{lambda }\OtherTok{\textless{}{-}} \DecValTok{5}
\NormalTok{x }\OtherTok{\textless{}{-}} \DecValTok{0}\SpecialCharTok{:}\DecValTok{15}
\NormalTok{y }\OtherTok{\textless{}{-}} \FunctionTok{ppois}\NormalTok{(x, }\AttributeTok{lambda =}\NormalTok{ lambda)}

\FunctionTok{plot}\NormalTok{(x, y, }\AttributeTok{type =} \StringTok{"s"}\NormalTok{, }\AttributeTok{lwd =} \DecValTok{2}\NormalTok{, }\AttributeTok{col =} \StringTok{"blue"}\NormalTok{,}
     \AttributeTok{main =} \StringTok{"Função de distribuição da Poisson(5)"}\NormalTok{,}
     \AttributeTok{xlab =} \StringTok{"x"}\NormalTok{, }\AttributeTok{ylab =} \StringTok{"F(x)"}\NormalTok{)}
\end{Highlighting}
\end{Shaded}

\begin{figure}

{\centering \includegraphics[width=0.8\linewidth]{introR_files/figure-latex/pois-fd-1} 

}

\caption{Função de distribuição da Poisson(5).}\label{fig:pois-fd}
\end{figure}

\textbf{Exemplo.} Numa região ocorrem, em média, 3 sismos por mês; o número mensal de sismos
segue uma Poisson. Sendo \(X \sim \text{Poisson}(\lambda = 3)\):

\begin{Shaded}
\begin{Highlighting}[]
\FunctionTok{dpois}\NormalTok{(}\DecValTok{2}\NormalTok{, }\AttributeTok{lambda =} \DecValTok{3}\NormalTok{)                        }\CommentTok{\# P(X = 2)}
\DocumentationTok{\#\# [1] 0.2240418}
\FunctionTok{ppois}\NormalTok{(}\DecValTok{4}\NormalTok{, }\AttributeTok{lambda =} \DecValTok{3}\NormalTok{, }\AttributeTok{lower.tail =} \ConstantTok{FALSE}\NormalTok{)    }\CommentTok{\# P(X \textgreater{} 4)}
\DocumentationTok{\#\# [1] 0.1847368}
\end{Highlighting}
\end{Shaded}

Para a probabilidade de pelo menos 1 sismo em \textbf{2 semanas}, usamos a aditividade: em 2
semanas (meio mês) o parâmetro é \(\lambda = 1{,}5\), e
\(P(\tilde{X} \geq 1) = 1 - P(\tilde{X} = 0)\):

\begin{Shaded}
\begin{Highlighting}[]
\DecValTok{1} \SpecialCharTok{{-}} \FunctionTok{dpois}\NormalTok{(}\DecValTok{0}\NormalTok{, }\AttributeTok{lambda =} \FloatTok{1.5}\NormalTok{)                  }\CommentTok{\# P(X \textgreater{}= 1) em 2 semanas}
\DocumentationTok{\#\# [1] 0.7768698}
\end{Highlighting}
\end{Shaded}

\subsection{Exercícios}\label{exercuxedcios-29}

\textbf{1.} Uma fábrica produz, em média, 4 defeitos por dia (Poisson).

\begin{enumerate}
\def\labelenumi{(\alph{enumi})}
\tightlist
\item
  Qual a probabilidade de exatamente 5 defeitos num dia?
\item
  Qual a probabilidade de 3 ou menos defeitos num dia?
\end{enumerate}

\textbf{2.} Simule o número de defeitos em 30 dias com \texttt{rpois()} e \(\lambda = 4\).

\begin{enumerate}
\def\labelenumi{(\alph{enumi})}
\tightlist
\item
  Gere uma amostra de dimensão 30.
\item
  Calcule a média e o desvio padrão da amostra.
\item
  Compare com os valores teóricos.
\end{enumerate}

\textbf{3.} Uma linha de apoio recebe, em média, 20 pedidos por cada 10 minutos (Poisson).

\begin{enumerate}
\def\labelenumi{(\alph{enumi})}
\tightlist
\item
  Calcule a probabilidade de, em 10 minutos, chegarem exatamente 20 pedidos.
\item
  Simule para \(n_1 = 5\), \(n_2 = 10\), \(n_3 = 100\) e \(n_4 = 1000\) repetições; para cada
  caso, determine a percentagem de vezes em que chegam exatamente 20 pedidos.
\item
  Para cada amostra, calcule a média e a variância e compare com \(E(X)\) e \(V(X)\).
\end{enumerate}

\textbf{4.} Uma loja recebe, em média, 12 clientes por hora (Poisson).

\begin{enumerate}
\def\labelenumi{(\alph{enumi})}
\tightlist
\item
  Qual a probabilidade de receber no máximo 10 clientes numa hora?
\item
  Qual a probabilidade de receber mais de 15?
\item
  Simule 500 horas: (i) gere a amostra; (ii) calcule a frequência relativa de horas
  com \(\leq 10\) clientes e compare com (a).
\item
  Calcule a média e o desvio padrão da amostra e compare com os teóricos (\(\lambda = 12\)).
\item
  Compare graficamente a distribuição teórica com a frequência relativa observada.
\end{enumerate}

\textbf{5.} Numa área de conservação, avistam-se em média 8 aves por hora (\(\lambda = 8\)).

\begin{enumerate}
\def\labelenumi{(\alph{enumi})}
\tightlist
\item
  Simule o número de aves em 1000 horas.
\item
  Construa a função de distribuição empírica (\texttt{ecdf()}).
\item
  Compare-a com a função de distribuição teórica (\texttt{ppois()}).
\end{enumerate}

\textbf{6.} Com \texttt{set.seed(543)}, gere 2400 observações de \(Y \sim \text{Poisson}(\lambda = 6)\).

\begin{enumerate}
\def\labelenumi{(\alph{enumi})}
\tightlist
\item
  Faça um histograma de frequências relativas e sobreponha a distribuição de \(Y\).
\item
  Com a função de distribuição empírica, estime \(P(Y > 5)\) e compare com o teórico.
\end{enumerate}

\textbf{7.} Para \(\lambda = 5\), represente \(P(X = x)\) para \(x\) de 0 a 15.

\begin{enumerate}
\def\labelenumi{(\alph{enumi})}
\tightlist
\item
  Use \texttt{dpois()} para as probabilidades.
\item
  Faça um gráfico de barras.
\end{enumerate}

\textbf{8.} Para \(\lambda = 50\), use a aproximação normal:

\begin{enumerate}
\def\labelenumi{(\alph{enumi})}
\tightlist
\item
  Calcule \(P(X \geq 55)\) com a Poisson.
\item
  Calcule a mesma probabilidade com a aproximação \(N(\mu = 50, \sigma^2 = 50)\).
\item
  Compare os resultados.
\end{enumerate}

\textbf{9.} O número de acidentes por dia numa autoestrada segue uma Poisson com \(\lambda = 2\).

\begin{enumerate}
\def\labelenumi{(\alph{enumi})}
\tightlist
\item
  Simule 1000 amostras de dimensão 30.
\item
  Calcule a média de cada amostra.
\item
  Faça o histograma das médias amostrais e sobreponha a densidade de uma normal de
  média \(\mu = 2\) e desvio padrão \(\sigma = \sqrt{\lambda / n}\).
\end{enumerate}

\textbf{10.} Num hospital, atendem-se em média 5 pacientes por hora (\(\lambda = 5\)). Com
\texttt{set.seed(456)}:

\begin{itemize}
\tightlist
\item
  Simule 1000 amostras de dimensão 50, 100 e 1000.
\item
  Para cada dimensão, calcule a média de cada amostra.
\item
  Faça o histograma das médias e sobreponha a densidade normal de média \(\lambda\) e
  desvio padrão \(\sqrt{\lambda / n}\).
\item
  Comente a aproximação à normal à medida que \(n\) cresce, relacionando com o Teorema do
  Limite Central.
\end{itemize}

\section{Distribuição uniforme contínua}\label{distribuiuxe7uxe3o-uniforme-contuxednua}

\textbf{Definição.} A variável contínua \(X\) tem \textbf{distribuição uniforme} no intervalo
\((a, b)\), com \(a < b\), se a sua densidade for constante nesse intervalo:
\[f_X(x) = \begin{cases} \frac{1}{b-a}, & a \leq x \leq b \\ 0, & \text{caso contrário} \end{cases}, \qquad
F_X(x) = \begin{cases} 0, & x \leq a \\ \frac{x-a}{b-a}, & a < x < b \\ 1, & x \geq b \end{cases}\]

\textbf{Notação e momentos.} \(X \sim \text{Uniforme}(a, b)\), com \(E(X) = \dfrac{a+b}{2}\) e
\(V(X) = \dfrac{(b-a)^2}{12}\).

\subsection{Cálculo de probabilidades}\label{cuxe1lculo-de-probabilidades-4}

Seja \(X \sim \text{Uniforme}(0, 1)\):

\begin{itemize}
\tightlist
\item
  \(P(X \leq 0{,}5) \to\) \texttt{punif(0.5,\ min\ =\ 0,\ max\ =\ 1)} \(= 0{,}5\)
\item
  \(P(X > 0{,}5) \to\) \texttt{punif(0.5,\ min\ =\ 0,\ max\ =\ 1,\ lower.tail\ =\ FALSE)} \(= 0{,}5\)
\end{itemize}

\subsection{Densidade, simulação e comparação}\label{densidade-simulauxe7uxe3o-e-comparauxe7uxe3o}

\begin{Shaded}
\begin{Highlighting}[]
\NormalTok{x\_vals }\OtherTok{\textless{}{-}} \FunctionTok{seq}\NormalTok{(}\DecValTok{0}\NormalTok{, }\DecValTok{1}\NormalTok{, }\AttributeTok{length.out =} \DecValTok{100}\NormalTok{)}
\FunctionTok{plot}\NormalTok{(x\_vals, }\FunctionTok{dunif}\NormalTok{(x\_vals, }\DecValTok{0}\NormalTok{, }\DecValTok{1}\NormalTok{), }\AttributeTok{type =} \StringTok{"l"}\NormalTok{, }\AttributeTok{col =} \StringTok{"red"}\NormalTok{, }\AttributeTok{lwd =} \DecValTok{2}\NormalTok{,}
     \AttributeTok{main =} \StringTok{"Densidade da Uniforme(0, 1)"}\NormalTok{, }\AttributeTok{xlab =} \StringTok{"x"}\NormalTok{, }\AttributeTok{ylab =} \StringTok{"Densidade"}\NormalTok{)}
\end{Highlighting}
\end{Shaded}

\begin{figure}

{\centering \includegraphics[width=0.8\linewidth]{introR_files/figure-latex/unif-dens-1} 

}

\caption{Densidade teórica da Uniforme(0, 1).}\label{fig:unif-dens}
\end{figure}

\begin{Shaded}
\begin{Highlighting}[]
\FunctionTok{set.seed}\NormalTok{(}\DecValTok{123}\NormalTok{)}
\NormalTok{dados }\OtherTok{\textless{}{-}} \FunctionTok{runif}\NormalTok{(}\DecValTok{10000}\NormalTok{, }\AttributeTok{min =} \DecValTok{0}\NormalTok{, }\AttributeTok{max =} \DecValTok{1}\NormalTok{)}

\FunctionTok{hist}\NormalTok{(dados, }\AttributeTok{probability =} \ConstantTok{TRUE}\NormalTok{,}
     \AttributeTok{main =} \StringTok{"Amostra da Uniforme(0, 1)"}\NormalTok{,}
     \AttributeTok{xlab =} \StringTok{"x"}\NormalTok{, }\AttributeTok{ylab =} \StringTok{"Densidade"}\NormalTok{,}
     \AttributeTok{col =} \StringTok{"lightblue"}\NormalTok{, }\AttributeTok{border =} \StringTok{"darkblue"}\NormalTok{)}
\FunctionTok{curve}\NormalTok{(}\FunctionTok{dunif}\NormalTok{(x, }\DecValTok{0}\NormalTok{, }\DecValTok{1}\NormalTok{), }\AttributeTok{add =} \ConstantTok{TRUE}\NormalTok{, }\AttributeTok{col =} \StringTok{"red"}\NormalTok{, }\AttributeTok{lwd =} \DecValTok{2}\NormalTok{)}
\end{Highlighting}
\end{Shaded}

\begin{figure}

{\centering \includegraphics[width=0.8\linewidth]{introR_files/figure-latex/unif-comp-1} 

}

\caption{Histograma de uma amostra da Uniforme(0, 1) com a densidade teórica sobreposta.}\label{fig:unif-comp}
\end{figure}

\begin{Shaded}
\begin{Highlighting}[]
\NormalTok{x\_vals }\OtherTok{\textless{}{-}} \FunctionTok{seq}\NormalTok{(}\DecValTok{0}\NormalTok{, }\DecValTok{1}\NormalTok{, }\AttributeTok{length.out =} \DecValTok{100}\NormalTok{)}
\FunctionTok{plot}\NormalTok{(x\_vals, }\FunctionTok{punif}\NormalTok{(x\_vals, }\DecValTok{0}\NormalTok{, }\DecValTok{1}\NormalTok{), }\AttributeTok{type =} \StringTok{"l"}\NormalTok{, }\AttributeTok{col =} \StringTok{"blue"}\NormalTok{, }\AttributeTok{lwd =} \DecValTok{2}\NormalTok{,}
     \AttributeTok{main =} \StringTok{"Função de distribuição da Uniforme(0, 1)"}\NormalTok{, }\AttributeTok{xlab =} \StringTok{"x"}\NormalTok{, }\AttributeTok{ylab =} \StringTok{"F(x)"}\NormalTok{)}
\end{Highlighting}
\end{Shaded}

\begin{figure}

{\centering \includegraphics[width=0.8\linewidth]{introR_files/figure-latex/unif-fd-1} 

}

\caption{Função de distribuição da Uniforme(0, 1).}\label{fig:unif-fd}
\end{figure}

\subsection{Exercícios}\label{exercuxedcios-30}

\textbf{1.} Simule 1000 valores de \(X \sim U(0, 1)\).

\begin{enumerate}
\def\labelenumi{(\alph{enumi})}
\tightlist
\item
  Calcule a média e a variância dos valores simulados.
\item
  Compare com os valores teóricos (\(E(X) = 0{,}5\), \(V(X) = 1/12\)).
\end{enumerate}

\textbf{2.} Simule 500 valores de \(X \sim U(-3, 7)\).

\begin{enumerate}
\def\labelenumi{(\alph{enumi})}
\tightlist
\item
  Faça o histograma dos valores.
\item
  Sobreponha a densidade teórica.
\end{enumerate}

\textbf{3.} O peso de uma barra de chocolate (que deveria pesar 100 g) tem distribuição
uniforme em \([85, 105]\) gramas.

\begin{enumerate}
\def\labelenumi{(\alph{enumi})}
\tightlist
\item
  Qual a probabilidade de uma barra pesar menos de 100 g?
\item
  Simule para \(n_1 = 5\), \(n_2 = 10\), \(n_3 = 100\) e \(n_4 = 1000\); para cada caso,
  determine a percentagem de barras com peso inferior a 100 g.
\end{enumerate}

\textbf{4.} Simule 10 000 valores de \(X \sim U(2, 8)\).

\begin{enumerate}
\def\labelenumi{(\alph{enumi})}
\tightlist
\item
  Estime \(P(X > 5)\).
\item
  Compare com o valor teórico (\texttt{punif()}).
\end{enumerate}

\textbf{5.} Seja \(Y \sim U(10, 20)\). Simule 1000 valores.

\begin{enumerate}
\def\labelenumi{(\alph{enumi})}
\tightlist
\item
  Calcule a proporção empírica de valores em \([12, 15]\).
\item
  Compare com o valor teórico (\texttt{punif()}).
\end{enumerate}

\textbf{6.} Simule 1000 amostras de dimensão 30 de \(X \sim U(-5, 5)\).

\begin{enumerate}
\def\labelenumi{(\alph{enumi})}
\tightlist
\item
  Calcule a média de cada amostra.
\item
  Faça o histograma das médias e sobreponha a densidade normal de média 0 e variância
  \(\frac{(5-(-5))^2}{12 \cdot n}\).
\end{enumerate}

\textbf{7.} Seja \(Z = 3X + 2\), com \(X \sim U(0, 1)\). Simule 1000 valores de \(X\) e
transforme-os em \(Z\).

\begin{enumerate}
\def\labelenumi{(\alph{enumi})}
\tightlist
\item
  Calcule a média e a variância de \(Z\).
\item
  Compare com \(E(Z) = 3E(X) + 2\) e \(V(Z) = 9V(X)\).
\end{enumerate}

\textbf{8.} Uma fábrica produz itens com peso uniforme em \([100, 120]\) g. Simule 2000 itens.

\begin{enumerate}
\def\labelenumi{(\alph{enumi})}
\tightlist
\item
  Calcule a proporção com peso inferior a 105 g.
\item
  Compare a densidade empírica com a teórica.
\end{enumerate}

\textbf{9.} Com \texttt{set.seed(123)}, gere amostras de dimensão \(n = 100, 1000, 10\,000, 100\,000\)
de \(W \sim U(0, 1)\). Para cada uma, calcule a média amostral e compare com o valor
esperado teórico. Comente a convergência (Lei dos Grandes Números).

\textbf{10.} O tempo de entrega de um \emph{drone} (minutos) segue \(X \sim U(10, 30)\). Com
\texttt{set.seed(1430)}, gere 8000 amostras de dimensão \(n = 100\).

\begin{enumerate}
\def\labelenumi{(\alph{enumi})}
\tightlist
\item
  Calcule a soma de cada amostra, \(S_n = \sum_{i=1}^n X_i\).
\item
  Faça o histograma das somas e sobreponha uma normal de média \(nE(X)\) e desvio padrão
  \(\sqrt{n\,V(X)}\).
\item
  Calcule a média de cada amostra, \(\bar{X}_n\).
\item
  Faça o histograma das médias e sobreponha uma normal de média \(E(X)\) e desvio padrão
  \(\sqrt{V(X)/n}\).
\end{enumerate}

\textbf{11.} Num \emph{data center}, o número de servidores que falham por hora segue uma
Poisson(\(\lambda = 4\)); o tempo de reparação de cada um segue \(U(2, 6)\) horas, e a
energia gasta é \(E_i = T_i^2\). Simule 1000 horas e, em cada uma, calcule o consumo total
\(E_{\text{total}} = \sum_{i=1}^N T_i^2\), onde \(N\) é o número de servidores falhados.

\section{Distribuição exponencial}\label{distribuiuxe7uxe3o-exponencial}

A distribuição exponencial modela \textbf{tempos} --- a duração de um equipamento, o intervalo
entre ocorrências consecutivas de um mesmo tipo de acontecimento (chegadas de clientes,
avarias, colisões).

\textbf{Definição.} A variável contínua \(X\) tem \textbf{distribuição exponencial} de parâmetro
\(\lambda > 0\) se
\[f_X(x) = \begin{cases} \lambda e^{-\lambda x}, & x > 0 \\ 0, & x \leq 0 \end{cases}, \qquad
F_X(x) = \begin{cases} 1 - e^{-\lambda x}, & x > 0 \\ 0, & x \leq 0 \end{cases}\]

\textbf{Notação e momentos.} \(X \sim \text{Exponencial}(\lambda)\), com \(E(X) = \dfrac{1}{\lambda}\)
e \(V(X) = \dfrac{1}{\lambda^2}\).

A exponencial é a única distribuição contínua \textbf{sem memória}: \(P(X > s + t \mid X > s) = P(X > t)\).
Por palavras: se um componente já durou \(s\) horas, a probabilidade de durar mais \(t\) é a
mesma que teria em novo. É esta propriedade que a liga naturalmente ao processo de
Poisson --- os intervalos entre ocorrências de uma Poisson(\(\lambda\)) são
exponenciais(\(\lambda\)).

\subsection{Cálculo de probabilidades}\label{cuxe1lculo-de-probabilidades-5}

Seja \(X \sim \text{Exponencial}(\lambda = 1)\):

\begin{itemize}
\tightlist
\item
  \(P(X \leq 0{,}5) \to\) \texttt{pexp(0.5,\ rate\ =\ 1)} \(= 0{,}3935\)
\item
  \(P(X > 0{,}5) \to\) \texttt{pexp(0.5,\ rate\ =\ 1,\ lower.tail\ =\ FALSE)} \(= 0{,}6065\)
\end{itemize}

\subsection{Densidade, simulação e comparação}\label{densidade-simulauxe7uxe3o-e-comparauxe7uxe3o-1}

\begin{Shaded}
\begin{Highlighting}[]
\NormalTok{x\_vals }\OtherTok{\textless{}{-}} \FunctionTok{seq}\NormalTok{(}\DecValTok{0}\NormalTok{, }\DecValTok{10}\NormalTok{, }\AttributeTok{length.out =} \DecValTok{100}\NormalTok{)}
\FunctionTok{plot}\NormalTok{(x\_vals, }\FunctionTok{dexp}\NormalTok{(x\_vals, }\AttributeTok{rate =} \DecValTok{1}\NormalTok{), }\AttributeTok{type =} \StringTok{"l"}\NormalTok{, }\AttributeTok{col =} \StringTok{"red"}\NormalTok{, }\AttributeTok{lwd =} \DecValTok{2}\NormalTok{,}
     \AttributeTok{main =} \StringTok{"Densidade da Exponencial(1)"}\NormalTok{, }\AttributeTok{xlab =} \StringTok{"x"}\NormalTok{, }\AttributeTok{ylab =} \StringTok{"Densidade"}\NormalTok{)}
\end{Highlighting}
\end{Shaded}

\begin{figure}

{\centering \includegraphics[width=0.8\linewidth]{introR_files/figure-latex/exp-dens-1} 

}

\caption{Densidade teórica da Exponencial(1).}\label{fig:exp-dens}
\end{figure}

\begin{Shaded}
\begin{Highlighting}[]
\FunctionTok{set.seed}\NormalTok{(}\DecValTok{123}\NormalTok{)}
\NormalTok{dados }\OtherTok{\textless{}{-}} \FunctionTok{rexp}\NormalTok{(}\DecValTok{10000}\NormalTok{, }\AttributeTok{rate =} \DecValTok{1}\NormalTok{)}

\FunctionTok{hist}\NormalTok{(dados, }\AttributeTok{probability =} \ConstantTok{TRUE}\NormalTok{,}
     \AttributeTok{main =} \StringTok{"Amostra da Exponencial(1)"}\NormalTok{,}
     \AttributeTok{xlab =} \StringTok{"x"}\NormalTok{, }\AttributeTok{ylab =} \StringTok{"Densidade"}\NormalTok{,}
     \AttributeTok{col =} \StringTok{"lightblue"}\NormalTok{, }\AttributeTok{border =} \StringTok{"darkblue"}\NormalTok{)}
\FunctionTok{curve}\NormalTok{(}\FunctionTok{dexp}\NormalTok{(x, }\AttributeTok{rate =} \DecValTok{1}\NormalTok{), }\AttributeTok{add =} \ConstantTok{TRUE}\NormalTok{, }\AttributeTok{col =} \StringTok{"red"}\NormalTok{, }\AttributeTok{lwd =} \DecValTok{2}\NormalTok{)}
\end{Highlighting}
\end{Shaded}

\begin{figure}

{\centering \includegraphics[width=0.8\linewidth]{introR_files/figure-latex/exp-comp-1} 

}

\caption{Histograma de uma amostra da Exponencial(1) com a densidade teórica sobreposta.}\label{fig:exp-comp}
\end{figure}

\begin{Shaded}
\begin{Highlighting}[]
\NormalTok{x\_vals }\OtherTok{\textless{}{-}} \FunctionTok{seq}\NormalTok{(}\DecValTok{0}\NormalTok{, }\DecValTok{10}\NormalTok{, }\AttributeTok{length.out =} \DecValTok{100}\NormalTok{)}
\FunctionTok{plot}\NormalTok{(x\_vals, }\FunctionTok{pexp}\NormalTok{(x\_vals, }\AttributeTok{rate =} \DecValTok{1}\NormalTok{), }\AttributeTok{type =} \StringTok{"l"}\NormalTok{, }\AttributeTok{col =} \StringTok{"red"}\NormalTok{, }\AttributeTok{lwd =} \DecValTok{2}\NormalTok{,}
     \AttributeTok{main =} \StringTok{"Função de distribuição da Exponencial(1)"}\NormalTok{, }\AttributeTok{xlab =} \StringTok{"x"}\NormalTok{, }\AttributeTok{ylab =} \StringTok{"F(x)"}\NormalTok{)}
\end{Highlighting}
\end{Shaded}

\begin{figure}

{\centering \includegraphics[width=0.8\linewidth]{introR_files/figure-latex/exp-fd-1} 

}

\caption{Função de distribuição da Exponencial(1).}\label{fig:exp-fd}
\end{figure}

\subsection{Exercícios}\label{exercuxedcios-31}

\textbf{1.} Os intervalos entre chamadas numa central seguem uma exponencial com \(\lambda = 2\)
(chamadas por minuto).

\begin{enumerate}
\def\labelenumi{(\alph{enumi})}
\tightlist
\item
  Simule 1000 intervalos.
\item
  Calcule o tempo médio e compare com \(1/\lambda\).
\item
  Faça o histograma e sobreponha a densidade teórica.
\end{enumerate}

\textbf{2.} Numa fila, o tempo entre chegadas segue uma exponencial com \(\lambda = 0{,}5\).

\begin{enumerate}
\def\labelenumi{(\alph{enumi})}
\tightlist
\item
  Calcule \(P(\text{tempo} > 3)\) teoricamente.
\item
  Simule 5000 intervalos e estime essa probabilidade.
\item
  Compare.
\end{enumerate}

\textbf{3.} O tempo de atendimento segue uma exponencial com \(\lambda = 0{,}25\) por minuto.

\begin{enumerate}
\def\labelenumi{(\alph{enumi})}
\tightlist
\item
  Simule os tempos para \(n = 50\) clientes.
\item
  Calcule o tempo total.
\item
  Após 1000 simulações, faça o histograma do tempo total e sobreponha uma normal de
  média \(n/\lambda\) e desvio padrão \(\sqrt{n/\lambda^2}\).
\end{enumerate}

\textbf{4.} Dois processos independentes têm tempos exponenciais com \(\lambda_1 = 3\) e
\(\lambda_2 = 5\).

\begin{enumerate}
\def\labelenumi{(\alph{enumi})}
\tightlist
\item
  Simule 1000 tempos de cada.
\item
  Some os tempos.
\item
  Faça o histograma da soma e discuta se ainda é exponencial.
\end{enumerate}

\textbf{5.} Simule 5000 amostras de dimensão 30 de \(X \sim \text{Exponencial}(\lambda = 1{,}5)\).

\begin{enumerate}
\def\labelenumi{(\alph{enumi})}
\tightlist
\item
  Calcule a média de cada amostra.
\item
  Faça o histograma das médias e sobreponha uma normal de média \(1/\lambda\) e desvio
  padrão \(1/(\lambda\sqrt{n})\).
\item
  Explique a ligação ao Teorema do Limite Central.
\end{enumerate}

\textbf{6.} O tempo até à primeira avaria de uma máquina segue uma exponencial com
\(\lambda = 0{,}1\).

\begin{enumerate}
\def\labelenumi{(\alph{enumi})}
\tightlist
\item
  Simule o tempo de avaria de 10 000 máquinas.
\item
  Estime a probabilidade de avariar em menos de 15 horas.
\item
  Compare com o valor teórico (função de distribuição).
\end{enumerate}

\textbf{7.} Numa estrada, os intervalos entre acidentes seguem uma exponencial com
\(\lambda = 0{,}8\) (acidentes por hora).

\begin{enumerate}
\def\labelenumi{(\alph{enumi})}
\tightlist
\item
  Simule os intervalos para 1000 horas.
\item
  Estime a probabilidade de um intervalo ser inferior a 2 horas.
\item
  Faça o histograma e compare com a densidade teórica.
\end{enumerate}

\textbf{8.} Considere \(X_1 \sim \text{Exponencial}(0{,}5)\) e \(X_2 = 2X_1 + 1\).

\begin{enumerate}
\def\labelenumi{(\alph{enumi})}
\tightlist
\item
  Simule 5000 pares \((X_1, X_2)\).
\item
  Calcule a média e a variância de \(X_2\).
\item
  Faça o gráfico de dispersão de \(X_1\) e \(X_2\) e comente a relação.
\end{enumerate}

\textbf{9.} Os intervalos entre chamadas seguem uma exponencial com \(\lambda = 2\).

\begin{enumerate}
\def\labelenumi{(\alph{enumi})}
\tightlist
\item
  Simule 5000 amostras de dimensão \(n = 5\).
\item
  Para cada amostra, calcule a soma \(S = \sum_{i=1}^n X_i\).
\item
  Compare o histograma de \(S\) com a densidade de uma \(\text{Gama}(k = 5, \theta = 1/\lambda)\).
\item
  Compare a média e a variância empíricas de \(S\) com \(E(S) = k\theta\) e \(V(S) = k\theta^2\).
\end{enumerate}

\textbf{10.} O tempo até à avaria de uma máquina segue uma exponencial com \(\lambda = 0{,}5\);
observa-se o tempo até 10 avarias consecutivas.

\begin{enumerate}
\def\labelenumi{(\alph{enumi})}
\tightlist
\item
  Simule 10 000 observações de \(S = \sum_{i=1}^{10} X_i\).
\item
  Compare o histograma com a densidade de uma \(\text{Gama}(k = 10, \theta = 1/\lambda)\).
\item
  Calcule \(P(S > 25)\) teoricamente (densidade gama) e compare com a estimativa
  empírica.
\end{enumerate}

\textbf{11.} O tempo de uma tarefa é a soma de \(n = 7\) tempos exponenciais com \(\lambda = 3\).

\begin{enumerate}
\def\labelenumi{(\alph{enumi})}
\tightlist
\item
  Simule 5000 valores de \(S = \sum_{i=1}^n X_i\).
\item
  Ajuste-os a uma \(\text{Gama}(k = 7, \theta = 1/\lambda)\).
\item
  Faça o histograma e compare com a densidade teórica.
\item
  Estime o percentil 90 de \(S\) e compare com o teórico (\texttt{qgamma()} e \texttt{quantile()}).
\end{enumerate}

\section{Distribuição normal}\label{distribuiuxe7uxe3o-normal}

A distribuição normal é a mais importante da estatística. Vejamos alguns exemplos, a
começar pela normal padrão \(N(\mu = 0, \sigma = 1)\), que é o que as funções assumem por
omissão:

\begin{Shaded}
\begin{Highlighting}[]
\FunctionTok{dnorm}\NormalTok{(}\SpecialCharTok{{-}}\DecValTok{1}\NormalTok{)      }\CommentTok{\# densidade em x = {-}1}
\DocumentationTok{\#\# [1] 0.2419707}
\FunctionTok{pnorm}\NormalTok{(}\SpecialCharTok{{-}}\DecValTok{1}\NormalTok{)      }\CommentTok{\# P(X \textless{}= {-}1)}
\DocumentationTok{\#\# [1] 0.1586553}
\FunctionTok{qnorm}\NormalTok{(}\FloatTok{0.975}\NormalTok{)   }\CommentTok{\# x tal que P(X \textless{}= x) = 0.975}
\DocumentationTok{\#\# [1] 1.959964}
\FunctionTok{rnorm}\NormalTok{(}\DecValTok{10}\NormalTok{)      }\CommentTok{\# amostra de 10 valores (varia a cada execução)}
\end{Highlighting}
\end{Shaded}

O valor de \texttt{dnorm(-1)} é a densidade da normal padrão no ponto \(x = -1\). A densidade da
normal \(N(\mu, \sigma)\) é
\[f(x) = \frac{1}{\sigma\sqrt{2\pi}}\, e^{-\frac{1}{2}\left(\frac{x - \mu}{\sigma}\right)^2},\]
pelo que o mesmo valor se obtém substituindo diretamente na fórmula:

\begin{Shaded}
\begin{Highlighting}[]
\NormalTok{mu }\OtherTok{\textless{}{-}} \DecValTok{0}\NormalTok{; sigma }\OtherTok{\textless{}{-}} \DecValTok{1}\NormalTok{; x }\OtherTok{\textless{}{-}} \SpecialCharTok{{-}}\DecValTok{1}
\NormalTok{(}\DecValTok{1} \SpecialCharTok{/}\NormalTok{ (sigma }\SpecialCharTok{*} \FunctionTok{sqrt}\NormalTok{(}\DecValTok{2} \SpecialCharTok{*}\NormalTok{ pi))) }\SpecialCharTok{*} \FunctionTok{exp}\NormalTok{(}\SpecialCharTok{{-}}\DecValTok{1}\SpecialCharTok{/}\DecValTok{2} \SpecialCharTok{*}\NormalTok{ ((x }\SpecialCharTok{{-}}\NormalTok{ mu) }\SpecialCharTok{/}\NormalTok{ sigma)}\SpecialCharTok{\^{}}\DecValTok{2}\NormalTok{)}
\DocumentationTok{\#\# [1] 0.242}
\end{Highlighting}
\end{Shaded}

As funções aceitam os argumentos \texttt{mean} e \texttt{sd} para outras médias e desvios padrão, que
podem ser passados de várias formas (por nome, abreviados ou por posição):

\begin{Shaded}
\begin{Highlighting}[]
\FunctionTok{qnorm}\NormalTok{(}\FloatTok{0.975}\NormalTok{, }\AttributeTok{mean =} \DecValTok{100}\NormalTok{, }\AttributeTok{sd =} \DecValTok{8}\NormalTok{)   }\CommentTok{\# todas estas três chamadas...}
\FunctionTok{qnorm}\NormalTok{(}\FloatTok{0.975}\NormalTok{, }\AttributeTok{m =} \DecValTok{100}\NormalTok{, }\AttributeTok{s =} \DecValTok{8}\NormalTok{)       }\CommentTok{\# ...são equivalentes}
\FunctionTok{qnorm}\NormalTok{(}\FloatTok{0.975}\NormalTok{, }\DecValTok{100}\NormalTok{, }\DecValTok{8}\NormalTok{)}
\DocumentationTok{\#\# [1] 115.6797}
\end{Highlighting}
\end{Shaded}

No R, o segundo parâmetro da normal é o \textbf{desvio padrão} (\(\sigma\)), e não a variância
(\(\sigma^2\)). Assim, para \(X \sim N(\mu = 100, \sigma^2 = 64)\) escreve-se
\texttt{sd\ =\ 8}, e não \texttt{sd\ =\ 64}. Confundir os dois é um erro comum e silencioso --- o R não
avisa, apenas devolve resultados errados.

Os cálculos de probabilidades que outrora exigiam tabelas fazem-se agora diretamente.
Seja \(X \sim N(\mu = 100, \sigma = 10)\):

\begin{Shaded}
\begin{Highlighting}[]
\FunctionTok{pnorm}\NormalTok{(}\DecValTok{95}\NormalTok{, }\DecValTok{100}\NormalTok{, }\DecValTok{10}\NormalTok{)                          }\CommentTok{\# P(X \textless{} 95)}
\DocumentationTok{\#\# [1] 0.3085375}
\FunctionTok{pnorm}\NormalTok{(}\DecValTok{110}\NormalTok{, }\DecValTok{100}\NormalTok{, }\DecValTok{10}\NormalTok{) }\SpecialCharTok{{-}} \FunctionTok{pnorm}\NormalTok{(}\DecValTok{90}\NormalTok{, }\DecValTok{100}\NormalTok{, }\DecValTok{10}\NormalTok{)    }\CommentTok{\# P(90 \textless{} X \textless{} 110)}
\DocumentationTok{\#\# [1] 0.6826895}
\FunctionTok{pnorm}\NormalTok{(}\DecValTok{95}\NormalTok{, }\DecValTok{100}\NormalTok{, }\DecValTok{10}\NormalTok{, }\AttributeTok{lower.tail =} \ConstantTok{FALSE}\NormalTok{)      }\CommentTok{\# P(X \textgreater{} 95)}
\DocumentationTok{\#\# [1] 0.6914625}
\end{Highlighting}
\end{Shaded}

\subsection{Densidade e função de distribuição}\label{densidade-e-funuxe7uxe3o-de-distribuiuxe7uxe3o}

\begin{Shaded}
\begin{Highlighting}[]
\FunctionTok{par}\NormalTok{(}\AttributeTok{mfrow =} \FunctionTok{c}\NormalTok{(}\DecValTok{1}\NormalTok{, }\DecValTok{2}\NormalTok{))}
\FunctionTok{plot}\NormalTok{(dnorm, }\AttributeTok{from =} \SpecialCharTok{{-}}\DecValTok{3}\NormalTok{, }\AttributeTok{to =} \DecValTok{3}\NormalTok{, }\AttributeTok{xlab =} \StringTok{"x"}\NormalTok{, }\AttributeTok{ylab =} \StringTok{"Densidade"}\NormalTok{)}
\FunctionTok{title}\NormalTok{(}\StringTok{"Densidade}\SpecialCharTok{\textbackslash{}n}\StringTok{N(0, 1)"}\NormalTok{)}
\FunctionTok{plot}\NormalTok{(pnorm, }\AttributeTok{from =} \SpecialCharTok{{-}}\DecValTok{3}\NormalTok{, }\AttributeTok{to =} \DecValTok{3}\NormalTok{, }\AttributeTok{xlab =} \StringTok{"x"}\NormalTok{, }\AttributeTok{ylab =} \StringTok{"F(x)"}\NormalTok{)}
\FunctionTok{title}\NormalTok{(}\StringTok{"Função de distribuição}\SpecialCharTok{\textbackslash{}n}\StringTok{N(0, 1)"}\NormalTok{)}
\end{Highlighting}
\end{Shaded}

\begin{figure}

{\centering \includegraphics[width=0.8\linewidth]{introR_files/figure-latex/norm-padrao-1} 

}

\caption{Densidade e função de distribuição da normal padrão.}\label{fig:norm-padrao}
\end{figure}

\begin{Shaded}
\begin{Highlighting}[]
\FunctionTok{par}\NormalTok{(}\AttributeTok{mfrow =} \FunctionTok{c}\NormalTok{(}\DecValTok{1}\NormalTok{, }\DecValTok{1}\NormalTok{))}
\end{Highlighting}
\end{Shaded}

O efeito dos parâmetros vê-se bem comparando várias normais: a média desloca a curva; o
desvio padrão controla a sua largura:

\begin{Shaded}
\begin{Highlighting}[]
\FunctionTok{plot}\NormalTok{(}\ControlFlowTok{function}\NormalTok{(x) }\FunctionTok{dnorm}\NormalTok{(x, }\DecValTok{100}\NormalTok{, }\DecValTok{8}\NormalTok{), }\DecValTok{60}\NormalTok{, }\DecValTok{140}\NormalTok{, }\AttributeTok{ylab =} \StringTok{"f(x)"}\NormalTok{, }\AttributeTok{xlab =} \StringTok{"x"}\NormalTok{)}
\FunctionTok{plot}\NormalTok{(}\ControlFlowTok{function}\NormalTok{(x) }\FunctionTok{dnorm}\NormalTok{(x, }\DecValTok{90}\NormalTok{, }\DecValTok{8}\NormalTok{), }\DecValTok{60}\NormalTok{, }\DecValTok{140}\NormalTok{, }\AttributeTok{add =} \ConstantTok{TRUE}\NormalTok{, }\AttributeTok{col =} \DecValTok{2}\NormalTok{)}
\FunctionTok{plot}\NormalTok{(}\ControlFlowTok{function}\NormalTok{(x) }\FunctionTok{dnorm}\NormalTok{(x, }\DecValTok{100}\NormalTok{, }\DecValTok{15}\NormalTok{), }\DecValTok{60}\NormalTok{, }\DecValTok{140}\NormalTok{, }\AttributeTok{add =} \ConstantTok{TRUE}\NormalTok{, }\AttributeTok{col =} \DecValTok{3}\NormalTok{)}
\FunctionTok{legend}\NormalTok{(}\StringTok{"topright"}\NormalTok{, }\AttributeTok{fill =} \DecValTok{1}\SpecialCharTok{:}\DecValTok{3}\NormalTok{,}
       \AttributeTok{legend =} \FunctionTok{c}\NormalTok{(}\StringTok{"N(100, 64)"}\NormalTok{, }\StringTok{"N(90, 64)"}\NormalTok{, }\StringTok{"N(100, 225)"}\NormalTok{))}
\end{Highlighting}
\end{Shaded}

\begin{figure}

{\centering \includegraphics[width=0.8\linewidth]{introR_files/figure-latex/norm-param-1} 

}

\caption{Três normais com médias e desvios padrão diferentes.}\label{fig:norm-param}
\end{figure}

\subsection{Exercícios}\label{exercuxedcios-32}

\textbf{1.} O peso de um mineral (g) segue uma normal de média \(\mu = 50\) e desvio padrão
\(\sigma = 5\).

\begin{enumerate}
\def\labelenumi{(\alph{enumi})}
\tightlist
\item
  Simule 1000 pesos.
\item
  Compare a média e o desvio padrão simulados com os teóricos.
\item
  Faça o histograma e sobreponha a densidade teórica.
\end{enumerate}

\textbf{2.} Sejam \(X_1 \sim N(5, 2)\) e \(X_2 \sim N(10, 3)\) independentes.

\begin{enumerate}
\def\labelenumi{(\alph{enumi})}
\tightlist
\item
  Simule 1000 pares.
\item
  Calcule \(S = X_1 + X_2\).
\item
  Faça o histograma de \(S\) e sobreponha uma normal de média \(\mu_1 + \mu_2\) e variância
  \(\sigma_1^2 + \sigma_2^2\).
\end{enumerate}

\textbf{3.} Sejam \(X_1, X_2 \sim N(0, 1)\) independentes.

\begin{enumerate}
\def\labelenumi{(\alph{enumi})}
\tightlist
\item
  Simule 5000 pares.
\item
  Calcule \(Q = X_1^2 + X_2^2\).
\item
  Compare o histograma de \(Q\) com a densidade de uma qui-quadrado com 2 graus de
  liberdade (\texttt{dchisq()}).
\end{enumerate}

\textbf{4.} Estude a média de amostras de uma normal padrão \(N(0, 1)\).

\begin{enumerate}
\def\labelenumi{(\alph{enumi})}
\tightlist
\item
  Extraia 1000 amostras de cada dimensão \(n = 5, 10, 30, 100\).
\item
  Para cada amostra, calcule a média.
\item
  Faça o histograma das médias e sobreponha a densidade de uma \(t\)-Student com \(n - 1\)
  graus de liberdade e a de uma normal de média 0 e variância \(1/n\).
\item
  Comente como a \(t\)-Student se aproxima da normal quando \(n\) aumenta, e o efeito de
  \(n\) na variabilidade das médias.
\end{enumerate}

\textbf{5.} Simule 1000 amostras de dimensão 20 de \(X \sim N(0, 1)\).

\begin{enumerate}
\def\labelenumi{(\alph{enumi})}
\tightlist
\item
  Calcule a soma dos quadrados de cada amostra.
\item
  Compare o histograma com a densidade de uma qui-quadrado com 20 graus de liberdade.
\item
  Interprete os resultados.
\end{enumerate}

\textbf{6.} Simule 5000 amostras de dimensão 30 de \(X \sim N(\mu = 10, \sigma = 4)\).

\begin{enumerate}
\def\labelenumi{(\alph{enumi})}
\tightlist
\item
  Calcule a média de cada amostra.
\item
  Faça o histograma das médias e sobreponha uma normal de média 10 e desvio padrão
  \(4/\sqrt{30}\).
\item
  Explique a ligação ao Teorema do Limite Central.
\end{enumerate}

\textbf{7.} Sejam \(X_1 \sim N(0, 1)\) e \(X_2 = 0{,}5 X_1 + Z\), com \(Z \sim N(0, 1)\)
independente.

\begin{enumerate}
\def\labelenumi{(\alph{enumi})}
\tightlist
\item
  Simule 5000 pares.
\item
  Calcule \(S = X_1 + X_2\).
\item
  Faça o histograma de \(S\).
\end{enumerate}

\textbf{8.} Simule 1000 amostras de dimensão 10 de \(X \sim N(0, 1)\). Para cada uma, calcule a
média \(\bar{X}\) e a variância \(S^2\), e construa
\(T = \dfrac{\bar{X}}{S / \sqrt{n}}\). Compare o histograma de \(T\) com a densidade de uma
\(t\)-Student com \(n - 1\) graus de liberdade.

\textbf{9.} Seja \(X \sim N(0, 1)\). Simule 1000 amostras de dimensão \(n = 10\). Para cada uma:

\begin{enumerate}
\def\labelenumi{(\alph{enumi})}
\tightlist
\item
  Calcule a variância amostral \(S^2\).
\item
  Construa a estatística \(Q = \dfrac{(n-1)S^2}{\sigma^2}\), com \(\sigma^2 = 1\).
\item
  Compare o histograma de \(Q\) com a densidade de uma \(\chi^2(n-1)\).
\item
  Compare a média e a variância empíricas de \(Q\) com \(E(Q) = n - 1\) e \(V(Q) = 2(n-1)\).
\end{enumerate}

Conhecidas as distribuições e as funções \texttt{r}\emph{nome} que geram amostras diretamente, uma
pergunta natural impõe-se: e quando a distribuição de que precisamos não tem uma função
\texttt{r} pronta no R? É esse o problema que os próximos capítulos resolvem, a começar pelo
método da transformada inversa.

\chapter{Exercícios extra}\label{exercicios-extra}

Antes de avançarmos para os métodos de geração de variáveis aleatórias, este capítulo
propõe uma bateria de exercícios que consolidam tudo o que vimos sobre distribuições e
simulação. Os primeiros são de aquecimento --- gerar amostras, calcular médias e
variâncias, comparar frequências empíricas com probabilidades teóricas. Os últimos
antecipam um dos resultados mais belos da estatística, o \textbf{Teorema do Limite Central}:
ao somar ou promediar muitas variáveis independentes, a distribuição do resultado
aproxima-se de uma normal, seja qual for a distribuição de partida.

Consolidar, através de exercícios aplicados, a geração de amostras com as funções
\texttt{r}\emph{nome}, o cálculo de probabilidades teóricas e empíricas, e a observação da Lei dos
Grandes Números e do Teorema do Limite Central por simulação.

\textbf{1.} Com \texttt{set.seed(123)}, simule 1000 lançamentos de uma moeda com probabilidade 0,5
de sair cara. Conte o número de caras e faça um gráfico de barras dos resultados.

\begin{Shaded}
\begin{Highlighting}[]
\FunctionTok{set.seed}\NormalTok{(}\DecValTok{123}\NormalTok{)}
\NormalTok{n }\OtherTok{\textless{}{-}} \DecValTok{1000}
\NormalTok{p }\OtherTok{\textless{}{-}} \FloatTok{0.5}

\NormalTok{amostra }\OtherTok{\textless{}{-}} \FunctionTok{rbinom}\NormalTok{(n, }\DecValTok{1}\NormalTok{, p)      }\CommentTok{\# 1000 provas de Bernoulli}
\FunctionTok{sum}\NormalTok{(amostra)                    }\CommentTok{\# número de caras}

\NormalTok{freq\_relativa }\OtherTok{\textless{}{-}} \FunctionTok{table}\NormalTok{(amostra) }\SpecialCharTok{/} \FunctionTok{length}\NormalTok{(amostra)}
\FunctionTok{barplot}\NormalTok{(freq\_relativa, }\AttributeTok{col =} \StringTok{"lightblue"}\NormalTok{)}
\end{Highlighting}
\end{Shaded}

\textbf{2.} Com \texttt{set.seed(123)}, gere 5000 observações de uma binomial com \(n = 10\) e
\(p = 0{,}3\). Calcule a média e a variância.

\begin{Shaded}
\begin{Highlighting}[]
\FunctionTok{set.seed}\NormalTok{(}\DecValTok{123}\NormalTok{)}
\NormalTok{amostra }\OtherTok{\textless{}{-}} \FunctionTok{rbinom}\NormalTok{(}\DecValTok{5000}\NormalTok{, }\DecValTok{10}\NormalTok{, }\FloatTok{0.3}\NormalTok{)}
\FunctionTok{mean}\NormalTok{(amostra)}
\FunctionTok{var}\NormalTok{(amostra)}
\end{Highlighting}
\end{Shaded}

\textbf{3.} Com \texttt{set.seed(123)}, gere 2300 observações de uma Poisson com \(\lambda = 4\).
Calcule a média e o desvio padrão.

\begin{Shaded}
\begin{Highlighting}[]
\FunctionTok{set.seed}\NormalTok{(}\DecValTok{123}\NormalTok{)}
\NormalTok{amostra }\OtherTok{\textless{}{-}} \FunctionTok{rpois}\NormalTok{(}\DecValTok{2300}\NormalTok{, }\DecValTok{4}\NormalTok{)}
\FunctionTok{mean}\NormalTok{(amostra)}
\FunctionTok{sd}\NormalTok{(amostra)}
\end{Highlighting}
\end{Shaded}

\textbf{4.} Num processo de qualidade, \(X\) é o número de produtos defeituosos num lote de 50,
com probabilidade 0,1 de um produto ser defeituoso. Com \texttt{set.seed(123)}, gere 10 000
observações de \(X\), conte os lotes com exatamente 5 defeituosos e compare a proporção
com \(P(X = 5)\) para \(X \sim \text{Binomial}(50, 0{,}1)\).

\begin{Shaded}
\begin{Highlighting}[]
\FunctionTok{set.seed}\NormalTok{(}\DecValTok{123}\NormalTok{)}
\NormalTok{amostra }\OtherTok{\textless{}{-}} \FunctionTok{rbinom}\NormalTok{(}\DecValTok{10000}\NormalTok{, }\DecValTok{50}\NormalTok{, }\FloatTok{0.1}\NormalTok{)}

\NormalTok{freq\_rela }\OtherTok{\textless{}{-}} \FunctionTok{mean}\NormalTok{(amostra }\SpecialCharTok{==} \DecValTok{5}\NormalTok{)          }\CommentTok{\# proporção empírica}
\NormalTok{prob\_teorica }\OtherTok{\textless{}{-}} \FunctionTok{dbinom}\NormalTok{(}\DecValTok{5}\NormalTok{, }\DecValTok{50}\NormalTok{, }\FloatTok{0.1}\NormalTok{)       }\CommentTok{\# probabilidade teórica}

\FunctionTok{data.frame}\NormalTok{(}\AttributeTok{empirica =}\NormalTok{ freq\_rela, }\AttributeTok{teorica =}\NormalTok{ prob\_teorica)}
\end{Highlighting}
\end{Shaded}

\textbf{5.} Com \texttt{set.seed(123)}, gere 5000 observações de \(X \sim \text{Binomial}(20, 0{,}7)\).

\begin{enumerate}
\def\labelenumi{(\alph{enumi})}
\tightlist
\item
  Faça um histograma de frequências relativas e sobreponha a distribuição de \(X\).
\item
  Use a função de distribuição empírica para estimar \(P(X \leq 10)\) e compare com o
  valor teórico.
\end{enumerate}

\begin{Shaded}
\begin{Highlighting}[]
\FunctionTok{set.seed}\NormalTok{(}\DecValTok{123}\NormalTok{)}
\NormalTok{amostra }\OtherTok{\textless{}{-}} \FunctionTok{rbinom}\NormalTok{(}\DecValTok{5000}\NormalTok{, }\DecValTok{20}\NormalTok{, }\FloatTok{0.7}\NormalTok{)}

\FunctionTok{hist}\NormalTok{(amostra, }\AttributeTok{freq =} \ConstantTok{FALSE}\NormalTok{, }\AttributeTok{col =} \StringTok{"lightblue"}\NormalTok{)}
\FunctionTok{lines}\NormalTok{(}\DecValTok{0}\SpecialCharTok{:}\DecValTok{20}\NormalTok{, }\FunctionTok{dbinom}\NormalTok{(}\DecValTok{0}\SpecialCharTok{:}\DecValTok{20}\NormalTok{, }\DecValTok{20}\NormalTok{, }\FloatTok{0.7}\NormalTok{), }\AttributeTok{col =} \StringTok{"magenta"}\NormalTok{, }\AttributeTok{type =} \StringTok{"p"}\NormalTok{)}

\NormalTok{Fn }\OtherTok{\textless{}{-}} \FunctionTok{ecdf}\NormalTok{(amostra)}
\FunctionTok{data.frame}\NormalTok{(}\AttributeTok{empirica =} \FunctionTok{Fn}\NormalTok{(}\DecValTok{10}\NormalTok{), }\AttributeTok{teorica =} \FunctionTok{pbinom}\NormalTok{(}\DecValTok{10}\NormalTok{, }\DecValTok{20}\NormalTok{, }\FloatTok{0.7}\NormalTok{))}
\end{Highlighting}
\end{Shaded}

\textbf{6.} Com \texttt{set.seed(543)}, gere 2400 observações de \(Y \sim \text{Poisson}(\lambda = 6)\).

\begin{enumerate}
\def\labelenumi{(\alph{enumi})}
\tightlist
\item
  Faça um histograma de frequências relativas e sobreponha a distribuição de \(Y\).
\item
  Use a função de distribuição empírica para estimar \(P(Y > 5)\) e compare com o valor
  teórico.
\end{enumerate}

\begin{Shaded}
\begin{Highlighting}[]
\FunctionTok{set.seed}\NormalTok{(}\DecValTok{543}\NormalTok{)}
\NormalTok{amostra }\OtherTok{\textless{}{-}} \FunctionTok{rpois}\NormalTok{(}\DecValTok{2400}\NormalTok{, }\DecValTok{6}\NormalTok{)}

\FunctionTok{hist}\NormalTok{(amostra, }\AttributeTok{freq =} \ConstantTok{FALSE}\NormalTok{, }\AttributeTok{col =} \StringTok{"lightblue"}\NormalTok{)}
\FunctionTok{lines}\NormalTok{(}\DecValTok{0}\SpecialCharTok{:}\DecValTok{15}\NormalTok{, }\FunctionTok{dpois}\NormalTok{(}\DecValTok{0}\SpecialCharTok{:}\DecValTok{15}\NormalTok{, }\DecValTok{6}\NormalTok{), }\AttributeTok{col =} \StringTok{"magenta"}\NormalTok{, }\AttributeTok{type =} \StringTok{"p"}\NormalTok{)}

\NormalTok{Fn }\OtherTok{\textless{}{-}} \FunctionTok{ecdf}\NormalTok{(amostra)}
\FunctionTok{data.frame}\NormalTok{(}\AttributeTok{empirica =} \DecValTok{1} \SpecialCharTok{{-}} \FunctionTok{Fn}\NormalTok{(}\DecValTok{5}\NormalTok{), }\AttributeTok{teorica =} \FunctionTok{ppois}\NormalTok{(}\DecValTok{5}\NormalTok{, }\DecValTok{6}\NormalTok{, }\AttributeTok{lower.tail =} \ConstantTok{FALSE}\NormalTok{))}
\end{Highlighting}
\end{Shaded}

\textbf{7.} Com \texttt{set.seed(345)}, gere 3450 observações de \(Z \sim U(0, 1)\). Estime
\(P(Z \leq 0{,}5)\) pela função de distribuição empírica e compare com o valor teórico.

\begin{Shaded}
\begin{Highlighting}[]
\FunctionTok{set.seed}\NormalTok{(}\DecValTok{345}\NormalTok{)}
\NormalTok{amostra }\OtherTok{\textless{}{-}} \FunctionTok{runif}\NormalTok{(}\DecValTok{3450}\NormalTok{, }\DecValTok{0}\NormalTok{, }\DecValTok{1}\NormalTok{)}

\FunctionTok{hist}\NormalTok{(amostra, }\AttributeTok{freq =} \ConstantTok{FALSE}\NormalTok{, }\AttributeTok{col =} \StringTok{"lightblue"}\NormalTok{)}
\FunctionTok{curve}\NormalTok{(}\FunctionTok{dunif}\NormalTok{(x, }\DecValTok{0}\NormalTok{, }\DecValTok{1}\NormalTok{), }\AttributeTok{col =} \StringTok{"red"}\NormalTok{, }\AttributeTok{add =} \ConstantTok{TRUE}\NormalTok{)}

\NormalTok{Fn }\OtherTok{\textless{}{-}} \FunctionTok{ecdf}\NormalTok{(amostra)}
\FunctionTok{data.frame}\NormalTok{(}\AttributeTok{empirica =} \FunctionTok{Fn}\NormalTok{(}\FloatTok{0.5}\NormalTok{), }\AttributeTok{teorica =} \FunctionTok{punif}\NormalTok{(}\FloatTok{0.5}\NormalTok{, }\DecValTok{0}\NormalTok{, }\DecValTok{1}\NormalTok{))}
\end{Highlighting}
\end{Shaded}

\textbf{8.} Com \texttt{set.seed(123)}, gere 3467 observações de \(W \sim N(0, 1)\).

\begin{enumerate}
\def\labelenumi{(\alph{enumi})}
\tightlist
\item
  Faça um histograma de frequências relativas e sobreponha a densidade de \(W\).
\item
  Use a função de distribuição empírica para estimar \(P(W > 1)\) e compare com o valor
  teórico.
\end{enumerate}

\textbf{9.} Com \texttt{set.seed(123)}, gere 1234 observações de \(V \sim \text{Exponencial}(\lambda = 0{,}5)\).

\begin{enumerate}
\def\labelenumi{(\alph{enumi})}
\tightlist
\item
  Faça um histograma de frequências relativas e sobreponha a densidade de \(V\).
\item
  Use a função de distribuição empírica para estimar \(P(V > 2)\) e compare com o valor
  teórico.
\end{enumerate}

\textbf{10.} O número de acertos num alvo em 30 tentativas, com probabilidade de acerto 0,4,
é \(X \sim \text{Binomial}(30, 0{,}4)\). Com \texttt{set.seed(123)}, gere uma amostra de dimensão
700.

\begin{enumerate}
\def\labelenumi{(\alph{enumi})}
\tightlist
\item
  Faça um histograma de frequências relativas e sobreponha a distribuição de \(X\).
\item
  Com a função de distribuição empírica, estime \(P(X > 15)\) e calcule também o valor
  teórico.
\end{enumerate}

\textbf{11.} Com \texttt{set.seed(123)}, gere amostras de dimensão crescente
\(n = 100, 1000, 10\,000, 100\,000\) de \(X \sim \text{Poisson}(\lambda = 3)\). Para cada
dimensão, calcule a média amostral e compare-a com o valor esperado teórico. Comente a
convergência (Lei dos Grandes Números).

\begin{Shaded}
\begin{Highlighting}[]
\FunctionTok{set.seed}\NormalTok{(}\DecValTok{123}\NormalTok{)}
\NormalTok{tamanhos }\OtherTok{\textless{}{-}} \FunctionTok{c}\NormalTok{(}\DecValTok{100}\NormalTok{, }\DecValTok{1000}\NormalTok{, }\DecValTok{10000}\NormalTok{, }\DecValTok{100000}\NormalTok{)}
\NormalTok{lambda }\OtherTok{\textless{}{-}} \DecValTok{3}

\NormalTok{media\_amostra }\OtherTok{\textless{}{-}} \FunctionTok{numeric}\NormalTok{(}\FunctionTok{length}\NormalTok{(tamanhos))}
\ControlFlowTok{for}\NormalTok{ (i }\ControlFlowTok{in} \FunctionTok{seq\_along}\NormalTok{(tamanhos)) \{}
\NormalTok{  media\_amostra[i] }\OtherTok{\textless{}{-}} \FunctionTok{mean}\NormalTok{(}\FunctionTok{rpois}\NormalTok{(tamanhos[i], }\AttributeTok{lambda =}\NormalTok{ lambda))}
\NormalTok{\}}

\FunctionTok{data.frame}\NormalTok{(}\AttributeTok{Tamanho =}\NormalTok{ tamanhos, }\AttributeTok{Media\_amostra =}\NormalTok{ media\_amostra, }\AttributeTok{Media\_teorica =}\NormalTok{ lambda)}
\end{Highlighting}
\end{Shaded}

\textbf{12.} Com \texttt{set.seed(123)}, gere amostras de dimensão crescente
\(n = 100, 1000, 10\,000, 100\,000\) de \(W \sim U(0, 1)\). Para cada dimensão, calcule a
média amostral e compare-a com o valor esperado teórico. Comente a convergência.

\textbf{13.} Um grupo de estudantes conduz um estudo sobre a satisfação com um novo curso.
Cada aluno indica se está satisfeito, com probabilidade 0,75. Com \texttt{set.seed(42)}, simule
amostras de dimensão crescente \(n = 100, 500, 1000, 5000, 10\,000\) de
\(X \sim \text{Binomial}\), onde \(X\) é o número de satisfeitos. Para cada dimensão,
calcule a proporção de satisfeitos e compare-a com a probabilidade teórica (0,75).

\textbf{14.} Com \texttt{set.seed(1058)}, gere 9060 amostras de dimensão 9 de
\(X \sim \text{Binomial}(41, 0{,}81)\). Calcule a média de cada amostra, obtendo uma
amostra de médias; calcule ainda \(E(X)\) e compare-o com a média das médias.

\begin{Shaded}
\begin{Highlighting}[]
\FunctionTok{set.seed}\NormalTok{(}\DecValTok{1058}\NormalTok{)}
\NormalTok{m }\OtherTok{\textless{}{-}} \DecValTok{9060}      \CommentTok{\# número de amostras}
\NormalTok{n }\OtherTok{\textless{}{-}} \DecValTok{9}         \CommentTok{\# dimensão de cada amostra}

\NormalTok{amostra }\OtherTok{\textless{}{-}} \FunctionTok{matrix}\NormalTok{(}\FunctionTok{rbinom}\NormalTok{(m }\SpecialCharTok{*}\NormalTok{ n, }\DecValTok{41}\NormalTok{, }\FloatTok{0.81}\NormalTok{), }\AttributeTok{nrow =}\NormalTok{ m, }\AttributeTok{ncol =}\NormalTok{ n)}
\NormalTok{media\_amostra }\OtherTok{\textless{}{-}} \FunctionTok{rowMeans}\NormalTok{(amostra)}

\FunctionTok{data.frame}\NormalTok{(}\AttributeTok{empirica =} \FunctionTok{mean}\NormalTok{(media\_amostra), }\AttributeTok{teorica =} \DecValTok{41} \SpecialCharTok{*} \FloatTok{0.81}\NormalTok{)}
\end{Highlighting}
\end{Shaded}

\textbf{15.} Num hospital, o tempo de atendimento segue uma exponencial de média 30 minutos.
Com \texttt{set.seed(456)}, simule 1000 amostras de dimensão 50, 100 e 1000. Para cada
dimensão, calcule a média de cada amostra, faça o histograma das médias e compare com a
densidade normal de média \(E(X)\) e desvio padrão \(\sqrt{V(X)/n}\). Comente a aplicação do
Teorema do Limite Central.

\textbf{16.} O tempo de espera (minutos) no balcão de informações de um banco segue
\(X \sim U(5, 20)\). Com \texttt{set.seed(1430)}, gere 8000 amostras de dimensão \(n = 100\).

\begin{enumerate}
\def\labelenumi{(\alph{enumi})}
\tightlist
\item
  Calcule a soma de cada amostra, \(S_n = \sum_{i=1}^n X_i\).
\item
  Faça o histograma das somas e sobreponha uma normal de média \(nE(X)\) e desvio padrão
  \(\sqrt{n\,V(X)}\).
\item
  Calcule a média de cada amostra, \(\bar{X}_n\).
\item
  Faça o histograma das médias e sobreponha uma normal de média \(E(X)\) e desvio padrão
  \(\sqrt{V(X)/n}\).
\end{enumerate}

\begin{Shaded}
\begin{Highlighting}[]
\FunctionTok{set.seed}\NormalTok{(}\DecValTok{1430}\NormalTok{)}
\NormalTok{m }\OtherTok{\textless{}{-}} \DecValTok{8000}\NormalTok{; n }\OtherTok{\textless{}{-}} \DecValTok{100}
\NormalTok{a }\OtherTok{\textless{}{-}} \DecValTok{5}\NormalTok{; b }\OtherTok{\textless{}{-}} \DecValTok{20}

\NormalTok{media }\OtherTok{\textless{}{-}}\NormalTok{ (a }\SpecialCharTok{+}\NormalTok{ b) }\SpecialCharTok{/} \DecValTok{2}
\NormalTok{variancia }\OtherTok{\textless{}{-}}\NormalTok{ (b }\SpecialCharTok{{-}}\NormalTok{ a)}\SpecialCharTok{\^{}}\DecValTok{2} \SpecialCharTok{/} \DecValTok{12}

\NormalTok{amostras }\OtherTok{\textless{}{-}} \FunctionTok{matrix}\NormalTok{(}\FunctionTok{runif}\NormalTok{(m }\SpecialCharTok{*}\NormalTok{ n, a, b), }\AttributeTok{nrow =}\NormalTok{ m, }\AttributeTok{ncol =}\NormalTok{ n)}

\CommentTok{\# Soma de cada amostra + curva normal}
\NormalTok{soma\_amostras }\OtherTok{\textless{}{-}} \FunctionTok{rowSums}\NormalTok{(amostras)}
\FunctionTok{hist}\NormalTok{(soma\_amostras, }\AttributeTok{freq =} \ConstantTok{FALSE}\NormalTok{, }\AttributeTok{col =} \StringTok{"lightblue"}\NormalTok{)}
\FunctionTok{curve}\NormalTok{(}\FunctionTok{dnorm}\NormalTok{(x, n }\SpecialCharTok{*}\NormalTok{ media, }\FunctionTok{sqrt}\NormalTok{(variancia }\SpecialCharTok{*}\NormalTok{ n)), }\AttributeTok{col =} \StringTok{"red"}\NormalTok{, }\AttributeTok{add =} \ConstantTok{TRUE}\NormalTok{, }\AttributeTok{lwd =} \DecValTok{2}\NormalTok{)}

\CommentTok{\# Média de cada amostra + curva normal}
\NormalTok{media\_amostras }\OtherTok{\textless{}{-}} \FunctionTok{rowMeans}\NormalTok{(amostras)}
\FunctionTok{hist}\NormalTok{(media\_amostras, }\AttributeTok{freq =} \ConstantTok{FALSE}\NormalTok{, }\AttributeTok{col =} \StringTok{"lightblue"}\NormalTok{)}
\FunctionTok{curve}\NormalTok{(}\FunctionTok{dnorm}\NormalTok{(x, media, }\FunctionTok{sqrt}\NormalTok{(variancia }\SpecialCharTok{/}\NormalTok{ n)), }\AttributeTok{col =} \StringTok{"red"}\NormalTok{, }\AttributeTok{add =} \ConstantTok{TRUE}\NormalTok{, }\AttributeTok{lwd =} \DecValTok{2}\NormalTok{)}
\end{Highlighting}
\end{Shaded}

\textbf{17.} O tempo de atendimento de doentes graves segue
\(X \sim \text{Exponencial}(\lambda = 0{,}21)\). Com \texttt{set.seed(1580)}, gere 1234 amostras
de dimensão \(n = 50\).

\begin{enumerate}
\def\labelenumi{(\alph{enumi})}
\tightlist
\item
  Calcule a soma de cada amostra, \(S_n = \sum_{i=1}^n X_i\).
\item
  Faça o histograma das somas e sobreponha uma normal de média \(nE(X)\) e desvio padrão
  \(\sqrt{n\,V(X)}\).
\item
  Calcule a soma padronizada \(\dfrac{S_n - E(S_n)}{\sqrt{V(S_n)}}\) e faça o histograma,
  sobrepondo uma normal padrão.
\item
  Calcule a média de cada amostra, \(\bar{X}_n\).
\item
  Faça o histograma das médias e sobreponha uma normal de média \(E(X)\) e desvio padrão
  \(\sqrt{V(X)/n}\).
\end{enumerate}

\textbf{18.} A altura (cm) dos alunos de uma escola segue \(X \sim N(\mu = 170, \sigma = 10)\).
Com \texttt{set.seed(678)}, gere 9876 amostras de dimensão \(n = 80\).

\begin{enumerate}
\def\labelenumi{(\alph{enumi})}
\tightlist
\item
  Calcule a soma de cada amostra, \(S_n\).
\item
  Faça o histograma das somas e sobreponha uma normal de média \(nE(X)\) e desvio padrão
  \(\sqrt{n\,V(X)}\).
\item
  Calcule a soma padronizada e faça o histograma, sobrepondo uma normal padrão.
\item
  Calcule a média de cada amostra, \(\bar{X}_n\).
\item
  Faça o histograma das médias e sobreponha uma normal de média \(E(X)\) e desvio padrão
  \(\sqrt{V(X)/n}\).
\item
  Faça o histograma da média padronizada
  \(\dfrac{\bar{X}_n - E(\bar{X}_n)}{\sqrt{V(\bar{X}_n)}}\) e sobreponha uma normal
  padrão.
\end{enumerate}

\textbf{19.} A chegada de clientes a uma loja numa hora, com uma taxa média de 20 clientes por
hora, segue \(X \sim \text{Poisson}(\lambda = 20)\). Com \texttt{set.seed(1222)}, gere 8050
amostras de dimensão 30.

\begin{enumerate}
\def\labelenumi{(\alph{enumi})}
\tightlist
\item
  Calcule a soma de cada amostra, \(S_n\).
\item
  Faça o histograma das somas e sobreponha uma normal de média \(nE(X)\) e desvio padrão
  \(\sqrt{n\,V(X)}\).
\item
  Calcule a soma padronizada e faça o histograma, sobrepondo uma normal padrão.
\item
  Calcule a média de cada amostra, \(\bar{X}_n\).
\item
  Faça o histograma das médias e sobreponha uma normal de média \(E(X)\) e desvio padrão
  \(\sqrt{V(X)/n}\).
\item
  Faça o histograma da média padronizada e sobreponha uma normal padrão.
\end{enumerate}

Os últimos exercícios mostram, por simulação, um dos resultados centrais da estatística.
No próximo capítulo mudamos de perspetiva: em vez de usar as funções \texttt{r}\emph{nome} já feitas,
aprendemos a \textbf{construí-las} --- a gerar nós próprios variáveis aleatórias de qualquer
distribuição, a começar pelo método da transformada inversa.

\chapter{Método da transformada inversa}\label{transformada}

Nos capítulos anteriores gerámos variáveis aleatórias com as funções \texttt{r}\emph{nome} do R ---
\texttt{rbinom}, \texttt{rpois}, \texttt{rnorm}, e as demais. Mas de onde vêm esses números? E o que fazer
quando precisamos de uma distribuição para a qual o R \textbf{não} oferece uma função pronta?
A resposta a ambas as perguntas começa com uma observação notável: se soubermos gerar
números uniformes em \((0, 1)\) --- a matéria-prima de todo o acaso no computador ---, podemos
transformá-los em amostras de \emph{qualquer} distribuição. O \textbf{método da transformada
inversa} é a primeira e mais fundamental dessas transformações.

Vale a pena dominá-lo por três razões: permite gerar amostras de \textbf{distribuições não
implementadas} no R; dá \textbf{flexibilidade} para simulações personalizadas; e, sobretudo,
esclarece os \textbf{fundamentos} da geração de variáveis aleatórias --- implementá-lo à mão é
compreender o que as funções \texttt{r}\emph{nome} fazem por dentro.

Neste capítulo aprendemos o resultado que está na base da geração de variáveis
aleatórias --- a transformada integral de probabilidade --- e aplicamo-lo para gerar, a
partir de números uniformes, variáveis discretas (arbitrárias, Bernoulli, binomial,
geométrica, Poisson) e contínuas (exponencial, uniforme, Weibull, Cauchy, e outras).

\section{O resultado fundamental}\label{o-resultado-fundamental}

Tudo assenta no seguinte teorema, conhecido como \textbf{transformada integral de
probabilidade}.

\textbf{Teorema.} Seja \(F\) uma função de distribuição qualquer e defina-se a sua \emph{inversa
generalizada} por
\[F^{-1}(u) = \inf\{x : F(x) \geq u\}, \qquad 0 < u < 1.\]
Se \(U \sim U(0, 1)\), então a variável aleatória \(X = F^{-1}(U)\) tem função de
distribuição \(F\). Reciprocamente, se \(X\) é contínua com função de distribuição \(F\),
então \(F(X) \sim U(0, 1)\).

A demonstração é curta no caso em que \(F\) é contínua e estritamente crescente (e, por
isso, invertível no sentido usual). A chave é a equivalência \(F^{-1}(u) \leq x \iff u \leq F(x)\).
Então
\[P(X \leq x) = P\big(F^{-1}(U) \leq x\big) = P\big(U \leq F(x)\big) = F(x),\]
onde a última igualdade decorre de \(U \sim U(0, 1)\) e \(0 \leq F(x) \leq 1\) (para uma
uniforme, \(P(U \leq a) = a\)). Logo \(X = F^{-1}(U)\) tem mesmo a distribuição \(F\).

A intuição é geométrica: a função de distribuição \(F\) transforma o eixo dos \(x\) na escala
de probabilidades \([0, 1]\); inverter esse mapa é sortear uma ``altura'' \(U\) uniforme em
\([0, 1]\) e ler qual o \(x\) que lhe corresponde. Regiões onde \(F\) sobe depressa (isto é,
onde a densidade é alta) recebem, naturalmente, mais valores.

Deste teorema resulta uma receita geral:

\textbf{Receita da transformada inversa.}

\begin{enumerate}
\def\labelenumi{\arabic{enumi}.}
\tightlist
\item
  Determine a função de distribuição \(F(x) = P(X \leq x)\).
\item
  Resolva a equação \(u = F(x)\) em ordem a \(x\), obtendo a inversa \(x = F^{-1}(u)\).
\item
  Gere \(U \sim U(0, 1)\) e devolva \(X = F^{-1}(U)\).
\end{enumerate}

Vejamos como isto se concretiza, primeiro para variáveis discretas e depois para
contínuas.

\section{Variáveis aleatórias discretas}\label{variuxe1veis-aleatuxf3rias-discretas}

Suponha que \(X\) toma os valores \(x_0 < x_1 < x_2 < \cdots\) com probabilidades
\(p_i = P(X = x_i)\), \(\sum_i p_i = 1\). A sua função de distribuição é a soma acumulada,
\(F(x_k) = \sum_{i \leq k} p_i\). Aplicar a inversa generalizada significa, aqui, gerar
\(U \sim U(0, 1)\) e devolver
\[X = x_i \quad \text{se} \quad F(x_{i-1}) \leq U < F(x_i),\]
isto é, o menor valor \(x_i\) cuja probabilidade acumulada já ultrapassa \(U\). Que este
procedimento produz a distribuição certa vê-se de imediato: como
\(P(a \leq U < b) = b - a\),
\[P(X = x_i) = P\big(F(x_{i-1}) \leq U < F(x_i)\big) = F(x_i) - F(x_{i-1}) = p_i.\]

\begin{Shaded}
\begin{Highlighting}[]
\CommentTok{\# Esquema do algoritmo}
\CommentTok{\# 1. Gerar u \textasciitilde{} U(0, 1)}
\CommentTok{\# 2. Percorrer os valores, acumulando probabilidades, até F(x\_i) \textgreater{}= u}
\CommentTok{\# 3. Devolver esse x\_i}
\end{Highlighting}
\end{Shaded}

\textbf{Exemplo 1 (distribuição arbitrária).} Seja \(X\) discreta com \(p_1 = 0{,}20\),
\(p_2 = 0{,}15\), \(p_3 = 0{,}25\), \(p_4 = 0{,}40\). A sua função de distribuição é
\[F(x) = \begin{cases} 0, & x < 1 \\ 0{,}20, & 1 \leq x < 2 \\ 0{,}35, & 2 \leq x < 3 \\ 0{,}60, & 3 \leq x < 4 \\ 1, & x \geq 4 \end{cases}\]

Gerando \(U\), devolvemos 1 se \(U < 0{,}20\); 2 se \(0{,}20 \leq U < 0{,}35\); 3 se
\(0{,}35 \leq U < 0{,}60\); e 4 se \(U \geq 0{,}60\):

\begin{Shaded}
\begin{Highlighting}[]
\NormalTok{gerar\_inversa }\OtherTok{\textless{}{-}} \ControlFlowTok{function}\NormalTok{() \{}
\NormalTok{  u }\OtherTok{\textless{}{-}} \FunctionTok{runif}\NormalTok{(}\DecValTok{1}\NormalTok{)}
  \ControlFlowTok{if}\NormalTok{ (u }\SpecialCharTok{\textless{}} \FloatTok{0.20}\NormalTok{)      }\FunctionTok{return}\NormalTok{(}\DecValTok{1}\NormalTok{)}
  \ControlFlowTok{else} \ControlFlowTok{if}\NormalTok{ (u }\SpecialCharTok{\textless{}} \FloatTok{0.35}\NormalTok{) }\FunctionTok{return}\NormalTok{(}\DecValTok{2}\NormalTok{)}
  \ControlFlowTok{else} \ControlFlowTok{if}\NormalTok{ (u }\SpecialCharTok{\textless{}} \FloatTok{0.60}\NormalTok{) }\FunctionTok{return}\NormalTok{(}\DecValTok{3}\NormalTok{)}
  \ControlFlowTok{else}               \FunctionTok{return}\NormalTok{(}\DecValTok{4}\NormalTok{)}
\NormalTok{\}}

\FunctionTok{set.seed}\NormalTok{(}\DecValTok{123}\NormalTok{)}
\NormalTok{valores }\OtherTok{\textless{}{-}} \FunctionTok{replicate}\NormalTok{(}\DecValTok{10000}\NormalTok{, }\FunctionTok{gerar\_inversa}\NormalTok{())}
\FunctionTok{table}\NormalTok{(valores) }\SpecialCharTok{/} \DecValTok{10000}     \CommentTok{\# frequências relativas \textasciitilde{} (0.20, 0.15, 0.25, 0.40)}
\DocumentationTok{\#\# valores}
\DocumentationTok{\#\#      1      2      3      4 }
\DocumentationTok{\#\# 0.1998 0.1499 0.2550 0.3953}
\end{Highlighting}
\end{Shaded}

\textbf{Exemplo 2 (uniforme discreta, algoritmo geral).} Para \(X\) uniforme em
\(\{1, 2, \ldots, 10\}\) (todas as probabilidades iguais a \(1/10\)), o algoritmo geral
percorre os valores acumulando probabilidade até ultrapassar \(U\):

\begin{Shaded}
\begin{Highlighting}[]
\NormalTok{gerar\_va\_inversa }\OtherTok{\textless{}{-}} \ControlFlowTok{function}\NormalTok{() \{}
\NormalTok{  u }\OtherTok{\textless{}{-}} \FunctionTok{runif}\NormalTok{(}\DecValTok{1}\NormalTok{)}
\NormalTok{  p }\OtherTok{\textless{}{-}} \DecValTok{1}\SpecialCharTok{/}\DecValTok{10}       \CommentTok{\# P(X = 1)}
\NormalTok{  F }\OtherTok{\textless{}{-}}\NormalTok{ p          }\CommentTok{\# função de distribuição acumulada}
\NormalTok{  X }\OtherTok{\textless{}{-}} \DecValTok{1}
  \ControlFlowTok{while}\NormalTok{ (u }\SpecialCharTok{\textgreater{}}\NormalTok{ F) \{  }\CommentTok{\# avança enquanto U ainda não foi ultrapassado}
\NormalTok{    X }\OtherTok{\textless{}{-}}\NormalTok{ X }\SpecialCharTok{+} \DecValTok{1}
\NormalTok{    F }\OtherTok{\textless{}{-}}\NormalTok{ F }\SpecialCharTok{+}\NormalTok{ p}
\NormalTok{  \}}
  \FunctionTok{return}\NormalTok{(X)}
\NormalTok{\}}

\FunctionTok{set.seed}\NormalTok{(}\DecValTok{123}\NormalTok{)}
\NormalTok{valores }\OtherTok{\textless{}{-}} \FunctionTok{replicate}\NormalTok{(}\DecValTok{5000}\NormalTok{, }\FunctionTok{gerar\_va\_inversa}\NormalTok{())}
\FunctionTok{barplot}\NormalTok{(}\FunctionTok{table}\NormalTok{(valores) }\SpecialCharTok{/} \DecValTok{5000}\NormalTok{, }\AttributeTok{col =} \StringTok{"lightblue"}\NormalTok{)}
\end{Highlighting}
\end{Shaded}

\begin{figure}

{\centering \includegraphics[width=0.8\linewidth]{introR_files/figure-latex/ti-unifdisc-1} 

}

\caption{Amostra de uma uniforme discreta gerada por transformada inversa.}\label{fig:ti-unifdisc}
\end{figure}

\textbf{Exemplo 3 (Bernoulli).} Para \(X \sim \text{Bernoulli}(p)\), a receita reduz-se a: gerar
\(U\) e devolver 1 se \(U \leq p\), e 0 caso contrário.

\begin{Shaded}
\begin{Highlighting}[]
\NormalTok{gerar\_bernoulli\_inversa }\OtherTok{\textless{}{-}} \ControlFlowTok{function}\NormalTok{(p) \{}
\NormalTok{  U }\OtherTok{\textless{}{-}} \FunctionTok{runif}\NormalTok{(}\DecValTok{1}\NormalTok{)}
  \ControlFlowTok{if}\NormalTok{ (U }\SpecialCharTok{\textless{}=}\NormalTok{ p) }\DecValTok{1} \ControlFlowTok{else} \DecValTok{0}
\NormalTok{\}}

\FunctionTok{set.seed}\NormalTok{(}\DecValTok{123}\NormalTok{)}
\NormalTok{valores }\OtherTok{\textless{}{-}} \FunctionTok{replicate}\NormalTok{(}\DecValTok{100}\NormalTok{, }\FunctionTok{gerar\_bernoulli\_inversa}\NormalTok{(}\FloatTok{0.8}\NormalTok{))}
\FunctionTok{mean}\NormalTok{(valores)     }\CommentTok{\# \textasciitilde{} 0.8}
\DocumentationTok{\#\# [1] 0.82}
\end{Highlighting}
\end{Shaded}

\textbf{Exemplo 4 (binomial).} Uma \(\text{Binomial}(n, p)\) é a soma de \(n\) Bernoullis i.i.d.,
pelo que basta somar \(n\) chamadas ao gerador anterior:

\begin{Shaded}
\begin{Highlighting}[]
\NormalTok{gerar\_binomial\_inversa }\OtherTok{\textless{}{-}} \ControlFlowTok{function}\NormalTok{(n, p) \{}
  \FunctionTok{sum}\NormalTok{(}\FunctionTok{replicate}\NormalTok{(n, }\FunctionTok{gerar\_bernoulli\_inversa}\NormalTok{(p)))}
\NormalTok{\}}

\FunctionTok{set.seed}\NormalTok{(}\DecValTok{123}\NormalTok{)}
\NormalTok{valores }\OtherTok{\textless{}{-}} \FunctionTok{replicate}\NormalTok{(}\DecValTok{10000}\NormalTok{, }\FunctionTok{gerar\_binomial\_inversa}\NormalTok{(}\DecValTok{10}\NormalTok{, }\FloatTok{0.5}\NormalTok{))}
\FunctionTok{hist}\NormalTok{(valores, }\AttributeTok{freq =} \ConstantTok{FALSE}\NormalTok{, }\AttributeTok{breaks =} \DecValTok{11}\NormalTok{, }\AttributeTok{col =} \StringTok{"lightblue"}\NormalTok{,}
     \AttributeTok{main =} \StringTok{"Binomial(10, 0.5) por transformada inversa"}\NormalTok{, }\AttributeTok{xlab =} \StringTok{"x"}\NormalTok{)}
\FunctionTok{points}\NormalTok{(}\DecValTok{0}\SpecialCharTok{:}\DecValTok{10}\NormalTok{, }\FunctionTok{dbinom}\NormalTok{(}\DecValTok{0}\SpecialCharTok{:}\DecValTok{10}\NormalTok{, }\DecValTok{10}\NormalTok{, }\FloatTok{0.5}\NormalTok{), }\AttributeTok{col =} \StringTok{"magenta"}\NormalTok{, }\AttributeTok{pch =} \DecValTok{19}\NormalTok{)}
\end{Highlighting}
\end{Shaded}

\begin{figure}

{\centering \includegraphics[width=0.8\linewidth]{introR_files/figure-latex/ti-binom-1} 

}

\caption{Amostra de uma Binomial(10, 0.5) obtida somando Bernoullis.}\label{fig:ti-binom}
\end{figure}

\textbf{Exemplo 5 (geométrica).} Recorde que a geométrica do R conta o número de \emph{insucessos}
\(Y\) antes do primeiro sucesso, com \(P(Y = y) = p(1-p)^y\) e
\(F(y) = 1 - (1-p)^{y+1}\), \(y \geq 0\). Invertendo \(u = F(y)\) obtém-se um algoritmo direto,
sem ciclos: resolvendo \(1 - (1-p)^{y+1} = u\) chega-se a
\[Y = \left\lceil \frac{\ln(1 - U)}{\ln(1 - p)} \right\rceil - 1.\]

\begin{Shaded}
\begin{Highlighting}[]
\NormalTok{gerar\_geometrica\_inversa }\OtherTok{\textless{}{-}} \ControlFlowTok{function}\NormalTok{(p) \{}
\NormalTok{  U }\OtherTok{\textless{}{-}} \FunctionTok{runif}\NormalTok{(}\DecValTok{1}\NormalTok{)}
  \FunctionTok{ceiling}\NormalTok{(}\FunctionTok{log}\NormalTok{(}\DecValTok{1} \SpecialCharTok{{-}}\NormalTok{ U) }\SpecialCharTok{/} \FunctionTok{log}\NormalTok{(}\DecValTok{1} \SpecialCharTok{{-}}\NormalTok{ p)) }\SpecialCharTok{{-}} \DecValTok{1}
\NormalTok{\}}

\FunctionTok{set.seed}\NormalTok{(}\DecValTok{123}\NormalTok{)}
\NormalTok{valores }\OtherTok{\textless{}{-}} \FunctionTok{replicate}\NormalTok{(}\DecValTok{10000}\NormalTok{, }\FunctionTok{gerar\_geometrica\_inversa}\NormalTok{(}\FloatTok{0.5}\NormalTok{))}
\FunctionTok{hist}\NormalTok{(valores, }\AttributeTok{freq =} \ConstantTok{FALSE}\NormalTok{, }\AttributeTok{col =} \StringTok{"lightblue"}\NormalTok{,}
     \AttributeTok{main =} \StringTok{"Geométrica(0.5) por transformada inversa"}\NormalTok{, }\AttributeTok{xlab =} \StringTok{"x"}\NormalTok{)}
\FunctionTok{points}\NormalTok{(}\DecValTok{0}\SpecialCharTok{:}\DecValTok{10}\NormalTok{, }\FunctionTok{dgeom}\NormalTok{(}\DecValTok{0}\SpecialCharTok{:}\DecValTok{10}\NormalTok{, }\FloatTok{0.5}\NormalTok{), }\AttributeTok{col =} \StringTok{"magenta"}\NormalTok{, }\AttributeTok{pch =} \DecValTok{19}\NormalTok{)}
\end{Highlighting}
\end{Shaded}

\begin{figure}

{\centering \includegraphics[width=0.8\linewidth]{introR_files/figure-latex/ti-geom-1} 

}

\caption{Amostra de uma geométrica (número de insucessos) por transformada inversa.}\label{fig:ti-geom}
\end{figure}

\textbf{Exemplo 6 (Poisson).} Para a Poisson, o algoritmo acumula probabilidades como no
Exemplo 2, mas usa a recursão que liga probabilidades consecutivas,
\[p_{i+1} = \frac{\lambda}{i + 1}\, p_i, \qquad p_0 = e^{-\lambda},\]
o que evita recalcular fatoriais e exponenciais a cada passo:

\begin{Shaded}
\begin{Highlighting}[]
\NormalTok{gerar\_poisson\_inversa }\OtherTok{\textless{}{-}} \ControlFlowTok{function}\NormalTok{(lambda) \{}
\NormalTok{  U }\OtherTok{\textless{}{-}} \FunctionTok{runif}\NormalTok{(}\DecValTok{1}\NormalTok{)}
\NormalTok{  p }\OtherTok{\textless{}{-}} \FunctionTok{exp}\NormalTok{(}\SpecialCharTok{{-}}\NormalTok{lambda)   }\CommentTok{\# P(X = 0)}
\NormalTok{  F }\OtherTok{\textless{}{-}}\NormalTok{ p              }\CommentTok{\# função de distribuição acumulada}
\NormalTok{  X }\OtherTok{\textless{}{-}} \DecValTok{0}
  \ControlFlowTok{while}\NormalTok{ (U }\SpecialCharTok{\textgreater{}}\NormalTok{ F) \{}
\NormalTok{    X }\OtherTok{\textless{}{-}}\NormalTok{ X }\SpecialCharTok{+} \DecValTok{1}
\NormalTok{    p }\OtherTok{\textless{}{-}}\NormalTok{ p }\SpecialCharTok{*}\NormalTok{ lambda }\SpecialCharTok{/}\NormalTok{ X   }\CommentTok{\# recursão: p\_i {-}\textgreater{} p\_\{i+1\}}
\NormalTok{    F }\OtherTok{\textless{}{-}}\NormalTok{ F }\SpecialCharTok{+}\NormalTok{ p}
\NormalTok{  \}}
  \FunctionTok{return}\NormalTok{(X)}
\NormalTok{\}}

\FunctionTok{set.seed}\NormalTok{(}\DecValTok{123}\NormalTok{)}
\NormalTok{valores }\OtherTok{\textless{}{-}} \FunctionTok{replicate}\NormalTok{(}\DecValTok{10000}\NormalTok{, }\FunctionTok{gerar\_poisson\_inversa}\NormalTok{(}\DecValTok{3}\NormalTok{))}
\FunctionTok{hist}\NormalTok{(valores, }\AttributeTok{freq =} \ConstantTok{FALSE}\NormalTok{, }\AttributeTok{breaks =} \DecValTok{12}\NormalTok{, }\AttributeTok{col =} \StringTok{"lightblue"}\NormalTok{,}
     \AttributeTok{main =} \StringTok{"Poisson(3) por transformada inversa"}\NormalTok{, }\AttributeTok{xlab =} \StringTok{"x"}\NormalTok{)}
\FunctionTok{points}\NormalTok{(}\DecValTok{0}\SpecialCharTok{:}\DecValTok{10}\NormalTok{, }\FunctionTok{dpois}\NormalTok{(}\DecValTok{0}\SpecialCharTok{:}\DecValTok{10}\NormalTok{, }\DecValTok{3}\NormalTok{), }\AttributeTok{col =} \StringTok{"magenta"}\NormalTok{, }\AttributeTok{pch =} \DecValTok{19}\NormalTok{)}
\end{Highlighting}
\end{Shaded}

\begin{figure}

{\centering \includegraphics[width=0.8\linewidth]{introR_files/figure-latex/ti-pois-1} 

}

\caption{Amostra de uma Poisson(3) por transformada inversa.}\label{fig:ti-pois}
\end{figure}

\subsection{Exercícios}\label{exercuxedcios-33}

\textbf{1.} Seja \(X\) com \(P(X = 1) = 0{,}2\), \(P(X = 2) = 0{,}5\), \(P(X = 3) = 0{,}3\).

\begin{itemize}
\tightlist
\item
  Calcule a função de distribuição de \(X\).
\item
  Gere 1000 amostras de \(X\) por transformada inversa.
\item
  Faça um gráfico de barras e compare com as probabilidades teóricas.
\end{itemize}

\textbf{2.} Uma moeda enviesada tem probabilidade 0,7 de cara e 0,3 de coroa.

\begin{itemize}
\tightlist
\item
  Defina a função de distribuição da experiência.
\item
  Simule 500 lançamentos por transformada inversa.
\item
  Compare as proporções de caras e coroas com as probabilidades esperadas.
\end{itemize}

\textbf{3.} Seja \(X\) com \(P(X = 1) = 0{,}1\), \(P(X = 2) = 0{,}3\), \(P(X = 3) = 0{,}4\),
\(P(X = 4) = 0{,}2\). Gere 2000 amostras por transformada inversa e compare a distribuição
empírica com a teórica.

\textbf{4.} Seja \(X\) geométrica com \(p = 0{,}2\) (número de provas até ao primeiro sucesso,
\(P(X = k) = (1 - p)^{k-1}p\)). Com \texttt{set.seed(1234)}, gere 1000 valores por transformada
inversa (terá de determinar a função de distribuição de \(X\)) e indique, com 4 casas
decimais, a proporção de valores superiores à soma da média com o desvio padrão
amostrais.

\section{Variáveis aleatórias contínuas}\label{variuxe1veis-aleatuxf3rias-contuxednuas}

Para variáveis contínuas, a receita é ainda mais direta, porque a função de distribuição
\(F\) é, em geral, uma função contínua e estritamente crescente no suporte --- logo,
invertível no sentido habitual.

\textbf{Proposição.} Seja \(X\) uma variável aleatória contínua com função de distribuição \(F\)
invertível, de inversa \(F^{-1}\), e seja \(U \sim U(0, 1)\). Então \(X = F^{-1}(U)\) tem
função de distribuição \(F\).

A demonstração é imediata a partir da equivalência \(F^{-1}(u) \leq x \iff u \leq F(x)\):
\[P\big(F^{-1}(U) \leq x\big) = P\big(U \leq F(x)\big) = F(x),\]
usando, na última igualdade, que \(U \sim U(0, 1)\) e \(0 \leq F(x) \leq 1\). Basta, então,
seguir os três passos da receita: determinar \(F\), invertê-la e aplicar a inversa a um
número uniforme.

\textbf{Exemplo 1 (densidade \(f(x) = 2x\)).} Seja \(X\) com densidade \(f(x) = 2x\) para
\(0 < x < 1\) (e 0 fora). A sua função de distribuição é
\[F(x) = \int_0^x 2t\, dt = x^2, \qquad 0 < x < 1.\]
Como \(F(x) = x^2\) é invertível em \((0, 1)\), com \(F^{-1}(u) = \sqrt{u}\), a receita diz que
\(X = \sqrt{U}\) tem a distribuição pretendida:

\begin{Shaded}
\begin{Highlighting}[]
\FunctionTok{set.seed}\NormalTok{(}\DecValTok{123}\NormalTok{)}
\NormalTok{n }\OtherTok{\textless{}{-}} \DecValTok{10000}
\NormalTok{simlist }\OtherTok{\textless{}{-}} \FunctionTok{sqrt}\NormalTok{(}\FunctionTok{runif}\NormalTok{(n))}

\FunctionTok{hist}\NormalTok{(simlist, }\AttributeTok{prob =} \ConstantTok{TRUE}\NormalTok{, }\AttributeTok{main =} \StringTok{"Densidade f(x) = 2x"}\NormalTok{, }\AttributeTok{xlab =} \StringTok{"x"}\NormalTok{, }\AttributeTok{col =} \StringTok{"lightblue"}\NormalTok{)}
\FunctionTok{curve}\NormalTok{(}\DecValTok{2} \SpecialCharTok{*}\NormalTok{ x, }\DecValTok{0}\NormalTok{, }\DecValTok{1}\NormalTok{, }\AttributeTok{add =} \ConstantTok{TRUE}\NormalTok{, }\AttributeTok{col =} \StringTok{"red"}\NormalTok{, }\AttributeTok{lwd =} \DecValTok{2}\NormalTok{)}
\end{Highlighting}
\end{Shaded}

\begin{figure}

{\centering \includegraphics[width=0.8\linewidth]{introR_files/figure-latex/ti-cont1-1} 

}

\caption{Amostra da densidade f(x) = 2x gerada por X = sqrt(U).}\label{fig:ti-cont1}
\end{figure}

\textbf{Exemplo 2 (uniforme \(U(a, b)\)).} A função de distribuição da uniforme é
\(F(x) = \frac{x - a}{b - a}\), cuja inversa é \(F^{-1}(u) = a + (b - a)u\). Logo,
\(X = a + (b - a)U\):

\begin{Shaded}
\begin{Highlighting}[]
\NormalTok{a }\OtherTok{\textless{}{-}} \SpecialCharTok{{-}}\DecValTok{2}\NormalTok{; b }\OtherTok{\textless{}{-}} \DecValTok{2}\NormalTok{; n }\OtherTok{\textless{}{-}} \DecValTok{10000}
\FunctionTok{set.seed}\NormalTok{(}\DecValTok{123}\NormalTok{)}
\NormalTok{simlist }\OtherTok{\textless{}{-}}\NormalTok{ a }\SpecialCharTok{+}\NormalTok{ (b }\SpecialCharTok{{-}}\NormalTok{ a) }\SpecialCharTok{*} \FunctionTok{runif}\NormalTok{(n)}

\FunctionTok{hist}\NormalTok{(simlist, }\AttributeTok{prob =} \ConstantTok{TRUE}\NormalTok{, }\AttributeTok{main =} \StringTok{"Uniforme({-}2, 2)"}\NormalTok{, }\AttributeTok{xlab =} \StringTok{"x"}\NormalTok{, }\AttributeTok{col =} \StringTok{"lightblue"}\NormalTok{)}
\FunctionTok{curve}\NormalTok{(}\FunctionTok{dunif}\NormalTok{(x, a, b), }\AttributeTok{col =} \StringTok{"red"}\NormalTok{, }\AttributeTok{add =} \ConstantTok{TRUE}\NormalTok{, }\AttributeTok{lwd =} \DecValTok{2}\NormalTok{)}
\end{Highlighting}
\end{Shaded}

\begin{figure}

{\centering \includegraphics[width=0.8\linewidth]{introR_files/figure-latex/ti-cont2-1} 

}

\caption{Amostra de uma Uniforme(-2, 2) por transformada inversa.}\label{fig:ti-cont2}
\end{figure}

\textbf{Exemplo 3 (exponencial).} Seja \(X\) exponencial de taxa 1, com função de distribuição
\[F(x) = 1 - e^{-x}, \qquad x > 0.\]
Fazendo \(u = F(x) = 1 - e^{-x}\) e resolvendo em ordem a \(x\):
\[1 - u = e^{-x} \;\Longrightarrow\; x = -\ln(1 - u).\]
Assim, \(X = -\ln(1 - U)\) gera uma exponencial de taxa 1. Uma pequena simplificação: como
\(1 - U\) também é uniforme em \((0, 1)\), podemos usar \(-\ln(U)\), com a mesma distribuição.
Além disso, se \(X\) é exponencial de taxa 1, então \(\frac{1}{\lambda}X\) é exponencial de
taxa \(\lambda\); logo, uma exponencial de taxa \(\lambda\) gera-se por
\[X = -\frac{1}{\lambda}\ln(U).\]

\begin{Shaded}
\begin{Highlighting}[]
\NormalTok{lambda }\OtherTok{\textless{}{-}} \DecValTok{3}\NormalTok{; n }\OtherTok{\textless{}{-}} \DecValTok{10000}
\FunctionTok{set.seed}\NormalTok{(}\DecValTok{123}\NormalTok{)}
\NormalTok{simlist }\OtherTok{\textless{}{-}} \SpecialCharTok{{-}}\FunctionTok{log}\NormalTok{(}\FunctionTok{runif}\NormalTok{(n)) }\SpecialCharTok{/}\NormalTok{ lambda}

\FunctionTok{hist}\NormalTok{(simlist, }\AttributeTok{probability =} \ConstantTok{TRUE}\NormalTok{, }\AttributeTok{main =} \StringTok{"Exponencial simulada vs teórica"}\NormalTok{,}
     \AttributeTok{xlab =} \StringTok{"x"}\NormalTok{, }\AttributeTok{ylab =} \StringTok{"Densidade"}\NormalTok{, }\AttributeTok{col =} \StringTok{"lightblue"}\NormalTok{, }\AttributeTok{border =} \StringTok{"black"}\NormalTok{)}
\FunctionTok{curve}\NormalTok{(}\FunctionTok{dexp}\NormalTok{(x, }\AttributeTok{rate =}\NormalTok{ lambda), }\AttributeTok{add =} \ConstantTok{TRUE}\NormalTok{, }\AttributeTok{col =} \StringTok{"red"}\NormalTok{, }\AttributeTok{lwd =} \DecValTok{2}\NormalTok{)}
\FunctionTok{legend}\NormalTok{(}\StringTok{"topright"}\NormalTok{, }\AttributeTok{legend =} \FunctionTok{c}\NormalTok{(}\StringTok{"Simulação"}\NormalTok{, }\StringTok{"Teórica"}\NormalTok{),}
       \AttributeTok{col =} \FunctionTok{c}\NormalTok{(}\StringTok{"lightblue"}\NormalTok{, }\StringTok{"red"}\NormalTok{), }\AttributeTok{lwd =} \DecValTok{2}\NormalTok{)}
\end{Highlighting}
\end{Shaded}

\begin{figure}

{\centering \includegraphics[width=0.8\linewidth]{introR_files/figure-latex/ti-cont3-1} 

}

\caption{Amostra de uma Exponencial(3) gerada por X = -ln(U)/lambda.}\label{fig:ti-cont3}
\end{figure}

A recíproca do teorema --- \(F(X) \sim U(0, 1)\) para \(X\) contínua --- é igualmente útil na
prática. É ela que justifica, por exemplo, que os valores-\(p\) de um teste de hipóteses
tenham distribuição uniforme quando a hipótese nula é verdadeira, e está na base de
vários testes de ajustamento de distribuições. Gerar e verificar são, afinal, as duas
faces da mesma transformação.

\subsection{Exercícios}\label{exercuxedcios-34}

\textbf{1.} Para \(X \sim \text{Exponencial}(\lambda = 2)\), com \(F_X(x) = 1 - e^{-\lambda x}\),
mostre que \(X = -\dfrac{\ln(1 - U)}{\lambda}\). Gere 1000 valores por transformada inversa
e visualize a distribuição num histograma.

\textbf{2.} Use a transformada inversa para gerar 500 valores de \(X \sim U(2, 5)\). A função de
distribuição de uma \(U(a, b)\) é \(F_X(x) = \dfrac{x - a}{b - a}\); mostre que
\(X = a + (b - a)U\). Visualize num histograma e compare com a densidade teórica.

\textbf{3.} Para a distribuição de \textbf{Cauchy} de parâmetros 0 e 1, com
\(F_X(x) = \dfrac{1}{\pi}\arctan\!\left(\dfrac{x - x_0}{\gamma}\right) + \dfrac{1}{2}\),
mostre que \(X = x_0 + \gamma\,\tan\!\big[\pi(U - 0{,}5)\big]\). Gere 1000 valores,
visualize num histograma e compare com a densidade teórica (\texttt{dcauchy()}).

\textbf{4.} Para a distribuição de \textbf{Weibull} \(W(\alpha, \beta)\), com densidade
\(f(x) = \alpha \beta^{-\alpha} x^{\alpha - 1} e^{-(x/\beta)^\alpha}\) (\(x > 0\)) e função
de distribuição \(F(x) = 1 - e^{-(x/\beta)^\alpha}\), mostre que
\(X = \beta\,[-\ln(U)]^{1/\alpha}\). Gere 10 000 valores de uma \(W(2, 3)\), visualize num
histograma e compare com a densidade teórica (\texttt{dweibull()}).

\textbf{5.} Obtenha o gerador de números aleatórios das distribuições seguintes pelo método
da transformada inversa:

\begin{enumerate}
\def\labelenumi{(\alph{enumi})}
\tightlist
\item
  \(F(x) = 1 - (x - 1)^2\), \(\;0 < x < 1\);
\item
  \(F(x) = x^\theta\), \(\;0 < x < 1\), \(\theta > 1\);
\item
  \(F(x) = 1 - e^{-x^2 / (2\tau^2)}\), \(\;x > 0\), \(\tau > 0\);
\item
  \(f(x) = \tfrac{1}{2}\sin(x)\), \(\;0 < x < \pi\);
\item
  \(f(x) = \dfrac{\sec^2(x)}{\sqrt{3}}\), \(\;0 < x < \pi/3\).
\end{enumerate}

\textbf{Soluções.}

\begin{enumerate}
\def\labelenumi{(\alph{enumi})}
\tightlist
\item
  \(F^{-1}(u) = 1 - \sqrt{1 - u}\) (ou, por simetria, \(1 - \sqrt{u}\)).
\item
  \(F^{-1}(u) = u^{1/\theta}\).
\item
  \(F^{-1}(u) = \tau\sqrt{-2\ln(1 - u)}\).
\item
  \(F(x) = \dfrac{1 - \cos(x)}{2}\) e \(F^{-1}(u) = \arccos(1 - 2u)\).
\item
  \(F(x) = \dfrac{\tan(x)}{\sqrt{3}}\) e \(F^{-1}(u) = \arctan(u\sqrt{3})\).
\end{enumerate}

\section{Quando usar o método?}\label{quando-usar-o-muxe9todo}

Sempre que for possível, o método da transformada inversa é a primeira escolha: é exato,
simples e usa um único número uniforme por observação. A sua aplicabilidade depende,
porém, de conseguirmos a inversa \(F^{-1}\):

\begin{itemize}
\tightlist
\item
  funciona diretamente quando \(F^{-1}(u)\) tem forma fechada (foi o caso de todos os
  exemplos deste capítulo);
\item
  em muitas distribuições, \(F\) existe mas \textbf{não é invertível analiticamente} (a normal é
  o exemplo mais conhecido: não há fórmula fechada para \(F^{-1}\));
\item
  noutras, nem sequer dispomos de \(F\) em forma fechada, apenas da densidade \(f\).
\end{itemize}

Nesses casos, há alternativas: aproximar numericamente \(F\) e \(F^{-1}\), ou recorrer a
métodos que dispensam a inversa e trabalham diretamente com a densidade. É a um desses
métodos --- o de \textbf{aceitação-rejeição} --- que dedicamos o próximo capítulo.

\chapter{Método da aceitação-rejeição}\label{aceitacao-rejeicao}

No capítulo anterior vimos que o método da transformação inversa nos permite
gerar observações de uma variável aleatória sempre que conhecemos --- e sabemos
inverter --- a sua função de distribuição. Essa é uma condição forte. Para muitas
distribuições de grande importância prática, a função de distribuição não tem
forma fechada (basta pensar na normal) ou, tendo-a, não é invertível de forma
explícita. Precisamos, por isso, de uma segunda ferramenta, mais geral, que
funcione mesmo quando só conhecemos a \textbf{função densidade} \(f\) e não a sua
primitiva. Essa ferramenta é o \textbf{método da aceitação-rejeição}.

A ideia é engenhosamente simples: se não sabemos amostrar diretamente de \(f\),
amostramos de uma distribuição \emph{auxiliar} que sabemos gerar com facilidade e,
em seguida, ``filtramos'' essas observações, aceitando umas e rejeitando outras,
de modo que as sobreviventes se distribuam exatamente segundo \(f\). É como
esculpir a distribuição desejada a partir de um bloco de material mais fácil de
obter.

\section{Contexto histórico}\label{contexto-histuxf3rico}

O método da aceitação-rejeição deve-se a \textbf{John von Neumann}, que o descreveu
em 1951 no artigo \emph{Various techniques used in connection with random digits},
publicado no âmbito dos primeiros trabalhos sobre o método de Monte Carlo. Von
Neumann integrava, juntamente com Stanislaw Ulam e Nicholas Metropolis, o grupo
de Los Alamos que, no final da década de 1940, desenvolveu a simulação
estocástica como instrumento de cálculo científico --- precisamente na altura em
que surgiam os primeiros computadores eletrónicos capazes de produzir e
manipular grandes quantidades de números pseudoaleatórios.

O que torna o método notável é a sua generalidade e a elegância do argumento
que o sustenta. Não exige o conhecimento da função de distribuição, nem sequer
que a densidade \(f\) esteja normalizada (mais à frente veremos por que razão
isto é útil): basta saber avaliar \(f\) ponto a ponto e dispor de uma
distribuição auxiliar cómoda. Mais de setenta anos depois, continua a ser um
dos pilares da simulação e o ponto de partida conceptual para técnicas
modernas, como os métodos de Monte Carlo por cadeias de Markov (MCMC).

\section{A ideia geométrica}\label{a-ideia-geomuxe9trica}

Antes da formalização, vale a pena fixar a intuição. Suponhamos que queremos
gerar pontos sob o gráfico de uma densidade \(f\). Se conseguíssemos distribuir
pontos \textbf{uniformemente} na região delimitada pelo eixo das abcissas e pela
curva \(y = f(x)\), então a abcissa \(X\) de um ponto escolhido ao acaso nessa
região teria exatamente densidade \(f\). Este é o princípio fundamental: gerar
uniformidade numa região a duas dimensões e ler a primeira coordenada.

O problema é que raramente sabemos amostrar uniformemente sob uma curva
arbitrária. A solução de von Neumann é cobrir essa região com outra, maior e
mais simples --- a região sob uma curva \(c\,g(x)\), onde \(g\) é uma densidade fácil
de simular e \(c\) é uma constante que garante que \(c\,g\) fica sempre \emph{acima} de
\(f\). Amostramos uniformemente sob a curva mais fácil e ficamos apenas com os
pontos que caem também sob \(f\). Os restantes são descartados.

Aprender a gerar observações de uma densidade \(f\) da qual não sabemos amostrar
diretamente, recorrendo a uma densidade auxiliar \(g\) de fácil simulação, através
do critério de aceitação-rejeição de von Neumann. No fim do capítulo, deverá ser
capaz de escolher uma proposta \(g\) adequada, determinar a constante \(c\) e
implementar o algoritmo em R.

\section{Formalização do método}\label{formalizauxe7uxe3o-do-muxe9todo}

Sejam \(f\) a densidade da qual pretendemos amostrar (a densidade \emph{alvo}) e \(g\)
uma densidade auxiliar, dita \textbf{densidade de proposta} ou \emph{instrumental}, que
sabemos simular. A única exigência é que exista uma constante \(c \geq 1\) tal que

\[
f(x) \leq c\,g(x), \qquad \text{para todo o } x \text{ com } f(x) > 0.
\]

Equivalentemente, exige-se que o rácio \(f(x)/g(x)\) seja limitado e que

\[
c = \sup_{x \,:\, f(x) > 0} \frac{f(x)}{g(x)} < \infty.
\]

Repare-se numa condição implícita mas essencial: o \textbf{suporte} de \(g\) tem de
conter o suporte de \(f\). Se existir algum ponto onde \(f(x) > 0\) mas \(g(x) = 0\),
o rácio explode e o método não é aplicável --- nunca poderíamos gerar
observações nessa zona porque a proposta lá não coloca massa.

A desigualdade \(f(x) \leq c\,g(x)\) tem de valer em \textbf{todo} o suporte de \(f\),
não apenas na região onde a maioria da massa se concentra. Se escolhermos um
\(c\) demasiado pequeno, a envolvente \(c\,g\) deixa de cobrir \(f\) nalgum ponto e as
observações produzidas seguem uma distribuição \emph{errada}, sem que qualquer
mensagem de erro nos alerte. É um erro silencioso e, por isso, particularmente
perigoso.

\subsection{O algoritmo}\label{o-algoritmo}

O procedimento de aceitação-rejeição consiste em repetir os seguintes passos:

\begin{enumerate}
\def\labelenumi{\arabic{enumi}.}
\tightlist
\item
  Gerar uma observação \(Y\) da densidade de proposta \(g\).
\item
  Gerar \(U \sim \mathcal{U}(0,1)\), independente de \(Y\).
\item
  Se
  \[
  U \leq \frac{f(Y)}{c\,g(Y)},
  \]
  \textbf{aceitar} e devolver \(X = Y\). Caso contrário, \textbf{rejeitar} \(Y\) e voltar ao
  passo 1.
\end{enumerate}

O quociente \(f(Y)/[c\,g(Y)]\) está sempre entre \(0\) e \(1\), pelo que a comparação
com um uniforme faz sentido. Intuitivamente, um valor \(Y\) é aceite com uma
probabilidade tanto maior quanto mais ``alta'' a densidade alvo estiver nesse
ponto, relativamente à envolvente. Onde \(f\) e \(c\,g\) quase coincidem, quase
todos os \(Y\) são aceites; onde \(f\) é muito menor que \(c\,g\), a maioria é
descartada.

\subsection{Demonstração de que o método é correto}\label{demonstrauxe7uxe3o-de-que-o-muxe9todo-uxe9-correto}

Vejamos por que razão as observações aceites têm efetivamente densidade \(f\).
Um valor devolvido pelo algoritmo é um valor de \(Y\) \emph{condicionado} ao
acontecimento ``aceitar''. Calculemos primeiro a probabilidade de aceitação num
único passo, para um dado \(Y = y\):

\[
P\!\left(U \leq \frac{f(Y)}{c\,g(Y)} \,\Big|\, Y = y\right)
= \frac{f(y)}{c\,g(y)} .
\]

Integrando sobre a distribuição de \(Y\), a probabilidade \emph{incondicional} de
aceitação em cada passo é

\[
p = P(\text{aceitar})
= \int \frac{f(y)}{c\,g(y)}\, g(y)\, dy
= \frac{1}{c}\int f(y)\, dy
= \frac{1}{c}.
\]

Agora, a função de distribuição do valor devolvido é a de \(Y\) condicionada à
aceitação. Para qualquer \(x\),

\[
P(X \leq x)
= P\!\left(Y \leq x \,\middle|\, \text{aceitar}\right)
= \frac{P(Y \leq x,\ \text{aceitar})}{P(\text{aceitar})}.
\]

O numerador é

\[
P(Y \leq x,\ \text{aceitar})
= \int_{-\infty}^{x} \frac{f(y)}{c\,g(y)}\, g(y)\, dy
= \frac{1}{c}\int_{-\infty}^{x} f(y)\, dy
= \frac{F(x)}{c},
\]

onde \(F\) é a função de distribuição associada a \(f\). Dividindo pelo
denominador \(p = 1/c\), obtém-se

\[
P(X \leq x) = \frac{F(x)/c}{1/c} = F(x).
\]

Ou seja, o valor aceite tem exatamente a função de distribuição \(F\), logo a
densidade \(f\). É este o resultado que legitima o método. \(\qquad\blacksquare\)

\subsection{\texorpdfstring{Eficiência: o papel da constante \(c\)}{Eficiência: o papel da constante c}}\label{eficiuxeancia-o-papel-da-constante-c}

A demonstração revelou um facto de enorme importância prática: a probabilidade
de aceitar em cada tentativa é \(p = 1/c\). O número \(N\) de tentativas necessárias
até obter a primeira aceitação segue, portanto, uma distribuição \textbf{geométrica}
de parâmetro \(1/c\), e o seu valor esperado é

\[
E[N] = \frac{1}{p} = c.
\]

Por outras palavras, em média são precisas \(c\) tentativas para produzir uma
única observação aceite. A constante \(c\) é assim, simultaneamente, o \emph{fator de
sobre-cobertura} geométrico e o \emph{custo médio} do algoritmo.

Como \(c \geq 1\) sempre (porque \(f\) e \(g\) são densidades e integram \(1\)), o
melhor cenário possível é \(c = 1\), que só ocorre quando \(f = g\). Na prática,
quanto mais a proposta \(g\) se ``assemelhar'' à alvo \(f\), menor será \(c\) e mais
eficiente o método. A arte de aplicar a aceitação-rejeição está, em boa medida,
na escolha de uma boa proposta.

Este é o principal critério para comparar propostas: entre duas densidades
auxiliares admissíveis, prefere-se a que conduz ao menor \(c\). Uma proposta
cómoda de simular mas que force um \(c\) muito grande pode ser, no cômputo geral,
pior do que uma proposta um pouco mais trabalhosa mas com \(c\) próximo de \(1\).

\section{Implementação genérica em R}\label{implementauxe7uxe3o-genuxe9rica-em-r}

Comecemos por traduzir o algoritmo numa função reutilizável. A função recebe o
número de observações desejadas, a densidade alvo \texttt{f}, uma função \texttt{g\_amostra}
que gera observações da proposta, a densidade da proposta \texttt{g\_dens} e a constante
\texttt{c}. Para além da amostra, devolve a taxa de aceitação observada, que nos
permite verificar empiricamente se ela se aproxima de \(1/c\).

\begin{Shaded}
\begin{Highlighting}[]
\NormalTok{aceita\_rejeita }\OtherTok{\textless{}{-}} \ControlFlowTok{function}\NormalTok{(n, f, g\_amostra, g\_dens, c) \{}
\NormalTok{  amostra    }\OtherTok{\textless{}{-}} \FunctionTok{numeric}\NormalTok{(n)}
\NormalTok{  aceites    }\OtherTok{\textless{}{-}} \DecValTok{0}
\NormalTok{  tentativas }\OtherTok{\textless{}{-}} \DecValTok{0}
  \ControlFlowTok{while}\NormalTok{ (aceites }\SpecialCharTok{\textless{}}\NormalTok{ n) \{}
\NormalTok{    y }\OtherTok{\textless{}{-}} \FunctionTok{g\_amostra}\NormalTok{(}\DecValTok{1}\NormalTok{)              }\CommentTok{\# observação da proposta}
\NormalTok{    u }\OtherTok{\textless{}{-}} \FunctionTok{runif}\NormalTok{(}\DecValTok{1}\NormalTok{)                  }\CommentTok{\# uniforme para o teste}
\NormalTok{    tentativas }\OtherTok{\textless{}{-}}\NormalTok{ tentativas }\SpecialCharTok{+} \DecValTok{1}
    \ControlFlowTok{if}\NormalTok{ (u }\SpecialCharTok{\textless{}=} \FunctionTok{f}\NormalTok{(y) }\SpecialCharTok{/}\NormalTok{ (c }\SpecialCharTok{*} \FunctionTok{g\_dens}\NormalTok{(y))) \{}
\NormalTok{      aceites }\OtherTok{\textless{}{-}}\NormalTok{ aceites }\SpecialCharTok{+} \DecValTok{1}
\NormalTok{      amostra[aceites] }\OtherTok{\textless{}{-}}\NormalTok{ y}
\NormalTok{    \}}
\NormalTok{  \}}
  \FunctionTok{list}\NormalTok{(}\AttributeTok{amostra =}\NormalTok{ amostra,}
       \AttributeTok{taxa\_aceitacao =}\NormalTok{ aceites }\SpecialCharTok{/}\NormalTok{ tentativas)}
\NormalTok{\}}
\end{Highlighting}
\end{Shaded}

A versão acima gera uma observação de cada vez, o que é claro do ponto de vista
didático mas relativamente lento em R, uma linguagem em que os ciclos explícitos
têm um custo apreciável. Numa aplicação em que a eficiência importe, é
preferível \textbf{vetorizar}: gerar um grande lote de candidatos e respetivos
uniformes de uma só vez, aceitar em bloco os que passam no teste, e repetir
apenas se ainda faltarem observações. A lógica é idêntica; muda só a forma de a
executar.

\section{Exemplo 1 --- uma densidade limitada com proposta uniforme}\label{exemplo-1-uma-densidade-limitada-com-proposta-uniforme}

Consideremos a densidade da distribuição \(\text{Beta}(2,2)\), dada por

\[
f(x) = 6\,x\,(1-x), \qquad 0 < x < 1,
\]

e nula fora de \((0,1)\). Esta densidade tem suporte limitado, pelo que uma
escolha natural para a proposta é a uniforme em \((0,1)\), ou seja, \(g(x) = 1\)
nesse intervalo. Como \(g\) é constante, determinar \(c\) reduz-se a encontrar o
máximo de \(f\):

\[
c = \sup_{0<x<1} \frac{f(x)}{g(x)} = \max_{0<x<1} 6\,x\,(1-x).
\]

Derivando e igualando a zero, o máximo ocorre em \(x = 1/2\), onde
\(f(1/2) = 6 \cdot \tfrac{1}{2} \cdot \tfrac{1}{2} = \tfrac{3}{2}\). Logo
\(c = 1{,}5\), e esperamos uma taxa de aceitação de \(1/c \approx 0{,}667\).

\begin{Shaded}
\begin{Highlighting}[]
\FunctionTok{set.seed}\NormalTok{(}\DecValTok{2024}\NormalTok{)}

\NormalTok{f\_beta }\OtherTok{\textless{}{-}} \ControlFlowTok{function}\NormalTok{(x) }\FunctionTok{ifelse}\NormalTok{(x }\SpecialCharTok{\textgreater{}} \DecValTok{0} \SpecialCharTok{\&}\NormalTok{ x }\SpecialCharTok{\textless{}} \DecValTok{1}\NormalTok{, }\DecValTok{6} \SpecialCharTok{*}\NormalTok{ x }\SpecialCharTok{*}\NormalTok{ (}\DecValTok{1} \SpecialCharTok{{-}}\NormalTok{ x), }\DecValTok{0}\NormalTok{)}
\NormalTok{g\_unif }\OtherTok{\textless{}{-}} \ControlFlowTok{function}\NormalTok{(x) }\FunctionTok{ifelse}\NormalTok{(x }\SpecialCharTok{\textgreater{}} \DecValTok{0} \SpecialCharTok{\&}\NormalTok{ x }\SpecialCharTok{\textless{}} \DecValTok{1}\NormalTok{, }\DecValTok{1}\NormalTok{, }\DecValTok{0}\NormalTok{)}

\NormalTok{resultado }\OtherTok{\textless{}{-}} \FunctionTok{aceita\_rejeita}\NormalTok{(}
  \AttributeTok{n         =} \DecValTok{10000}\NormalTok{,}
  \AttributeTok{f         =}\NormalTok{ f\_beta,}
  \AttributeTok{g\_amostra =} \ControlFlowTok{function}\NormalTok{(n) }\FunctionTok{runif}\NormalTok{(n),}
  \AttributeTok{g\_dens    =}\NormalTok{ g\_unif,}
  \AttributeTok{c         =} \FloatTok{1.5}
\NormalTok{)}

\NormalTok{resultado}\SpecialCharTok{$}\NormalTok{taxa\_aceitacao      }\CommentTok{\# deverá rondar 0.667}
\DocumentationTok{\#\# [1] 0.6674}
\end{Highlighting}
\end{Shaded}

A taxa de aceitação observada está muito próxima do valor teórico \(1/1{,}5\), o
que confirma o funcionamento do algoritmo. Comparemos agora o histograma das
observações geradas com a densidade teórica.

\begin{Shaded}
\begin{Highlighting}[]
\FunctionTok{hist}\NormalTok{(resultado}\SpecialCharTok{$}\NormalTok{amostra, }\AttributeTok{breaks =} \DecValTok{40}\NormalTok{, }\AttributeTok{probability =} \ConstantTok{TRUE}\NormalTok{,}
     \AttributeTok{col =} \StringTok{"grey85"}\NormalTok{, }\AttributeTok{border =} \StringTok{"white"}\NormalTok{,}
     \AttributeTok{main =} \StringTok{""}\NormalTok{, }\AttributeTok{xlab =} \StringTok{"x"}\NormalTok{, }\AttributeTok{ylab =} \StringTok{"densidade"}\NormalTok{)}
\FunctionTok{curve}\NormalTok{(f\_beta, }\AttributeTok{from =} \DecValTok{0}\NormalTok{, }\AttributeTok{to =} \DecValTok{1}\NormalTok{, }\AttributeTok{add =} \ConstantTok{TRUE}\NormalTok{, }\AttributeTok{lwd =} \DecValTok{2}\NormalTok{, }\AttributeTok{col =} \StringTok{"\#10656d"}\NormalTok{)}
\FunctionTok{legend}\NormalTok{(}\StringTok{"topright"}\NormalTok{, }\AttributeTok{legend =} \StringTok{"densidade Beta(2,2)"}\NormalTok{,}
       \AttributeTok{lwd =} \DecValTok{2}\NormalTok{, }\AttributeTok{col =} \StringTok{"\#10656d"}\NormalTok{, }\AttributeTok{bty =} \StringTok{"n"}\NormalTok{)}
\end{Highlighting}
\end{Shaded}

\begin{figure}

{\centering \includegraphics[width=0.8\linewidth]{introR_files/figure-latex/ar-beta-hist-1} 

}

\caption{Histograma de 10\,000 observações geradas pelo método da aceitação-rejeição, sobreposto à densidade teórica da Beta(2,2).}\label{fig:ar-beta-hist}
\end{figure}

O acordo entre o histograma e a curva teórica é excelente. Repare-se ainda que,
tendo já implementado o método, poderíamos ter obtido a mesma amostra com
\texttt{rbeta(10000,\ 2,\ 2)}; o exercício serve para \emph{validar} a nossa implementação
contra uma referência de confiança, uma boa prática sempre que se programa um
gerador de raiz.

\section{Exemplo 2 --- a normal padrão a partir da exponencial}\label{exemplo-2-a-normal-padruxe3o-a-partir-da-exponencial}

O exemplo anterior era confortável porque a densidade alvo tinha suporte
limitado. Vejamos agora um caso com suporte ilimitado e uma proposta menos
óbvia, que ilustra bem a subtileza do método: gerar observações da \textbf{normal
padrão} \(N(0,1)\).

A estratégia clássica passa por gerar primeiro o \textbf{valor absoluto} de uma
normal padrão --- que segue a chamada distribuição \emph{meia-normal}, com densidade

\[
f(x) = \sqrt{\frac{2}{\pi}}\; e^{-x^2/2}, \qquad x \geq 0,
\]

e, no fim, atribuir-lhe um sinal aleatório, positivo ou negativo com igual
probabilidade. Como proposta para a meia-normal usamos a exponencial de
parâmetro \(1\), \(g(x) = e^{-x}\) para \(x \geq 0\), que sabemos simular facilmente
pela transformação inversa.

Determinemos \(c\). O rácio das densidades é

\[
\frac{f(x)}{g(x)} = \sqrt{\frac{2}{\pi}}\; e^{-x^2/2 + x}
= \sqrt{\frac{2}{\pi}}\; e^{\,1/2}\, e^{-(x-1)^2/2},
\]

onde completámos o quadrado no expoente. O fator \(e^{-(x-1)^2/2}\) é máximo em
\(x = 1\), onde vale \(1\). Portanto

\[
c = \sqrt{\frac{2}{\pi}}\; e^{\,1/2} = \sqrt{\frac{2e}{\pi}} \approx 1{,}3155.
\]

A taxa de aceitação esperada é \(1/c \approx 0{,}760\) --- bastante boa. E a
condição de aceitação simplifica-se de forma elegante:

\[
U \leq \frac{f(Y)}{c\,g(Y)} = e^{-(Y-1)^2/2}.
\]

\begin{Shaded}
\begin{Highlighting}[]
\FunctionTok{set.seed}\NormalTok{(}\DecValTok{2024}\NormalTok{)}

\NormalTok{c\_norm }\OtherTok{\textless{}{-}} \FunctionTok{sqrt}\NormalTok{(}\DecValTok{2} \SpecialCharTok{*} \FunctionTok{exp}\NormalTok{(}\DecValTok{1}\NormalTok{) }\SpecialCharTok{/}\NormalTok{ pi)}
\NormalTok{c\_norm}
\DocumentationTok{\#\# [1] 1.315}

\NormalTok{gerar\_normal\_ar }\OtherTok{\textless{}{-}} \ControlFlowTok{function}\NormalTok{(n) \{}
\NormalTok{  x }\OtherTok{\textless{}{-}} \FunctionTok{numeric}\NormalTok{(n)}
\NormalTok{  aceites }\OtherTok{\textless{}{-}} \DecValTok{0}
\NormalTok{  tentativas }\OtherTok{\textless{}{-}} \DecValTok{0}
  \ControlFlowTok{while}\NormalTok{ (aceites }\SpecialCharTok{\textless{}}\NormalTok{ n) \{}
\NormalTok{    y }\OtherTok{\textless{}{-}} \FunctionTok{rexp}\NormalTok{(}\DecValTok{1}\NormalTok{, }\AttributeTok{rate =} \DecValTok{1}\NormalTok{)              }\CommentTok{\# proposta: Exp(1)}
\NormalTok{    u }\OtherTok{\textless{}{-}} \FunctionTok{runif}\NormalTok{(}\DecValTok{1}\NormalTok{)}
\NormalTok{    tentativas }\OtherTok{\textless{}{-}}\NormalTok{ tentativas }\SpecialCharTok{+} \DecValTok{1}
    \ControlFlowTok{if}\NormalTok{ (u }\SpecialCharTok{\textless{}=} \FunctionTok{exp}\NormalTok{(}\SpecialCharTok{{-}}\NormalTok{(y }\SpecialCharTok{{-}} \DecValTok{1}\NormalTok{)}\SpecialCharTok{\^{}}\DecValTok{2} \SpecialCharTok{/} \DecValTok{2}\NormalTok{)) \{     }\CommentTok{\# teste simplificado}
\NormalTok{      aceites }\OtherTok{\textless{}{-}}\NormalTok{ aceites }\SpecialCharTok{+} \DecValTok{1}
\NormalTok{      s }\OtherTok{\textless{}{-}} \FunctionTok{sample}\NormalTok{(}\FunctionTok{c}\NormalTok{(}\SpecialCharTok{{-}}\DecValTok{1}\NormalTok{, }\DecValTok{1}\NormalTok{), }\DecValTok{1}\NormalTok{)          }\CommentTok{\# sinal aleatório}
\NormalTok{      x[aceites] }\OtherTok{\textless{}{-}}\NormalTok{ s }\SpecialCharTok{*}\NormalTok{ y}
\NormalTok{    \}}
\NormalTok{  \}}
  \FunctionTok{list}\NormalTok{(}\AttributeTok{amostra =}\NormalTok{ x, }\AttributeTok{taxa\_aceitacao =}\NormalTok{ aceites }\SpecialCharTok{/}\NormalTok{ tentativas)}
\NormalTok{\}}

\NormalTok{normais }\OtherTok{\textless{}{-}} \FunctionTok{gerar\_normal\_ar}\NormalTok{(}\DecValTok{10000}\NormalTok{)}
\NormalTok{normais}\SpecialCharTok{$}\NormalTok{taxa\_aceitacao                  }\CommentTok{\# deverá rondar 0.760}
\DocumentationTok{\#\# [1] 0.7596}
\end{Highlighting}
\end{Shaded}

\begin{Shaded}
\begin{Highlighting}[]
\FunctionTok{hist}\NormalTok{(normais}\SpecialCharTok{$}\NormalTok{amostra, }\AttributeTok{breaks =} \DecValTok{50}\NormalTok{, }\AttributeTok{probability =} \ConstantTok{TRUE}\NormalTok{,}
     \AttributeTok{col =} \StringTok{"grey85"}\NormalTok{, }\AttributeTok{border =} \StringTok{"white"}\NormalTok{,}
     \AttributeTok{main =} \StringTok{""}\NormalTok{, }\AttributeTok{xlab =} \StringTok{"x"}\NormalTok{, }\AttributeTok{ylab =} \StringTok{"densidade"}\NormalTok{)}
\FunctionTok{curve}\NormalTok{(}\FunctionTok{dnorm}\NormalTok{(x), }\AttributeTok{from =} \SpecialCharTok{{-}}\DecValTok{4}\NormalTok{, }\AttributeTok{to =} \DecValTok{4}\NormalTok{, }\AttributeTok{add =} \ConstantTok{TRUE}\NormalTok{, }\AttributeTok{lwd =} \DecValTok{2}\NormalTok{, }\AttributeTok{col =} \StringTok{"\#10656d"}\NormalTok{)}
\FunctionTok{legend}\NormalTok{(}\StringTok{"topright"}\NormalTok{, }\AttributeTok{legend =} \StringTok{"densidade N(0,1)"}\NormalTok{,}
       \AttributeTok{lwd =} \DecValTok{2}\NormalTok{, }\AttributeTok{col =} \StringTok{"\#10656d"}\NormalTok{, }\AttributeTok{bty =} \StringTok{"n"}\NormalTok{)}
\end{Highlighting}
\end{Shaded}

\begin{figure}

{\centering \includegraphics[width=0.8\linewidth]{introR_files/figure-latex/ar-normal-hist-1} 

}

\caption{Observações da normal padrão geradas por aceitação-rejeição a partir de uma proposta exponencial, com a densidade N(0,1) sobreposta.}\label{fig:ar-normal-hist}
\end{figure}

Uma verificação numérica adicional: a média e o desvio-padrão amostrais devem
aproximar-se de \(0\) e \(1\), respetivamente.

\begin{Shaded}
\begin{Highlighting}[]
\FunctionTok{c}\NormalTok{(}\AttributeTok{media =} \FunctionTok{mean}\NormalTok{(normais}\SpecialCharTok{$}\NormalTok{amostra), }\AttributeTok{desvio =} \FunctionTok{sd}\NormalTok{(normais}\SpecialCharTok{$}\NormalTok{amostra))}
\DocumentationTok{\#\#      media     desvio }
\DocumentationTok{\#\# {-}0.0002611  1.0000975}
\end{Highlighting}
\end{Shaded}

Este exemplo tem valor histórico: variantes deste esquema --- combinando a
exponencial como proposta com transformações inteligentes --- estiveram na base
dos primeiros geradores eficientes de números normais. Hoje o R usa por defeito
o método da inversão (\texttt{rnorm}), mais rápido, mas o argumento da
aceitação-rejeição permanece o mais transparente para compreender \emph{porque é que}
tais geradores funcionam.

\section{O caso de densidades não normalizadas}\label{o-caso-de-densidades-nuxe3o-normalizadas}

Uma virtude do método, que merece destaque, é a seguinte: para aplicar a
aceitação-rejeição não precisamos de conhecer a \emph{constante de normalização} de
\(f\). Suponhamos que só sabemos que a densidade alvo é proporcional a uma função
conhecida \(\tilde f\), isto é, \(f(x) = \tilde f(x)/Z\), com \(Z = \int \tilde f\)
possivelmente desconhecido ou difícil de calcular. Se encontrarmos \(\tilde c\)
tal que \(\tilde f(x) \leq \tilde c\, g(x)\), o critério de aceitação
\(U \leq \tilde f(Y)/[\tilde c\, g(Y)]\) produz observações de \(f\) na mesma, porque
a constante \(Z\) se cancela no quociente. Esta propriedade é decisiva na
estatística bayesiana, onde as distribuições \emph{a posteriori} são frequentemente
conhecidas apenas a menos de uma constante.

\section{Síntese}\label{suxedntese}

O método da aceitação-rejeição alarga consideravelmente o alcance da simulação:
permite gerar observações de praticamente qualquer densidade, exigindo apenas
que a saibamos avaliar ponto a ponto e que disponhamos de uma proposta cómoda
que a domine, a menos de uma constante \(c\). O preço a pagar é o descarte de
parte das observações --- em média, \(c\) tentativas por cada aceite --- o que
transforma a escolha da proposta no ponto crítico de toda a aplicação: uma boa
proposta é a que fica ``colada'' à densidade alvo, mantendo \(c\) próximo de \(1\).

\section{Exercícios}\label{exercuxedcios-35}

Nos exercícios seguintes, implemente o método da aceitação-rejeição, escolha
uma constante \(c\) adequada (justificando-a) e valide sempre o resultado
comparando o histograma das observações geradas com a densidade teórica.
Reporte também a taxa de aceitação e confronte-a com o valor esperado \(1/c\).

\textbf{1.} Gere variáveis aleatórias com distribuição exponencial
\(\text{Exp}(\lambda = 1)\) usando o método da aceitação-rejeição. Use uma
função de proposta uniforme no intervalo \([0, 5]\) e visualize o histograma das
observações geradas. Comente o efeito de truncar o suporte da exponencial em
\(5\).

\textbf{2.} Gere variáveis aleatórias com distribuição normal \(N(0,1)\) usando uma
distribuição de proposta Cauchy padrão \(C(0,1)\). Gere \(1000\) observações e
compare-as com a densidade da normal padrão. Determine analiticamente a
constante \(c\).

\textbf{3.} Use o método da aceitação-rejeição para gerar observações de uma
variável \(\text{Beta}(2, 3)\), recorrendo a uma proposta uniforme
\(\mathcal{U}(0, 1)\). Calcule o valor exato de \(c\) localizando o máximo da
densidade.

\textbf{4.} Use o método da aceitação-rejeição para gerar variáveis aleatórias
\(X \sim \Gamma(3, 2)\) (forma \(3\), taxa \(2\)) usando uma proposta exponencial
\(\text{Exp}(2)\). Explique por que razão a proposta deve ter média maior ou igual
à da distribuição alvo para que \(c\) seja finito.

\textbf{5.} Use o método da aceitação-rejeição para gerar variáveis log-normais
\(X \sim \text{LogNormal}(0, 1)\) usando uma proposta normal \(N(0, 1)\). Discuta as
dificuldades que surgem por a log-normal ter cauda mais pesada do que a normal.

\textbf{6.} Use o método da aceitação-rejeição para gerar variáveis de Weibull
\(W(\text{forma}=2, \text{escala}=1)\) usando uma proposta exponencial. Determine
\(c\) numericamente maximizando o rácio das densidades com \texttt{optimize()}.

\textbf{7.} Considere a densidade \(f(x) = \frac{1}{2}\,(1 + \cos x)\) no intervalo
\([-\pi, \pi]\) e nula fora dele. Gere observações desta distribuição com uma
proposta uniforme em \([-\pi, \pi]\) e verifique o resultado graficamente.

\textbf{8.} Implemente uma versão \textbf{vetorizada} da função \texttt{aceita\_rejeita} que gere
os candidatos em lotes (por exemplo, de tamanho \(2n\) de cada vez) e aceite em
bloco os que passam no teste, repetindo apenas se necessário. Compare o tempo de
execução com a versão baseada num ciclo \texttt{while}, usando \texttt{system.time()}, para
gerar \(10^5\) observações da Beta(2,2).

\textbf{9.} Estude, por simulação, como a \textbf{eficiência} do método depende da
constante \(c\). Para gerar observações da meia-normal a partir da exponencial,
registe a taxa de aceitação empírica e compare-a com \(1/c\). Em seguida, repita a
experiência propositadamente com um \(c\) demasiado \emph{grande} (por exemplo, o dobro
do correto) e confirme que a distribuição gerada continua correta, mas que a
taxa de aceitação cai. O que aconteceria com um \(c\) demasiado \emph{pequeno}?

\textbf{10.} A densidade \(f(x) = \frac{2}{\pi}\sqrt{1 - x^2}\) em \([-1, 1]\) (a
distribuição do \emph{semicírculo}, ou lei de Wigner no intervalo unitário) surge no
estudo de matrizes aleatórias. Gere \(5000\) observações desta distribuição por
aceitação-rejeição com proposta uniforme em \([-1,1]\), determine \(c\)
analiticamente e confirme graficamente o resultado.

\textbf{11.} \emph{(Densidade não normalizada.)} Pretende-se amostrar de uma densidade
proporcional a \(\tilde f(x) = e^{-x^2/2}\,(1 + \sin^2 3x)\) em \(\mathbb{R}\), cuja
constante de normalização não é conhecida em forma fechada. Usando uma proposta
normal com variância suficientemente grande, mostre que consegue gerar
observações de \(f\) sem calcular essa constante. Estime, a partir da taxa de
aceitação, o valor da constante de normalização.

\chapter{Teoremas limites}\label{teoremas-limites}

Ao longo dos capítulos anteriores, gerámos repetidamente amostras aleatórias e
calculámos médias, proporções e outras estatísticas. Sempre que o fizemos,
apoiámo-nos --- por vezes de forma implícita --- em dois resultados que constituem
o alicerce de toda a inferência estatística e de toda a simulação: a \textbf{lei dos
grandes números} e o \textbf{teorema limite central}. São eles que explicam por que
razão a média de muitas observações se aproxima do valor esperado, e por que
razão a distribuição normal surge, uma e outra vez, onde menos se espera. Este
capítulo dedica-se a compreendê-los, a demonstrá-los nos casos acessíveis e,
sobretudo, a \emph{vê-los em ação} através da simulação.

Os teoremas limites clássicos da probabilidade referem-se a sequências de
variáveis aleatórias independentes e identicamente distribuídas (i.i.d.).
Consideremos uma tal sequência \(X_1, X_2, \ldots\) com média comum
\(E(X_i) = \mu < \infty\), e definamos a soma parcial

\[
S_n = X_1 + X_2 + \cdots + X_n.
\]

A média amostral é \(\bar{X}_n = S_n / n\). A \textbf{lei dos grandes números} afirma
que \(\bar{X}_n\) converge para \(\mu\) quando \(n \to +\infty\); o \textbf{teorema limite
central} vai mais longe e descreve \emph{como} se comporta o desvio de \(\bar{X}_n\)
em torno de \(\mu\), revelando que, depois de convenientemente reescalado, esse
desvio é aproximadamente normal, seja qual for a distribuição de partida.

\section{Contexto histórico}\label{contexto-histuxf3rico-1}

A primeira versão da lei dos grandes números deve-se a \textbf{Jacob Bernoulli}, que
a demonstrou para proporções de sucessos em provas de Bernoulli na sua obra
\emph{Ars Conjectandi}, publicada postumamente em 1713. Bernoulli chamou-lhe o seu
``teorema de ouro'' e dedicou-lhe vinte anos de trabalho: era a prova matemática
de uma convicção antiga, a de que a frequência relativa de um acontecimento
tende a estabilizar em torno da sua probabilidade à medida que se acumulam
observações. O nome ``lei dos grandes números'' foi cunhado mais tarde, em 1837,
por \textbf{Siméon Denis Poisson}.

A distinção entre as versões \emph{fraca} e \emph{forte} da lei --- que abaixo tornaremos
precisa --- foi clarificada no início do século XX: \textbf{Émile Borel} estabeleceu em
1909 a lei forte para o caso binomial, e \textbf{Andrei Kolmogorov} deu-lhe, nos anos
1930, a formulação geral que hoje usamos, na sequência da sua axiomatização da
teoria da probabilidade.

O teorema limite central tem uma história igualmente ilustre. \textbf{Abraham de
Moivre} obteve, em 1733, a aproximação normal da distribuição binomial;
\textbf{Pierre-Simon Laplace} generalizou-a substancialmente em 1810; e a versão
geral, com condições rigorosas, foi estabelecida por \textbf{Aleksandr Lyapunov}
(1901) e, na forma que hoje é mais citada para variáveis i.i.d., por
\textbf{Lindeberg} e \textbf{Lévy}. A designação \emph{teorema limite central} --- no sentido de
teorema \emph{central} da teoria --- foi popularizada por George Pólya em 1920.

\section{Modos de convergência}\label{modos-de-converguxeancia}

Falar de convergência de variáveis aleatórias exige cuidado, porque uma
sucessão de variáveis aleatórias pode convergir em sentidos diferentes. Três
modos interessam-nos aqui.

Diz-se que \(Y_n\) converge \textbf{em probabilidade} para \(Y\), e escreve-se
\(Y_n \xrightarrow{P} Y\), se para todo o \(\varepsilon > 0\)

\[
\lim_{n \to \infty} P\big(|Y_n - Y| > \varepsilon\big) = 0.
\]

Diz-se que \(Y_n\) converge \textbf{quase certamente} (ou com probabilidade \(1\)) para
\(Y\), e escreve-se \(Y_n \xrightarrow{\text{q.c.}} Y\), se

\[
P\Big(\lim_{n \to \infty} Y_n = Y\Big) = 1.
\]

Diz-se que \(Y_n\) converge \textbf{em distribuição} para \(Y\), e escreve-se
\(Y_n \xrightarrow{d} Y\), se \(F_{Y_n}(y) \to F_Y(y)\) em todos os pontos \(y\) de
continuidade de \(F_Y\).

Estes três modos não são equivalentes. A convergência quase certa é mais forte
do que a convergência em probabilidade (implica-a), e esta, por sua vez, é mais
forte do que a convergência em distribuição. A lei \emph{fraca} dos grandes números
diz respeito à convergência em probabilidade; a lei \emph{forte}, à convergência
quase certa; o teorema limite central, à convergência em distribuição. É esta a
distinção que dá nome aos resultados.

\section{Lei fraca dos grandes números}\label{lei-fraca-dos-grandes-nuxfameros}

A lei fraca dos grandes números é também conhecida como \emph{teorema de Bernoulli}.
De acordo com a lei, a média dos resultados obtidos num grande número de provas
está próxima da média da população, com probabilidade cada vez mais alta à
medida que o número de provas aumenta.

Enunciar e demonstrar a lei fraca dos grandes números no caso de variância
finita, compreender o seu significado e ilustrá-lo por simulação.

Sejam \(X_1, \ldots, X_n\) variáveis aleatórias i.i.d., cada uma com média \(\mu\) e
variância \(\sigma^2 < \infty\). Então, para todo o \(\varepsilon > 0\),

\[
\lim_{n \to \infty} P\big(|\bar{X}_n - \mu| > \varepsilon\big) = 0,
\qquad \text{ou seja,} \qquad
\bar{X}_n \xrightarrow{P} \mu.
\]

\textbf{Demonstração.} A média amostral tem valor esperado

\[
E(\bar{X}_n) = \frac{1}{n}\sum_{i=1}^{n} E(X_i) = \mu,
\]

e, pela independência das observações, variância

\[
\operatorname{Var}(\bar{X}_n)
= \frac{1}{n^2}\sum_{i=1}^{n}\operatorname{Var}(X_i)
= \frac{\sigma^2}{n}.
\]

Aplicando a \textbf{desigualdade de Chebyshev} à variável \(\bar{X}_n\), cuja média é
\(\mu\) e cuja variância é \(\sigma^2/n\), obtém-se

\[
P\big(|\bar{X}_n - \mu| > \varepsilon\big)
\leq \frac{\operatorname{Var}(\bar{X}_n)}{\varepsilon^2}
= \frac{\sigma^2}{n\,\varepsilon^2}.
\]

Fixado \(\varepsilon > 0\), o membro direito tende para zero quando
\(n \to \infty\), o que estabelece a convergência em probabilidade. \(\qquad\blacksquare\)

A demonstração acima usa a hipótese de variância finita porque recorre à
desigualdade de Chebyshev. Existe, porém, uma versão mais geral --- o \emph{teorema de
Khinchin} --- que garante a lei fraca exigindo apenas que a média \(\mu\) exista e
seja finita, dispensando a hipótese sobre a variância. A demonstração desse
resultado é mais delicada e recorre a funções características.

\section{Lei forte dos grandes números}\label{lei-forte-dos-grandes-nuxfameros}

A lei fraca garante que, para \(n\) grande, é \emph{improvável} que \(\bar{X}_n\) se
afaste de \(\mu\). A lei \textbf{forte} afirma algo qualitativamente mais poderoso: que
a própria sucessão \(\bar{X}_1, \bar{X}_2, \ldots\) converge para \(\mu\), encarada
como uma trajetória, com probabilidade \(1\).

Sejam \(X_1, X_2, \ldots\) variáveis aleatórias i.i.d. com \(E|X_1| < \infty\) e
média \(\mu\). Então

\[
\bar{X}_n \xrightarrow{\text{q.c.}} \mu
\qquad \text{quando} \qquad n \to \infty.
\]

Este é o \textbf{teorema de Kolmogorov}. Note-se que exige apenas a existência da
média --- nem sequer se assume variância finita. A demonstração está para além do
âmbito deste livro, mas o seu significado é fundamental e vale a pena fixá-lo.

A diferença entre as duas leis é subtil mas importante. A lei fraca diz que,
para cada \(n\) suficientemente grande, a probabilidade de \(\bar{X}_n\) estar longe
de \(\mu\) é pequena --- mas não exclui que, ao longo da sucessão, ocorram
ocasionalmente afastamentos grandes. A lei forte garante que, para \emph{quase todas}
as trajetórias possíveis, a sucessão acaba por entrar e permanecer arbitrariamente
perto de \(\mu\): os afastamentos grandes deixam de ocorrer a partir de certa
ordem. É esta a garantia que legitima, em última análise, a interpretação
frequencista da probabilidade e o método de Monte Carlo.

\subsection{Ilustração por simulação}\label{ilustrauxe7uxe3o-por-simulauxe7uxe3o}

Vejamos a lei dos grandes números em funcionamento. Simulemos o lançamento
repetido de um dado equilibrado, cuja média teórica é
\(\mu = (1+2+3+4+5+6)/6 = 3{,}5\), e observemos a média acumulada
\(\bar{X}_n\) a estabilizar em torno desse valor à medida que \(n\) cresce.

\begin{Shaded}
\begin{Highlighting}[]
\FunctionTok{set.seed}\NormalTok{(}\DecValTok{2024}\NormalTok{)}

\NormalTok{n }\OtherTok{\textless{}{-}} \DecValTok{5000}
\NormalTok{lancamentos   }\OtherTok{\textless{}{-}} \FunctionTok{sample}\NormalTok{(}\DecValTok{1}\SpecialCharTok{:}\DecValTok{6}\NormalTok{, }\AttributeTok{size =}\NormalTok{ n, }\AttributeTok{replace =} \ConstantTok{TRUE}\NormalTok{)}
\NormalTok{media\_acumulada }\OtherTok{\textless{}{-}} \FunctionTok{cumsum}\NormalTok{(lancamentos) }\SpecialCharTok{/} \FunctionTok{seq\_along}\NormalTok{(lancamentos)}
\end{Highlighting}
\end{Shaded}

\begin{Shaded}
\begin{Highlighting}[]
\FunctionTok{plot}\NormalTok{(media\_acumulada, }\AttributeTok{type =} \StringTok{"l"}\NormalTok{, }\AttributeTok{col =} \StringTok{"\#10656d"}\NormalTok{, }\AttributeTok{lwd =} \FloatTok{1.2}\NormalTok{,}
     \AttributeTok{ylim =} \FunctionTok{c}\NormalTok{(}\FloatTok{2.5}\NormalTok{, }\FloatTok{4.5}\NormalTok{), }\AttributeTok{log =} \StringTok{"x"}\NormalTok{,}
     \AttributeTok{xlab =} \StringTok{"número de lançamentos (escala logarítmica)"}\NormalTok{,}
     \AttributeTok{ylab =} \FunctionTok{expression}\NormalTok{(}\FunctionTok{bar}\NormalTok{(X)[n]))}
\FunctionTok{abline}\NormalTok{(}\AttributeTok{h =} \FloatTok{3.5}\NormalTok{, }\AttributeTok{lty =} \DecValTok{2}\NormalTok{, }\AttributeTok{col =} \StringTok{"grey40"}\NormalTok{)}
\end{Highlighting}
\end{Shaded}

\begin{figure}

{\centering \includegraphics[width=0.8\linewidth]{introR_files/figure-latex/tl-lgn-grafico-1} 

}

\caption{Média acumulada dos lançamentos de um dado equilibrado em função do número de lançamentos. A linha tracejada assinala a média teórica, 3,5.}\label{fig:tl-lgn-grafico}
\end{figure}

O gráfico mostra com clareza o fenómeno: nos primeiros lançamentos a média
acumulada oscila bastante, mas rapidamente se aproxima de \(3{,}5\) e aí
permanece. A escala logarítmica no eixo horizontal ajuda a ver que a
estabilização não é fruto do acaso de uma única trajetória. Para o confirmar,
sobreponhamos várias trajetórias independentes.

\begin{Shaded}
\begin{Highlighting}[]
\FunctionTok{set.seed}\NormalTok{(}\DecValTok{2024}\NormalTok{)}

\FunctionTok{plot}\NormalTok{(}\ConstantTok{NULL}\NormalTok{, }\AttributeTok{xlim =} \FunctionTok{c}\NormalTok{(}\DecValTok{1}\NormalTok{, n), }\AttributeTok{ylim =} \FunctionTok{c}\NormalTok{(}\FloatTok{2.5}\NormalTok{, }\FloatTok{4.5}\NormalTok{), }\AttributeTok{log =} \StringTok{"x"}\NormalTok{,}
     \AttributeTok{xlab =} \StringTok{"número de lançamentos (escala logarítmica)"}\NormalTok{,}
     \AttributeTok{ylab =} \FunctionTok{expression}\NormalTok{(}\FunctionTok{bar}\NormalTok{(X)[n]))}
\ControlFlowTok{for}\NormalTok{ (k }\ControlFlowTok{in} \DecValTok{1}\SpecialCharTok{:}\DecValTok{10}\NormalTok{) \{}
\NormalTok{  x }\OtherTok{\textless{}{-}} \FunctionTok{sample}\NormalTok{(}\DecValTok{1}\SpecialCharTok{:}\DecValTok{6}\NormalTok{, }\AttributeTok{size =}\NormalTok{ n, }\AttributeTok{replace =} \ConstantTok{TRUE}\NormalTok{)}
  \FunctionTok{lines}\NormalTok{(}\FunctionTok{cumsum}\NormalTok{(x) }\SpecialCharTok{/} \FunctionTok{seq\_along}\NormalTok{(x), }\AttributeTok{col =} \FunctionTok{adjustcolor}\NormalTok{(}\StringTok{"\#10656d"}\NormalTok{, }\FloatTok{0.5}\NormalTok{))}
\NormalTok{\}}
\FunctionTok{abline}\NormalTok{(}\AttributeTok{h =} \FloatTok{3.5}\NormalTok{, }\AttributeTok{lty =} \DecValTok{2}\NormalTok{, }\AttributeTok{col =} \StringTok{"grey30"}\NormalTok{, }\AttributeTok{lwd =} \DecValTok{2}\NormalTok{)}
\end{Highlighting}
\end{Shaded}

\begin{figure}

{\centering \includegraphics[width=0.8\linewidth]{introR_files/figure-latex/tl-lgn-varias-1} 

}

\caption{Dez trajetórias independentes da média acumulada. Todas convergem para a média teórica, ilustrando a lei forte dos grandes números.}\label{fig:tl-lgn-varias}
\end{figure}

Todas as trajetórias, apesar de partirem de valores iniciais dispersos,
convergem para a mesma reta horizontal. É precisamente esta a mensagem da lei
forte: não é uma trajetória feliz que converge, são (quase) todas.

\section{Teorema limite central}\label{teorema-limite-central}

A lei dos grandes números diz-nos \emph{para onde} converge a média amostral, mas
nada diz sobre a \emph{rapidez} dessa convergência nem sobre a \emph{forma} das flutuações
em torno do limite. É a essa pergunta que responde o teorema limite central, um
dos resultados mais surpreendentes e úteis de toda a matemática.

Sejam \(X_1, X_2, \ldots\) variáveis aleatórias i.i.d. com média \(\mu\) e variância
\(\sigma^2\), ambas finitas, com \(\sigma^2 > 0\). Então a soma padronizada
converge em distribuição para a normal padrão:

\[
Z_n = \frac{S_n - n\mu}{\sigma\sqrt{n}}
     = \frac{\bar{X}_n - \mu}{\sigma/\sqrt{n}}
     \;\xrightarrow{d}\; N(0, 1)
\qquad \text{quando} \qquad n \to \infty.
\]

Equivalentemente, para \(n\) grande, a média amostral é aproximadamente normal:

\[
\bar{X}_n \;\approx\; N\!\left(\mu,\; \frac{\sigma^2}{n}\right).
\]

O que é notável --- quase paradoxal --- é que esta conclusão \textbf{não depende da
distribuição das \(X_i\)}. Sejam elas discretas ou contínuas, simétricas ou
enviesadas, a distribuição da média padronizada converge sempre para a mesma
curva em forma de sino. É esta universalidade que explica a onipresença da
normal na estatística: sempre que uma quantidade resulta da soma de muitos
contributos pequenos, independentes e de magnitude comparável, é natural que se
comporte de forma aproximadamente normal.

O teorema garante convergência quando \(n \to \infty\), mas na prática usamo-lo
como \emph{aproximação} para \(n\) finito. A qualidade dessa aproximação depende
fortemente da forma da distribuição de partida: para distribuições próximas da
simetria, uns poucos termos bastam; para distribuições muito enviesadas ou com
caudas pesadas, pode ser preciso um \(n\) considerável. A regra prática
``\(n \geq 30\)'', frequente nos manuais, é apenas uma orientação grosseira e falha
para distribuições fortemente assimétricas.

\subsection{Ilustração por simulação}\label{ilustrauxe7uxe3o-por-simulauxe7uxe3o-1}

Para ver o teorema em ação, partamos de uma distribuição deliberadamente
assimétrica --- a exponencial de parâmetro \(1\), que tem média \(\mu = 1\) e
desvio-padrão \(\sigma = 1\), e cuja densidade nada tem de ``sino''. Vamos gerar
muitas amostras de tamanho \(n\), calcular a média de cada uma, e observar a
distribuição dessas médias à medida que \(n\) aumenta.

\begin{Shaded}
\begin{Highlighting}[]
\FunctionTok{set.seed}\NormalTok{(}\DecValTok{2024}\NormalTok{)}

\NormalTok{simular\_medias }\OtherTok{\textless{}{-}} \ControlFlowTok{function}\NormalTok{(n, }\AttributeTok{replicas =} \DecValTok{10000}\NormalTok{) \{}
\NormalTok{  medias }\OtherTok{\textless{}{-}} \FunctionTok{replicate}\NormalTok{(replicas, }\FunctionTok{mean}\NormalTok{(}\FunctionTok{rexp}\NormalTok{(n, }\AttributeTok{rate =} \DecValTok{1}\NormalTok{)))}
  \CommentTok{\# padronização: (media {-} mu) / (sigma / sqrt(n)), com mu = sigma = 1}
\NormalTok{  (medias }\SpecialCharTok{{-}} \DecValTok{1}\NormalTok{) }\SpecialCharTok{/}\NormalTok{ (}\DecValTok{1} \SpecialCharTok{/} \FunctionTok{sqrt}\NormalTok{(n))}
\NormalTok{\}}

\NormalTok{z1  }\OtherTok{\textless{}{-}} \FunctionTok{simular\_medias}\NormalTok{(}\DecValTok{1}\NormalTok{)}
\NormalTok{z5  }\OtherTok{\textless{}{-}} \FunctionTok{simular\_medias}\NormalTok{(}\DecValTok{5}\NormalTok{)}
\NormalTok{z30 }\OtherTok{\textless{}{-}} \FunctionTok{simular\_medias}\NormalTok{(}\DecValTok{30}\NormalTok{)}
\end{Highlighting}
\end{Shaded}

\begin{Shaded}
\begin{Highlighting}[]
\NormalTok{op }\OtherTok{\textless{}{-}} \FunctionTok{par}\NormalTok{(}\AttributeTok{mfrow =} \FunctionTok{c}\NormalTok{(}\DecValTok{1}\NormalTok{, }\DecValTok{3}\NormalTok{), }\AttributeTok{mar =} \FunctionTok{c}\NormalTok{(}\DecValTok{4}\NormalTok{, }\DecValTok{4}\NormalTok{, }\DecValTok{2}\NormalTok{, }\DecValTok{1}\NormalTok{))}
\ControlFlowTok{for}\NormalTok{ (par\_z }\ControlFlowTok{in} \FunctionTok{list}\NormalTok{(}\FunctionTok{list}\NormalTok{(z1, }\StringTok{"n = 1"}\NormalTok{), }\FunctionTok{list}\NormalTok{(z5, }\StringTok{"n = 5"}\NormalTok{),}
                    \FunctionTok{list}\NormalTok{(z30, }\StringTok{"n = 30"}\NormalTok{))) \{}
  \FunctionTok{hist}\NormalTok{(par\_z[[}\DecValTok{1}\NormalTok{]], }\AttributeTok{breaks =} \DecValTok{40}\NormalTok{, }\AttributeTok{probability =} \ConstantTok{TRUE}\NormalTok{,}
       \AttributeTok{col =} \StringTok{"grey85"}\NormalTok{, }\AttributeTok{border =} \StringTok{"white"}\NormalTok{,}
       \AttributeTok{xlim =} \FunctionTok{c}\NormalTok{(}\SpecialCharTok{{-}}\DecValTok{4}\NormalTok{, }\DecValTok{4}\NormalTok{), }\AttributeTok{main =}\NormalTok{ par\_z[[}\DecValTok{2}\NormalTok{]], }\AttributeTok{xlab =} \StringTok{"z"}\NormalTok{)}
  \FunctionTok{curve}\NormalTok{(}\FunctionTok{dnorm}\NormalTok{(x), }\AttributeTok{add =} \ConstantTok{TRUE}\NormalTok{, }\AttributeTok{lwd =} \DecValTok{2}\NormalTok{, }\AttributeTok{col =} \StringTok{"\#10656d"}\NormalTok{)}
\NormalTok{\}}
\end{Highlighting}
\end{Shaded}

\begin{figure}

{\centering \includegraphics[width=0.8\linewidth]{introR_files/figure-latex/tl-tlc-grafico-1} 

}

\caption{Distribuição da média amostral padronizada de observações exponenciais, para amostras de tamanho crescente. A curva a cheio é a densidade normal padrão. À medida que n aumenta, a assimetria da exponencial esbate-se e a distribuição das médias aproxima-se da normal.}\label{fig:tl-tlc-grafico}
\end{figure}

\begin{Shaded}
\begin{Highlighting}[]
\FunctionTok{par}\NormalTok{(op)}
\end{Highlighting}
\end{Shaded}

A sequência de histogramas conta toda a história. Para \(n = 1\), a distribuição
das médias é simplesmente a própria exponencial padronizada --- visivelmente
assimétrica, com uma cauda longa à direita. Para \(n = 5\), a assimetria já é bem
menor. Para \(n = 30\), o histograma é praticamente indistinguível da curva normal
sobreposta. Esta convergência, partindo de uma distribuição tão pouco parecida
com um sino, é a demonstração visual da universalidade do teorema.

Podemos confirmar a normalidade de forma mais fina com um gráfico
quantil-quantil, que compara os quantis empíricos com os da normal padrão. Se a
aproximação for boa, os pontos alinham-se sobre a reta.

\begin{Shaded}
\begin{Highlighting}[]
\FunctionTok{qqnorm}\NormalTok{(z30, }\AttributeTok{main =} \StringTok{""}\NormalTok{, }\AttributeTok{col =} \FunctionTok{adjustcolor}\NormalTok{(}\StringTok{"\#10656d"}\NormalTok{, }\FloatTok{0.4}\NormalTok{), }\AttributeTok{pch =} \DecValTok{19}\NormalTok{, }\AttributeTok{cex =} \FloatTok{0.4}\NormalTok{)}
\FunctionTok{qqline}\NormalTok{(z30, }\AttributeTok{lwd =} \DecValTok{2}\NormalTok{, }\AttributeTok{col =} \StringTok{"grey30"}\NormalTok{)}
\end{Highlighting}
\end{Shaded}

\begin{figure}

{\centering \includegraphics[width=0.8\linewidth]{introR_files/figure-latex/tl-qq-1} 

}

\caption{Gráfico quantil-quantil das médias padronizadas para n = 30. O bom alinhamento sobre a reta confirma a proximidade à normal.}\label{fig:tl-qq}
\end{figure}

\subsection{O caso de de Moivre-Laplace}\label{o-caso-de-de-moivre-laplace}

O caso histórico original --- a aproximação normal da binomial --- é um caso
particular do teorema limite central, obtido quando as \(X_i\) são variáveis de
Bernoulli. Se \(X_i \sim \text{Bernoulli}(p)\), então \(S_n \sim \text{Binomial}(n, p)\)
e o teorema garante que, para \(n\) grande,

\[
\frac{S_n - np}{\sqrt{np(1-p)}} \;\xrightarrow{d}\; N(0,1).
\]

É este resultado que justifica a clássica aproximação da binomial pela normal.
Ilustremo-lo comparando a distribuição exata de uma binomial com a densidade
normal correspondente.

\begin{Shaded}
\begin{Highlighting}[]
\FunctionTok{set.seed}\NormalTok{(}\DecValTok{2024}\NormalTok{)}

\NormalTok{n\_ensaios }\OtherTok{\textless{}{-}} \DecValTok{50}
\NormalTok{p }\OtherTok{\textless{}{-}} \FloatTok{0.3}
\NormalTok{x }\OtherTok{\textless{}{-}} \DecValTok{0}\SpecialCharTok{:}\NormalTok{n\_ensaios}
\NormalTok{prob }\OtherTok{\textless{}{-}} \FunctionTok{dbinom}\NormalTok{(x, n\_ensaios, p)}

\FunctionTok{plot}\NormalTok{(x, prob, }\AttributeTok{type =} \StringTok{"h"}\NormalTok{, }\AttributeTok{lwd =} \DecValTok{2}\NormalTok{, }\AttributeTok{col =} \StringTok{"grey70"}\NormalTok{,}
     \AttributeTok{xlab =} \StringTok{"número de sucessos"}\NormalTok{, }\AttributeTok{ylab =} \StringTok{"probabilidade"}\NormalTok{)}
\FunctionTok{curve}\NormalTok{(}\FunctionTok{dnorm}\NormalTok{(x, }\AttributeTok{mean =}\NormalTok{ n\_ensaios }\SpecialCharTok{*}\NormalTok{ p,}
            \AttributeTok{sd =} \FunctionTok{sqrt}\NormalTok{(n\_ensaios }\SpecialCharTok{*}\NormalTok{ p }\SpecialCharTok{*}\NormalTok{ (}\DecValTok{1} \SpecialCharTok{{-}}\NormalTok{ p))),}
      \AttributeTok{add =} \ConstantTok{TRUE}\NormalTok{, }\AttributeTok{lwd =} \DecValTok{2}\NormalTok{, }\AttributeTok{col =} \StringTok{"\#10656d"}\NormalTok{)}
\end{Highlighting}
\end{Shaded}

\begin{figure}

{\centering \includegraphics[width=0.8\linewidth]{introR_files/figure-latex/tl-binomial-1} 

}

\caption{Distribuição binomial(50; 0,3) (barras) e aproximação normal (curva). Para n moderado, o acordo já é notável.}\label{fig:tl-binomial}
\end{figure}

\section{Síntese}\label{suxedntese-1}

A lei dos grandes números e o teorema limite central formam a espinha dorsal da
estatística inferencial e justificam a própria legitimidade da simulação. A
primeira garante que a média de muitas observações converge para o valor
esperado --- na versão fraca, em probabilidade; na versão forte, quase certamente.
O segundo descreve as flutuações em torno desse limite, mostrando que, depois de
reescaladas por \(\sqrt{n}\), elas são universalmente normais, independentemente da
distribuição de partida. Juntos, estes resultados dizem-nos não só \emph{que} as
estimativas de Monte Carlo convergem, mas \emph{a que ritmo} e \emph{com que incerteza} ---
tema que retomaremos no próximo capítulo, ao aproximar integrais.

\section{Exercícios}\label{exercuxedcios-36}

\textbf{1.} Simule o lançamento de uma moeda equilibrada \(n = 10\,000\) vezes e
represente graficamente a proporção acumulada de caras em função do número de
lançamentos. Confirme visualmente a convergência para \(0{,}5\) e comente a
rapidez com que a proporção se estabiliza.

\textbf{2.} Repita o exercício anterior com uma moeda \emph{viciada}, em que a
probabilidade de cara é \(0{,}7\). Sobreponha a reta horizontal correspondente ao
valor teórico e verifique a convergência.

\textbf{3.} Gere \(n = 2000\) observações de uma distribuição de Poisson de parâmetro
\(\lambda = 4\) e represente a média acumulada. Compare a trajetória com a média
teórica e discuta se a lei se aplica (a Poisson tem variância finita?).

\textbf{4.} Ilustre a lei dos grandes números para a média amostral de uma
distribuição \textbf{uniforme} \(\mathcal{U}(0,1)\). Qual o valor para o qual a média
acumulada deve convergir? Confirme por simulação.

\textbf{5.} A distribuição de \textbf{Cauchy} padrão não tem média finita. Gere a média
acumulada de \(n = 5000\) observações de uma Cauchy (\texttt{rcauchy}) e represente-a
graficamente. A média acumulada converge? Explique o resultado à luz das
hipóteses da lei dos grandes números.

\textbf{6.} Verifique numericamente a lei fraca dos grandes números: para
\(X_i \sim \mathcal{U}(0,1)\) (média \(1/2\), variância \(1/12\)) e \(\varepsilon = 0{,}05\),
estime por simulação a probabilidade \(P(|\bar{X}_n - 1/2| > \varepsilon)\) para
\(n = 10, 50, 100, 500\) e confirme que decresce. Compare com o majorante de
Chebyshev \(\sigma^2/(n\varepsilon^2)\).

\textbf{7.} Ilustre o teorema limite central partindo de uma distribuição
\textbf{uniforme} \(\mathcal{U}(0,1)\). Gere \(10\,000\) médias de amostras de tamanho
\(n = 2, 5, 12\) e sobreponha a densidade normal apropriada. A partir de que \(n\) a
aproximação lhe parece satisfatória? (A soma de doze uniformes esteve, de facto,
na base de geradores antigos de números normais.)

\textbf{8.} Repita a ilustração do teorema limite central partindo de uma
distribuição fortemente enviesada, como a \textbf{log-normal} \(\text{LogNormal}(0, 1)\).
Compare a rapidez de convergência com a do caso exponencial estudado no texto.
Qual necessita de um \(n\) maior para a aproximação normal ser aceitável?

\textbf{9.} Para a distribuição binomial \(\text{Binomial}(n, p)\) com \(p = 0{,}1\),
compare graficamente a distribuição exata com a aproximação normal, para
\(n = 10, 30, 100\). Confirme que, com \(p\) pequeno, é necessário um \(n\) maior do
que o habitual para que a aproximação seja razoável, e relacione isto com a
assimetria da binomial.

\textbf{10.} \emph{(Cobertura de um intervalo de confiança.)} Uma consequência prática do
teorema limite central é a validade aproximada do intervalo de confiança
\(\bar{X}_n \pm z_{1-\alpha/2}\,s/\sqrt{n}\) para a média. Gere \(5000\) amostras de
tamanho \(n = 40\) de uma exponencial de média \(1\), construa em cada uma o
intervalo de confiança a \(95\%\) para a média, e estime a proporção de intervalos
que efetivamente contêm o valor \(1\). A cobertura empírica aproxima-se de \(95\%\)?

\textbf{11.} \emph{(Efeito do tamanho da amostra na cobertura.)} Repita o exercício
anterior para \(n = 10, 40, 200\) e discuta como a cobertura empírica se aproxima
do nível nominal de \(95\%\) à medida que \(n\) cresce. Relacione o resultado com a
qualidade da aproximação normal para a média de observações exponenciais.

\textbf{12.} Combine as duas leis num único gráfico: para observações
\(\mathcal{U}(0,1)\), represente simultaneamente (a) a média acumulada de uma
trajetória a convergir para \(1/2\) e (b), num segundo painel, o histograma das
médias de muitas amostras de tamanho \(n = 30\) com a normal sobreposta. Comente
como cada painel ilustra um dos dois teoremas centrais do capítulo.

\chapter{Cálculo de integrais por Monte Carlo}\label{calculo-integrais}

Chegámos a um dos pontos altos deste livro, onde a simulação deixa de ser apenas
uma forma de gerar dados artificiais e se torna um verdadeiro instrumento de
cálculo. Muitos integrais que surgem na estatística e nas ciências não têm
primitiva expressável em forma fechada --- pense-se no integral da densidade
normal --- e mesmo os métodos numéricos determinísticos clássicos, como a regra
dos trapézios ou de Simpson, tornam-se impraticáveis quando o número de
dimensões cresce. É aqui que entra o \textbf{método de Monte Carlo}: usar números
aleatórios para aproximar o valor de um integral. A ideia é tão simples quanto
poderosa, e assenta diretamente nos teoremas limites que estudámos no capítulo
anterior.

\section{Contexto histórico}\label{contexto-histuxf3rico-2}

O método de Monte Carlo nasceu no final da década de 1940, no Laboratório
Nacional de Los Alamos, das mãos de \textbf{Stanislaw Ulam}, \textbf{John von Neumann} e
\textbf{Nicholas Metropolis}. Ulam, convalescente de uma doença e a jogar paciências,
perguntou-se qual seria a probabilidade de um jogo sair bem sucedido; percebendo
que o cálculo combinatório exato era inviável, ocorreu-lhe que seria mais
simples jogar muitas partidas e contar. Dessa intuição --- substituir o cálculo
exato pela repetição aleatória --- nasceu uma técnica que von Neumann rapidamente
adaptou aos problemas de difusão de neutrões então em estudo. O nome, sugerido
por Metropolis, é uma alusão ao famoso casino do Mónaco, evocando a natureza
aleatória do procedimento.

O método tem, contudo, um precursor ilustre e bem mais antigo: a \textbf{agulha de
Buffon}. Em 1777, Georges-Louis Leclerc, conde de Buffon, mostrou que a
probabilidade de uma agulha lançada ao acaso cruzar uma das linhas paralelas de
um soalho depende de \(\pi\) --- o que permite \emph{estimar} \(\pi\) atirando agulhas e
contando cruzamentos. É, em espírito, o primeiro método de Monte Carlo da
história, com quase dois séculos de antecedência sobre Los Alamos.

\section{A ideia fundamental}\label{a-ideia-fundamental}

Seja \(g(x)\) uma função e suponhamos que se pretende calcular

\[
\theta = \int_{0}^{1} g(x)\, dx.
\]

A chave do método está em reconhecer este integral como um \textbf{valor esperado}.
Consideremos \(U\) uma variável aleatória com distribuição uniforme no intervalo
\((0,1)\), cuja densidade é \(1\) nesse intervalo. Por definição de valor esperado,

\[
E[g(U)] = \int_{0}^{1} g(u) \cdot 1 \, du = \int_{0}^{1} g(x)\, dx = \theta.
\]

Ou seja, o integral que queremos calcular é exatamente a média da variável
aleatória \(g(U)\):

\[
\theta = E[g(U)].
\]

Esta reformulação é decisiva, porque sabemos aproximar valores esperados: basta
gerar muitas observações e calcular a média amostral. Se \(U_1, \ldots, U_k\) são
variáveis aleatórias independentes e uniformemente distribuídas em \((0,1)\),
então \(g(U_1), \ldots, g(U_k)\) são i.i.d. com média \(\theta\). Pela \textbf{lei forte
dos grandes números}, com probabilidade \(1\),

\[
\hat{\theta}_k = \frac{1}{k}\sum_{i=1}^{k} g(U_i)
\;\longrightarrow\; E[g(U)] = \theta,
\qquad \text{quando} \qquad k \to \infty.
\]

Aprender a aproximar integrais definidos --- em uma ou várias dimensões, em
intervalos limitados ou ilimitados --- através da média de avaliações de uma
função em pontos aleatórios, e saber quantificar a incerteza dessa aproximação.

Em termos práticos, para aproximar \(\theta\) geramos uma grande quantidade de
números aleatórios \(u_1, \ldots, u_k\) em \((0,1)\) e calculamos a média dos
valores \(g(u_i)\):

\[
\theta \approx \frac{1}{k}\sum_{i=1}^{k} g(u_i).
\]

É esta a essência do método de Monte Carlo para o cálculo de integrais.

\section{O erro da aproximação}\label{o-erro-da-aproximauxe7uxe3o}

Uma estimativa sem uma medida da sua incerteza tem pouco valor. Felizmente, o
mesmo enquadramento probabilístico que legitima o método fornece também o erro.
O estimador \(\hat{\theta}_k\) é uma média de variáveis i.i.d., pelo que é
\textbf{não enviesado} (\(E[\hat{\theta}_k] = \theta\)) e tem variância

\[
\operatorname{Var}(\hat{\theta}_k) = \frac{\sigma_g^2}{k},
\qquad \text{onde} \qquad
\sigma_g^2 = \operatorname{Var}\big(g(U)\big).
\]

O \textbf{erro-padrão} da estimativa é, portanto, \(\sigma_g / \sqrt{k}\), que decresce
proporcionalmente a \(1/\sqrt{k}\). Como \(\sigma_g^2\) é em geral desconhecido,
estimamo-lo pela variância amostral \(s_g^2\) dos valores \(g(u_i)\), obtendo o
erro-padrão estimado \(s_g/\sqrt{k}\). E, pelo \textbf{teorema limite central}, um
intervalo de confiança aproximado a \(95\%\) para \(\theta\) é

\[
\hat{\theta}_k \pm 1{,}96\,\frac{s_g}{\sqrt{k}}.
\]

A convergência a uma taxa de \(1/\sqrt{k}\) é lenta: para reduzir o erro para
metade é preciso \textbf{quadruplicar} o número de pontos. Esta é a principal
fraqueza do método de Monte Carlo em problemas de baixa dimensão, onde os
métodos determinísticos são muito mais eficientes. A grande virtude do Monte
Carlo revela-se noutro lado: a taxa \(1/\sqrt{k}\) \textbf{não depende da dimensão} do
integral. Um método determinístico baseado numa grelha vê o seu custo explodir
exponencialmente com o número de dimensões (a ``maldição da dimensionalidade''),
ao passo que o Monte Carlo mantém a mesma taxa. É por isso que, em dimensão
elevada, o Monte Carlo se torna imbatível.

\section{\texorpdfstring{Cálculo de \(\int_{a}^{b} g(x)\, dx\)}{Cálculo de \textbackslash int\_\{a\}\^{}\{b\} g(x)\textbackslash, dx}}\label{cuxe1lculo-de-int_ab-gx-dx}

Suponhamos agora que queremos integrar \(g\) num intervalo qualquer \([a, b]\). Basta
uma mudança de variável para nos reconduzir ao caso \((0,1)\). Definamos

\[
y = \frac{x - a}{b - a}, \qquad \text{de modo que} \qquad dy = \frac{dx}{b - a}.
\]

Quando \(x\) percorre \([a, b]\), \(y\) percorre \([0, 1]\), e \(x = a + (b-a)y\). O
integral transforma-se em

\[
\theta = \int_{0}^{1} g\big(a + (b-a)y\big)\,(b-a)\, dy
       = \int_{0}^{1} h(y)\, dy,
\]

onde \(h(y) = (b - a)\, g\big(a + (b - a)y\big)\). Reduzimo-lo, assim, ao caso
padrão, e podemos aplicar de novo o método de Monte Carlo: geram-se números
aleatórios uniformes em \((0,1)\), avalia-se \(h\) nesses pontos e toma-se a média:

\[
\theta \approx \frac{1}{k}\sum_{i=1}^{k} h(y_i).
\]

Existe uma forma equivalente e talvez mais intuitiva de o exprimir: gerar
diretamente pontos uniformes no intervalo \([a, b]\) e multiplicar a média das
avaliações pela amplitude do intervalo,

\[
\theta \approx (b - a)\cdot \frac{1}{k}\sum_{i=1}^{k} g(u_i),
\qquad u_i \sim \mathcal{U}(a, b).
\]

\subsection{Exemplo --- um integral elementar}\label{exemplo-um-integral-elementar}

Ilustremos o método com um integral cujo valor exato conhecemos. Consideremos

\[
\theta = \int_{0}^{1} x^2\, dx = \frac{1}{3}.
\]

Geramos \(k = 10\,000\) números aleatórios em \([0,1]\), avaliamos \(g(u_i) = u_i^2\)
e tomamos a média, que deverá aproximar-se de \(1/3\).

\begin{Shaded}
\begin{Highlighting}[]
\FunctionTok{set.seed}\NormalTok{(}\DecValTok{2024}\NormalTok{)}

\NormalTok{g }\OtherTok{\textless{}{-}} \ControlFlowTok{function}\NormalTok{(x) x}\SpecialCharTok{\^{}}\DecValTok{2}
\NormalTok{k }\OtherTok{\textless{}{-}} \DecValTok{10000}
\NormalTok{u }\OtherTok{\textless{}{-}} \FunctionTok{runif}\NormalTok{(k, }\AttributeTok{min =} \DecValTok{0}\NormalTok{, }\AttributeTok{max =} \DecValTok{1}\NormalTok{)}

\NormalTok{theta\_hat }\OtherTok{\textless{}{-}} \FunctionTok{mean}\NormalTok{(}\FunctionTok{g}\NormalTok{(u))              }\CommentTok{\# estimativa de Monte Carlo}
\NormalTok{theta\_hat}
\DocumentationTok{\#\# [1] 0.3323}

\FunctionTok{integrate}\NormalTok{(g, }\AttributeTok{lower =} \DecValTok{0}\NormalTok{, }\AttributeTok{upper =} \DecValTok{1}\NormalTok{)   }\CommentTok{\# valor de referência}
\DocumentationTok{\#\# 0.3333 with absolute error \textless{} 3.7e{-}15}
\end{Highlighting}
\end{Shaded}

Acrescentemos a estimativa do erro, que raramente deve faltar num cálculo de
Monte Carlo:

\begin{Shaded}
\begin{Highlighting}[]
\NormalTok{erro\_padrao }\OtherTok{\textless{}{-}} \FunctionTok{sd}\NormalTok{(}\FunctionTok{g}\NormalTok{(u)) }\SpecialCharTok{/} \FunctionTok{sqrt}\NormalTok{(k)}
\NormalTok{ic }\OtherTok{\textless{}{-}}\NormalTok{ theta\_hat }\SpecialCharTok{+} \FunctionTok{c}\NormalTok{(}\SpecialCharTok{{-}}\DecValTok{1}\NormalTok{, }\DecValTok{1}\NormalTok{) }\SpecialCharTok{*} \FloatTok{1.96} \SpecialCharTok{*}\NormalTok{ erro\_padrao}
\FunctionTok{c}\NormalTok{(}\AttributeTok{estimativa =}\NormalTok{ theta\_hat, }\AttributeTok{erro\_padrao =}\NormalTok{ erro\_padrao,}
  \AttributeTok{ic\_inf =}\NormalTok{ ic[}\DecValTok{1}\NormalTok{], }\AttributeTok{ic\_sup =}\NormalTok{ ic[}\DecValTok{2}\NormalTok{])}
\DocumentationTok{\#\#  estimativa erro\_padrao      ic\_inf      ic\_sup }
\DocumentationTok{\#\#    0.332306    0.002967    0.326491    0.338121}
\end{Highlighting}
\end{Shaded}

O valor exato \(1/3 \approx 0{,}3333\) cai, como esperado, dentro do intervalo de
confiança. Vejamos agora o integral no intervalo \([1, 2]\),

\[
\int_{1}^{2} x^2\, dx = \frac{7}{3},
\]

aplicando a transformação para \([0,1]\) e, em alternativa, a forma direta com
pontos em \([1, 2]\).

\begin{Shaded}
\begin{Highlighting}[]
\FunctionTok{set.seed}\NormalTok{(}\DecValTok{2024}\NormalTok{)}

\NormalTok{a }\OtherTok{\textless{}{-}} \DecValTok{1}\NormalTok{; b }\OtherTok{\textless{}{-}} \DecValTok{2}
\NormalTok{k }\OtherTok{\textless{}{-}} \DecValTok{10000}

\CommentTok{\# Forma 1: pontos em (0,1) e função transformada h(y)}
\NormalTok{u }\OtherTok{\textless{}{-}} \FunctionTok{runif}\NormalTok{(k)}
\FunctionTok{mean}\NormalTok{((b }\SpecialCharTok{{-}}\NormalTok{ a) }\SpecialCharTok{*} \FunctionTok{g}\NormalTok{(a }\SpecialCharTok{+}\NormalTok{ (b }\SpecialCharTok{{-}}\NormalTok{ a) }\SpecialCharTok{*}\NormalTok{ u))}
\DocumentationTok{\#\# [1] 2.332}

\CommentTok{\# Forma 2: pontos diretamente em (a, b)}
\NormalTok{u2 }\OtherTok{\textless{}{-}} \FunctionTok{runif}\NormalTok{(k, a, b)}
\NormalTok{(b }\SpecialCharTok{{-}}\NormalTok{ a) }\SpecialCharTok{*} \FunctionTok{mean}\NormalTok{(}\FunctionTok{g}\NormalTok{(u2))}
\DocumentationTok{\#\# [1] 2.322}

\CommentTok{\# Valor de referência}
\FunctionTok{integrate}\NormalTok{(g, }\AttributeTok{lower =} \DecValTok{1}\NormalTok{, }\AttributeTok{upper =} \DecValTok{2}\NormalTok{)}
\DocumentationTok{\#\# 2.333 with absolute error \textless{} 2.6e{-}14}
\end{Highlighting}
\end{Shaded}

As duas formas coincidem, como devem, e ambas se aproximam de \(7/3 \approx 2{,}333\).

\subsection{A convergência em imagem}\label{a-converguxeancia-em-imagem}

Vale a pena \emph{ver} a estimativa a convergir à medida que acumulamos pontos. O
gráfico seguinte mostra a estimativa acumulada de \(\int_0^1 x^2\,dx\) a
estabilizar em torno de \(1/3\), ladeada por uma banda de erro que se estreita ao
ritmo de \(1/\sqrt{k}\).

\begin{Shaded}
\begin{Highlighting}[]
\FunctionTok{set.seed}\NormalTok{(}\DecValTok{2024}\NormalTok{)}

\NormalTok{k }\OtherTok{\textless{}{-}} \DecValTok{5000}
\NormalTok{u }\OtherTok{\textless{}{-}} \FunctionTok{runif}\NormalTok{(k)}
\NormalTok{valores }\OtherTok{\textless{}{-}}\NormalTok{ u}\SpecialCharTok{\^{}}\DecValTok{2}
\NormalTok{estim\_acum }\OtherTok{\textless{}{-}} \FunctionTok{cumsum}\NormalTok{(valores) }\SpecialCharTok{/} \FunctionTok{seq\_along}\NormalTok{(valores)}
\CommentTok{\# erro{-}padrão acumulado}
\NormalTok{var\_acum }\OtherTok{\textless{}{-}}\NormalTok{ (}\FunctionTok{cumsum}\NormalTok{(valores}\SpecialCharTok{\^{}}\DecValTok{2}\NormalTok{) }\SpecialCharTok{/} \FunctionTok{seq\_along}\NormalTok{(valores) }\SpecialCharTok{{-}}\NormalTok{ estim\_acum}\SpecialCharTok{\^{}}\DecValTok{2}\NormalTok{)}
\NormalTok{ep\_acum  }\OtherTok{\textless{}{-}} \FunctionTok{sqrt}\NormalTok{(var\_acum }\SpecialCharTok{/} \FunctionTok{seq\_along}\NormalTok{(valores))}

\FunctionTok{plot}\NormalTok{(estim\_acum, }\AttributeTok{type =} \StringTok{"l"}\NormalTok{, }\AttributeTok{col =} \StringTok{"\#10656d"}\NormalTok{, }\AttributeTok{lwd =} \FloatTok{1.2}\NormalTok{,}
     \AttributeTok{ylim =} \FunctionTok{c}\NormalTok{(}\FloatTok{0.25}\NormalTok{, }\FloatTok{0.42}\NormalTok{),}
     \AttributeTok{xlab =} \StringTok{"número de pontos"}\NormalTok{, }\AttributeTok{ylab =} \FunctionTok{expression}\NormalTok{(}\FunctionTok{hat}\NormalTok{(theta)[k]))}
\FunctionTok{lines}\NormalTok{(estim\_acum }\SpecialCharTok{+} \DecValTok{2} \SpecialCharTok{*}\NormalTok{ ep\_acum, }\AttributeTok{col =} \StringTok{"grey70"}\NormalTok{, }\AttributeTok{lty =} \DecValTok{3}\NormalTok{)}
\FunctionTok{lines}\NormalTok{(estim\_acum }\SpecialCharTok{{-}} \DecValTok{2} \SpecialCharTok{*}\NormalTok{ ep\_acum, }\AttributeTok{col =} \StringTok{"grey70"}\NormalTok{, }\AttributeTok{lty =} \DecValTok{3}\NormalTok{)}
\FunctionTok{abline}\NormalTok{(}\AttributeTok{h =} \DecValTok{1}\SpecialCharTok{/}\DecValTok{3}\NormalTok{, }\AttributeTok{lty =} \DecValTok{2}\NormalTok{, }\AttributeTok{col =} \StringTok{"grey30"}\NormalTok{)}
\end{Highlighting}
\end{Shaded}

\begin{figure}

{\centering \includegraphics[width=0.8\linewidth]{introR_files/figure-latex/mc-convergencia-1} 

}

\caption{Estimativa acumulada do integral de x² em [0,1] em função do número de pontos, com banda de dois erros-padrão. A estimativa converge para o valor exato (linha tracejada) e a incerteza estreita-se ao ritmo de $1/\sqrt{k}$.}\label{fig:mc-convergencia}
\end{figure}

\section{\texorpdfstring{Cálculo de \(\int_{0}^{\infty} g(x)\, dx\)}{Cálculo de \textbackslash int\_\{0\}\^{}\{\textbackslash infty\} g(x)\textbackslash, dx}}\label{cuxe1lculo-de-int_0infty-gx-dx}

Para integrais em intervalos ilimitados, uma mudança de variável reconduz-nos
igualmente ao intervalo \((0,1)\). Para \(\int_{0}^{\infty} g(x)\, dx\), aplica-se a
substituição

\[
y = \frac{1}{x + 1},
\qquad \text{de modo que} \qquad
dy = -\frac{dx}{(x+1)^2} = -y^2\, dx.
\]

Quando \(x\) vai de \(0\) a \(+\infty\), \(y\) vai de \(1\) a \(0\), e \(x = \tfrac{1}{y} - 1\).
Trocando os limites (o que absorve o sinal), obtém-se

\[
\theta = \int_{0}^{1} h(y)\, dy,
\qquad \text{onde} \qquad
h(y) = \frac{g\!\left(\tfrac{1}{y} - 1\right)}{y^2}.
\]

Nem sempre esta substituição em particular é a melhor. Uma alternativa muitas
vezes mais eficiente para caudas exponenciais consiste em reconhecer o integral
como um valor esperado relativamente a uma densidade \emph{não uniforme} que
saibamos simular. Por exemplo, \(\int_0^\infty g(x)\,dx = \int_0^\infty
\frac{g(x)}{\lambda e^{-\lambda x}}\,\lambda e^{-\lambda x}\,dx =
E\!\left[\frac{g(X)}{\lambda e^{-\lambda X}}\right]\) com \(X \sim \text{Exp}(\lambda)\).
Esta é a ideia da \emph{amostragem por importância}, em que a escolha da densidade de
amostragem pode reduzir drasticamente a variância do estimador.

\section{\texorpdfstring{O método ``acerta-ou-falha'' e a estimativa de \(\pi\)}{O método ``acerta-ou-falha'' e a estimativa de \textbackslash pi}}\label{o-muxe9todo-acerta-ou-falha-e-a-estimativa-de-pi}

O método que temos usado --- a \textbf{média amostral} --- não é a única via. Há uma
variante geométrica, o método \emph{acerta-ou-falha} (em inglês, \emph{hit-or-miss}), que
recupera diretamente a intuição de Buffon e é particularmente vívido.

A ideia: para estimar a área de uma região contida num retângulo conhecido,
atiram-se pontos ao acaso sobre o retângulo e conta-se a fração que cai \emph{dentro}
da região. Essa fração, multiplicada pela área do retângulo, estima a área
procurada. O exemplo canónico é a estimativa de \(\pi\): um quarto de círculo de
raio \(1\) tem área \(\pi/4\) e está inscrito no quadrado unitário. Atirando pontos
ao acaso no quadrado, a proporção dos que caem dentro do círculo (isto é, com
\(x^2 + y^2 \leq 1\)) aproxima \(\pi/4\).

\begin{Shaded}
\begin{Highlighting}[]
\FunctionTok{set.seed}\NormalTok{(}\DecValTok{2024}\NormalTok{)}

\NormalTok{k }\OtherTok{\textless{}{-}} \DecValTok{100000}
\NormalTok{x }\OtherTok{\textless{}{-}} \FunctionTok{runif}\NormalTok{(k)}
\NormalTok{y }\OtherTok{\textless{}{-}} \FunctionTok{runif}\NormalTok{(k)}
\NormalTok{dentro }\OtherTok{\textless{}{-}}\NormalTok{ x}\SpecialCharTok{\^{}}\DecValTok{2} \SpecialCharTok{+}\NormalTok{ y}\SpecialCharTok{\^{}}\DecValTok{2} \SpecialCharTok{\textless{}=} \DecValTok{1}

\NormalTok{pi\_hat }\OtherTok{\textless{}{-}} \DecValTok{4} \SpecialCharTok{*} \FunctionTok{mean}\NormalTok{(dentro)}
\NormalTok{pi\_hat}
\DocumentationTok{\#\# [1] 3.139}
\end{Highlighting}
\end{Shaded}

\begin{Shaded}
\begin{Highlighting}[]
\FunctionTok{set.seed}\NormalTok{(}\DecValTok{2024}\NormalTok{)}
\NormalTok{n\_plot }\OtherTok{\textless{}{-}} \DecValTok{3000}
\NormalTok{xp }\OtherTok{\textless{}{-}} \FunctionTok{runif}\NormalTok{(n\_plot); yp }\OtherTok{\textless{}{-}} \FunctionTok{runif}\NormalTok{(n\_plot)}
\NormalTok{dentro\_p }\OtherTok{\textless{}{-}}\NormalTok{ xp}\SpecialCharTok{\^{}}\DecValTok{2} \SpecialCharTok{+}\NormalTok{ yp}\SpecialCharTok{\^{}}\DecValTok{2} \SpecialCharTok{\textless{}=} \DecValTok{1}
\FunctionTok{plot}\NormalTok{(xp, yp, }\AttributeTok{pch =} \DecValTok{19}\NormalTok{, }\AttributeTok{cex =} \FloatTok{0.3}\NormalTok{, }\AttributeTok{asp =} \DecValTok{1}\NormalTok{,}
     \AttributeTok{col =} \FunctionTok{ifelse}\NormalTok{(dentro\_p, }\StringTok{"\#10656d"}\NormalTok{, }\StringTok{"grey70"}\NormalTok{),}
     \AttributeTok{xlab =} \StringTok{"x"}\NormalTok{, }\AttributeTok{ylab =} \StringTok{"y"}\NormalTok{)}
\FunctionTok{curve}\NormalTok{(}\FunctionTok{sqrt}\NormalTok{(}\DecValTok{1} \SpecialCharTok{{-}}\NormalTok{ x}\SpecialCharTok{\^{}}\DecValTok{2}\NormalTok{), }\AttributeTok{from =} \DecValTok{0}\NormalTok{, }\AttributeTok{to =} \DecValTok{1}\NormalTok{, }\AttributeTok{add =} \ConstantTok{TRUE}\NormalTok{, }\AttributeTok{lwd =} \DecValTok{2}\NormalTok{, }\AttributeTok{col =} \StringTok{"grey20"}\NormalTok{)}
\end{Highlighting}
\end{Shaded}

\begin{figure}

{\centering \includegraphics[width=0.8\linewidth]{introR_files/figure-latex/mc-pi-grafico-1} 

}

\caption{Estimativa de \(\pi\) pelo método acerta-ou-falha: pontos aleatórios no quadrado unitário. Os que caem dentro do quarto de círculo (a azul) representam uma fração próxima de \(\pi\)/4.}\label{fig:mc-pi-grafico}
\end{figure}

O método acerta-ou-falha é um caso particular do método da média amostral,
aplicado à função indicadora \(g(x,y) = \mathbf{1}\{x^2 + y^2 \leq 1\}\). Por essa
razão, a sua variância é tipicamente \emph{maior} do que a de estimadores baseados na
média de uma função suave, pelo que, quando existe uma alternativa por média
amostral, esta costuma ser preferível. O método acerta-ou-falha vale sobretudo
pela transparência geométrica e pelo seu valor histórico e pedagógico.

\section{Integrais multidimensionais}\label{integrais-multidimensionais}

É na dimensão elevada que o método de Monte Carlo mostra todo o seu valor.
Suponhamos que \(g\) é uma função de \(n\) argumentos e que se pretende calcular

\[
\theta = \int_{0}^{1}\!\int_{0}^{1}\!\cdots\!\int_{0}^{1}
g(x_1, \ldots, x_n)\, dx_1\, dx_2 \cdots dx_n.
\]

Tal como no caso unidimensional, este integral é um valor esperado:
\(\theta = E\big[g(U_1, \ldots, U_n)\big]\), onde \(U_1, \ldots, U_n\) são uniformes
independentes em \((0,1)\). Para o estimar, geram-se \(k\) conjuntos independentes,
cada um com \(n\) coordenadas uniformes independentes,

\begin{align*}
U_{1}^{1}, & \ldots, U_{n}^{1} \\
U_{1}^{2}, & \ldots, U_{n}^{2} \\
           & \;\;\vdots \\
U_{1}^{k}, & \ldots, U_{n}^{k},
\end{align*}

e as variáveis \(g(U_{1}^{i}, \ldots, U_{n}^{i})\), \(i = 1, \ldots, k\), são i.i.d.
com média \(\theta\). Estima-se então

\[
\theta \approx \frac{1}{k}\sum_{i=1}^{k} g\big(U_{1}^{i}, \ldots, U_{n}^{i}\big).
\]

\subsection{Exemplo --- um integral duplo}\label{exemplo-um-integral-duplo}

Estimemos

\[
\theta = \int_{0}^{1}\!\int_{0}^{1} (x^2 + y^2)\, dx\, dy,
\]

cujo valor exato é \(\tfrac{1}{3} + \tfrac{1}{3} = \tfrac{2}{3}\).

\begin{Shaded}
\begin{Highlighting}[]
\FunctionTok{set.seed}\NormalTok{(}\DecValTok{2024}\NormalTok{)}

\NormalTok{k }\OtherTok{\textless{}{-}} \DecValTok{100000}
\NormalTok{x }\OtherTok{\textless{}{-}} \FunctionTok{runif}\NormalTok{(k)}
\NormalTok{y }\OtherTok{\textless{}{-}} \FunctionTok{runif}\NormalTok{(k)}
\NormalTok{g2 }\OtherTok{\textless{}{-}}\NormalTok{ x}\SpecialCharTok{\^{}}\DecValTok{2} \SpecialCharTok{+}\NormalTok{ y}\SpecialCharTok{\^{}}\DecValTok{2}

\NormalTok{theta\_hat }\OtherTok{\textless{}{-}} \FunctionTok{mean}\NormalTok{(g2)}
\NormalTok{erro\_padrao }\OtherTok{\textless{}{-}} \FunctionTok{sd}\NormalTok{(g2) }\SpecialCharTok{/} \FunctionTok{sqrt}\NormalTok{(k)}

\FunctionTok{c}\NormalTok{(}\AttributeTok{estimativa =}\NormalTok{ theta\_hat, }\AttributeTok{exato =} \DecValTok{2}\SpecialCharTok{/}\DecValTok{3}\NormalTok{, }\AttributeTok{erro\_padrao =}\NormalTok{ erro\_padrao)}
\DocumentationTok{\#\#  estimativa       exato erro\_padrao }
\DocumentationTok{\#\#    0.665648    0.666667    0.001332}
\end{Highlighting}
\end{Shaded}

A estimativa aproxima-se de \(2/3\) com uma incerteza controlada --- e, o que é
crucial, a mesma abordagem funcionaria sem alterações significativas para um
integral em \(10\) ou \(100\) dimensões, onde qualquer método de grelha seria
impensável.

\section{Síntese}\label{suxedntese-2}

O método de Monte Carlo transforma o problema de calcular um integral no
problema de estimar um valor esperado, resolvido pela média de avaliações da
função em pontos aleatórios. A sua justificação é a lei dos grandes números; a
quantificação do seu erro, o teorema limite central. Converge lentamente, ao
ritmo de \(1/\sqrt{k}\), mas com a vantagem decisiva de essa taxa ser independente
da dimensão --- o que faz dele o método de eleição para integrais de dimensão
elevada, omnipresentes na estatística moderna, na física e nas finanças. Nos
casos aqui tratados usámos a distribuição uniforme como base; a generalização
para outras distribuições de amostragem, com vista a reduzir a variância, abre a
porta a técnicas mais avançadas, como a amostragem por importância.

\section{Exercícios}\label{exercuxedcios-37}

Use simulação para aproximar os integrais seguintes. Sempre que possível,
compare a sua estimativa com o valor exato e apresente também o erro-padrão e um
intervalo de confiança a \(95\%\).

\textbf{1.} \(\displaystyle\int_{0}^{1} \exp(e^x)\, dx\).

\textbf{2.} \(\displaystyle\int_{0}^{1} (1 - x^2)^{3/2}\, dx\).

\textbf{3.} \(\displaystyle\int_{-2}^{2} e^{x + x^2}\, dx\).

\textbf{4.} \(\displaystyle\int_{0}^{\pi/2} \cos(x)\, dx\). (Valor exato: \(1\).)

\textbf{5.} \(\displaystyle\int_{0}^{\infty} x\,(1 + x^2)^{-2}\, dx\). (Valor exato:
\(1/2\).)

\textbf{6.} \(\displaystyle\int_{-\infty}^{\infty} e^{-x^2}\, dx\). (Valor exato:
\(\sqrt{\pi}\).) Sugestão: separe o integral em duas metades ou explore a simetria.

\textbf{7.} \(\displaystyle\int_{0}^{1}\!\int_{0}^{1} e^{(x+y)^2}\, dy\, dx\).

\textbf{8.} \(\displaystyle\int_{0}^{1}\!\int_{0}^{1} (x^2 + y^2)\, dx\, dy\). (Valor
exato: \(2/3\).)

\textbf{9.} \emph{(Estimativa de \(\pi\) e taxa de convergência.)} Estime \(\pi\) pelo método
acerta-ou-falha e represente graficamente o valor absoluto do erro
\(|\hat{\pi}_k - \pi|\) em função de \(k\), em escala logarítmica em ambos os eixos.
Confirme visualmente que o erro decresce aproximadamente como \(1/\sqrt{k}\)
(declive \(-1/2\)).

\textbf{10.} \emph{(A agulha de Buffon.)} Simule o lançamento de \(10\,000\) agulhas de
comprimento \(\ell = 1\) sobre um soalho de linhas paralelas afastadas de \(d = 2\).
Sabendo que a probabilidade de uma agulha cruzar uma linha é
\(\frac{2\ell}{\pi d}\), use a proporção de cruzamentos para estimar \(\pi\). Compare
a eficiência com o método acerta-ou-falha do exercício anterior.

\textbf{11.} \emph{(Erro em função de \(k\).)} Para o integral \(\int_0^1 e^x\, dx = e - 1\),
estime \(\theta\) com \(k = 10^2, 10^3, 10^4, 10^5, 10^6\) e registe o erro absoluto
de cada estimativa. Represente o erro em função de \(k\) numa escala log-log e
confirme a taxa \(1/\sqrt{k}\).

\textbf{12.} \emph{(Redução de variância por variáveis antitéticas.)} Para estimar
\(\int_0^1 e^x\, dx\), compare o estimador de Monte Carlo simples com o estimador
\emph{antitético}, que usa os pares \(\big(g(U_i), g(1 - U_i)\big)\). Estime a variância
de cada estimador por repetição e quantifique a redução obtida. Explique por que
razão a técnica é eficaz quando \(g\) é monótona.

\textbf{13.} \emph{(Alta dimensão.)} Estime o volume da bola unitária em dimensão \(n\),
\(\displaystyle V_n = \int_{[-1,1]^n} \mathbf{1}\{x_1^2 + \cdots + x_n^2 \leq 1\}\,
dx_1\cdots dx_n\), pelo método acerta-ou-falha, para \(n = 2, 3, 5, 10\). Compare com
a fórmula exata \(V_n = \pi^{n/2}/\Gamma(n/2 + 1)\) e comente como a proporção de
acertos (e, portanto, a eficiência) se degrada à medida que a dimensão aumenta.

\chapter{Resoluções dos exercícios}\label{solucoes}

Este capítulo final reúne as resoluções dos exercícios propostos ao longo do
livro. Mais do que um mero solucionário, pretende ser um capítulo de estudo:
sempre que um exercício admite mais do que um caminho, apresentamos as
alternativas mais instrutivas; sempre que a resposta envolve um resultado
teórico, explicitamos o raciocínio antes do código; e sempre que uma solução
comum contém uma armadilha subtil, assinalamo-la com uma caixa de \textbf{Atenção}.

\textbf{Como ler este capítulo.} O código das resoluções está, por defeito, com a
opção \texttt{eval=FALSE}: é para ser executado pelo leitor, não para produzir
resultados impressos aqui. Aconselhamos vivamente a que o experimente e o
modifique --- trocar parâmetros, aumentar o número de simulações, alterar
sementes. As resoluções seguem a ordem temática do livro; dentro de cada tema,
a numeração acompanha a dos enunciados. Quando existe mais do que uma forma
correta de resolver um problema, indicamo-la com o comentário \texttt{\#\ Alternativa}.

\section{Vetores}\label{vetores}

\textbf{1.} Criação de vetores com \texttt{seq()}, \texttt{rep()} e o operador \texttt{:}.

\begin{Shaded}
\begin{Highlighting}[]
\CommentTok{\# (a) 1, 2, ..., 20}
\DecValTok{1}\SpecialCharTok{:}\DecValTok{20}

\CommentTok{\# (b) 20, 19, ..., 1}
\DecValTok{20}\SpecialCharTok{:}\DecValTok{1}

\CommentTok{\# (c) 1, ..., 20, 19, ..., 1}
\FunctionTok{c}\NormalTok{(}\DecValTok{1}\SpecialCharTok{:}\DecValTok{20}\NormalTok{, }\DecValTok{19}\SpecialCharTok{:}\DecValTok{1}\NormalTok{)}

\CommentTok{\# (d) 10, 20, ..., 90, 10}
\FunctionTok{c}\NormalTok{(}\FunctionTok{seq}\NormalTok{(}\DecValTok{10}\NormalTok{, }\DecValTok{90}\NormalTok{, }\AttributeTok{by =} \DecValTok{10}\NormalTok{), }\DecValTok{10}\NormalTok{)}

\CommentTok{\# (e) dez 1, vinte 2, trinta 3}
\FunctionTok{c}\NormalTok{(}\FunctionTok{rep}\NormalTok{(}\DecValTok{1}\NormalTok{, }\DecValTok{10}\NormalTok{), }\FunctionTok{rep}\NormalTok{(}\DecValTok{2}\NormalTok{, }\DecValTok{20}\NormalTok{), }\FunctionTok{rep}\NormalTok{(}\DecValTok{3}\NormalTok{, }\DecValTok{30}\NormalTok{))}
\end{Highlighting}
\end{Shaded}

\textbf{2.} Vetor de caracteres \texttt{"nome\ 1",\ ...,\ "nome\ 20"} com \texttt{paste()}.

\begin{Shaded}
\begin{Highlighting}[]
\FunctionTok{paste}\NormalTok{(}\StringTok{"nome"}\NormalTok{, }\DecValTok{1}\SpecialCharTok{:}\DecValTok{20}\NormalTok{)}
\end{Highlighting}
\end{Shaded}

\textbf{3. e 4.} Repetições com \texttt{each} e \texttt{times}. A diferença entre os dois argumentos
é essencial: \texttt{each} repete cada elemento em bloco; \texttt{times} repete o vetor
inteiro.

\begin{Shaded}
\begin{Highlighting}[]
\CommentTok{\# 3. "A" "A" "B" "B" ... "E" "E"}
\FunctionTok{rep}\NormalTok{(LETTERS[}\DecValTok{1}\SpecialCharTok{:}\DecValTok{5}\NormalTok{], }\AttributeTok{each =} \DecValTok{2}\NormalTok{)}

\CommentTok{\# 4. "a" "b" ... "e" "a" "b" ... "e"}
\FunctionTok{rep}\NormalTok{(letters[}\DecValTok{1}\SpecialCharTok{:}\DecValTok{5}\NormalTok{], }\AttributeTok{times =} \DecValTok{2}\NormalTok{)}
\end{Highlighting}
\end{Shaded}

\textbf{5.} Repetições com contagens diferentes.

\begin{Shaded}
\begin{Highlighting}[]
\FunctionTok{c}\NormalTok{(}\FunctionTok{rep}\NormalTok{(}\StringTok{"uva"}\NormalTok{, }\DecValTok{10}\NormalTok{), }\FunctionTok{rep}\NormalTok{(}\StringTok{"maçã"}\NormalTok{, }\DecValTok{9}\NormalTok{), }\FunctionTok{rep}\NormalTok{(}\StringTok{"laranja"}\NormalTok{, }\DecValTok{6}\NormalTok{), }\FunctionTok{rep}\NormalTok{(}\StringTok{"banana"}\NormalTok{, }\DecValTok{1}\NormalTok{))}
\end{Highlighting}
\end{Shaded}

\textbf{6.} Nomear elementos de um vetor e aceder por nome.

\begin{Shaded}
\begin{Highlighting}[]
\NormalTok{notas }\OtherTok{\textless{}{-}} \FunctionTok{c}\NormalTok{(}\DecValTok{7}\NormalTok{, }\DecValTok{8}\NormalTok{, }\DecValTok{9}\NormalTok{)}
\FunctionTok{names}\NormalTok{(notas) }\OtherTok{\textless{}{-}} \FunctionTok{c}\NormalTok{(}\StringTok{"Ana"}\NormalTok{, }\StringTok{"João"}\NormalTok{, }\StringTok{"Pedro"}\NormalTok{)}
\NormalTok{notas[}\StringTok{"João"}\NormalTok{]}
\end{Highlighting}
\end{Shaded}

\textbf{7.} Ordenar e localizar extremos.

\begin{Shaded}
\begin{Highlighting}[]
\NormalTok{x }\OtherTok{\textless{}{-}} \FunctionTok{sample}\NormalTok{(}\DecValTok{1}\SpecialCharTok{:}\DecValTok{100}\NormalTok{, }\AttributeTok{size =} \DecValTok{15}\NormalTok{)}
\FunctionTok{sort}\NormalTok{(x)                    }\CommentTok{\# crescente}
\FunctionTok{sort}\NormalTok{(x, }\AttributeTok{decreasing =} \ConstantTok{TRUE}\NormalTok{) }\CommentTok{\# decrescente}
\FunctionTok{min}\NormalTok{(x); }\FunctionTok{max}\NormalTok{(x)}
\end{Highlighting}
\end{Shaded}

\textbf{8.} Estatísticas de resumo de um vetor.

\begin{Shaded}
\begin{Highlighting}[]
\NormalTok{x }\OtherTok{\textless{}{-}} \FunctionTok{sample}\NormalTok{(}\DecValTok{1}\SpecialCharTok{:}\DecValTok{50}\NormalTok{, }\AttributeTok{size =} \DecValTok{20}\NormalTok{)}
\FunctionTok{sum}\NormalTok{(x); }\FunctionTok{mean}\NormalTok{(x); }\FunctionTok{sd}\NormalTok{(x); }\FunctionTok{prod}\NormalTok{(x)}
\end{Highlighting}
\end{Shaded}

\textbf{9.} Indexação lógica: extrair e substituir por condição.

\begin{Shaded}
\begin{Highlighting}[]
\NormalTok{x }\OtherTok{\textless{}{-}} \FunctionTok{sample}\NormalTok{(}\DecValTok{1}\SpecialCharTok{:}\DecValTok{100}\NormalTok{, }\AttributeTok{size =} \DecValTok{10}\NormalTok{)}
\NormalTok{x[x }\SpecialCharTok{\textgreater{}} \DecValTok{50}\NormalTok{]        }\CommentTok{\# elementos maiores que 50}
\NormalTok{x[x }\SpecialCharTok{\textless{}} \DecValTok{30}\NormalTok{] }\OtherTok{\textless{}{-}} \DecValTok{0}   \CommentTok{\# substitui in loco os menores que 30}
\end{Highlighting}
\end{Shaded}

\textbf{10.} Combinar condições com \texttt{\&}.

\begin{Shaded}
\begin{Highlighting}[]
\NormalTok{x }\OtherTok{\textless{}{-}} \FunctionTok{sample}\NormalTok{(}\DecValTok{1}\SpecialCharTok{:}\DecValTok{10}\NormalTok{)}
\NormalTok{x[x }\SpecialCharTok{\textgreater{}} \DecValTok{5} \SpecialCharTok{\&}\NormalTok{ x }\SpecialCharTok{\%\%} \DecValTok{2} \SpecialCharTok{==} \DecValTok{0}\NormalTok{]   }\CommentTok{\# maiores que 5 E pares}
\end{Highlighting}
\end{Shaded}

\textbf{11.} Indexação por posição.

\begin{Shaded}
\begin{Highlighting}[]
\NormalTok{x }\OtherTok{\textless{}{-}} \FunctionTok{sample}\NormalTok{(}\DecValTok{1}\SpecialCharTok{:}\DecValTok{10}\NormalTok{)}
\NormalTok{x[}\FunctionTok{c}\NormalTok{(}\DecValTok{2}\NormalTok{, }\DecValTok{4}\NormalTok{, }\DecValTok{6}\NormalTok{)] }\OtherTok{\textless{}{-}}\NormalTok{ x[}\FunctionTok{c}\NormalTok{(}\DecValTok{2}\NormalTok{, }\DecValTok{4}\NormalTok{, }\DecValTok{6}\NormalTok{)] }\SpecialCharTok{*} \DecValTok{2}
\NormalTok{x[}\FunctionTok{length}\NormalTok{(x)] }\OtherTok{\textless{}{-}} \DecValTok{100}      \CommentTok{\# seq\_along/length evita erros de índice}
\end{Highlighting}
\end{Shaded}

\textbf{12.} Médias na presença de valores especiais.

\begin{Shaded}
\begin{Highlighting}[]
\CommentTok{\# (a) com NA}
\NormalTok{x }\OtherTok{\textless{}{-}} \FunctionTok{c}\NormalTok{(}\DecValTok{1}\NormalTok{, }\DecValTok{0}\NormalTok{, }\ConstantTok{NA}\NormalTok{, }\DecValTok{5}\NormalTok{, }\DecValTok{7}\NormalTok{)}
\FunctionTok{mean}\NormalTok{(x, }\AttributeTok{na.rm =} \ConstantTok{TRUE}\NormalTok{)}

\CommentTok{\# (b) com infinitos: primeiro convertê{-}los em NA}
\NormalTok{y }\OtherTok{\textless{}{-}} \FunctionTok{c}\NormalTok{(}\SpecialCharTok{{-}}\ConstantTok{Inf}\NormalTok{, }\DecValTok{0}\NormalTok{, }\DecValTok{1}\NormalTok{, }\ConstantTok{NA}\NormalTok{, }\DecValTok{2}\NormalTok{, }\DecValTok{7}\NormalTok{, }\ConstantTok{Inf}\NormalTok{)}
\FunctionTok{mean}\NormalTok{(y[}\SpecialCharTok{!}\FunctionTok{is.infinite}\NormalTok{(y)], }\AttributeTok{na.rm =} \ConstantTok{TRUE}\NormalTok{)}
\end{Highlighting}
\end{Shaded}

\texttt{mean(y,\ na.rm\ =\ TRUE)} \textbf{não} remove \texttt{-Inf} nem \texttt{Inf}: \texttt{na.rm} só ignora \texttt{NA}.
Com infinitos presentes, a média devolve \texttt{NaN}. É preciso filtrá-los
explicitamente com \texttt{is.infinite()}, como acima.

\textbf{13.} Vetores construídos a partir de expressões.

\begin{Shaded}
\begin{Highlighting}[]
\CommentTok{\# (a)}
\NormalTok{x }\OtherTok{\textless{}{-}} \FunctionTok{seq}\NormalTok{(}\DecValTok{2}\NormalTok{, }\DecValTok{6}\NormalTok{, }\AttributeTok{by =} \FloatTok{0.1}\NormalTok{)}
\FunctionTok{exp}\NormalTok{(x) }\SpecialCharTok{*} \FunctionTok{sin}\NormalTok{(x)}

\CommentTok{\# (b) (3, 3\^{}2/2, 3\^{}3/3, ..., 3\^{}30/30)}
\NormalTok{i }\OtherTok{\textless{}{-}} \DecValTok{1}\SpecialCharTok{:}\DecValTok{30}
\DecValTok{3}\SpecialCharTok{\^{}}\NormalTok{i }\SpecialCharTok{/}\NormalTok{ i}
\end{Highlighting}
\end{Shaded}

\textbf{14.} Somatórios com vetorização.

\begin{Shaded}
\begin{Highlighting}[]
\CommentTok{\# (a)}
\NormalTok{i }\OtherTok{\textless{}{-}} \DecValTok{1}\SpecialCharTok{:}\DecValTok{1000}\NormalTok{; }\FunctionTok{sum}\NormalTok{(}\DecValTok{9} \SpecialCharTok{/} \DecValTok{10}\SpecialCharTok{\^{}}\NormalTok{i)}

\CommentTok{\# (b)}
\NormalTok{i }\OtherTok{\textless{}{-}} \DecValTok{10}\SpecialCharTok{:}\DecValTok{100}\NormalTok{; }\FunctionTok{sum}\NormalTok{(i}\SpecialCharTok{\^{}}\DecValTok{3} \SpecialCharTok{+} \DecValTok{4} \SpecialCharTok{*}\NormalTok{ i}\SpecialCharTok{\^{}}\DecValTok{2}\NormalTok{)}

\CommentTok{\# (c)}
\NormalTok{i }\OtherTok{\textless{}{-}} \DecValTok{1}\SpecialCharTok{:}\DecValTok{25}\NormalTok{; }\FunctionTok{sum}\NormalTok{(}\DecValTok{2}\SpecialCharTok{\^{}}\NormalTok{i }\SpecialCharTok{/}\NormalTok{ i }\SpecialCharTok{+} \DecValTok{3}\SpecialCharTok{\^{}}\NormalTok{i }\SpecialCharTok{/}\NormalTok{ i}\SpecialCharTok{\^{}}\DecValTok{2}\NormalTok{)}
\end{Highlighting}
\end{Shaded}

\textbf{15.} Consumo de energia (aplicação integradora).

\begin{Shaded}
\begin{Highlighting}[]
\FunctionTok{set.seed}\NormalTok{(}\DecValTok{2025}\NormalTok{)}
\NormalTok{consumo }\OtherTok{\textless{}{-}} \FunctionTok{runif}\NormalTok{(}\AttributeTok{n =} \DecValTok{30}\NormalTok{, }\AttributeTok{min =} \DecValTok{10}\NormalTok{, }\AttributeTok{max =} \DecValTok{40}\NormalTok{)   }\CommentTok{\# (a)}
\NormalTok{media   }\OtherTok{\textless{}{-}} \FunctionTok{mean}\NormalTok{(consumo)                       }\CommentTok{\# (b)}
\NormalTok{acima\_da\_media }\OtherTok{\textless{}{-}}\NormalTok{ consumo }\SpecialCharTok{\textgreater{}}\NormalTok{ media              }\CommentTok{\# (c)}
\FunctionTok{sum}\NormalTok{(acima\_da\_media)                            }\CommentTok{\# (d) quantos dias}
\FunctionTok{which}\NormalTok{(consumo }\SpecialCharTok{\textgreater{}} \DecValTok{35}\NormalTok{)                            }\CommentTok{\# (e) posições de pico}
\NormalTok{consumo[}\FunctionTok{which}\NormalTok{(consumo }\SpecialCharTok{\textgreater{}} \DecValTok{35}\NormalTok{)]                   }\CommentTok{\#     valores de pico}
\NormalTok{consumo[consumo }\SpecialCharTok{\textless{}} \DecValTok{15}\NormalTok{] }\OtherTok{\textless{}{-}} \ConstantTok{NA}                    \CommentTok{\# (f) falhas de medição}
\FunctionTok{mean}\NormalTok{(consumo, }\AttributeTok{na.rm =} \ConstantTok{TRUE}\NormalTok{)                    }\CommentTok{\# (g) média corrigida}
\NormalTok{consumo[}\FunctionTok{seq}\NormalTok{(}\DecValTok{2}\NormalTok{, }\DecValTok{30}\NormalTok{, }\AttributeTok{by =} \DecValTok{2}\NormalTok{)]                    }\CommentTok{\# (h) dias pares}
\end{Highlighting}
\end{Shaded}

\textbf{16.} Produção agrícola (aplicação integradora).

\begin{Shaded}
\begin{Highlighting}[]
\NormalTok{producao }\OtherTok{\textless{}{-}} \FunctionTok{c}\NormalTok{(}\DecValTok{80}\NormalTok{, }\DecValTok{95}\NormalTok{, }\DecValTok{70}\NormalTok{, }\DecValTok{68}\NormalTok{, }\DecValTok{105}\NormalTok{, }\DecValTok{120}\NormalTok{, }\DecValTok{73}\NormalTok{, }\DecValTok{89}\NormalTok{, }\DecValTok{110}\NormalTok{, }\DecValTok{66}\NormalTok{, }\DecValTok{92}\NormalTok{, }\DecValTok{78}\NormalTok{)}
\NormalTok{media\_total }\OtherTok{\textless{}{-}} \FunctionTok{mean}\NormalTok{(producao)                          }\CommentTok{\# (b)}
\FunctionTok{which}\NormalTok{(producao }\SpecialCharTok{\textless{}}\NormalTok{ media\_total)                          }\CommentTok{\# (c) abaixo da média}
\NormalTok{producao[}\FunctionTok{c}\NormalTok{(}\DecValTok{4}\NormalTok{, }\DecValTok{10}\NormalTok{, }\DecValTok{12}\NormalTok{)] }\OtherTok{\textless{}{-}}\NormalTok{ producao[}\FunctionTok{c}\NormalTok{(}\DecValTok{4}\NormalTok{, }\DecValTok{10}\NormalTok{, }\DecValTok{12}\NormalTok{)] }\SpecialCharTok{*} \FloatTok{1.10} \CommentTok{\# (d) correção +10\%}
\FunctionTok{mean}\NormalTok{(producao)                                         }\CommentTok{\# (e) nova média}
\NormalTok{alta }\OtherTok{\textless{}{-}}\NormalTok{ producao }\SpecialCharTok{\textgreater{}} \DecValTok{100}                                 \CommentTok{\# (f)}
\FunctionTok{sum}\NormalTok{(alta); }\FunctionTok{sum}\NormalTok{(alta) }\SpecialCharTok{/} \FunctionTok{length}\NormalTok{(producao)                }\CommentTok{\# (g) contagem e proporção}
\NormalTok{ordenado }\OtherTok{\textless{}{-}} \FunctionTok{sort}\NormalTok{(producao)                             }\CommentTok{\# (h)}
\FunctionTok{tail}\NormalTok{(ordenado, }\DecValTok{3}\NormalTok{)                                       }\CommentTok{\# (i) as 3 maiores}
\end{Highlighting}
\end{Shaded}

Aumentar em 10\% escreve-se de forma mais limpa como \texttt{producao\ *\ 1.10} do que
\texttt{producao\ +\ producao\ *\ 0.10}. As duas expressões dão o mesmo resultado, mas a
primeira lê-se de imediato como ``acrescentar dez por cento''.

\section{Fatores}\label{fatores-1}

\textbf{1.} Criar um fator com níveis explícitos.

\begin{Shaded}
\begin{Highlighting}[]
\NormalTok{tamanhos }\OtherTok{\textless{}{-}} \FunctionTok{c}\NormalTok{(}\StringTok{"pequeno"}\NormalTok{, }\StringTok{"médio"}\NormalTok{, }\StringTok{"grande"}\NormalTok{, }\StringTok{"pequeno"}\NormalTok{, }\StringTok{"grande"}\NormalTok{)}
\NormalTok{f }\OtherTok{\textless{}{-}} \FunctionTok{factor}\NormalTok{(tamanhos, }\AttributeTok{levels =} \FunctionTok{c}\NormalTok{(}\StringTok{"pequeno"}\NormalTok{, }\StringTok{"médio"}\NormalTok{, }\StringTok{"grande"}\NormalTok{))}
\FunctionTok{levels}\NormalTok{(f)}
\end{Highlighting}
\end{Shaded}

\textbf{2.} Tabela de frequências de um fator.

\begin{Shaded}
\begin{Highlighting}[]
\NormalTok{cores }\OtherTok{\textless{}{-}} \FunctionTok{factor}\NormalTok{(}\FunctionTok{c}\NormalTok{(}\StringTok{"vermelho"}\NormalTok{, }\StringTok{"azul"}\NormalTok{, }\StringTok{"preto"}\NormalTok{, }\StringTok{"vermelho"}\NormalTok{, }\StringTok{"branco"}\NormalTok{, }\StringTok{"azul"}\NormalTok{))}
\FunctionTok{table}\NormalTok{(cores)}
\end{Highlighting}
\end{Shaded}

\textbf{3.} Renomear níveis.

\begin{Shaded}
\begin{Highlighting}[]
\NormalTok{avaliacoes }\OtherTok{\textless{}{-}} \FunctionTok{factor}\NormalTok{(}\FunctionTok{c}\NormalTok{(}\StringTok{"bom"}\NormalTok{, }\StringTok{"médio"}\NormalTok{, }\StringTok{"ruim"}\NormalTok{, }\StringTok{"bom"}\NormalTok{, }\StringTok{"ruim"}\NormalTok{))}
\FunctionTok{levels}\NormalTok{(avaliacoes) }\OtherTok{\textless{}{-}} \FunctionTok{c}\NormalTok{(}\StringTok{"excelente"}\NormalTok{, }\StringTok{"razoável"}\NormalTok{, }\StringTok{"insatisfatório"}\NormalTok{)}
\NormalTok{avaliacoes}
\end{Highlighting}
\end{Shaded}

\textbf{4.} Testar e converter para fator.

\begin{Shaded}
\begin{Highlighting}[]
\NormalTok{estado\_civil }\OtherTok{\textless{}{-}} \FunctionTok{c}\NormalTok{(}\StringTok{"solteiro"}\NormalTok{, }\StringTok{"casado"}\NormalTok{, }\StringTok{"divorciado"}\NormalTok{)}
\FunctionTok{is.factor}\NormalTok{(estado\_civil)              }\CommentTok{\# FALSE}
\NormalTok{estado\_civil }\OtherTok{\textless{}{-}} \FunctionTok{factor}\NormalTok{(estado\_civil)}
\FunctionTok{is.factor}\NormalTok{(estado\_civil)              }\CommentTok{\# TRUE}
\end{Highlighting}
\end{Shaded}

\textbf{5.} Fator ordenado e ordenação.

\begin{Shaded}
\begin{Highlighting}[]
\NormalTok{escolaridade }\OtherTok{\textless{}{-}} \FunctionTok{factor}\NormalTok{(}
  \FunctionTok{c}\NormalTok{(}\StringTok{"mestrado"}\NormalTok{, }\StringTok{"doutoramento"}\NormalTok{, }\StringTok{"licenciatura"}\NormalTok{, }\StringTok{"secundário"}\NormalTok{),}
  \AttributeTok{levels =} \FunctionTok{c}\NormalTok{(}\StringTok{"secundário"}\NormalTok{, }\StringTok{"licenciatura"}\NormalTok{, }\StringTok{"mestrado"}\NormalTok{, }\StringTok{"doutoramento"}\NormalTok{),}
  \AttributeTok{ordered =} \ConstantTok{TRUE}\NormalTok{)}
\FunctionTok{sort}\NormalTok{(escolaridade)}
\end{Highlighting}
\end{Shaded}

\textbf{6.} Converter um fator em códigos numéricos.

\begin{Shaded}
\begin{Highlighting}[]
\NormalTok{pressao }\OtherTok{\textless{}{-}} \FunctionTok{factor}\NormalTok{(}\FunctionTok{c}\NormalTok{(}\StringTok{"baixo"}\NormalTok{, }\StringTok{"alto"}\NormalTok{, }\StringTok{"médio"}\NormalTok{, }\StringTok{"baixo"}\NormalTok{, }\StringTok{"alto"}\NormalTok{),}
                  \AttributeTok{levels =} \FunctionTok{c}\NormalTok{(}\StringTok{"baixo"}\NormalTok{, }\StringTok{"médio"}\NormalTok{, }\StringTok{"alto"}\NormalTok{))}
\FunctionTok{as.numeric}\NormalTok{(pressao)   }\CommentTok{\# 1, 3, 2, 1, 3 — segue a ordem dos níveis}
\end{Highlighting}
\end{Shaded}

\texttt{as.numeric()} de um fator devolve os \textbf{códigos internos dos níveis}, não o
texto. Se o fator contiver números guardados como texto (por exemplo
\texttt{factor(c("10","20"))}), \texttt{as.numeric()} devolve \texttt{1,\ 2} e não \texttt{10,\ 20}. Nesse
caso use \texttt{as.numeric(as.character(f))}.

\textbf{7.} Substituir um nível específico.

\begin{Shaded}
\begin{Highlighting}[]
\NormalTok{cores }\OtherTok{\textless{}{-}} \FunctionTok{factor}\NormalTok{(}\FunctionTok{c}\NormalTok{(}\StringTok{"verde"}\NormalTok{, }\StringTok{"amarelo"}\NormalTok{, }\StringTok{"vermelho"}\NormalTok{, }\StringTok{"verde"}\NormalTok{, }\StringTok{"amarelo"}\NormalTok{))}
\FunctionTok{levels}\NormalTok{(cores)[}\FunctionTok{levels}\NormalTok{(cores) }\SpecialCharTok{==} \StringTok{"amarelo"}\NormalTok{] }\OtherTok{\textless{}{-}} \StringTok{"laranja"}
\NormalTok{cores}
\end{Highlighting}
\end{Shaded}

\textbf{8.} Inquérito de satisfação (fator ordenado).

\begin{Shaded}
\begin{Highlighting}[]
\NormalTok{respostas }\OtherTok{\textless{}{-}} \FunctionTok{c}\NormalTok{(}
  \StringTok{"Satisfeito"}\NormalTok{,}\StringTok{"Muito satisfeito"}\NormalTok{,}\StringTok{"Neutro"}\NormalTok{,}\StringTok{"Insatisfeito"}\NormalTok{,}\StringTok{"Satisfeito"}\NormalTok{,}
  \StringTok{"Muito insatisfeito"}\NormalTok{,}\StringTok{"Satisfeito"}\NormalTok{,}\StringTok{"Neutro"}\NormalTok{,}\StringTok{"Insatisfeito"}\NormalTok{,}\StringTok{"Satisfeito"}\NormalTok{,}
  \StringTok{"Muito satisfeito"}\NormalTok{,}\StringTok{"Neutro"}\NormalTok{,}\StringTok{"Insatisfeito"}\NormalTok{,}\StringTok{"Satisfeito"}\NormalTok{,}\StringTok{"Neutro"}\NormalTok{,}
  \StringTok{"Satisfeito"}\NormalTok{,}\StringTok{"Muito satisfeito"}\NormalTok{,}\StringTok{"Muito insatisfeito"}\NormalTok{,}\StringTok{"Neutro"}\NormalTok{,}\StringTok{"Satisfeito"}\NormalTok{)}

\CommentTok{\# (b) fator ordenado}
\NormalTok{rf }\OtherTok{\textless{}{-}} \FunctionTok{factor}\NormalTok{(respostas,}
  \AttributeTok{levels =} \FunctionTok{c}\NormalTok{(}\StringTok{"Muito insatisfeito"}\NormalTok{,}\StringTok{"Insatisfeito"}\NormalTok{,}\StringTok{"Neutro"}\NormalTok{,}
             \StringTok{"Satisfeito"}\NormalTok{,}\StringTok{"Muito satisfeito"}\NormalTok{),}
  \AttributeTok{ordered =} \ConstantTok{TRUE}\NormalTok{)}
\FunctionTok{str}\NormalTok{(rf)}

\CommentTok{\# (c) número de respostas (Muito) insatisfeito}
\FunctionTok{sum}\NormalTok{(rf }\SpecialCharTok{==} \StringTok{"Insatisfeito"} \SpecialCharTok{|}\NormalTok{ rf }\SpecialCharTok{==} \StringTok{"Muito insatisfeito"}\NormalTok{)}

\CommentTok{\# (d) categoria mais frequente}
\FunctionTok{which.max}\NormalTok{(}\FunctionTok{table}\NormalTok{(rf))}

\CommentTok{\# (e) frequências relativas}
\FunctionTok{table}\NormalTok{(rf) }\SpecialCharTok{/} \FunctionTok{length}\NormalTok{(rf)}

\CommentTok{\# (f) renomear "Neutro" para "Indiferente"}
\FunctionTok{levels}\NormalTok{(rf)[}\FunctionTok{levels}\NormalTok{(rf) }\SpecialCharTok{==} \StringTok{"Neutro"}\NormalTok{] }\OtherTok{\textless{}{-}} \StringTok{"Indiferente"}
\FunctionTok{levels}\NormalTok{(rf)}
\end{Highlighting}
\end{Shaded}

\textbf{9.} Desafios associados a TDAH (fator + gráfico de barras).

\begin{Shaded}
\begin{Highlighting}[]
\NormalTok{tdah }\OtherTok{\textless{}{-}} \FunctionTok{c}\NormalTok{(}\StringTok{"Foco"}\NormalTok{,}\StringTok{"Procrastinação"}\NormalTok{,}\StringTok{"Gestão do tempo"}\NormalTok{,}\StringTok{"Organização"}\NormalTok{,}\StringTok{"Nenhum"}\NormalTok{,}
          \StringTok{"Foco"}\NormalTok{,}\StringTok{"Foco"}\NormalTok{,}\StringTok{"Organização"}\NormalTok{,}\StringTok{"Organização"}\NormalTok{,}\StringTok{"Organização"}\NormalTok{,}
          \StringTok{"Gestão do tempo"}\NormalTok{,}\StringTok{"Gestão do tempo"}\NormalTok{,}\StringTok{"Procrastinação"}\NormalTok{,}
          \StringTok{"Procrastinação"}\NormalTok{,}\StringTok{"Procrastinação"}\NormalTok{,}\StringTok{"Procrastinação"}\NormalTok{,}\StringTok{"Foco"}\NormalTok{,}\StringTok{"Foco"}\NormalTok{,}\StringTok{"Foco"}\NormalTok{)}
\NormalTok{tf }\OtherTok{\textless{}{-}} \FunctionTok{factor}\NormalTok{(tdah)}

\FunctionTok{table}\NormalTok{(tf)                       }\CommentTok{\# (b) frequências absolutas}
\FunctionTok{table}\NormalTok{(tf) }\SpecialCharTok{/} \FunctionTok{length}\NormalTok{(tdah)        }\CommentTok{\#     frequências relativas}

\CommentTok{\# (d) renomear "Nenhum"}
\FunctionTok{levels}\NormalTok{(tf)[}\FunctionTok{levels}\NormalTok{(tf) }\SpecialCharTok{==} \StringTok{"Nenhum"}\NormalTok{] }\OtherTok{\textless{}{-}} \StringTok{"Sem dificuldade declarada"}

\CommentTok{\# (e) gráfico de barras}
\FunctionTok{barplot}\NormalTok{(}\FunctionTok{table}\NormalTok{(tf), }\AttributeTok{ylab =} \StringTok{"Frequência absoluta"}\NormalTok{, }\AttributeTok{main =} \StringTok{"Desafios (TDAH)"}\NormalTok{)}
\end{Highlighting}
\end{Shaded}

\section{Listas}\label{listas-1}

\textbf{1.} Criar uma lista e aceder aos elementos com \texttt{\$} e \texttt{{[}{[}\ {]}{]}}.

\begin{Shaded}
\begin{Highlighting}[]
\NormalTok{dados\_estudante }\OtherTok{\textless{}{-}} \FunctionTok{list}\NormalTok{(}
  \AttributeTok{nome  =} \StringTok{"João"}\NormalTok{,}
  \AttributeTok{idade =} \DecValTok{21}\NormalTok{,}
  \AttributeTok{notas =} \FunctionTok{c}\NormalTok{(}\AttributeTok{Estatistica =} \DecValTok{85}\NormalTok{, }\AttributeTok{Matematica =} \DecValTok{90}\NormalTok{, }\AttributeTok{Computacao =} \DecValTok{95}\NormalTok{))}

\NormalTok{dados\_estudante}\SpecialCharTok{$}\NormalTok{nome}
\NormalTok{dados\_estudante}\SpecialCharTok{$}\NormalTok{idade}
\NormalTok{dados\_estudante}\SpecialCharTok{$}\NormalTok{notas[}\StringTok{"Computacao"}\NormalTok{]}
\end{Highlighting}
\end{Shaded}

\textbf{2.} Adicionar e modificar elementos.

\begin{Shaded}
\begin{Highlighting}[]
\NormalTok{dados\_estudante}\SpecialCharTok{$}\NormalTok{status }\OtherTok{\textless{}{-}} \StringTok{"Aprovado"}
\NormalTok{dados\_estudante}\SpecialCharTok{$}\NormalTok{notas[}\StringTok{"Estatistica"}\NormalTok{] }\OtherTok{\textless{}{-}} \DecValTok{88}
\end{Highlighting}
\end{Shaded}

\textbf{3. e 4.} Listas de listas (uma turma de estudantes).

\begin{Shaded}
\begin{Highlighting}[]
\NormalTok{estudante1 }\OtherTok{\textless{}{-}} \FunctionTok{list}\NormalTok{(}\AttributeTok{nome =} \StringTok{"Maria"}\NormalTok{,  }\AttributeTok{idade =} \DecValTok{22}\NormalTok{, }\AttributeTok{notas =} \FunctionTok{c}\NormalTok{(}\DecValTok{78}\NormalTok{,}\DecValTok{85}\NormalTok{,}\DecValTok{90}\NormalTok{), }\AttributeTok{status =} \StringTok{"Aprovado"}\NormalTok{)}
\NormalTok{estudante2 }\OtherTok{\textless{}{-}} \FunctionTok{list}\NormalTok{(}\AttributeTok{nome =} \StringTok{"Carlos"}\NormalTok{, }\AttributeTok{idade =} \DecValTok{23}\NormalTok{, }\AttributeTok{notas =} \FunctionTok{c}\NormalTok{(}\DecValTok{70}\NormalTok{,}\DecValTok{75}\NormalTok{,}\DecValTok{80}\NormalTok{), }\AttributeTok{status =} \StringTok{"Aprovado"}\NormalTok{)}
\NormalTok{turma }\OtherTok{\textless{}{-}} \FunctionTok{list}\NormalTok{(estudante1, estudante2)}

\NormalTok{turma[[}\DecValTok{2}\NormalTok{]]}\SpecialCharTok{$}\NormalTok{nome            }\CommentTok{\# (4a) nome do segundo estudante}
\FunctionTok{mean}\NormalTok{(turma[[}\DecValTok{1}\NormalTok{]]}\SpecialCharTok{$}\NormalTok{notas)     }\CommentTok{\# (4b) média das notas do primeiro}
\NormalTok{turma[[}\DecValTok{2}\NormalTok{]]}\SpecialCharTok{$}\NormalTok{status }\OtherTok{\textless{}{-}} \StringTok{"Reprovado"}  \CommentTok{\# (4c) alterar o estado}
\end{Highlighting}
\end{Shaded}

Repare na diferença entre \texttt{{[}\ {]}} e \texttt{{[}{[}\ {]}{]}} numa lista: \texttt{turma{[}2{]}} devolve uma
\textbf{sublista} (ainda uma lista, de comprimento 1), enquanto \texttt{turma{[}{[}2{]}{]}} devolve
o \textbf{conteúdo} do segundo elemento (a lista do estudante). Para aceder aos
campos com \texttt{\$}, precisa do conteúdo, logo de \texttt{{[}{[}\ {]}{]}}.

\textbf{5.} Listas encaixadas (várias turmas).

\begin{Shaded}
\begin{Highlighting}[]
\NormalTok{estudante3 }\OtherTok{\textless{}{-}} \FunctionTok{list}\NormalTok{(}\AttributeTok{nome =} \StringTok{"Ana"}\NormalTok{,  }\AttributeTok{idade =} \DecValTok{19}\NormalTok{, }\AttributeTok{notas =} \FunctionTok{c}\NormalTok{(}\DecValTok{70}\NormalTok{,}\DecValTok{80}\NormalTok{,}\DecValTok{96}\NormalTok{), }\AttributeTok{status =} \StringTok{"Aprovado"}\NormalTok{)}
\NormalTok{estudante4 }\OtherTok{\textless{}{-}} \FunctionTok{list}\NormalTok{(}\AttributeTok{nome =} \StringTok{"João"}\NormalTok{, }\AttributeTok{idade =} \DecValTok{23}\NormalTok{, }\AttributeTok{notas =} \FunctionTok{c}\NormalTok{(}\DecValTok{85}\NormalTok{,}\DecValTok{80}\NormalTok{,}\DecValTok{95}\NormalTok{), }\AttributeTok{status =} \StringTok{"Aprovado"}\NormalTok{)}
\NormalTok{turma1 }\OtherTok{\textless{}{-}} \FunctionTok{list}\NormalTok{(estudante1, estudante2)}
\NormalTok{turma2 }\OtherTok{\textless{}{-}} \FunctionTok{list}\NormalTok{(estudante3, estudante4)}
\NormalTok{cadeira }\OtherTok{\textless{}{-}} \FunctionTok{list}\NormalTok{(turma1, turma2)}

\NormalTok{cadeira[[}\DecValTok{2}\NormalTok{]][[}\DecValTok{1}\NormalTok{]]}\SpecialCharTok{$}\NormalTok{nome          }\CommentTok{\# (a) nome do 1.º estudante da 2.ª turma}
\FunctionTok{mean}\NormalTok{(cadeira[[}\DecValTok{1}\NormalTok{]][[}\DecValTok{2}\NormalTok{]]}\SpecialCharTok{$}\NormalTok{notas)   }\CommentTok{\# (b) média do 2.º estudante da 1.ª turma}
\end{Highlighting}
\end{Shaded}

\textbf{6.} Produtos de uma fábrica (aplicação com listas paralelas).

\begin{Shaded}
\begin{Highlighting}[]
\NormalTok{produtos }\OtherTok{\textless{}{-}} \FunctionTok{list}\NormalTok{(}
  \AttributeTok{nome       =} \FunctionTok{c}\NormalTok{(}\StringTok{"Telemóvel"}\NormalTok{, }\StringTok{"PC"}\NormalTok{, }\StringTok{"Rato"}\NormalTok{, }\StringTok{"Fone"}\NormalTok{),}
  \AttributeTok{quantidade =} \FunctionTok{c}\NormalTok{(}\DecValTok{1000}\NormalTok{, }\DecValTok{1500}\NormalTok{, }\DecValTok{2000}\NormalTok{, }\DecValTok{2500}\NormalTok{),}
  \AttributeTok{custo      =} \FunctionTok{c}\NormalTok{(}\DecValTok{266}\NormalTok{, }\DecValTok{283}\NormalTok{, }\DecValTok{16}\NormalTok{, }\DecValTok{26}\NormalTok{),}
  \AttributeTok{venda      =} \FunctionTok{c}\NormalTok{(}\DecValTok{800}\NormalTok{, }\DecValTok{850}\NormalTok{, }\DecValTok{50}\NormalTok{, }\DecValTok{80}\NormalTok{))}

\CommentTok{\# margem de lucro mensal por produto}
\NormalTok{produtos}\SpecialCharTok{$}\NormalTok{margem }\OtherTok{\textless{}{-}}\NormalTok{ (produtos}\SpecialCharTok{$}\NormalTok{venda }\SpecialCharTok{{-}}\NormalTok{ produtos}\SpecialCharTok{$}\NormalTok{custo) }\SpecialCharTok{*}\NormalTok{ produtos}\SpecialCharTok{$}\NormalTok{quantidade}

\CommentTok{\# produto com maior lucro}
\NormalTok{produtos}\SpecialCharTok{$}\NormalTok{nome[}\FunctionTok{which.max}\NormalTok{(produtos}\SpecialCharTok{$}\NormalTok{margem)]}

\CommentTok{\# aumentar o preço de venda do 2.º produto em 5\% e recalcular}
\NormalTok{produtos}\SpecialCharTok{$}\NormalTok{venda[}\DecValTok{2}\NormalTok{] }\OtherTok{\textless{}{-}}\NormalTok{ produtos}\SpecialCharTok{$}\NormalTok{venda[}\DecValTok{2}\NormalTok{] }\SpecialCharTok{*} \FloatTok{1.05}
\NormalTok{produtos}\SpecialCharTok{$}\NormalTok{margem[}\DecValTok{2}\NormalTok{] }\OtherTok{\textless{}{-}}\NormalTok{ (produtos}\SpecialCharTok{$}\NormalTok{venda[}\DecValTok{2}\NormalTok{] }\SpecialCharTok{{-}}\NormalTok{ produtos}\SpecialCharTok{$}\NormalTok{custo[}\DecValTok{2}\NormalTok{]) }\SpecialCharTok{*}\NormalTok{ produtos}\SpecialCharTok{$}\NormalTok{quantidade[}\DecValTok{2}\NormalTok{]}
\end{Highlighting}
\end{Shaded}

\textbf{7.} Registos de utilização (listas encaixadas com vetores).

\begin{Shaded}
\begin{Highlighting}[]
\NormalTok{logs }\OtherTok{\textless{}{-}} \FunctionTok{list}\NormalTok{(}
  \AttributeTok{e1 =} \FunctionTok{list}\NormalTok{(}\AttributeTok{nome =} \StringTok{"Ana"}\NormalTok{,   }\AttributeTok{mensagens =} \FunctionTok{c}\NormalTok{(}\DecValTok{12}\NormalTok{, }\DecValTok{15}\NormalTok{, }\DecValTok{9}\NormalTok{),    }\AttributeTok{tempo\_total =} \DecValTok{35}\NormalTok{),}
  \AttributeTok{e2 =} \FunctionTok{list}\NormalTok{(}\AttributeTok{nome =} \StringTok{"Bruno"}\NormalTok{, }\AttributeTok{mensagens =} \FunctionTok{c}\NormalTok{(}\DecValTok{7}\NormalTok{, }\DecValTok{6}\NormalTok{, }\DecValTok{5}\NormalTok{, }\DecValTok{12}\NormalTok{),  }\AttributeTok{tempo\_total =} \DecValTok{42}\NormalTok{),}
  \AttributeTok{e3 =} \FunctionTok{list}\NormalTok{(}\AttributeTok{nome =} \StringTok{"Carla"}\NormalTok{, }\AttributeTok{mensagens =} \FunctionTok{c}\NormalTok{(}\DecValTok{10}\NormalTok{, }\DecValTok{20}\NormalTok{),        }\AttributeTok{tempo\_total =} \DecValTok{28}\NormalTok{))}

\NormalTok{logs}\SpecialCharTok{$}\NormalTok{e2}\SpecialCharTok{$}\NormalTok{mensagens[}\DecValTok{3}\NormalTok{]                       }\CommentTok{\# 3.ª mensagem do Bruno}
\FunctionTok{sapply}\NormalTok{(logs, }\ControlFlowTok{function}\NormalTok{(e) }\FunctionTok{sum}\NormalTok{(e}\SpecialCharTok{$}\NormalTok{mensagens)) }\CommentTok{\# total de mensagens por pessoa}
\FunctionTok{sapply}\NormalTok{(logs, }\ControlFlowTok{function}\NormalTok{(e) e}\SpecialCharTok{$}\NormalTok{tempo\_total }\SpecialCharTok{/} \FunctionTok{length}\NormalTok{(e}\SpecialCharTok{$}\NormalTok{mensagens))  }\CommentTok{\# tempo médio/sessão}
\end{Highlighting}
\end{Shaded}

Sempre que precisar de aplicar o mesmo cálculo a todos os elementos de uma
lista, prefira \texttt{sapply()}/\texttt{lapply()} a repetir o código manualmente para cada
elemento. Além de mais curto, o código fica imune a erros de cópia e adapta-se
sozinho se a lista crescer.

\section{\texorpdfstring{Matrizes e \emph{arrays}}{Matrizes e arrays}}\label{matrizes-e-arrays}

\textbf{1.} Criar uma matriz e somar uma coluna (três formas equivalentes).

\begin{Shaded}
\begin{Highlighting}[]
\NormalTok{A }\OtherTok{\textless{}{-}} \FunctionTok{matrix}\NormalTok{(}\DecValTok{1}\SpecialCharTok{:}\DecValTok{12}\NormalTok{, }\AttributeTok{nrow =} \DecValTok{3}\NormalTok{, }\AttributeTok{ncol =} \DecValTok{4}\NormalTok{)}
\FunctionTok{sum}\NormalTok{(A[, }\DecValTok{2}\NormalTok{])        }\CommentTok{\# forma direta}
\FunctionTok{colSums}\NormalTok{(A)[}\DecValTok{2}\NormalTok{]      }\CommentTok{\# Alternativa}
\FunctionTok{apply}\NormalTok{(A, }\DecValTok{2}\NormalTok{, sum)[}\DecValTok{2}\NormalTok{] }\CommentTok{\# Alternativa}
\end{Highlighting}
\end{Shaded}

\textbf{2.} Operações elemento a elemento.

\begin{Shaded}
\begin{Highlighting}[]
\NormalTok{A }\OtherTok{\textless{}{-}} \FunctionTok{matrix}\NormalTok{(}\FunctionTok{sample}\NormalTok{(}\DecValTok{1}\SpecialCharTok{:}\DecValTok{10}\NormalTok{, }\DecValTok{6}\NormalTok{, }\AttributeTok{replace =} \ConstantTok{TRUE}\NormalTok{), }\AttributeTok{nrow =} \DecValTok{2}\NormalTok{)}
\NormalTok{B }\OtherTok{\textless{}{-}} \FunctionTok{matrix}\NormalTok{(}\FunctionTok{sample}\NormalTok{(}\DecValTok{1}\SpecialCharTok{:}\DecValTok{10}\NormalTok{, }\DecValTok{6}\NormalTok{, }\AttributeTok{replace =} \ConstantTok{TRUE}\NormalTok{), }\AttributeTok{nrow =} \DecValTok{2}\NormalTok{)}
\NormalTok{A }\SpecialCharTok{+}\NormalTok{ B    }\CommentTok{\# soma}
\NormalTok{A }\SpecialCharTok{*}\NormalTok{ B    }\CommentTok{\# produto de Hadamard (elemento a elemento)}
\end{Highlighting}
\end{Shaded}

\textbf{3.} Produto matriz-vetor.

\begin{Shaded}
\begin{Highlighting}[]
\NormalTok{M }\OtherTok{\textless{}{-}} \FunctionTok{matrix}\NormalTok{(}\FunctionTok{sample}\NormalTok{(}\DecValTok{1}\SpecialCharTok{:}\DecValTok{9}\NormalTok{, }\DecValTok{9}\NormalTok{, }\AttributeTok{replace =} \ConstantTok{TRUE}\NormalTok{), }\AttributeTok{nrow =} \DecValTok{3}\NormalTok{)}
\NormalTok{v }\OtherTok{\textless{}{-}} \FunctionTok{c}\NormalTok{(}\DecValTok{1}\NormalTok{, }\DecValTok{2}\NormalTok{, }\DecValTok{3}\NormalTok{)}
\NormalTok{M }\SpecialCharTok{\%*\%}\NormalTok{ v   }\CommentTok{\# M v (coluna)}
\NormalTok{v }\SpecialCharTok{\%*\%}\NormalTok{ M   }\CommentTok{\# v\^{}T M (linha)}
\end{Highlighting}
\end{Shaded}

\texttt{*} é o produto \textbf{elemento a elemento}; o produto matricial faz-se com \texttt{\%*\%}.
Confundi-los é um dos erros mais frequentes em álgebra linear com R.

\textbf{4.} Transposta e somas por linha.

\begin{Shaded}
\begin{Highlighting}[]
\NormalTok{N }\OtherTok{\textless{}{-}} \FunctionTok{matrix}\NormalTok{(}\DecValTok{1}\SpecialCharTok{:}\DecValTok{8}\NormalTok{, }\AttributeTok{nrow =} \DecValTok{4}\NormalTok{)}
\NormalTok{Nt }\OtherTok{\textless{}{-}} \FunctionTok{t}\NormalTok{(N)}
\FunctionTok{rowSums}\NormalTok{(Nt)}
\end{Highlighting}
\end{Shaded}

\textbf{5.} Determinante e inversa.

\begin{Shaded}
\begin{Highlighting}[]
\NormalTok{P }\OtherTok{\textless{}{-}} \FunctionTok{matrix}\NormalTok{(}\FunctionTok{sample}\NormalTok{(}\DecValTok{1}\SpecialCharTok{:}\DecValTok{9}\NormalTok{, }\DecValTok{9}\NormalTok{, }\AttributeTok{replace =} \ConstantTok{TRUE}\NormalTok{), }\AttributeTok{nrow =} \DecValTok{3}\NormalTok{)}
\FunctionTok{det}\NormalTok{(P)}
\ControlFlowTok{if}\NormalTok{ (}\FunctionTok{det}\NormalTok{(P) }\SpecialCharTok{!=} \DecValTok{0}\NormalTok{) }\FunctionTok{solve}\NormalTok{(P)   }\CommentTok{\# a inversa só existe se det != 0}
\end{Highlighting}
\end{Shaded}

\textbf{6.} Extração de submatrizes.

\begin{Shaded}
\begin{Highlighting}[]
\NormalTok{Q }\OtherTok{\textless{}{-}} \FunctionTok{matrix}\NormalTok{(}\FunctionTok{sample}\NormalTok{(}\DecValTok{1}\SpecialCharTok{:}\DecValTok{25}\NormalTok{, }\DecValTok{25}\NormalTok{, }\AttributeTok{replace =} \ConstantTok{TRUE}\NormalTok{), }\AttributeTok{nrow =} \DecValTok{5}\NormalTok{)}
\NormalTok{Q[}\DecValTok{2}\SpecialCharTok{:}\DecValTok{4}\NormalTok{, }\DecValTok{2}\SpecialCharTok{:}\DecValTok{4}\NormalTok{]}
\end{Highlighting}
\end{Shaded}

\textbf{7.} Somas de linha e de coluna.

\begin{Shaded}
\begin{Highlighting}[]
\NormalTok{A }\OtherTok{\textless{}{-}} \FunctionTok{matrix}\NormalTok{(}\DecValTok{1}\SpecialCharTok{:}\DecValTok{9}\NormalTok{, }\AttributeTok{nrow =} \DecValTok{3}\NormalTok{)}
\FunctionTok{sum}\NormalTok{(A[}\DecValTok{1}\NormalTok{, ])   }\CommentTok{\# 1.ª linha}
\FunctionTok{sum}\NormalTok{(A[, }\DecValTok{3}\NormalTok{])   }\CommentTok{\# 3.ª coluna}
\end{Highlighting}
\end{Shaded}

\textbf{8.} Potências de matrizes e manipulação de colunas.

\begin{Shaded}
\begin{Highlighting}[]
\FunctionTok{library}\NormalTok{(expm)   }\CommentTok{\# necessário para o operador \%\^{}\%}
\NormalTok{A }\OtherTok{\textless{}{-}} \FunctionTok{rbind}\NormalTok{(}\FunctionTok{c}\NormalTok{(}\DecValTok{1}\NormalTok{, }\DecValTok{1}\NormalTok{, }\DecValTok{3}\NormalTok{), }\FunctionTok{c}\NormalTok{(}\DecValTok{5}\NormalTok{, }\DecValTok{2}\NormalTok{, }\DecValTok{6}\NormalTok{), }\FunctionTok{c}\NormalTok{(}\SpecialCharTok{{-}}\DecValTok{2}\NormalTok{, }\SpecialCharTok{{-}}\DecValTok{1}\NormalTok{, }\SpecialCharTok{{-}}\DecValTok{3}\NormalTok{))}
\NormalTok{A }\SpecialCharTok{\%*\%}\NormalTok{ A }\SpecialCharTok{\%*\%}\NormalTok{ A   }\CommentTok{\# A\^{}3}
\NormalTok{A }\SpecialCharTok{\%\^{}\%} \DecValTok{3}         \CommentTok{\# o mesmo, com expm}

\NormalTok{A[, }\DecValTok{3}\NormalTok{] }\OtherTok{\textless{}{-}}\NormalTok{ A[, }\DecValTok{1}\NormalTok{] }\SpecialCharTok{+}\NormalTok{ A[, }\DecValTok{2}\NormalTok{]   }\CommentTok{\# substitui a 3.ª coluna}

\NormalTok{M }\OtherTok{\textless{}{-}} \FunctionTok{cbind}\NormalTok{(}\FunctionTok{c}\NormalTok{(}\DecValTok{20}\NormalTok{,}\DecValTok{34}\NormalTok{,}\DecValTok{99}\NormalTok{,}\DecValTok{12}\NormalTok{,}\DecValTok{10}\NormalTok{,}\DecValTok{14}\NormalTok{), }\FunctionTok{c}\NormalTok{(}\DecValTok{22}\NormalTok{,}\DecValTok{55}\NormalTok{,}\DecValTok{97}\NormalTok{,}\DecValTok{16}\NormalTok{,}\DecValTok{53}\NormalTok{,}\DecValTok{21}\NormalTok{), }\FunctionTok{c}\NormalTok{(}\DecValTok{23}\NormalTok{,}\DecValTok{57}\NormalTok{,}\DecValTok{71}\NormalTok{,}\DecValTok{19}\NormalTok{,}\DecValTok{24}\NormalTok{,}\DecValTok{28}\NormalTok{))}
\NormalTok{M[M }\SpecialCharTok{\%\%} \DecValTok{2} \SpecialCharTok{==} \DecValTok{0}\NormalTok{] }\OtherTok{\textless{}{-}} \DecValTok{0}         \CommentTok{\# anula as entradas pares}
\end{Highlighting}
\end{Shaded}

\textbf{9.} Verificar simetria.

\begin{Shaded}
\begin{Highlighting}[]
\NormalTok{M }\OtherTok{\textless{}{-}} \FunctionTok{matrix}\NormalTok{(}\FunctionTok{c}\NormalTok{(}\DecValTok{1}\NormalTok{,}\DecValTok{2}\NormalTok{,}\DecValTok{3}\NormalTok{, }\DecValTok{2}\NormalTok{,}\DecValTok{4}\NormalTok{,}\DecValTok{5}\NormalTok{, }\DecValTok{3}\NormalTok{,}\DecValTok{5}\NormalTok{,}\DecValTok{6}\NormalTok{), }\AttributeTok{nrow =} \DecValTok{3}\NormalTok{)}
\FunctionTok{all}\NormalTok{(M }\SpecialCharTok{==} \FunctionTok{t}\NormalTok{(M))     }\CommentTok{\# TRUE se simétrica}
\end{Highlighting}
\end{Shaded}

\textbf{10.} Traço de uma matriz.

\begin{Shaded}
\begin{Highlighting}[]
\NormalTok{A }\OtherTok{\textless{}{-}} \FunctionTok{matrix}\NormalTok{(}\DecValTok{1}\SpecialCharTok{:}\DecValTok{16}\NormalTok{, }\AttributeTok{nrow =} \DecValTok{4}\NormalTok{)}
\FunctionTok{sum}\NormalTok{(}\FunctionTok{diag}\NormalTok{(A))}
\end{Highlighting}
\end{Shaded}

\textbf{11.} Selecionar colunas pares.

\begin{Shaded}
\begin{Highlighting}[]
\NormalTok{M }\OtherTok{\textless{}{-}} \FunctionTok{matrix}\NormalTok{(}\DecValTok{1}\SpecialCharTok{:}\DecValTok{25}\NormalTok{, }\AttributeTok{nrow =} \DecValTok{5}\NormalTok{, }\AttributeTok{byrow =} \ConstantTok{TRUE}\NormalTok{)}
\NormalTok{M[, }\FunctionTok{seq}\NormalTok{(}\DecValTok{2}\NormalTok{, }\DecValTok{5}\NormalTok{, }\AttributeTok{by =} \DecValTok{2}\NormalTok{)]}
\end{Highlighting}
\end{Shaded}

\textbf{12.} Desempenho de alunos (matriz com nomes).

\begin{Shaded}
\begin{Highlighting}[]
\NormalTok{M }\OtherTok{\textless{}{-}} \FunctionTok{matrix}\NormalTok{(}\FunctionTok{c}\NormalTok{(}\DecValTok{12}\NormalTok{,}\DecValTok{8}\NormalTok{,}\DecValTok{5}\NormalTok{, }\DecValTok{20}\NormalTok{,}\DecValTok{10}\NormalTok{,}\DecValTok{4}\NormalTok{, }\DecValTok{15}\NormalTok{,}\DecValTok{9}\NormalTok{,}\DecValTok{5}\NormalTok{), }\AttributeTok{nrow =} \DecValTok{3}\NormalTok{, }\AttributeTok{byrow =} \ConstantTok{TRUE}\NormalTok{)}
\FunctionTok{colnames}\NormalTok{(M) }\OtherTok{\textless{}{-}} \FunctionTok{c}\NormalTok{(}\StringTok{"Perguntas"}\NormalTok{, }\StringTok{"TempoResposta"}\NormalTok{, }\StringTok{"Satisfacao"}\NormalTok{)}
\FunctionTok{rownames}\NormalTok{(M) }\OtherTok{\textless{}{-}} \FunctionTok{c}\NormalTok{(}\StringTok{"AlunoA"}\NormalTok{, }\StringTok{"AlunoB"}\NormalTok{, }\StringTok{"AlunoC"}\NormalTok{)}

\NormalTok{M }\OtherTok{\textless{}{-}} \FunctionTok{rbind}\NormalTok{(M, }\AttributeTok{AlunoD =} \FunctionTok{c}\NormalTok{(}\DecValTok{18}\NormalTok{, }\DecValTok{11}\NormalTok{, }\DecValTok{3}\NormalTok{))     }\CommentTok{\# (b)}
\NormalTok{M }\OtherTok{\textless{}{-}} \FunctionTok{cbind}\NormalTok{(M, }\AttributeTok{HorasUso =} \FunctionTok{c}\NormalTok{(}\DecValTok{5}\NormalTok{, }\DecValTok{7}\NormalTok{, }\DecValTok{6}\NormalTok{, }\DecValTok{8}\NormalTok{))  }\CommentTok{\# (c)}
\FunctionTok{colMeans}\NormalTok{(M)                              }\CommentTok{\# (d)}
\FunctionTok{rownames}\NormalTok{(M)[}\FunctionTok{which.max}\NormalTok{(M[, }\StringTok{"Perguntas"}\NormalTok{])] }\CommentTok{\# (e) quem fez mais perguntas}
\end{Highlighting}
\end{Shaded}

\textbf{13.} Imagem representada por uma matriz.

\begin{Shaded}
\begin{Highlighting}[]
\NormalTok{I }\OtherTok{\textless{}{-}} \FunctionTok{matrix}\NormalTok{(}\FunctionTok{c}\NormalTok{(}
  \DecValTok{0}\NormalTok{,}\DecValTok{0}\NormalTok{,}\DecValTok{0}\NormalTok{,}\DecValTok{0}\NormalTok{,}\DecValTok{0}\NormalTok{,}\DecValTok{0}\NormalTok{,}\DecValTok{0}\NormalTok{,}\DecValTok{0}\NormalTok{,}\DecValTok{0}\NormalTok{,}\DecValTok{0}\NormalTok{,  }\DecValTok{0}\NormalTok{,}\DecValTok{9}\NormalTok{,}\DecValTok{9}\NormalTok{,}\DecValTok{9}\NormalTok{,}\DecValTok{0}\NormalTok{,}\DecValTok{0}\NormalTok{,}\DecValTok{9}\NormalTok{,}\DecValTok{9}\NormalTok{,}\DecValTok{9}\NormalTok{,}\DecValTok{0}\NormalTok{,  }\DecValTok{0}\NormalTok{,}\DecValTok{9}\NormalTok{,}\DecValTok{0}\NormalTok{,}\DecValTok{9}\NormalTok{,}\DecValTok{0}\NormalTok{,}\DecValTok{0}\NormalTok{,}\DecValTok{9}\NormalTok{,}\DecValTok{0}\NormalTok{,}\DecValTok{9}\NormalTok{,}\DecValTok{0}\NormalTok{,}
  \DecValTok{0}\NormalTok{,}\DecValTok{9}\NormalTok{,}\DecValTok{9}\NormalTok{,}\DecValTok{9}\NormalTok{,}\DecValTok{0}\NormalTok{,}\DecValTok{0}\NormalTok{,}\DecValTok{9}\NormalTok{,}\DecValTok{9}\NormalTok{,}\DecValTok{9}\NormalTok{,}\DecValTok{0}\NormalTok{,  }\DecValTok{0}\NormalTok{,}\DecValTok{9}\NormalTok{,}\DecValTok{0}\NormalTok{,}\DecValTok{0}\NormalTok{,}\DecValTok{0}\NormalTok{,}\DecValTok{0}\NormalTok{,}\DecValTok{0}\NormalTok{,}\DecValTok{0}\NormalTok{,}\DecValTok{9}\NormalTok{,}\DecValTok{0}\NormalTok{,  }\DecValTok{0}\NormalTok{,}\DecValTok{9}\NormalTok{,}\DecValTok{0}\NormalTok{,}\DecValTok{0}\NormalTok{,}\DecValTok{0}\NormalTok{,}\DecValTok{0}\NormalTok{,}\DecValTok{0}\NormalTok{,}\DecValTok{0}\NormalTok{,}\DecValTok{9}\NormalTok{,}\DecValTok{0}\NormalTok{,}
  \DecValTok{0}\NormalTok{,}\DecValTok{9}\NormalTok{,}\DecValTok{0}\NormalTok{,}\DecValTok{0}\NormalTok{,}\DecValTok{0}\NormalTok{,}\DecValTok{0}\NormalTok{,}\DecValTok{0}\NormalTok{,}\DecValTok{0}\NormalTok{,}\DecValTok{9}\NormalTok{,}\DecValTok{0}\NormalTok{,  }\DecValTok{0}\NormalTok{,}\DecValTok{9}\NormalTok{,}\DecValTok{0}\NormalTok{,}\DecValTok{0}\NormalTok{,}\DecValTok{0}\NormalTok{,}\DecValTok{0}\NormalTok{,}\DecValTok{0}\NormalTok{,}\DecValTok{0}\NormalTok{,}\DecValTok{9}\NormalTok{,}\DecValTok{0}\NormalTok{,  }\DecValTok{0}\NormalTok{,}\DecValTok{9}\NormalTok{,}\DecValTok{9}\NormalTok{,}\DecValTok{9}\NormalTok{,}\DecValTok{9}\NormalTok{,}\DecValTok{9}\NormalTok{,}\DecValTok{9}\NormalTok{,}\DecValTok{9}\NormalTok{,}\DecValTok{9}\NormalTok{,}\DecValTok{0}\NormalTok{,}
  \DecValTok{0}\NormalTok{,}\DecValTok{0}\NormalTok{,}\DecValTok{0}\NormalTok{,}\DecValTok{0}\NormalTok{,}\DecValTok{0}\NormalTok{,}\DecValTok{0}\NormalTok{,}\DecValTok{0}\NormalTok{,}\DecValTok{0}\NormalTok{,}\DecValTok{0}\NormalTok{,}\DecValTok{0}\NormalTok{), }\AttributeTok{nrow =} \DecValTok{10}\NormalTok{, }\AttributeTok{byrow =} \ConstantTok{TRUE}\NormalTok{)}

\FunctionTok{image}\NormalTok{(I, }\AttributeTok{col =} \FunctionTok{gray.colors}\NormalTok{(}\DecValTok{10}\NormalTok{))          }\CommentTok{\# (b) visualizar}
\FunctionTok{image}\NormalTok{(}\DecValTok{9} \SpecialCharTok{{-}}\NormalTok{ I, }\AttributeTok{col =} \FunctionTok{gray.colors}\NormalTok{(}\DecValTok{10}\NormalTok{))       }\CommentTok{\# (c) inverter cores}
\FunctionTok{mean}\NormalTok{(I)                                   }\CommentTok{\# (d) brilho médio}
\NormalTok{I }\OtherTok{\textless{}{-}} \FunctionTok{rbind}\NormalTok{(I, }\FunctionTok{rep}\NormalTok{(}\DecValTok{0}\NormalTok{, }\FunctionTok{ncol}\NormalTok{(I))); }\FunctionTok{dim}\NormalTok{(I)    }\CommentTok{\# (e) nova linha de zeros}
\end{Highlighting}
\end{Shaded}

\textbf{14.} \emph{Arrays} a três dimensões.

\begin{Shaded}
\begin{Highlighting}[]
\NormalTok{A }\OtherTok{\textless{}{-}} \FunctionTok{array}\NormalTok{(}\DecValTok{1}\SpecialCharTok{:}\DecValTok{27}\NormalTok{, }\AttributeTok{dim =} \FunctionTok{c}\NormalTok{(}\DecValTok{3}\NormalTok{, }\DecValTok{3}\NormalTok{, }\DecValTok{3}\NormalTok{))}
\NormalTok{A[}\DecValTok{2}\NormalTok{, }\DecValTok{3}\NormalTok{, }\DecValTok{1}\NormalTok{]     }\CommentTok{\# elemento}
\NormalTok{A[, , }\DecValTok{2}\NormalTok{]       }\CommentTok{\# segunda "camada"}
\NormalTok{A[}\DecValTok{3}\NormalTok{, , }\DecValTok{2}\NormalTok{]      }\CommentTok{\# 3.ª linha da 2.ª camada}
\end{Highlighting}
\end{Shaded}

No material original, este último acesso aparecia como \texttt{arr{[}3,\ ,\ 2{]}}, mas o
\emph{array} chama-se \texttt{A} (ou \texttt{array\_3d}). Referir um objeto inexistente devolve o
erro \texttt{object\ \textquotesingle{}arr\textquotesingle{}\ not\ found}. Verifique sempre que o nome usado é o do objeto
efetivamente criado.

\textbf{15.} Operações com dois \emph{arrays}.

\begin{Shaded}
\begin{Highlighting}[]
\FunctionTok{set.seed}\NormalTok{(}\DecValTok{123}\NormalTok{)}
\NormalTok{a1 }\OtherTok{\textless{}{-}} \FunctionTok{array}\NormalTok{(}\FunctionTok{sample}\NormalTok{(}\DecValTok{1}\SpecialCharTok{:}\DecValTok{10}\NormalTok{, }\DecValTok{8}\NormalTok{, }\AttributeTok{replace =} \ConstantTok{TRUE}\NormalTok{), }\AttributeTok{dim =} \FunctionTok{c}\NormalTok{(}\DecValTok{2}\NormalTok{, }\DecValTok{2}\NormalTok{, }\DecValTok{2}\NormalTok{))}
\NormalTok{a2 }\OtherTok{\textless{}{-}} \FunctionTok{array}\NormalTok{(}\FunctionTok{sample}\NormalTok{(}\DecValTok{1}\SpecialCharTok{:}\DecValTok{10}\NormalTok{, }\DecValTok{8}\NormalTok{, }\AttributeTok{replace =} \ConstantTok{TRUE}\NormalTok{), }\AttributeTok{dim =} \FunctionTok{c}\NormalTok{(}\DecValTok{2}\NormalTok{, }\DecValTok{2}\NormalTok{, }\DecValTok{2}\NormalTok{))}
\NormalTok{a1 }\SpecialCharTok{+}\NormalTok{ a2                     }\CommentTok{\# soma}
\NormalTok{a1 }\SpecialCharTok{*}\NormalTok{ a2                     }\CommentTok{\# produto elemento a elemento}
\FunctionTok{mean}\NormalTok{(a1 }\SpecialCharTok{+}\NormalTok{ a2)               }\CommentTok{\# média da soma}
\NormalTok{a1[, , }\DecValTok{1}\NormalTok{] }\SpecialCharTok{\%*\%}\NormalTok{ a2[, , }\DecValTok{2}\NormalTok{]     }\CommentTok{\# produto matricial de duas camadas}
\end{Highlighting}
\end{Shaded}

\textbf{16.} Cadeia de Markov (iteração de uma matriz de transição).

\begin{Shaded}
\begin{Highlighting}[]
\FunctionTok{library}\NormalTok{(expm)}
\NormalTok{M }\OtherTok{\textless{}{-}} \FunctionTok{matrix}\NormalTok{(}\FunctionTok{c}\NormalTok{(}\FloatTok{0.80}\NormalTok{, }\FloatTok{0.20}\NormalTok{, }\FloatTok{0.10}\NormalTok{, }\FloatTok{0.90}\NormalTok{), }\AttributeTok{nrow =} \DecValTok{2}\NormalTok{)  }\CommentTok{\# colunas somam 1}
\NormalTok{v }\OtherTok{\textless{}{-}} \FunctionTok{c}\NormalTok{(}\DecValTok{150}\NormalTok{, }\DecValTok{150}\NormalTok{)}
\NormalTok{M }\SpecialCharTok{\%*\%}\NormalTok{ v            }\CommentTok{\# 1.ª iteração}
\NormalTok{(M }\SpecialCharTok{\%\^{}\%} \DecValTok{10}\NormalTok{) }\SpecialCharTok{\%*\%}\NormalTok{ v   }\CommentTok{\# 10.ª iteração}
\end{Highlighting}
\end{Shaded}

\section{\texorpdfstring{\emph{Data frames}}{Data frames}}\label{data-frames}

\textbf{1.} Criar, acrescentar linha e ordenar.

\begin{Shaded}
\begin{Highlighting}[]
\NormalTok{estudantes }\OtherTok{\textless{}{-}} \FunctionTok{data.frame}\NormalTok{(}
  \AttributeTok{nome  =} \FunctionTok{c}\NormalTok{(}\StringTok{"Ana"}\NormalTok{, }\StringTok{"João"}\NormalTok{, }\StringTok{"Maria"}\NormalTok{, }\StringTok{"Inês"}\NormalTok{, }\StringTok{"Gonçalo"}\NormalTok{),}
  \AttributeTok{idade =} \FunctionTok{c}\NormalTok{(}\DecValTok{18}\NormalTok{, }\DecValTok{19}\NormalTok{, }\DecValTok{20}\NormalTok{, }\DecValTok{18}\NormalTok{, }\DecValTok{20}\NormalTok{),}
  \AttributeTok{curso =} \FunctionTok{c}\NormalTok{(}\StringTok{"Matemática"}\NormalTok{,}\StringTok{"Psicologia"}\NormalTok{,}\StringTok{"Geologia"}\NormalTok{,}\StringTok{"Gestão"}\NormalTok{,}\StringTok{"Física"}\NormalTok{),}
  \AttributeTok{nota  =} \FunctionTok{c}\NormalTok{(}\DecValTok{16}\NormalTok{, }\DecValTok{7}\NormalTok{, }\DecValTok{14}\NormalTok{, }\DecValTok{6}\NormalTok{, }\DecValTok{16}\NormalTok{))}
\NormalTok{estudantes }\OtherTok{\textless{}{-}} \FunctionTok{rbind}\NormalTok{(estudantes, }\FunctionTok{list}\NormalTok{(}\StringTok{"Filipa"}\NormalTok{, }\DecValTok{20}\NormalTok{, }\StringTok{"Matemática"}\NormalTok{, }\DecValTok{8}\NormalTok{))}
\NormalTok{estudantes[}\FunctionTok{order}\NormalTok{(estudantes}\SpecialCharTok{$}\NormalTok{nota), ]}
\end{Highlighting}
\end{Shaded}

Desde a versão 4.0 do R, \texttt{data.frame()} \textbf{já não} converte automaticamente
texto em fatores (o antigo \texttt{stringsAsFactors\ =\ TRUE} deixou de ser o padrão).
Por isso a coluna \texttt{curso} acima é de caracteres, e podemos atribuir-lhe novos
valores sem tropeçar em ``níveis'' inexistentes.

\textbf{2.} Modificar e recodificar por condição.

\begin{Shaded}
\begin{Highlighting}[]
\NormalTok{estudantes}\SpecialCharTok{$}\NormalTok{nota }\OtherTok{\textless{}{-}}\NormalTok{ estudantes}\SpecialCharTok{$}\NormalTok{nota }\SpecialCharTok{+} \FloatTok{0.5}
\NormalTok{estudantes}\SpecialCharTok{$}\NormalTok{curso[estudantes}\SpecialCharTok{$}\NormalTok{nota }\SpecialCharTok{\textgreater{}} \DecValTok{8}\NormalTok{] }\OtherTok{\textless{}{-}} \StringTok{"Estatística"}   \CommentTok{\# Alternativa a subset}
\end{Highlighting}
\end{Shaded}

\textbf{3.} Filtrar linhas (\texttt{subset} vs.~indexação).

\begin{Shaded}
\begin{Highlighting}[]
\FunctionTok{subset}\NormalTok{(mtcars, cyl }\SpecialCharTok{==} \DecValTok{6} \SpecialCharTok{\&}\NormalTok{ hp }\SpecialCharTok{\textgreater{}} \DecValTok{100}\NormalTok{)          }\CommentTok{\# forma legível}
\NormalTok{mtcars[mtcars}\SpecialCharTok{$}\NormalTok{cyl }\SpecialCharTok{==} \DecValTok{6} \SpecialCharTok{\&}\NormalTok{ mtcars}\SpecialCharTok{$}\NormalTok{hp }\SpecialCharTok{\textgreater{}} \DecValTok{100}\NormalTok{, ]  }\CommentTok{\# Alternativa com indexação}
\end{Highlighting}
\end{Shaded}

\textbf{4.} Criar uma nova coluna.

\begin{Shaded}
\begin{Highlighting}[]
\NormalTok{mtcars}\SpecialCharTok{$}\NormalTok{Peso\_KG }\OtherTok{\textless{}{-}}\NormalTok{ mtcars}\SpecialCharTok{$}\NormalTok{wt }\SpecialCharTok{*} \FloatTok{453.592}
\end{Highlighting}
\end{Shaded}

\textbf{5.} Agregar por grupo com \texttt{aggregate()}.

\begin{Shaded}
\begin{Highlighting}[]
\FunctionTok{aggregate}\NormalTok{(hp }\SpecialCharTok{\textasciitilde{}}\NormalTok{ cyl, }\AttributeTok{data =}\NormalTok{ mtcars, }\AttributeTok{FUN =}\NormalTok{ mean)}
\FunctionTok{aggregate}\NormalTok{(hp }\SpecialCharTok{\textasciitilde{}}\NormalTok{ cyl, }\AttributeTok{data =}\NormalTok{ mtcars, }\AttributeTok{FUN =}\NormalTok{ min)}
\FunctionTok{aggregate}\NormalTok{(hp }\SpecialCharTok{\textasciitilde{}}\NormalTok{ cyl, }\AttributeTok{data =}\NormalTok{ mtcars, }\AttributeTok{FUN =}\NormalTok{ max)}
\end{Highlighting}
\end{Shaded}

\textbf{6.} Aplicar uma função a todas as colunas.

\begin{Shaded}
\begin{Highlighting}[]
\FunctionTok{sapply}\NormalTok{(mtcars, sd)}
\end{Highlighting}
\end{Shaded}

\textbf{7.} Converter uma coluna em fator com rótulos.

\begin{Shaded}
\begin{Highlighting}[]
\NormalTok{mtcars}\SpecialCharTok{$}\NormalTok{am }\OtherTok{\textless{}{-}} \FunctionTok{factor}\NormalTok{(mtcars}\SpecialCharTok{$}\NormalTok{am, }\AttributeTok{levels =} \FunctionTok{c}\NormalTok{(}\DecValTok{0}\NormalTok{, }\DecValTok{1}\NormalTok{),}
                    \AttributeTok{labels =} \FunctionTok{c}\NormalTok{(}\StringTok{"Automático"}\NormalTok{, }\StringTok{"Manual"}\NormalTok{))}
\FunctionTok{table}\NormalTok{(mtcars}\SpecialCharTok{$}\NormalTok{am)}
\end{Highlighting}
\end{Shaded}

\textbf{8.} Ordenar e extrair extremos.

\begin{Shaded}
\begin{Highlighting}[]
\NormalTok{ord }\OtherTok{\textless{}{-}}\NormalTok{ mtcars[}\FunctionTok{order}\NormalTok{(mtcars}\SpecialCharTok{$}\NormalTok{mpg, }\AttributeTok{decreasing =} \ConstantTok{TRUE}\NormalTok{), ]}
\FunctionTok{head}\NormalTok{(ord, }\DecValTok{5}\NormalTok{)   }\CommentTok{\# 5 mais económicos}
\FunctionTok{tail}\NormalTok{(ord, }\DecValTok{5}\NormalTok{)   }\CommentTok{\# 5 menos económicos}
\end{Highlighting}
\end{Shaded}

\textbf{9.} Combinar \emph{data frames} com \texttt{merge()}.

\begin{Shaded}
\begin{Highlighting}[]
\NormalTok{mtcars}\SpecialCharTok{$}\NormalTok{car }\OtherTok{\textless{}{-}} \FunctionTok{rownames}\NormalTok{(mtcars)}
\NormalTok{carros1 }\OtherTok{\textless{}{-}}\NormalTok{ mtcars[, }\FunctionTok{c}\NormalTok{(}\StringTok{"car"}\NormalTok{, }\StringTok{"mpg"}\NormalTok{, }\StringTok{"cyl"}\NormalTok{)]}
\NormalTok{carros2 }\OtherTok{\textless{}{-}}\NormalTok{ mtcars[, }\FunctionTok{c}\NormalTok{(}\StringTok{"car"}\NormalTok{, }\StringTok{"hp"}\NormalTok{, }\StringTok{"wt"}\NormalTok{)]}
\FunctionTok{merge}\NormalTok{(carros1, carros2, }\AttributeTok{by =} \StringTok{"car"}\NormalTok{)}
\end{Highlighting}
\end{Shaded}

\textbf{10. a 14.} Exploração do conjunto \texttt{airquality}.

\begin{Shaded}
\begin{Highlighting}[]
\FunctionTok{data}\NormalTok{(airquality)}
\FunctionTok{head}\NormalTok{(airquality); }\FunctionTok{tail}\NormalTok{(airquality)               }\CommentTok{\# 10}
\FunctionTok{colSums}\NormalTok{(}\FunctionTok{is.na}\NormalTok{(airquality))                        }\CommentTok{\# 11 valores em falta por coluna}
\FunctionTok{mean}\NormalTok{(airquality}\SpecialCharTok{$}\NormalTok{Solar.R, }\AttributeTok{na.rm =} \ConstantTok{TRUE}\NormalTok{)            }\CommentTok{\# 12}
\FunctionTok{subset}\NormalTok{(airquality, Wind }\SpecialCharTok{\textgreater{}} \DecValTok{10} \SpecialCharTok{\&}\NormalTok{ Temp }\SpecialCharTok{\textgreater{}} \DecValTok{80}\NormalTok{)         }\CommentTok{\# 13 filtro}
\NormalTok{airquality}\SpecialCharTok{$}\NormalTok{Temp\_Wind }\OtherTok{\textless{}{-}}\NormalTok{ airquality}\SpecialCharTok{$}\NormalTok{Temp }\SpecialCharTok{+}\NormalTok{ airquality}\SpecialCharTok{$}\NormalTok{Wind  }\CommentTok{\# 14 nova coluna}
\end{Highlighting}
\end{Shaded}

\textbf{15.} Manipular o conjunto \texttt{iris}.

\begin{Shaded}
\begin{Highlighting}[]
\FunctionTok{data}\NormalTok{(iris)}
\NormalTok{setosa }\OtherTok{\textless{}{-}} \FunctionTok{subset}\NormalTok{(iris, Species }\SpecialCharTok{==} \StringTok{"setosa"}\NormalTok{)}
\NormalTok{setosa}\SpecialCharTok{$}\NormalTok{Petal.Area }\OtherTok{\textless{}{-}}\NormalTok{ setosa}\SpecialCharTok{$}\NormalTok{Petal.Length }\SpecialCharTok{*}\NormalTok{ setosa}\SpecialCharTok{$}\NormalTok{Petal.Width}
\NormalTok{setosa }\OtherTok{\textless{}{-}}\NormalTok{ setosa[setosa}\SpecialCharTok{$}\NormalTok{Petal.Area }\SpecialCharTok{\textgreater{}} \FloatTok{0.2}\NormalTok{, ]}
\NormalTok{setosa[}\FunctionTok{order}\NormalTok{(setosa}\SpecialCharTok{$}\NormalTok{Petal.Length, }\AttributeTok{decreasing =} \ConstantTok{TRUE}\NormalTok{), ]}
\end{Highlighting}
\end{Shaded}

\textbf{16.} Limpeza e agregação em \texttt{airquality}.

\begin{Shaded}
\begin{Highlighting}[]
\NormalTok{limpo }\OtherTok{\textless{}{-}} \FunctionTok{na.omit}\NormalTok{(airquality)}
\NormalTok{limpo}\SpecialCharTok{$}\NormalTok{Temp.Celsius }\OtherTok{\textless{}{-}}\NormalTok{ (limpo}\SpecialCharTok{$}\NormalTok{Temp }\SpecialCharTok{{-}} \DecValTok{32}\NormalTok{) }\SpecialCharTok{/} \FloatTok{1.8}
\FunctionTok{aggregate}\NormalTok{(Ozone }\SpecialCharTok{\textasciitilde{}}\NormalTok{ Month, }\AttributeTok{data =}\NormalTok{ limpo, }\AttributeTok{FUN =}\NormalTok{ mean)}
\FunctionTok{plot}\NormalTok{(limpo}\SpecialCharTok{$}\NormalTok{Temp.Celsius, limpo}\SpecialCharTok{$}\NormalTok{Ozone, }\AttributeTok{col =}\NormalTok{ limpo}\SpecialCharTok{$}\NormalTok{Month, }\AttributeTok{pch =} \DecValTok{16}\NormalTok{,}
     \AttributeTok{xlab =} \StringTok{"Temperatura (ºC)"}\NormalTok{, }\AttributeTok{ylab =} \StringTok{"Ozono"}\NormalTok{)}
\end{Highlighting}
\end{Shaded}

\textbf{17.} Um pequeno conjunto ao estilo \emph{Titanic}.

\begin{Shaded}
\begin{Highlighting}[]
\NormalTok{titanic }\OtherTok{\textless{}{-}} \FunctionTok{data.frame}\NormalTok{(}
  \AttributeTok{nome =} \FunctionTok{c}\NormalTok{(}\StringTok{"Ana"}\NormalTok{,}\StringTok{"Bruno"}\NormalTok{,}\StringTok{"Clara"}\NormalTok{,}\StringTok{"Daniel"}\NormalTok{,}\StringTok{"Eva"}\NormalTok{,}\StringTok{"Filipe"}\NormalTok{),}
  \AttributeTok{sexo =} \FunctionTok{c}\NormalTok{(}\StringTok{"F"}\NormalTok{,}\StringTok{"M"}\NormalTok{,}\StringTok{"F"}\NormalTok{,}\StringTok{"M"}\NormalTok{,}\StringTok{"F"}\NormalTok{,}\StringTok{"M"}\NormalTok{),}
  \AttributeTok{idade =} \FunctionTok{c}\NormalTok{(}\DecValTok{22}\NormalTok{,}\DecValTok{35}\NormalTok{,}\DecValTok{18}\NormalTok{,}\DecValTok{50}\NormalTok{,}\DecValTok{28}\NormalTok{,}\DecValTok{40}\NormalTok{),}
  \AttributeTok{classe =} \FunctionTok{c}\NormalTok{(}\StringTok{"1ª"}\NormalTok{,}\StringTok{"3ª"}\NormalTok{,}\StringTok{"3ª"}\NormalTok{,}\StringTok{"1ª"}\NormalTok{,}\StringTok{"2ª"}\NormalTok{,}\StringTok{"2ª"}\NormalTok{),}
  \AttributeTok{sobreviveu =} \FunctionTok{c}\NormalTok{(}\ConstantTok{TRUE}\NormalTok{,}\ConstantTok{FALSE}\NormalTok{,}\ConstantTok{TRUE}\NormalTok{,}\ConstantTok{TRUE}\NormalTok{,}\ConstantTok{FALSE}\NormalTok{,}\ConstantTok{FALSE}\NormalTok{))}

\FunctionTok{sum}\NormalTok{(titanic}\SpecialCharTok{$}\NormalTok{sexo }\SpecialCharTok{==} \StringTok{"F"} \SpecialCharTok{\&}\NormalTok{ titanic}\SpecialCharTok{$}\NormalTok{sobreviveu)     }\CommentTok{\# (b) mulheres sobreviventes}
\FunctionTok{mean}\NormalTok{(titanic}\SpecialCharTok{$}\NormalTok{idade[titanic}\SpecialCharTok{$}\NormalTok{sobreviveu])           }\CommentTok{\# (c) idade média dos sobreviventes}
\FunctionTok{sum}\NormalTok{(titanic}\SpecialCharTok{$}\NormalTok{classe }\SpecialCharTok{==} \StringTok{"3ª"}\NormalTok{)                       }\CommentTok{\# (d) passageiros de 3.ª}
\FunctionTok{mean}\NormalTok{(titanic}\SpecialCharTok{$}\NormalTok{sobreviveu[titanic}\SpecialCharTok{$}\NormalTok{classe }\SpecialCharTok{==} \StringTok{"3ª"}\NormalTok{])  }\CommentTok{\#     taxa de sobrevivência na 3.ª}
\end{Highlighting}
\end{Shaded}

A média de um vetor lógico é a \textbf{proporção de \texttt{TRUE}}, porque \texttt{TRUE} conta como
1 e \texttt{FALSE} como 0. Por isso \texttt{mean(titanic\$sobreviveu{[}...{]})} dá diretamente a
taxa de sobrevivência, sem necessidade de contar e dividir à mão.

\section{Leitura de dados e tabelas de frequência}\label{leitura-de-dados-e-tabelas-de-frequuxeancia}

Este conjunto de exercícios parte de um ficheiro externo (\texttt{dados\_instagram.xlsx})
e ilustra a construção de tabelas de frequência e de gráficos exploratórios. A
leitura faz-se com o pacote \texttt{readxl}.

\begin{Shaded}
\begin{Highlighting}[]
\FunctionTok{library}\NormalTok{(readxl)}
\NormalTok{dados }\OtherTok{\textless{}{-}} \FunctionTok{read\_excel}\NormalTok{(}\StringTok{"dados\_instagram.xlsx"}\NormalTok{)   }\CommentTok{\# ajuste o caminho ao seu ficheiro}
\end{Highlighting}
\end{Shaded}

\textbf{Variáveis qualitativas (género, país).} Para variáveis categóricas, a tabela
de frequências e o gráfico de barras (ou circular) são o ponto de partida.

\begin{Shaded}
\begin{Highlighting}[]
\NormalTok{freq\_abs }\OtherTok{\textless{}{-}} \FunctionTok{table}\NormalTok{(dados}\SpecialCharTok{$}\NormalTok{Genero)}
\NormalTok{tab }\OtherTok{\textless{}{-}} \FunctionTok{as.data.frame}\NormalTok{(freq\_abs)}
\FunctionTok{names}\NormalTok{(tab) }\OtherTok{\textless{}{-}} \FunctionTok{c}\NormalTok{(}\StringTok{"Género"}\NormalTok{, }\StringTok{"Freq\_Abs"}\NormalTok{)}
\NormalTok{tab}\SpecialCharTok{$}\NormalTok{Freq\_Rel }\OtherTok{\textless{}{-}} \FunctionTok{round}\NormalTok{(tab}\SpecialCharTok{$}\NormalTok{Freq\_Abs }\SpecialCharTok{/} \FunctionTok{sum}\NormalTok{(tab}\SpecialCharTok{$}\NormalTok{Freq\_Abs), }\DecValTok{2}\NormalTok{)}

\FunctionTok{barplot}\NormalTok{(freq\_abs, }\AttributeTok{col =} \StringTok{"lightblue"}\NormalTok{)}
\FunctionTok{pie}\NormalTok{(tab}\SpecialCharTok{$}\NormalTok{Freq\_Rel,}
    \AttributeTok{labels =} \FunctionTok{paste}\NormalTok{(tab}\SpecialCharTok{$}\NormalTok{Género, }\FunctionTok{round}\NormalTok{(tab}\SpecialCharTok{$}\NormalTok{Freq\_Rel }\SpecialCharTok{*} \DecValTok{100}\NormalTok{, }\DecValTok{1}\NormalTok{), }\StringTok{"\%"}\NormalTok{))}
\end{Highlighting}
\end{Shaded}

\textbf{Variáveis quantitativas (índice de influência, nº de publicações, seguidores,
likes).} Para variáveis contínuas, a tabela de frequências agrupa os dados em
classes. Uma escolha comum para o número de classes é a \textbf{regra de Sturges}.
Como o procedimento é o mesmo para todas as variáveis, vale a pena encapsulá-lo
numa função --- evita repetir o bloco quatro vezes.

\begin{Shaded}
\begin{Highlighting}[]
\NormalTok{tabela\_frequencias }\OtherTok{\textless{}{-}} \ControlFlowTok{function}\NormalTok{(x) \{}
\NormalTok{  x }\OtherTok{\textless{}{-}}\NormalTok{ x[}\SpecialCharTok{!}\FunctionTok{is.na}\NormalTok{(x)]}
\NormalTok{  k }\OtherTok{\textless{}{-}} \FunctionTok{ceiling}\NormalTok{(}\DecValTok{1} \SpecialCharTok{+} \FloatTok{3.322} \SpecialCharTok{*} \FunctionTok{log10}\NormalTok{(}\FunctionTok{length}\NormalTok{(x)))       }\CommentTok{\# nº de classes (Sturges)}
\NormalTok{  classes }\OtherTok{\textless{}{-}} \FunctionTok{cut}\NormalTok{(x, }\AttributeTok{breaks =} \FunctionTok{seq}\NormalTok{(}\FunctionTok{min}\NormalTok{(x), }\FunctionTok{max}\NormalTok{(x), }\AttributeTok{length.out =}\NormalTok{ k }\SpecialCharTok{+} \DecValTok{1}\NormalTok{),}
                 \AttributeTok{include.lowest =} \ConstantTok{TRUE}\NormalTok{)}
\NormalTok{  fa }\OtherTok{\textless{}{-}} \FunctionTok{table}\NormalTok{(classes)}
  \FunctionTok{data.frame}\NormalTok{(}
    \AttributeTok{Intervalo   =} \FunctionTok{names}\NormalTok{(fa),}
    \AttributeTok{Freq\_Abs    =} \FunctionTok{as.integer}\NormalTok{(fa),}
    \AttributeTok{Freq\_Rel    =} \FunctionTok{round}\NormalTok{(}\FunctionTok{as.integer}\NormalTok{(fa) }\SpecialCharTok{/} \FunctionTok{length}\NormalTok{(x), }\DecValTok{3}\NormalTok{),}
    \AttributeTok{Freq\_Abs\_Ac =} \FunctionTok{cumsum}\NormalTok{(}\FunctionTok{as.integer}\NormalTok{(fa)),}
    \AttributeTok{Freq\_Rel\_Ac =} \FunctionTok{round}\NormalTok{(}\FunctionTok{cumsum}\NormalTok{(}\FunctionTok{as.integer}\NormalTok{(fa)) }\SpecialCharTok{/} \FunctionTok{length}\NormalTok{(x), }\DecValTok{3}\NormalTok{))}
\NormalTok{\}}

\FunctionTok{tabela\_frequencias}\NormalTok{(dados}\SpecialCharTok{$}\NormalTok{indice\_influencia)}
\FunctionTok{hist}\NormalTok{(dados}\SpecialCharTok{$}\NormalTok{indice\_influencia, }\AttributeTok{breaks =} \StringTok{"Sturges"}\NormalTok{, }\AttributeTok{col =} \StringTok{"lightblue"}\NormalTok{)}
\end{Highlighting}
\end{Shaded}

No material original, a tabela da variável \texttt{media\_likes} reutilizava por
distração as contagens de \texttt{num\_seguidores} (\texttt{tabela\_frequencia\_ns\_df}) no cálculo
das frequências relativas e acumuladas. Encapsular o procedimento numa função,
como acima, elimina de raiz este tipo de erro de copiar-colar: cada variável é
processada de forma independente.

\textbf{Medidas por grupo.} Para comparar grupos, \texttt{tapply()} aplica uma função a uma
variável quantitativa separada por uma variável qualitativa.

\begin{Shaded}
\begin{Highlighting}[]
\FunctionTok{tapply}\NormalTok{(dados}\SpecialCharTok{$}\NormalTok{indice\_influencia, dados}\SpecialCharTok{$}\NormalTok{Genero, mean)}
\FunctionTok{tapply}\NormalTok{(dados}\SpecialCharTok{$}\NormalTok{media\_likes,       dados}\SpecialCharTok{$}\NormalTok{Genero, median)}
\FunctionTok{quantile}\NormalTok{(dados}\SpecialCharTok{$}\NormalTok{num\_seguidores, }\AttributeTok{probs =} \FunctionTok{c}\NormalTok{(}\FloatTok{0.10}\NormalTok{, }\FloatTok{0.25}\NormalTok{, }\FloatTok{0.75}\NormalTok{, }\FloatTok{0.90}\NormalTok{))}
\end{Highlighting}
\end{Shaded}

\section{Gráficos com o R base}\label{gruxe1ficos-com-o-r-base}

\textbf{1.} Gráfico de barras de frequências (absolutas e relativas).

\begin{Shaded}
\begin{Highlighting}[]
\FunctionTok{data}\NormalTok{(iris)}
\FunctionTok{barplot}\NormalTok{(}\FunctionTok{table}\NormalTok{(iris}\SpecialCharTok{$}\NormalTok{Species), }\AttributeTok{col =} \FunctionTok{c}\NormalTok{(}\StringTok{"blue"}\NormalTok{,}\StringTok{"green"}\NormalTok{,}\StringTok{"yellow"}\NormalTok{),}
        \AttributeTok{main =} \StringTok{"Contagem de espécies"}\NormalTok{, }\AttributeTok{xlab =} \StringTok{"Espécie"}\NormalTok{,}
        \AttributeTok{ylab =} \StringTok{"Frequência absoluta"}\NormalTok{)}
\FunctionTok{barplot}\NormalTok{(}\FunctionTok{table}\NormalTok{(iris}\SpecialCharTok{$}\NormalTok{Species) }\SpecialCharTok{/} \FunctionTok{length}\NormalTok{(iris}\SpecialCharTok{$}\NormalTok{Species),}
        \AttributeTok{col =} \FunctionTok{c}\NormalTok{(}\StringTok{"blue"}\NormalTok{,}\StringTok{"green"}\NormalTok{,}\StringTok{"yellow"}\NormalTok{), }\AttributeTok{ylab =} \StringTok{"Frequência relativa"}\NormalTok{)}
\end{Highlighting}
\end{Shaded}

\textbf{2.} Barras de uma média por grupo.

\begin{Shaded}
\begin{Highlighting}[]
\FunctionTok{data}\NormalTok{(airquality)}
\NormalTok{media\_mes }\OtherTok{\textless{}{-}} \FunctionTok{tapply}\NormalTok{(airquality}\SpecialCharTok{$}\NormalTok{Temp, airquality}\SpecialCharTok{$}\NormalTok{Month, mean, }\AttributeTok{na.rm =} \ConstantTok{TRUE}\NormalTok{)}
\FunctionTok{barplot}\NormalTok{(media\_mes, }\AttributeTok{col =} \StringTok{"lightblue"}\NormalTok{,}
        \AttributeTok{names.arg =} \FunctionTok{c}\NormalTok{(}\StringTok{"Maio"}\NormalTok{,}\StringTok{"Junho"}\NormalTok{,}\StringTok{"Julho"}\NormalTok{,}\StringTok{"Agosto"}\NormalTok{,}\StringTok{"Setembro"}\NormalTok{),}
        \AttributeTok{main =} \StringTok{"Temperatura média por mês"}\NormalTok{, }\AttributeTok{ylab =} \StringTok{"Temperatura (ºF)"}\NormalTok{)}
\end{Highlighting}
\end{Shaded}

\textbf{3. e 4.} Gráficos circulares.

\begin{Shaded}
\begin{Highlighting}[]
\NormalTok{fr }\OtherTok{\textless{}{-}} \FunctionTok{table}\NormalTok{(iris}\SpecialCharTok{$}\NormalTok{Species) }\SpecialCharTok{/} \FunctionTok{length}\NormalTok{(iris}\SpecialCharTok{$}\NormalTok{Species)}
\FunctionTok{pie}\NormalTok{(fr, }\AttributeTok{labels =} \FunctionTok{paste}\NormalTok{(}\FunctionTok{names}\NormalTok{(fr), }\FunctionTok{round}\NormalTok{(fr }\SpecialCharTok{*} \DecValTok{100}\NormalTok{, }\DecValTok{1}\NormalTok{), }\StringTok{"\%"}\NormalTok{),}
    \AttributeTok{col =} \FunctionTok{c}\NormalTok{(}\StringTok{"blue"}\NormalTok{,}\StringTok{"green"}\NormalTok{,}\StringTok{"yellow"}\NormalTok{))}
\end{Highlighting}
\end{Shaded}

\textbf{5.} Histograma com curva de densidade sobreposta.

\begin{Shaded}
\begin{Highlighting}[]
\FunctionTok{data}\NormalTok{(iris)}
\FunctionTok{hist}\NormalTok{(iris}\SpecialCharTok{$}\NormalTok{Sepal.Length, }\AttributeTok{col =} \StringTok{"skyblue"}\NormalTok{, }\AttributeTok{breaks =} \DecValTok{15}\NormalTok{, }\AttributeTok{freq =} \ConstantTok{FALSE}\NormalTok{,}
     \AttributeTok{main =} \StringTok{"Comprimento da sépala"}\NormalTok{, }\AttributeTok{xlab =} \StringTok{"cm"}\NormalTok{)}
\FunctionTok{lines}\NormalTok{(}\FunctionTok{density}\NormalTok{(iris}\SpecialCharTok{$}\NormalTok{Sepal.Length), }\AttributeTok{col =} \StringTok{"red"}\NormalTok{, }\AttributeTok{lwd =} \DecValTok{2}\NormalTok{)}
\end{Highlighting}
\end{Shaded}

\textbf{6.} Dois histogramas sobrepostos (com transparência).

\begin{Shaded}
\begin{Highlighting}[]
\FunctionTok{set.seed}\NormalTok{(}\DecValTok{123}\NormalTok{)}
\NormalTok{A }\OtherTok{\textless{}{-}} \FunctionTok{rnorm}\NormalTok{(}\DecValTok{100}\NormalTok{, }\DecValTok{150}\NormalTok{, }\DecValTok{10}\NormalTok{); B }\OtherTok{\textless{}{-}} \FunctionTok{rnorm}\NormalTok{(}\DecValTok{100}\NormalTok{, }\DecValTok{160}\NormalTok{, }\DecValTok{15}\NormalTok{)}
\FunctionTok{hist}\NormalTok{(A, }\AttributeTok{col =} \FunctionTok{rgb}\NormalTok{(}\DecValTok{1}\NormalTok{,}\DecValTok{0}\NormalTok{,}\DecValTok{0}\NormalTok{,}\FloatTok{0.5}\NormalTok{), }\AttributeTok{breaks =} \DecValTok{10}\NormalTok{, }\AttributeTok{xlim =} \FunctionTok{c}\NormalTok{(}\DecValTok{100}\NormalTok{, }\DecValTok{200}\NormalTok{),}
     \AttributeTok{main =} \StringTok{"Distribuição de alturas"}\NormalTok{, }\AttributeTok{xlab =} \StringTok{"Altura (cm)"}\NormalTok{)}
\FunctionTok{hist}\NormalTok{(B, }\AttributeTok{col =} \FunctionTok{rgb}\NormalTok{(}\DecValTok{0}\NormalTok{,}\DecValTok{0}\NormalTok{,}\DecValTok{1}\NormalTok{,}\FloatTok{0.5}\NormalTok{), }\AttributeTok{breaks =} \DecValTok{10}\NormalTok{, }\AttributeTok{add =} \ConstantTok{TRUE}\NormalTok{)}
\FunctionTok{legend}\NormalTok{(}\StringTok{"topright"}\NormalTok{, }\FunctionTok{c}\NormalTok{(}\StringTok{"Grupo A"}\NormalTok{,}\StringTok{"Grupo B"}\NormalTok{),}
       \AttributeTok{fill =} \FunctionTok{c}\NormalTok{(}\FunctionTok{rgb}\NormalTok{(}\DecValTok{1}\NormalTok{,}\DecValTok{0}\NormalTok{,}\DecValTok{0}\NormalTok{,}\FloatTok{0.5}\NormalTok{), }\FunctionTok{rgb}\NormalTok{(}\DecValTok{0}\NormalTok{,}\DecValTok{0}\NormalTok{,}\DecValTok{1}\NormalTok{,}\FloatTok{0.5}\NormalTok{)))}
\end{Highlighting}
\end{Shaded}

\textbf{7.} Histograma com linhas de referência (média e mediana).

\begin{Shaded}
\begin{Highlighting}[]
\FunctionTok{set.seed}\NormalTok{(}\DecValTok{1234}\NormalTok{)}
\NormalTok{notas }\OtherTok{\textless{}{-}} \FunctionTok{rnorm}\NormalTok{(}\DecValTok{200}\NormalTok{, }\AttributeTok{mean =} \DecValTok{14}\NormalTok{, }\AttributeTok{sd =} \DecValTok{4}\NormalTok{)}
\FunctionTok{hist}\NormalTok{(notas, }\AttributeTok{col =} \StringTok{"lightblue"}\NormalTok{, }\AttributeTok{breaks =} \DecValTok{15}\NormalTok{, }\AttributeTok{main =} \StringTok{"Distribuição de notas"}\NormalTok{)}
\FunctionTok{abline}\NormalTok{(}\AttributeTok{v =} \FunctionTok{mean}\NormalTok{(notas),   }\AttributeTok{col =} \StringTok{"red"}\NormalTok{,  }\AttributeTok{lwd =} \DecValTok{2}\NormalTok{, }\AttributeTok{lty =} \DecValTok{2}\NormalTok{)}
\FunctionTok{abline}\NormalTok{(}\AttributeTok{v =} \FunctionTok{median}\NormalTok{(notas), }\AttributeTok{col =} \StringTok{"blue"}\NormalTok{, }\AttributeTok{lwd =} \DecValTok{2}\NormalTok{, }\AttributeTok{lty =} \DecValTok{2}\NormalTok{)}
\end{Highlighting}
\end{Shaded}

\textbf{8.} Vários gráficos numa só figura com \texttt{par(mfrow=...)}.

\begin{Shaded}
\begin{Highlighting}[]
\FunctionTok{data}\NormalTok{(airquality)}
\FunctionTok{par}\NormalTok{(}\AttributeTok{mfrow =} \FunctionTok{c}\NormalTok{(}\DecValTok{3}\NormalTok{, }\DecValTok{2}\NormalTok{))}
\ControlFlowTok{for}\NormalTok{ (m }\ControlFlowTok{in} \DecValTok{5}\SpecialCharTok{:}\DecValTok{9}\NormalTok{) \{}
  \FunctionTok{hist}\NormalTok{(airquality}\SpecialCharTok{$}\NormalTok{Ozone[airquality}\SpecialCharTok{$}\NormalTok{Month }\SpecialCharTok{==}\NormalTok{ m],}
       \AttributeTok{main =} \FunctionTok{paste}\NormalTok{(}\StringTok{"Ozono — mês"}\NormalTok{, m), }\AttributeTok{xlab =} \StringTok{"Ozono"}\NormalTok{,}
       \AttributeTok{col =} \StringTok{"lightblue"}\NormalTok{, }\AttributeTok{breaks =} \DecValTok{10}\NormalTok{)}
\NormalTok{\}}
\FunctionTok{par}\NormalTok{(}\AttributeTok{mfrow =} \FunctionTok{c}\NormalTok{(}\DecValTok{1}\NormalTok{, }\DecValTok{1}\NormalTok{))}
\end{Highlighting}
\end{Shaded}

\texttt{hist()} não tem argumento \texttt{na.rm}; os \texttt{NA} são simplesmente ignorados no
traçado. Passar \texttt{na.rm\ =\ TRUE} a \texttt{hist()}, como surgia no material original, não
provoca erro mas também não tem qualquer efeito --- é um argumento inexistente
para esta função.

\textbf{9. e 10.} Diagramas de caixa (\emph{boxplots}) por grupo.

\begin{Shaded}
\begin{Highlighting}[]
\FunctionTok{boxplot}\NormalTok{(Petal.Length }\SpecialCharTok{\textasciitilde{}}\NormalTok{ Species, }\AttributeTok{data =}\NormalTok{ iris,}
        \AttributeTok{col =} \FunctionTok{c}\NormalTok{(}\StringTok{"lightblue"}\NormalTok{,}\StringTok{"lightgreen"}\NormalTok{,}\StringTok{"lightpink"}\NormalTok{),}
        \AttributeTok{main =} \StringTok{"Comprimento da pétala por espécie"}\NormalTok{)}

\FunctionTok{boxplot}\NormalTok{(Ozone }\SpecialCharTok{\textasciitilde{}}\NormalTok{ Month, }\AttributeTok{data =}\NormalTok{ airquality,}
        \AttributeTok{main =} \StringTok{"Ozono por mês"}\NormalTok{, }\AttributeTok{xlab =} \StringTok{"Mês"}\NormalTok{, }\AttributeTok{ylab =} \StringTok{"Ozono"}\NormalTok{)}
\FunctionTok{abline}\NormalTok{(}\AttributeTok{h =} \FunctionTok{median}\NormalTok{(airquality}\SpecialCharTok{$}\NormalTok{Ozone, }\AttributeTok{na.rm =} \ConstantTok{TRUE}\NormalTok{), }\AttributeTok{col =} \StringTok{"red"}\NormalTok{, }\AttributeTok{lwd =} \DecValTok{2}\NormalTok{, }\AttributeTok{lty =} \DecValTok{2}\NormalTok{)}
\end{Highlighting}
\end{Shaded}

\textbf{11. e 12.} Diagramas de dispersão e séries.

\begin{Shaded}
\begin{Highlighting}[]
\NormalTok{altura }\OtherTok{\textless{}{-}} \FunctionTok{runif}\NormalTok{(}\DecValTok{30}\NormalTok{, }\DecValTok{150}\NormalTok{, }\DecValTok{200}\NormalTok{)}
\NormalTok{peso }\OtherTok{\textless{}{-}}\NormalTok{ altura }\SpecialCharTok{*} \FloatTok{0.4} \SpecialCharTok{+} \FunctionTok{runif}\NormalTok{(}\DecValTok{30}\NormalTok{, }\SpecialCharTok{{-}}\DecValTok{5}\NormalTok{, }\DecValTok{5}\NormalTok{)}
\FunctionTok{plot}\NormalTok{(altura, peso, }\AttributeTok{xlab =} \StringTok{"Altura (cm)"}\NormalTok{, }\AttributeTok{ylab =} \StringTok{"Peso (kg)"}\NormalTok{)}

\NormalTok{anos }\OtherTok{\textless{}{-}} \DecValTok{2000}\SpecialCharTok{:}\DecValTok{2019}
\NormalTok{pop  }\OtherTok{\textless{}{-}} \DecValTok{100000} \SpecialCharTok{*} \FunctionTok{cumprod}\NormalTok{(}\FloatTok{1.02} \SpecialCharTok{+} \FunctionTok{runif}\NormalTok{(}\DecValTok{20}\NormalTok{, }\SpecialCharTok{{-}}\FloatTok{0.01}\NormalTok{, }\FloatTok{0.01}\NormalTok{))}
\FunctionTok{plot}\NormalTok{(anos, pop, }\AttributeTok{type =} \StringTok{"l"}\NormalTok{, }\AttributeTok{col =} \StringTok{"green"}\NormalTok{,}
     \AttributeTok{xlab =} \StringTok{"Ano"}\NormalTok{, }\AttributeTok{ylab =} \StringTok{"População"}\NormalTok{)}
\end{Highlighting}
\end{Shaded}

\section{\texorpdfstring{Estruturas condicionais (\texttt{if} \ldots{} \texttt{else})}{Estruturas condicionais (if \ldots{} else)}}\label{estruturas-condicionais-if-else}

Vários destes exercícios leem números do teclado com \texttt{readline()}. Esse comando
só funciona numa sessão interativa; num documento compilado deve usar-se
\texttt{eval=FALSE}, como aqui.

\textbf{1. a 5.} Decisões simples e encadeadas.

\begin{Shaded}
\begin{Highlighting}[]
\CommentTok{\# 1. sinal de um número}
\NormalTok{num }\OtherTok{\textless{}{-}} \FunctionTok{as.numeric}\NormalTok{(}\FunctionTok{readline}\NormalTok{(}\StringTok{"Introduza um número: "}\NormalTok{))}
\ControlFlowTok{if}\NormalTok{ (num }\SpecialCharTok{\textgreater{}} \DecValTok{0}\NormalTok{) }\FunctionTok{print}\NormalTok{(}\StringTok{"O número é positivo"}\NormalTok{)}

\CommentTok{\# 2. par ou ímpar}
\ControlFlowTok{if}\NormalTok{ (num }\SpecialCharTok{\%\%} \DecValTok{2} \SpecialCharTok{==} \DecValTok{0}\NormalTok{) }\FunctionTok{print}\NormalTok{(}\StringTok{"Par"}\NormalTok{) }\ControlFlowTok{else} \FunctionTok{print}\NormalTok{(}\StringTok{"Ímpar"}\NormalTok{)}

\CommentTok{\# 4. faixa etária}
\NormalTok{idade }\OtherTok{\textless{}{-}} \FunctionTok{as.numeric}\NormalTok{(}\FunctionTok{readline}\NormalTok{(}\StringTok{"Idade: "}\NormalTok{))}
\ControlFlowTok{if}\NormalTok{ (idade }\SpecialCharTok{\textless{}} \DecValTok{0}\NormalTok{)               }\FunctionTok{print}\NormalTok{(}\StringTok{"Idade inválida"}\NormalTok{) }\ControlFlowTok{else}
\ControlFlowTok{if}\NormalTok{ (idade }\SpecialCharTok{\textless{}=} \DecValTok{12}\NormalTok{)             }\FunctionTok{print}\NormalTok{(}\StringTok{"Criança"}\NormalTok{)        }\ControlFlowTok{else}
\ControlFlowTok{if}\NormalTok{ (idade }\SpecialCharTok{\textless{}=} \DecValTok{17}\NormalTok{)             }\FunctionTok{print}\NormalTok{(}\StringTok{"Adolescente"}\NormalTok{)    }\ControlFlowTok{else}
\ControlFlowTok{if}\NormalTok{ (idade }\SpecialCharTok{\textless{}=} \DecValTok{64}\NormalTok{)             }\FunctionTok{print}\NormalTok{(}\StringTok{"Adulto"}\NormalTok{)         }\ControlFlowTok{else}
                             \FunctionTok{print}\NormalTok{(}\StringTok{"Idoso"}\NormalTok{)}

\CommentTok{\# 5. desconto por escalão de preço}
\NormalTok{preco }\OtherTok{\textless{}{-}} \FunctionTok{as.numeric}\NormalTok{(}\FunctionTok{readline}\NormalTok{(}\StringTok{"Preço: "}\NormalTok{))}
\ControlFlowTok{if}\NormalTok{ (preco }\SpecialCharTok{\textless{}} \DecValTok{0}\NormalTok{)          }\FunctionTok{print}\NormalTok{(}\StringTok{"Preço inválido"}\NormalTok{) }\ControlFlowTok{else}
\ControlFlowTok{if}\NormalTok{ (preco }\SpecialCharTok{\textless{}} \DecValTok{100}\NormalTok{)        }\FunctionTok{print}\NormalTok{(preco }\SpecialCharTok{*} \FloatTok{0.95}\NormalTok{)     }\ControlFlowTok{else}
\ControlFlowTok{if}\NormalTok{ (preco }\SpecialCharTok{\textless{}=} \DecValTok{500}\NormalTok{)       }\FunctionTok{print}\NormalTok{(preco }\SpecialCharTok{*} \FloatTok{0.90}\NormalTok{)     }\ControlFlowTok{else}
                        \FunctionTok{print}\NormalTok{(preco }\SpecialCharTok{*} \FloatTok{0.85}\NormalTok{)}
\end{Highlighting}
\end{Shaded}

\textbf{6.} Localização de um ponto no plano.

\begin{Shaded}
\begin{Highlighting}[]
\NormalTok{x }\OtherTok{\textless{}{-}} \FunctionTok{as.numeric}\NormalTok{(}\FunctionTok{readline}\NormalTok{(}\StringTok{"x: "}\NormalTok{)); y }\OtherTok{\textless{}{-}} \FunctionTok{as.numeric}\NormalTok{(}\FunctionTok{readline}\NormalTok{(}\StringTok{"y: "}\NormalTok{))}
\ControlFlowTok{if}\NormalTok{ (x }\SpecialCharTok{==} \DecValTok{0} \SpecialCharTok{\&\&}\NormalTok{ y }\SpecialCharTok{==} \DecValTok{0}\NormalTok{)      }\FunctionTok{print}\NormalTok{(}\StringTok{"Origem"}\NormalTok{)            }\ControlFlowTok{else}
\ControlFlowTok{if}\NormalTok{ (x }\SpecialCharTok{==} \DecValTok{0}\NormalTok{)                }\FunctionTok{print}\NormalTok{(}\StringTok{"Eixo Y"}\NormalTok{)            }\ControlFlowTok{else}
\ControlFlowTok{if}\NormalTok{ (y }\SpecialCharTok{==} \DecValTok{0}\NormalTok{)                }\FunctionTok{print}\NormalTok{(}\StringTok{"Eixo X"}\NormalTok{)            }\ControlFlowTok{else}
\ControlFlowTok{if}\NormalTok{ (x }\SpecialCharTok{\textgreater{}} \DecValTok{0} \SpecialCharTok{\&\&}\NormalTok{ y }\SpecialCharTok{\textgreater{}} \DecValTok{0}\NormalTok{)        }\FunctionTok{print}\NormalTok{(}\StringTok{"1.º quadrante"}\NormalTok{)     }\ControlFlowTok{else}
\ControlFlowTok{if}\NormalTok{ (x }\SpecialCharTok{\textless{}} \DecValTok{0} \SpecialCharTok{\&\&}\NormalTok{ y }\SpecialCharTok{\textgreater{}} \DecValTok{0}\NormalTok{)        }\FunctionTok{print}\NormalTok{(}\StringTok{"2.º quadrante"}\NormalTok{)     }\ControlFlowTok{else}
\ControlFlowTok{if}\NormalTok{ (x }\SpecialCharTok{\textless{}} \DecValTok{0} \SpecialCharTok{\&\&}\NormalTok{ y }\SpecialCharTok{\textless{}} \DecValTok{0}\NormalTok{)        }\FunctionTok{print}\NormalTok{(}\StringTok{"3.º quadrante"}\NormalTok{)     }\ControlFlowTok{else}
                           \FunctionTok{print}\NormalTok{(}\StringTok{"4.º quadrante"}\NormalTok{)}
\end{Highlighting}
\end{Shaded}

No material original os quadrantes estavam trocados: a condição \texttt{x\textless{}0\ \&\ y\textgreater{}0}
aparecia rotulada como ``primeiro quadrante'' e \texttt{x\textgreater{}0\ \&\ y\textgreater{}0} como ``segundo''. A
convenção correta é: 1.º quadrante \(x>0,\,y>0\); 2.º \(x<0,\,y>0\); 3.º
\(x<0,\,y<0\); 4.º \(x>0,\,y<0\). A versão acima corrige a atribuição.

\textbf{7.} Ano bissexto.

\begin{Shaded}
\begin{Highlighting}[]
\NormalTok{ano }\OtherTok{\textless{}{-}} \FunctionTok{as.numeric}\NormalTok{(}\FunctionTok{readline}\NormalTok{(}\StringTok{"Ano: "}\NormalTok{))}
\ControlFlowTok{if}\NormalTok{ ((ano }\SpecialCharTok{\%\%} \DecValTok{4} \SpecialCharTok{==} \DecValTok{0} \SpecialCharTok{\&\&}\NormalTok{ ano }\SpecialCharTok{\%\%} \DecValTok{100} \SpecialCharTok{!=} \DecValTok{0}\NormalTok{) }\SpecialCharTok{||}\NormalTok{ ano }\SpecialCharTok{\%\%} \DecValTok{400} \SpecialCharTok{==} \DecValTok{0}\NormalTok{)}
  \FunctionTok{print}\NormalTok{(}\StringTok{"Ano bissexto"}\NormalTok{) }\ControlFlowTok{else} \FunctionTok{print}\NormalTok{(}\StringTok{"Não é bissexto"}\NormalTok{)}
\end{Highlighting}
\end{Shaded}

\textbf{8.} Decompor um número em centenas, dezenas e unidades.

\begin{Shaded}
\begin{Highlighting}[]
\NormalTok{num }\OtherTok{\textless{}{-}} \FunctionTok{as.numeric}\NormalTok{(}\FunctionTok{readline}\NormalTok{(}\StringTok{"Número (0–999): "}\NormalTok{))}
\ControlFlowTok{if}\NormalTok{ (num }\SpecialCharTok{\textgreater{}=} \DecValTok{0} \SpecialCharTok{\&\&}\NormalTok{ num }\SpecialCharTok{\textless{}=} \DecValTok{999}\NormalTok{) \{}
\NormalTok{  c }\OtherTok{\textless{}{-}}\NormalTok{ num }\SpecialCharTok{\%/\%} \DecValTok{100}\NormalTok{; d }\OtherTok{\textless{}{-}}\NormalTok{ (num }\SpecialCharTok{\%\%} \DecValTok{100}\NormalTok{) }\SpecialCharTok{\%/\%} \DecValTok{10}\NormalTok{; u }\OtherTok{\textless{}{-}}\NormalTok{ num }\SpecialCharTok{\%\%} \DecValTok{10}
  \FunctionTok{print}\NormalTok{(}\FunctionTok{paste}\NormalTok{(c, }\StringTok{"centenas"}\NormalTok{, d, }\StringTok{"dezenas"}\NormalTok{, u, }\StringTok{"unidades"}\NormalTok{))}
\NormalTok{\} }\ControlFlowTok{else} \FunctionTok{print}\NormalTok{(}\StringTok{"Fora do intervalo"}\NormalTok{)}
\end{Highlighting}
\end{Shaded}

\textbf{9.} Numeração romana (1 a 5).

\begin{Shaded}
\begin{Highlighting}[]
\NormalTok{num }\OtherTok{\textless{}{-}} \FunctionTok{as.numeric}\NormalTok{(}\FunctionTok{readline}\NormalTok{(}\StringTok{"Número (1–5): "}\NormalTok{))}
\NormalTok{romanos }\OtherTok{\textless{}{-}} \FunctionTok{c}\NormalTok{(}\StringTok{"I"}\NormalTok{, }\StringTok{"II"}\NormalTok{, }\StringTok{"III"}\NormalTok{, }\StringTok{"IV"}\NormalTok{, }\StringTok{"V"}\NormalTok{)}
\ControlFlowTok{if}\NormalTok{ (num }\SpecialCharTok{\textgreater{}=} \DecValTok{1} \SpecialCharTok{\&\&}\NormalTok{ num }\SpecialCharTok{\textless{}=} \DecValTok{5}\NormalTok{) }\FunctionTok{print}\NormalTok{(romanos[num]) }\ControlFlowTok{else} \FunctionTok{print}\NormalTok{(}\StringTok{"Fora do intervalo"}\NormalTok{)}
\end{Highlighting}
\end{Shaded}

Em vez de cinco ramos \texttt{if}/\texttt{else\ if} para converter 1--5 em numeração romana,
um único vetor indexado --- \texttt{romanos{[}num{]}} --- resolve o problema numa linha. Sempre
que uma cadeia de \texttt{if}s apenas escolhe um de vários valores fixos consoante um
inteiro, um vetor de consulta é mais limpo.

\textbf{10.} Menor de vários números (versão vetorizada).

\begin{Shaded}
\begin{Highlighting}[]
\CommentTok{\# Alternativa concisa à leitura número a número:}
\NormalTok{nums }\OtherTok{\textless{}{-}} \FunctionTok{as.numeric}\NormalTok{(}\FunctionTok{strsplit}\NormalTok{(}\FunctionTok{readline}\NormalTok{(}\StringTok{"Números separados por espaço: "}\NormalTok{), }\StringTok{" "}\NormalTok{)[[}\DecValTok{1}\NormalTok{]])}
\FunctionTok{min}\NormalTok{(nums)}
\end{Highlighting}
\end{Shaded}

\textbf{11. a 13.} \texttt{ifelse()} vetorizado.

\begin{Shaded}
\begin{Highlighting}[]
\FunctionTok{ifelse}\NormalTok{(}\DecValTok{1}\SpecialCharTok{:}\DecValTok{10} \SpecialCharTok{\%\%} \DecValTok{2} \SpecialCharTok{==} \DecValTok{0}\NormalTok{, }\StringTok{"par"}\NormalTok{, }\StringTok{"ímpar"}\NormalTok{)                    }\CommentTok{\# 11}

\NormalTok{notas }\OtherTok{\textless{}{-}} \FunctionTok{c}\NormalTok{(}\DecValTok{50}\NormalTok{, }\DecValTok{85}\NormalTok{, }\DecValTok{70}\NormalTok{, }\DecValTok{40}\NormalTok{, }\DecValTok{60}\NormalTok{, }\DecValTok{90}\NormalTok{)                        }\CommentTok{\# 12}
\FunctionTok{ifelse}\NormalTok{(notas }\SpecialCharTok{\textgreater{}=} \DecValTok{60}\NormalTok{, }\StringTok{"Aprovado"}\NormalTok{,}
       \FunctionTok{ifelse}\NormalTok{(notas }\SpecialCharTok{\textgreater{}=} \DecValTok{50}\NormalTok{, }\StringTok{"Aprovado com restrição"}\NormalTok{, }\StringTok{"Reprovado"}\NormalTok{))}

\NormalTok{idades }\OtherTok{\textless{}{-}} \FunctionTok{c}\NormalTok{(}\DecValTok{1}\NormalTok{, }\DecValTok{2}\NormalTok{, }\DecValTok{14}\NormalTok{, }\DecValTok{20}\NormalTok{, }\DecValTok{35}\NormalTok{, }\DecValTok{67}\NormalTok{, }\DecValTok{80}\NormalTok{)                     }\CommentTok{\# 13}
\FunctionTok{ifelse}\NormalTok{(idades }\SpecialCharTok{\textless{}=} \DecValTok{12}\NormalTok{, }\StringTok{"Criança"}\NormalTok{,}
       \FunctionTok{ifelse}\NormalTok{(idades }\SpecialCharTok{\textless{}=} \DecValTok{18}\NormalTok{, }\StringTok{"Adolescente"}\NormalTok{, }\StringTok{"Adulto"}\NormalTok{))}
\end{Highlighting}
\end{Shaded}

\texttt{if} avalia \textbf{uma única} condição (um valor lógico); aplicá-lo a um vetor usa
apenas o primeiro elemento e, nas versões recentes do R, dá erro. Para decidir
elemento a elemento ao longo de um vetor, o instrumento certo é \texttt{ifelse()}.

\section{\texorpdfstring{Ciclos \texttt{for}}{Ciclos for}}\label{ciclos-for}

\textbf{1. a 6.} Percursos e acumuladores.

\begin{Shaded}
\begin{Highlighting}[]
\ControlFlowTok{for}\NormalTok{ (i }\ControlFlowTok{in} \DecValTok{0}\SpecialCharTok{:}\DecValTok{10}\NormalTok{) }\FunctionTok{print}\NormalTok{(i)                               }\CommentTok{\# 1}

\ControlFlowTok{for}\NormalTok{ (i }\ControlFlowTok{in} \DecValTok{1}\SpecialCharTok{:}\DecValTok{5}\NormalTok{) }\FunctionTok{print}\NormalTok{(}\FunctionTok{paste}\NormalTok{(}\StringTok{"O quadrado de"}\NormalTok{, i, }\StringTok{"é"}\NormalTok{, i}\SpecialCharTok{\^{}}\DecValTok{2}\NormalTok{))  }\CommentTok{\# 2}

\NormalTok{soma }\OtherTok{\textless{}{-}} \DecValTok{0}                                              \CommentTok{\# 3}
\ControlFlowTok{for}\NormalTok{ (num }\ControlFlowTok{in} \FunctionTok{c}\NormalTok{(}\DecValTok{2}\NormalTok{,}\DecValTok{4}\NormalTok{,}\DecValTok{6}\NormalTok{,}\DecValTok{8}\NormalTok{,}\DecValTok{10}\NormalTok{)) soma }\OtherTok{\textless{}{-}}\NormalTok{ soma }\SpecialCharTok{+}\NormalTok{ num}
\NormalTok{soma}

\NormalTok{notas }\OtherTok{\textless{}{-}} \FunctionTok{c}\NormalTok{(}\DecValTok{85}\NormalTok{, }\DecValTok{90}\NormalTok{, }\DecValTok{78}\NormalTok{, }\DecValTok{92}\NormalTok{, }\DecValTok{88}\NormalTok{)                         }\CommentTok{\# 4}
\NormalTok{soma }\OtherTok{\textless{}{-}} \DecValTok{0}\NormalTok{; }\ControlFlowTok{for}\NormalTok{ (n }\ControlFlowTok{in}\NormalTok{ notas) soma }\OtherTok{\textless{}{-}}\NormalTok{ soma }\SpecialCharTok{+}\NormalTok{ n}
\NormalTok{soma }\SpecialCharTok{/} \FunctionTok{length}\NormalTok{(notas)}

\ControlFlowTok{for}\NormalTok{ (num }\ControlFlowTok{in} \FunctionTok{c}\NormalTok{(}\DecValTok{1}\NormalTok{,}\DecValTok{3}\NormalTok{,}\DecValTok{4}\NormalTok{,}\DecValTok{6}\NormalTok{,}\DecValTok{9}\NormalTok{,}\DecValTok{11}\NormalTok{,}\DecValTok{14}\NormalTok{))                        }\CommentTok{\# 5}
  \ControlFlowTok{if}\NormalTok{ (num }\SpecialCharTok{\%\%} \DecValTok{2} \SpecialCharTok{!=} \DecValTok{0}\NormalTok{) }\FunctionTok{print}\NormalTok{(}\FunctionTok{paste}\NormalTok{(num, }\StringTok{"é ímpar"}\NormalTok{))}

\NormalTok{n }\OtherTok{\textless{}{-}} \DecValTok{7}                                                 \CommentTok{\# 6 tabuada}
\ControlFlowTok{for}\NormalTok{ (i }\ControlFlowTok{in} \DecValTok{1}\SpecialCharTok{:}\DecValTok{10}\NormalTok{) }\FunctionTok{print}\NormalTok{(}\FunctionTok{paste}\NormalTok{(n, }\StringTok{"x"}\NormalTok{, i, }\StringTok{"="}\NormalTok{, n }\SpecialCharTok{*}\NormalTok{ i))}
\end{Highlighting}
\end{Shaded}

\textbf{7.} Construir um vetor dentro de um ciclo.

\begin{Shaded}
\begin{Highlighting}[]
\NormalTok{n }\OtherTok{\textless{}{-}} \DecValTok{5}
\NormalTok{v }\OtherTok{\textless{}{-}} \FunctionTok{numeric}\NormalTok{(n)                 }\CommentTok{\# pré{-}alocar é melhor do que crescer com c()}
\ControlFlowTok{for}\NormalTok{ (i }\ControlFlowTok{in} \FunctionTok{seq\_len}\NormalTok{(n)) v[i] }\OtherTok{\textless{}{-}}\NormalTok{ i}
\NormalTok{v}
\end{Highlighting}
\end{Shaded}

Fazer crescer um vetor com \texttt{v\ \textless{}-\ c(v,\ i)} dentro de um ciclo obriga o R a copiar
o vetor inteiro a cada passo, o que se torna lento para muitas iterações.
\textbf{Pré-alocar} com \texttt{numeric(n)} e preencher por índice é substancialmente mais
eficiente. Use \texttt{seq\_len(n)} (ou \texttt{seq\_along(x)}) em vez de \texttt{1:n}, pois estas
lidam corretamente com o caso \texttt{n\ =\ 0}.

\textbf{8. e 9.} Números perfeitos (um número igual à soma dos seus divisores próprios).

\begin{Shaded}
\begin{Highlighting}[]
\CommentTok{\# 8. testar um número}
\NormalTok{n }\OtherTok{\textless{}{-}} \DecValTok{28}
\NormalTok{soma }\OtherTok{\textless{}{-}} \DecValTok{0}
\ControlFlowTok{for}\NormalTok{ (i }\ControlFlowTok{in} \DecValTok{1}\SpecialCharTok{:}\NormalTok{(n }\SpecialCharTok{{-}} \DecValTok{1}\NormalTok{)) }\ControlFlowTok{if}\NormalTok{ (n }\SpecialCharTok{\%\%}\NormalTok{ i }\SpecialCharTok{==} \DecValTok{0}\NormalTok{) soma }\OtherTok{\textless{}{-}}\NormalTok{ soma }\SpecialCharTok{+}\NormalTok{ i}
\ControlFlowTok{if}\NormalTok{ (soma }\SpecialCharTok{==}\NormalTok{ n) }\FunctionTok{print}\NormalTok{(}\FunctionTok{paste}\NormalTok{(n, }\StringTok{"é perfeito"}\NormalTok{)) }\ControlFlowTok{else} \FunctionTok{print}\NormalTok{(}\FunctionTok{paste}\NormalTok{(n, }\StringTok{"não é perfeito"}\NormalTok{))}

\CommentTok{\# 9. todos os perfeitos até 1000 (encapsulado numa função)}
\NormalTok{e\_perfeito }\OtherTok{\textless{}{-}} \ControlFlowTok{function}\NormalTok{(n) \{}
\NormalTok{  soma }\OtherTok{\textless{}{-}} \FunctionTok{sum}\NormalTok{((}\DecValTok{1}\SpecialCharTok{:}\NormalTok{(n }\SpecialCharTok{{-}} \DecValTok{1}\NormalTok{))[n }\SpecialCharTok{\%\%}\NormalTok{ (}\DecValTok{1}\SpecialCharTok{:}\NormalTok{(n }\SpecialCharTok{{-}} \DecValTok{1}\NormalTok{)) }\SpecialCharTok{==} \DecValTok{0}\NormalTok{])}
\NormalTok{  soma }\SpecialCharTok{==}\NormalTok{ n}
\NormalTok{\}}
\ControlFlowTok{for}\NormalTok{ (n }\ControlFlowTok{in} \DecValTok{2}\SpecialCharTok{:}\DecValTok{1000}\NormalTok{) }\ControlFlowTok{if}\NormalTok{ (}\FunctionTok{e\_perfeito}\NormalTok{(n)) }\FunctionTok{print}\NormalTok{(n)   }\CommentTok{\# 6, 28, 496}
\end{Highlighting}
\end{Shaded}

\textbf{10. e 11.} Sequências e médias interativas.

\begin{Shaded}
\begin{Highlighting}[]
\CommentTok{\# 10. imprimir 0, m, 2m, ... até n}
\ControlFlowTok{for}\NormalTok{ (i }\ControlFlowTok{in} \FunctionTok{seq}\NormalTok{(}\DecValTok{0}\NormalTok{, n, }\AttributeTok{by =}\NormalTok{ m)) }\FunctionTok{print}\NormalTok{(i)}

\CommentTok{\# 11. média de n notas lidas do teclado}
\NormalTok{n }\OtherTok{\textless{}{-}} \FunctionTok{as.numeric}\NormalTok{(}\FunctionTok{readline}\NormalTok{(}\StringTok{"Quantas notas? "}\NormalTok{))}
\NormalTok{soma }\OtherTok{\textless{}{-}} \DecValTok{0}
\ControlFlowTok{for}\NormalTok{ (i }\ControlFlowTok{in} \FunctionTok{seq\_len}\NormalTok{(n)) soma }\OtherTok{\textless{}{-}}\NormalTok{ soma }\SpecialCharTok{+} \FunctionTok{as.numeric}\NormalTok{(}\FunctionTok{readline}\NormalTok{(}\StringTok{"Nota: "}\NormalTok{))}
\NormalTok{soma }\SpecialCharTok{/}\NormalTok{ n}
\end{Highlighting}
\end{Shaded}

\textbf{12.} Estatísticas por linha e por coluna de uma matriz.

\begin{Shaded}
\begin{Highlighting}[]
\NormalTok{A }\OtherTok{\textless{}{-}} \FunctionTok{matrix}\NormalTok{(}\DecValTok{1}\SpecialCharTok{:}\DecValTok{16}\NormalTok{, }\AttributeTok{nrow =} \DecValTok{4}\NormalTok{)}

\CommentTok{\# média de cada coluna (com ciclo, para praticar; colMeans faria o mesmo)}
\NormalTok{media\_col }\OtherTok{\textless{}{-}} \FunctionTok{numeric}\NormalTok{(}\FunctionTok{ncol}\NormalTok{(A))}
\ControlFlowTok{for}\NormalTok{ (j }\ControlFlowTok{in} \FunctionTok{seq\_len}\NormalTok{(}\FunctionTok{ncol}\NormalTok{(A))) media\_col[j] }\OtherTok{\textless{}{-}} \FunctionTok{mean}\NormalTok{(A[, j])}

\CommentTok{\# média dos números pares de cada coluna}
\ControlFlowTok{for}\NormalTok{ (j }\ControlFlowTok{in} \FunctionTok{seq\_len}\NormalTok{(}\FunctionTok{ncol}\NormalTok{(A))) \{}
\NormalTok{  pares }\OtherTok{\textless{}{-}}\NormalTok{ A[, j][A[, j] }\SpecialCharTok{\%\%} \DecValTok{2} \SpecialCharTok{==} \DecValTok{0}\NormalTok{]}
  \FunctionTok{cat}\NormalTok{(}\StringTok{"Coluna"}\NormalTok{, j, }\StringTok{"{-} média dos pares:"}\NormalTok{, }\FunctionTok{mean}\NormalTok{(pares), }\StringTok{"}\SpecialCharTok{\textbackslash{}n}\StringTok{"}\NormalTok{)}
\NormalTok{\}}
\end{Highlighting}
\end{Shaded}

\section{\texorpdfstring{Ciclos \texttt{while}}{Ciclos while}}\label{ciclos-while}

\textbf{1. a 5.} Contadores e acumuladores controlados por condição.

\begin{Shaded}
\begin{Highlighting}[]
\CommentTok{\# 1. imprimir de 0 até k}
\NormalTok{k }\OtherTok{\textless{}{-}} \DecValTok{5}\NormalTok{; contador }\OtherTok{\textless{}{-}} \DecValTok{0}
\ControlFlowTok{while}\NormalTok{ (contador }\SpecialCharTok{\textless{}=}\NormalTok{ k) \{ }\FunctionTok{print}\NormalTok{(contador); contador }\OtherTok{\textless{}{-}}\NormalTok{ contador }\SpecialCharTok{+} \DecValTok{1}\NormalTok{ \}}

\CommentTok{\# 3. somatório dos primeiros n inteiros}
\NormalTok{n }\OtherTok{\textless{}{-}} \DecValTok{10}\NormalTok{; soma }\OtherTok{\textless{}{-}} \DecValTok{0}\NormalTok{; i }\OtherTok{\textless{}{-}} \DecValTok{1}
\ControlFlowTok{while}\NormalTok{ (i }\SpecialCharTok{\textless{}=}\NormalTok{ n) \{ soma }\OtherTok{\textless{}{-}}\NormalTok{ soma }\SpecialCharTok{+}\NormalTok{ i; i }\OtherTok{\textless{}{-}}\NormalTok{ i }\SpecialCharTok{+} \DecValTok{1}\NormalTok{ \}}
\NormalTok{soma                       }\CommentTok{\# n(n+1)/2 = 55}

\CommentTok{\# 5. somar números até introduzir um negativo}
\NormalTok{soma }\OtherTok{\textless{}{-}} \DecValTok{0}
\NormalTok{n }\OtherTok{\textless{}{-}} \FunctionTok{as.numeric}\NormalTok{(}\FunctionTok{readline}\NormalTok{(}\StringTok{"Número (negativo para sair): "}\NormalTok{))}
\ControlFlowTok{while}\NormalTok{ (n }\SpecialCharTok{\textgreater{}=} \DecValTok{0}\NormalTok{) \{ soma }\OtherTok{\textless{}{-}}\NormalTok{ soma }\SpecialCharTok{+}\NormalTok{ n; n }\OtherTok{\textless{}{-}} \FunctionTok{as.numeric}\NormalTok{(}\FunctionTok{readline}\NormalTok{(}\StringTok{"Número: "}\NormalTok{)) \}}
\NormalTok{soma}
\end{Highlighting}
\end{Shaded}

\textbf{9.} Fatorial.

\begin{Shaded}
\begin{Highlighting}[]
\NormalTok{n }\OtherTok{\textless{}{-}} \DecValTok{6}\NormalTok{; fat }\OtherTok{\textless{}{-}} \DecValTok{1}\NormalTok{; i }\OtherTok{\textless{}{-}} \DecValTok{1}
\ControlFlowTok{while}\NormalTok{ (i }\SpecialCharTok{\textless{}=}\NormalTok{ n) \{ fat }\OtherTok{\textless{}{-}}\NormalTok{ fat }\SpecialCharTok{*}\NormalTok{ i; i }\OtherTok{\textless{}{-}}\NormalTok{ i }\SpecialCharTok{+} \DecValTok{1}\NormalTok{ \}}
\NormalTok{fat                        }\CommentTok{\# 720  (compare com factorial(6))}
\end{Highlighting}
\end{Shaded}

\textbf{10.} Soma dos algarismos de um número.

\begin{Shaded}
\begin{Highlighting}[]
\NormalTok{n }\OtherTok{\textless{}{-}} \DecValTok{1234}\NormalTok{; soma }\OtherTok{\textless{}{-}} \DecValTok{0}
\ControlFlowTok{while}\NormalTok{ (n }\SpecialCharTok{\textgreater{}} \DecValTok{0}\NormalTok{) \{ soma }\OtherTok{\textless{}{-}}\NormalTok{ soma }\SpecialCharTok{+}\NormalTok{ n }\SpecialCharTok{\%\%} \DecValTok{10}\NormalTok{; n }\OtherTok{\textless{}{-}}\NormalTok{ n }\SpecialCharTok{\%/\%} \DecValTok{10}\NormalTok{ \}}
\NormalTok{soma                       }\CommentTok{\# 1+2+3+4 = 10}
\end{Highlighting}
\end{Shaded}

\textbf{11.} Sucessão de Fibonacci até um valor máximo (uso de \texttt{break}).

\begin{Shaded}
\begin{Highlighting}[]
\NormalTok{valor\_max }\OtherTok{\textless{}{-}} \DecValTok{100}
\NormalTok{f1 }\OtherTok{\textless{}{-}} \DecValTok{1}\NormalTok{; f2 }\OtherTok{\textless{}{-}} \DecValTok{1}
\FunctionTok{print}\NormalTok{(f1); }\ControlFlowTok{if}\NormalTok{ (valor\_max }\SpecialCharTok{\textgreater{}} \DecValTok{1}\NormalTok{) }\FunctionTok{print}\NormalTok{(f2)}
\ControlFlowTok{repeat}\NormalTok{ \{}
\NormalTok{  f }\OtherTok{\textless{}{-}}\NormalTok{ f1 }\SpecialCharTok{+}\NormalTok{ f2}
  \ControlFlowTok{if}\NormalTok{ (f }\SpecialCharTok{\textgreater{}}\NormalTok{ valor\_max) }\ControlFlowTok{break}
  \FunctionTok{print}\NormalTok{(f)}
\NormalTok{  f1 }\OtherTok{\textless{}{-}}\NormalTok{ f2; f2 }\OtherTok{\textless{}{-}}\NormalTok{ f}
\NormalTok{\}}
\end{Highlighting}
\end{Shaded}

\texttt{while\ (TRUE)\ \{\ ...\ if\ (cond)\ break\ \}} e \texttt{repeat\ \{\ ...\ if\ (cond)\ break\ \}} são
equivalentes; \texttt{repeat} exprime de forma mais direta a ideia de ``repetir até uma
condição de paragem''. Em qualquer ciclo indefinido, garanta sempre que existe um
caminho que conduz ao \texttt{break} --- caso contrário o ciclo nunca termina.

\section{Funções}\label{funuxe7uxf5es}

\textbf{1. e 2.} Funções elementares.

\begin{Shaded}
\begin{Highlighting}[]
\NormalTok{quadrado }\OtherTok{\textless{}{-}} \ControlFlowTok{function}\NormalTok{(x) x}\SpecialCharTok{\^{}}\DecValTok{2}
\NormalTok{hipotenusa }\OtherTok{\textless{}{-}} \ControlFlowTok{function}\NormalTok{(a, b) }\FunctionTok{sqrt}\NormalTok{(a}\SpecialCharTok{\^{}}\DecValTok{2} \SpecialCharTok{+}\NormalTok{ b}\SpecialCharTok{\^{}}\DecValTok{2}\NormalTok{)}
\end{Highlighting}
\end{Shaded}

\textbf{3.} Divisão segura (função completa --- o material original omitia o cabeçalho).

\begin{Shaded}
\begin{Highlighting}[]
\NormalTok{divisao\_segura }\OtherTok{\textless{}{-}} \ControlFlowTok{function}\NormalTok{(numerador, denominador) \{}
  \ControlFlowTok{if}\NormalTok{ (denominador }\SpecialCharTok{==} \DecValTok{0}\NormalTok{) }\FunctionTok{return}\NormalTok{(}\StringTok{"Divisão por zero!"}\NormalTok{)}
\NormalTok{  numerador }\SpecialCharTok{/}\NormalTok{ denominador}
\NormalTok{\}}
\FunctionTok{divisao\_segura}\NormalTok{(}\DecValTok{10}\NormalTok{, }\DecValTok{2}\NormalTok{)   }\CommentTok{\# 5}
\FunctionTok{divisao\_segura}\NormalTok{(}\DecValTok{10}\NormalTok{, }\DecValTok{0}\NormalTok{)   }\CommentTok{\# "Divisão por zero!"}
\end{Highlighting}
\end{Shaded}

\textbf{4. e 5.} Compor funções (maior e menor de dois ou três números).

\begin{Shaded}
\begin{Highlighting}[]
\NormalTok{maior }\OtherTok{\textless{}{-}} \ControlFlowTok{function}\NormalTok{(x, y) }\ControlFlowTok{if}\NormalTok{ (x }\SpecialCharTok{\textgreater{}}\NormalTok{ y) x }\ControlFlowTok{else}\NormalTok{ y}
\NormalTok{menor }\OtherTok{\textless{}{-}} \ControlFlowTok{function}\NormalTok{(x, y) }\ControlFlowTok{if}\NormalTok{ (x }\SpecialCharTok{\textless{}}\NormalTok{ y) x }\ControlFlowTok{else}\NormalTok{ y}
\NormalTok{maior3 }\OtherTok{\textless{}{-}} \ControlFlowTok{function}\NormalTok{(x, y, z) }\FunctionTok{maior}\NormalTok{(}\FunctionTok{maior}\NormalTok{(x, y), z)}
\end{Highlighting}
\end{Shaded}

\textbf{6.} Função que devolve várias estatísticas.

\begin{Shaded}
\begin{Highlighting}[]
\NormalTok{analise\_vetor }\OtherTok{\textless{}{-}} \ControlFlowTok{function}\NormalTok{(x) \{}
  \FunctionTok{c}\NormalTok{(}\AttributeTok{soma =} \FunctionTok{sum}\NormalTok{(x), }\AttributeTok{media =} \FunctionTok{mean}\NormalTok{(x), }\AttributeTok{dp =} \FunctionTok{sd}\NormalTok{(x), }\AttributeTok{maximo =} \FunctionTok{max}\NormalTok{(x))}
\NormalTok{\}}
\FunctionTok{analise\_vetor}\NormalTok{(}\FunctionTok{c}\NormalTok{(}\DecValTok{4}\NormalTok{, }\DecValTok{8}\NormalTok{, }\DecValTok{15}\NormalTok{, }\DecValTok{16}\NormalTok{, }\DecValTok{23}\NormalTok{, }\DecValTok{42}\NormalTok{))}
\end{Highlighting}
\end{Shaded}

\textbf{7. e 8.} Classificação e índice de massa corporal.

\begin{Shaded}
\begin{Highlighting}[]
\NormalTok{classificar }\OtherTok{\textless{}{-}} \ControlFlowTok{function}\NormalTok{(x) \{}
  \ControlFlowTok{if}\NormalTok{ (x }\SpecialCharTok{==} \DecValTok{0}\NormalTok{) }\StringTok{"zero"} \ControlFlowTok{else} \ControlFlowTok{if}\NormalTok{ (x }\SpecialCharTok{\textgreater{}} \DecValTok{0}\NormalTok{) }\StringTok{"positivo"} \ControlFlowTok{else} \StringTok{"negativo"}
\NormalTok{\}}

\NormalTok{imc }\OtherTok{\textless{}{-}} \ControlFlowTok{function}\NormalTok{(peso, altura) peso }\SpecialCharTok{/}\NormalTok{ altura}\SpecialCharTok{\^{}}\DecValTok{2}
\CommentTok{\# a categoria pode obter{-}se com cut(), evitando uma longa cadeia de if/else:}
\NormalTok{categoria\_imc }\OtherTok{\textless{}{-}} \ControlFlowTok{function}\NormalTok{(v) \{}
  \FunctionTok{cut}\NormalTok{(v, }\AttributeTok{breaks =} \FunctionTok{c}\NormalTok{(}\SpecialCharTok{{-}}\ConstantTok{Inf}\NormalTok{, }\FloatTok{18.5}\NormalTok{, }\DecValTok{25}\NormalTok{, }\DecValTok{30}\NormalTok{, }\DecValTok{35}\NormalTok{, }\DecValTok{40}\NormalTok{, }\ConstantTok{Inf}\NormalTok{),}
      \AttributeTok{labels =} \FunctionTok{c}\NormalTok{(}\StringTok{"Baixo peso"}\NormalTok{,}\StringTok{"Normal"}\NormalTok{,}\StringTok{"Excesso de peso"}\NormalTok{,}
                 \StringTok{"Obesidade I"}\NormalTok{,}\StringTok{"Obesidade II"}\NormalTok{,}\StringTok{"Obesidade III"}\NormalTok{),}
      \AttributeTok{right =} \ConstantTok{FALSE}\NormalTok{)}
\NormalTok{\}}
\FunctionTok{categoria\_imc}\NormalTok{(}\FunctionTok{imc}\NormalTok{(}\DecValTok{70}\NormalTok{, }\FloatTok{1.75}\NormalTok{))}
\end{Highlighting}
\end{Shaded}

\textbf{9.} Conversor de moedas.

\begin{Shaded}
\begin{Highlighting}[]
\NormalTok{converter }\OtherTok{\textless{}{-}} \ControlFlowTok{function}\NormalTok{(euros) \{}
  \FunctionTok{c}\NormalTok{(}\AttributeTok{dolar =}\NormalTok{ euros }\SpecialCharTok{*} \FloatTok{1.10}\NormalTok{, }\AttributeTok{libra =}\NormalTok{ euros }\SpecialCharTok{*} \FloatTok{0.85}\NormalTok{,}
    \AttributeTok{iene  =}\NormalTok{ euros }\SpecialCharTok{*} \DecValTok{130}\NormalTok{,  }\AttributeTok{real  =}\NormalTok{ euros }\SpecialCharTok{*} \FloatTok{5.20}\NormalTok{)}
\NormalTok{\}}
\FunctionTok{converter}\NormalTok{(}\DecValTok{100}\NormalTok{)}
\end{Highlighting}
\end{Shaded}

\textbf{10. e 11.} Intervalo de confiança para a média.

\begin{Shaded}
\begin{Highlighting}[]
\NormalTok{IC }\OtherTok{\textless{}{-}} \ControlFlowTok{function}\NormalTok{(amostra, alpha) \{}
\NormalTok{  media }\OtherTok{\textless{}{-}} \FunctionTok{mean}\NormalTok{(amostra); n }\OtherTok{\textless{}{-}} \FunctionTok{length}\NormalTok{(amostra)}
\NormalTok{  erro  }\OtherTok{\textless{}{-}} \FunctionTok{qnorm}\NormalTok{(}\DecValTok{1} \SpecialCharTok{{-}}\NormalTok{ alpha }\SpecialCharTok{/} \DecValTok{2}\NormalTok{) }\SpecialCharTok{*} \FunctionTok{sd}\NormalTok{(amostra) }\SpecialCharTok{/} \FunctionTok{sqrt}\NormalTok{(n)}
  \FunctionTok{c}\NormalTok{(}\AttributeTok{inferior =}\NormalTok{ media }\SpecialCharTok{{-}}\NormalTok{ erro, }\AttributeTok{superior =}\NormalTok{ media }\SpecialCharTok{+}\NormalTok{ erro)}
\NormalTok{\}}
\FunctionTok{IC}\NormalTok{(}\DecValTok{1}\SpecialCharTok{:}\DecValTok{15}\NormalTok{, }\FloatTok{0.05}\NormalTok{)}

\CommentTok{\# 11. vários níveis de confiança de uma só vez (vetorizado)}
\NormalTok{ICs }\OtherTok{\textless{}{-}} \ControlFlowTok{function}\NormalTok{(amostra, alphas) \{}
\NormalTok{  media }\OtherTok{\textless{}{-}} \FunctionTok{mean}\NormalTok{(amostra); n }\OtherTok{\textless{}{-}} \FunctionTok{length}\NormalTok{(amostra)}
\NormalTok{  ep }\OtherTok{\textless{}{-}} \FunctionTok{sd}\NormalTok{(amostra) }\SpecialCharTok{/} \FunctionTok{sqrt}\NormalTok{(n)}
\NormalTok{  z  }\OtherTok{\textless{}{-}} \FunctionTok{qnorm}\NormalTok{(}\DecValTok{1} \SpecialCharTok{{-}}\NormalTok{ alphas }\SpecialCharTok{/} \DecValTok{2}\NormalTok{)}
  \FunctionTok{data.frame}\NormalTok{(}\AttributeTok{Alpha =}\NormalTok{ alphas,}
             \AttributeTok{Inferior =}\NormalTok{ media }\SpecialCharTok{{-}}\NormalTok{ z }\SpecialCharTok{*}\NormalTok{ ep,}
             \AttributeTok{Superior =}\NormalTok{ media }\SpecialCharTok{+}\NormalTok{ z }\SpecialCharTok{*}\NormalTok{ ep)}
\NormalTok{\}}
\FunctionTok{set.seed}\NormalTok{(}\DecValTok{123}\NormalTok{)}
\FunctionTok{ICs}\NormalTok{(}\FunctionTok{rnorm}\NormalTok{(}\DecValTok{100}\NormalTok{, }\DecValTok{50}\NormalTok{, }\DecValTok{10}\NormalTok{), }\FunctionTok{c}\NormalTok{(}\FloatTok{0.001}\NormalTok{, }\FloatTok{0.01}\NormalTok{, }\FloatTok{0.05}\NormalTok{, }\FloatTok{0.10}\NormalTok{))}
\end{Highlighting}
\end{Shaded}

\textbf{12.} Moda (valor mais frequente).

\begin{Shaded}
\begin{Highlighting}[]
\NormalTok{calcula\_moda }\OtherTok{\textless{}{-}} \ControlFlowTok{function}\NormalTok{(x) \{}
\NormalTok{  fr }\OtherTok{\textless{}{-}} \FunctionTok{table}\NormalTok{(x)}
\NormalTok{  modas }\OtherTok{\textless{}{-}} \FunctionTok{names}\NormalTok{(fr)[fr }\SpecialCharTok{==} \FunctionTok{max}\NormalTok{(fr)]}
  \ControlFlowTok{if}\NormalTok{ (}\FunctionTok{is.numeric}\NormalTok{(x)) }\FunctionTok{as.numeric}\NormalTok{(modas) }\ControlFlowTok{else}\NormalTok{ modas}
\NormalTok{\}}
\FunctionTok{calcula\_moda}\NormalTok{(}\FunctionTok{c}\NormalTok{(}\DecValTok{1}\NormalTok{, }\DecValTok{2}\NormalTok{, }\DecValTok{2}\NormalTok{, }\DecValTok{3}\NormalTok{, }\DecValTok{3}\NormalTok{, }\DecValTok{4}\NormalTok{, }\DecValTok{4}\NormalTok{))     }\CommentTok{\# 2 3 4 (distribuição multimodal)}
\end{Highlighting}
\end{Shaded}

\textbf{13. e 14.} Distância euclidiana e análise de matrizes.

\begin{Shaded}
\begin{Highlighting}[]
\NormalTok{distancia }\OtherTok{\textless{}{-}} \ControlFlowTok{function}\NormalTok{(x, y) }\FunctionTok{sqrt}\NormalTok{(}\FunctionTok{sum}\NormalTok{((x }\SpecialCharTok{{-}}\NormalTok{ y)}\SpecialCharTok{\^{}}\DecValTok{2}\NormalTok{))}
\FunctionTok{distancia}\NormalTok{(}\FunctionTok{c}\NormalTok{(}\DecValTok{2}\NormalTok{, }\DecValTok{3}\NormalTok{), }\FunctionTok{c}\NormalTok{(}\DecValTok{6}\NormalTok{, }\DecValTok{9}\NormalTok{))              }\CommentTok{\# 7.211}

\NormalTok{analisa\_matriz }\OtherTok{\textless{}{-}} \ControlFlowTok{function}\NormalTok{(A) \{}
\NormalTok{  res }\OtherTok{\textless{}{-}} \FunctionTok{list}\NormalTok{(}\AttributeTok{soma =} \FunctionTok{sum}\NormalTok{(A), }\AttributeTok{media =} \FunctionTok{mean}\NormalTok{(A))}
  \ControlFlowTok{if}\NormalTok{ (}\FunctionTok{nrow}\NormalTok{(A) }\SpecialCharTok{==} \FunctionTok{ncol}\NormalTok{(A)) \{}
\NormalTok{    res}\SpecialCharTok{$}\NormalTok{traco }\OtherTok{\textless{}{-}} \FunctionTok{sum}\NormalTok{(}\FunctionTok{diag}\NormalTok{(A)); res}\SpecialCharTok{$}\NormalTok{det }\OtherTok{\textless{}{-}} \FunctionTok{det}\NormalTok{(A)}
\NormalTok{  \}}
\NormalTok{  res}
\NormalTok{\}}
\end{Highlighting}
\end{Shaded}

\section{\texorpdfstring{\emph{Scripts} e a função \texttt{source()}}{Scripts e a função source()}}\label{scripts-e-a-funuxe7uxe3o-source}

Estes exercícios exercitam a organização do código em ficheiros separados,
carregados com \texttt{source()}. Cada \emph{script} define funções que depois se usam. Por
exemplo, um ficheiro \texttt{estatistica.R} com

\begin{Shaded}
\begin{Highlighting}[]
\CommentTok{\# ficheiro estatistica.R}
\NormalTok{ordena     }\OtherTok{\textless{}{-}} \ControlFlowTok{function}\NormalTok{(x) }\FunctionTok{sort}\NormalTok{(x)}
\NormalTok{media      }\OtherTok{\textless{}{-}} \ControlFlowTok{function}\NormalTok{(x) }\FunctionTok{mean}\NormalTok{(x)}
\NormalTok{variancia  }\OtherTok{\textless{}{-}} \ControlFlowTok{function}\NormalTok{(x) }\FunctionTok{var}\NormalTok{(x)}
\NormalTok{desv\_padrao }\OtherTok{\textless{}{-}} \ControlFlowTok{function}\NormalTok{(x) }\FunctionTok{sd}\NormalTok{(x)}
\end{Highlighting}
\end{Shaded}

é carregado e usado assim:

\begin{Shaded}
\begin{Highlighting}[]
\FunctionTok{source}\NormalTok{(}\StringTok{"estatistica.R"}\NormalTok{)            }\CommentTok{\# ajuste o caminho ao seu ficheiro}
\NormalTok{amostra }\OtherTok{\textless{}{-}} \FunctionTok{c}\NormalTok{(}\DecValTok{1}\NormalTok{, }\DecValTok{2}\NormalTok{, }\DecValTok{3}\NormalTok{, }\DecValTok{4}\NormalTok{, }\DecValTok{5}\NormalTok{, }\DecValTok{6}\NormalTok{, }\DecValTok{7}\NormalTok{)}
\FunctionTok{variancia}\NormalTok{(amostra)}
\FunctionTok{desv\_padrao}\NormalTok{(amostra)}
\end{Highlighting}
\end{Shaded}

De igual modo se organizam os \emph{scripts} \texttt{quadrado.R} (perímetro e área do
quadrado), \texttt{circulo.R} (área e perímetro do círculo), \texttt{analise.R} (mínimo,
máximo, amplitude), \texttt{categorias.R} (contagem de categorias), \texttt{escore\_z.R}
(cálculo do valor padronizado \(z\)) e \texttt{binomial.R}, este último com a função de
probabilidade da binomial calculada de raiz:

\begin{Shaded}
\begin{Highlighting}[]
\CommentTok{\# ficheiro binomial.R}
\NormalTok{probabilidade\_binomial }\OtherTok{\textless{}{-}} \ControlFlowTok{function}\NormalTok{(n, p, x) }\FunctionTok{choose}\NormalTok{(n, x) }\SpecialCharTok{*}\NormalTok{ p}\SpecialCharTok{\^{}}\NormalTok{x }\SpecialCharTok{*}\NormalTok{ (}\DecValTok{1} \SpecialCharTok{{-}}\NormalTok{ p)}\SpecialCharTok{\^{}}\NormalTok{(n }\SpecialCharTok{{-}}\NormalTok{ x)}

\CommentTok{\# uso, comparando com a função interna do R}
\FunctionTok{source}\NormalTok{(}\StringTok{"binomial.R"}\NormalTok{)}
\FunctionTok{probabilidade\_binomial}\NormalTok{(}\DecValTok{5}\NormalTok{, }\FloatTok{0.6}\NormalTok{, }\DecValTok{3}\NormalTok{)}
\FunctionTok{dbinom}\NormalTok{(}\AttributeTok{x =} \DecValTok{3}\NormalTok{, }\AttributeTok{size =} \DecValTok{5}\NormalTok{, }\AttributeTok{prob =} \FloatTok{0.6}\NormalTok{)   }\CommentTok{\# deve coincidir}
\end{Highlighting}
\end{Shaded}

Separar funções reutilizáveis em \emph{scripts} próprios e carregá-las com \texttt{source()}
mantém o código de análise limpo e as funções disponíveis para vários projetos.
É o primeiro passo, informal, na direção da criação de \emph{pacotes} de R.

\section{\texorpdfstring{A família \texttt{apply}}{A família apply}}\label{a-famuxedlia-apply}

\textbf{1. a 5.} \texttt{lapply}, \texttt{sapply}, \texttt{apply} e \texttt{tapply}.

\begin{Shaded}
\begin{Highlighting}[]
\FunctionTok{sapply}\NormalTok{(}\FunctionTok{c}\NormalTok{(}\StringTok{"cachorro"}\NormalTok{,}\StringTok{"gato"}\NormalTok{,}\StringTok{"elefante"}\NormalTok{,}\StringTok{"leão"}\NormalTok{), nchar)   }\CommentTok{\# 2}

\NormalTok{M }\OtherTok{\textless{}{-}} \FunctionTok{matrix}\NormalTok{(}\DecValTok{1}\SpecialCharTok{:}\DecValTok{9}\NormalTok{, }\AttributeTok{nrow =} \DecValTok{3}\NormalTok{)}
\FunctionTok{apply}\NormalTok{(M, }\DecValTok{1}\NormalTok{, sum)     }\CommentTok{\# soma por linha (MARGIN = 1)}
\FunctionTok{apply}\NormalTok{(M, }\DecValTok{2}\NormalTok{, mean)    }\CommentTok{\# média por coluna (MARGIN = 2)}

\FunctionTok{tapply}\NormalTok{(mtcars}\SpecialCharTok{$}\NormalTok{mpg, mtcars}\SpecialCharTok{$}\NormalTok{cyl, mean)   }\CommentTok{\# 4 média de mpg por nº de cilindros}
\end{Highlighting}
\end{Shaded}

A diferença entre \texttt{lapply()} e \texttt{sapply()} é apenas o formato do resultado:
\texttt{lapply()} devolve sempre uma \textbf{lista}; \texttt{sapply()} tenta \textbf{simplificar} para
vetor ou matriz quando possível. Para o argumento \texttt{MARGIN} de \texttt{apply()}, a regra
é \texttt{1\ =\ linhas}, \texttt{2\ =\ colunas}.

\textbf{6.} Aplicar uma função personalizada a cada elemento.

\begin{Shaded}
\begin{Highlighting}[]
\FunctionTok{sapply}\NormalTok{(}\DecValTok{1}\SpecialCharTok{:}\DecValTok{10}\NormalTok{, }\ControlFlowTok{function}\NormalTok{(x) }\ControlFlowTok{if}\NormalTok{ (x }\SpecialCharTok{\%\%} \DecValTok{2} \SpecialCharTok{==} \DecValTok{0}\NormalTok{) }\StringTok{"par"} \ControlFlowTok{else} \StringTok{"ímpar"}\NormalTok{)}
\end{Highlighting}
\end{Shaded}

\textbf{7. e 8.} Estatísticas por grupo e normalização de colunas.

\begin{Shaded}
\begin{Highlighting}[]
\FunctionTok{tapply}\NormalTok{(iris}\SpecialCharTok{$}\NormalTok{Petal.Length, iris}\SpecialCharTok{$}\NormalTok{Species, mean)}
\FunctionTok{tapply}\NormalTok{(iris}\SpecialCharTok{$}\NormalTok{Sepal.Width,  iris}\SpecialCharTok{$}\NormalTok{Species, var)}

\NormalTok{normaliza }\OtherTok{\textless{}{-}} \ControlFlowTok{function}\NormalTok{(x) (x }\SpecialCharTok{{-}} \FunctionTok{min}\NormalTok{(x)) }\SpecialCharTok{/}\NormalTok{ (}\FunctionTok{max}\NormalTok{(x) }\SpecialCharTok{{-}} \FunctionTok{min}\NormalTok{(x))}
\NormalTok{dados\_norm }\OtherTok{\textless{}{-}} \FunctionTok{as.data.frame}\NormalTok{(}\FunctionTok{apply}\NormalTok{(mtcars, }\DecValTok{2}\NormalTok{, normaliza))}
\end{Highlighting}
\end{Shaded}

\textbf{9. e 10.} Aplicações a matrizes de dados simulados.

\begin{Shaded}
\begin{Highlighting}[]
\FunctionTok{set.seed}\NormalTok{(}\DecValTok{123}\NormalTok{)}
\NormalTok{notas }\OtherTok{\textless{}{-}} \FunctionTok{as.data.frame}\NormalTok{(}\FunctionTok{matrix}\NormalTok{(}\FunctionTok{runif}\NormalTok{(}\DecValTok{30} \SpecialCharTok{*} \DecValTok{4}\NormalTok{, }\DecValTok{0}\NormalTok{, }\DecValTok{100}\NormalTok{), }\AttributeTok{nrow =} \DecValTok{30}\NormalTok{))}
\FunctionTok{colnames}\NormalTok{(notas) }\OtherTok{\textless{}{-}} \FunctionTok{paste0}\NormalTok{(}\StringTok{"Exame"}\NormalTok{, }\DecValTok{1}\SpecialCharTok{:}\DecValTok{4}\NormalTok{)}
\FunctionTok{apply}\NormalTok{(notas, }\DecValTok{1}\NormalTok{, mean)     }\CommentTok{\# média de cada aluno (por linha)}
\FunctionTok{sapply}\NormalTok{(notas, min)        }\CommentTok{\# nota mínima por exame (por coluna)}
\FunctionTok{sapply}\NormalTok{(notas, max)}
\end{Highlighting}
\end{Shaded}

\section{\texorpdfstring{O operador \emph{pipe}}{O operador pipe}}\label{o-operador-pipe}

O \emph{pipe} (\texttt{\%\textgreater{}\%} do \texttt{magrittr}, ou \texttt{\textbar{}\textgreater{}} nativo do R) encadeia operações, passando
o resultado de cada passo como primeiro argumento do seguinte.

\textbf{1. a 3.} Encadeamentos numéricos.

\begin{Shaded}
\begin{Highlighting}[]
\FunctionTok{library}\NormalTok{(magrittr)}
\DecValTok{1}\SpecialCharTok{:}\DecValTok{10} \SpecialCharTok{\%\textgreater{}\%} \FunctionTok{sum}\NormalTok{() }\SpecialCharTok{\%\textgreater{}\%} \FunctionTok{divide\_by}\NormalTok{(}\DecValTok{3}\NormalTok{) }\SpecialCharTok{\%\textgreater{}\%} \FunctionTok{round}\NormalTok{(}\DecValTok{1}\NormalTok{)   }\CommentTok{\# 1}

\FunctionTok{set.seed}\NormalTok{(}\DecValTok{123}\NormalTok{)}
\FunctionTok{rnorm}\NormalTok{(}\DecValTok{100}\NormalTok{) }\SpecialCharTok{\%\textgreater{}\%} \FunctionTok{extract}\NormalTok{(. }\SpecialCharTok{\textgreater{}} \DecValTok{0}\NormalTok{) }\SpecialCharTok{\%\textgreater{}\%} \FunctionTok{mean}\NormalTok{() }\SpecialCharTok{\%\textgreater{}\%} \FunctionTok{round}\NormalTok{(}\DecValTok{2}\NormalTok{)   }\CommentTok{\# 2 (só os positivos)}
\end{Highlighting}
\end{Shaded}

No exercício 3, o encadeamento \texttt{2\ +\ 2} → colocar num vetor com \texttt{6} e \texttt{NA} →
\texttt{mean(na.rm\ =\ TRUE)} → comparar com \texttt{5} devolve \texttt{TRUE}: a média de \texttt{c(4,\ 6)},
ignorando o \texttt{NA}, é de facto \texttt{5}. O \emph{pipe} torna esta leitura sequencial
imediata, lendo-se o código da esquerda para a direita como uma frase.

\textbf{4. a 7.} Fluxos de manipulação de dados com o \texttt{dplyr}.

\begin{Shaded}
\begin{Highlighting}[]
\FunctionTok{library}\NormalTok{(dplyr)}
\NormalTok{mtcars }\SpecialCharTok{\%\textgreater{}\%} \FunctionTok{select}\NormalTok{(mpg, hp, wt) }\SpecialCharTok{\%\textgreater{}\%} \FunctionTok{filter}\NormalTok{(hp }\SpecialCharTok{\textgreater{}} \DecValTok{100}\NormalTok{) }\SpecialCharTok{\%\textgreater{}\%} \FunctionTok{arrange}\NormalTok{(wt)   }\CommentTok{\# 4}

\NormalTok{airquality }\SpecialCharTok{\%\textgreater{}\%} \FunctionTok{na.omit}\NormalTok{() }\SpecialCharTok{\%\textgreater{}\%} \FunctionTok{filter}\NormalTok{(Wind }\SpecialCharTok{\textgreater{}} \DecValTok{10}\NormalTok{) }\SpecialCharTok{\%\textgreater{}\%}                    \CommentTok{\# 6}
  \FunctionTok{summarise}\NormalTok{(}\AttributeTok{ozono\_medio =} \FunctionTok{mean}\NormalTok{(Ozone))}

\NormalTok{mtcars }\SpecialCharTok{\%\textgreater{}\%} \FunctionTok{mutate}\NormalTok{(}\AttributeTok{hp\_wt =}\NormalTok{ hp }\SpecialCharTok{/}\NormalTok{ wt) }\SpecialCharTok{\%\textgreater{}\%} \FunctionTok{filter}\NormalTok{(hp\_wt }\SpecialCharTok{\textgreater{}} \DecValTok{30}\NormalTok{) }\SpecialCharTok{\%\textgreater{}\%}         \CommentTok{\# 7}
  \FunctionTok{select}\NormalTok{(mpg, hp\_wt, gear)}
\end{Highlighting}
\end{Shaded}

\section{Exercícios de revisão}\label{exercuxedcios-de-revisuxe3o}

Estes problemas integradores combinam \emph{data frames}, funções e ciclos.

\textbf{1.} Satisfação de clientes: média por linha e classificação.

\begin{Shaded}
\begin{Highlighting}[]
\FunctionTok{set.seed}\NormalTok{(}\DecValTok{123}\NormalTok{)}
\NormalTok{satisfacao }\OtherTok{\textless{}{-}} \FunctionTok{data.frame}\NormalTok{(}
  \AttributeTok{S1 =} \FunctionTok{sample}\NormalTok{(}\DecValTok{1}\SpecialCharTok{:}\DecValTok{10}\NormalTok{, }\DecValTok{10}\NormalTok{, }\AttributeTok{replace =} \ConstantTok{TRUE}\NormalTok{),}
  \AttributeTok{S2 =} \FunctionTok{sample}\NormalTok{(}\DecValTok{1}\SpecialCharTok{:}\DecValTok{10}\NormalTok{, }\DecValTok{10}\NormalTok{, }\AttributeTok{replace =} \ConstantTok{TRUE}\NormalTok{),}
  \AttributeTok{S3 =} \FunctionTok{sample}\NormalTok{(}\DecValTok{1}\SpecialCharTok{:}\DecValTok{10}\NormalTok{, }\DecValTok{10}\NormalTok{, }\AttributeTok{replace =} \ConstantTok{TRUE}\NormalTok{))}

\NormalTok{satisfacao}\SpecialCharTok{$}\NormalTok{media }\OtherTok{\textless{}{-}} \FunctionTok{rowMeans}\NormalTok{(satisfacao)                       }\CommentTok{\# (b)}
\NormalTok{satisfacao}\SpecialCharTok{$}\NormalTok{grau  }\OtherTok{\textless{}{-}} \FunctionTok{ifelse}\NormalTok{(satisfacao}\SpecialCharTok{$}\NormalTok{media }\SpecialCharTok{\textgreater{}=} \DecValTok{7}\NormalTok{,              }\CommentTok{\# (c)}
                           \StringTok{"Satisfeito"}\NormalTok{, }\StringTok{"Insatisfeito"}\NormalTok{)}
\end{Highlighting}
\end{Shaded}

\textbf{2.} Metas de colaboradores: função que verifica o número de dias.

\begin{Shaded}
\begin{Highlighting}[]
\NormalTok{verifica\_meta }\OtherTok{\textless{}{-}} \ControlFlowTok{function}\NormalTok{(df) \{}
  \CommentTok{\# dias com \textgreater{}= 20 tarefas, por colaborador; meta = pelo menos 3 dias}
  \FunctionTok{ifelse}\NormalTok{(}\FunctionTok{rowSums}\NormalTok{(df }\SpecialCharTok{\textgreater{}=} \DecValTok{20}\NormalTok{) }\SpecialCharTok{\textgreater{}=} \DecValTok{3}\NormalTok{, }\StringTok{"Atingiu Meta"}\NormalTok{, }\StringTok{"Não Atingiu Meta"}\NormalTok{)}
\NormalTok{\}}
\end{Highlighting}
\end{Shaded}

\texttt{rowSums(df\ \textgreater{}=\ 20)} conta, em cada linha, quantas entradas satisfazem a condição,
porque \texttt{df\ \textgreater{}=\ 20} produz uma matriz de \texttt{TRUE}/\texttt{FALSE} (isto é, 1/0). Substitui de
forma elegante um ciclo duplo sobre linhas e colunas.

\textbf{3.} Médias de amostras, interrompendo ao ultrapassar um limite (\texttt{while} + \texttt{break}).

\begin{Shaded}
\begin{Highlighting}[]
\NormalTok{medias }\OtherTok{\textless{}{-}} \FunctionTok{rowMeans}\NormalTok{(amostras)}
\NormalTok{i }\OtherTok{\textless{}{-}} \DecValTok{1}
\ControlFlowTok{while}\NormalTok{ (i }\SpecialCharTok{\textless{}=} \FunctionTok{length}\NormalTok{(medias)) \{}
  \FunctionTok{cat}\NormalTok{(}\StringTok{"Amostra"}\NormalTok{, i, }\StringTok{":"}\NormalTok{, medias[i], }\StringTok{"}\SpecialCharTok{\textbackslash{}n}\StringTok{"}\NormalTok{)}
  \ControlFlowTok{if}\NormalTok{ (medias[i] }\SpecialCharTok{\textgreater{}} \DecValTok{75}\NormalTok{) \{ }\FunctionTok{cat}\NormalTok{(}\StringTok{"Excedeu o limite!}\SpecialCharTok{\textbackslash{}n}\StringTok{"}\NormalTok{); }\ControlFlowTok{break}\NormalTok{ \}}
\NormalTok{  i }\OtherTok{\textless{}{-}}\NormalTok{ i }\SpecialCharTok{+} \DecValTok{1}
\NormalTok{\}}
\end{Highlighting}
\end{Shaded}

\textbf{4.} Lista \texttt{listas\_transformadas} com \(x_j = j^2 + i\).

\begin{Shaded}
\begin{Highlighting}[]
\NormalTok{listas\_transformadas }\OtherTok{\textless{}{-}} \FunctionTok{vector}\NormalTok{(}\StringTok{"list"}\NormalTok{, }\DecValTok{7}\NormalTok{)}
\ControlFlowTok{for}\NormalTok{ (i }\ControlFlowTok{in} \DecValTok{1}\SpecialCharTok{:}\DecValTok{7}\NormalTok{) listas\_transformadas[[i]] }\OtherTok{\textless{}{-}}\NormalTok{ (}\DecValTok{1}\SpecialCharTok{:}\NormalTok{i)}\SpecialCharTok{\^{}}\DecValTok{2} \SpecialCharTok{+}\NormalTok{ i}
\FunctionTok{length}\NormalTok{(listas\_transformadas[[}\DecValTok{7}\NormalTok{]])                  }\CommentTok{\# (b) = 7}

\NormalTok{soma\_elementos }\OtherTok{\textless{}{-}} \ControlFlowTok{function}\NormalTok{(lst) }\FunctionTok{sapply}\NormalTok{(lst, sum)   }\CommentTok{\# (c)}
\FunctionTok{soma\_elementos}\NormalTok{(listas\_transformadas)               }\CommentTok{\# (d)}
\end{Highlighting}
\end{Shaded}

\textbf{5. e 6.} Operações vetorizadas com \texttt{sample()} e \texttt{rnorm()}.

\begin{Shaded}
\begin{Highlighting}[]
\FunctionTok{set.seed}\NormalTok{(}\DecValTok{1}\NormalTok{)}
\NormalTok{x }\OtherTok{\textless{}{-}} \FunctionTok{sample}\NormalTok{(}\DecValTok{0}\SpecialCharTok{:}\DecValTok{999}\NormalTok{, }\DecValTok{60}\NormalTok{, }\AttributeTok{replace =} \ConstantTok{TRUE}\NormalTok{)}
\NormalTok{y }\OtherTok{\textless{}{-}} \FunctionTok{sample}\NormalTok{(}\DecValTok{0}\SpecialCharTok{:}\DecValTok{999}\NormalTok{, }\DecValTok{60}\NormalTok{, }\AttributeTok{replace =} \ConstantTok{TRUE}\NormalTok{)}
\NormalTok{y[y }\SpecialCharTok{\textless{}} \DecValTok{500} \SpecialCharTok{\&}\NormalTok{ x }\SpecialCharTok{\%\%} \DecValTok{3} \SpecialCharTok{==} \DecValTok{0}\NormalTok{]                 }\CommentTok{\# (5b)}
\NormalTok{dif }\OtherTok{\textless{}{-}} \FunctionTok{abs}\NormalTok{(x }\SpecialCharTok{{-}}\NormalTok{ y); }\FunctionTok{c}\NormalTok{(}\AttributeTok{media =} \FunctionTok{mean}\NormalTok{(dif), }\AttributeTok{dp =} \FunctionTok{sd}\NormalTok{(dif))   }\CommentTok{\# (5c{-}d)}

\NormalTok{x }\OtherTok{\textless{}{-}} \FunctionTok{rnorm}\NormalTok{(}\DecValTok{80}\NormalTok{, }\AttributeTok{mean =} \DecValTok{5}\NormalTok{, }\AttributeTok{sd =} \DecValTok{2}\NormalTok{)         }\CommentTok{\# 6}
\NormalTok{xbar }\OtherTok{\textless{}{-}} \FunctionTok{mean}\NormalTok{(x)}
\NormalTok{resultado }\OtherTok{\textless{}{-}} \FunctionTok{sqrt}\NormalTok{(}\FunctionTok{abs}\NormalTok{(x}\SpecialCharTok{\^{}}\DecValTok{3} \SpecialCharTok{{-}}\NormalTok{ xbar}\SpecialCharTok{\^{}}\DecValTok{2}\NormalTok{))}
\FunctionTok{c}\NormalTok{(}\AttributeTok{max =} \FunctionTok{max}\NormalTok{(resultado), }\AttributeTok{min =} \FunctionTok{min}\NormalTok{(resultado), }\AttributeTok{media =} \FunctionTok{mean}\NormalTok{(resultado))}
\end{Highlighting}
\end{Shaded}

\textbf{7.} Temperaturas semanais de 12 cidades (função + \texttt{while}).

\begin{Shaded}
\begin{Highlighting}[]
\NormalTok{media\_semanal }\OtherTok{\textless{}{-}} \FunctionTok{rowMeans}\NormalTok{(temperaturas\_cidades)}
\NormalTok{classificacao }\OtherTok{\textless{}{-}} \FunctionTok{ifelse}\NormalTok{(media\_semanal }\SpecialCharTok{\textgreater{}} \DecValTok{25}\NormalTok{, }\StringTok{"Acima de 25ºC"}\NormalTok{, }\StringTok{"Abaixo de 25ºC"}\NormalTok{)}
\NormalTok{temperaturas\_cidades}\SpecialCharTok{$}\NormalTok{media }\OtherTok{\textless{}{-}}\NormalTok{ media\_semanal}
\NormalTok{temperaturas\_cidades}\SpecialCharTok{$}\NormalTok{classe }\OtherTok{\textless{}{-}}\NormalTok{ classificacao}
\end{Highlighting}
\end{Shaded}

\textbf{8.} \emph{Bit flips} na memória (simulação com a binomial/Poisson). A modelação do
número de \emph{bits} alterados num bloco de memória por unidade de tempo faz-se
naturalmente com uma binomial (muitos \emph{bits}, cada um com probabilidade ínfima
de se alterar) ou, no limite, com a Poisson. A resolução detalhada segue o mesmo
padrão dos exercícios de simulação da secção seguinte: gerar, contar, comparar
com a probabilidade teórica.

\section{\texorpdfstring{Amostragem com \texttt{sample()}}{Amostragem com sample()}}\label{amostragem-com-sample}

\textbf{1. a 6.} Amostragem com e sem reposição, com e sem probabilidades.

\begin{Shaded}
\begin{Highlighting}[]
\FunctionTok{sample}\NormalTok{(}\DecValTok{1}\SpecialCharTok{:}\DecValTok{20}\NormalTok{, }\AttributeTok{size =} \DecValTok{5}\NormalTok{, }\AttributeTok{replace =} \ConstantTok{FALSE}\NormalTok{)   }\CommentTok{\# 1 sem reposição}
\FunctionTok{sample}\NormalTok{(}\DecValTok{1}\SpecialCharTok{:}\DecValTok{10}\NormalTok{, }\AttributeTok{size =} \DecValTok{10}\NormalTok{, }\AttributeTok{replace =} \ConstantTok{TRUE}\NormalTok{)   }\CommentTok{\# 2 com reposição}
\FunctionTok{sample}\NormalTok{(}\FunctionTok{c}\NormalTok{(}\StringTok{"A"}\NormalTok{,}\StringTok{"B"}\NormalTok{,}\StringTok{"C"}\NormalTok{,}\StringTok{"D"}\NormalTok{,}\StringTok{"E"}\NormalTok{), }\AttributeTok{size =} \DecValTok{3}\NormalTok{,  }\CommentTok{\# 3 probabilidades desiguais}
       \AttributeTok{prob =} \FunctionTok{c}\NormalTok{(}\FloatTok{0.1}\NormalTok{, }\FloatTok{0.2}\NormalTok{, }\FloatTok{0.3}\NormalTok{, }\FloatTok{0.25}\NormalTok{, }\FloatTok{0.15}\NormalTok{))}
\FunctionTok{sample}\NormalTok{(}\DecValTok{1}\SpecialCharTok{:}\DecValTok{10}\NormalTok{)                              }\CommentTok{\# 4 permutação (size e replace por defeito)}
\FunctionTok{sample}\NormalTok{(}\DecValTok{1}\SpecialCharTok{:}\DecValTok{1000}\NormalTok{, }\AttributeTok{size =} \DecValTok{50}\NormalTok{)                 }\CommentTok{\# 6 amostra simples de uma população}
\end{Highlighting}
\end{Shaded}

\textbf{7.} Soma de dois dados (10 000 lançamentos) --- duas formas.

\begin{Shaded}
\begin{Highlighting}[]
\FunctionTok{set.seed}\NormalTok{(}\DecValTok{1}\NormalTok{)}
\NormalTok{d1 }\OtherTok{\textless{}{-}} \FunctionTok{sample}\NormalTok{(}\DecValTok{1}\SpecialCharTok{:}\DecValTok{6}\NormalTok{, }\DecValTok{10000}\NormalTok{, }\AttributeTok{replace =} \ConstantTok{TRUE}\NormalTok{)}
\NormalTok{d2 }\OtherTok{\textless{}{-}} \FunctionTok{sample}\NormalTok{(}\DecValTok{1}\SpecialCharTok{:}\DecValTok{6}\NormalTok{, }\DecValTok{10000}\NormalTok{, }\AttributeTok{replace =} \ConstantTok{TRUE}\NormalTok{)}
\FunctionTok{hist}\NormalTok{(d1 }\SpecialCharTok{+}\NormalTok{ d2, }\AttributeTok{breaks =} \DecValTok{0}\SpecialCharTok{:}\DecValTok{12}\NormalTok{, }\AttributeTok{col =} \StringTok{"lightblue"}\NormalTok{)   }\CommentTok{\# distribuição triangular em 7}
\end{Highlighting}
\end{Shaded}

\textbf{8.} Amostragem estratificada.

\begin{Shaded}
\begin{Highlighting}[]
\FunctionTok{set.seed}\NormalTok{(}\DecValTok{123}\NormalTok{)}
\NormalTok{turmas }\OtherTok{\textless{}{-}} \FunctionTok{c}\NormalTok{(}\FunctionTok{rep}\NormalTok{(}\StringTok{"A"}\NormalTok{, }\DecValTok{50}\NormalTok{), }\FunctionTok{rep}\NormalTok{(}\StringTok{"B"}\NormalTok{, }\DecValTok{100}\NormalTok{), }\FunctionTok{rep}\NormalTok{(}\StringTok{"C"}\NormalTok{, }\DecValTok{50}\NormalTok{))}
\CommentTok{\# amostrar proporcionalmente: 5 de A, 10 de B, 5 de C}
\NormalTok{idx }\OtherTok{\textless{}{-}} \FunctionTok{c}\NormalTok{(}\FunctionTok{sample}\NormalTok{(}\FunctionTok{which}\NormalTok{(turmas }\SpecialCharTok{==} \StringTok{"A"}\NormalTok{), }\DecValTok{5}\NormalTok{),}
         \FunctionTok{sample}\NormalTok{(}\FunctionTok{which}\NormalTok{(turmas }\SpecialCharTok{==} \StringTok{"B"}\NormalTok{), }\DecValTok{10}\NormalTok{),}
         \FunctionTok{sample}\NormalTok{(}\FunctionTok{which}\NormalTok{(turmas }\SpecialCharTok{==} \StringTok{"C"}\NormalTok{), }\DecValTok{5}\NormalTok{))}
\NormalTok{turmas[idx]}
\end{Highlighting}
\end{Shaded}

\textbf{9. e 10.} Cartões de bingo e afetação a grupos.

\begin{Shaded}
\begin{Highlighting}[]
\FunctionTok{replicate}\NormalTok{(}\DecValTok{5}\NormalTok{, }\FunctionTok{sample}\NormalTok{(}\DecValTok{1}\SpecialCharTok{:}\DecValTok{75}\NormalTok{, }\DecValTok{24}\NormalTok{, }\AttributeTok{replace =} \ConstantTok{FALSE}\NormalTok{))   }\CommentTok{\# 9 cinco cartões}

\FunctionTok{set.seed}\NormalTok{(}\DecValTok{1}\NormalTok{)                                        }\CommentTok{\# 10 tratamento vs. controlo}
\NormalTok{grupos }\OtherTok{\textless{}{-}} \FunctionTok{sample}\NormalTok{(}\FunctionTok{c}\NormalTok{(}\FunctionTok{rep}\NormalTok{(}\StringTok{"Tratamento"}\NormalTok{, }\DecValTok{20}\NormalTok{), }\FunctionTok{rep}\NormalTok{(}\StringTok{"Controlo"}\NormalTok{, }\DecValTok{10}\NormalTok{)))}
\end{Highlighting}
\end{Shaded}

Para dividir aleatoriamente unidades em grupos de tamanhos fixos, baralhar um
vetor de rótulos com \texttt{sample()} é mais seguro do que sortear índices e usar
\texttt{setdiff()}: garante automaticamente que cada unidade fica em exatamente um
grupo.

\section{\texorpdfstring{A semente \texttt{set.seed()}}{A semente set.seed()}}\label{a-semente-set.seed}

Estes exercícios ilustram um princípio central da simulação reprodutível.

\begin{Shaded}
\begin{Highlighting}[]
\FunctionTok{set.seed}\NormalTok{(}\DecValTok{123}\NormalTok{); }\FunctionTok{sample}\NormalTok{(}\DecValTok{1}\SpecialCharTok{:}\DecValTok{50}\NormalTok{, }\DecValTok{5}\NormalTok{)   }\CommentTok{\# produz sempre a mesma amostra...}
\FunctionTok{set.seed}\NormalTok{(}\DecValTok{123}\NormalTok{); }\FunctionTok{sample}\NormalTok{(}\DecValTok{1}\SpecialCharTok{:}\DecValTok{50}\NormalTok{, }\DecValTok{5}\NormalTok{)   }\CommentTok{\# ...se a semente for a mesma}
\FunctionTok{set.seed}\NormalTok{(}\DecValTok{456}\NormalTok{); }\FunctionTok{sample}\NormalTok{(}\DecValTok{1}\SpecialCharTok{:}\DecValTok{50}\NormalTok{, }\DecValTok{5}\NormalTok{)   }\CommentTok{\# semente diferente {-}\textgreater{} amostra diferente}
\end{Highlighting}
\end{Shaded}

\texttt{set.seed()} fixa o estado do gerador de números pseudoaleatórios. Repor a mesma
semente \textbf{antes} de gerar reproduz exatamente a mesma sequência --- indispensável
para que outra pessoa (ou o próprio, mais tarde) obtenha resultados idênticos. É
por isso que quase todos os exemplos de simulação deste livro começam por
\texttt{set.seed()}. Note-se que a semente deve ser reposta imediatamente antes do bloco
que se quer reproduzir: chamadas intermédias ao gerador alteram o estado.

\section{Probabilidade e variáveis aleatórias}\label{probabilidade-e-variuxe1veis-aleatuxf3rias}

Esta secção mistura cálculo analítico com verificação em R. Sempre que possível,
deduzimos o resultado exato e confirmamo-lo com \texttt{integrate()} ou com as funções
de probabilidade do R.

\textbf{1.} Variável discreta com \(P(X=x)\) dada por \((0{,}1,\,0{,}2,\,0{,}3,\,0{,}2,\,0{,}2)\)
para \(x = 0, 1, 2, 3, 4\).

As probabilidades somam \(1\), logo constituem uma função massa de probabilidade
válida. Os momentos são
\[E(X) = \sum_x x\,P(X=x) = 2{,}2, \qquad
\operatorname{Var}(X) = \sum_x (x - 2{,}2)^2 P(X=x) = 1{,}56.\]

\begin{Shaded}
\begin{Highlighting}[]
\NormalTok{x  }\OtherTok{\textless{}{-}} \DecValTok{0}\SpecialCharTok{:}\DecValTok{4}
\NormalTok{px }\OtherTok{\textless{}{-}} \FunctionTok{c}\NormalTok{(}\FloatTok{0.1}\NormalTok{, }\FloatTok{0.2}\NormalTok{, }\FloatTok{0.3}\NormalTok{, }\FloatTok{0.2}\NormalTok{, }\FloatTok{0.2}\NormalTok{)}
\FunctionTok{sum}\NormalTok{(px)                        }\CommentTok{\# (a) = 1}
\NormalTok{EX  }\OtherTok{\textless{}{-}} \FunctionTok{sum}\NormalTok{(x }\SpecialCharTok{*}\NormalTok{ px)             }\CommentTok{\# (b) = 2.2}
\NormalTok{VarX }\OtherTok{\textless{}{-}} \FunctionTok{sum}\NormalTok{((x }\SpecialCharTok{{-}}\NormalTok{ EX)}\SpecialCharTok{\^{}}\DecValTok{2} \SpecialCharTok{*}\NormalTok{ px)   }\CommentTok{\#      = 1.56}
\FunctionTok{sum}\NormalTok{(px[x }\SpecialCharTok{\textgreater{}} \DecValTok{2}\NormalTok{])                 }\CommentTok{\# (c) P(X \textgreater{} 2) = 0.4}
\end{Highlighting}
\end{Shaded}

No material original, a alínea (c) calculava \texttt{sum(P\_X{[}X\ \textless{}=\ 2{]})}, que dá
\(P(X \le 2) = 0{,}6\), e rotulava-a como ``probabilidade de haver mais de 2
reclamações''. As duas quantidades são complementares: \(P(X > 2) = 1 - P(X \le 2)
= 0{,}4\). Atenção à fronteira: ``mais de 2'' exclui o valor 2.

\textbf{3.} Densidade definida por troços (por exemplo, um rendimento).

\begin{Shaded}
\begin{Highlighting}[]
\NormalTok{f }\OtherTok{\textless{}{-}} \ControlFlowTok{function}\NormalTok{(x) }\FunctionTok{ifelse}\NormalTok{(x }\SpecialCharTok{\textgreater{}=} \DecValTok{0} \SpecialCharTok{\&}\NormalTok{ x }\SpecialCharTok{\textless{}=} \DecValTok{2}\NormalTok{, x}\SpecialCharTok{/}\DecValTok{10} \SpecialCharTok{+} \DecValTok{1}\SpecialCharTok{/}\DecValTok{10}\NormalTok{,}
                 \FunctionTok{ifelse}\NormalTok{(x }\SpecialCharTok{\textgreater{}} \DecValTok{2} \SpecialCharTok{\&}\NormalTok{ x }\SpecialCharTok{\textless{}=} \DecValTok{6}\NormalTok{, }\SpecialCharTok{{-}}\DecValTok{3}\SpecialCharTok{/}\DecValTok{40} \SpecialCharTok{*}\NormalTok{ x }\SpecialCharTok{+} \DecValTok{9}\SpecialCharTok{/}\DecValTok{20}\NormalTok{, }\DecValTok{0}\NormalTok{))}
\FunctionTok{integrate}\NormalTok{(f, }\DecValTok{0}\NormalTok{, }\DecValTok{6}\NormalTok{)                     }\CommentTok{\# verifica que integra 1}
\FunctionTok{integrate}\NormalTok{(}\ControlFlowTok{function}\NormalTok{(x) x }\SpecialCharTok{*} \FunctionTok{f}\NormalTok{(x), }\DecValTok{0}\NormalTok{, }\DecValTok{6}\NormalTok{)  }\CommentTok{\# E(X)}
\FunctionTok{integrate}\NormalTok{(f, }\FloatTok{4.5}\NormalTok{, }\DecValTok{6}\NormalTok{)                   }\CommentTok{\# P(X \textgreater{} 4.5)}
\end{Highlighting}
\end{Shaded}

\textbf{4.} Uniforme contínua em \([10, 20]\): \(f(x) = 1/10\).

Analiticamente, \(E(X) = 15\), \(\operatorname{Var}(X) = (20-10)^2/12 = 100/12
\approx 8{,}33\), \(P(X < 15) = 0{,}5\) e \(P(12 \le X \le 18) = 6/10 = 0{,}6\).

\begin{Shaded}
\begin{Highlighting}[]
\NormalTok{f }\OtherTok{\textless{}{-}} \ControlFlowTok{function}\NormalTok{(x) }\FunctionTok{ifelse}\NormalTok{(x }\SpecialCharTok{\textgreater{}=} \DecValTok{10} \SpecialCharTok{\&}\NormalTok{ x }\SpecialCharTok{\textless{}=} \DecValTok{20}\NormalTok{, }\DecValTok{1}\SpecialCharTok{/}\DecValTok{10}\NormalTok{, }\DecValTok{0}\NormalTok{)}
\FunctionTok{integrate}\NormalTok{(f, }\DecValTok{10}\NormalTok{, }\DecValTok{15}\NormalTok{)    }\CommentTok{\# 0.5}
\FunctionTok{integrate}\NormalTok{(f, }\DecValTok{12}\NormalTok{, }\DecValTok{18}\NormalTok{)    }\CommentTok{\# 0.6}
\FunctionTok{integrate}\NormalTok{(}\ControlFlowTok{function}\NormalTok{(x) x }\SpecialCharTok{*} \FunctionTok{f}\NormalTok{(x), }\DecValTok{10}\NormalTok{, }\DecValTok{20}\NormalTok{)          }\CommentTok{\# E(X) = 15}
\FunctionTok{integrate}\NormalTok{(}\ControlFlowTok{function}\NormalTok{(x) x}\SpecialCharTok{\^{}}\DecValTok{2} \SpecialCharTok{*} \FunctionTok{f}\NormalTok{(x), }\DecValTok{10}\NormalTok{, }\DecValTok{20}\NormalTok{)}\SpecialCharTok{$}\NormalTok{value }\SpecialCharTok{{-}}
  \FunctionTok{integrate}\NormalTok{(}\ControlFlowTok{function}\NormalTok{(x) x }\SpecialCharTok{*} \FunctionTok{f}\NormalTok{(x), }\DecValTok{10}\NormalTok{, }\DecValTok{20}\NormalTok{)}\SpecialCharTok{$}\NormalTok{value}\SpecialCharTok{\^{}}\DecValTok{2} \CommentTok{\# Var(X)}
\end{Highlighting}
\end{Shaded}

\textbf{5.} Densidade triangular \(f(y) = 0{,}125\,y\) em \([0, 4]\).

A função de distribuição é \(F(y) = \int_0^y 0{,}125\,t\,dt = 0{,}0625\,y^2\) em
\([0,4]\) (note-se que \(0{,}0625 = 1/16\), e \(F(4) = 1\)). Daqui,
\(P(Y > 3) = 1 - F(3) = 7/16\); \(P(1 < Y < 3) = F(3) - F(1) = 8/16 = 0{,}5\); e
\(P(1<Y<3\mid Y>1) = \tfrac{8/16}{15/16} = 8/15\).

\begin{Shaded}
\begin{Highlighting}[]
\NormalTok{f }\OtherTok{\textless{}{-}} \ControlFlowTok{function}\NormalTok{(y) }\FunctionTok{ifelse}\NormalTok{(y }\SpecialCharTok{\textgreater{}=} \DecValTok{0} \SpecialCharTok{\&}\NormalTok{ y }\SpecialCharTok{\textless{}=} \DecValTok{4}\NormalTok{, }\FloatTok{0.125} \SpecialCharTok{*}\NormalTok{ y, }\DecValTok{0}\NormalTok{)}
\NormalTok{F }\OtherTok{\textless{}{-}} \ControlFlowTok{function}\NormalTok{(y) }\FunctionTok{ifelse}\NormalTok{(y }\SpecialCharTok{\textless{}} \DecValTok{0}\NormalTok{, }\DecValTok{0}\NormalTok{, }\FunctionTok{ifelse}\NormalTok{(y }\SpecialCharTok{\textless{}=} \DecValTok{4}\NormalTok{, }\FloatTok{0.0625} \SpecialCharTok{*}\NormalTok{ y}\SpecialCharTok{\^{}}\DecValTok{2}\NormalTok{, }\DecValTok{1}\NormalTok{))}
\DecValTok{1} \SpecialCharTok{{-}} \FunctionTok{F}\NormalTok{(}\DecValTok{3}\NormalTok{)                            }\CommentTok{\# P(Y \textgreater{} 3) = 7/16}
\FunctionTok{F}\NormalTok{(}\DecValTok{3}\NormalTok{) }\SpecialCharTok{{-}} \FunctionTok{F}\NormalTok{(}\DecValTok{1}\NormalTok{)                         }\CommentTok{\# P(1 \textless{} Y \textless{} 3) = 0.5}
\NormalTok{(}\FunctionTok{F}\NormalTok{(}\DecValTok{3}\NormalTok{) }\SpecialCharTok{{-}} \FunctionTok{F}\NormalTok{(}\DecValTok{1}\NormalTok{)) }\SpecialCharTok{/}\NormalTok{ (}\DecValTok{1} \SpecialCharTok{{-}} \FunctionTok{F}\NormalTok{(}\DecValTok{1}\NormalTok{))          }\CommentTok{\# P(1\textless{}Y\textless{}3 | Y\textgreater{}1) = 8/15}
\end{Highlighting}
\end{Shaded}

\textbf{6.} Exponencial de parâmetro \(\lambda = 3\): \(f(x) = 3e^{-3x}\), \(x \ge 0\).

A função de distribuição é \(F(x) = 1 - e^{-3x}\) (note-se o sinal \textbf{negativo} no
expoente). Logo \(P(X > 2) = e^{-6}\), e assim por diante.

\begin{Shaded}
\begin{Highlighting}[]
\NormalTok{f }\OtherTok{\textless{}{-}} \ControlFlowTok{function}\NormalTok{(x) }\FunctionTok{ifelse}\NormalTok{(x }\SpecialCharTok{\textgreater{}=} \DecValTok{0}\NormalTok{, }\DecValTok{3} \SpecialCharTok{*} \FunctionTok{exp}\NormalTok{(}\SpecialCharTok{{-}}\DecValTok{3} \SpecialCharTok{*}\NormalTok{ x), }\DecValTok{0}\NormalTok{)}
\FunctionTok{integrate}\NormalTok{(f, }\DecValTok{2}\NormalTok{, }\ConstantTok{Inf}\NormalTok{)                }\CommentTok{\# P(X \textgreater{} 2) = e\^{}{-}6}
\FunctionTok{integrate}\NormalTok{(f, }\FloatTok{0.5}\NormalTok{, }\FloatTok{1.2}\NormalTok{)              }\CommentTok{\# P(0.5 \textless{} X \textless{} 1.2)}
\NormalTok{(}\FunctionTok{integrate}\NormalTok{(f, }\DecValTok{1}\NormalTok{, }\DecValTok{3}\NormalTok{)}\SpecialCharTok{$}\NormalTok{value) }\SpecialCharTok{/}\NormalTok{ (}\FunctionTok{integrate}\NormalTok{(f, }\DecValTok{1}\NormalTok{, }\ConstantTok{Inf}\NormalTok{)}\SpecialCharTok{$}\NormalTok{value)  }\CommentTok{\# P(1\textless{}X\textless{}3 | X\textgreater{}1)}
\end{Highlighting}
\end{Shaded}

A função de distribuição da exponencial é \(F(x) = 1 - e^{-\lambda x}\), com sinal
\textbf{negativo} no expoente --- nunca \(1 - e^{\lambda x}\), que seria negativa e
cresceria sem limite. Este é um erro de sinal fácil de cometer e que conduz a
probabilidades absurdas.

\section{Distribuição uniforme discreta}\label{distribuiuxe7uxe3o-uniforme-discreta-1}

\textbf{1. e 2.} Simulação e verificação de propriedades.

\begin{Shaded}
\begin{Highlighting}[]
\FunctionTok{set.seed}\NormalTok{(}\DecValTok{42}\NormalTok{)}
\NormalTok{valores }\OtherTok{\textless{}{-}} \DecValTok{1}\SpecialCharTok{:}\DecValTok{10}
\NormalTok{sim }\OtherTok{\textless{}{-}} \FunctionTok{sample}\NormalTok{(valores, }\DecValTok{1000}\NormalTok{, }\AttributeTok{replace =} \ConstantTok{TRUE}\NormalTok{)}
\FunctionTok{table}\NormalTok{(sim) }\SpecialCharTok{/} \DecValTok{1000}                  \CommentTok{\# frequências próximas de 0.1 cada}

\CommentTok{\# esperança e variância teóricas de uma uniforme discreta em \{a,...,b\}}
\CommentTok{\# E(X) = (a+b)/2 ;  Var(X) = ((b{-}a+1)\^{}2 {-} 1)/12}
\FunctionTok{mean}\NormalTok{(sim); }\FunctionTok{var}\NormalTok{(sim)}
\end{Highlighting}
\end{Shaded}

Para a uniforme discreta no conjunto \(\{a, a+1, \ldots, b\}\), a esperança é
\((a+b)/2\) e a variância é \(\big((b-a+1)^2 - 1\big)/12\). Repare que a fórmula da
variância \textbf{não} coincide com a da uniforme contínua: o termo \(-1\) e o uso de
\((b-a+1)\) resultam de os valores serem inteiros discretos e não um intervalo
contínuo. A função \texttt{var()} do R, aplicada ao conjunto dos valores, usa o divisor
\(n-1\) e serve apenas como verificação aproximada via simulação.

\textbf{3. a 5.} Somas e transformações.

\begin{Shaded}
\begin{Highlighting}[]
\FunctionTok{set.seed}\NormalTok{(}\DecValTok{42}\NormalTok{)}
\NormalTok{X }\OtherTok{\textless{}{-}} \FunctionTok{sample}\NormalTok{(}\DecValTok{1}\SpecialCharTok{:}\DecValTok{6}\NormalTok{, }\DecValTok{5000}\NormalTok{, }\AttributeTok{replace =} \ConstantTok{TRUE}\NormalTok{)}
\NormalTok{Y }\OtherTok{\textless{}{-}} \FunctionTok{sample}\NormalTok{(}\DecValTok{1}\SpecialCharTok{:}\DecValTok{6}\NormalTok{, }\DecValTok{5000}\NormalTok{, }\AttributeTok{replace =} \ConstantTok{TRUE}\NormalTok{)}
\NormalTok{Z }\OtherTok{\textless{}{-}}\NormalTok{ X }\SpecialCharTok{+}\NormalTok{ Y}
\FunctionTok{mean}\NormalTok{(Z); }\FunctionTok{var}\NormalTok{(Z); }\FunctionTok{mean}\NormalTok{(Z }\SpecialCharTok{\textgreater{}} \DecValTok{8}\NormalTok{)                }\CommentTok{\# soma de dois dados}

\NormalTok{W }\OtherTok{\textless{}{-}} \FunctionTok{sample}\NormalTok{(}\SpecialCharTok{{-}}\DecValTok{5}\SpecialCharTok{:}\DecValTok{5}\NormalTok{, }\DecValTok{5000}\NormalTok{, }\AttributeTok{replace =} \ConstantTok{TRUE}\NormalTok{)     }\CommentTok{\# transformação não linear}
\NormalTok{V }\OtherTok{\textless{}{-}} \DecValTok{3} \SpecialCharTok{*}\NormalTok{ W}\SpecialCharTok{\^{}}\DecValTok{3} \SpecialCharTok{{-}} \DecValTok{2} \SpecialCharTok{*}\NormalTok{ W}\SpecialCharTok{\^{}}\DecValTok{2} \SpecialCharTok{+}\NormalTok{ W}
\FunctionTok{mean}\NormalTok{(V); }\FunctionTok{mean}\NormalTok{(V }\SpecialCharTok{\textgreater{}} \DecValTok{0}\NormalTok{)}
\end{Highlighting}
\end{Shaded}

\section{Distribuição de Bernoulli}\label{distribuiuxe7uxe3o-de-bernoulli-1}

\textbf{1. e 2.} Frequências e momentos empíricos vs.~teóricos.

\begin{Shaded}
\begin{Highlighting}[]
\FunctionTok{set.seed}\NormalTok{(}\DecValTok{42}\NormalTok{)}
\NormalTok{p }\OtherTok{\textless{}{-}} \FloatTok{0.7}
\NormalTok{X }\OtherTok{\textless{}{-}} \FunctionTok{rbinom}\NormalTok{(}\DecValTok{1000}\NormalTok{, }\AttributeTok{size =} \DecValTok{1}\NormalTok{, }\AttributeTok{prob =}\NormalTok{ p)   }\CommentTok{\# Bernoulli(p) = Binomial(1, p)}
\FunctionTok{table}\NormalTok{(X) }\SpecialCharTok{/} \DecValTok{1000}                          \CommentTok{\# \textasciitilde{} (1{-}p, p)}

\CommentTok{\# momentos: E(X) = p ; Var(X) = p(1{-}p)}
\NormalTok{p2 }\OtherTok{\textless{}{-}} \FloatTok{0.4}
\NormalTok{X2 }\OtherTok{\textless{}{-}} \FunctionTok{rbinom}\NormalTok{(}\DecValTok{10000}\NormalTok{, }\DecValTok{1}\NormalTok{, p2)}
\FunctionTok{c}\NormalTok{(}\AttributeTok{media\_emp =} \FunctionTok{mean}\NormalTok{(X2), }\AttributeTok{media\_teo =}\NormalTok{ p2,}
  \AttributeTok{var\_emp =} \FunctionTok{var}\NormalTok{(X2), }\AttributeTok{var\_teo =}\NormalTok{ p2 }\SpecialCharTok{*}\NormalTok{ (}\DecValTok{1} \SpecialCharTok{{-}}\NormalTok{ p2))}
\end{Highlighting}
\end{Shaded}

\textbf{3.} Soma de 5 Bernoulli independentes é uma Binomial(5, p).

\begin{Shaded}
\begin{Highlighting}[]
\FunctionTok{set.seed}\NormalTok{(}\DecValTok{42}\NormalTok{)}
\NormalTok{X }\OtherTok{\textless{}{-}} \FunctionTok{matrix}\NormalTok{(}\FunctionTok{rbinom}\NormalTok{(}\DecValTok{5000} \SpecialCharTok{*} \DecValTok{5}\NormalTok{, }\DecValTok{1}\NormalTok{, }\FloatTok{0.5}\NormalTok{), }\AttributeTok{nrow =} \DecValTok{5000}\NormalTok{)}
\NormalTok{S }\OtherTok{\textless{}{-}} \FunctionTok{rowSums}\NormalTok{(X)}
\FunctionTok{table}\NormalTok{(S) }\SpecialCharTok{/} \DecValTok{5000}                          \CommentTok{\# empírico}
\FunctionTok{dbinom}\NormalTok{(}\DecValTok{0}\SpecialCharTok{:}\DecValTok{5}\NormalTok{, }\AttributeTok{size =} \DecValTok{5}\NormalTok{, }\AttributeTok{prob =} \FloatTok{0.5}\NormalTok{)        }\CommentTok{\# teórico}
\end{Highlighting}
\end{Shaded}

\section{Distribuição binomial}\label{distribuiuxe7uxe3o-binomial-1}

\textbf{1.} \(X \sim \text{Binomial}(10;\,0{,}5)\).

\begin{Shaded}
\begin{Highlighting}[]
\FunctionTok{dbinom}\NormalTok{(}\DecValTok{6}\NormalTok{, }\DecValTok{10}\NormalTok{, }\FloatTok{0.5}\NormalTok{)          }\CommentTok{\# (a) P(X = 6)}
\FunctionTok{pbinom}\NormalTok{(}\DecValTok{4}\NormalTok{, }\DecValTok{10}\NormalTok{, }\FloatTok{0.5}\NormalTok{)          }\CommentTok{\# (b) P(X \textless{}= 4)}
\FunctionTok{pbinom}\NormalTok{(}\DecValTok{7}\NormalTok{, }\DecValTok{10}\NormalTok{, }\FloatTok{0.5}\NormalTok{)          }\CommentTok{\# (c) P(X \textless{}= 7)}
\end{Highlighting}
\end{Shaded}

\textbf{2.} \(X \sim \text{Binomial}(12;\,1/6)\) --- teoria e simulação.

\begin{Shaded}
\begin{Highlighting}[]
\FunctionTok{pbinom}\NormalTok{(}\DecValTok{5}\NormalTok{, }\DecValTok{12}\NormalTok{, }\DecValTok{1}\SpecialCharTok{/}\DecValTok{6}\NormalTok{)                        }\CommentTok{\# (a) P(X \textless{}= 5)}
\FunctionTok{set.seed}\NormalTok{(}\DecValTok{1}\NormalTok{)}
\NormalTok{amostra }\OtherTok{\textless{}{-}} \FunctionTok{rbinom}\NormalTok{(}\DecValTok{1000}\NormalTok{, }\DecValTok{12}\NormalTok{, }\DecValTok{1}\SpecialCharTok{/}\DecValTok{6}\NormalTok{)          }\CommentTok{\# (b)}
\FunctionTok{barplot}\NormalTok{(}\FunctionTok{table}\NormalTok{(amostra) }\SpecialCharTok{/} \FunctionTok{length}\NormalTok{(amostra), }\AttributeTok{col =} \StringTok{"lightblue"}\NormalTok{)  }\CommentTok{\# (c)}
\FunctionTok{points}\NormalTok{(}\DecValTok{0}\SpecialCharTok{:}\DecValTok{12}\NormalTok{, }\FunctionTok{dbinom}\NormalTok{(}\DecValTok{0}\SpecialCharTok{:}\DecValTok{12}\NormalTok{, }\DecValTok{12}\NormalTok{, }\DecValTok{1}\SpecialCharTok{/}\DecValTok{6}\NormalTok{), }\AttributeTok{col =} \StringTok{"magenta"}\NormalTok{, }\AttributeTok{pch =} \DecValTok{19}\NormalTok{)}
\NormalTok{Fn }\OtherTok{\textless{}{-}} \FunctionTok{ecdf}\NormalTok{(amostra); }\FunctionTok{Fn}\NormalTok{(}\DecValTok{5}\NormalTok{); }\FunctionTok{pbinom}\NormalTok{(}\DecValTok{5}\NormalTok{, }\DecValTok{12}\NormalTok{, }\DecValTok{1}\SpecialCharTok{/}\DecValTok{6}\NormalTok{)   }\CommentTok{\# (d) empírica vs teórica}
\DecValTok{12} \SpecialCharTok{*}\NormalTok{ (}\DecValTok{1}\SpecialCharTok{/}\DecValTok{6}\NormalTok{); }\DecValTok{12} \SpecialCharTok{*}\NormalTok{ (}\DecValTok{1}\SpecialCharTok{/}\DecValTok{6}\NormalTok{) }\SpecialCharTok{*}\NormalTok{ (}\DecValTok{5}\SpecialCharTok{/}\DecValTok{6}\NormalTok{)            }\CommentTok{\# (e) média e variância teóricas}
\FunctionTok{mean}\NormalTok{(amostra); }\FunctionTok{var}\NormalTok{(amostra)               }\CommentTok{\# (f) empíricas}
\end{Highlighting}
\end{Shaded}

A \textbf{função de distribuição empírica} \texttt{ecdf(amostra)} devolve uma função: \texttt{Fn(5)}
estima \(P(X \le 5)\) pela proporção de observações \(\le 5\). À medida que a
dimensão da amostra cresce, aproxima-se da função de distribuição teórica
\texttt{pbinom} --- uma manifestação direta da lei dos grandes números.

\textbf{3.} Convergência da frequência relativa para a probabilidade teórica.

\begin{Shaded}
\begin{Highlighting}[]
\FunctionTok{set.seed}\NormalTok{(}\DecValTok{123}\NormalTok{)}
\NormalTok{prob\_teorica }\OtherTok{\textless{}{-}} \FunctionTok{dbinom}\NormalTok{(}\DecValTok{3}\NormalTok{, }\DecValTok{7}\NormalTok{, }\FloatTok{0.5}\NormalTok{)}
\ControlFlowTok{for}\NormalTok{ (n }\ControlFlowTok{in} \FunctionTok{c}\NormalTok{(}\DecValTok{5}\NormalTok{, }\DecValTok{10}\NormalTok{, }\DecValTok{100}\NormalTok{, }\DecValTok{1000}\NormalTok{)) \{}
\NormalTok{  amostra }\OtherTok{\textless{}{-}} \FunctionTok{rbinom}\NormalTok{(n, }\DecValTok{7}\NormalTok{, }\FloatTok{0.5}\NormalTok{)}
  \FunctionTok{cat}\NormalTok{(}\StringTok{"n ="}\NormalTok{, n, }\StringTok{"{-}\textgreater{} P(X=3) empírica:"}\NormalTok{, }\FunctionTok{mean}\NormalTok{(amostra }\SpecialCharTok{==} \DecValTok{3}\NormalTok{),}
      \StringTok{" (teórica:"}\NormalTok{, }\FunctionTok{round}\NormalTok{(prob\_teorica, }\DecValTok{3}\NormalTok{), }\StringTok{")}\SpecialCharTok{\textbackslash{}n}\StringTok{"}\NormalTok{)}
\NormalTok{\}}
\end{Highlighting}
\end{Shaded}

\textbf{4.} Função massa de probabilidade e moda.

\begin{Shaded}
\begin{Highlighting}[]
\NormalTok{k }\OtherTok{\textless{}{-}} \DecValTok{0}\SpecialCharTok{:}\DecValTok{10}
\NormalTok{probs }\OtherTok{\textless{}{-}} \FunctionTok{dbinom}\NormalTok{(k, }\DecValTok{10}\NormalTok{, }\FloatTok{0.25}\NormalTok{)}
\FunctionTok{barplot}\NormalTok{(probs, }\AttributeTok{names.arg =}\NormalTok{ k)}
\NormalTok{k[}\FunctionTok{which.max}\NormalTok{(probs)]        }\CommentTok{\# valor mais provável (moda)}
\end{Highlighting}
\end{Shaded}

\textbf{5. a 8.} Probabilidades de intervalos, caudas e comparação empírica.

\begin{Shaded}
\begin{Highlighting}[]
\FunctionTok{sum}\NormalTok{(}\FunctionTok{dbinom}\NormalTok{(}\DecValTok{8}\SpecialCharTok{:}\DecValTok{12}\NormalTok{, }\DecValTok{15}\NormalTok{, }\FloatTok{0.6}\NormalTok{))               }\CommentTok{\# 5 P(8 \textless{}= X \textless{}= 12)}
\DecValTok{1} \SpecialCharTok{{-}} \FunctionTok{pbinom}\NormalTok{(}\DecValTok{9}\NormalTok{, }\DecValTok{25}\NormalTok{, }\FloatTok{0.3}\NormalTok{)                    }\CommentTok{\# 7 P(X \textgreater{}= 10)}

\FunctionTok{set.seed}\NormalTok{(}\DecValTok{123}\NormalTok{)                             }\CommentTok{\# 8 histograma + fmp + ecdf}
\NormalTok{amostra }\OtherTok{\textless{}{-}} \FunctionTok{rbinom}\NormalTok{(}\DecValTok{700}\NormalTok{, }\DecValTok{30}\NormalTok{, }\FloatTok{0.4}\NormalTok{)}
\FunctionTok{barplot}\NormalTok{(}\FunctionTok{table}\NormalTok{(amostra) }\SpecialCharTok{/} \FunctionTok{length}\NormalTok{(amostra), }\AttributeTok{col =} \StringTok{"lightblue"}\NormalTok{)}
\FunctionTok{points}\NormalTok{(}\DecValTok{0}\SpecialCharTok{:}\DecValTok{30}\NormalTok{, }\FunctionTok{dbinom}\NormalTok{(}\DecValTok{0}\SpecialCharTok{:}\DecValTok{30}\NormalTok{, }\DecValTok{30}\NormalTok{, }\FloatTok{0.4}\NormalTok{), }\AttributeTok{col =} \StringTok{"magenta"}\NormalTok{, }\AttributeTok{pch =} \DecValTok{19}\NormalTok{)}
\end{Highlighting}
\end{Shaded}

\section{Distribuição geométrica}\label{distribuiuxe7uxe3o-geomuxe9trica-1}

Há duas convenções para a distribuição geométrica. Na convenção matemática mais
comum, \(X\) conta o \textbf{número de provas até ao primeiro sucesso} (com suporte
\(\{1, 2, \ldots\}\)). O R, porém, usa a outra convenção: \texttt{rgeom}, \texttt{dgeom} e
\texttt{pgeom} referem-se ao \textbf{número de insucessos antes do primeiro sucesso} (suporte
\(\{0, 1, 2, \ldots\}\)). Assim, se \(Y\) é a variável do R e \(X\) a da convenção
matemática, tem-se \(X = Y + 1\). Tenha esta relação presente ao comparar médias:
a média de \texttt{rgeom} é \((1-p)/p\), e não \(1/p\).

\textbf{1. e 2.} Simulação, função massa de probabilidade e probabilidades acumuladas.

\begin{Shaded}
\begin{Highlighting}[]
\FunctionTok{set.seed}\NormalTok{(}\DecValTok{1}\NormalTok{)}
\NormalTok{amostra }\OtherTok{\textless{}{-}} \FunctionTok{rgeom}\NormalTok{(}\DecValTok{1000}\NormalTok{, }\AttributeTok{prob =} \DecValTok{1}\SpecialCharTok{/}\DecValTok{6}\NormalTok{)          }\CommentTok{\# nº de insucessos (convenção do R)}
\FunctionTok{barplot}\NormalTok{(}\FunctionTok{table}\NormalTok{(amostra) }\SpecialCharTok{/} \FunctionTok{length}\NormalTok{(amostra), }\AttributeTok{col =} \StringTok{"lightblue"}\NormalTok{)}
\FunctionTok{points}\NormalTok{(}\DecValTok{0}\SpecialCharTok{:}\FunctionTok{max}\NormalTok{(amostra), }\FunctionTok{dgeom}\NormalTok{(}\DecValTok{0}\SpecialCharTok{:}\FunctionTok{max}\NormalTok{(amostra), }\DecValTok{1}\SpecialCharTok{/}\DecValTok{6}\NormalTok{), }\AttributeTok{col =} \StringTok{"magenta"}\NormalTok{, }\AttributeTok{pch =} \DecValTok{19}\NormalTok{)}

\FunctionTok{set.seed}\NormalTok{(}\DecValTok{123}\NormalTok{)}
\NormalTok{amostra }\OtherTok{\textless{}{-}} \FunctionTok{rgeom}\NormalTok{(}\DecValTok{5000}\NormalTok{, }\AttributeTok{prob =} \FloatTok{0.2}\NormalTok{)}
\FunctionTok{mean}\NormalTok{(amostra }\SpecialCharTok{\textless{}=} \DecValTok{5}\NormalTok{)          }\CommentTok{\# empírica}
\FunctionTok{pgeom}\NormalTok{(}\DecValTok{5}\NormalTok{, }\AttributeTok{prob =} \FloatTok{0.2}\NormalTok{)        }\CommentTok{\# teórica P(Y \textless{}= 5)}
\end{Highlighting}
\end{Shaded}

\textbf{3.} Efeito do parâmetro \(p\) na forma da distribuição.

\begin{Shaded}
\begin{Highlighting}[]
\FunctionTok{par}\NormalTok{(}\AttributeTok{mfrow =} \FunctionTok{c}\NormalTok{(}\DecValTok{1}\NormalTok{, }\DecValTok{3}\NormalTok{))}
\ControlFlowTok{for}\NormalTok{ (p }\ControlFlowTok{in} \FunctionTok{c}\NormalTok{(}\FloatTok{0.1}\NormalTok{, }\FloatTok{0.5}\NormalTok{, }\FloatTok{0.9}\NormalTok{))}
  \FunctionTok{barplot}\NormalTok{(}\FunctionTok{dgeom}\NormalTok{(}\DecValTok{0}\SpecialCharTok{:}\DecValTok{20}\NormalTok{, p), }\AttributeTok{names.arg =} \DecValTok{0}\SpecialCharTok{:}\DecValTok{20}\NormalTok{,}
          \AttributeTok{main =} \FunctionTok{paste}\NormalTok{(}\StringTok{"p ="}\NormalTok{, p), }\AttributeTok{col =} \StringTok{"lightblue"}\NormalTok{)}
\FunctionTok{par}\NormalTok{(}\AttributeTok{mfrow =} \FunctionTok{c}\NormalTok{(}\DecValTok{1}\NormalTok{, }\DecValTok{1}\NormalTok{))}
\end{Highlighting}
\end{Shaded}

Quanto \textbf{menor} \(p\), mais dispersa a distribuição (é preciso esperar mais pelo
sucesso); quanto \textbf{maior} \(p\), mais concentrada junto de zero.

\textbf{4. a 6.} Média, aplicações e comparação teórica/empírica.

\begin{Shaded}
\begin{Highlighting}[]
\FunctionTok{mean}\NormalTok{(}\FunctionTok{rgeom}\NormalTok{(}\DecValTok{10000}\NormalTok{, }\FloatTok{0.3}\NormalTok{))     }\CommentTok{\# 4  \textasciitilde{} (1{-}p)/p = 0.7/0.3 \textasciitilde{} 2.33}

\FunctionTok{set.seed}\NormalTok{(}\DecValTok{123}\NormalTok{)               }\CommentTok{\# 6}
\NormalTok{p }\OtherTok{\textless{}{-}} \FloatTok{0.3}
\FunctionTok{pgeom}\NormalTok{(}\DecValTok{5}\NormalTok{, p)                 }\CommentTok{\# (a) P(no máximo 5 insucessos)}
\DecValTok{1} \SpecialCharTok{{-}} \FunctionTok{pgeom}\NormalTok{(}\DecValTok{8}\NormalTok{, p)             }\CommentTok{\# (b) P(mais de 8)}
\NormalTok{sim }\OtherTok{\textless{}{-}} \FunctionTok{rgeom}\NormalTok{(}\DecValTok{500}\NormalTok{, p)}
\FunctionTok{c}\NormalTok{(}\AttributeTok{media\_emp =} \FunctionTok{mean}\NormalTok{(sim), }\AttributeTok{media\_teo =}\NormalTok{ (}\DecValTok{1} \SpecialCharTok{{-}}\NormalTok{ p) }\SpecialCharTok{/}\NormalTok{ p,}
  \AttributeTok{dp\_emp =} \FunctionTok{sd}\NormalTok{(sim),      }\AttributeTok{dp\_teo =} \FunctionTok{sqrt}\NormalTok{((}\DecValTok{1} \SpecialCharTok{{-}}\NormalTok{ p) }\SpecialCharTok{/}\NormalTok{ p}\SpecialCharTok{\^{}}\DecValTok{2}\NormalTok{))}
\end{Highlighting}
\end{Shaded}

\section{Distribuição de Poisson}\label{distribuiuxe7uxe3o-de-poisson-1}

\textbf{1.} \(X \sim \text{Poisson}(\lambda = 4)\).

\begin{Shaded}
\begin{Highlighting}[]
\FunctionTok{dpois}\NormalTok{(}\DecValTok{5}\NormalTok{, }\DecValTok{4}\NormalTok{)     }\CommentTok{\# (a) P(X = 5)}
\FunctionTok{ppois}\NormalTok{(}\DecValTok{3}\NormalTok{, }\DecValTok{4}\NormalTok{)     }\CommentTok{\# (b) P(X \textless{}= 3)}
\end{Highlighting}
\end{Shaded}

\textbf{2.} Momentos empíricos (recorde que, na Poisson, \(E(X) = \operatorname{Var}(X) = \lambda\)).

\begin{Shaded}
\begin{Highlighting}[]
\FunctionTok{set.seed}\NormalTok{(}\DecValTok{1}\NormalTok{)}
\NormalTok{amostra }\OtherTok{\textless{}{-}} \FunctionTok{rpois}\NormalTok{(}\DecValTok{30}\NormalTok{, }\DecValTok{4}\NormalTok{)}
\FunctionTok{c}\NormalTok{(}\AttributeTok{media =} \FunctionTok{mean}\NormalTok{(amostra), }\AttributeTok{dp =} \FunctionTok{sd}\NormalTok{(amostra))   }\CommentTok{\# \textasciitilde{} 4 e \textasciitilde{} 2}
\end{Highlighting}
\end{Shaded}

\textbf{3.} Convergência da frequência e dos momentos com o tamanho da amostra.

\begin{Shaded}
\begin{Highlighting}[]
\FunctionTok{set.seed}\NormalTok{(}\DecValTok{123}\NormalTok{)}
\NormalTok{lambda }\OtherTok{\textless{}{-}} \DecValTok{20}
\ControlFlowTok{for}\NormalTok{ (n }\ControlFlowTok{in} \FunctionTok{c}\NormalTok{(}\DecValTok{5}\NormalTok{, }\DecValTok{10}\NormalTok{, }\DecValTok{100}\NormalTok{, }\DecValTok{1000}\NormalTok{)) \{}
\NormalTok{  sim }\OtherTok{\textless{}{-}} \FunctionTok{rpois}\NormalTok{(n, lambda)}
  \FunctionTok{cat}\NormalTok{(}\StringTok{"n ="}\NormalTok{, n, }\StringTok{"| média:"}\NormalTok{, }\FunctionTok{round}\NormalTok{(}\FunctionTok{mean}\NormalTok{(sim), }\DecValTok{2}\NormalTok{),}
      \StringTok{"| variância:"}\NormalTok{, }\FunctionTok{round}\NormalTok{(}\FunctionTok{var}\NormalTok{(sim), }\DecValTok{2}\NormalTok{), }\StringTok{"}\SpecialCharTok{\textbackslash{}n}\StringTok{"}\NormalTok{)}
\NormalTok{\}}
\end{Highlighting}
\end{Shaded}

\textbf{4.} \(X \sim \text{Poisson}(12)\) --- probabilidades e representação.

\begin{Shaded}
\begin{Highlighting}[]
\FunctionTok{ppois}\NormalTok{(}\DecValTok{10}\NormalTok{, }\DecValTok{12}\NormalTok{)                       }\CommentTok{\# (a) P(X \textless{}= 10)}
\FunctionTok{ppois}\NormalTok{(}\DecValTok{15}\NormalTok{, }\DecValTok{12}\NormalTok{, }\AttributeTok{lower.tail =} \ConstantTok{FALSE}\NormalTok{)   }\CommentTok{\# (b) P(X \textgreater{} 15)}
\FunctionTok{set.seed}\NormalTok{(}\DecValTok{1}\NormalTok{)}
\NormalTok{amostra }\OtherTok{\textless{}{-}} \FunctionTok{rpois}\NormalTok{(}\DecValTok{500}\NormalTok{, }\DecValTok{12}\NormalTok{)}
\FunctionTok{mean}\NormalTok{(amostra }\SpecialCharTok{\textless{}=} \DecValTok{10}\NormalTok{)                  }\CommentTok{\# (c) frequência empírica}
\FunctionTok{c}\NormalTok{(}\AttributeTok{media =} \FunctionTok{mean}\NormalTok{(amostra), }\AttributeTok{dp =} \FunctionTok{sd}\NormalTok{(amostra))   }\CommentTok{\# (d)  — nota: \textquotesingle{}\textless{}{-}\textquotesingle{}, não \textquotesingle{}\textless{}\textquotesingle{}}
\FunctionTok{barplot}\NormalTok{(}\FunctionTok{table}\NormalTok{(amostra) }\SpecialCharTok{/} \FunctionTok{length}\NormalTok{(amostra), }\AttributeTok{col =} \StringTok{"lightblue"}\NormalTok{)   }\CommentTok{\# (e)}
\FunctionTok{points}\NormalTok{(}\DecValTok{0}\SpecialCharTok{:}\DecValTok{25}\NormalTok{, }\FunctionTok{dpois}\NormalTok{(}\DecValTok{0}\SpecialCharTok{:}\DecValTok{25}\NormalTok{, }\DecValTok{12}\NormalTok{), }\AttributeTok{col =} \StringTok{"magenta"}\NormalTok{, }\AttributeTok{pch =} \DecValTok{19}\NormalTok{)}
\end{Highlighting}
\end{Shaded}

No material original surgia \texttt{desv\_padrao\_empirico\ \textless{}\ sd(amostra)}, com um único
\texttt{\textless{}}. Isto \textbf{não} é uma atribuição --- é uma comparação lógica, cujo resultado é
descartado, deixando a variável por criar. A atribuição correta é \texttt{\textless{}-}. É um
lapso comum e difícil de detetar, porque não gera erro.

\textbf{5. e 6.} Função de distribuição empírica vs.~teórica.

\begin{Shaded}
\begin{Highlighting}[]
\FunctionTok{set.seed}\NormalTok{(}\DecValTok{123}\NormalTok{)}
\NormalTok{aves }\OtherTok{\textless{}{-}} \FunctionTok{rpois}\NormalTok{(}\DecValTok{1000}\NormalTok{, }\DecValTok{8}\NormalTok{)}
\FunctionTok{plot}\NormalTok{(}\FunctionTok{ecdf}\NormalTok{(aves), }\AttributeTok{verticals =} \ConstantTok{TRUE}\NormalTok{, }\AttributeTok{do.points =} \ConstantTok{FALSE}\NormalTok{, }\AttributeTok{col =} \StringTok{"blue"}\NormalTok{,}
     \AttributeTok{main =} \StringTok{"F. distribuição empírica vs teórica"}\NormalTok{)}
\FunctionTok{points}\NormalTok{(}\DecValTok{0}\SpecialCharTok{:}\FunctionTok{max}\NormalTok{(aves), }\FunctionTok{ppois}\NormalTok{(}\DecValTok{0}\SpecialCharTok{:}\FunctionTok{max}\NormalTok{(aves), }\DecValTok{8}\NormalTok{), }\AttributeTok{col =} \StringTok{"red"}\NormalTok{, }\AttributeTok{pch =} \DecValTok{19}\NormalTok{)}
\end{Highlighting}
\end{Shaded}

\textbf{9. e 10.} Teorema limite central aplicado à Poisson (a média de amostras é
aproximadamente normal).

\begin{Shaded}
\begin{Highlighting}[]
\FunctionTok{set.seed}\NormalTok{(}\DecValTok{1}\NormalTok{)}
\NormalTok{m }\OtherTok{\textless{}{-}} \DecValTok{1000}\NormalTok{; n }\OtherTok{\textless{}{-}} \DecValTok{30}
\NormalTok{medias }\OtherTok{\textless{}{-}} \FunctionTok{rowMeans}\NormalTok{(}\FunctionTok{matrix}\NormalTok{(}\FunctionTok{rpois}\NormalTok{(m }\SpecialCharTok{*}\NormalTok{ n, }\DecValTok{2}\NormalTok{), }\AttributeTok{nrow =}\NormalTok{ m))}
\FunctionTok{hist}\NormalTok{(medias, }\AttributeTok{probability =} \ConstantTok{TRUE}\NormalTok{, }\AttributeTok{col =} \StringTok{"lightblue"}\NormalTok{)}
\FunctionTok{curve}\NormalTok{(}\FunctionTok{dnorm}\NormalTok{(x, }\AttributeTok{mean =} \DecValTok{2}\NormalTok{, }\AttributeTok{sd =} \FunctionTok{sqrt}\NormalTok{(}\DecValTok{2} \SpecialCharTok{/}\NormalTok{ n)), }\AttributeTok{col =} \StringTok{"red"}\NormalTok{, }\AttributeTok{lwd =} \DecValTok{2}\NormalTok{, }\AttributeTok{add =} \ConstantTok{TRUE}\NormalTok{)}
\end{Highlighting}
\end{Shaded}

\section{Distribuição uniforme contínua}\label{distribuiuxe7uxe3o-uniforme-contuxednua-1}

\textbf{1.} Momentos: \(E(X) = (a+b)/2\), \(\operatorname{Var}(X) = (b-a)^2/12\).

\begin{Shaded}
\begin{Highlighting}[]
\FunctionTok{set.seed}\NormalTok{(}\DecValTok{1}\NormalTok{)}
\NormalTok{amostra }\OtherTok{\textless{}{-}} \FunctionTok{runif}\NormalTok{(}\DecValTok{1000}\NormalTok{, }\DecValTok{0}\NormalTok{, }\DecValTok{1}\NormalTok{)}
\FunctionTok{c}\NormalTok{(}\AttributeTok{media =} \FunctionTok{mean}\NormalTok{(amostra), }\AttributeTok{var =} \FunctionTok{var}\NormalTok{(amostra))   }\CommentTok{\# \textasciitilde{} 0.5 e \textasciitilde{} 1/12}
\end{Highlighting}
\end{Shaded}

\textbf{2.} Histograma com densidade teórica.

\begin{Shaded}
\begin{Highlighting}[]
\FunctionTok{set.seed}\NormalTok{(}\DecValTok{1}\NormalTok{)}
\NormalTok{amostra }\OtherTok{\textless{}{-}} \FunctionTok{runif}\NormalTok{(}\DecValTok{500}\NormalTok{, }\SpecialCharTok{{-}}\DecValTok{3}\NormalTok{, }\DecValTok{7}\NormalTok{)}
\FunctionTok{hist}\NormalTok{(amostra, }\AttributeTok{probability =} \ConstantTok{TRUE}\NormalTok{, }\AttributeTok{col =} \StringTok{"lightblue"}\NormalTok{)}
\FunctionTok{curve}\NormalTok{(}\FunctionTok{dunif}\NormalTok{(x, }\SpecialCharTok{{-}}\DecValTok{3}\NormalTok{, }\DecValTok{7}\NormalTok{), }\AttributeTok{col =} \StringTok{"red"}\NormalTok{, }\AttributeTok{lwd =} \DecValTok{2}\NormalTok{, }\AttributeTok{add =} \ConstantTok{TRUE}\NormalTok{)}
\end{Highlighting}
\end{Shaded}

\textbf{3. a 5, 7, 8.} Probabilidades, transformações lineares e comparação empírica.

\begin{Shaded}
\begin{Highlighting}[]
\FunctionTok{punif}\NormalTok{(}\DecValTok{100}\NormalTok{, }\DecValTok{85}\NormalTok{, }\DecValTok{105}\NormalTok{)                              }\CommentTok{\# 3 P(X \textless{} 100)}
\FunctionTok{set.seed}\NormalTok{(}\DecValTok{1}\NormalTok{)}
\NormalTok{amostra }\OtherTok{\textless{}{-}} \FunctionTok{runif}\NormalTok{(}\DecValTok{10000}\NormalTok{, }\DecValTok{2}\NormalTok{, }\DecValTok{8}\NormalTok{); }\FunctionTok{mean}\NormalTok{(amostra }\SpecialCharTok{\textgreater{}} \DecValTok{5}\NormalTok{) }\CommentTok{\# 4 vs punif(5,2,8,lower.tail=FALSE)}
\FunctionTok{mean}\NormalTok{(}\FunctionTok{runif}\NormalTok{(}\DecValTok{1000}\NormalTok{, }\DecValTok{10}\NormalTok{, }\DecValTok{20}\NormalTok{) }\SpecialCharTok{\%in\%} \DecValTok{12}\SpecialCharTok{:}\DecValTok{15}\NormalTok{)             }\CommentTok{\# 5 (aproximação da faixa [12,15])}

\CommentTok{\# 7 transformação linear Z = 3X + 2, com X \textasciitilde{} U(0,1)}
\NormalTok{z }\OtherTok{\textless{}{-}} \DecValTok{3} \SpecialCharTok{*} \FunctionTok{runif}\NormalTok{(}\DecValTok{1000}\NormalTok{) }\SpecialCharTok{+} \DecValTok{2}
\FunctionTok{c}\NormalTok{(}\AttributeTok{media =} \FunctionTok{mean}\NormalTok{(z), }\AttributeTok{media\_teo =} \DecValTok{3} \SpecialCharTok{*} \FloatTok{0.5} \SpecialCharTok{+} \DecValTok{2}\NormalTok{,      }\CommentTok{\# E(Z) = 3·E(X)+2 = 3.5}
  \AttributeTok{var =} \FunctionTok{var}\NormalTok{(z),   }\AttributeTok{var\_teo =} \DecValTok{9} \SpecialCharTok{*}\NormalTok{ (}\DecValTok{1}\SpecialCharTok{/}\DecValTok{12}\NormalTok{))          }\CommentTok{\# Var(Z) = 9·Var(X) = 9/12}
\end{Highlighting}
\end{Shaded}

\textbf{6, 9, 10.} Lei dos grandes números e teorema limite central.

\begin{Shaded}
\begin{Highlighting}[]
\FunctionTok{set.seed}\NormalTok{(}\DecValTok{123}\NormalTok{)                                     }\CommentTok{\# 9 média converge para 0.5}
\FunctionTok{sapply}\NormalTok{(}\FunctionTok{c}\NormalTok{(}\DecValTok{100}\NormalTok{, }\DecValTok{1000}\NormalTok{, }\DecValTok{10000}\NormalTok{, }\DecValTok{100000}\NormalTok{),}
       \ControlFlowTok{function}\NormalTok{(n) }\FunctionTok{mean}\NormalTok{(}\FunctionTok{runif}\NormalTok{(n)))}

\FunctionTok{set.seed}\NormalTok{(}\DecValTok{1}\NormalTok{)                                       }\CommentTok{\# 10 TLC para a soma/média}
\NormalTok{m }\OtherTok{\textless{}{-}} \DecValTok{8000}\NormalTok{; n }\OtherTok{\textless{}{-}} \DecValTok{100}
\NormalTok{amostras }\OtherTok{\textless{}{-}} \FunctionTok{matrix}\NormalTok{(}\FunctionTok{runif}\NormalTok{(m }\SpecialCharTok{*}\NormalTok{ n, }\DecValTok{10}\NormalTok{, }\DecValTok{30}\NormalTok{), }\AttributeTok{nrow =}\NormalTok{ m)}
\NormalTok{media }\OtherTok{\textless{}{-}}\NormalTok{ (}\DecValTok{10} \SpecialCharTok{+} \DecValTok{30}\NormalTok{) }\SpecialCharTok{/} \DecValTok{2}\NormalTok{; variancia }\OtherTok{\textless{}{-}}\NormalTok{ (}\DecValTok{30} \SpecialCharTok{{-}} \DecValTok{10}\NormalTok{)}\SpecialCharTok{\^{}}\DecValTok{2} \SpecialCharTok{/} \DecValTok{12}
\FunctionTok{hist}\NormalTok{(}\FunctionTok{rowMeans}\NormalTok{(amostras), }\AttributeTok{probability =} \ConstantTok{TRUE}\NormalTok{, }\AttributeTok{col =} \StringTok{"lightblue"}\NormalTok{)}
\FunctionTok{curve}\NormalTok{(}\FunctionTok{dnorm}\NormalTok{(x, media, }\FunctionTok{sqrt}\NormalTok{(variancia }\SpecialCharTok{/}\NormalTok{ n)), }\AttributeTok{col =} \StringTok{"red"}\NormalTok{, }\AttributeTok{lwd =} \DecValTok{2}\NormalTok{, }\AttributeTok{add =} \ConstantTok{TRUE}\NormalTok{)}
\end{Highlighting}
\end{Shaded}

\section{Distribuição normal}\label{distribuiuxe7uxe3o-normal-1}

\textbf{1.} Amostra e densidade teórica.

\begin{Shaded}
\begin{Highlighting}[]
\FunctionTok{set.seed}\NormalTok{(}\DecValTok{1}\NormalTok{)}
\NormalTok{amostra }\OtherTok{\textless{}{-}} \FunctionTok{rnorm}\NormalTok{(}\DecValTok{1000}\NormalTok{, }\AttributeTok{mean =} \DecValTok{50}\NormalTok{, }\AttributeTok{sd =} \DecValTok{5}\NormalTok{)}
\FunctionTok{hist}\NormalTok{(amostra, }\AttributeTok{probability =} \ConstantTok{TRUE}\NormalTok{, }\AttributeTok{col =} \StringTok{"lightblue"}\NormalTok{)}
\FunctionTok{curve}\NormalTok{(}\FunctionTok{dnorm}\NormalTok{(x, }\DecValTok{50}\NormalTok{, }\DecValTok{5}\NormalTok{), }\AttributeTok{col =} \StringTok{"red"}\NormalTok{, }\AttributeTok{lwd =} \DecValTok{2}\NormalTok{, }\AttributeTok{add =} \ConstantTok{TRUE}\NormalTok{)}
\end{Highlighting}
\end{Shaded}

\textbf{2.} Soma de normais independentes é normal: se \(X \sim N(\mu_1, \sigma_1^2)\) e
\(Y \sim N(\mu_2, \sigma_2^2)\) forem independentes, \(X + Y \sim N(\mu_1 + \mu_2,\,
\sigma_1^2 + \sigma_2^2)\).

\begin{Shaded}
\begin{Highlighting}[]
\FunctionTok{set.seed}\NormalTok{(}\DecValTok{1}\NormalTok{)}
\NormalTok{soma }\OtherTok{\textless{}{-}} \FunctionTok{rnorm}\NormalTok{(}\DecValTok{1000}\NormalTok{, }\DecValTok{5}\NormalTok{, }\DecValTok{2}\NormalTok{) }\SpecialCharTok{+} \FunctionTok{rnorm}\NormalTok{(}\DecValTok{1000}\NormalTok{, }\DecValTok{10}\NormalTok{, }\DecValTok{3}\NormalTok{)}
\FunctionTok{hist}\NormalTok{(soma, }\AttributeTok{probability =} \ConstantTok{TRUE}\NormalTok{, }\AttributeTok{col =} \StringTok{"lightblue"}\NormalTok{)}
\FunctionTok{curve}\NormalTok{(}\FunctionTok{dnorm}\NormalTok{(x, }\AttributeTok{mean =} \DecValTok{15}\NormalTok{, }\AttributeTok{sd =} \FunctionTok{sqrt}\NormalTok{(}\DecValTok{2}\SpecialCharTok{\^{}}\DecValTok{2} \SpecialCharTok{+} \DecValTok{3}\SpecialCharTok{\^{}}\DecValTok{2}\NormalTok{)), }\AttributeTok{col =} \StringTok{"red"}\NormalTok{, }\AttributeTok{lwd =} \DecValTok{2}\NormalTok{, }\AttributeTok{add =} \ConstantTok{TRUE}\NormalTok{)}
\end{Highlighting}
\end{Shaded}

As variâncias somam-se (para variáveis independentes), \textbf{não} os
desvios-padrão: aqui \(\sigma = \sqrt{2^2 + 3^2} = \sqrt{13}\), e não \(2 + 3 = 5\).
Confundir as duas coisas é um erro clássico.

\textbf{3.} Soma de quadrados de normais padrão dá uma qui-quadrado: se \(Z_1, Z_2 \sim
N(0,1)\) independentes, \(Z_1^2 + Z_2^2 \sim \chi^2_2\).

\begin{Shaded}
\begin{Highlighting}[]
\FunctionTok{set.seed}\NormalTok{(}\DecValTok{1}\NormalTok{)}
\NormalTok{Q }\OtherTok{\textless{}{-}} \FunctionTok{rnorm}\NormalTok{(}\DecValTok{5000}\NormalTok{)}\SpecialCharTok{\^{}}\DecValTok{2} \SpecialCharTok{+} \FunctionTok{rnorm}\NormalTok{(}\DecValTok{5000}\NormalTok{)}\SpecialCharTok{\^{}}\DecValTok{2}
\FunctionTok{hist}\NormalTok{(Q, }\AttributeTok{probability =} \ConstantTok{TRUE}\NormalTok{, }\AttributeTok{col =} \StringTok{"lightblue"}\NormalTok{)}
\FunctionTok{curve}\NormalTok{(}\FunctionTok{dchisq}\NormalTok{(x, }\AttributeTok{df =} \DecValTok{2}\NormalTok{), }\AttributeTok{col =} \StringTok{"red"}\NormalTok{, }\AttributeTok{lwd =} \DecValTok{2}\NormalTok{, }\AttributeTok{add =} \ConstantTok{TRUE}\NormalTok{)}
\end{Highlighting}
\end{Shaded}

\textbf{4 a 9.} Distribuições amostrais fundamentais (a base da inferência): a média
padronizada segue uma \(t\) de Student; \((n-1)S^2/\sigma^2\) segue uma
qui-quadrado; a soma de quadrados de \(n\) normais padrão é \(\chi^2_n\).

\begin{Shaded}
\begin{Highlighting}[]
\FunctionTok{set.seed}\NormalTok{(}\DecValTok{1}\NormalTok{)}
\NormalTok{n }\OtherTok{\textless{}{-}} \DecValTok{20}\NormalTok{; m }\OtherTok{\textless{}{-}} \DecValTok{1000}

\CommentTok{\# soma de n quadrados \textasciitilde{} qui{-}quadrado com n graus de liberdade (ex. 5)}
\NormalTok{sq }\OtherTok{\textless{}{-}} \FunctionTok{rowSums}\NormalTok{(}\FunctionTok{matrix}\NormalTok{(}\FunctionTok{rnorm}\NormalTok{(m }\SpecialCharTok{*}\NormalTok{ n), }\AttributeTok{nrow =}\NormalTok{ m)}\SpecialCharTok{\^{}}\DecValTok{2}\NormalTok{)}
\FunctionTok{hist}\NormalTok{(sq, }\AttributeTok{probability =} \ConstantTok{TRUE}\NormalTok{, }\AttributeTok{col =} \StringTok{"lightblue"}\NormalTok{)}
\FunctionTok{curve}\NormalTok{(}\FunctionTok{dchisq}\NormalTok{(x, }\AttributeTok{df =}\NormalTok{ n), }\AttributeTok{col =} \StringTok{"red"}\NormalTok{, }\AttributeTok{lwd =} \DecValTok{2}\NormalTok{, }\AttributeTok{add =} \ConstantTok{TRUE}\NormalTok{)}

\CommentTok{\# estatística t: sqrt(n)·média/desvio \textasciitilde{} t\_\{n{-}1\} (ex. 8)}
\NormalTok{n }\OtherTok{\textless{}{-}} \DecValTok{10}
\NormalTok{A }\OtherTok{\textless{}{-}} \FunctionTok{matrix}\NormalTok{(}\FunctionTok{rnorm}\NormalTok{(m }\SpecialCharTok{*}\NormalTok{ n), }\AttributeTok{nrow =}\NormalTok{ m)}
\NormalTok{Tstat }\OtherTok{\textless{}{-}} \FunctionTok{sqrt}\NormalTok{(n) }\SpecialCharTok{*} \FunctionTok{rowMeans}\NormalTok{(A) }\SpecialCharTok{/} \FunctionTok{apply}\NormalTok{(A, }\DecValTok{1}\NormalTok{, sd)}
\FunctionTok{hist}\NormalTok{(Tstat, }\AttributeTok{probability =} \ConstantTok{TRUE}\NormalTok{, }\AttributeTok{col =} \StringTok{"lightblue"}\NormalTok{)}
\FunctionTok{curve}\NormalTok{(}\FunctionTok{dt}\NormalTok{(x, }\AttributeTok{df =}\NormalTok{ n }\SpecialCharTok{{-}} \DecValTok{1}\NormalTok{), }\AttributeTok{col =} \StringTok{"red"}\NormalTok{, }\AttributeTok{lwd =} \DecValTok{2}\NormalTok{, }\AttributeTok{add =} \ConstantTok{TRUE}\NormalTok{)}
\end{Highlighting}
\end{Shaded}

\section{Distribuição exponencial}\label{distribuiuxe7uxe3o-exponencial-1}

\textbf{1. e 2.} Amostra, tempo médio e probabilidades de cauda.

\begin{Shaded}
\begin{Highlighting}[]
\FunctionTok{set.seed}\NormalTok{(}\DecValTok{1}\NormalTok{)}
\NormalTok{amostra }\OtherTok{\textless{}{-}} \FunctionTok{rexp}\NormalTok{(}\DecValTok{1000}\NormalTok{, }\AttributeTok{rate =} \DecValTok{2}\NormalTok{)}
\FunctionTok{mean}\NormalTok{(amostra)                       }\CommentTok{\# \textasciitilde{} 1/2 (média teórica = 1/rate)}
\FunctionTok{hist}\NormalTok{(amostra, }\AttributeTok{probability =} \ConstantTok{TRUE}\NormalTok{, }\AttributeTok{col =} \StringTok{"lightblue"}\NormalTok{)}
\FunctionTok{curve}\NormalTok{(}\FunctionTok{dexp}\NormalTok{(x, }\AttributeTok{rate =} \DecValTok{2}\NormalTok{), }\AttributeTok{col =} \StringTok{"red"}\NormalTok{, }\AttributeTok{lwd =} \DecValTok{2}\NormalTok{, }\AttributeTok{add =} \ConstantTok{TRUE}\NormalTok{)}

\FunctionTok{pexp}\NormalTok{(}\DecValTok{3}\NormalTok{, }\AttributeTok{rate =} \FloatTok{0.5}\NormalTok{, }\AttributeTok{lower.tail =} \ConstantTok{FALSE}\NormalTok{)   }\CommentTok{\# P(X \textgreater{} 3) = e\^{}\{{-}1.5\}}
\end{Highlighting}
\end{Shaded}

\textbf{3, 9, 10, 11.} Soma de exponenciais independentes com a mesma taxa é uma Gama:
\(\sum_{i=1}^{n} X_i \sim \Gamma(n, \lambda)\).

\begin{Shaded}
\begin{Highlighting}[]
\FunctionTok{set.seed}\NormalTok{(}\DecValTok{1}\NormalTok{)}
\NormalTok{n }\OtherTok{\textless{}{-}} \DecValTok{5}\NormalTok{; m }\OtherTok{\textless{}{-}} \DecValTok{5000}
\NormalTok{soma }\OtherTok{\textless{}{-}} \FunctionTok{rowSums}\NormalTok{(}\FunctionTok{matrix}\NormalTok{(}\FunctionTok{rexp}\NormalTok{(m }\SpecialCharTok{*}\NormalTok{ n, }\AttributeTok{rate =} \DecValTok{2}\NormalTok{), }\AttributeTok{nrow =}\NormalTok{ m))}
\FunctionTok{hist}\NormalTok{(soma, }\AttributeTok{probability =} \ConstantTok{TRUE}\NormalTok{, }\AttributeTok{col =} \StringTok{"lightblue"}\NormalTok{)}
\FunctionTok{curve}\NormalTok{(}\FunctionTok{dgamma}\NormalTok{(x, }\AttributeTok{shape =}\NormalTok{ n, }\AttributeTok{rate =} \DecValTok{2}\NormalTok{), }\AttributeTok{col =} \StringTok{"red"}\NormalTok{, }\AttributeTok{lwd =} \DecValTok{2}\NormalTok{, }\AttributeTok{add =} \ConstantTok{TRUE}\NormalTok{)}
\FunctionTok{c}\NormalTok{(}\AttributeTok{media\_emp =} \FunctionTok{mean}\NormalTok{(soma), }\AttributeTok{media\_teo =}\NormalTok{ n }\SpecialCharTok{/} \DecValTok{2}\NormalTok{,     }\CommentTok{\# E = n/λ}
  \AttributeTok{var\_emp =} \FunctionTok{var}\NormalTok{(soma),   }\AttributeTok{var\_teo =}\NormalTok{ n }\SpecialCharTok{/} \DecValTok{2}\SpecialCharTok{\^{}}\DecValTok{2}\NormalTok{)       }\CommentTok{\# Var = n/λ²}
\end{Highlighting}
\end{Shaded}

\textbf{4 a 8.} Somas de processos, transformações lineares e teorema limite central.

\begin{Shaded}
\begin{Highlighting}[]
\CommentTok{\# 8 transformação linear Y = 2X + 1}
\FunctionTok{set.seed}\NormalTok{(}\DecValTok{1}\NormalTok{)}
\NormalTok{X }\OtherTok{\textless{}{-}} \FunctionTok{rexp}\NormalTok{(}\DecValTok{5000}\NormalTok{, }\AttributeTok{rate =} \FloatTok{0.5}\NormalTok{)}
\NormalTok{Y }\OtherTok{\textless{}{-}} \DecValTok{2} \SpecialCharTok{*}\NormalTok{ X }\SpecialCharTok{+} \DecValTok{1}
\FunctionTok{c}\NormalTok{(}\AttributeTok{media =} \FunctionTok{mean}\NormalTok{(Y), }\AttributeTok{media\_teo =} \DecValTok{2} \SpecialCharTok{*}\NormalTok{ (}\DecValTok{1}\SpecialCharTok{/}\FloatTok{0.5}\NormalTok{) }\SpecialCharTok{+} \DecValTok{1}\NormalTok{,   }\CommentTok{\# E(Y) = 2·2 + 1 = 5}
  \AttributeTok{var =} \FunctionTok{var}\NormalTok{(Y),   }\AttributeTok{var\_teo =} \DecValTok{4} \SpecialCharTok{*}\NormalTok{ (}\DecValTok{1}\SpecialCharTok{/}\FloatTok{0.5}\SpecialCharTok{\^{}}\DecValTok{2}\NormalTok{))         }\CommentTok{\# Var(Y) = 4·4 = 16}

\CommentTok{\# 5 TLC: média de n exponenciais é aproximadamente normal}
\NormalTok{n }\OtherTok{\textless{}{-}} \DecValTok{30}
\NormalTok{medias }\OtherTok{\textless{}{-}} \FunctionTok{rowMeans}\NormalTok{(}\FunctionTok{matrix}\NormalTok{(}\FunctionTok{rexp}\NormalTok{(}\DecValTok{5000} \SpecialCharTok{*}\NormalTok{ n, }\AttributeTok{rate =} \FloatTok{1.5}\NormalTok{), }\AttributeTok{nrow =} \DecValTok{5000}\NormalTok{))}
\FunctionTok{hist}\NormalTok{(medias, }\AttributeTok{probability =} \ConstantTok{TRUE}\NormalTok{, }\AttributeTok{col =} \StringTok{"lightblue"}\NormalTok{)}
\FunctionTok{curve}\NormalTok{(}\FunctionTok{dnorm}\NormalTok{(x, }\DecValTok{1}\SpecialCharTok{/}\FloatTok{1.5}\NormalTok{, }\DecValTok{1}\SpecialCharTok{/}\NormalTok{(}\FloatTok{1.5} \SpecialCharTok{*} \FunctionTok{sqrt}\NormalTok{(n))), }\AttributeTok{col =} \StringTok{"red"}\NormalTok{, }\AttributeTok{lwd =} \DecValTok{2}\NormalTok{, }\AttributeTok{add =} \ConstantTok{TRUE}\NormalTok{)}
\end{Highlighting}
\end{Shaded}

Muitos destes exercícios ilustram um mesmo padrão: \textbf{somas} de exponenciais dão
Gamas exatas, enquanto \textbf{médias} de qualquer distribuição, para \(n\) grande, dão
aproximadamente normais (teorema limite central). Distinguir o resultado exato
do aproximado é essencial: a soma de 5 exponenciais é \emph{exatamente} \(\Gamma(5,
\lambda)\), mas a sua média só é \emph{aproximadamente} normal.

\section{Simulação de variáveis aleatórias}\label{simulauxe7uxe3o-de-variuxe1veis-aleatuxf3rias}

Esta secção reúne exercícios que geram amostras de distribuições conhecidas,
comparam frequências empíricas com probabilidades teóricas (via \texttt{ecdf}) e
ilustram os teoremas limites. Como o padrão se repete, apresentamos um exemplo
representativo de cada tipo.

\textbf{Comparar frequências empíricas com teóricas.}

\begin{Shaded}
\begin{Highlighting}[]
\FunctionTok{set.seed}\NormalTok{(}\DecValTok{123}\NormalTok{)}
\NormalTok{amostra }\OtherTok{\textless{}{-}} \FunctionTok{rbinom}\NormalTok{(}\DecValTok{5000}\NormalTok{, }\AttributeTok{size =} \DecValTok{20}\NormalTok{, }\AttributeTok{prob =} \FloatTok{0.7}\NormalTok{)}
\FunctionTok{hist}\NormalTok{(amostra, }\AttributeTok{probability =} \ConstantTok{TRUE}\NormalTok{, }\AttributeTok{col =} \StringTok{"lightblue"}\NormalTok{)}
\FunctionTok{points}\NormalTok{(}\DecValTok{0}\SpecialCharTok{:}\DecValTok{20}\NormalTok{, }\FunctionTok{dbinom}\NormalTok{(}\DecValTok{0}\SpecialCharTok{:}\DecValTok{20}\NormalTok{, }\DecValTok{20}\NormalTok{, }\FloatTok{0.7}\NormalTok{), }\AttributeTok{col =} \StringTok{"magenta"}\NormalTok{)   }\CommentTok{\# fmp teórica}

\NormalTok{Fn }\OtherTok{\textless{}{-}} \FunctionTok{ecdf}\NormalTok{(amostra)}
\FunctionTok{data.frame}\NormalTok{(}\AttributeTok{empirica =} \FunctionTok{Fn}\NormalTok{(}\DecValTok{10}\NormalTok{), }\AttributeTok{teorica =} \FunctionTok{pbinom}\NormalTok{(}\DecValTok{10}\NormalTok{, }\DecValTok{20}\NormalTok{, }\FloatTok{0.7}\NormalTok{))  }\CommentTok{\# P(X \textless{}= 10)}
\end{Highlighting}
\end{Shaded}

\textbf{Convergência da média amostral (lei dos grandes números).}

\begin{Shaded}
\begin{Highlighting}[]
\FunctionTok{set.seed}\NormalTok{(}\DecValTok{123}\NormalTok{)}
\NormalTok{lambda }\OtherTok{\textless{}{-}} \DecValTok{3}
\FunctionTok{data.frame}\NormalTok{(}
  \AttributeTok{n =} \FunctionTok{c}\NormalTok{(}\DecValTok{100}\NormalTok{, }\DecValTok{1000}\NormalTok{, }\DecValTok{10000}\NormalTok{, }\DecValTok{100000}\NormalTok{),}
  \AttributeTok{media =} \FunctionTok{sapply}\NormalTok{(}\FunctionTok{c}\NormalTok{(}\DecValTok{100}\NormalTok{, }\DecValTok{1000}\NormalTok{, }\DecValTok{10000}\NormalTok{, }\DecValTok{100000}\NormalTok{),}
                 \ControlFlowTok{function}\NormalTok{(n) }\FunctionTok{mean}\NormalTok{(}\FunctionTok{rpois}\NormalTok{(n, lambda))),}
  \AttributeTok{teorica =}\NormalTok{ lambda)}
\end{Highlighting}
\end{Shaded}

\textbf{Teorema limite central (soma e média padronizadas).} Este é o exercício mais
completo do conjunto: gerar muitas amostras, calcular soma e média de cada uma, e
sobrepor a curva normal apropriada.

\begin{Shaded}
\begin{Highlighting}[]
\FunctionTok{set.seed}\NormalTok{(}\DecValTok{1580}\NormalTok{)}
\NormalTok{m }\OtherTok{\textless{}{-}} \DecValTok{1234}\NormalTok{; n }\OtherTok{\textless{}{-}} \DecValTok{50}\NormalTok{; lambda }\OtherTok{\textless{}{-}} \FloatTok{0.21}
\NormalTok{media\_teo }\OtherTok{\textless{}{-}} \DecValTok{1} \SpecialCharTok{/}\NormalTok{ lambda; var\_teo }\OtherTok{\textless{}{-}} \DecValTok{1} \SpecialCharTok{/}\NormalTok{ lambda}\SpecialCharTok{\^{}}\DecValTok{2}
\NormalTok{amostras }\OtherTok{\textless{}{-}} \FunctionTok{matrix}\NormalTok{(}\FunctionTok{rexp}\NormalTok{(m }\SpecialCharTok{*}\NormalTok{ n, lambda), }\AttributeTok{nrow =}\NormalTok{ m)}

\CommentTok{\# soma \textasciitilde{} N(n·μ, n·σ²)}
\NormalTok{soma }\OtherTok{\textless{}{-}} \FunctionTok{rowSums}\NormalTok{(amostras)}
\FunctionTok{hist}\NormalTok{(soma, }\AttributeTok{probability =} \ConstantTok{TRUE}\NormalTok{, }\AttributeTok{col =} \StringTok{"lightblue"}\NormalTok{)}
\FunctionTok{curve}\NormalTok{(}\FunctionTok{dnorm}\NormalTok{(x, n }\SpecialCharTok{*}\NormalTok{ media\_teo, }\FunctionTok{sqrt}\NormalTok{(var\_teo }\SpecialCharTok{*}\NormalTok{ n)), }\AttributeTok{col =} \StringTok{"red"}\NormalTok{, }\AttributeTok{lwd =} \DecValTok{2}\NormalTok{, }\AttributeTok{add =} \ConstantTok{TRUE}\NormalTok{)}

\CommentTok{\# soma padronizada \textasciitilde{} N(0,1)}
\NormalTok{z }\OtherTok{\textless{}{-}}\NormalTok{ (soma }\SpecialCharTok{{-}} \FunctionTok{mean}\NormalTok{(soma)) }\SpecialCharTok{/} \FunctionTok{sd}\NormalTok{(soma)}
\FunctionTok{hist}\NormalTok{(z, }\AttributeTok{probability =} \ConstantTok{TRUE}\NormalTok{, }\AttributeTok{col =} \StringTok{"lightblue"}\NormalTok{)}
\FunctionTok{curve}\NormalTok{(}\FunctionTok{dnorm}\NormalTok{(x), }\AttributeTok{col =} \StringTok{"red"}\NormalTok{, }\AttributeTok{lwd =} \DecValTok{2}\NormalTok{, }\AttributeTok{add =} \ConstantTok{TRUE}\NormalTok{)}
\end{Highlighting}
\end{Shaded}

\section{Método da transformação inversa (caso discreto)}\label{muxe9todo-da-transformauxe7uxe3o-inversa-caso-discreto}

A ideia é dividir o intervalo \([0,1]\) em subintervalos com comprimentos iguais
às probabilidades da distribuição; um \(U \sim \mathcal{U}(0,1)\) cai num deles com
a probabilidade correta.

\textbf{1 a 3.} Variáveis discretas a partir da função de distribuição acumulada.

\begin{Shaded}
\begin{Highlighting}[]
\FunctionTok{set.seed}\NormalTok{(}\DecValTok{1}\NormalTok{)}
\CommentTok{\# distribuição com P(X=1)=0.2, P(X=2)=0.5, P(X=3)=0.3  =\textgreater{}  F = (0.2, 0.7, 1.0)}
\NormalTok{gerar }\OtherTok{\textless{}{-}} \ControlFlowTok{function}\NormalTok{(n, valores, prob) \{}
\NormalTok{  Fx }\OtherTok{\textless{}{-}} \FunctionTok{cumsum}\NormalTok{(prob)}
\NormalTok{  u }\OtherTok{\textless{}{-}} \FunctionTok{runif}\NormalTok{(n)}
\NormalTok{  valores[}\FunctionTok{findInterval}\NormalTok{(u, Fx) }\SpecialCharTok{+} \DecValTok{1}\NormalTok{]   }\CommentTok{\# findInterval localiza o subintervalo}
\NormalTok{\}}
\NormalTok{sim }\OtherTok{\textless{}{-}} \FunctionTok{gerar}\NormalTok{(}\DecValTok{1000}\NormalTok{, }\AttributeTok{valores =} \DecValTok{1}\SpecialCharTok{:}\DecValTok{3}\NormalTok{, }\AttributeTok{prob =} \FunctionTok{c}\NormalTok{(}\FloatTok{0.2}\NormalTok{, }\FloatTok{0.5}\NormalTok{, }\FloatTok{0.3}\NormalTok{))}
\FunctionTok{table}\NormalTok{(sim) }\SpecialCharTok{/} \DecValTok{1000}
\end{Highlighting}
\end{Shaded}

Em vez de um ciclo \texttt{for} com uma cadeia de \texttt{if}/\texttt{else\ if} para localizar o
subintervalo (como no material original), \texttt{findInterval(u,\ Fx)} faz o mesmo de
forma \textbf{vetorizada} e para qualquer número de categorias. Além de mais curto, é
muito mais rápido para amostras grandes.

\textbf{4.} Geométrica pela transformação inversa. Como \(F(x) = 1 - (1-p)^x\), a inversa
dá \(x = \big\lceil \ln(1-u) / \ln(1-p) \big\rceil\).

\begin{Shaded}
\begin{Highlighting}[]
\FunctionTok{set.seed}\NormalTok{(}\DecValTok{1}\NormalTok{)}
\NormalTok{p }\OtherTok{\textless{}{-}} \FloatTok{0.2}\NormalTok{; n }\OtherTok{\textless{}{-}} \DecValTok{1000}
\NormalTok{u }\OtherTok{\textless{}{-}} \FunctionTok{runif}\NormalTok{(n)}
\NormalTok{x }\OtherTok{\textless{}{-}} \FunctionTok{ceiling}\NormalTok{(}\FunctionTok{log}\NormalTok{(}\DecValTok{1} \SpecialCharTok{{-}}\NormalTok{ u) }\SpecialCharTok{/} \FunctionTok{log}\NormalTok{(}\DecValTok{1} \SpecialCharTok{{-}}\NormalTok{ p))   }\CommentTok{\# nº de provas até ao 1.º sucesso}
\FunctionTok{table}\NormalTok{(x) }\SpecialCharTok{/}\NormalTok{ n}
\end{Highlighting}
\end{Shaded}

\section{Método da transformação inversa (caso contínuo)}\label{muxe9todo-da-transformauxe7uxe3o-inversa-caso-contuxednuo}

Se \(F\) é a função de distribuição e \(U \sim \mathcal{U}(0,1)\), então
\(X = F^{-1}(U)\) tem função de distribuição \(F\). Basta, pois, inverter \(F\).

\textbf{1.} Exponencial: \(F(x) = 1 - e^{-\lambda x} \Rightarrow F^{-1}(u) = -\ln(1-u)/\lambda\).

\begin{Shaded}
\begin{Highlighting}[]
\NormalTok{gerar\_exp }\OtherTok{\textless{}{-}} \ControlFlowTok{function}\NormalTok{(n, lambda) }\SpecialCharTok{{-}}\FunctionTok{log}\NormalTok{(}\DecValTok{1} \SpecialCharTok{{-}} \FunctionTok{runif}\NormalTok{(n)) }\SpecialCharTok{/}\NormalTok{ lambda}
\FunctionTok{set.seed}\NormalTok{(}\DecValTok{123}\NormalTok{)}
\NormalTok{x }\OtherTok{\textless{}{-}} \FunctionTok{gerar\_exp}\NormalTok{(}\DecValTok{1000}\NormalTok{, }\DecValTok{2}\NormalTok{)}
\FunctionTok{hist}\NormalTok{(x, }\AttributeTok{probability =} \ConstantTok{TRUE}\NormalTok{, }\AttributeTok{col =} \StringTok{"lightblue"}\NormalTok{)}
\FunctionTok{curve}\NormalTok{(}\FunctionTok{dexp}\NormalTok{(x, }\DecValTok{2}\NormalTok{), }\AttributeTok{col =} \StringTok{"red"}\NormalTok{, }\AttributeTok{lwd =} \DecValTok{2}\NormalTok{, }\AttributeTok{add =} \ConstantTok{TRUE}\NormalTok{)}
\end{Highlighting}
\end{Shaded}

\textbf{2.} Uniforme \(\mathcal{U}(a,b)\): \(F^{-1}(u) = a + (b-a)u\).

\begin{Shaded}
\begin{Highlighting}[]
\NormalTok{gerar\_unif }\OtherTok{\textless{}{-}} \ControlFlowTok{function}\NormalTok{(n, a, b) a }\SpecialCharTok{+}\NormalTok{ (b }\SpecialCharTok{{-}}\NormalTok{ a) }\SpecialCharTok{*} \FunctionTok{runif}\NormalTok{(n)}
\end{Highlighting}
\end{Shaded}

\textbf{3.} Cauchy: \(F^{-1}(u) = x_0 + \gamma \tan\!\big(\pi(u - 1/2)\big)\).

\begin{Shaded}
\begin{Highlighting}[]
\NormalTok{gerar\_cauchy }\OtherTok{\textless{}{-}} \ControlFlowTok{function}\NormalTok{(n, }\AttributeTok{x0 =} \DecValTok{0}\NormalTok{, }\AttributeTok{gama =} \DecValTok{1}\NormalTok{) x0 }\SpecialCharTok{+}\NormalTok{ gama }\SpecialCharTok{*} \FunctionTok{tan}\NormalTok{(pi }\SpecialCharTok{*}\NormalTok{ (}\FunctionTok{runif}\NormalTok{(n) }\SpecialCharTok{{-}} \FloatTok{0.5}\NormalTok{))}
\FunctionTok{set.seed}\NormalTok{(}\DecValTok{1}\NormalTok{)}
\NormalTok{x }\OtherTok{\textless{}{-}} \FunctionTok{gerar\_cauchy}\NormalTok{(}\DecValTok{10000}\NormalTok{)}
\FunctionTok{hist}\NormalTok{(x, }\AttributeTok{probability =} \ConstantTok{TRUE}\NormalTok{, }\AttributeTok{breaks =} \DecValTok{100}\NormalTok{, }\AttributeTok{xlim =} \FunctionTok{c}\NormalTok{(}\SpecialCharTok{{-}}\DecValTok{10}\NormalTok{, }\DecValTok{10}\NormalTok{), }\AttributeTok{col =} \StringTok{"lightblue"}\NormalTok{)}
\FunctionTok{curve}\NormalTok{(}\FunctionTok{dcauchy}\NormalTok{(x), }\AttributeTok{col =} \StringTok{"red"}\NormalTok{, }\AttributeTok{lwd =} \DecValTok{2}\NormalTok{, }\AttributeTok{add =} \ConstantTok{TRUE}\NormalTok{)}
\end{Highlighting}
\end{Shaded}

\textbf{4.} Weibull: \(F(x) = 1 - e^{-(x/\beta)^\alpha} \Rightarrow
F^{-1}(u) = \beta\,\big(-\ln(1-u)\big)^{1/\alpha}\).

\begin{Shaded}
\begin{Highlighting}[]
\NormalTok{gerar\_weibull }\OtherTok{\textless{}{-}} \ControlFlowTok{function}\NormalTok{(n, alfa, beta) beta }\SpecialCharTok{*}\NormalTok{ (}\SpecialCharTok{{-}}\FunctionTok{log}\NormalTok{(}\DecValTok{1} \SpecialCharTok{{-}} \FunctionTok{runif}\NormalTok{(n)))}\SpecialCharTok{\^{}}\NormalTok{(}\DecValTok{1} \SpecialCharTok{/}\NormalTok{ alfa)}
\FunctionTok{set.seed}\NormalTok{(}\DecValTok{1}\NormalTok{)}
\NormalTok{x }\OtherTok{\textless{}{-}} \FunctionTok{gerar\_weibull}\NormalTok{(}\DecValTok{10000}\NormalTok{, }\AttributeTok{alfa =} \DecValTok{2}\NormalTok{, }\AttributeTok{beta =} \DecValTok{3}\NormalTok{)}
\FunctionTok{hist}\NormalTok{(x, }\AttributeTok{probability =} \ConstantTok{TRUE}\NormalTok{, }\AttributeTok{col =} \StringTok{"lightblue"}\NormalTok{)}
\FunctionTok{curve}\NormalTok{(}\FunctionTok{dweibull}\NormalTok{(x, }\DecValTok{2}\NormalTok{, }\DecValTok{3}\NormalTok{), }\AttributeTok{col =} \StringTok{"red"}\NormalTok{, }\AttributeTok{lwd =} \DecValTok{2}\NormalTok{, }\AttributeTok{add =} \ConstantTok{TRUE}\NormalTok{)}
\end{Highlighting}
\end{Shaded}

\section{Método da aceitação-rejeição}\label{muxe9todo-da-aceitauxe7uxe3o-rejeiuxe7uxe3o}

\textbf{A constante \(c\) tem de ser o supremo verdadeiro de \(f/g\).} Vários dos
exercícios seguintes revelam o erro mais perigoso deste método: escolher um \(c\)
demasiado pequeno (ou uma proposta cujo rácio \(f/g\) é ilimitado). Nesses casos o
algoritmo produz observações de uma distribuição \textbf{errada}, sem qualquer aviso.
Determinamos sempre \(c\) analiticamente --- ou com \texttt{optimize()} --- e nunca ``à vista''.

\textbf{1.} Exponencial \(\text{Exp}(1)\) com proposta uniforme em \([0,5]\). Como
\(f(x) = e^{-x}\) e \(g(x) = 1/5\), o rácio \(f/g = 5e^{-x}\) é máximo em \(x = 0\), logo
\(c = 5\). (Gera-se, na verdade, uma exponencial \textbf{truncada} em \(5\); como
\(P(X > 5) = e^{-5} \approx 0{,}007\), a diferença é desprezável.)

\begin{Shaded}
\begin{Highlighting}[]
\FunctionTok{set.seed}\NormalTok{(}\DecValTok{123}\NormalTok{)}
\NormalTok{n }\OtherTok{\textless{}{-}} \DecValTok{1000}\NormalTok{; amostra }\OtherTok{\textless{}{-}} \FunctionTok{numeric}\NormalTok{(}\DecValTok{0}\NormalTok{)}
\ControlFlowTok{while}\NormalTok{ (}\FunctionTok{length}\NormalTok{(amostra) }\SpecialCharTok{\textless{}}\NormalTok{ n) \{}
\NormalTok{  y }\OtherTok{\textless{}{-}} \FunctionTok{runif}\NormalTok{(}\DecValTok{1}\NormalTok{, }\DecValTok{0}\NormalTok{, }\DecValTok{5}\NormalTok{); u }\OtherTok{\textless{}{-}} \FunctionTok{runif}\NormalTok{(}\DecValTok{1}\NormalTok{)}
  \ControlFlowTok{if}\NormalTok{ (u }\SpecialCharTok{\textless{}=} \FunctionTok{exp}\NormalTok{(}\SpecialCharTok{{-}}\NormalTok{y) }\SpecialCharTok{/}\NormalTok{ (}\DecValTok{5} \SpecialCharTok{*}\NormalTok{ (}\DecValTok{1}\SpecialCharTok{/}\DecValTok{5}\NormalTok{))) amostra }\OtherTok{\textless{}{-}} \FunctionTok{c}\NormalTok{(amostra, y)  }\CommentTok{\# aceita se u \textless{}= e\^{}\{{-}y\}}
\NormalTok{\}}
\FunctionTok{hist}\NormalTok{(amostra, }\AttributeTok{probability =} \ConstantTok{TRUE}\NormalTok{, }\AttributeTok{breaks =} \DecValTok{30}\NormalTok{, }\AttributeTok{col =} \StringTok{"lightblue"}\NormalTok{)}
\FunctionTok{curve}\NormalTok{(}\FunctionTok{dexp}\NormalTok{(x, }\DecValTok{1}\NormalTok{), }\AttributeTok{col =} \StringTok{"red"}\NormalTok{, }\AttributeTok{lwd =} \DecValTok{2}\NormalTok{, }\AttributeTok{add =} \ConstantTok{TRUE}\NormalTok{)}
\end{Highlighting}
\end{Shaded}

\textbf{2.} Normal padrão com proposta Cauchy. O rácio \(f/g = \text{dnorm}/\text{dcauchy}\)
é máximo em \(x = \pm 1\) (não em \(0\)), onde vale
\[c = \sqrt{\tfrac{\pi}{2}}\;\cdot 2\,e^{-1/2} \approx 1{,}520.\]

\begin{Shaded}
\begin{Highlighting}[]
\FunctionTok{set.seed}\NormalTok{(}\DecValTok{123}\NormalTok{)}
\NormalTok{c }\OtherTok{\textless{}{-}} \FunctionTok{optimize}\NormalTok{(}\ControlFlowTok{function}\NormalTok{(x) }\FunctionTok{dnorm}\NormalTok{(x) }\SpecialCharTok{/} \FunctionTok{dcauchy}\NormalTok{(x), }\FunctionTok{c}\NormalTok{(}\DecValTok{0}\NormalTok{, }\DecValTok{5}\NormalTok{), }\AttributeTok{maximum =} \ConstantTok{TRUE}\NormalTok{)}\SpecialCharTok{$}\NormalTok{objective}
\NormalTok{n }\OtherTok{\textless{}{-}} \DecValTok{1000}\NormalTok{; amostra }\OtherTok{\textless{}{-}} \FunctionTok{numeric}\NormalTok{(}\DecValTok{0}\NormalTok{)}
\ControlFlowTok{while}\NormalTok{ (}\FunctionTok{length}\NormalTok{(amostra) }\SpecialCharTok{\textless{}}\NormalTok{ n) \{}
\NormalTok{  y }\OtherTok{\textless{}{-}} \FunctionTok{rcauchy}\NormalTok{(}\DecValTok{1}\NormalTok{); u }\OtherTok{\textless{}{-}} \FunctionTok{runif}\NormalTok{(}\DecValTok{1}\NormalTok{)}
  \ControlFlowTok{if}\NormalTok{ (u }\SpecialCharTok{\textless{}=} \FunctionTok{dnorm}\NormalTok{(y) }\SpecialCharTok{/}\NormalTok{ (c }\SpecialCharTok{*} \FunctionTok{dcauchy}\NormalTok{(y))) amostra }\OtherTok{\textless{}{-}} \FunctionTok{c}\NormalTok{(amostra, y)}
\NormalTok{\}}
\FunctionTok{hist}\NormalTok{(amostra, }\AttributeTok{probability =} \ConstantTok{TRUE}\NormalTok{, }\AttributeTok{breaks =} \DecValTok{30}\NormalTok{, }\AttributeTok{col =} \StringTok{"lightgreen"}\NormalTok{)}
\FunctionTok{curve}\NormalTok{(}\FunctionTok{dnorm}\NormalTok{(x), }\AttributeTok{col =} \StringTok{"blue"}\NormalTok{, }\AttributeTok{lwd =} \DecValTok{2}\NormalTok{, }\AttributeTok{add =} \ConstantTok{TRUE}\NormalTok{)}
\end{Highlighting}
\end{Shaded}

O material original usava \(c = 1{,}5\) para este exemplo. Mas o verdadeiro
supremo é \(\approx 1{,}520\): com \(c = 1{,}5\), a envolvente \(c\,g\) fica \emph{abaixo}
de \(f\) na vizinhança de \(x = \pm 1\), e as observações geradas seguem uma
distribuição enviesada. É exatamente o erro silencioso contra o qual avisámos.

\textbf{3.} Beta\((2,3)\) com proposta uniforme. O máximo de \(\text{dbeta}(x;2,3)\)
ocorre em \(x = 1/3\), onde \(c = \text{dbeta}(1/3;2,3) = 16/9 \approx 1{,}778\).

\begin{Shaded}
\begin{Highlighting}[]
\FunctionTok{set.seed}\NormalTok{(}\DecValTok{123}\NormalTok{)}
\NormalTok{c }\OtherTok{\textless{}{-}} \FunctionTok{dbeta}\NormalTok{(}\DecValTok{1}\SpecialCharTok{/}\DecValTok{3}\NormalTok{, }\DecValTok{2}\NormalTok{, }\DecValTok{3}\NormalTok{)          }\CommentTok{\# = 16/9, o máximo exato (não dbeta(0.4, ...))}
\NormalTok{n }\OtherTok{\textless{}{-}} \DecValTok{1000}\NormalTok{; amostra }\OtherTok{\textless{}{-}} \FunctionTok{numeric}\NormalTok{(}\DecValTok{0}\NormalTok{)}
\ControlFlowTok{while}\NormalTok{ (}\FunctionTok{length}\NormalTok{(amostra) }\SpecialCharTok{\textless{}}\NormalTok{ n) \{}
\NormalTok{  y }\OtherTok{\textless{}{-}} \FunctionTok{runif}\NormalTok{(}\DecValTok{1}\NormalTok{); u }\OtherTok{\textless{}{-}} \FunctionTok{runif}\NormalTok{(}\DecValTok{1}\NormalTok{)}
  \ControlFlowTok{if}\NormalTok{ (u }\SpecialCharTok{\textless{}=} \FunctionTok{dbeta}\NormalTok{(y, }\DecValTok{2}\NormalTok{, }\DecValTok{3}\NormalTok{) }\SpecialCharTok{/}\NormalTok{ c) amostra }\OtherTok{\textless{}{-}} \FunctionTok{c}\NormalTok{(amostra, y)}
\NormalTok{\}}
\FunctionTok{hist}\NormalTok{(amostra, }\AttributeTok{probability =} \ConstantTok{TRUE}\NormalTok{, }\AttributeTok{breaks =} \DecValTok{30}\NormalTok{, }\AttributeTok{col =} \StringTok{"lightcoral"}\NormalTok{)}
\FunctionTok{curve}\NormalTok{(}\FunctionTok{dbeta}\NormalTok{(x, }\DecValTok{2}\NormalTok{, }\DecValTok{3}\NormalTok{), }\AttributeTok{col =} \StringTok{"purple"}\NormalTok{, }\AttributeTok{lwd =} \DecValTok{2}\NormalTok{, }\AttributeTok{add =} \ConstantTok{TRUE}\NormalTok{)}
\end{Highlighting}
\end{Shaded}

\textbf{4.} Gama\((3, 2)\) com proposta exponencial. Aqui há uma subtileza importante:
com \(\text{Exp}(2)\), o rácio
\[\frac{\text{dgamma}(x;3,2)}{\text{dexp}(x;2)} = \frac{4x^2 e^{-2x}}{2e^{-2x}} = 2x^2\]
é \textbf{ilimitado} --- não existe \(c\) finito, e o método é inaplicável. A proposta
tem de decair mais devagar do que a alvo, ou seja, ter taxa \textbf{inferior} a \(2\). A
escolha ótima iguala as médias: \(\text{Exp}(2/3)\), para a qual \(c \approx 1{,}83\).

\begin{Shaded}
\begin{Highlighting}[]
\FunctionTok{set.seed}\NormalTok{(}\DecValTok{123}\NormalTok{)}
\NormalTok{alfa }\OtherTok{\textless{}{-}} \DecValTok{3}\NormalTok{; beta }\OtherTok{\textless{}{-}} \DecValTok{2}
\NormalTok{lambda\_prop }\OtherTok{\textless{}{-}}\NormalTok{ beta }\SpecialCharTok{/}\NormalTok{ alfa          }\CommentTok{\# = 2/3 : proposta com a mesma média da alvo}
\NormalTok{c }\OtherTok{\textless{}{-}} \FunctionTok{optimize}\NormalTok{(}\ControlFlowTok{function}\NormalTok{(x) }\FunctionTok{dgamma}\NormalTok{(x, alfa, }\AttributeTok{rate =}\NormalTok{ beta) }\SpecialCharTok{/} \FunctionTok{dexp}\NormalTok{(x, lambda\_prop),}
              \FunctionTok{c}\NormalTok{(}\FloatTok{1e{-}4}\NormalTok{, }\DecValTok{50}\NormalTok{), }\AttributeTok{maximum =} \ConstantTok{TRUE}\NormalTok{)}\SpecialCharTok{$}\NormalTok{objective}
\NormalTok{n }\OtherTok{\textless{}{-}} \DecValTok{1000}\NormalTok{; amostra }\OtherTok{\textless{}{-}} \FunctionTok{numeric}\NormalTok{(}\DecValTok{0}\NormalTok{)}
\ControlFlowTok{while}\NormalTok{ (}\FunctionTok{length}\NormalTok{(amostra) }\SpecialCharTok{\textless{}}\NormalTok{ n) \{}
\NormalTok{  y }\OtherTok{\textless{}{-}} \FunctionTok{rexp}\NormalTok{(}\DecValTok{1}\NormalTok{, lambda\_prop); u }\OtherTok{\textless{}{-}} \FunctionTok{runif}\NormalTok{(}\DecValTok{1}\NormalTok{)}
  \ControlFlowTok{if}\NormalTok{ (u }\SpecialCharTok{\textless{}=} \FunctionTok{dgamma}\NormalTok{(y, alfa, }\AttributeTok{rate =}\NormalTok{ beta) }\SpecialCharTok{/}\NormalTok{ (c }\SpecialCharTok{*} \FunctionTok{dexp}\NormalTok{(y, lambda\_prop)))}
\NormalTok{    amostra }\OtherTok{\textless{}{-}} \FunctionTok{c}\NormalTok{(amostra, y)}
\NormalTok{\}}
\FunctionTok{hist}\NormalTok{(amostra, }\AttributeTok{probability =} \ConstantTok{TRUE}\NormalTok{, }\AttributeTok{breaks =} \DecValTok{30}\NormalTok{, }\AttributeTok{col =} \StringTok{"lightblue"}\NormalTok{)}
\FunctionTok{curve}\NormalTok{(}\FunctionTok{dgamma}\NormalTok{(x, alfa, }\AttributeTok{rate =}\NormalTok{ beta), }\AttributeTok{col =} \StringTok{"red"}\NormalTok{, }\AttributeTok{lwd =} \DecValTok{2}\NormalTok{, }\AttributeTok{add =} \ConstantTok{TRUE}\NormalTok{)}
\end{Highlighting}
\end{Shaded}

\textbf{5.} Log-normal com proposta normal. Também aqui a proposta é \textbf{inadmissível}:
a cauda da log-normal é mais pesada do que a da normal, pelo que o rácio \(f/g\) é
ilimitado. A solução elegante não é rejeitar de todo, mas usar a definição: se
\(Z \sim N(0,1)\), então \(e^Z \sim \text{LogNormal}(0,1)\). Basta exponenciar.

\begin{Shaded}
\begin{Highlighting}[]
\FunctionTok{set.seed}\NormalTok{(}\DecValTok{123}\NormalTok{)}
\NormalTok{amostra }\OtherTok{\textless{}{-}} \FunctionTok{exp}\NormalTok{(}\FunctionTok{rnorm}\NormalTok{(}\DecValTok{1000}\NormalTok{))        }\CommentTok{\# gera log{-}normal sem qualquer rejeição}
\FunctionTok{hist}\NormalTok{(amostra, }\AttributeTok{probability =} \ConstantTok{TRUE}\NormalTok{, }\AttributeTok{breaks =} \DecValTok{40}\NormalTok{, }\AttributeTok{col =} \StringTok{"lightyellow"}\NormalTok{)}
\FunctionTok{curve}\NormalTok{(}\FunctionTok{dlnorm}\NormalTok{(x), }\AttributeTok{col =} \StringTok{"darkred"}\NormalTok{, }\AttributeTok{lwd =} \DecValTok{2}\NormalTok{, }\AttributeTok{add =} \ConstantTok{TRUE}\NormalTok{)}
\end{Highlighting}
\end{Shaded}

\textbf{6.} Weibull\((\text{forma}=2, \text{escala}=1)\) com proposta \(\text{Exp}(1)\). O
rácio \(2x\,e^{-x^2 + x}\) é máximo em \(x = 1\), onde vale exatamente \(c = 2\).

\begin{Shaded}
\begin{Highlighting}[]
\FunctionTok{set.seed}\NormalTok{(}\DecValTok{123}\NormalTok{)}
\NormalTok{c }\OtherTok{\textless{}{-}} \DecValTok{2}                             \CommentTok{\# valor exato (o material usava 1.5, pequeno demais)}
\NormalTok{n }\OtherTok{\textless{}{-}} \DecValTok{1000}\NormalTok{; amostra }\OtherTok{\textless{}{-}} \FunctionTok{numeric}\NormalTok{(}\DecValTok{0}\NormalTok{)}
\ControlFlowTok{while}\NormalTok{ (}\FunctionTok{length}\NormalTok{(amostra) }\SpecialCharTok{\textless{}}\NormalTok{ n) \{}
\NormalTok{  y }\OtherTok{\textless{}{-}} \FunctionTok{rexp}\NormalTok{(}\DecValTok{1}\NormalTok{, }\DecValTok{1}\NormalTok{); u }\OtherTok{\textless{}{-}} \FunctionTok{runif}\NormalTok{(}\DecValTok{1}\NormalTok{)}
  \ControlFlowTok{if}\NormalTok{ (u }\SpecialCharTok{\textless{}=} \FunctionTok{dweibull}\NormalTok{(y, }\DecValTok{2}\NormalTok{, }\DecValTok{1}\NormalTok{) }\SpecialCharTok{/}\NormalTok{ (c }\SpecialCharTok{*} \FunctionTok{dexp}\NormalTok{(y, }\DecValTok{1}\NormalTok{))) amostra }\OtherTok{\textless{}{-}} \FunctionTok{c}\NormalTok{(amostra, y)}
\NormalTok{\}}
\FunctionTok{hist}\NormalTok{(amostra, }\AttributeTok{probability =} \ConstantTok{TRUE}\NormalTok{, }\AttributeTok{breaks =} \DecValTok{30}\NormalTok{, }\AttributeTok{col =} \StringTok{"lightpink"}\NormalTok{)}
\FunctionTok{curve}\NormalTok{(}\FunctionTok{dweibull}\NormalTok{(x, }\DecValTok{2}\NormalTok{, }\DecValTok{1}\NormalTok{), }\AttributeTok{col =} \StringTok{"purple"}\NormalTok{, }\AttributeTok{lwd =} \DecValTok{2}\NormalTok{, }\AttributeTok{add =} \ConstantTok{TRUE}\NormalTok{)}
\end{Highlighting}
\end{Shaded}

\textbf{7.} Densidade \(f(x) = \tfrac12(1 + \cos x)\) em \([-\pi, \pi]\), com proposta
uniforme \(g = 1/(2\pi)\). O rácio \(f/g = \pi(1 + \cos x)\) é máximo em \(x = 0\),
onde \(c = 2\pi\).

\begin{Shaded}
\begin{Highlighting}[]
\FunctionTok{set.seed}\NormalTok{(}\DecValTok{1}\NormalTok{)}
\NormalTok{f }\OtherTok{\textless{}{-}} \ControlFlowTok{function}\NormalTok{(x) }\FloatTok{0.5} \SpecialCharTok{*}\NormalTok{ (}\DecValTok{1} \SpecialCharTok{+} \FunctionTok{cos}\NormalTok{(x))}
\NormalTok{c }\OtherTok{\textless{}{-}} \DecValTok{2} \SpecialCharTok{*}\NormalTok{ pi}
\NormalTok{n }\OtherTok{\textless{}{-}} \DecValTok{5000}\NormalTok{; amostra }\OtherTok{\textless{}{-}} \FunctionTok{numeric}\NormalTok{(}\DecValTok{0}\NormalTok{)}
\ControlFlowTok{while}\NormalTok{ (}\FunctionTok{length}\NormalTok{(amostra) }\SpecialCharTok{\textless{}}\NormalTok{ n) \{}
\NormalTok{  y }\OtherTok{\textless{}{-}} \FunctionTok{runif}\NormalTok{(}\DecValTok{1}\NormalTok{, }\SpecialCharTok{{-}}\NormalTok{pi, pi); u }\OtherTok{\textless{}{-}} \FunctionTok{runif}\NormalTok{(}\DecValTok{1}\NormalTok{)}
  \ControlFlowTok{if}\NormalTok{ (u }\SpecialCharTok{\textless{}=} \FunctionTok{f}\NormalTok{(y) }\SpecialCharTok{/}\NormalTok{ (c }\SpecialCharTok{*}\NormalTok{ (}\DecValTok{1} \SpecialCharTok{/}\NormalTok{ (}\DecValTok{2} \SpecialCharTok{*}\NormalTok{ pi)))) amostra }\OtherTok{\textless{}{-}} \FunctionTok{c}\NormalTok{(amostra, y)}
\NormalTok{\}}
\FunctionTok{hist}\NormalTok{(amostra, }\AttributeTok{probability =} \ConstantTok{TRUE}\NormalTok{, }\AttributeTok{breaks =} \DecValTok{40}\NormalTok{, }\AttributeTok{col =} \StringTok{"lightblue"}\NormalTok{)}
\FunctionTok{curve}\NormalTok{(f, }\SpecialCharTok{{-}}\NormalTok{pi, pi, }\AttributeTok{col =} \StringTok{"red"}\NormalTok{, }\AttributeTok{lwd =} \DecValTok{2}\NormalTok{, }\AttributeTok{add =} \ConstantTok{TRUE}\NormalTok{)}
\end{Highlighting}
\end{Shaded}

\textbf{8.} Versão vetorizada e comparação de tempos.

\begin{Shaded}
\begin{Highlighting}[]
\NormalTok{ar\_vetorizado }\OtherTok{\textless{}{-}} \ControlFlowTok{function}\NormalTok{(n, f, c) \{           }\CommentTok{\# proposta uniforme(0,1)}
\NormalTok{  out }\OtherTok{\textless{}{-}} \FunctionTok{numeric}\NormalTok{(}\DecValTok{0}\NormalTok{)}
  \ControlFlowTok{while}\NormalTok{ (}\FunctionTok{length}\NormalTok{(out) }\SpecialCharTok{\textless{}}\NormalTok{ n) \{}
\NormalTok{    y }\OtherTok{\textless{}{-}} \FunctionTok{runif}\NormalTok{(}\DecValTok{2} \SpecialCharTok{*}\NormalTok{ n); u }\OtherTok{\textless{}{-}} \FunctionTok{runif}\NormalTok{(}\DecValTok{2} \SpecialCharTok{*}\NormalTok{ n)}
\NormalTok{    out }\OtherTok{\textless{}{-}} \FunctionTok{c}\NormalTok{(out, y[u }\SpecialCharTok{\textless{}=} \FunctionTok{f}\NormalTok{(y) }\SpecialCharTok{/}\NormalTok{ c])}
\NormalTok{  \}}
\NormalTok{  out[}\DecValTok{1}\SpecialCharTok{:}\NormalTok{n]}
\NormalTok{\}}
\NormalTok{f\_beta }\OtherTok{\textless{}{-}} \ControlFlowTok{function}\NormalTok{(x) }\FunctionTok{dbeta}\NormalTok{(x, }\DecValTok{2}\NormalTok{, }\DecValTok{2}\NormalTok{)}
\FunctionTok{system.time}\NormalTok{(}\FunctionTok{ar\_vetorizado}\NormalTok{(}\FloatTok{1e5}\NormalTok{, f\_beta, }\AttributeTok{c =} \FloatTok{1.5}\NormalTok{))   }\CommentTok{\# compare com a versão em while}
\end{Highlighting}
\end{Shaded}

\textbf{9.} Estudo da eficiência. Com o \(c\) correto, a taxa de aceitação empírica
aproxima-se de \(1/c\). Com um \(c\) \textbf{maior} do que o necessário, a distribuição
gerada continua correta, mas a taxa de aceitação cai (desperdício). Com um \(c\)
\textbf{menor} do que o supremo, a distribuição fica \textbf{errada} --- o caso a evitar.

\begin{Shaded}
\begin{Highlighting}[]
\NormalTok{taxa }\OtherTok{\textless{}{-}} \ControlFlowTok{function}\NormalTok{(c) \{}
\NormalTok{  aceites }\OtherTok{\textless{}{-}} \DecValTok{0}\NormalTok{; tentativas }\OtherTok{\textless{}{-}} \DecValTok{0}
  \ControlFlowTok{while}\NormalTok{ (aceites }\SpecialCharTok{\textless{}} \DecValTok{5000}\NormalTok{) \{}
\NormalTok{    y }\OtherTok{\textless{}{-}} \FunctionTok{rexp}\NormalTok{(}\DecValTok{1}\NormalTok{); u }\OtherTok{\textless{}{-}} \FunctionTok{runif}\NormalTok{(}\DecValTok{1}\NormalTok{); tentativas }\OtherTok{\textless{}{-}}\NormalTok{ tentativas }\SpecialCharTok{+} \DecValTok{1}
    \ControlFlowTok{if}\NormalTok{ (u }\SpecialCharTok{\textless{}=} \FunctionTok{sqrt}\NormalTok{(}\DecValTok{2}\SpecialCharTok{/}\NormalTok{pi) }\SpecialCharTok{*} \FunctionTok{exp}\NormalTok{(}\SpecialCharTok{{-}}\NormalTok{y}\SpecialCharTok{\^{}}\DecValTok{2}\SpecialCharTok{/}\DecValTok{2}\NormalTok{) }\SpecialCharTok{/}\NormalTok{ (c }\SpecialCharTok{*} \FunctionTok{dexp}\NormalTok{(y))) aceites }\OtherTok{\textless{}{-}}\NormalTok{ aceites }\SpecialCharTok{+} \DecValTok{1}
\NormalTok{  \}}
\NormalTok{  aceites }\SpecialCharTok{/}\NormalTok{ tentativas}
\NormalTok{\}}
\NormalTok{c\_correto }\OtherTok{\textless{}{-}} \FunctionTok{sqrt}\NormalTok{(}\DecValTok{2} \SpecialCharTok{*} \FunctionTok{exp}\NormalTok{(}\DecValTok{1}\NormalTok{) }\SpecialCharTok{/}\NormalTok{ pi)   }\CommentTok{\# \textasciitilde{} 1.3155}
\FunctionTok{c}\NormalTok{(}\FunctionTok{taxa}\NormalTok{(c\_correto), }\DecValTok{1} \SpecialCharTok{/}\NormalTok{ c\_correto,    }\CommentTok{\# empírica \textasciitilde{} teórica}
  \FunctionTok{taxa}\NormalTok{(}\DecValTok{2} \SpecialCharTok{*}\NormalTok{ c\_correto))               }\CommentTok{\# c grande demais: taxa cai para metade}
\end{Highlighting}
\end{Shaded}

\textbf{10.} Distribuição do semicírculo \(f(x) = \tfrac{2}{\pi}\sqrt{1 - x^2}\) em
\([-1, 1]\), com proposta uniforme \(g = 1/2\). O rácio \(f/g = \tfrac{4}{\pi}\sqrt{1-x^2}\)
é máximo em \(x = 0\), onde \(c = 4/\pi \approx 1{,}273\).

\begin{Shaded}
\begin{Highlighting}[]
\FunctionTok{set.seed}\NormalTok{(}\DecValTok{1}\NormalTok{)}
\NormalTok{f }\OtherTok{\textless{}{-}} \ControlFlowTok{function}\NormalTok{(x) (}\DecValTok{2} \SpecialCharTok{/}\NormalTok{ pi) }\SpecialCharTok{*} \FunctionTok{sqrt}\NormalTok{(}\DecValTok{1} \SpecialCharTok{{-}}\NormalTok{ x}\SpecialCharTok{\^{}}\DecValTok{2}\NormalTok{)}
\NormalTok{c }\OtherTok{\textless{}{-}} \DecValTok{4} \SpecialCharTok{/}\NormalTok{ pi}
\NormalTok{n }\OtherTok{\textless{}{-}} \DecValTok{5000}\NormalTok{; amostra }\OtherTok{\textless{}{-}} \FunctionTok{numeric}\NormalTok{(}\DecValTok{0}\NormalTok{)}
\ControlFlowTok{while}\NormalTok{ (}\FunctionTok{length}\NormalTok{(amostra) }\SpecialCharTok{\textless{}}\NormalTok{ n) \{}
\NormalTok{  y }\OtherTok{\textless{}{-}} \FunctionTok{runif}\NormalTok{(}\DecValTok{1}\NormalTok{, }\SpecialCharTok{{-}}\DecValTok{1}\NormalTok{, }\DecValTok{1}\NormalTok{); u }\OtherTok{\textless{}{-}} \FunctionTok{runif}\NormalTok{(}\DecValTok{1}\NormalTok{)}
  \ControlFlowTok{if}\NormalTok{ (u }\SpecialCharTok{\textless{}=} \FunctionTok{f}\NormalTok{(y) }\SpecialCharTok{/}\NormalTok{ (c }\SpecialCharTok{*} \FloatTok{0.5}\NormalTok{)) amostra }\OtherTok{\textless{}{-}} \FunctionTok{c}\NormalTok{(amostra, y)}
\NormalTok{\}}
\FunctionTok{hist}\NormalTok{(amostra, }\AttributeTok{probability =} \ConstantTok{TRUE}\NormalTok{, }\AttributeTok{breaks =} \DecValTok{40}\NormalTok{, }\AttributeTok{col =} \StringTok{"lightblue"}\NormalTok{)}
\FunctionTok{curve}\NormalTok{(f, }\SpecialCharTok{{-}}\DecValTok{1}\NormalTok{, }\DecValTok{1}\NormalTok{, }\AttributeTok{col =} \StringTok{"red"}\NormalTok{, }\AttributeTok{lwd =} \DecValTok{2}\NormalTok{, }\AttributeTok{add =} \ConstantTok{TRUE}\NormalTok{)}
\end{Highlighting}
\end{Shaded}

\textbf{11.} Densidade não normalizada. Para \(\tilde f(x) = e^{-x^2/2}(1 + \sin^2 3x)\),
usamos uma proposta normal de variância grande. A taxa de aceitação estima a
constante de normalização: se aceitamos com probabilidade \(1/c\) relativamente a
\(\tilde f\), então \(\int \tilde f = c \cdot p_{\text{aceite}} \cdot \sqrt{2\pi}\,\sigma\)
(a menos da constante da proposta). Na prática, estima-se \(\int \tilde f\)
diretamente por Monte Carlo, \(\int \tilde f = E[\tilde f(X)/g(X)]\).

\begin{Shaded}
\begin{Highlighting}[]
\FunctionTok{set.seed}\NormalTok{(}\DecValTok{1}\NormalTok{)}
\NormalTok{f\_til }\OtherTok{\textless{}{-}} \ControlFlowTok{function}\NormalTok{(x) }\FunctionTok{exp}\NormalTok{(}\SpecialCharTok{{-}}\NormalTok{x}\SpecialCharTok{\^{}}\DecValTok{2}\SpecialCharTok{/}\DecValTok{2}\NormalTok{) }\SpecialCharTok{*}\NormalTok{ (}\DecValTok{1} \SpecialCharTok{+} \FunctionTok{sin}\NormalTok{(}\DecValTok{3}\SpecialCharTok{*}\NormalTok{x)}\SpecialCharTok{\^{}}\DecValTok{2}\NormalTok{)}
\NormalTok{X }\OtherTok{\textless{}{-}} \FunctionTok{rnorm}\NormalTok{(}\FloatTok{1e5}\NormalTok{, }\DecValTok{0}\NormalTok{, }\DecValTok{2}\NormalTok{)                       }\CommentTok{\# proposta larga}
\NormalTok{Z\_hat }\OtherTok{\textless{}{-}} \FunctionTok{mean}\NormalTok{(}\FunctionTok{f\_til}\NormalTok{(X) }\SpecialCharTok{/} \FunctionTok{dnorm}\NormalTok{(X, }\DecValTok{0}\NormalTok{, }\DecValTok{2}\NormalTok{))    }\CommentTok{\# estimativa de ∫ f\_til}
\NormalTok{Z\_hat}
\end{Highlighting}
\end{Shaded}

\section{Teoremas limites}\label{teoremas-limites-1}

Os exercícios deste capítulo são, na sua maioria, ilustrações por simulação.
Apresentamos as resoluções essenciais e o valor teórico para que serve cada uma.

\textbf{1. e 2.} Proporção acumulada de caras (moeda equilibrada e viciada).

\begin{Shaded}
\begin{Highlighting}[]
\FunctionTok{set.seed}\NormalTok{(}\DecValTok{1}\NormalTok{)}
\NormalTok{lancamentos }\OtherTok{\textless{}{-}} \FunctionTok{sample}\NormalTok{(}\FunctionTok{c}\NormalTok{(}\DecValTok{0}\NormalTok{, }\DecValTok{1}\NormalTok{), }\DecValTok{10000}\NormalTok{, }\AttributeTok{replace =} \ConstantTok{TRUE}\NormalTok{, }\AttributeTok{prob =} \FunctionTok{c}\NormalTok{(}\FloatTok{0.5}\NormalTok{, }\FloatTok{0.5}\NormalTok{))}
\FunctionTok{plot}\NormalTok{(}\FunctionTok{cumsum}\NormalTok{(lancamentos) }\SpecialCharTok{/} \FunctionTok{seq\_along}\NormalTok{(lancamentos), }\AttributeTok{type =} \StringTok{"l"}\NormalTok{,}
     \AttributeTok{ylim =} \FunctionTok{c}\NormalTok{(}\DecValTok{0}\NormalTok{, }\DecValTok{1}\NormalTok{), }\AttributeTok{xlab =} \StringTok{"lançamentos"}\NormalTok{, }\AttributeTok{ylab =} \StringTok{"proporção de caras"}\NormalTok{)}
\FunctionTok{abline}\NormalTok{(}\AttributeTok{h =} \FloatTok{0.5}\NormalTok{, }\AttributeTok{col =} \StringTok{"red"}\NormalTok{, }\AttributeTok{lty =} \DecValTok{2}\NormalTok{)   }\CommentTok{\# converge para p = 0.5 (para 0.7 no ex. 2)}
\end{Highlighting}
\end{Shaded}

\textbf{3. e 4.} Média acumulada (Poisson e uniforme) --- converge para \(\lambda\) e para
\(1/2\), respetivamente.

\begin{Shaded}
\begin{Highlighting}[]
\FunctionTok{set.seed}\NormalTok{(}\DecValTok{1}\NormalTok{)}
\NormalTok{x }\OtherTok{\textless{}{-}} \FunctionTok{rpois}\NormalTok{(}\DecValTok{2000}\NormalTok{, }\DecValTok{4}\NormalTok{)}
\FunctionTok{plot}\NormalTok{(}\FunctionTok{cumsum}\NormalTok{(x) }\SpecialCharTok{/} \FunctionTok{seq\_along}\NormalTok{(x), }\AttributeTok{type =} \StringTok{"l"}\NormalTok{); }\FunctionTok{abline}\NormalTok{(}\AttributeTok{h =} \DecValTok{4}\NormalTok{, }\AttributeTok{col =} \StringTok{"red"}\NormalTok{, }\AttributeTok{lty =} \DecValTok{2}\NormalTok{)}
\end{Highlighting}
\end{Shaded}

\textbf{5.} Cauchy: a média acumulada \textbf{não} converge, porque a Cauchy não tem média
finita --- as hipóteses da lei dos grandes números falham.

\begin{Shaded}
\begin{Highlighting}[]
\FunctionTok{set.seed}\NormalTok{(}\DecValTok{1}\NormalTok{)}
\NormalTok{x }\OtherTok{\textless{}{-}} \FunctionTok{rcauchy}\NormalTok{(}\DecValTok{5000}\NormalTok{)}
\FunctionTok{plot}\NormalTok{(}\FunctionTok{cumsum}\NormalTok{(x) }\SpecialCharTok{/} \FunctionTok{seq\_along}\NormalTok{(x), }\AttributeTok{type =} \StringTok{"l"}\NormalTok{,}
     \AttributeTok{main =} \StringTok{"Média acumulada de uma Cauchy (não converge)"}\NormalTok{)}
\end{Highlighting}
\end{Shaded}

\textbf{6.} Lei fraca vs.~majorante de Chebyshev, para \(X_i \sim \mathcal{U}(0,1)\).

\begin{Shaded}
\begin{Highlighting}[]
\FunctionTok{set.seed}\NormalTok{(}\DecValTok{1}\NormalTok{)}
\NormalTok{eps }\OtherTok{\textless{}{-}} \FloatTok{0.05}
\ControlFlowTok{for}\NormalTok{ (n }\ControlFlowTok{in} \FunctionTok{c}\NormalTok{(}\DecValTok{10}\NormalTok{, }\DecValTok{50}\NormalTok{, }\DecValTok{100}\NormalTok{, }\DecValTok{500}\NormalTok{)) \{}
\NormalTok{  p\_emp }\OtherTok{\textless{}{-}} \FunctionTok{mean}\NormalTok{(}\FunctionTok{replicate}\NormalTok{(}\DecValTok{5000}\NormalTok{, }\FunctionTok{abs}\NormalTok{(}\FunctionTok{mean}\NormalTok{(}\FunctionTok{runif}\NormalTok{(n)) }\SpecialCharTok{{-}} \FloatTok{0.5}\NormalTok{) }\SpecialCharTok{\textgreater{}}\NormalTok{ eps))}
\NormalTok{  cheby }\OtherTok{\textless{}{-}}\NormalTok{ (}\DecValTok{1}\SpecialCharTok{/}\DecValTok{12}\NormalTok{) }\SpecialCharTok{/}\NormalTok{ (n }\SpecialCharTok{*}\NormalTok{ eps}\SpecialCharTok{\^{}}\DecValTok{2}\NormalTok{)   }\CommentTok{\# σ²/(n ε²), com σ² = 1/12}
  \FunctionTok{cat}\NormalTok{(}\StringTok{"n ="}\NormalTok{, n, }\StringTok{"| P empírica:"}\NormalTok{, }\FunctionTok{round}\NormalTok{(p\_emp, }\DecValTok{4}\NormalTok{),}
      \StringTok{"| majorante Chebyshev:"}\NormalTok{, }\FunctionTok{round}\NormalTok{(cheby, }\DecValTok{4}\NormalTok{), }\StringTok{"}\SpecialCharTok{\textbackslash{}n}\StringTok{"}\NormalTok{)}
\NormalTok{\}}
\end{Highlighting}
\end{Shaded}

\textbf{7. e 8.} Teorema limite central a partir da uniforme e da log-normal (esta
converge mais devagar, por ser fortemente assimétrica).

\begin{Shaded}
\begin{Highlighting}[]
\FunctionTok{set.seed}\NormalTok{(}\DecValTok{1}\NormalTok{)}
\FunctionTok{par}\NormalTok{(}\AttributeTok{mfrow =} \FunctionTok{c}\NormalTok{(}\DecValTok{1}\NormalTok{, }\DecValTok{3}\NormalTok{))}
\ControlFlowTok{for}\NormalTok{ (n }\ControlFlowTok{in} \FunctionTok{c}\NormalTok{(}\DecValTok{2}\NormalTok{, }\DecValTok{5}\NormalTok{, }\DecValTok{12}\NormalTok{)) \{}
\NormalTok{  medias }\OtherTok{\textless{}{-}} \FunctionTok{rowMeans}\NormalTok{(}\FunctionTok{matrix}\NormalTok{(}\FunctionTok{runif}\NormalTok{(}\DecValTok{10000} \SpecialCharTok{*}\NormalTok{ n), }\AttributeTok{nrow =} \DecValTok{10000}\NormalTok{))}
  \FunctionTok{hist}\NormalTok{(medias, }\AttributeTok{probability =} \ConstantTok{TRUE}\NormalTok{, }\AttributeTok{main =} \FunctionTok{paste}\NormalTok{(}\StringTok{"n ="}\NormalTok{, n), }\AttributeTok{col =} \StringTok{"lightblue"}\NormalTok{)}
  \FunctionTok{curve}\NormalTok{(}\FunctionTok{dnorm}\NormalTok{(x, }\FloatTok{0.5}\NormalTok{, }\FunctionTok{sqrt}\NormalTok{((}\DecValTok{1}\SpecialCharTok{/}\DecValTok{12}\NormalTok{) }\SpecialCharTok{/}\NormalTok{ n)), }\AttributeTok{col =} \StringTok{"red"}\NormalTok{, }\AttributeTok{lwd =} \DecValTok{2}\NormalTok{, }\AttributeTok{add =} \ConstantTok{TRUE}\NormalTok{)}
\NormalTok{\}}
\FunctionTok{par}\NormalTok{(}\AttributeTok{mfrow =} \FunctionTok{c}\NormalTok{(}\DecValTok{1}\NormalTok{, }\DecValTok{1}\NormalTok{))}
\end{Highlighting}
\end{Shaded}

\textbf{10. e 11.} Cobertura de um intervalo de confiança. A proporção de intervalos
\(\bar{x} \pm 1{,}96\,s/\sqrt{n}\) que contêm o verdadeiro valor deve aproximar-se
de \(95\%\) à medida que \(n\) cresce.

\begin{Shaded}
\begin{Highlighting}[]
\FunctionTok{set.seed}\NormalTok{(}\DecValTok{1}\NormalTok{)}
\NormalTok{cobertura }\OtherTok{\textless{}{-}} \ControlFlowTok{function}\NormalTok{(n, }\AttributeTok{reps =} \DecValTok{5000}\NormalTok{) \{}
\NormalTok{  contido }\OtherTok{\textless{}{-}} \FunctionTok{replicate}\NormalTok{(reps, \{}
\NormalTok{    x }\OtherTok{\textless{}{-}} \FunctionTok{rexp}\NormalTok{(n, }\AttributeTok{rate =} \DecValTok{1}\NormalTok{)               }\CommentTok{\# média verdadeira = 1}
\NormalTok{    erro }\OtherTok{\textless{}{-}} \FloatTok{1.96} \SpecialCharTok{*} \FunctionTok{sd}\NormalTok{(x) }\SpecialCharTok{/} \FunctionTok{sqrt}\NormalTok{(n)}
    \FunctionTok{abs}\NormalTok{(}\FunctionTok{mean}\NormalTok{(x) }\SpecialCharTok{{-}} \DecValTok{1}\NormalTok{) }\SpecialCharTok{\textless{}=}\NormalTok{ erro}
\NormalTok{  \})}
  \FunctionTok{mean}\NormalTok{(contido)}
\NormalTok{\}}
\FunctionTok{sapply}\NormalTok{(}\FunctionTok{c}\NormalTok{(}\DecValTok{10}\NormalTok{, }\DecValTok{40}\NormalTok{, }\DecValTok{200}\NormalTok{), cobertura)        }\CommentTok{\# aproxima{-}se de 0.95 quando n cresce}
\end{Highlighting}
\end{Shaded}

\section{Cálculo de integrais por Monte Carlo}\label{cuxe1lculo-de-integrais-por-monte-carlo}

Para cada integral, a estimativa é a média das avaliações da função em pontos
aleatórios, com erro-padrão \(s/\sqrt{k}\). Indicamos o valor exato ou de
referência quando é conhecido.

\textbf{1.} \(\int_0^1 e^{e^x}\,dx \approx 6{,}317\) (sem forma fechada).

\begin{Shaded}
\begin{Highlighting}[]
\FunctionTok{set.seed}\NormalTok{(}\DecValTok{1}\NormalTok{); u }\OtherTok{\textless{}{-}} \FunctionTok{runif}\NormalTok{(}\FloatTok{1e6}\NormalTok{)}
\FunctionTok{mean}\NormalTok{(}\FunctionTok{exp}\NormalTok{(}\FunctionTok{exp}\NormalTok{(u)))          }\CommentTok{\# \textasciitilde{} 6.317}
\end{Highlighting}
\end{Shaded}

\textbf{2.} \(\int_0^1 (1-x^2)^{3/2}\,dx = \dfrac{3\pi}{16} \approx 0{,}589\).

\begin{Shaded}
\begin{Highlighting}[]
\FunctionTok{set.seed}\NormalTok{(}\DecValTok{1}\NormalTok{); u }\OtherTok{\textless{}{-}} \FunctionTok{runif}\NormalTok{(}\FloatTok{1e6}\NormalTok{)}
\FunctionTok{mean}\NormalTok{((}\DecValTok{1} \SpecialCharTok{{-}}\NormalTok{ u}\SpecialCharTok{\^{}}\DecValTok{2}\NormalTok{)}\SpecialCharTok{\^{}}\NormalTok{(}\DecValTok{3}\SpecialCharTok{/}\DecValTok{2}\NormalTok{))      }\CommentTok{\# \textasciitilde{} 0.589 = 3pi/16}
\end{Highlighting}
\end{Shaded}

\textbf{3.} \(\int_{-2}^{2} e^{x + x^2}\,dx \approx 93{,}16\). Usa-se a forma para
intervalo \([a,b]\): multiplicar a média por \((b-a) = 4\).

\begin{Shaded}
\begin{Highlighting}[]
\FunctionTok{set.seed}\NormalTok{(}\DecValTok{1}\NormalTok{); u }\OtherTok{\textless{}{-}} \FunctionTok{runif}\NormalTok{(}\FloatTok{1e6}\NormalTok{, }\SpecialCharTok{{-}}\DecValTok{2}\NormalTok{, }\DecValTok{2}\NormalTok{)}
\DecValTok{4} \SpecialCharTok{*} \FunctionTok{mean}\NormalTok{(}\FunctionTok{exp}\NormalTok{(u }\SpecialCharTok{+}\NormalTok{ u}\SpecialCharTok{\^{}}\DecValTok{2}\NormalTok{))     }\CommentTok{\# \textasciitilde{} 93.16}
\end{Highlighting}
\end{Shaded}

\textbf{4.} \(\int_0^{\pi/2} \cos x\,dx = 1\).

\begin{Shaded}
\begin{Highlighting}[]
\FunctionTok{set.seed}\NormalTok{(}\DecValTok{1}\NormalTok{); u }\OtherTok{\textless{}{-}} \FunctionTok{runif}\NormalTok{(}\FloatTok{1e6}\NormalTok{, }\DecValTok{0}\NormalTok{, pi}\SpecialCharTok{/}\DecValTok{2}\NormalTok{)}
\NormalTok{(pi}\SpecialCharTok{/}\DecValTok{2}\NormalTok{) }\SpecialCharTok{*} \FunctionTok{mean}\NormalTok{(}\FunctionTok{cos}\NormalTok{(u))      }\CommentTok{\# \textasciitilde{} 1}
\end{Highlighting}
\end{Shaded}

\textbf{5.} \(\int_0^{\infty} x(1+x^2)^{-2}\,dx = \dfrac12\). Substituindo
\(y = 1/(x+1)\), o integral torna-se \(\int_0^1 h(y)\,dy\) com
\(h(y) = g(1/y - 1)/y^2\).

\begin{Shaded}
\begin{Highlighting}[]
\FunctionTok{set.seed}\NormalTok{(}\DecValTok{1}\NormalTok{)}
\NormalTok{g }\OtherTok{\textless{}{-}} \ControlFlowTok{function}\NormalTok{(x) x }\SpecialCharTok{*}\NormalTok{ (}\DecValTok{1} \SpecialCharTok{+}\NormalTok{ x}\SpecialCharTok{\^{}}\DecValTok{2}\NormalTok{)}\SpecialCharTok{\^{}}\NormalTok{(}\SpecialCharTok{{-}}\DecValTok{2}\NormalTok{)}
\NormalTok{u }\OtherTok{\textless{}{-}} \FunctionTok{runif}\NormalTok{(}\FloatTok{1e6}\NormalTok{)}
\NormalTok{h }\OtherTok{\textless{}{-}} \FunctionTok{g}\NormalTok{(}\DecValTok{1}\SpecialCharTok{/}\NormalTok{u }\SpecialCharTok{{-}} \DecValTok{1}\NormalTok{) }\SpecialCharTok{/}\NormalTok{ u}\SpecialCharTok{\^{}}\DecValTok{2}
\FunctionTok{mean}\NormalTok{(h)                    }\CommentTok{\# \textasciitilde{} 0.5}
\end{Highlighting}
\end{Shaded}

\textbf{6.} \(\int_{-\infty}^{\infty} e^{-x^2}\,dx = \sqrt{\pi} \approx 1{,}772\). Uma
via simples: reconhecer o integrando como proporcional à densidade de uma
\(N(0, 1/2)\) e usar amostragem por importância.

\begin{Shaded}
\begin{Highlighting}[]
\FunctionTok{set.seed}\NormalTok{(}\DecValTok{1}\NormalTok{)}
\NormalTok{X }\OtherTok{\textless{}{-}} \FunctionTok{rnorm}\NormalTok{(}\FloatTok{1e6}\NormalTok{)                       }\CommentTok{\# proposta N(0,1)}
\FunctionTok{mean}\NormalTok{(}\FunctionTok{exp}\NormalTok{(}\SpecialCharTok{{-}}\NormalTok{X}\SpecialCharTok{\^{}}\DecValTok{2}\NormalTok{) }\SpecialCharTok{/} \FunctionTok{dnorm}\NormalTok{(X))            }\CommentTok{\# \textasciitilde{} sqrt(pi) = 1.7725}
\end{Highlighting}
\end{Shaded}

\textbf{7.} \(\int_0^1\!\int_0^1 e^{(x+y)^2}\,dy\,dx \approx 4{,}899\) (integral duplo).

\begin{Shaded}
\begin{Highlighting}[]
\FunctionTok{set.seed}\NormalTok{(}\DecValTok{1}\NormalTok{); k }\OtherTok{\textless{}{-}} \FloatTok{1e6}
\NormalTok{x }\OtherTok{\textless{}{-}} \FunctionTok{runif}\NormalTok{(k); y }\OtherTok{\textless{}{-}} \FunctionTok{runif}\NormalTok{(k)}
\FunctionTok{mean}\NormalTok{(}\FunctionTok{exp}\NormalTok{((x }\SpecialCharTok{+}\NormalTok{ y)}\SpecialCharTok{\^{}}\DecValTok{2}\NormalTok{))       }\CommentTok{\# \textasciitilde{} 4.899}
\end{Highlighting}
\end{Shaded}

\textbf{8.} \(\int_0^1\!\int_0^1 (x^2 + y^2)\,dx\,dy = \dfrac23\).

\begin{Shaded}
\begin{Highlighting}[]
\FunctionTok{set.seed}\NormalTok{(}\DecValTok{1}\NormalTok{); k }\OtherTok{\textless{}{-}} \FloatTok{1e6}
\FunctionTok{mean}\NormalTok{(}\FunctionTok{runif}\NormalTok{(k)}\SpecialCharTok{\^{}}\DecValTok{2} \SpecialCharTok{+} \FunctionTok{runif}\NormalTok{(k)}\SpecialCharTok{\^{}}\DecValTok{2}\NormalTok{)   }\CommentTok{\# \textasciitilde{} 0.667}
\end{Highlighting}
\end{Shaded}

\textbf{9. e 10.} Estimativa de \(\pi\) (acerta-ou-falha e agulha de Buffon).

\begin{Shaded}
\begin{Highlighting}[]
\FunctionTok{set.seed}\NormalTok{(}\DecValTok{1}\NormalTok{); k }\OtherTok{\textless{}{-}} \FloatTok{1e6}
\CommentTok{\# acerta{-}ou{-}falha: fração de pontos dentro do quarto de círculo}
\NormalTok{x }\OtherTok{\textless{}{-}} \FunctionTok{runif}\NormalTok{(k); y }\OtherTok{\textless{}{-}} \FunctionTok{runif}\NormalTok{(k)}
\DecValTok{4} \SpecialCharTok{*} \FunctionTok{mean}\NormalTok{(x}\SpecialCharTok{\^{}}\DecValTok{2} \SpecialCharTok{+}\NormalTok{ y}\SpecialCharTok{\^{}}\DecValTok{2} \SpecialCharTok{\textless{}=} \DecValTok{1}\NormalTok{)             }\CommentTok{\# \textasciitilde{} pi}

\CommentTok{\# agulha de Buffon: l = 1, d = 2  =\textgreater{}  P(cruzar) = 2l/(pi d)  =\textgreater{}  pi \textasciitilde{} 2l/(dp)}
\NormalTok{centro }\OtherTok{\textless{}{-}} \FunctionTok{runif}\NormalTok{(k, }\DecValTok{0}\NormalTok{, }\DecValTok{1}\NormalTok{)             }\CommentTok{\# distância do centro à linha mais próxima (d/2 = 1)}
\NormalTok{angulo }\OtherTok{\textless{}{-}} \FunctionTok{runif}\NormalTok{(k, }\DecValTok{0}\NormalTok{, pi}\SpecialCharTok{/}\DecValTok{2}\NormalTok{)}
\NormalTok{cruza  }\OtherTok{\textless{}{-}}\NormalTok{ centro }\SpecialCharTok{\textless{}=} \FloatTok{0.5} \SpecialCharTok{*} \FunctionTok{sin}\NormalTok{(angulo)}
\DecValTok{2} \SpecialCharTok{*} \DecValTok{1} \SpecialCharTok{/}\NormalTok{ (}\DecValTok{2} \SpecialCharTok{*} \FunctionTok{mean}\NormalTok{(cruza))            }\CommentTok{\# \textasciitilde{} pi}
\end{Highlighting}
\end{Shaded}

\textbf{11.} Taxa de convergência \(1/\sqrt{k}\). Para \(\int_0^1 e^x\,dx = e - 1\), o erro
absoluto decresce, em média, como \(k^{-1/2}\) (declive \(-1/2\) em escala log-log).

\begin{Shaded}
\begin{Highlighting}[]
\FunctionTok{set.seed}\NormalTok{(}\DecValTok{1}\NormalTok{)}
\NormalTok{ks }\OtherTok{\textless{}{-}} \DecValTok{10}\SpecialCharTok{\^{}}\NormalTok{(}\DecValTok{2}\SpecialCharTok{:}\DecValTok{6}\NormalTok{)}
\NormalTok{erros }\OtherTok{\textless{}{-}} \FunctionTok{sapply}\NormalTok{(ks, }\ControlFlowTok{function}\NormalTok{(k) }\FunctionTok{abs}\NormalTok{(}\FunctionTok{mean}\NormalTok{(}\FunctionTok{exp}\NormalTok{(}\FunctionTok{runif}\NormalTok{(k))) }\SpecialCharTok{{-}}\NormalTok{ (}\FunctionTok{exp}\NormalTok{(}\DecValTok{1}\NormalTok{) }\SpecialCharTok{{-}} \DecValTok{1}\NormalTok{)))}
\FunctionTok{plot}\NormalTok{(}\FunctionTok{log10}\NormalTok{(ks), }\FunctionTok{log10}\NormalTok{(erros), }\AttributeTok{type =} \StringTok{"b"}\NormalTok{,}
     \AttributeTok{xlab =} \StringTok{"log10(k)"}\NormalTok{, }\AttributeTok{ylab =} \StringTok{"log10(erro)"}\NormalTok{)   }\CommentTok{\# declive \textasciitilde{} {-}0.5}
\end{Highlighting}
\end{Shaded}

\textbf{12.} Redução de variância por variáveis antitéticas, em \(\int_0^1 e^x\,dx\). O
estimador antitético usa os pares \(\big(g(U), g(1-U)\big)\); como \(g(x) = e^x\) é
monótona, \(g(U)\) e \(g(1-U)\) são negativamente correlacionadas, o que reduz a
variância da média.

\begin{Shaded}
\begin{Highlighting}[]
\FunctionTok{set.seed}\NormalTok{(}\DecValTok{1}\NormalTok{)}
\NormalTok{simples    }\OtherTok{\textless{}{-}} \FunctionTok{replicate}\NormalTok{(}\DecValTok{1000}\NormalTok{, }\FunctionTok{mean}\NormalTok{(}\FunctionTok{exp}\NormalTok{(}\FunctionTok{runif}\NormalTok{(}\DecValTok{1000}\NormalTok{))))}
\NormalTok{antitetico }\OtherTok{\textless{}{-}} \FunctionTok{replicate}\NormalTok{(}\DecValTok{1000}\NormalTok{, \{ u }\OtherTok{\textless{}{-}} \FunctionTok{runif}\NormalTok{(}\DecValTok{500}\NormalTok{)}
                                \FunctionTok{mean}\NormalTok{(}\FunctionTok{c}\NormalTok{(}\FunctionTok{exp}\NormalTok{(u), }\FunctionTok{exp}\NormalTok{(}\DecValTok{1} \SpecialCharTok{{-}}\NormalTok{ u))) \})}
\FunctionTok{c}\NormalTok{(}\AttributeTok{var\_simples =} \FunctionTok{var}\NormalTok{(simples), }\AttributeTok{var\_antitetico =} \FunctionTok{var}\NormalTok{(antitetico))  }\CommentTok{\# menor no 2.º}
\end{Highlighting}
\end{Shaded}

\textbf{13.} Volume da bola unitária em dimensão \(n\), por acerta-ou-falha. A eficiência
degrada-se com \(n\): a fração de pontos do hipercubo \([-1,1]^n\) que caem dentro da
bola tende rapidamente para zero.

\begin{Shaded}
\begin{Highlighting}[]
\FunctionTok{set.seed}\NormalTok{(}\DecValTok{1}\NormalTok{)}
\NormalTok{volume\_bola }\OtherTok{\textless{}{-}} \ControlFlowTok{function}\NormalTok{(n, }\AttributeTok{k =} \FloatTok{1e5}\NormalTok{) \{}
\NormalTok{  X }\OtherTok{\textless{}{-}} \FunctionTok{matrix}\NormalTok{(}\FunctionTok{runif}\NormalTok{(k }\SpecialCharTok{*}\NormalTok{ n, }\SpecialCharTok{{-}}\DecValTok{1}\NormalTok{, }\DecValTok{1}\NormalTok{), }\AttributeTok{ncol =}\NormalTok{ n)}
\NormalTok{  dentro }\OtherTok{\textless{}{-}} \FunctionTok{rowSums}\NormalTok{(X}\SpecialCharTok{\^{}}\DecValTok{2}\NormalTok{) }\SpecialCharTok{\textless{}=} \DecValTok{1}
  \DecValTok{2}\SpecialCharTok{\^{}}\NormalTok{n }\SpecialCharTok{*} \FunctionTok{mean}\NormalTok{(dentro)                       }\CommentTok{\# volume do cubo (2\^{}n) × fração dentro}
\NormalTok{\}}
\FunctionTok{sapply}\NormalTok{(}\DecValTok{2}\SpecialCharTok{:}\DecValTok{5}\NormalTok{, volume\_bola)}
\CommentTok{\# comparar com o valor exato: pi\^{}(n/2) / gamma(n/2 + 1)}
\FunctionTok{sapply}\NormalTok{(}\DecValTok{2}\SpecialCharTok{:}\DecValTok{5}\NormalTok{, }\ControlFlowTok{function}\NormalTok{(n) pi}\SpecialCharTok{\^{}}\NormalTok{(n}\SpecialCharTok{/}\DecValTok{2}\NormalTok{) }\SpecialCharTok{/} \FunctionTok{gamma}\NormalTok{(n}\SpecialCharTok{/}\DecValTok{2} \SpecialCharTok{+} \DecValTok{1}\NormalTok{))}
\end{Highlighting}
\end{Shaded}

O último exercício ilustra tanto a força como o limite do método de Monte Carlo.
A força: o mesmo código de três linhas funciona em qualquer dimensão. O limite:
como a fração de acertos colapsa com \(n\), o método acerta-ou-falha torna-se
ineficiente em dimensão alta --- nessa altura, técnicas como a amostragem por
importância passam a ser indispensáveis. Fica, assim, o mote para um estudo mais
avançado da simulação estatística, para além deste primeiro livro.

\end{document}
